\documentclass[11pt,oneside,openany,a4paper]{book}
\usepackage[T1]{fontenc}
\usepackage[utf8]{inputenc}
\usepackage{lmodern,microtype}
\usepackage[a4paper,left=26mm,right=26mm,top=25mm,bottom=25mm,headheight=16pt,headsep=8mm]{geometry}
\usepackage{amsmath,amssymb,amsthm,mathtools,bm,mathrsfs}
\usepackage{booktabs,array,tabularx,longtable,ragged2e}
\usepackage{enumitem,needspace,fancyhdr,xcolor,titlesec}
\usepackage[most]{tcolorbox}
\usepackage{hyperref,xurl,bookmark}
\definecolor{inkblue}{RGB}{30,57,82}
\definecolor{softblue}{RGB}{240,245,248}
\definecolor{rulegray}{RGB}{205,213,219}
\hypersetup{colorlinks=true,linkcolor=inkblue,citecolor=inkblue,urlcolor=inkblue,
 pdftitle={Shared-Sketch HSS Approximation: Proof Development, Obstructions, and Current Route},
 pdfauthor={},pdfsubject={Uniform refitting, adaptive basis learning, obstructions, and source-coupled regression}}
\pagestyle{fancy}\fancyhf{}
\fancyhead[L]{\small Shared-sketch HSS approximation}
\fancyhead[R]{\small\thepage}
\renewcommand{\headrulewidth}{.3pt}
\renewcommand{\chaptermark}[1]{\markboth{\thechapter\quad #1}{}}
\fancypagestyle{plain}{\fancyhf{}\fancyfoot[C]{\thepage}\renewcommand{\headrulewidth}{0pt}}
\titleformat{\chapter}[display]{\normalfont\huge\bfseries\color{inkblue}\raggedright}
 {\large\normalfont\scshape Chapter \thechapter}{9pt}{\Huge}
\titlespacing*{\chapter}{0pt}{0pt}{20pt}
\titleformat{\section}{\normalfont\Large\bfseries\color{inkblue}}{\thesection}{.7em}{}
\setlength{\parindent}{0pt}
\setlength{\parskip}{4pt plus 1pt minus .4pt}
\setlength{\emergencystretch}{3em}
\setlist{topsep=4pt,itemsep=2pt,parsep=0pt}
\allowdisplaybreaks[2]
\clubpenalty=10000\widowpenalty=10000\displaywidowpenalty=10000
\setcounter{tocdepth}{1}\setcounter{secnumdepth}{2}
\numberwithin{equation}{chapter}
\newtheorem{theorem}{Theorem}[chapter]
\newtheorem{lemma}[theorem]{Lemma}
\newtheorem{proposition}[theorem]{Proposition}
\newtheorem{corollary}[theorem]{Corollary}
\theoremstyle{definition}
\newtheorem{definition}[theorem]{Definition}
\newtheorem{assumption}[theorem]{Assumption}
\theoremstyle{remark}
\newtheorem{remark}[theorem]{Remark}
\newtheorem{example}[theorem]{Example}
\newtheorem{auditnote}[theorem]{Qualification}
\newtheorem{checkpoint}[theorem]{Historical open checkpoint}
\newtcolorbox{statusbox}[1][]{enhanced,breakable,colback=softblue,colframe=rulegray,
 boxrule=.5pt,arc=.6mm,left=3mm,right=3mm,top=2mm,bottom=2mm,#1}
\newtcolorbox{scopebox}{enhanced,breakable,colback=softblue,colframe=rulegray,
 boxrule=.4pt,arc=.5mm,left=2.8mm,right=2.8mm,top=2mm,bottom=2mm,
 before skip=0pt,after skip=13pt,fontupper=\small}
\newenvironment{chaptersynopsis}{\begingroup\small\textbf{Synopsis.}\ }{\par\endgroup\medskip}
\newcommand{\proofstep}[1]{\par\addvspace{.45\baselineskip}\Needspace{3\baselineskip}\textbf{#1}\par\nobreak}
\newcolumntype{Y}{>{\RaggedRight\arraybackslash}X}
\providecommand{\tightlist}{\setlength{\itemsep}{0pt}\setlength{\parskip}{0pt}}
\newcommand{\R}{\mathbb R}\newcommand{\E}{\mathbb E}\newcommand{\Prb}{\mathbb P}\newcommand{\PP}{\mathbb P}
\newcommand{\T}{\mathsf T}\newcommand{\Id}{I}
\newcommand{\F}[1]{\left\lVert#1\right\rVert_F}
\newcommand{\N}[1]{\left\lVert#1\right\rVert}
\newcommand{\op}[1]{\left\lVert#1\right\rVert_{\rm op}}
\newcommand{\ip}[2]{\left\langle#1,#2\right\rangle_F}
\DeclareMathOperator{\tr}{tr}\DeclareMathOperator{\rank}{rank}
\DeclareMathOperator{\diag}{blockdiag}\DeclareMathOperator{\range}{range}
\DeclareMathOperator{\ran}{range}\DeclareMathOperator{\OPT}{OPT}
\DeclareMathOperator{\HSS}{HSS}\DeclareMathOperator{\dist}{dist}
\DeclareMathOperator{\Var}{Var}\DeclareMathOperator{\supp}{supp}
\DeclareMathOperator{\orth}{orth}\DeclareMathOperator*{\argmin}{arg\,min}
\newcommand{\HH}{\mathcal H}\newcommand{\SSS}{\mathcal S}
\newcommand{\UC}{\mathcal U}\newcommand{\VC}{\mathcal V}\newcommand{\err}{\mathcal E}
\newcommand{\Sn}[2]{\left\lVert#1\right\rVert_{S_{#2}}}
\newcommand{\Cov}{\mathsf C}\newcommand{\Loss}{\mathsf L}\newcommand{\Fil}{\mathcal F}
\DeclareMathOperator{\diver}{div}\DeclareMathOperator{\vecop}{vec}\DeclareMathOperator{\Gammafun}{Gamma}
\newcommand{\J}{\mathcal J}\newcommand{\transpose}{\mathsf T}
\newcommand{\sourcefile}[1]{{\small\nolinkurl{#1}}}
\pdfstringdefDisableCommands{\def\rho{rho}\def\Phi{Phi}\def\vartheta{theta}\def\mathcal#1{#1}\def\mathsf#1{#1}\def\texorpdfstring#1#2{#2}}
% A single public bibliography; private proof dependencies are internal references.
\makeatletter
\renewenvironment{thebibliography}[1]
 {\list{\@biblabel{\@arabic\c@enumiv}}%
 {\settowidth\labelwidth{\@biblabel{#1}}\leftmargin\labelwidth
 \advance\leftmargin\labelsep\@openbib@code\usecounter{enumiv}
 \let\p@enumiv\@empty\renewcommand\theenumiv{\@arabic\c@enumiv}}%
 \sloppy\clubpenalty4000\widowpenalty4000\sfcode`\.\@m}
 {\def\@noitemerr{\@latex@warning{Empty bibliography}}\endlist}
\makeatother

\begin{document}
\frontmatter
\hypersetup{pageanchor=false}
\begin{titlepage}
\vspace*{10mm}
{\color{inkblue}\large\scshape Technical manuscript\par}
\vspace{15mm}
{\color{inkblue}\Huge\bfseries Shared-Sketch\par HSS Approximation\par}
\vspace{7mm}
{\LARGE Uniform refitting, adaptive basis learning,\par and source-coupled regression\par}
\vspace{11mm}
{\large Proof development, algorithmic obstructions,\\ and current route\par}
\vspace{3mm}
{\normalsize Updated through 20 September 2026\par}
\vspace{12mm}
\begin{statusbox}
\textbf{Established endpoint.} Uniform independent fixed-space refitting, synchronized exact capture, and the signed unlifted source budget are proved. Restricted fixed-core and independent-source stability theorems retain their stated assumptions.

\textbf{Remaining endpoint.} The current synchronized learner still needs a general joint physical-error estimate. A polynomial-time reduction to final HSS rank $k$ is also required. The conjecture remains unresolved.
\end{statusbox}
\vfill
{\itshape Real matrices. Exact arithmetic. Squared Frobenius error.\par
A fixed binary hierarchy. Explicit conditioning and rank conventions.}
\end{titlepage}
\setcounter{page}{1}\hypersetup{pageanchor=true}

% Replacement front matter for hss_complete_proof_chain.tex.
% Include from the Abstract heading through the material before the ToC.
% Original bibliography keys; imported chapter labels have prefix upd:.

\chapter*{Abstract}
\addcontentsline{toc}{chapter}{Abstract}

This book records a proof chain and its obstructions for approximating a
fixed matrix from shared Gaussian matrix--vector products on a fixed
hierarchical semi-separable (HSS) tree.  The intended theorem returns an
HSS-rank-at-most-$k$ matrix with expected squared Frobenius error
$O(L)\OPT_k(A)$, using $O(k)$ total products with the matrix and its
transpose and polynomial additional work.  Here $L$ is the tree depth
and $\OPT_k(A)$ is the best squared HSS-rank-$k$ error on the original
target tree.  This theorem is not established in the book.

The completed linear component is independent fixed-space refitting.
For prescribed nested bases of rank at most $s$, with physical leaves
of size at most $2s$, the proved coefficient decoder uses $8s$ fresh
products and has expected squared-error factor $1607$ relative to the
best matrix in that fixed coefficient space, uniformly in depth.  The
original sharper-oversampling variants remain valid.  Independence of
the refit streams from the bases is essential; the theorem does not
learn an accurate space.

The current learning candidate is a completed synchronized two-view
union.  It uses per-view width $r=2k$, retained width $R=4k$, Gaussian
query width $q=32k$ per stream, and physical leaves of size $8k$.
Both views receive the same completed child unions.  A raw fit of rank
at most $R$ is retained in full; otherwise the lifted leading-$2k$
raw and annihilated ranges are united before one canonical completion,
with query-visible directions preceding kernel directions.  Exact
capture is proved for fixed inputs whose outgoing and incoming cut
ranks are at most $2k$.  For general inputs, a fixed-comparator right
separator yields an unlifted signed-source budget
$24qL\lVert A-B\rVert_F^2$ and an exact projected physical identity
containing source, inherited return, and invisible-design terms.  The
joint adaptive physical estimate needed for a general basis-quality
theorem remains open.

Restricted nonlinear results are retained with their precise models.
Inverse-optimal completion around fixed nested physical cores gives
physical-plus-retained source gains below $6.733$, or below $3.020$
when every ordinary core has half the retained rank.  Later theorems
give a depth-independent moment bound for independent fixed sources
at $r\geq4096$, and extend it to compatible target transfers using
fixed guard pools.  Pure fresh resets and noncentral Gaussian projector
identities supply further uniform or one-step estimates.  These are
not source-decomposition theorems for arbitrary HSS inputs.

The counterexamples concern distinct algorithms and logical claims.
An older full-range, inverse-optimally completed learner has an actual
expected approximation failure on fixed noisy inputs.  A separately
trained two-view union has a reachable clean-parent witness with
squared signal loss $1/3$, but its positive-probability neighborhood
depends on the noise scale and gives no uniform expected-error lower
bound by itself.  Neither result is a counterexample to the current
synchronized rule.  Its saved experiment contains 624 cases from 80
independent probe pairs, with paired noise levels and no nodewise
resampling, final refit, or rank-$k$ reduction.  These are finite
diagnostics.  The remaining end-to-end requirements are the general
physical coefficient-space bound and an efficient same-rank output
reduction; unrestricted best-HSS projection is only a query-only
reduction.

\chapter*{Main results and reading map}
\addcontentsline{toc}{chapter}{Main results and reading map}

\section*{The current target and conditional query ledger}

For the fixed target tree, with original physical leaves of size at
most $2k$, define
\[
 \OPT_k(A)=\inf_{B\in\HSS_k}\lVert A-B\rVert_F^2.
\]
The benchmark in Amsel et al.\ \cite{AmselHSS} has $O(kL)$ products
and an $O(L)$ expected squared-error factor.  The aim is to remove the
depth factor from the query count while keeping rank $k$, this error
factor, and polynomial additional work.  The reused-sketch method of
Levitt--Martinsson\ \cite{LM} provides the relevant linear-query
algorithmic precedent, but its stated analysis does not provide that
arbitrary-input guarantee for adaptive coarser-level spaces.  These
are statements about the cited results, not a claim that a bounded
literature check establishes the absence of every later theorem.

Let $\mathcal S(A)$ be the linear coefficient space of the learned
nested bases.  The current missing basis estimate is
\[
 \E\operatorname{dist}_F(A,\mathcal S(A))^2
 \leq C_bL\OPT_k(A).
\]
For synchronized retained width $R=4k$, training takes $64k$ products.
An independent fixed-space refit then takes $8R=32k$, for a
\emph{conditional} total of $96k$.  Its error factor $1607$ is proved;
the displayed basis estimate is not.  If it were proved, exact offline
projection of the refitted matrix $C$ to a best HSS-rank-$k$ matrix
$\widehat B$ would give
\[
 \lVert A-\widehat B\rVert_F
 \leq2\lVert A-C\rVert_F+\sqrt{\OPT_k(A)}.
\]
This last inequality settles only query accounting.  No polynomial-time
same-rank reduction with the needed constant-factor error has been
established here.  An extra depth factor in recompression would change
the intended approximation bound.  The learner's coarser $8k$ leaves
also do not change the original tree used to define the final target.

\section*{Where the current results are stated}

\begin{description}[style=nextline,leftmargin=0pt]
\item[Independent refitting: proved.]
Chapters~\ref{ch:drift}--\ref{ch:uniform} contain the original
depth-uniform proof, including retained-coefficient cancellation and
source assembly.  Section~\ref{upd:sec:route-refit} states its present
interface using rank $s$, so the substitution $s=4k$ in the query
ledger is explicit.  The old factors $15/4$ at $36s$ queries and
$1+\varepsilon$ at $O(s/\varepsilon)$ queries remain fixed-space
results, not basis-learning guarantees.

\item[Synchronized exact capture: proved on the stated rank-slack family.]
Chapter~\ref{upd:chap:route-synchronized} specifies all full-range,
truncated, zero-rank, and completion branches.  It proves almost-sure
full stacked design rank and true-cut capture for every fixed input
with both exterior-cut ranks at most $2k$.  A clean parent also has
exact comparator capture when its common child template contains that
cut and its compressed design is injective.  These conclusions do not
bound the inverse of a general noisy adaptive design.

\item[General-input separator and signed identity: proved; physical closure open.]
Chapter~\ref{upd:chap:route-signed} constructs the separator from the
fixed comparator, proves the $24qL$ unlifted budget, and preserves the
source/return cross term and invisible component in the physical
identity.  Multiplying by the adaptive left inverse requires a joint
estimate.  Separate expected bounds cannot be factored, and a query
budget alone is not a coefficient-space approximation theorem.

\item[Restricted source stability: proved under distinct assumptions.]
Chapter~\ref{ch:core} retains the $6.733$ and $3.020$ fixed-core
theorems.  Chapter~\ref{upd:chap:route-sources} also records independent
physical envelopes, the one-residual-source gain below $13.384$, the
once-only actual/averaged second-moment factor $11520$, and the
$r\geq4096$ finite moment ladder with fixed-guard compatible transfers.
Source independence, local equivariance, and the permitted conditioning
records are substantive hypotheses.  They are not asserted for
arbitrary HSS residual histories.

\item[Fresh bridges and optional residual correction: qualified results.]
Chapter~\ref{upd:chap:route-fresh} states the reserved-support freshness
condition, uniform pure-reset bounds, and the noncentral one-step
bridge with its explicit inherited metric cost.  Its optional
$4k+2$-product correction depends on a small \emph{global rank-$k$
tail} of the pilot residual; HSS structure alone does not supply that
condition or an accurate pilot.

\item[Actual numerical checks: finite evidence.]
Chapter~\ref{upd:chap:route-evidence} records 192 exact controls,
384 noisy cases, and 48 clean-parent cases: 624 cases from 80
independent probe pairs.  Means and standard errors use independent
whole trees within a fixed cell; paired cases are not pooled as
independent samples.  Noisy ratios use specified comparator residuals,
not evaluated OPT.  A 400-node signed-identity check reuses three of
those pairs.  No all-depth physical bound or final rank-$k$ theorem is
inferred from these observations.
\end{description}

\section*{What the obstructions establish}

The annihilate-first failure in Chapter~\ref{ch:anchor} is an
expected basis-risk counterexample for that specified learner.  The
later old full-range rule, with retained rank $8k$, wide query width
$64k$, response width $4k$, and leaves of size $16k$, also has an
actual full-projector expected-error counterexample; see
Section~\ref{upd:sec:later-fullrange}.  Its growing predesign inverse
on an exactly recoverable input is a separate, weaker obstruction and
must not be confused with that approximation failure.

The maintained \emph{separate-view} union runs its annihilated and raw
hierarchies separately, with retained width $2k$ in each, then unites
their nested spaces.  Section~\ref{upd:sec:later-separateviews} proves
a reachable loss of $1/3$ at a clean parent although the child union
contains the signal.  For each fixed positive noise level the witness
has a positive-probability probe neighborhood; that probability is
noise-dependent.  This disproves a deterministic clean-parent
inheritance claim, not the union's all-input expected bound by itself.
The synchronized rule feeds unions into both recursions and has a
full-raw branch, so it is a different algorithm.  The historical
$64k$ conditional ledger for the separate-view union must likewise
not be substituted for the current synchronized $96k$ ledger.

\section*{Reading the historical proof chain}

The original parts are preserved to show which ideas survived and how
later statements changed their scope.
\begin{enumerate}
 \item Chapters~\ref{ch:foundations}--\ref{ch:predictable} develop the
 model, raw reuse, orthogonal probing, source registers, storage, and
 preliminary moment certificates.  Their intermediate checkpoints are
 not the final status of the refit problem.
 \item Chapters~\ref{ch:centered}--\ref{ch:uniform} close the uniform
 fixed-space refit through cancellation and residual assembly.
 \item Chapters~\ref{ch:bicriteria}--\ref{ch:slack} study selection
 reductions, leaf safeguards, adaptive erasure, and deterministic rank
 slack.  Offline projection statements are query-only unless a
 computational guarantee is supplied explicitly.
 \item Chapters~\ref{ch:comm}--\ref{ch:anchor} introduce original-data
 commutators, actual-rule inverse obstructions, local physical control,
 and the annihilate-first failure and separate-view repair.
 \item Chapters~\ref{ch:tangent}--\ref{ch:core} establish the tangent
 and local nonlinear regression calculus, inverse-optimal completion,
 and fixed-core stability.  Fixed-height limits, local-model bounds,
 and arbitrary-input approximation remain different statements.
 \item Chapter~\ref{ch:synthesis} and Appendices~\ref{ch:routeledger}--\ref{ch:external} retain the
 earlier synthesis, route ledger, and external ingredients.  The
 imported current-route chapters and later-obstruction chapter provide
 the updated status where subsequent results supersede a historical
 checkpoint.
\end{enumerate}

The record is frozen at the instruction to stop research on
20 September 2026.  Internal ``audits'' are recorded mathematical
reviews, not formal verification or external peer review.  The
stop-point discard-accounting and retained-rank observations are
segregated in Section~\ref{upd:sec:route-unreviewed}; they are not
silently promoted to proved theorems.  The current unresolved endpoints
are the joint adaptive physical budget and a computationally justified
same-rank output reduction.

\clearpage

\section*{Revision scope and source conventions}
Parts~I--VI preserve the original development, with dated qualifications at
the affected chapters. Part~VII integrates the subsequent obstruction
catalogue and the current synchronized proof route. It takes precedence
over earlier descriptions of the current learner or remaining endpoint.
The original input manuscript is preserved as
\sourcefile{backups/hss_complete_proof_chain.pre_2026-09-20_update.tex}.
When an added chapter cites \sourcefile{hss_complete_proof_chain.tex} as a
source, it refers to that preserved original development, reproduced here
in Parts~I--VI and the original appendices.

The historical chapters retain their detailed internal arguments.
The added consolidation gives explicit statements, constructions, and
proofs or proof outlines, and identifies the supporting proof and audit
notes where further details are retained. Those private records are listed
in Appendix~\ref{upd:appendix:sources}; the public bibliography remains
separate. Internal review is not external peer review or formal verification.
This update consolidates the frozen research record and does not resume
proof search or add experiments.
\clearpage
\tableofcontents
\clearpage

\mainmatter

\part{Model, reuse, and the preliminary refit chain}

% ===== FOUNDATIONS =====
\begingroup
\let\R\relax
\newcommand{\R}{\mathbb R}
\let\E\relax
\newcommand{\E}{\mathbb E}
\let\Prb\relax
\newcommand{\Prb}{\mathbb P}
\let\T\relax
\newcommand{\T}{\mathsf T}
\let\F\relax
\newcommand{\F}[1]{\left\lVert#1\right\rVert_{F}}
\let\N\relax
\newcommand{\N}[1]{\left\lVert#1\right\rVert}
\let\ip\relax
\newcommand{\ip}[2]{\left\langle#1,#2\right\rangle_F}
\let\tr\relax
\newcommand{\tr}{\operatorname{tr}}
\let\rank\relax
\newcommand{\rank}{\operatorname{rank}}
\let\diag\relax
\newcommand{\diag}{\operatorname{blockdiag}}
\let\range\relax
\newcommand{\range}{\operatorname{range}}
\let\OPT\relax
\newcommand{\OPT}{\operatorname{OPT}}
\let\HSS\relax
\newcommand{\HSS}{\operatorname{HSS}}
\let\diver\relax
\newcommand{\diver}{\operatorname{div}}
\let\Var\relax
\newcommand{\Var}{\operatorname{Var}}
\let\UC\relax
\newcommand{\UC}{\mathcal U}
\let\VC\relax
\newcommand{\VC}{\mathcal V}
\let\err\relax
\newcommand{\err}{\mathcal E}
\let\proofstep\relax
\newcommand{\proofstep}[1]{\par\addvspace{.45\baselineskip}\Needspace{3\baselineskip}\textbf{#1}\par\nobreak}
\let\op\relax
\newcommand{\op}[1]{\left\lVert#1\right\rVert_{\mathrm{op}}}
\let\Sn\relax
\newcommand{\Sn}[2]{\left\lVert#1\right\rVert_{S_{#2}}}
\let\Cov\relax
\newcommand{\Cov}{\mathsf C}
\let\Loss\relax
\newcommand{\Loss}{\mathsf L}
\let\Fil\relax
\newcommand{\Fil}{\mathcal F}
\let\Id\relax
\newcommand{\Id}{I}
\let\supp\relax
\newcommand{\supp}{\operatorname{supp}}
\let\vecop\relax
\newcommand{\vecop}{\operatorname{vec}}
\let\Gammafun\relax
\newcommand{\Gammafun}{\operatorname{Gamma}}
\chapter{Model and raw shared-sketch learning}\label{ch:foundations}
\begin{scopebox}
\textbf{Update, 20 September 2026.} The raw shared-sketch learner in this chapter is a historical algorithm, distinct from the completed synchronized learner of Chapter~\ref{upd:chap:current-route}. Its exact-input and first-level results retain their stated scope. Later local models and synchronized identities do not retrospectively supply the missing all-input induction for this original rule.
\end{scopebox}

\begin{scopebox}
The raw learner analyzed in this chapter is distinct from the later commutator and union learners. Its exact-input and local approximation results do not establish the all-input target. The fixed-space refit obstruction encountered in the preliminary chain is resolved in Part~II; the basis-learning premise remains open.
\end{scopebox}
\section{Problem, conventions, and attribution}\label{base:sec:model}
\subsection{The approximation target}
Fix integers $k\ge1$, $L\ge1$, and
\begin{equation}\label{base:eq:dimensions}
 N=2^{L+1}k,\qquad m=2k,\qquad s=8k.
\end{equation}
The binary tree is fixed in advance. Its leaves are consecutive sets of
$m$ indices. Each parent is the union of two adjacent children. At
compression level $\ell=0,\ldots,L-1$, the original node sets have size
$2^{\ell+1}k$, while their current compressed coordinates have size $2k$.
The root is not compressed; its final core has size $2k\times2k$.

\begin{definition}[HSS class]
For every nonroot node $I$ of this tree, require
\begin{equation}\label{base:eq:cut-ranks}
 \rank B[I,I^c]\le k,\qquad \rank B[I^c,I]\le k.
\end{equation}
The matrices satisfying all these constraints form $\HSS(L,k)$.
Set
\begin{equation}\label{base:eq:opt}
 \OPT(A)=\min_{B\in\HSS(L,k)}\F{A-B}^2.
\end{equation}
\end{definition}
These cut constraints are the usual nested off-diagonal rank condition.
Equivalently, matrices in this class have a telescoping representation with
$2k\times k$ orthonormal node bases, arbitrary $2k\times2k$ block-diagonal
discrepancies, and a $2k\times2k$ root core. Here is a direct reason for the
nestedness needed in this equivalence. At a parent node, the restriction of
its exterior interaction to either child lies in that child's interaction
space. Thus a parent basis can be expressed in the direct sum of the two
child bases. Bases of dimension below $k$ may be padded orthonormally.
Conversely, tracing a telescoping representation across a node cut gives a
factorization through that node's $k$-column basis. Diagonal discrepancies
inside the node do not cross its cut.

The minimum in~\eqref{base:eq:opt} exists. Each rank constraint is closed, so the
class is closed; it contains zero. A minimizing sequence can be restricted
to $\F{A-B}\le\F A+1$, hence to a bounded subset of a finite-dimensional
space. A convergent subsequence and continuity attain the minimum. Fix one
minimizer $B_*$ when a comparator is needed; the algorithm never uses it.

\textbf{All approximation factors in this chapter multiply squared
Frobenius error.} The desired statement is
\begin{equation}\label{base:eq:target}
 \E\F{A-\widehat B}^2\le C L\,\OPT(A),
 \quad \widehat B\in\HSS(L,k),
 \quad \#\text{queries}\le C_qk,
\end{equation}
with absolute constants independent of $N,k,L,A$.

\subsection{Access model and notation}
A forward query returns $Ax$ for one chosen vector $x\in\R^N$; a transpose
query returns $A^{\T} x$. Multiplication by a matrix with $s$ columns counts
as $s$ queries, not as one query. Subsequent queries may depend on all
previous answers. No individual matrix entries are available to the
reconstruction algorithm.

The transpose is $M^{\T}$ and the Moore--Penrose inverse is $M^\dagger$.
The norms $\F M$ and $\N M$ are Frobenius and spectral norms. All inner
products between matrices are Frobenius inner products. A standard Gaussian
matrix has independent $N(0,1)$ entries. The notation $[M]_k$ denotes a best
rank-at-most-$k$ approximation in Frobenius norm.

A basis-selection rule chooses a rank-$k$ leading eigenspace of $EE^{\T}$,
with a fixed measurable rule for ties and for choosing a frame inside that
eigenspace. Basing this rule on the Gram, rather than on a particular
factorization of it, makes nullified samples and regression innovations
consistent. No spectral gap is assumed for the pathwise results. For the
exact-input theorem and its applications below, specialize this rule to
the canonical completion in Section~\ref{base:sec:exact-recovery}.

\subsection{What is inherited and what is analyzed here}
Amsel, Chen, Duman Keles, Halikias, C.~Musco, Ch.~Musco, and Persson
\cite{AmselHSS} prove an $O(L)$ expected squared-error factor using
$O(kL)$ queries with fresh sketches across levels. Their stated research
question is whether depth-independent sketch reuse can retain that factor.
Levitt and Martinsson~\cite{LM} supply the black-box acquisition,
block-nullification, discrepancy, and telescoping reuse construction.
The algorithm below uses leading singular subspaces and replaces the final
reused regression with exact adaptive core queries.

The organizing method separates an exact state update from its quantitative cost.
Section~\ref{base:new:fp} proves the finite-profile comparison motivating that method.
Linear descent through an encoding has the classical Douglas factorization
antecedent~\cite{Douglas}; nonlinear singular-vector selection instead requires
a specific invariance proof. Gaussian products, inverse moments and integration
by parts are classical. No claim of publication priority or formal proof verification
is made for the application-specific development below.

\subsection{Comparison with arXiv:2505.16937}\label{base:subsec:paper-comparison}
Table~\ref{base:tab:paper-comparison} uses version~2 (6 September 2025) of
Amsel et al.~\cite{AmselHSS}, with the same $N=2^{L+1}k$ and squared
Frobenius error convention. Write $F=\F{A-\widehat B}^2$. For their sketch
width $s_{\rm p}\ge3k+2$, Theorems~3.1 and~4.1 give
\begin{align}
 g(s_{\rm p})&=\left(1+
 \frac{2e(s_{\rm p}-2k)}{\sqrt{(s_{\rm p}-3k)^2-1}}\right)^2,
 & d(s_{\rm p})&=\frac{2k}{s_{\rm p}-2k-1},\nonumber\\
 \beta(s_{\rm p})&=2g(s_{\rm p})(1+d(s_{\rm p})).
 \label{base:eq:paper-beta}
\end{align}
Thus $s_{\rm p}=5k$ gives an absolute constant $\beta(s_{\rm p})$.
Their additional runtime is $O(Nks_{\rm p}L+Ns_{\rm p}^2)$.

The fixed-space error and the best-HSS error are different comparators:
$d_{\mathcal S}(A)^2\ge\OPT(A)$. Neither exact recovery on exact HSS inputs nor
an upper bound relative to $d_{\mathcal S}(A)$ implies the desired best-HSS
approximation statement. The complete query ledger is in
Section~\ref{base:new:ledger}.

The published greedy lower bound approaching a factor two is not an
$\Omega(L)$ lower bound and not an impossibility theorem for $O(k)$ queries.
The published fresh-level construction uses four width-$s_{\rm p}$ sketches
per level plus $2k$ core columns; both types are counted here.

\section{The candidate algorithm and the query ledger}\label{base:sec:algorithm}
\subsection{Acquisition and local compression}
Draw independent standard Gaussian $\Omega_0,\Psi_0\in\R^{N\times s}$ and
acquire
\begin{equation}\label{base:eq:acquisition}
 Y_0=A\Omega_0,\qquad Z_0=A^{\T}\Psi_0.
\end{equation}
At level $\ell$, all four current matrices have $n_\ell=N/2^\ell$ rows.
Partition them into $m$-row blocks. For each block $I$, compute
\begin{align}
 B_{Y,I}&=Y_{\ell,I}\Omega_{\ell,I}^{\dagger},&
 B_{Z,I}&=(\Psi_{\ell,I}^{\dagger})^{\T} Z_{\ell,I}^{\T},
                                                        \label{base:eq:local-regressions}\\
 E_{Y,I}&=Y_{\ell,I}-B_{Y,I}\Omega_{\ell,I},&
 E_{Z,I}&=Z_{\ell,I}-B_{Z,I}^{\T}\Psi_{\ell,I}.\label{base:eq:local-innovations}
\end{align}
Choose $U_{\ell,I}$ and $V_{\ell,I}$ from $E_{Y,I}$ and $E_{Z,I}$ by the
basis-selection rule. Define $P_{\ell,I}=U_{\ell,I}U_{\ell,I}^{\T}$ and
$Q_{\ell,I}=V_{\ell,I}V_{\ell,I}^{\T}$, and set
\begin{equation}\label{base:eq:local-D}
 D_{\ell,I}=(I-P_{\ell,I})B_{Y,I}
              +P_{\ell,I}B_{Z,I}(I-Q_{\ell,I}).
\end{equation}
Assemble $U_\ell,V_\ell,D_\ell$ block diagonally, and recycle
\begin{align}
 \Omega_{\ell+1}&=V_\ell^{\T}\Omega_\ell,&
 \Psi_{\ell+1}&=U_\ell^{\T}\Psi_\ell,\label{base:eq:test-recycle}\\
 Y_{\ell+1}&=U_\ell^{\T}(Y_\ell-D_\ell\Omega_\ell),&
 Z_{\ell+1}&=V_\ell^{\T}(Z_\ell-D_\ell^{\T}\Psi_\ell).
                                                   \label{base:eq:data-recycle}
\end{align}
Pair adjacent $k$-row child outputs to form the next $2k$-row blocks.
There are no new matrix queries in these operations.

\begin{remark}[Rank-deficient designs]
The use of the full projector $I-\Omega^\dagger\Omega$ makes the algorithm
well-defined even if a local design loses row rank. On the full-row-rank
locus it is equivalent to using all $s-m$ nullspace columns. At deficient
rank, it is a specified extension using the entire nullspace, not an
assertion that the subsequent quotient formulas remain valid. The gauge,
HSS membership, core recovery, and global error identities below hold for
this extension. The quotient and regression theorems explicitly require
full row rank where they use it. Full row rank and inverse integrability
through arbitrary depth are not inferred from Gaussian initialization.
\end{remark}

\subsection{Gauge, true compressed matrices, and exact root recovery}
Define cumulative isometries
\begin{equation}\label{base:eq:cumulative}
 \UC_0=\VC_0=I_N,\quad
 \UC_{\ell+1}=\UC_\ell U_\ell,\quad
 \VC_{\ell+1}=\VC_\ell V_\ell.
\end{equation}
Equation~\eqref{base:eq:local-D} gives the pathwise gauge
\begin{equation}\label{base:eq:gauge}
 U_\ell^{\T} D_\ell V_\ell=0.
\end{equation}
For analysis, define
\begin{equation}\label{base:eq:true-compress}
 A_0=A,\qquad
 A_{\ell+1}=U_\ell^{\T}(A_\ell-D_\ell)V_\ell
            =U_\ell^{\T} A_\ell V_\ell.
\end{equation}
Thus
\begin{equation}\label{base:eq:actual-core}
 A_\ell=\UC_\ell^{\T} A\VC_\ell.
\end{equation}
The algorithm has noisy recycled measurements of $A_\ell$, not its entries.
At the final level, however, $\VC_L$ has only $2k$ columns. Query those
columns of $A$ and form
\begin{equation}\label{base:eq:root}
 A_L=\UC_L^{\T}(A\VC_L).
\end{equation}
The output is
\begin{equation}\label{base:eq:output}
 \widehat B=
 \sum_{\ell=0}^{L-1}\UC_\ell D_\ell\VC_\ell^{\T}
 +\UC_L A_L\VC_L^{\T}.
\end{equation}
This is the telescoping HSS representation. Its cut ranks are at most $k$
by the nested factorization argument in Section~\ref{base:sec:model}.

\begin{proposition}[Acquisition cost]\label{base:prop:queries}
The algorithm uses $10k$ forward queries and $8k$ transpose queries,
independently of $L$. It returns a matrix in $\HSS(L,k)$. These assertions
do not require the target error bound or a full-rank hypothesis.
\end{proposition}
\begin{proof}
Equation~\eqref{base:eq:acquisition} uses $s=8k$ columns in each direction.
Only~\eqref{base:eq:root} adds queries, namely $2k$ forward columns. All remaining
operations use the answers and the computed factors. The output's
representation proves HSS membership.\end{proof}
The computation is polynomial in $N,k,L$. A blockwise implementation should
store the factors rather than form the dense matrix~\eqref{base:eq:output}.
A dense implementation of the displayed algorithm is a verification implementation, not a claim
of optimal arithmetic complexity.

\section{Exact recovery at every finite depth}\label{base:sec:exact-recovery}
\subsection{Make the completion rule explicit}
For a subspace $S\subseteq\R^{2k}$ of dimension at most $k$, define a fixed
canonical $k$-dimensional extension as follows. Obtain a canonical frame
of $S$ by applying Gram--Schmidt, in coordinate order, to the projections
of the standard coordinate vectors onto $S$, skipping zero vectors.
Append coordinate vectors after projecting out the current span until its
dimension is $k$. Finally select the frame of the resulting projector by
the same coordinate rule.

When an innovation has rank at most $k$, use this extension of its column
space. When its rank is greater than $k$, use any fixed measurable leading
rank-$k$ projector, with a fixed rule at a tied cutoff, and choose the
frame canonically from that projector. This is a valid specialization of
the Gram-based basis-selection rule of Section~\ref{base:sec:model}. Crucially, in the rank-at-most-
$k$ case the output depends only on the column space, not on its singular
values or a random rotation of its positive singular vectors.

\begin{theorem}[Exact recovery at arbitrary finite depth]\label{base:aug:thm:exact}
Use the canonical completion rule above. For each fixed
$A\in\HSS(L,k)$, the raw $18k$-query reconstruction equals $A$ almost surely.
All the local designs reached in this exact-input case have full row rank
almost surely. No full-rank-cut or spectral-gap assumption is needed.
\end{theorem}
\begin{proof}
First construct a deterministic comparison trajectory using explicit
access to $A$ only for this proof. At each level choose $U_I$ from the
column space of the true exterior row $A_\ell[I,I^c]$ and choose $V_I$
from the column space of $A_\ell[I^c,I]^{\T}$, using the canonical extension.
These spaces have dimension at most $k$. Set
\[
 D_I=A_\ell[I,I]-P_IA_\ell[I,I]Q_I,
 \qquad A_{\ell+1}=U_\ell^{\T}A_\ell V_\ell.
\]
The HSS cut constraints persist under these aligned compressions.
Every exterior interaction lies in its retained row and column spaces, so
\begin{equation}\label{base:aug:eq:det-exact-step}
 A_\ell=D_\ell+U_\ell A_{\ell+1}V_\ell^{\T}.
\end{equation}
Thus the entire comparison trajectory, including its frames, is a
deterministic function of the fixed matrix $A$.

Now evaluate the original Gaussian sketches on these deterministic
cumulative isometries. At any one level, the compressed designs are standard
Gaussian, with independent row blocks, since multiplication on the left by
a fixed matrix with orthonormal rows preserves that law. This assertion
does not claim independence across levels.

For a block $I$, condition on its $2k\times8k$ local design $\Omega_I$.
It has full row rank almost surely. An orthonormal matrix $N_I$ spanning
its nullspace has $6k$ columns. The exterior matrix
$\Omega_{I^c}N_I$ is, conditionally, a standard Gaussian matrix. Therefore
\[
 \range\big(A_\ell[I,I^c]\Omega_{I^c}N_I\big)
   =\range A_\ell[I,I^c]
\]
almost surely: a fixed matrix of rank at most $k$ retains its column space
after multiplication by $6k$ Gaussian vectors. The analogous transpose
statement holds as well. There are only finitely many blocks and levels,
so all these statements hold simultaneously on an event of probability one.

On this event, induction identifies the actual and comparison trajectories.
At the first level the nullified innovations have the comparison column
spaces, hence the same canonical bases. In each local regression the
exterior contamination is killed by $I-P_I$ on the forward side and by
$I-Q_I$ on the transpose side. Consequently the computed discrepancy is
exactly the $D_I$ in~\eqref{base:aug:eq:det-exact-step}. The recycled data are then
exact measurements of $A_{\ell+1}$ with the comparison designs. This proves
the induction at every level. The final exact core queries complete the
reconstruction.
\end{proof}

\begin{remark}
This proves the required $\OPT(A)=0$ endpoint for this explicit implementation.
It does not prove full row rank for arbitrary non-HSS inputs, uniform inverse
integrability, or a perturbation bound with constants independent of $A,N,L$.
In particular, exact recovery and pointwise continuity cannot replace a
residual-relative quantitative estimate near the HSS class.
\end{remark}

\section{Exact global error accounting}\label{base:sec:global}
\subsection{Pythagoras and recycled noise}
Put
\begin{align}
 S_\ell&=\sum_{j<\ell}\UC_jD_j\VC_j^{\T},&
 \Pi_\ell&=\UC_\ell\UC_\ell^{\T},&
 \T heta_\ell&=\VC_\ell\VC_\ell^{\T},\label{base:eq:global-projectors}\\
 R_\ell&=A-S_\ell-\Pi_\ell A\T heta_\ell.&&&\label{base:eq:accumulated-R}
\end{align}
The gauge at each earlier level implies
$\UC_\ell^{\T} S_\ell\VC_\ell=0$, and hence
\begin{equation}\label{base:eq:retained-zero}
 \Pi_\ell R_\ell\T heta_\ell=0.
\end{equation}
Define the true local error
\begin{equation}\label{base:eq:local-error}
 \err_\ell=A_\ell-D_\ell-U_\ell A_{\ell+1}V_\ell^{\T}.
\end{equation}
Then
\begin{equation}\label{base:eq:local-zero}
 U_\ell^{\T}\err_\ell V_\ell=0,
 \qquad
 R_{\ell+1}=R_\ell+\UC_\ell\err_\ell\VC_\ell^{\T}.
\end{equation}

\begin{proposition}[Exact error and measurement identities]\label{base:prop:pythagoras}
For every realization,
\begin{equation}\label{base:eq:pythagoras}
 \F{A-\widehat B}^2=\sum_{\ell=0}^{L-1}\F{\err_\ell}^2.
\end{equation}
Moreover,
\begin{align}
 Y_\ell-A_\ell\Omega_\ell
   &=\UC_\ell^{\T} R_\ell\Omega_0,\label{base:eq:forward-noise}\\
 Z_\ell-A_\ell^{\T}\Psi_\ell
   &=\VC_\ell^{\T} R_\ell^{\T}\Psi_0.\label{base:eq:transpose-noise}
\end{align}
\end{proposition}
\begin{proof}
The old residual in~\eqref{base:eq:local-zero} is orthogonal to every matrix of
the form $\UC_\ell M\VC_\ell^{\T}$ by~\eqref{base:eq:retained-zero}.
Induction therefore gives $\F{R_L}^2=\sum_\ell\F{\err_\ell}^2$.
Exact core recovery makes $R_L=A-\widehat B$.
For the measurement identity, induct in~\eqref{base:eq:data-recycle} to obtain
$Y_\ell=\UC_\ell^{\T}(A-S_\ell)\Omega_0$ and use
$\Omega_\ell=\VC_\ell^{\T}\Omega_0$ and~\eqref{base:eq:actual-core}.
The transpose calculation is identical.\end{proof}

\subsection{Which error can still influence the algorithm?}
By~\eqref{base:eq:retained-zero}, the accumulated error is the orthogonal sum
\begin{equation}\label{base:eq:three-sectors}
 R_\ell=
 \Pi_\ell R_\ell(I-\T heta_\ell)
 +(I-\Pi_\ell)R_\ell\T heta_\ell
 +(I-\Pi_\ell)R_\ell(I-\T heta_\ell).
\end{equation}
The first two sectors remain visible in the forward and transpose streams.
The last sector contributes to neither~\eqref{base:eq:forward-noise}
nor~\eqref{base:eq:transpose-noise}. Since the retained spaces only shrink, it
cannot affect later basis choices. It still counts toward final error.

Define the active and inactive energies
\begin{align}
 \Lambda_\ell&=\F{\UC_\ell^{\T} R_\ell}^2
                   +\F{\VC_\ell^{\T} R_\ell^{\T}}^2,\label{base:eq:active}\\
 J_\ell&=\F{(I-\Pi_\ell)R_\ell(I-\T heta_\ell)}^2.\label{base:eq:inactive}
\end{align}
Then $\F{R_\ell}^2=\Lambda_\ell+J_\ell$.
The old active energy discarded at the next step is
\begin{align}
 d_\ell={}&\F{(\Pi_\ell-\Pi_{\ell+1})R_\ell(I-\T heta_\ell)}^2\notag\\
 &+\F{(I-\Pi_\ell)R_\ell(\T heta_\ell-\T heta_{\ell+1})}^2.
                                                        \label{base:eq:d-energy}
\end{align}
The newly generated active and inactive energies are
\begin{align}
 t_\ell&=\F{U_\ell^{\T}\err_\ell}^2+\F{\err_\ell V_\ell}^2,
                                                        \label{base:eq:t-energy}\\
 v_\ell&=\F{(I-U_\ell U_\ell^{\T})\err_\ell
                      (I-V_\ell V_\ell^{\T})}^2.           \label{base:eq:v-energy}
\end{align}

\begin{theorem}[Exact active/inactive balance]\label{base:thm:balance}
Pathwise, at every level,
\begin{align}
 \Lambda_{\ell+1}&=\Lambda_\ell-d_\ell+t_\ell,\label{base:eq:active-balance}\\
 J_{\ell+1}&=J_\ell+d_\ell+v_\ell,\label{base:eq:inactive-balance}\\
 \F{\err_\ell}^2&=t_\ell+v_\ell.\label{base:eq:local-split}
\end{align}
In particular,
\begin{equation}\label{base:eq:discard-le-born}
 \sum_{\ell=0}^{L-1}d_\ell
 =\sum_{\ell=0}^{L-1}t_\ell-\Lambda_L
 \le\sum_{\ell=0}^{L-1}t_\ell.
\end{equation}
\end{theorem}
\begin{proof}
For the old forward-active sector, shrinking the retained row space removes
exactly the first term in~\eqref{base:eq:d-energy}. Its surviving columns lie
outside $\range\VC_\ell$. The new forward-active contribution lies inside
that range, so the two contributions are orthogonal. The transpose-active
calculation uses the analogous row orthogonality. This proves
\eqref{base:eq:active-balance}.

The removed portions of the old active sectors enter the inactive sector.
They are orthogonal to its old content and to the new inactive component:
respectively they occupy the old-retained/old-discarded,
old-discarded/old-retained, and newly-discarded/newly-discarded sectors.
This proves~\eqref{base:eq:inactive-balance}.
Finally,~\eqref{base:eq:local-zero} makes the three nonzero projection sectors of
$\err_\ell$ orthogonal, proving~\eqref{base:eq:local-split}.
Sum the active balance and use $\Lambda_0=0$ to get
\eqref{base:eq:discard-le-born}.\end{proof}

This bookkeeping is stronger than repeatedly charging the full accumulated
error as new contamination. It is nevertheless only an identity. It does
not assert that the new active energy has a subunit expected return gain.
That missing estimate is isolated in Section~\ref{base:sec:closure}.

\section{An exact quotient state for adaptive basis selection}\label{base:sec:quotient}
\subsection{Local encoding and its kernel}
In this section a local block has full-row-rank designs
$\Omega,\Psi\in\R^{m\times s}$ and data $Y,Z\in\R^{m\times s}$.
Suppress all level and node indices and define
\begin{align}
 B_Y&=Y\Omega^\dagger,&
 B_Z&=(\Psi^\dagger)^{\T} Z^{\T},\label{base:eq:BY-BZ}\\
 E_Y&=Y-B_Y\Omega,&
 E_Z&=Z-B_Z^{\T}\Psi,&
 \Delta&=B_Y-B_Z.\label{base:eq:encoding}
\end{align}
The innovations obey
\begin{equation}\label{base:eq:innovation-orth}
 E_Y\Omega^{\T}=0,\qquad E_Z\Psi^{\T}=0.
\end{equation}
If $N_\Omega$ is an orthonormal basis of $\ker\Omega$, then
$E_Y=(YN_\Omega)N_\Omega^{\T}$ and
$E_YE_Y^{\T}=(YN_\Omega)(YN_\Omega)^{\T}$.
Thus this encoding produces the algorithm's nullification-based choices.

\begin{proposition}[Common-coefficient invariance]\label{base:prop:quotient}
For fixed full-row-rank $\Omega,\Psi$, the kernel of the linear map
$(Y,Z)\mapsto(E_Y,E_Z,\Delta)$ is exactly
\begin{equation}\label{base:eq:quotient-kernel}
 \{(C\Omega,C^{\T}\Psi):C\in\R^{m\times m}\}.
\end{equation}
Adding $(C\Omega,C^{\T}\Psi)$ to the observations does not change $U,V$.
The change in the discrepancy is $C-PCQ$, and the changes in the compressed
observations are
\begin{equation}\label{base:eq:common-transport}
 (U^{\T}CV)\Omega',\qquad (U^{\T}CV)^{\T}\Psi'.
\end{equation}
\end{proposition}
\begin{proof}
Full row rank gives $\Omega\Omega^\dagger=I$ and
$\Psi\Psi^\dagger=I$. Thus the common shift adds $C$ to both regressions
and changes neither innovation nor mismatch. Conversely, zero innovations
imply $Y=B_Y\Omega$ and $Z=B_Z^{\T}\Psi$; zero mismatch makes their coefficients
equal. This proves the kernel description.

The basis rule is unchanged because both innovation Grams are unchanged.
Equation~\eqref{base:eq:local-D} changes by
$(I-P)C+PC(I-Q)=C-PCQ$. Substitution in the compressed data gives
$U^{\T}(PCQ)\Omega=(U^{\T}CV)\Omega'$ and its transpose counterpart.
\end{proof}

The autonomous basis-selection state is therefore
\begin{equation}\label{base:eq:quotient-state}
 \mathfrak s=(\Omega,\Psi,E_Y,E_Z,\Delta).
\end{equation}
The common coefficient is not recoverable from this state and must be kept
as a separate reconstruction register. The state is finite and local, but
its exact transport is nonlinear because the selected bases depend on its
innovations. No uniform moment contraction follows from the quotient alone.

\subsection{Exact compression formulas}\label{base:subsec:quotient-compression}
Complete the selected bases to orthogonal matrices
$[U\ U_\perp]$ and $[V\ V_\perp]$. Set
\begin{equation}\label{base:eq:delta-blocks}
 \begin{bmatrix}\Delta_{11}&\Delta_{12}\\
                 \Delta_{21}&\Delta_{22}\end{bmatrix}
 =\begin{bmatrix}U^{\T}\\ U_\perp^{\T}\end{bmatrix}
    \Delta[\,V\ V_\perp\,].
\end{equation}
Each block is $k\times k$. Define
\begin{align}
 G_\Omega&=V^{\T}\Omega,& H_\Omega&=V_\perp^{\T}\Omega,&
 C_\Omega&=H_\Omega G_\Omega^\dagger,&
 J_\Omega&=H_\Omega(I-G_\Omega^\dagger G_\Omega),\label{base:eq:Omega-split}\\
 G_\Psi&=U^{\T}\Psi,& H_\Psi&=U_\perp^{\T}\Psi,&
 C_\Psi&=H_\Psi G_\Psi^\dagger,&
 J_\Psi&=H_\Psi(I-G_\Psi^\dagger G_\Psi).\label{base:eq:Psi-split}
\end{align}
These retained designs have full row rank because the parent designs do.

\begin{theorem}[Compression in quotient coordinates]\label{base:thm:quotient-compression}
The state of the compressed observations is exactly
\begin{align}
 E_Y'&=U^{\T}E_Y+\Delta_{12}J_\Omega,\label{base:eq:q-EY}\\
 E_Z'&=V^{\T}E_Z-\Delta_{21}^{\T} J_\Psi,\label{base:eq:q-EZ}\\
 \Delta'&=\Delta_{11}+\Delta_{12}C_\Omega
                         +C_\Psi^{\T}\Delta_{21}.\label{base:eq:q-Delta}
\end{align}
A sufficient side register for reconstruction is $B_Z$, transported by
\begin{align}
 D&=B_Z-PB_ZQ+(I-P)\Delta,\label{base:eq:q-D-side}\\
 B_Z'&=U^{\T}B_ZV-C_\Psi^{\T}\Delta_{21}.\label{base:eq:q-BZ-side}
\end{align}
\end{theorem}
\begin{proof}
\proofstep{1. Expand the two compressed observations.}
Using $B_Y=B_Z+\Delta$ and the discrepancy formula gives
\begin{align}
 Y'&=U^{\T}B_ZV\,G_\Omega+U^{\T}\Delta\Omega+U^{\T}E_Y,
                                                        \label{base:eq:q-Yexpand}\\
 Z'&=V^{\T}B_Y^{\T}U\,G_\Psi-V^{\T}\Delta^{\T}\Psi+V^{\T}E_Z.
                                                        \label{base:eq:q-Zexpand}
\end{align}
\proofstep{2. Remove the local regressions.}
Multiply~\eqref{base:eq:q-Yexpand} by $I-G_\Omega^\dagger G_\Omega$.
The first term disappears. The old innovation is unchanged because it is
orthogonal to the whole row space of $\Omega$, hence to that of $G_\Omega$.
Also $\Omega(I-G_\Omega^\dagger G_\Omega)=V_\perp J_\Omega$.
This proves~\eqref{base:eq:q-EY}; transposition proves~\eqref{base:eq:q-EZ}.
\proofstep{3. Subtract the new coefficients.}
The regression coefficients extracted from~\eqref{base:eq:q-Yexpand}--\eqref{base:eq:q-Zexpand}
are
\begin{align*}
 B_Y'&=U^{\T}B_ZV+U^{\T}\Delta\Omega G_\Omega^\dagger,\\
 B_Z'&=U^{\T}B_YV-(\Psi G_\Psi^\dagger)^{\T}\Delta V.
\end{align*}
Use $\Omega G_\Omega^\dagger=V+V_\perp C_\Omega$ and
$\Psi G_\Psi^\dagger=U+U_\perp C_\Psi$.
The difference is~\eqref{base:eq:q-Delta}, and the second expression reduces to
\eqref{base:eq:q-BZ-side}. Equation~\eqref{base:eq:q-D-side} follows directly from
$B_Y=B_Z+\Delta$.\end{proof}

\begin{corollary}[Orthogonal innovation injection]\label{base:cor:q-energy}
Pathwise,
\begin{align}
 \F{E_Y'}^2&=\F{U^{\T}E_Y}^2+\F{\Delta_{12}J_\Omega}^2,\label{base:eq:q-energy-Y}\\
 \F{E_Z'}^2&=\F{V^{\T}E_Z}^2+\F{\Delta_{21}^{\T} J_\Psi}^2.\label{base:eq:q-energy-Z}
\end{align}
The mismatch block $\Delta_{22}$ does not occur in the outgoing quotient
state. If $\Delta=0$, compression only contracts the two innovations and
keeps zero mismatch.
\end{corollary}
\begin{proof}
The rows of $J_\Omega$ lie in the row space of $\Omega$, while
$E_Y\Omega^{\T}=0$. Thus the two terms in~\eqref{base:eq:q-EY} are orthogonal.
The transpose argument is the same. The other assertions are read directly
from~\eqref{base:eq:q-EY}--\eqref{base:eq:q-Delta}.\end{proof}

\subsection{Exact child merging}\label{base:subsec:quotient-merge}
Take two child states with $k$-row designs. Stack their designs and
innovations:
\begin{equation}\label{base:eq:merge-stack}
 \Omega=\begin{bmatrix}\Omega_1\\\Omega_2\end{bmatrix},\quad
 \Psi=\begin{bmatrix}\Psi_1\\\Psi_2\end{bmatrix},\quad
 \mathsf E_Y=\begin{bmatrix}E_{Y,1}\\E_{Y,2}\end{bmatrix},\quad
 \mathsf E_Z=\begin{bmatrix}E_{Z,1}\\E_{Z,2}\end{bmatrix}.
\end{equation}
Assume the stacked parent designs have full row rank.

\begin{theorem}[Merge in quotient coordinates]\label{base:thm:quotient-merge}
The parent state is
\begin{align}
 E_Y&=\mathsf E_Y(I-\Omega^\dagger\Omega),\label{base:eq:merge-EY}\\
 E_Z&=\mathsf E_Z(I-\Psi^\dagger\Psi),\label{base:eq:merge-EZ}\\
 \Delta&=\diag(\Delta_1,\Delta_2)
          +\mathsf E_Y\Omega^\dagger
          -(\Psi^\dagger)^{\T}\mathsf E_Z^{\T}.\label{base:eq:merge-Delta}
\end{align}
The reconstruction register is
\begin{equation}\label{base:eq:merge-BZ}
 B_Z=\diag(B_{Z,1},B_{Z,2})+(\Psi^\dagger)^{\T}\mathsf E_Z^{\T}.
\end{equation}
\end{theorem}
\begin{proof}
The stacked observations are
\begin{align*}
 Y&=\diag(B_{Y,1},B_{Y,2})\Omega+\mathsf E_Y,\\
 Z&=\diag(B_{Z,1},B_{Z,2})^{\T}\Psi+\mathsf E_Z.
\end{align*}
Compute their full-row-rank regressions and subtract them to get
\eqref{base:eq:merge-EY}--\eqref{base:eq:merge-BZ}.\end{proof}

Theorems~\ref{base:thm:quotient-compression} and~\ref{base:thm:quotient-merge}
prove closure under the actual alternating operations. Common coefficients
remain common under stacking, so the quotient invariance extends through
all future full-rank steps. The feedback that still requires a quantitative
bound is
\[
 \text{innovation}\ \xrightarrow{\text{merge regression}}\
 \text{mismatch}\ \xrightarrow{\text{compression}}\
 \text{new innovation}.
\]
This is not a linear contraction theorem, and the observation-space
innovation norm is not the true approximation-error norm.

\section{Moment budgets for the quotient state}\label{base:sec:quotient-moments}
\subsection{Initialization by the optimal residual}
Fix $B_*$ from~\eqref{base:eq:opt} and let $E_*=A-B_*$. Encode the observations
of $A$ and $B_*$ using the same initial designs. Their leaf-level state
difference is the encoding of $E_*$, because the encoding is linear in
$(Y,Z)$ with the designs fixed. Write this difference as
$(\delta E_{Y,I},\delta E_{Z,I},\delta\Delta_I)$.
Set
\begin{equation}\label{base:eq:q-alpha}
 q=s-m=6k,\qquad \alpha=\frac{m}{s-m-1}=\frac{2k}{6k-1}\le\frac25.
\end{equation}

\begin{proposition}[Initial residual-relative state budget]\label{base:prop:init-state}
For a leaf node $I$, put $R_I=E_*[I,I^c]$ and $C_I=E_*[I^c,I]$. Then
\begin{align}
 \E\F{\delta E_{Y,I}}^2&=q\F{R_I}^2,\label{base:eq:init-y}\\
 \E\F{\delta E_{Z,I}}^2&=q\F{C_I}^2,\label{base:eq:init-z}\\
 \E\F{\delta\Delta_I}^2&=\alpha(\F{R_I}^2+\F{C_I}^2).\label{base:eq:init-delta}
\end{align}
Consequently,
\begin{equation}\label{base:eq:init-budget}
 \E\sum_I\left[
 \frac{\F{\delta E_{Y,I}}^2+\F{\delta E_{Z,I}}^2}{q}
 +\frac{\F{\delta\Delta_I}^2}{\alpha}\right]
 \le4\,\OPT(A).
\end{equation}
\end{proposition}
\begin{proof}
The diagonal block of $E_*$ cancels from both innovations and from the
mismatch. In particular,
$\delta E_{Y,I}=R_I\Omega_{0,I^c}(I-\Omega_{0,I}^\dagger\Omega_{0,I})$.
Condition on $\Omega_{0,I}$. Its null projector has rank $q$, and the exterior
Gaussian matrix is independent. The Gaussian product identity of
Appendix~\ref{base:app:gaussian} gives~\eqref{base:eq:init-y}; the transpose calculation
gives~\eqref{base:eq:init-z}.

The two mismatch terms are
$R_I\Omega_{0,I^c}\Omega_{0,I}^\dagger$ and
$(\Psi_{0,I}^\dagger)^{\T}\Psi_{0,I^c}^{\T} C_I$.
They are independent and mean zero. Conditioning first on the local designs
and then using their inverse moments gives~\eqref{base:eq:init-delta}.
Each entry of $E_*$ occurs at most once in the sum of squared exterior
block-row norms and once in the analogous column sum. Thus
$\sum_I(\F{R_I}^2+\F{C_I}^2)\le2\F{E_*}^2$, which proves
\eqref{base:eq:init-budget}.\end{proof}

The budget concerns the state difference \emph{before} basis selection.
The two trajectories can then choose different bases. A Lipschitz or
residual-relative stability theorem for their actual nonlinear transport
has not been established; it cannot be replaced by assuming a common
trajectory for $A$ and $B_*$.

\subsection{A frozen-coefficient compression benchmark}
The following calculation isolates a legitimate local Gaussian reference
model. It is not applied to the adaptive hierarchy without further work.
Fix $U,V,\Delta$ independently of fresh Gaussian designs
$\Omega,\Psi\in\R^{2k\times s}$. Incoming innovations may depend on these
designs, but must satisfy~\eqref{base:eq:innovation-orth} and have finite second
moments. Define
\begin{equation}\label{base:eq:gamma-quotient}
 \gamma=\frac{k}{s-k-1},\qquad
 w=\frac{1-\gamma}{s-k},\qquad
 \mathscr V(E_Y,E_Z,\Delta)
 =w(\F{E_Y}^2+\F{E_Z}^2)+\F\Delta^2.
\end{equation}
Assume $s>2k+1$, so $w>0$.

\begin{proposition}[Frozen quotient energy]\label{base:prop:frozen-quotient}
Under these assumptions, applying~\eqref{base:eq:q-EY}--\eqref{base:eq:q-Delta} gives
\begin{align}
 \E\F{\Delta'}^2
   &=\F{\Delta_{11}}^2
       +\gamma(\F{\Delta_{12}}^2+\F{\Delta_{21}}^2),\label{base:eq:frozen-delta}\\
 \E\F{\Delta_{12}J_\Omega}^2&=(s-k)\F{\Delta_{12}}^2,\label{base:eq:frozen-inject}\\
 \E\mathscr V(E_Y',E_Z',\Delta')
   &\le\E\mathscr V(E_Y,E_Z,\Delta).\label{base:eq:frozen-energy}
\end{align}
At $s=8k$, $\gamma=k/(7k-1)\le1/6$.
\end{proposition}
\begin{proof}
The retained and discarded Gaussian row blocks in
\eqref{base:eq:Omega-split} are independent; the same holds in the other stream.
The matrices $C_\Omega,C_\Psi$ are mean zero and independent across streams.
Conditional Gaussian products and the inverse moment of a $k\times s$
design give~\eqref{base:eq:frozen-delta}. The null projector of the retained
design has rank $s-k$, giving~\eqref{base:eq:frozen-inject} and its transpose
counterpart. Use Corollary~\ref{base:cor:q-energy} and the contraction of
$U^{\T},V^{\T}$. For each off-diagonal mismatch block the total outgoing
coefficient is $w(s-k)+\gamma=1$. The $\Delta_{11}$ coefficient is one and
$\Delta_{22}$ is discarded. This proves~\eqref{base:eq:frozen-energy}.\end{proof}

The bound prices compression only. Merge can create mismatch from
innovations, and the actual $U,V,\Delta$ are not independent of the designs.
Neither issue is addressed by multiplying~\eqref{base:eq:frozen-energy} over
levels.

\section{The first compression and its recycled measurements}\label{base:sec:finite}
This section proves unconditional bounds on the original Gaussian
probability space. It does not use a smoothness hypothesis on the SVD.

\subsection{A row-dimension-sensitive Gaussian subspace bound}
\begin{lemma}[Local sample subspace]\label{base:lem:sample-subspace}
Let $M\in\R^{m\times n}$ be fixed and $G_0\in\R^{n\times q}$ standard
Gaussian. Let $P$ be a leading rank-$k$ left singular projector of $MG_0$.
If $q-m\ge4$, then
\begin{equation}\label{base:eq:sample-bound}
 \E\F{(I-P)M}^2
 \le\frac{e^2q(q-2/m)}{(q-m)^2-1}\F{M-[M]_k}^2.
\end{equation}
For $m=2k$, $q=6k$, the prefactor is at most
\begin{equation}\label{base:eq:gamma-local}
 \Gamma_0=\frac94e^2<17.
\end{equation}
\end{lemma}
\begin{proof}
Let $\Sigma=MM^{\T}$. The sample Gram has the distribution
$\Sigma^{1/2}W\Sigma^{1/2}$, where $W=GG^{\T}$ for a standard
$m\times q$ Gaussian $G$. This representation remains valid when
$\Sigma$ is singular. Let $\mu=\lambda_{\min}(W)$ and let $P_*$ be optimal
for the true covariance. Optimality of the sample projector gives
\begin{align*}
 \tr((I-P)\Sigma)
 &\le\mu^{-1}\tr((I-P)\Sigma^{1/2}W\Sigma^{1/2})\\
 &\le\mu^{-1}\tr((I-P_*)\Sigma^{1/2}W\Sigma^{1/2}).
\end{align*}
Rotational invariance gives
$\E[\mu^{-1}W]=\beta I_m$, with
$\beta=m^{-1}\E[\F G^2\N{G^\dagger}^2]$.
Write $G=\rho\widehat G$, where $\rho^2\sim\chi^2_{mq}$ is independent of
$\widehat G$. Since $\E\rho^{-2}=1/(mq-2)$,
\[
 \beta=(q-2/m)\E\N{G^\dagger}^2.
\]
The explicit spectral inverse bound in
Lemma~\ref{base:lem:gaussian-spectral} gives
$\E\N{G^\dagger}^2\le e^2q/((q-m)^2-1)$.
The best true projection error is $\F{M-[M]_k}^2$, proving
\eqref{base:eq:sample-bound}. Substitution yields
$(36k^2-6)/(16k^2-1)\le9/4$, proving~\eqref{base:eq:gamma-local}.
\end{proof}

\subsection{Same-level nullification and regression use independent coordinates}
At the initial level, abbreviate $\Omega_0,\Psi_0$ by $\Omega,\Psi$.
For a leaf $I$, write $R_I=A[I,I^c]$, $C_I=A[I^c,I]$, and define
\begin{equation}\label{base:eq:r-c-local}
 r_I=\F{(I-P_I)R_I}^2,\qquad
 c_I=\F{C_I(I-Q_I)}^2.
\end{equation}

\begin{lemma}[Same-level regression identity]\label{base:lem:same-level}
For the bases selected by the actual initial sketches,
\begin{align}
 \E\F{(I-P_I)R_I\Omega_{I^c}\Omega_I^\dagger}^2
     &=\alpha\E r_I,\label{base:eq:reg-r}\\
 \E\F{(\Psi_I^\dagger)^{\T}\Psi_{I^c}^{\T} C_I(I-Q_I)}^2
     &=\alpha\E c_I.\label{base:eq:reg-c}
\end{align}
\end{lemma}
\begin{proof}
Condition on $\Omega_I$. Complete an orthonormal null basis $N_I$ to
$[N_I\ T_I]$, where $T_I$ spans its row space. Conditional on the local
design, $\Omega_{I^c}N_I$ and $\Omega_{I^c}T_I$ are independent standard
Gaussian matrices. The selected $P_I$ depends on the first through
$Y_I N_I=R_I\Omega_{I^c}N_I$. The regression uses the second because
$\range\Omega_I^\dagger\subseteq\range T_I$.
Therefore the conditional expected squared norm in~\eqref{base:eq:reg-r} equals
$r_I\F{\Omega_I^\dagger}^2$. The conditional law of the nullified sample
is independent of $\Omega_I$, so $\E[r_I\mid\Omega_I]$ is constant.
Average the Gaussian inverse moment to obtain~\eqref{base:eq:reg-r}.
The transpose argument proves~\eqref{base:eq:reg-c}.\end{proof}

Let
\begin{align}
 F_I&=(I-P_I)R_I\Omega_{I^c}\Omega_I^\dagger,\label{base:eq:F-local}\\
 G_I&=(\Psi_I^\dagger)^{\T}\Psi_{I^c}^{\T} C_I(I-Q_I).\label{base:eq:G-local}
\end{align}
Expanding the discrepancy gives
\begin{equation}\label{base:eq:D-expansion}
 D_I=A_{II}-P_IA_{II}Q_I+F_I+P_IG_I.
\end{equation}
The last two terms have orthogonal output spaces, so
\begin{equation}\label{base:eq:diag-bound}
 \E\F{F_I+P_IG_I}^2\le\alpha(\E r_I+\E c_I).
\end{equation}
This orthogonality is pathwise and does not require $U_I,V_I$ to be
independent.

\begin{theorem}[First-compression error]\label{base:thm:first-error}
For arbitrary fixed $A$,
\begin{equation}\label{base:eq:first-error}
 \E\F{\err_0}^2
 \le\frac{63e^2}{10}\OPT(A)\le47\OPT(A).
\end{equation}
\end{theorem}
\begin{proof}
By nullification and Lemma~\ref{base:lem:sample-subspace},
$\E r_I\le\Gamma_0\F{R_I-[R_I]_k}^2$, and similarly for $c_I$.
The corresponding block rows and columns of $B_*$ have rank at most $k$.
Their residuals partition off-diagonal entries, giving
\begin{equation}\label{base:eq:sum-r-c}
 \sum_I\E r_I\le\Gamma_0\OPT(A),\qquad
 \sum_I\E c_I\le\Gamma_0\OPT(A).
\end{equation}
For $I\ne J$,
$A_{IJ}-P_IA_{IJ}Q_J=(I-P_I)A_{IJ}+P_IA_{IJ}(I-Q_J)$,
an orthogonal sum. Thus the off-diagonal error is at most
$\sum_I(r_I+c_I)$. Add~\eqref{base:eq:diag-bound} to get
$\E\F{\err_0}^2\le2(1+\alpha)\Gamma_0\OPT(A)$.
Since $\alpha\le2/5$ and $\Gamma_0=9e^2/4$, the coefficient is at most
$63e^2/10$.\end{proof}

\subsection{The actual first recycled error}
Define the block-diagonal matrix $T_0$ by
\begin{equation}\label{base:eq:T0}
 (T_0)_I=(\Psi_I^\dagger)^{\T} Z_I^{\T}(I-Q_I),
 \qquad
 H_0=(A-T_0)(I-V_0V_0^{\T}).
\end{equation}
The pair $(H_0,V_0)$ depends only on $\Psi$, and
\begin{equation}\label{base:eq:H0-orth}
 H_0V_0=0.
\end{equation}
Its off-diagonal blocks are $A_{IJ}(I-Q_J)$ and its diagonal blocks are
$-G_I$. Hence
\begin{equation}\label{base:eq:H0-bound}
 \E\F{H_0}^2=(1+\alpha)\sum_I\E c_I
 \le\frac{63e^2}{20}\OPT(A)\le24\OPT(A).
\end{equation}
Since $U_0^{\T}D_0=U_0^{\T}T_0$ and $T_0V_0=0$, direct substitution gives
\begin{equation}\label{base:eq:first-noise-exact}
 Y_1=A_1\Omega_1+U_0^{\T}H_0\Omega.
\end{equation}
Although $U_0$ depends on $\Omega$, it can be removed by contraction before
expectation. Since $H_0$ is independent of $\Omega$,
\begin{equation}\label{base:eq:first-noise-bound}
 \frac1s\E\F{Y_1-A_1\Omega_1}^2
 \le\E\F{H_0}^2\le24\OPT(A).
\end{equation}
The transpose construction
\begin{equation}\label{base:eq:Ht0}
 (\widetilde T_0)_I=(\Omega_I^\dagger)^{\T}Y_I^{\T}(I-P_I),\qquad
 \widetilde H_0=(A^{\T}-\widetilde T_0)(I-U_0U_0^{\T})
\end{equation}
depends only on $\Omega$, annihilates $U_0$, and proves
\begin{equation}\label{base:eq:first-transpose-noise-bound}
 \frac1s\E\F{Z_1-A_1^{\T}\Psi_1}^2\le24\OPT(A).
\end{equation}
These are residual-relative expectation bounds without an exceptional-event
term proportional to $\F A^2$.

\section{A genuine first cross-level reuse bound}\label{base:sec:second}
The properties retained in~\eqref{base:eq:H0-orth} are essential. Conditional on
$\Psi$, the matrices
\begin{equation}\label{base:eq:retained-discarded-Gaussian}
 V_0^{\T}\Omega,\qquad V_{0,\perp}^{\T}\Omega
\end{equation}
are independent standard Gaussian matrices, and
$H_0\Omega=H_0V_{0,\perp}(V_{0,\perp}^{\T}\Omega)$.
The remaining left compression can depend on all the randomness. The next
lemma is designed to permit exactly that dependence.

\subsection{An adaptive-left lemma}
\begin{lemma}[Adaptive left compression and regression]\label{base:lem:adaptive-left}
Fix $C,K\in\R^{p\times n}$ and $H\in\R^{p\times h}$, with
$\rank K\le k$, $n\ge k$, and $p\ge m=2k$. Put $E=C-K$ and
$w=\F E^2+\F H^2$. Independently draw standard Gaussian matrices
\[
 O\in\R^{m\times s},\quad G\in\R^{n\times s},\quad
 W\in\R^{h\times s},\qquad s=8k.
\]
Let $N$ be an orthonormal basis for $\ker O$ and let
$J\in\R^{p\times m}$ be an arbitrary measurable isometry depending on all
these matrices and any additional randomness. Let $P$ be a leading rank-$k$
left singular projector of
\begin{equation}\label{base:eq:adaptive-left-sample}
 S=J^{\T}(CG+HW)N.
\end{equation}
There are constants $\alpha\le2/5$, $\beta<12$, $\theta<12$ such that
\begin{align}
 \left(\E\F{(I-P)J^{\T}C}^2\right)^{1/2}
       &\le\F E+2\sqrt{\beta w},\label{base:eq:adaptive-left-proj}\\
 \E\F{(I-P)J^{\T}(CG+HW)O^\dagger}^2
       &\le(2\sqrt{\beta\theta}+\sqrt\alpha)^2w.
                                                        \label{base:eq:adaptive-left-reg}
\end{align}
\end{lemma}
\begin{proof}
\proofstep{1. Remove the adaptive contractions pathwise.}
Choose an $n\times k$ isometry $R$ containing the row space of $K$.
Set $B=R^{\T}GN$ and $\Delta_s=(EG+HW)N$. The first term of
$S=J^{\T}KR B+J^{\T}\Delta_s$ has rank at most $k$.
Optimality of $P$ gives $\F{(I-P)S}\le\F{\Delta_s}$.
Since $B$ has full row rank almost surely,
\begin{equation}\label{base:eq:adaptive-left-key}
 \F{(I-P)J^{\T}KR}
 \le2\F{\Delta_s}\N{B^\dagger}.
\end{equation}
Neither $J$ nor $P$ remains on the right-hand side.

\proofstep{2. Integrate the coupled sample and inverse.}
Let $q=s-m=6k$, $\mu=\E\N{G_{k\times q}^\dagger}^2$, and $\beta=q\mu$.
Decompose $GN$ along $R$ and its orthogonal complement. Cross terms have
zero conditional expectation, and rotational/radial invariance gives
\begin{align}
 \E[\F{\Delta_s}^2\N{B^\dagger}^2]
 ={}&\mu(q-2/k)\F{ER}^2\notag\\
 &+q\mu\F{E(I-RR^{\T})}^2+q\mu\F H^2\le\beta w.
                                                        \label{base:eq:adaptive-left-mixed}
\end{align}
For the first term, $\E[BB^{\T}\N{B^\dagger}^2]$ is scalar and its scalar
coefficient is $(q-2/k)\mu$ by the same radius calculation as in
Lemma~\ref{base:lem:sample-subspace}. Combining~\eqref{base:eq:adaptive-left-key}
with $(I-P)J^{\T}C=(I-P)J^{\T}E+(I-P)J^{\T}K$ and applying $L_2$ Minkowski gives
\eqref{base:eq:adaptive-left-proj}.

\proofstep{3. Separate null and row-space coordinates for regression.}
Complete $N$ by an orthonormal row-space basis $T$ of $O$. Pathwise,
\begin{align}
 \F{(I-P)J^{\T}(CG+HW)O^\dagger}
 \le{}&2\F{\Delta_s}\N{B^\dagger}\,
                \N{R^{\T}GO^\dagger}\notag\\
 &+\F{(EG+HW)O^\dagger}.\label{base:eq:adaptive-left-reg-path}
\end{align}
Conditional on $O$, the first factor uses only the $N$ coordinates of $G,W$,
while $R^{\T}GO^\dagger$ uses only the $T$ coordinates of $G$.
The null-coordinate law does not depend on $O$. Therefore
\[
 \E[\F{\Delta_s}^2\N{B^\dagger}^2\N{R^{\T}GO^\dagger}^2]
 \le\beta\theta w,
\]
where one may take
\[
 \theta=\E\N{G_{k\times2k}}^2\,
              \E\N{G_{2k\times8k}^\dagger}^2.
\]
The other term satisfies
$\E\F{(EG+HW)O^\dagger}^2=\alpha w$.
Applying Minkowski to~\eqref{base:eq:adaptive-left-reg-path} proves
\eqref{base:eq:adaptive-left-reg}.

\proofstep{4. Evaluate sufficient constants.}
Appendix~\ref{base:app:gaussian} gives
\[
 \beta\le\frac{36e^2k^2}{25k^2-1}\le\frac32e^2<12,
\]
and
\[
 \theta\le\frac{8e^2(4+2\sqrt2)}{35}<12,
 \qquad \alpha=\frac{2k}{6k-1}\le\frac25.
\]
For the bound on $\theta$, use
$\E\N{G_{k\times2k}}^2\le(\sqrt k+\sqrt{2k})^2+1$ and the spectral
inverse bound.\end{proof}

\subsection{Application to the second compression}
\begin{theorem}[First two compressions]\label{base:thm:two-compressions}
For every fixed $A$ and every $L\ge2$, the initial $16k$ queries, reused
without fresh sampling, give
\begin{equation}\label{base:eq:two-compressions}
 \E\bigl[\F{\err_0}^2+\F{\err_1}^2\bigr]
 \le33000\,\OPT(A).
\end{equation}
The number $33000$ is a loose explicit bound, not an optimized constant.
For $L>2$, the theorem bounds only these two local errors.
\end{theorem}
\begin{proof}
\proofstep{1. Condition on one stream, not on the selected left bases.}
A parent node $J$ consists of two original leaves, hence has $4k$ original
indices. Conditional on $\Psi$, define fixed matrices
\begin{align*}
 C_J&=A[J,J^c]V_{0,J^c},&
 K_J&=B_*[J,J^c]V_{0,J^c},\\
 E_J&=C_J-K_J,&
 H_J&=H_0[J,:]V_{0,\perp}.
\end{align*}
Here $V_{0,J^c}$ denotes the exterior part of the block-diagonal isometry.
The cut-rank property gives $\rank K_J\le k$.
The restriction $U_{0,J}$ is a $4k\times2k$ isometry and may depend on all
forward Gaussian coordinates. By~\eqref{base:eq:first-noise-exact}, the
nullified sample used at this parent is exactly of the form
\eqref{base:eq:adaptive-left-sample}, with $J$ in that lemma equal to $U_{0,J}$.
The true off-diagonal block row is $U_{0,J}^{\T}C_J$.

\proofstep{2. Sum the residual-relative budgets.}
The parent row sets partition the indices, so orthogonal contraction gives
\begin{equation}\label{base:eq:parent-budgets}
 \sum_J\F{E_J}^2\le\OPT(A),\qquad
 \E\sum_J\F{H_J}^2=\E\F{H_0}^2\le24\OPT(A).
\end{equation}
Apply~\eqref{base:eq:adaptive-left-proj} in the direct sum over parents, followed
by Minkowski. The expected sum of second-level row projection losses is at
most
\begin{equation}\label{base:eq:second-projection}
 (1+2\sqrt{300})^2\OPT(A).
\end{equation}
The transpose proof conditions on $\Omega$ and uses $\widetilde H_0$ from
\eqref{base:eq:Ht0}, giving the same column bound.

\proofstep{3. Include the actual reused regression errors.}
The diagonal contribution from the forward stream at a parent is exactly
\[
 (I-P_{1,J})U_{0,J}^{\T}(C_JG+H_JW)O^\dagger.
\]
Lemma~\ref{base:lem:adaptive-left} and~\eqref{base:eq:parent-budgets} bound the sum of
its squared norms in expectation by
\begin{equation}\label{base:eq:second-regression}
 25(24+\sqrt{2/5})^2\OPT(A).
\end{equation}
The transpose contribution has the same bound. The two diagonal terms
have complementary output spaces, and each off-diagonal residual has the
same row/column orthogonal splitting as at the first level. Therefore
\begin{equation}\label{base:eq:second-total}
 \E\F{\err_1}^2
 \le2\bigl[(1+2\sqrt{300})^2
            +25(24+\sqrt{2/5})^2\bigr]\OPT(A).
\end{equation}
Adding $47\OPT(A)$ from Theorem~\ref{base:thm:first-error} gives a coefficient
strictly below $33000$.\end{proof}

\begin{corollary}[What is complete at small depth]
With exact core recovery, the candidate algorithm has expected error
$\le47\OPT(A)$ when $L=1$, and $\le33000\OPT(A)$ when $L=2$, using $18k$
queries. At fixed depth, $O(kL)$ is already $O(k)$; the value of the second
case is its treatment of genuine cross-level dependence, not a
fixed-depth query-complexity separation.
\end{corollary}

\subsection{A complete fresh-between-levels comparison}
If fresh Gaussian sketches are acquired at every level, the current
$A_\ell$ is fixed conditional on the past. Theorem~\ref{base:thm:first-error}
then applies conditionally at each step. Also
\begin{equation}\label{base:eq:opt-nonincreasing}
 \OPT_{L-\ell-1,k}(A_{\ell+1})
 \le\OPT_{L-\ell,k}(A_\ell)
\end{equation}
pathwise: compress an optimal comparator with the same aligned isometries.
Every coarser cut rank remains at most $k$, and the Frobenius error contracts.
The gauge removes $D_\ell$ from both compressed matrices.
Thus the modified acquisition has the complete guarantee
\begin{equation}\label{base:eq:fresh-endpoint}
 \E\F{A-\widehat B}^2\le47L\OPT(A),
 \qquad \#\text{queries}=16kL+2k.
\end{equation}
Each compressed query is simulated by lifting its vector with $\VC_\ell$
and projecting the original answer with $\UC_\ell^{\T}$, or vice versa.
This is a fresh-sketch benchmark, not the requested depth-independent
endpoint.

\section{Adaptive inverse control and its limits}\label{base:sec:inverse}
\subsection{A valid bound for arbitrary adaptive compression}
\begin{lemma}[Inverse of an adaptive compression]\label{base:lem:inverse-compression}
Let $G\in\R^{d\times s}$ be standard Gaussian with $s>d+1$.
Let $W\in\R^{d\times r}$ be an arbitrary measurable isometry, possibly
depending on $G$ and other randomness. Then
\begin{equation}\label{base:eq:inverse-compression}
 \E\F{(W^{\T}G)^\dagger}^2\le\frac{d}{s-d-1}.
\end{equation}
\end{lemma}
\begin{proof}
For a positive-definite $K$, the block-inverse formula in an orthogonal
basis extending $W$ gives
\[
 (W^{\T}KW)^{-1}\preceq W^{\T}K^{-1}W.
\]
Indeed, the right side is the inverse of a Schur complement, which dominates
the inverse of the corresponding diagonal block. Take $K=GG^{\T}$.
Then $\F{(W^{\T}G)^\dagger}^2\le\tr K^{-1}$ pathwise, regardless of the
choice of $W$. Average the inverse moment from
Lemma~\ref{base:lem:gaussian-basic}.\end{proof}

For a prescribed block at the third compression (level $\ell=2$), the
local forward design is an adaptive orthogonal compression of a
$4k\times8k$ block of $\Omega_1$. Conditional on $\Psi_0$, that ambient
block is standard Gaussian. Thus
\begin{equation}\label{base:eq:third-inverse}
 \E\F{\Omega_{2,I}^\dagger}^2
 \le\frac{4k}{4k-1}\le\frac43.
\end{equation}
The transpose bound follows by conditioning on $\Omega_0$.
This controls an inverse by itself; it does not control its product with
an adaptive residual coefficient.

\subsection{Why generic nestedness is insufficient}
\begin{proposition}[Infinite inverse moment under nested compression]
\label{base:prop:nested-obstruction}
There is a measurable binary hierarchy of rank-$k$ orthogonal compressions
of a standard Gaussian $G\in\R^{8k\times8k}$ whose cumulative isometry
$W(G)\in\R^{8k\times2k}$ satisfies: $W^{\T}G$ has full row rank almost surely,
but
\begin{equation}\label{base:eq:nested-infinite}
 \E\F{(W^{\T}G)^\dagger}^2=\infty.
\end{equation}
\end{proposition}
\begin{proof}
Let $u$ be a measurable unit left singular vector for the smallest singular
value of $G$. Split its coordinates into four groups of size $2k$. At each
group retain a $k$-dimensional subspace containing the restriction of $u$.
Merge adjacent groups and retain, in each merged compressed half, a
$k$-dimensional subspace containing the resulting restricted coordinates of
$u$. Pad with a deterministic orthonormal completion whenever needed. This
constructs a nested isometry $W$ whose range contains $u$.

For $K=GG^{\T}$, compression cannot lower its smallest Rayleigh quotient, and
the range of $W$ contains a minimizing vector. Therefore
\[
 \lambda_{\min}(W^{\T}KW)=\lambda_{\min}(K),
 \quad\text{so}\quad
 \sigma_{\min}(W^{\T}G)=\sigma_{\min}(G).
\]
The square Gaussian matrix is invertible almost surely. But the last
diagonal entry of $K^{-1}$ is the reciprocal squared distance of its last
row from the span of the preceding rows. Conditional on those rows, the
distance has distribution $|Z|$ for $Z\sim N(0,1)$, and $\E Z^{-2}=\infty$.
Since $\N{K^{-1}}\ge(K^{-1})_{dd}$, the expected squared reciprocal of the
smallest singular value is infinite. This proves~\eqref{base:eq:nested-infinite}.
\end{proof}

This construction is not asserted to be produced by the SVD rule for a
fixed input $A$. It rules out a proof using only initial Gaussianity,
orthonormality, and nesting. With ambient dimension larger than the sketch
width, unrestricted adaptive nesting can even retain a vector in the left
kernel. The actual algorithm's coefficient relations must do additional work.

\begin{example}[A scalar noise-energy bound cannot price a correlated inverse]
\label{base:ex:correlated-inverse}
Let $g\in\R^{1\times s}$ be standard Gaussian and let
$p_\varepsilon=\Prb\{\N g\le\varepsilon\}>0$. Put
\[
 f(g)=p_\varepsilon^{-1/2}\,
      \mathbf1_{\{\N g\le\varepsilon\}}\frac{g}{\N g}.
\]
Then $\E\N{f(g)}^2=1$, but
\[
 \E|f(g)g^\dagger|^2
 =\E[\N g^{-2}\mid\N g\le\varepsilon]\ge\varepsilon^{-2}.
\]
Separate second moments of the noise and of a Gaussian inverse therefore
do not imply a dimension-free bound for their correlated product.
\end{example}

\section{Wick--Hermite interpretation and an adaptive regression identity}
\label{base:sec:stein}
\subsection{What the framework contributes here}
For a Gaussian vector $g$ and normalized Hermite polynomials $h_j$,
\[
 xh_j(x)=\sqrt{j+1}\,h_{j+1}(x)+\sqrt j\,h_{j-1}(x).
\]
Thus multiplication is represented exactly by creation and annihilation in
Hermite coordinates. For a self-adjoint random matrix $Z$ with the required
moments, multiplication by $Z$ on $\R^D\otimes L^2$ gives the vacuum identity
\begin{equation}\label{base:eq:vacuum}
 \E\tr Z^{2q}
 =\sum_{a=1}^D\N{\mathcal M_Z^q(e_a\otimes1)}_{L^2}^{2}.
\end{equation}
This follows directly by expanding the squared norm; it is a representation,
not an estimate. In the present problem the base randomness is genuinely
Gaussian. No non-Gaussian transfer or comparison distribution is used.

The quotient results provide exact relations worth retaining before a norm
is taken. The finite-depth proofs show which Gaussian coordinates can be
integrated out after adaptive contractions have been removed. Neither result
asserts that SVD projectors or pseudoinverses are finite-degree polynomials.
Consequently a finite-grade polynomial moment theorem does not immediately
bound the complete algorithm.

A small diagnostic illustrates the distinction. For an adaptive matrix
$R(\Omega)$ with $R^{\T}R\in L^2$,
\begin{align}
 \E\F{R\Omega}^2
 ={}&s\E\F R^2
   +\E\tr\bigl[R^{\T}R(\Omega\Omega^{\T}-sI)\bigr]\notag\\
 ={}&s\E\F R^2
   +\E\tr\bigl[\Pi_2(R^{\T}R)(\Omega\Omega^{\T}-sI)\bigr].
                                                        \label{base:eq:second-chaos}
\end{align}
Here $\Pi_2$ is entrywise orthogonal projection onto the second Hermite
chaos of the original Gaussian entries. The centered quadratic matrix has
only that Hermite degree, which proves the second equality. This one
expectation only sees $\Pi_2(R^{\T}R)$, but later products or inverses can see
information discarded by that projection. The quotient's continuation test
must precede a quantitative relaxation.

\subsection{The active return has an orthogonal-regression structure}
Let $X_\ell=\UC_\ell^{\T}R_\ell$. For a current block $I$, let $E_I$ embed
its $2k$ coordinates into the current compressed column space and set
\begin{equation}\label{base:eq:return-coeff}
 C_{\ell,I}=(I-P_{\ell,I})X_\ell[I,:],\qquad
 W_{\ell,I}=\VC_\ell E_I.
\end{equation}
Then $W_{\ell,I}^{\T}W_{\ell,I}=I_{2k}$ and
$C_{\ell,I}W_{\ell,I}=0$ by~\eqref{base:eq:retained-zero}.
The old forward-noise contribution returning to a retained column space is
\begin{equation}\label{base:eq:active-return}
 C_{\ell,I}\Omega_0
       (W_{\ell,I}^{\T}\Omega_0)^\dagger Q_{\ell,I}.
\end{equation}
The corresponding transpose expression exchanges the two streams.
Summing the squared coefficient norms $\F{C_{\ell,I}}^2$ and their
transpose analogues gives exactly $d_\ell$.

For fixed $C,W,P$ with $CW=0$, $W^{\T}W=I_m$ and $P$ a rank-$r$ projector,
independence of orthogonal Gaussian coordinates yields
\begin{equation}\label{base:eq:fresh-return}
 \E\F{C\Xi(W^{\T}\Xi)^\dagger P}^2
 =\frac{r}{s-m-1}\F C^2.
\end{equation}
Condition first on $M=W^{\T}\Xi$ and then use
$\E(MM^{\T})^{-1}=I_m/(s-m-1)$ to prove this formula.
For~\eqref{base:eq:active-return}, the fixed-coefficient benchmark is
\begin{equation}\label{base:eq:return-fifth}
 \frac{k}{s-2k-1}=\frac{k}{6k-1}\le\frac15.
\end{equation}
It concerns the \emph{returning rank-$k$ part}, not the whole regression
error. The HSS coefficients in~\eqref{base:eq:return-coeff} are adaptive, so
\eqref{base:eq:fresh-return} is not directly applicable.

\subsection{Regularization and the signed adaptation correction}
The next identity handles adaptive coefficients under explicit smoothness
conditions. It is an application-specific Gaussian integration-by-parts
calculation, not a replacement for those conditions. Gaussian Stein
identities and their differentiability requirements are classical; see,
for example, Bellec--Zhang~\cite{BellecZhang}.

Let $\Xi\in\R^{n\times s}$ be standard Gaussian. Suppose
$C=C(\Xi)\in\R^{p\times n}$, $W=W(\Xi)\in\R^{n\times m}$, and
$P=P(\Xi)\in\R^{m\times m}$ satisfy
\begin{equation}\label{base:eq:adaptive-reg-hyp}
 W^{\T}W=I_m,\quad CW=0,\quad P=P^{\T}=P^2,\quad \rank P=r,
 \qquad a=s-m-1>0.
\end{equation}
For $\lambda>0$, define
\begin{equation}\label{base:eq:regularized-design}
 M=W^{\T}\Xi,\quad K_\lambda=MM^{\T}+\lambda I_m,\quad
 X_\lambda=M^{\T}K_\lambda^{-1},\quad c=\F C^2,
\end{equation}
and the $n\times s$ vector field
\begin{equation}\label{base:eq:Stein-field}
 Z_\lambda=C^{\T}C\,\Xi X_\lambda P X_\lambda^{\T}
                   -\frac ca WPX_\lambda^{\T}.
\end{equation}
The \emph{frozen divergence} differentiates only explicit occurrences of
$\Xi$, holding $C,W,P$ at their current values. The total divergence also
differentiates their dependence on $\Xi$. Put
\begin{equation}\label{base:eq:Stein-charge}
 \mathfrak A_\lambda
 =\diver_\Xi Z_\lambda
       -\diver_\Xi^{\mathrm{frozen}}Z_\lambda.
\end{equation}
Assume these derivatives and Gaussian integration by parts are legitimate
and the displayed expectations are absolutely integrable. One sufficient
condition is globally $C^1$ coefficient maps with their entries and first
derivatives of at most polynomial growth. Smooth compactly supported
localizations are another route, provided all resulting derivative terms
are retained.

\begin{theorem}[Regularized adaptive return identity]\label{base:thm:Stein}
Under these hypotheses,
\begin{equation}\label{base:eq:Stein-identity}
 \E\F{C\Xi X_\lambda P}^2
 =\frac ra\E c+\E\mathfrak A_\lambda-\E\mathfrak R_\lambda,
\end{equation}
where
\begin{align}
 \mathfrak R_\lambda=\frac{\lambda c}{a}\bigl[
 &\tr(PK_\lambda^{-1})
 +\tr(K_\lambda^{-1})\tr(PK_\lambda^{-1})\notag\\
 &+(a+1)\tr(PK_\lambda^{-2})\bigr]\ge0.
                                                     \label{base:eq:Stein-penalty}
\end{align}
\end{theorem}
\begin{proof}
\proofstep{1. Compute the Gaussian inner product.}
Using $P^2=P$ and $K_\lambda^{-1}MM^{\T}=I-\lambda K_\lambda^{-1}$,
\begin{equation}\label{base:eq:Stein-inner}
 \ip{\Xi}{Z_\lambda}
 =\F{C\Xi X_\lambda P}^2
  -\frac ca\bigl[r-\lambda\tr(PK_\lambda^{-1})\bigr].
\end{equation}
\proofstep{2. Compute the frozen divergence.}
In coordinates extending $W$, $C$ is supported in the orthogonal-complement
rows, while $X_\lambda$ depends only on the retained Gaussian rows. Hence
\begin{align}
 \diver_\Xi^{\mathrm{frozen}}
   (C^{\T}C\,\Xi X_\lambda P X_\lambda^{\T})
  &=c\tr(PX_\lambda^{\T}X_\lambda)\notag\\
  &=c\tr(PK_\lambda^{-1})
       -\lambda c\tr(PK_\lambda^{-2}).\label{base:eq:Stein-div-first}
\end{align}
For the other field, differentiation in the retained coordinates reduces
to the divergence of $PK_\lambda^{-1}M$ with respect to $M$. The direct
term is $s\tr(PK_\lambda^{-1})$. Differentiating the two occurrences of
$M$ in $MM^{\T}$ subtracts respectively
\[
 \tr(PK_\lambda^{-1})\tr(M^{\T}K_\lambda^{-1}M),
 \qquad
 \tr(PK_\lambda^{-1}MM^{\T}K_\lambda^{-1}).
\]
Since these traces are $m-\lambda\tr K_\lambda^{-1}$ and
$\tr(PK_\lambda^{-1})-\lambda\tr(PK_\lambda^{-2})$ in the appropriate
factors, we get
\begin{align}
 \diver_\Xi^{\mathrm{frozen}}(WPX_\lambda^{\T})
 ={}&a\tr(PK_\lambda^{-1})\notag\\
 &+\lambda\tr(K_\lambda^{-1})\tr(PK_\lambda^{-1})
   +\lambda\tr(PK_\lambda^{-2}).\label{base:eq:Stein-div-second}
\end{align}
Combining~\eqref{base:eq:Stein-div-first}--\eqref{base:eq:Stein-div-second},
\begin{align}
 \diver_\Xi^{\mathrm{frozen}}Z_\lambda
 =-\frac{\lambda c}{a}\bigl[
 &\tr(K_\lambda^{-1})\tr(PK_\lambda^{-1})\notag\\
 &+(a+1)\tr(PK_\lambda^{-2})\bigr].\label{base:eq:frozen-divergence}
\end{align}
\proofstep{3. Apply integration by parts.}
For a scalar standard Gaussian coordinate, integration of
$(f(x)e^{-x^2/2})'$ gives $\E[xf(x)]=\E[f'(x)]$ when the boundary and
integrability conditions hold. Apply this to every entry of $Z_\lambda$:
$\E\ip\Xi{Z_\lambda}=\E\diver_\Xi Z_\lambda$.
Insert~\eqref{base:eq:Stein-inner}, \eqref{base:eq:Stein-charge}, and
\eqref{base:eq:frozen-divergence} to obtain~\eqref{base:eq:Stein-identity}.
All traces in~\eqref{base:eq:Stein-penalty} are nonnegative.\end{proof}

\begin{corollary}[A sufficient limit bound]\label{base:cor:Stein-limit}
Suppose the identity holds for every $\lambda>0$, $\E c<\infty$, and
\[
 \sup_{0<\lambda\le1}\E\mathfrak A_\lambda\le\beta\E c.
\]
Then
\begin{equation}\label{base:eq:Stein-limit}
 \E\F{C\Xi(W^{\T}\Xi)^\dagger P}^2
 \le\left(\frac r{s-m-1}+\beta\right)\E c.
\end{equation}
\end{corollary}
\begin{proof}
Drop the nonnegative penalty. Since $X_\lambda\to M^\dagger$ pointwise
as $\lambda\downarrow0$, including at deficient rank, Fatou's lemma proves
the claim. No integrability of the unregularized risk is presupposed.
\end{proof}

For the actual algorithm, derivatives of selected singular subspaces may
fail the stated conditions at ties or near rank loss. Neither a negligible
set of ties nor the pointwise existence of a pseudoinverse justifies all
boundary limits. A valid localization argument and a uniform bound for its
extra charges remain required. The regularizer here changes the observable
in the proof, not the algorithm's discrepancy formula.

\subsection{The adaptation correction has no universal favorable sign}
Keep $C,W$ fixed and let $P$ select the $r$ weakest left singular directions
of $M=W^{\T}\Xi$, with $0<r<m$. Conditional on $M$,
\[
 \E[\F{C\Xi M^\dagger P}^2\mid M]
 =\F C^2\tr\bigl(P(MM^{\T})^{-1}\bigr).
\]
Almost surely, the selected sum of the largest $r$ inverse eigenvalues is
strictly larger than $r/m$ times their total. Thus the expected risk is
strictly larger than the fresh coefficient $r\F C^2/(s-m-1)$.
This example does not describe the HSS selection rule, but rules out
assuming a nonpositive adaptation penalty merely because a projector is
being selected.

\section{A sufficient all-depth closure theorem}\label{base:sec:closure}
The preceding exact identities can now be separated cleanly from the
missing probability estimate. At level $\ell$, let $\mathcal I_\ell$ be
the original node partition processed there, and define
\begin{equation}\label{base:eq:eta}
 \eta_\ell=\sum_{I\in\mathcal I_\ell}
 \bigl(\F{(A-B_*)[I,I^c]}^2+\F{(A-B_*)[I^c,I]}^2\bigr).
\end{equation}
Each row and column partition charges each entry at most once, so
\begin{equation}\label{base:eq:H-budget}
 H:=\sum_{\ell=0}^{L-1}\eta_\ell\le2L\OPT(A).
\end{equation}
Set
\[
 T=\sum_\ell\E t_\ell,\qquad
 D=\sum_\ell\E d_\ell,\qquad
 V=\sum_\ell\E v_\ell.
\]

\begin{theorem}[Conditional global completion]\label{base:thm:closure}
Suppose these quantities are finite and the actual reuse algorithm satisfies
\begin{equation}\label{base:eq:missing-T}
 T\le\rho D+bH,\qquad 0\le\rho<1,
\end{equation}
and
\begin{equation}\label{base:eq:missing-V}
 V\le aH+cD,
\end{equation}
with absolute nonnegative $a,b,c$ independent of $A,N,k,L$.
Then the $18k$-query output satisfies
\begin{equation}\label{base:eq:conditional-endpoint}
 \E\F{A-\widehat B}^2
 \le2\left(a+\frac{(1+c)b}{1-\rho}\right)L\OPT(A).
\end{equation}
\end{theorem}
\begin{proof}
Equation~\eqref{base:eq:discard-le-born} gives $D\le T$, so
$T\le\rho T+bH$ and $T\le bH/(1-\rho)$.
Hence $T+V\le[a+(1+c)b/(1-\rho)]H$.
The exact error identity and~\eqref{base:eq:local-split} identify the expected
error with $T+V$. Apply~\eqref{base:eq:H-budget}.\end{proof}

An alternative sufficient formulation avoids assuming integrability at the
start. Let $\mathcal T=\sum t_\ell$, $\mathcal D=\sum d_\ell$,
$\mathcal V=\sum v_\ell$. Establish the same inequalities uniformly for
$T_M=\E\min(\mathcal T,M)$, $D_M=\E\min(\mathcal D,M)$,
$V_M=\E\min(\mathcal V,M)$ for every $M>0$. The pathwise inequality
$\mathcal D\le\mathcal T$ gives $D_M\le T_M$, and monotone convergence then
proves finiteness and the endpoint. These truncated inequalities themselves
are not established here.

\subsection{How the adaptive return identity could enter}
For fixed selected bases, split the discrepancy into the part computed from
true compressed observations $A_\ell\Omega_\ell,A_\ell^{\T}\Psi_\ell$ and
the part computed from the recycled noise. The direct old-noise active
returns are precisely~\eqref{base:eq:active-return} and their transpose
counterparts. Write the sum of their expected squared norms as $F_{\rm act}$.
If the smoothness, localization, and limit steps in Section~\ref{base:sec:stein}
are justified for those coefficients, its identity would yield
\begin{equation}\label{base:eq:active-charge}
 F_{\rm act}\le\frac15D+\mathcal A,
\end{equation}
where $\mathcal A$ is the sum of the actual signed adaptation corrections.
Let $S$ be the sum of the true-data active errors evaluated at the bases
selected from the noisy observations. The direct-sum triangle inequality
gives $T\le2F_{\rm act}+2S$.
For illustration, bounds
\[
 \mathcal A\le\tfrac15D+b_0H,\qquad S\le a_0H
\]
would imply $T\le\tfrac45D+2(a_0+b_0)H$.
Neither displayed bound is proved for the full algorithm. The basis-selection
cost in $S$ is different from the direct regression return, and the inactive
error still needs~\eqref{base:eq:missing-V}.

\subsection{What the quotient route still needs}
The quotient state avoids differentiating the bases merely to describe
continuation. To yield~\eqref{base:eq:missing-T}--\eqref{base:eq:missing-V}, however,
it needs an actual-law estimate for the combined merge--compression map,
not only the frozen compression price of Proposition~\ref{base:prop:frozen-quotient}.
Its observation-space energy also needs a justified comparison to the true
error sectors. Initialization~\eqref{base:eq:init-budget} does not supply either
step. Common coefficients that are invisible to future basis decisions
must remain in reconstruction accounting.

A proof through either route must also resolve rank and inverse-integrability
issues on the states the algorithm actually reaches. An event of small
probability cannot simply be charged by $\F A^2$: when $\OPT(A)$ is zero or
arbitrarily small, such an additive term does not fit the target. The exact-input endpoint is now proved separately in
Theorem~\ref{base:aug:thm:exact}; it does not supply a quantitative
residual-relative bound near the HSS class.

\begin{statusbox}
\textbf{Precisely what remains unproved.}
There is no theorem here giving an absolute subunit return coefficient for
the actual two-stream merge--compression process at every depth, with its
basis-selection and inactive-error costs controlled by the optimal
residual. The exact representation, the first-two-level estimates, the
smooth adaptive identity, and the frozen Gaussian benchmark are not a
substitute for this quantitative closure.
\end{statusbox}

\section{Independent validation and amplification}\label{base:sec:validation}
\subsection{Selecting candidates with eight queries}
The next result does not depend on the HSS structure. It applies to any
finite collection of candidate matrices computed before the validation
sketch is drawn.

\begin{lemma}[Expected error of Gaussian validation]\label{base:aug:lem:validation}
Fix $A$ and candidates $B_1,\ldots,B_r$ independently of a standard Gaussian
$G\in\R^{N\times q}$, where $q>4$. Query $AG$ and choose
\begin{equation}\label{base:aug:eq:selection}
 \widehat j\in\arg\min_{1\le j\le r}
       \widehat F_j,
 \qquad \widehat F_j=q^{-1}\F{AG-B_jG}^2.
\end{equation}
Let $F_j=\F{A-B_j}^2$. Conditional on all the candidates,
\begin{equation}\label{base:aug:eq:val-bound}
 \E_G F_{\widehat j}\le c_{r,q}\min_jF_j,
 \qquad
 c_{r,q}=1+(r-1)\sqrt{\frac{q(q+2)}{(q-2)(q-4)}}.
\end{equation}
For $q=8$, one may use $c_{2,8}<3$, $c_{3,8}<5$, and $c_{4,8}<7$.
\end{lemma}
\begin{proof}
If any $F_j=0$, its score is zero, and every candidate with a nonzero
residual has a positive score almost surely. Hence the selected error is
zero. Otherwise put $T_j=\widehat F_j/F_j>0$. The singular value
decomposition of the fixed residual gives
\[
 T_j\overset{d}=\sum_a w_aX_a,\qquad
 w_a\ge0,\quad\sum_aw_a=1,\quad X_a\sim\chi_q^2/q,
\]
with independent $X_a$ within this representation. It follows that
\begin{equation}\label{base:aug:eq:validation-moments}
 \E T_j^2=1+\frac2q\sum_aw_a^2\le1+\frac2q,
 \qquad
 \E T_j^{-2}\le\frac{q^2}{(q-2)(q-4)}.
\end{equation}
The inverse inequality is Jensen's inequality for $x\mapsto x^{-2}$,
followed by the inverse second moment of a chi-square variable. Both
claims are independent of the residual's dimension and stable rank.

Fix an index $j_*$ minimizing the true errors. On the event that $j$ is
selected, score minimality gives
$F_j\le F_{j_*}T_{j_*}/T_j$. Hence pathwise
\[
 F_{\widehat j}\le F_{j_*}
      \left(1+\sum_{j\ne j_*}\frac{T_{j_*}}{T_j}\right).
\]
For each summand use Cauchy--Schwarz and
\eqref{base:aug:eq:validation-moments}. The scores may be correlated because they
use the same $G$; this proof does not assume otherwise. This proves
\eqref{base:aug:eq:val-bound}.
\end{proof}

Random candidate matrices are permitted by conditioning on them first.
The validation sketch must be independent of every candidate. It is not a
training sketch for their bases, discrepancies, or final cores. Computing
$B_jG$ uses the explicit HSS factors and adds no queries to $A$.

\subsection{Amplification of fractional moments and weak tails}
Let $\widehat B^{(1)},\ldots,\widehat B^{(r)}$ be independent raw
reconstructions of the same fixed $A$, including independent initial
sketches for every run. Let $F_1,\ldots,F_r$ be their squared errors.
After the candidates are complete, perform the independent validation in
Lemma~\ref{base:aug:lem:validation}.

\begin{theorem}[Fractional-moment amplification]\label{base:aug:thm:fractional}
Suppose $0<p\le1$ and a raw run satisfies $\E F^p<\infty$. For any integer
$r\ge1/p$, the wrapper satisfies
\begin{equation}\label{base:aug:eq:fractional-result}
 \E F_{\widehat j}\le c_{r,q}(\E F^p)^{1/p}.
\end{equation}
For the raw HSS method and $q=8$, its query count is $18rk+8$ and its output
still belongs to $\HSS(L,k)$.
\end{theorem}
\begin{proof}
For nonnegative numbers, their minimum is at most their geometric mean.
Independence therefore gives
\[
 \E\min_jF_j
 \le\E\prod_{j=1}^r F_j^{1/r}
 =(\E F^{1/r})^r
 \le(\E F^p)^{1/p}.
\]
The last inequality is monotonicity of $L_p$ means, equivalently Jensen's
inequality, since $1/r\le p$. Now condition on the candidates and apply
Lemma~\ref{base:aug:lem:validation}. This proof does not assume that $\E F$ exists.
The query count follows from $r$ raw runs and $q$ forward validation
queries. Selection, rather than averaging matrices, preserves HSS rank $k$.
\end{proof}

\begin{corollary}[Two particularly useful sufficient estimates]
If the raw algorithm satisfies
\[
 \E\F{A-\widehat B}\le C\sqrt{L\OPT(A)},
\]
uniformly over inputs, two runs and eight validation queries give
\[
 \#\text{queries}\le44k,\qquad
 \E\F{A-\widehat B_{\rm sel}}^2\le3C^2L\OPT(A).
\]
If instead one can prove the weaker fractional-moment estimate
\[
 \E\bigl(\F{A-\widehat B}^{1/2}\bigr)
       \le C\bigl(L\OPT(A)\bigr)^{1/4},
\]
four runs and eight validation queries give at most $80k$ queries and
expected squared error at most $7C^4L\OPT(A)$.
Both antecedents remain unproved for the general reuse algorithm.
\end{corollary}

\paragraph{The weak-tail hypothesis.}
For one raw run, write $F=\F{A-\widehat B}^2$. A different sufficient
target is the existence of an absolute $K\ge1$ such that, for every fixed
input with $\OPT(A)>0$,
\begin{equation}\label{base:aug:eq:weak-target}
 \Prb\{F>tL\OPT(A)\}\le Kt^{-1/2}\qquad(t\ge1).
\end{equation}
This hypothesis is not proved for the actual HSS recursion. The following
theorem establishes the implication, not its antecedent.

\begin{theorem}[Weak-tail amplification]\label{base:aug:thm:tail}
Suppose the raw reconstruction satisfies~\eqref{base:aug:eq:weak-target} for an
absolute $K\ge1$ and every fixed input with $\OPT(A)>0$. Use three
independent runs and $q=8$ validation vectors. Then
\[
 \#\text{queries}=54k+8\le62k,\qquad
 \E\F{A-\widehat B_{\rm sel}}^2\le15K^2L\OPT(A).
\]
For $\OPT(A)=0$, the same conclusion holds for the canonical implementation
by Theorem~\ref{base:aug:thm:exact}.
\end{theorem}
\begin{proof}
Write $M=L\OPT(A)>0$. Independence gives, for $t\ge1$,
\[
 \Prb\{\min_{j\le3}F_j>tM\}
 \le\min\{1,K^3t^{-3/2}\}.
\]
Integrate the survival probability and split at $K^2\ge1$:
\[
 \E\min_{j\le3}F_j
 \le M\left[K^2+\int_{K^2}^{\infty}K^3t^{-3/2}\,dt\right]
 =3K^2M.
\]
Use $c_{3,8}<5$ in Lemma~\ref{base:aug:lem:validation} to finish. Tonelli's theorem
justifies the survival integral without presupposing integrability; its
finite upper bound proves integrability of the selected error.
\end{proof}

The weak-tail theorem allows the raw expected squared error to be infinite.
It therefore does not imply, and does not need, the subunit
expected-return inequality for the original single-run process.

\subsection{Why this can circumvent an inverse-moment obstruction}
For a scalar $Z\sim N(0,1)$, the loss $F=Z^{-2}$ satisfies
\[
 \E F=\infty,\qquad
 \Prb\{F>t\}=\Prb\{|Z|<t^{-1/2}\}
 \le\sqrt{2/\pi}\,t^{-1/2}.
\]
Three independent copies have an integrable minimum, since its survival
probability decays as $t^{-3/2}$. Two copies still have a nonintegrable
minimum: their survival probability is asymptotic to $(2/\pi)t^{-1}$.
This example explains why a second-moment failure for a raw inverse need
not rule out a constant-query-multiple wrapper. It is a scalar diagnostic,
not an identification of the actual HSS error distribution.

\subsection{What a fixed-probability success theorem would not prove}
A guarantee $\Prb\{F\le CL\OPT(A)\}\ge0.9$ alone is insufficient.
For example, a scalar loss equal to $M$ with probability $0.9$ and $TM$
otherwise meets such a guarantee for every $T$. The expected minimum of
$r$ independent copies is $M[1+(T-1)0.1^r]$, unbounded as $T\to\infty$
for every fixed $r$. A residual-relative tail or fractional moment, not
merely a constant-probability event, is needed.

\section{Additional obstructions and their scope}\label{base:sec:additional-obstructions}
\subsection{Fractional moments still require actual-law information}
The following construction prevents an incorrect shortcut from the preceding
square-Gaussian example to a generic all-depth fractional-moment bound.

\begin{proposition}[No universal fractional inverse bound]\label{base:aug:prop:all-fractional}
For a standard Gaussian $G\in\R^{d\times s}$ with $d>s$, there is a
measurable unit vector $w(G)$ such that $w^{\T}G$ has row rank one almost
surely, but
\[
 \E\N{(w^{\T}G)^\dagger}^{2p}=\infty
 \qquad\text{for every }p>0.
\]
\end{proposition}
\begin{proof}
Choose a measurable unit vector $u_0\in\ker G^{\T}$ and a unit left
singular vector $u_1$ for $\sigma_{\max}(G)>0$. They are orthogonal. Put
\[
 R=\F G,\qquad \varepsilon=e^{-R^4},\qquad
 w=\sqrt{1-\varepsilon^2}\,u_0+\varepsilon u_1.
\]
Then $\N{w^{\T}G}=\varepsilon\sigma_{\max}(G)>0$ almost surely. Since
$\sigma_{\max}(G)\le R$,
\[
 \N{(w^{\T}G)^\dagger}^{2p}
 \ge e^{2pR^4}R^{-2p}.
\]
The chi-distributed radius $R$ has density proportional to
$R^{ds-1}e^{-R^2/2}$. The displayed lower bound is nonintegrable at
infinity for every $p>0$.
\end{proof}

This is not the SVD-based HSS rule and does not refute its desired
weak-tail estimate. It shows why the algorithm's actual dependence cannot
be discarded even when aiming only for fractional moments.

\subsection{Uniform sketch fitting cannot supply a free certificate}
Suppose the two initial design matrices have fewer than $N$ columns.
For every realization choose nonzero
$u\in\ker\Psi_0^{\T}$ and $v\in\ker\Omega_0^{\T}$. Then
\[
 M=uv^{\T}\in\HSS(L,1),\qquad
 M\Omega_0=0,\qquad M^{\T}\Psi_0=0.
\]
Thus the two-sided sketch cannot satisfy a positive uniform lower norm
bound over the entire HSS class, or over its secants. Simply minimizing a
global training residual cannot be justified by that uniform embedding
claim. The constructed $M$ depends on the sketches: this is \emph{not} a
lower bound against recovery of each fixed input, nor does it rule out a
carefully regularized or rank-selecting global reconstruction method.

\subsection{Two other changes do not resolve the proof gap}
\textbf{Orthogonal column mixing.} For a common orthogonal $s\times s$
matrix $R$, replacing $(Y,\Omega)$ by $(YR,\Omega R)$ leaves
$Y\Omega^\dagger$ and the innovation Gram unchanged. The basis and
discrepancy therefore remain unchanged, and the outgoing forward stream
is merely right-multiplied by $R$. The same holds in the transpose stream.
Applying such common random rotations at successive levels does not
refresh the estimator's effective randomness.

\textbf{Positive ridge in the discrepancy.} Retain the nullification-based
bases but replace the regression inverse by
$\Omega^{\T}(\Omega\Omega^{\T}+\lambda I)^{-1}$ with $\lambda>0$.
Already for $A=I_N\in\HSS(L,k)$, the zero innovations give the same
canonical $U=V$ at the leaves. With $P=UU^{\T}$ and a discarded frame
$W$, the first local discarded--discarded error is
\[
 W^{\T}(I-D_\lambda)W
 =\lambda W^{\T}(\Omega\Omega^{\T}+\lambda I)^{-1}W\succ0.
\]
It cannot be repaired by later retained-space updates. Consequently this
unconditionally positive-ridge modification fails the zero-optimum
requirement. This argument does not exclude a more sophisticated
regularization that vanishes appropriately on exact inputs.


\endgroup

% ===== PROBES =====
\begingroup
\let\R\relax
\newcommand{\R}{\mathbb R}
\let\E\relax
\newcommand{\E}{\mathbb E}
\let\Prb\relax
\newcommand{\Prb}{\mathbb P}
\let\T\relax
\newcommand{\T}{\mathsf T}
\let\F\relax
\newcommand{\F}[1]{\left\lVert#1\right\rVert_{F}}
\let\N\relax
\newcommand{\N}[1]{\left\lVert#1\right\rVert}
\let\ip\relax
\newcommand{\ip}[2]{\left\langle#1,#2\right\rangle_F}
\let\tr\relax
\newcommand{\tr}{\operatorname{tr}}
\let\rank\relax
\newcommand{\rank}{\operatorname{rank}}
\let\diag\relax
\newcommand{\diag}{\operatorname{blockdiag}}
\let\range\relax
\newcommand{\range}{\operatorname{range}}
\let\OPT\relax
\newcommand{\OPT}{\operatorname{OPT}}
\let\HSS\relax
\newcommand{\HSS}{\operatorname{HSS}}
\let\diver\relax
\newcommand{\diver}{\operatorname{div}}
\let\Var\relax
\newcommand{\Var}{\operatorname{Var}}
\let\UC\relax
\newcommand{\UC}{\mathcal U}
\let\VC\relax
\newcommand{\VC}{\mathcal V}
\let\err\relax
\newcommand{\err}{\mathcal E}
\let\proofstep\relax
\newcommand{\proofstep}[1]{\par\addvspace{.45\baselineskip}\Needspace{3\baselineskip}\textbf{#1}\par\nobreak}
\let\op\relax
\newcommand{\op}[1]{\left\lVert#1\right\rVert_{\mathrm{op}}}
\let\Sn\relax
\newcommand{\Sn}[2]{\left\lVert#1\right\rVert_{S_{#2}}}
\let\Cov\relax
\newcommand{\Cov}{\mathsf C}
\let\Loss\relax
\newcommand{\Loss}{\mathsf L}
\let\Fil\relax
\newcommand{\Fil}{\mathcal F}
\let\Id\relax
\newcommand{\Id}{I}
\let\supp\relax
\newcommand{\supp}{\operatorname{supp}}
\let\vecop\relax
\newcommand{\vecop}{\operatorname{vec}}
\let\Gammafun\relax
\newcommand{\Gammafun}{\operatorname{Gamma}}
\chapter{Orthogonal probing and exact propagation}\label{ch:probes}
\begin{scopebox}
This begins the independent-basis refit chain. Intermediate geometry-dependent and logarithmic certificates are retained because they supply the state identities and counterexamples used later. The final uniform refit theorem is in Chapter~\ref{ch:uniform}.
\end{scopebox}
\section{Orthogonal probing after basis learning}\label{base:sec:ortho-refit}
\subsection{Freeze the geometry and separate the two errors}
Use the convention $N=2^{L+1}k$, with $2k$ indices in each leaf.
The class $\HSS(L,k)$ imposes rank at most $k$ on every nonroot exterior
row and column block, and
\[
 \OPT(A)=\min_{B\in\HSS(L,k)}\F{A-B}^2.
\]
Run only the basis-learning stage of Section~\ref{base:sec:algorithm}: $8k$ forward and $8k$ transpose
Gaussian queries, followed by its $L$ recycled compressions. Retain the
block-diagonal isometries $U_\ell,V_\ell$, $0\le\ell<L$, and discard the
estimated discrepancies. No final-core queries are made at this stage.
Each local block of $U_\ell,V_\ell$ has size $2k\times k$.

Condition on these bases. They define a linear space $\mathcal S$ of
telescoping matrices
\begin{equation}\label{base:orf:eq:space}
 B_\ell=D_\ell+U_\ell B_{\ell+1}V_\ell^{\T},\qquad
 U_{\ell,I}^{\T}D_{\ell,I}V_{\ell,I}=0,
\end{equation}
where each $D_\ell$ is block diagonal and $B_L$ is an arbitrary
$2k\times2k$ core. Then $\mathcal S\subseteq\HSS(L,k)$.
The gauge in~\eqref{base:orf:eq:space} makes the lifted discrepancy spaces and core
space mutually orthogonal. In particular,
\begin{equation}\label{base:orf:eq:dimension}
 \dim\mathcal S
 =\sum_{\ell=0}^{L-1}\frac{N/2^\ell}{2k}(3k^2)+4k^2
 =3Nk-2k^2.
\end{equation}

Write $P_{\mathcal S}$ for Frobenius orthogonal projection and set
$d_{\mathcal S}(A)^2=\F{A-P_{\mathcal S}A}^2$. For every output
$\widehat B\in\mathcal S$,
\begin{equation}\label{base:orf:eq:split}
 \F{A-\widehat B}^2
 =d_{\mathcal S}(A)^2+\F{P_{\mathcal S}A-\widehat B}^2.
\end{equation}
Thus a coefficient refit cannot repair a poor learned space. Conversely,
coefficient instability is not evidence that the learned space is poor.

For clarity, the exact projection can be described without solving an
optimization problem: set $A_0=A$, $A_{\ell+1}=U_\ell^{\T}A_\ell V_\ell$,
use discrepancies
\[
 D_{\ell,I}^{*}=A_\ell[I,I]
   -P_{\ell,I}A_\ell[I,I]Q_{\ell,I},\qquad
 P_{\ell,I}=U_{\ell,I}U_{\ell,I}^{\T},\quad
 Q_{\ell,I}=V_{\ell,I}V_{\ell,I}^{\T},
\]
and use core $A_L$. This is a proof and diagnostic formula; the refitting
algorithm does not have dense access to $A$ or compute $P_{\mathcal S}A$.

\subsection{Construct the probes from the top down}
Fix a width $q\ge2k$. We construct a forward probe matrix $G\in\R^{N\times q}$
using the fixed column bases. Its compressed designs will satisfy
\begin{equation}\label{base:orf:eq:frames}
 G_{\ell,I}G_{\ell,I}^{\T}=I_{2k}\quad(0\le\ell<L),\qquad
 G_LG_L^{\T}=I_{2k}.
\end{equation}
Consequently every local pseudoinverse is just the transpose.

Start with any $2k\times q$ row-orthonormal root matrix. Its two $k$-row
blocks prescribe the retained probes $R_I$ of the root's two children.
At a nonroot node $I$, suppose its retained probe $R_I\in\R^{k\times q}$
is already prescribed and has orthonormal rows. Choose a row-orthonormal
$W_I\in\R^{k\times q}$ in its orthogonal complement, and set
\begin{equation}\label{base:orf:eq:lifting}
 G_I^{\rm loc}=V_I R_I+V_{I,\perp}W_I.
\end{equation}
Here $[V_I\ V_{I,\perp}]$ is an orthogonal $2k\times2k$ matrix.
The two $k$-row halves of $G_I^{\rm loc}$ prescribe the retained probes of
the two children; at an original leaf, its $2k$ rows are the corresponding
rows of $G$. This construction requires no matrix--vector queries.

Equation~\eqref{base:orf:eq:frames} follows by induction, as does
$V_I^{\T}G_I^{\rm loc}=R_I$. Construct $H$ in the same way from the row
bases $U_I$. Only then query $Y_0=AG$ and $Z_0=A^{\T}H$, costing $2q$ queries.

\begin{lemma}[Isotropic random completion]\label{base:orf:lem:isotropy}
Choose each $W_I$ uniformly from the row-orthonormal frames in the prescribed
orthogonal complement, independently conditional on previously constructed
ancestors. Then, conditional on the fixed bases,
\begin{equation}\label{base:orf:eq:isotropy}
 \E GG^{\T}=I_N,\qquad \E HH^{\T}=I_N.
\end{equation}
Every realization also has $\F G^2=\F H^2=N$.
\end{lemma}
\begin{proof}
In the orthonormal multiresolution coordinates of the fixed column bases,
the rows of $G$ consist of the root frame and the discarded frames $W_I$.
Each block has identity row Gram matrix. Each nonroot block has conditional
mean zero. Cross moments between an ancestor and descendant therefore
vanish; those between different branches vanish by conditional independence.
The expected Gram matrix in these coordinates is identity. Transforming
back gives the first identity; the row-basis construction gives the second.
The deterministic norm statement follows from the leaf identities
in~\eqref{base:orf:eq:frames}.
\end{proof}

\textbf{A $4k$-column refinement.}
For $q\ge4k$, process siblings $I,J$ together. Their prescribed retained
frames satisfy $[R_I;R_J][R_I;R_J]^{\T}=I_{2k}$. Choose
$[W_I;W_J]$ as a uniformly random $2k$-row orthonormal frame in the
orthogonal complement of $[R_I;R_J]$. Thus
\begin{equation}\label{base:orf:eq:siblings}
 W_I R_I^{\T}=W_I R_J^{\T}=W_J R_I^{\T}=W_J R_J^{\T}=0,
 \qquad W_IW_J^{\T}=0.
\end{equation}
The two expanded child frames in~\eqref{base:orf:eq:lifting} are mutually
orthogonal. Lemma~\ref{base:orf:lem:isotropy} still applies: the sibling cross
moment is now zero in every realization, and the joint discarded frame
has conditional mean zero.

\subsection{An inverse-free sequential refit}
Use the fixed bases, not new singular-vector decisions. At level $\ell$
and each $2k$-row block $I$, form
\begin{align}
 C_{Y,I}&=Y_{\ell,I}G_{\ell,I}^{\T},&
 C_{Z,I}&=\bigl(Z_{\ell,I}H_{\ell,I}^{\T}\bigr)^{\T},\nonumber\\
 D_{\ell,I}&=(I-P_{\ell,I})C_{Y,I}
       +P_{\ell,I}C_{Z,I}(I-Q_{\ell,I}).\label{base:orf:eq:refit}
\end{align}
Assemble $D_\ell$ block diagonally and update
\begin{align}
 Y_{\ell+1}&=U_\ell^{\T}(Y_\ell-D_\ell G_\ell),&
 G_{\ell+1}&=V_\ell^{\T}G_\ell,\nonumber\\
 Z_{\ell+1}&=V_\ell^{\T}(Z_\ell-D_\ell^{\T}H_\ell),&
 H_{\ell+1}&=U_\ell^{\T}H_\ell.\label{base:orf:eq:transport}
\end{align}
Merge consecutive retained blocks as usual. At the root use
\begin{equation}\label{base:orf:eq:root}
 \widehat C=\tfrac12\left(Y_LG_L^{\T}
                  +(Z_LH_L^{\T})^{\T}\right),
\end{equation}
and reconstruct by~\eqref{base:orf:eq:space}. There are no additional core queries.

\begin{proposition}[Exactness on the fixed space]\label{base:orf:prop:exact}
For every fixed choice of bases and every probe realization satisfying
\eqref{base:orf:eq:frames}, this refit recovers every $B\in\mathcal S$ exactly.
\end{proposition}
\begin{proof}
At the first level, every exterior interaction of $B$ is killed on the
left by $I-P_I$ and on the right by $I-Q_I$. Since
$G_IG_I^{\T}=H_IH_I^{\T}=I$, equation~\eqref{base:orf:eq:refit} therefore recovers
exactly $B[I,I]-P_IB[I,I]Q_I$. The recycled observations are exact
observations of $U^{\T}BV$ with its fixed bases. Induction proves the
statement at every level, and either summand in~\eqref{base:orf:eq:root} equals
the true core.
\end{proof}

\begin{theorem}[A uniform, but insufficient, all-depth bound]\label{base:orf:thm:bound}
Conditional on any fixed learned bases, use the isotropic completion rule
of Lemma~\ref{base:orf:lem:isotropy}, either with $q=2k$ or with the sibling refinement
at $q=4k$. Then
\begin{equation}\label{base:orf:eq:mainbound}
 \E_{\rm refit}\F{A-\widehat B}^2
 \le(3\cdot2^L-1)\,d_{\mathcal S}(A)^2.
\end{equation}
For every realization, irrespective of the distribution of the probes,
\begin{equation}\label{base:orf:eq:pathbound}
 \F{A-\widehat B}^2
 \le\bigl[1+N(3\cdot2^L-2)\bigr]d_{\mathcal S}(A)^2.
\end{equation}
The same two bounds hold for joint least-squares fitting of all coefficients
in $\mathcal S$ from $(AG,A^{\T}H)$.
\end{theorem}
\begin{proof}
First regard the sequential construction as a linear map $\mathcal R$ from
an arbitrary pair of data matrices $(Y_0,Z_0)$ to a matrix in $\mathcal S$.
Let $s_\ell=\F{Y_\ell}^2+\F{Z_\ell}^2$. At a given level, define
\[
 a^2=\F{(I-P)Y}^2+\F{(I-Q)Z}^2,\qquad
 r^2=\F{U^{\T}Y}^2+\F{V^{\T}Z}^2,
\]
so $a^2+r^2=s_\ell$. All these norms sum over the level's blocks.
Right multiplication by each local probe transpose is a contraction.
The two row sectors of $D$ in~\eqref{base:orf:eq:refit} are orthogonal, giving
\begin{equation}\label{base:orf:eq:Denergy}
 \F D^2\le a^2\le s_\ell.
\end{equation}
Only its retained-row/discarded-column sector enters $U^{\T}DG$; only
its discarded-row/retained-column sector enters $V^{\T}D^{\T}H$.
Using the local row-isometry identities,
\[
 \F{U^{\T}DG}^2+\F{V^{\T}D^{\T}H}^2\le a^2.
\]
The triangle inequality in the direct sum of the two measurement spaces
therefore yields
\begin{equation}\label{base:orf:eq:two}
 s_{\ell+1}\le(r+a)^2\le2s_\ell.
\end{equation}
The root estimate has squared norm at most $s_L/2$. Orthogonality of the
gauge decomposition and~\eqref{base:orf:eq:Denergy}--\eqref{base:orf:eq:two} give
\begin{equation}\label{base:orf:eq:Rbound}
 \F{\mathcal R(Y_0,Z_0)}^2
 \le\left(\sum_{\ell=0}^{L-1}2^\ell+2^{L-1}\right)s_0
 =c_Ls_0,\qquad c_L=3\cdot2^{L-1}-1.
\end{equation}

Put $E=A-P_{\mathcal S}A$. Exactness and linearity imply
$\widehat B=P_{\mathcal S}A+\mathcal R(EG,E^{\T}H)$. By
\eqref{base:orf:eq:split} and~\eqref{base:orf:eq:Rbound},
\[
 \F{A-\widehat B}^2\le\F E^2
       +c_L\bigl(\F{EG}^2+\F{E^{\T}H}^2\bigr).
\]
Conditional isotropy gives expected measurement energy $2\F E^2$,
proving~\eqref{base:orf:eq:mainbound}. The deterministic bounds
$\F G^2=\F H^2=N$ give~\eqref{base:orf:eq:pathbound}.

Finally define $\mathcal X(M)=(MG,M^{\T}H)$ on $\mathcal S$.
Proposition~\ref{base:orf:prop:exact} and~\eqref{base:orf:eq:Rbound} supply a left inverse
of norm at most $\sqrt{c_L}$. Thus $\mathcal X$ is injective and its
Moore--Penrose left inverse has norm at most $\sqrt{c_L}$. Repeating the
same argument proves both bounds for joint least squares.
\end{proof}

\begin{corollary}[Query ledger and exact-input endpoint]\label{base:orf:cor:queries}
The complete construction uses $20k$ queries for $q=2k$, or $24k$ for
$q=4k$: $16k$ for learning and $2q$ for refitting. It returns a matrix in
$\HSS(L,k)$ and has finite moments of every positive order for every
fixed input and finite depth. With the canonical rank-deficient completion
rule of Theorem~\ref{base:aug:thm:exact}, it recovers each fixed exact HSS
input almost surely.
\end{corollary}
\begin{proof}
The query count follows from the specified acquisitions. Since
$d_{\mathcal S}(A)\le\F A$, the deterministic bound~\eqref{base:orf:eq:pathbound}
is uniform over all learned bases and all refit randomness. This proves
the moment statement even if the discarded training coefficients have
infinite moments. On an exact HSS input, Theorem~\ref{base:aug:thm:exact}
shows that the learned space contains $A$ almost surely; apply
Proposition~\ref{base:orf:prop:exact}.
\end{proof}

\subsection{What the wider probes cancel}
The $q=4k$ construction provides a useful algebraic improvement beyond
the coarse bound in Theorem~\ref{base:orf:thm:bound}.

\begin{lemma}[No immediate coefficient feedback]\label{base:orf:lem:feedback}
With fixed bases and the sibling condition~\eqref{base:orf:eq:siblings}, perturb
the discrepancies at one level by arbitrary local gauge matrices
$\delta D_I$ satisfying $U_I^{\T}\delta D_IV_I=0$. Their induced changes
in the next level's two regression coefficients $C_Y,C_Z$ are zero.
This assertion concerns the immediate next level only.
\end{lemma}
\begin{proof}
In retained/discarded coordinates write
$U_I^{\T}\delta D_I=\delta D_{21,I}V_{I,\perp}^{\T}$.
The induced forward data perturbation after compression is
$-\delta D_{21,I}W_I$. The next local design is $[R_I;R_J]$ for the
sibling pair, and
\[
 W_I[R_I;R_J]^{\T}=0.
\]
It therefore contributes zero to the next forward regression. The
transpose statement follows from the corresponding row-basis construction
and the $12$ sector of $\delta D_I$. The paired-block argument applies
to all nodes at that level.
\end{proof}

This cancels precisely one feedback mechanism rather than estimating an
adaptive inverse. The perturbation remains in the measurement residual
and may become visible at a later ancestor. Later sections quantify this delayed return. Immediate cancellation alone does not justify an all-depth estimate.

\subsection{What is still needed for the requested target}
This first inverse-free construction avoids the adaptive inverse obstruction in its
final coefficient stage. It also eliminates immediate coefficient feedback
when $q=4k$. By itself it does not establish either of the two estimates needed by a
straightforward separation argument:
\begin{align*}
 \E_{\rm training}d_{\mathcal S}(A)^2&\le C_b L\OPT(A),\\
 \E_{\rm refit}\F{A-\widehat B}^2&\le C_r d_{\mathcal S}(A)^2
       \quad\text{uniformly over fixed nested bases.}
\end{align*}
These would imply the target with $24k$ queries. The first is an
all-depth basis-learning problem; the second is a coefficient-transport
problem with bounded local operators and no inverse tails. Alternatively,
a combined analysis could avoid this separation. Merely proving an
$O(L)$ factor in \emph{both} displays would give $O(L^2)$, not the target.



\section{What the \texorpdfstring{F$^+$}{F+} idea transfers, and what it does not}\label{base:new:fp}
The motivating model is a centered rank-one covariance sum. Let $Y_i=y_iy_i^\T\succeq0$ be rank one in dimension $r$, $\E Y_i=I$, and $X_i=Y_i-I$. The original setting takes independent identically distributed samples, for example $y_i=U^\T(g_i^{(1)}\otimes g_i^{(2)})$ on a fixed subspace. The proof below needs conditional versions of its moment hypotheses, not literal independence.

\begin{theorem}[Conditional finite-profile comparison]\label{base:new:fp-thm}
Suppose $Y_i$ are adapted, $\E[X_i\mid\Fil_{i-1}]=0$, and deterministic budgets satisfy
\[
 \E[X_i^2\mid\Fil_{i-1}]\preceq b_2I,\qquad
 \E[(Y_i+I)^j\mid\Fil_{i-1}]\preceq b_jI\quad(3\le j\le2Q),
 \qquad b_2>0.
\]
Put $S_n=\sum_{i\le n}X_i$, $M_{n,p}=r^{-1}\E\tr S_n^{2p}$ and $M_{n,0}=1$. Define
\[
 \mathcal I_{2a}(u)=u_a,\quad \mathcal I_{2a+1}(u)=\sqrt{u_au_{a+1}},\quad
 \Phi_p(u)=\sum_{j=2}^{2p}\binom{2p}j b_j\mathcal I_{2p-j}(u).
\]
Then $M_{0,p}=0$ for $p\ge1$ and
\begin{equation}\label{base:new:fp-rec}
 M_{n,p}\le M_{n-1,p}+\Phi_p(M_{n-1}).
\end{equation}
Let $\kappa_a=(2a-1)!!$, $\omega_{2a}=\kappa_a$, $\omega_{2a+1}=\sqrt{\kappa_a\kappa_{a+1}}$, and
\[
 C_{p,j}=\frac{(2p)!\,\omega_{2p-j}}{2p(2p-j)!\kappa_p},\qquad
 \phi_Q(s)=\sum_{j=2}^{2Q}\frac{C_{Q,j}b_j}{j!}s^{1-j}.
\]
The increasing inverse $T_Q(t)$ of $a\mapsto\int_0^a\phi_Q(s)^{-1}ds$ satisfies
\begin{equation}\label{base:new:fp-int}
 (\E\tr S_n^{2Q})^{1/(2Q)}\le r^{1/(2Q)}\kappa_Q^{1/(2Q)}T_Q(n).
\end{equation}
The budgets need not be the moments of a realizable scalar law.
\end{theorem}
\begin{proof}
Condition before expanding the noncommutative words in $(A+X)^{2p}$, with $A=S_{n-1}$. One insertion vanishes. For two insertions, diagonalizing $A$ and using weighted arithmetic--geometric mean bounds the trace by $\tr(|A|^{2p-2}X^2)$. For three or more insertions, expand $X=Y-I$. Every word with $\ell\ge1$ copies of $Y=\tau uu^\T$ factors cyclically as $\tau^\ell\prod_j u^\T A^{a_j}u$. Spectral H\"older bounds its absolute value by $\tr(|A|^{\sum a_j}Y^\ell)$. Summing the selections gives $(Y+I)^j$. Odd absolute moments are bounded by Cauchy--Schwarz under normalized spectral trace, proving \eqref{base:new:fp-rec}.

Let $H_0=1$, $H_p(0)=0$, $H_p'=\Phi_p(H)$. This triangular system is defined by successive integration; all components are nondecreasing. Hence $H(n+1)\ge H(n)+\Phi(H(n))$ and $M_{n,p}\le H_p(n)$. Put $B_p(t)=\kappa_pT_Q(t)^{2p}$. Direct factorial cancellation gives
\[
 C_{p,2\ell}=2^{\ell-1}(p-1)!/(p-\ell)!,\quad
 C_{p,2\ell+1}=\frac{2^\ell(p-1)!}{(p-\ell-1)!\sqrt{2(p-\ell)-1}}.
\]
Thus $C_{p,j}\le C_{p+1,j}$ and $C_{p,j}\le(2p)^{(j-2)/2}$. It follows that $B_p'\ge\Phi_p(B)$ for $p\le Q$. Induction on $p$, integration, and $B_p(0)=H_p(0)=0$ prove the claim. Positivity of $\phi_Q$, its behavior $b_2/(2s)$ at infinity, and integrability of $1/\phi_Q$ at zero justify its inverse and the zero initial value.
\end{proof}
The Hermite interpretation is exact at a terminal scalar Gaussian envelope: multiplication by a standard Gaussian in the normalized Hermite basis has off-diagonal coefficients $\sqrt{j+1}$ and $\sqrt j$. Wick pairings give $\kappa_Q$. This classical multiplication representation is not itself the reason the HSS hierarchy is stable. HSS does not append a fresh centered rank-one summand at every level. Its observations, bases and regressions reuse randomness.

The lesson is to retain centering, lower-order budgets and the signed state before scalarization. Zero initialization alone does not prevent exponential growth: $e_{n+1}=2e_n+b$, $e_0=0$, still gives $(2^n-1)b$. Nor may the conditional hypotheses of Theorem~\ref{base:new:fp-thm} be inferred from unconditional isotropy.

\section{The Wishart route: a correct interface, not a freshness argument}\label{base:new:wishart}
Wishart identities can enter at an intermediate stage if an actual conditional Gaussian law is proved. Here is the precise regression module used in screening that route.
\begin{proposition}[Conditionally Gaussian regression]
Conditional on $\mathcal A$, suppose $m$ independent rows $(x_i,z_i)$ are jointly Gaussian with covariance
\[
 \begin{pmatrix}K_{11}&K_{12}\\K_{21}&K_{22}\end{pmatrix},\qquad K_{11}\succ0.
\]
Let $X$ stack $x_i^\T$, let $Z$ stack $z_i^\T$, put $B=K_{11}^{-1}K_{12}$ and $S=K_{22}-K_{21}K_{11}^{-1}K_{12}$. For $m>k+1$,
\[
 \E[\F{X^\dagger Z-B}^2\mid\mathcal A]
 =\frac{\tr(K_{11}^{-1})\tr S}{m-k-1}.
\]
\end{proposition}
\begin{proof}
Gaussian regression gives $Z=XB+E$, with $E$ independent of $X$ and row covariance $S$. Conditioning on $X$ gives error $\tr S\tr[(X^\T X)^{-1}]$. Write $X=GK_{11}^{1/2}$ and use $\E(G^\T G)^{-1}=I/(m-k-1)$, derived in the Gaussian appendix. Average only after this conditional identity.
\end{proof}
An arbitrary outer average of the product above cannot be factored. In a hidden-covariance model the conditional law can be a mixture of Wisharts, not a Wishart with the mean covariance. A schematic diagnostic is $W\mid K\sim W_m(K)$: the total covariance contains both the conditional sampling covariance and $\operatorname{Cov}(mK)$. Replacing random $K$ by $\E K$ drops the latter term. Known-covariance Schur updates, weighted Gaussian-row comparisons, inverse-Wishart means, and fourth-moment lower-edge estimates address a \emph{verified} transition, not whether a recycled HSS transition has that law.

This route did not close the raw-reuse target. In the orthogonal refit the local inverse is already removed, so a stronger inverse theorem is not the main missing estimate. The relevant remaining object is the signed source--response coupling. Global least-squares refitting was also considered: Theorem~\ref{base:orf:thm:bound} proves injectivity and an exponential-depth bound for it, but no uniform lower singular-value estimate on the whole fixed space was established.

\section{An exact energy balance for compression and merge}\label{base:new:energy}
For row-orthonormal local designs $\Omega,\Psi$, retain the quotient of Section~\ref{base:sec:quotient}. Its natural energy is
\[
 \mathcal V=\F{E_Y}^2+\F{E_Z}^2+\F\Delta^2.
\]
Because $C_\Omega=C_\Psi=0$, $J_\Omega=V_\perp^\T\Omega$ and $J_\Psi=U_\perp^\T\Psi$, the exact compression equations reduce to
\[
 E_Y'=U^\T E_Y+\Delta_{12}V_\perp^\T\Omega,\quad
 E_Z'=V^\T E_Z-\Delta_{21}^\T U_\perp^\T\Psi,\quad
 \Delta'=\Delta_{11}.
\]
The two terms in each innovation are orthogonal. Therefore
\begin{equation}\label{base:new:Vcomp}
 \mathcal V'=\mathcal V-\F{U_\perp^\T E_Y}^2-\F{V_\perp^\T E_Z}^2-\F{\Delta_{22}}^2.
\end{equation}
For two children, let $F_Y=[E_{Y,1};E_{Y,2}]\Omega^\T$ and $F_Z=[E_{Z,1};E_{Z,2}]\Psi^\T$ in the parent coordinates. Orthogonal regression removes their squared norms from innovation energy; the mismatch adds them back. Their diagonal child blocks vanish, so they are orthogonal to the old block-diagonal mismatch. Thus
\begin{equation}\label{base:new:Vmerge}
 \mathcal V_{\rm parent}=\mathcal V_1+\mathcal V_2-2\ip{F_Y}{F_Z^\T}.
\end{equation}
At level $\ell$, write $\mathcal D_\ell\ge0$ for compression dissipation and $\mathcal C_\ell$ for the sum of these merge inner products. Then
\begin{equation}\label{base:new:Vlevel}
 \mathcal V_{\ell+1}=\mathcal V_\ell-\mathcal D_\ell-2\mathcal C_\ell.
\end{equation}
This identity is the first effective F$^+$-style transfer: retain the coupled signed expression until its internal squares cancel.

\paragraph{Why one scalar weight cannot close it pathwise.}
For $\mathcal V_a=\F{E_Y}^2+\F{E_Z}^2+a\F\Delta^2$, compression is nonexpansive on all admissible states only if $a\ge1$. Merge has worst increment $(a-1)(\F{F_Y}^2+\F{F_Z}^2)-2a\ip{F_Y}{F_Z^\T}$; the choice $F_Z^\T=-F_Y$ requires $a\le1/2$. No constant $a$ satisfies both. This is a counterexample to a universal scalar-potential proof, not to the actual random refit.

\subsection{A positive ledger that includes reconstruction}
Define
\[
 a_\ell=\sum_I\left(\F{U_{\ell,I,\perp}^\T Y_{\ell,I}}^2+
 \F{V_{\ell,I,\perp}^\T Z_{\ell,I}}^2\right),\qquad
 \mathscr W=\sum_{\ell<L}a_\ell+\F{Y_L}^2+\F{Z_L}^2.
\]
Each observation direction is charged when discarded, or at the root, not once for every level at which it survives.
\begin{proposition}[Output and signed-feedback interfaces]\label{base:new:ledger-interface}
For arbitrary fixed bases and arbitrary data,
\[
 \F{\mathcal R(Y_0,Z_0)}^2\le\mathscr W,\qquad
 \sum_{\ell<L}\mathcal D_\ell+\mathcal V_L\le2\mathscr W.
\]
For $E\perp\mathcal S$, $\E\mathcal V_0=2\F E^2$. Therefore $\E\mathscr W\le C_W\F E^2$ implies
\[
 \E\F{A-\widehat B}^2\le(1+C_W)d_{\mathcal S}(A)^2,\qquad
 -2\sum_\ell\E\mathcal C_\ell\le(2C_W-2)\F E^2.
\]
\end{proposition}
\begin{proof}
The two row sectors of each discrepancy are orthogonal and each local probe transpose is a contraction. Thus its squared coefficient norm is at most the corresponding discarded observation energy. Lifted discrepancy spaces and root space are orthogonal; the averaged root costs at most half the terminal observation energy. This proves the first inequality.
For the second, in retained/discarded block coordinates the local dissipation is
\[
 \F{U_\perp^\T Y}^2+\F{V_\perp^\T Z}^2-
 \F{(C_Y)_{21}}^2-\F{(C_Z)_{12}}^2
 -2\ip{(C_Y)_{22}}{(C_Z)_{22}}.
\]
Bound the last inner product by the sum of squares and use the local contraction. The terminal quotient energy is at most twice terminal observation energy. Finally
$\mathcal V_I=\F{Y_I}^2+\F{Z_I}^2-2\ip{B_{Y,I}}{B_{Z,I}}$.
The streams are independent and their mean regressions both equal $E[I,I]$. For $E\perp\mathcal S$, each compressed local diagonal block is zero, including the leaves. Isotropy gives the stated initialization. Sum \eqref{base:new:Vlevel}.
\end{proof}
The first compression--merge cycle has $\E\mathcal C_0\ge0$: immediate injected components are killed by sibling orthogonality, leaving forward-only and transpose-only regressions whose matching block means are $U_i^\T E[I_i,I_j]V_j$. Their expected inner product is the square of this deterministic mean. Later registers depend on both streams, so this argument cannot be iterated without additional work.

\section{The first all-depth closures: selectors and absolute path profiles}\label{base:new:early}
The following routes are retained because they explain the later changes of state.

\subsection{Coordinate-retaining bases}
Suppose every local basis retains its first $k$ coordinates. Regard the physical coordinates as $k$-blocks. A block $x\ne0$ is discarded toward $p(x)=x-2^{\nu_2(x)}$. A single input block produces two nonbranching, alternating-stream paths. A forward message $MG_z$ at $x$ emits
\[
 A_r=MG_zG_{p(x)}^\T
\]
and becomes $-A_r^\T H_x$ in the transpose stream at $p(x)$. The transpose rule is the transpose analogue. Later diagonal emissions vanish by sibling orthogonality.

For $q=4k$ a discarded $k$-frame in a $2k$-complement contributes a factor $1/2$ to its conditional second moment. Integrating the finest remaining frame in an ordered path product, then removing an adjacent row-isometry exactly, gives
\[
 \E\F{A_r}^2\le2^{-r}\F B^2.
\]
Repeated occurrences of a frame remain in the ordered product; they are not integrated as independent copies. Independent physical-block rotations and signs make the expected quadratic form block-diagonal in the input blocks, so this one-block calculation suffices. Equivalently, average a unit block as $uv^\T$; the independent forward and transpose products give cancelling factors of $k$.
The edge output of one path plus its root contribution is at most
$\sum_{r<d}2^{-r}+2^{-(d-1)}/4\le2$. Only the initial forward source can emit a nonroot diagonal block, at cost at most one. Combining the two paths in the other sectors gives refit noise at most $9\F E^2$ and hence
\begin{equation}\label{base:new:selector}
 \E\F{A-\widehat B}^2\le10d_{\mathcal S}(A)^2.
\end{equation}
The observation ledger gives $\E\mathscr W\le12\F E^2$, hence cumulative adverse merge correlation at most $22\F E^2$. These constants are not optimal; the sharper geometric results below cover different families.

\subsection{General one-pass feedback, but not iterated independence}
For lifted retained/discarded row bases $u_j,\widetilde u_i$, and column probe frames $W_j^G=\widetilde v_j^TG$, $R_i^G=v_i^TG$, the mixed coefficient arrays satisfy
\[
 a=a_0-K_Gb,\qquad b=b_0-K_Ha,
\]
\[
 (K_Gb)_i=\sum_{j\prec i}(\widetilde u_i^Tu_j)b_jW_j^G(R_i^G)^T.
\]
Only strict ancestors occur; immediate-parent terms vanish. For deterministic $b$,
\[
 \E K_Gb=0,\qquad \E\F{K_Gb}^2\le\rho\F b^2.
\]
To prove this, integrate the finest discarded source frame to kill cross terms and use
$\E[W_j^T W_j\mid\text{coarser frames}]=(k/d)P_j$ together with Bessel's inequality for the lifted spatial overlaps. The same bound holds for inputs independent of the entire relevant stream. Independent $G,H$ give a two-pass factor $\rho^2$ for a deterministic input. The solution arrays depend on both streams, so these inequalities cannot be restarted repeatedly on their actual values.

\subsection{Absolute ascending-path profiles}
Let $p_n^Y(j)$ be the sum of products of operator norms of the alternating column/row ancestor overlaps over all $n$-switch paths from $j$, and define $p_n^Z$ by interchanging orientations. Immediate-parent switches are omitted. Put
\[
 A_j^Y=\sum_{t\ge0}\rho^{t/2}p_{2t}^Y(j),\quad
 B_j^Y=\sum_{t\ge0}\rho^{(t+1)/2}p_{2t+1}^Y(j),\quad
 \Phi_j^Y=(A_j^Y)^2+(B_j^Y)^2.
\]
The sums are finite. An actual forward source may depend on all of $G$; only alternating $H$-returns are used for probabilistic contraction, while $G$ factors are bounded pathwise. Conditional source orthogonality, proved below, then gives
\[
 \E\mathscr W(EG,E^TH)\le4\F E^2+
 2\sum_j\Phi_j^Y\F{\widetilde u_j^TE}^2+
 2\sum_j\Phi_j^Z\F{E\widetilde v_j}^2.
\]
With absolute ancestor sums $s_U,s_V$ and $\chi=\sqrt\rho s_Us_V<1$, this is bounded by $C_W\F E^2$, where
\[
 C_W=4+\frac{2[2+\rho(s_U^2+s_V^2)]}{(1-\sqrt\rho s_Us_V)^2}.
\]
For balanced child contractions $1/\sqrt2$, $s_U,s_V\le1+1/\sqrt2$. Width $18k$ gives $1+C_W<70$ and $2C_W-2<135$, at $36k$ refit queries. This was the first genuinely mixing all-depth route; later storage arguments achieve better constants at width $4k$.

The obstruction is absolute summation. A scalar retained vector can have $h$ future discarded overlaps all of size $1/\sqrt{h+1}$, so the absolute sum grows like $\sqrt h$ even though the squared sum is bounded. Failure of this envelope is not evidence of growth of the actual response.


\endgroup

% ===== SOURCES =====
\begingroup
\let\R\relax
\newcommand{\R}{\mathbb R}
\let\E\relax
\newcommand{\E}{\mathbb E}
\let\Prb\relax
\newcommand{\Prb}{\mathbb P}
\let\T\relax
\newcommand{\T}{\mathsf T}
\let\F\relax
\newcommand{\F}[1]{\left\lVert#1\right\rVert_{F}}
\let\N\relax
\newcommand{\N}[1]{\left\lVert#1\right\rVert}
\let\ip\relax
\newcommand{\ip}[2]{\left\langle#1,#2\right\rangle_F}
\let\tr\relax
\newcommand{\tr}{\operatorname{tr}}
\let\rank\relax
\newcommand{\rank}{\operatorname{rank}}
\let\diag\relax
\newcommand{\diag}{\operatorname{blockdiag}}
\let\range\relax
\newcommand{\range}{\operatorname{range}}
\let\OPT\relax
\newcommand{\OPT}{\operatorname{OPT}}
\let\HSS\relax
\newcommand{\HSS}{\operatorname{HSS}}
\let\diver\relax
\newcommand{\diver}{\operatorname{div}}
\let\Var\relax
\newcommand{\Var}{\operatorname{Var}}
\let\UC\relax
\newcommand{\UC}{\mathcal U}
\let\VC\relax
\newcommand{\VC}{\mathcal V}
\let\err\relax
\newcommand{\err}{\mathcal E}
\let\proofstep\relax
\newcommand{\proofstep}[1]{\par\addvspace{.45\baselineskip}\Needspace{3\baselineskip}\textbf{#1}\par\nobreak}
\let\op\relax
\newcommand{\op}[1]{\left\lVert#1\right\rVert_{\mathrm{op}}}
\let\Sn\relax
\newcommand{\Sn}[2]{\left\lVert#1\right\rVert_{S_{#2}}}
\let\Cov\relax
\newcommand{\Cov}{\mathsf C}
\let\Loss\relax
\newcommand{\Loss}{\mathsf L}
\let\Fil\relax
\newcommand{\Fil}{\mathcal F}
\let\Id\relax
\newcommand{\Id}{I}
\let\supp\relax
\newcommand{\supp}{\operatorname{supp}}
\let\vecop\relax
\newcommand{\vecop}{\operatorname{vec}}
\let\Gammafun\relax
\newcommand{\Gammafun}{\operatorname{Gamma}}
\chapter{Sources, adapted restarts, and conditional amplifiers}\label{ch:sources}
\begin{scopebox}
Conditional amplification, source-weighted stability, and whole-space approximation are different statements. Open fixed-space interactions mentioned in these preliminary reductions are closed by Chapters~\ref{ch:feedback}--\ref{ch:uniform}; none of that closure asserts basis quality.
\end{scopebox}
\section{The original reverse-frame law and the exact source}\label{base:new:source}
From now on $q\ge4k$, $d=q-2k$ and $\rho=k/d\le1/2$. Fix the bases independently of the two refit probes. Along a source-to-root path, let $(g_s,w_s)$ be the current retained/discarded $k$-frames. With the current child placed first, write the column completion
\[
 Q_{V,s}=\begin{pmatrix}A_{V,s}&B_{V,s}\\C_{V,s}&D_{V,s}\end{pmatrix}.
\]
The next sibling $t_{s+1}$ is uniform among row-orthonormal $k$-frames in
$\operatorname{row}[g_s;w_s]^\perp$. The next pair is
\begin{equation}\label{base:new:reverse}
 g_{s+1}=A_{V,s}^Tg_s+C_{V,s}^Tt_{s+1},\qquad
 w_{s+1}=B_{V,s}^Tg_s+D_{V,s}^Tt_{s+1}.
\end{equation}
The transpose stream has the same rule with $U$ and independent draws. At the root draw its final sibling without another compression.

This is a reverse description of the top-down probe, not a new sampling algorithm. A Haar orthonormal triple consisting of a retained pair and the current discarded frame is exchangeable under a fixed orthogonal change inside the retained pair. Conditioning on the current pair leaves the unused sibling uniform in its complement. Subtrees already integrated out do not change this law. The root may first be globally right-rotated by an independent Haar orthogonal matrix: all regression products and output norms are invariant under this rotation, and the resulting root is Haar. This justifies the reverse description even when the implemented root was a fixed frame.

\begin{lemma}[Conditional frame moments]\label{base:new:frame-mom}
If $t$ is a uniform $k$-frame in a rank-$d$ projector $P$, then $\E t=0$ and
\[
 \E[t^THt]=\frac{\tr H}{d}P,\qquad
 \E[tBt^T]=\frac{\tr(PB)}d I_k.
\]
For a past-measurable $Z$, $K=Zt^T$ has conditional covariance
\[
 \E[\vecop K(\vecop K)^T\mid\mathcal F]=d^{-1}I_k\otimes ZPZ^T.
\]
\end{lemma}
\begin{proof}
Left and right orthogonal invariance force the covariance to be scalar on the available subspace and on the frame rows. Taking traces fixes the coefficients. Conditional centering follows by $t\mapsto-t$. The vectorization formula is the entrywise form. These arguments use no Gaussian replacement.
\end{proof}

\subsection{Orthogonality of different initial sources}
Let $\widetilde u_j$ denote each lifted discarded row basis, and $u_*$ the root row basis. They form an orthonormal multiresolution basis, so an arbitrary forward observation $Y=Y(G)$ has the expansion
\[
 Y=\sum_j\widetilde u_jx_j+u_*x_*,\qquad x_j=\widetilde u_j^TY.
\]
Pure retained compression of a source produces no discrepancy before its own discarded node. There it emits $a_j=x_j(R_j^G)^T$ and starts a transpose correction $-v_ja_j^TW_j^H$. All subsequent corrections from that source remain linear in this initial discarded $H$-frame conditional on coarser $H$-frames. Later operations take place only at ancestors and do not introduce a second copy of that initial frame into the linear decoder.

Integrate the finest initial frame in a cross-source product. Its conditional mean is zero, while a coarser source does not use it. At incomparable nodes use conditional independence or the independent block-row sign symmetry of a paired discarded draw. It follows that, conditional on all of $G$, distinct complete source corrections have zero expected inner product; each is mean zero and orthogonal in expectation to the free-compression output. This applies both to the coefficient-output vector and to the discarded-observation ledger. It also proves unbiasedness of the fixed-space refit:
\begin{equation}\label{base:new:unbiased}
 \E\widehat B=P_{\mathcal S}A.
\end{equation}
Unbiasedness by itself is not a variance bound.

\subsection{The actual source: chain, exterior, and interior components}
For $E\perp\mathcal S$, let $e=\widetilde u_j^\T E$ be a discarded row source. Its coefficients in the current compressed local retained and discarded column directions vanish. This does \emph{not} imply that coefficients in finer discarded directions strictly inside node $j$ vanish.

The full orthogonal spatial decomposition uses retained bases $v_{s_n}$ of exterior sibling subtrees, their discarded directions, and every finer discarded column direction inside $j$:
\[
 e=\sum_{n=1}^mB_nv_{s_n}^\T
   +\sum_{\alpha\in\mathrm{out}}B_\alpha\widetilde v_\alpha^\T
   +\sum_{\alpha\in\mathrm{in}}B_\alpha\widetilde v_\alpha^\T.
\]
Put $\beta=\sum_n\F{B_n}^2$, $\gamma_{\rm out}=\sum_{\rm out}\F{B_\alpha}^2$, $\gamma_{\rm in}=\sum_{\rm in}\F{B_\alpha}^2$, and $\gamma=\gamma_{\rm out}+\gamma_{\rm in}$. Parseval gives $\beta+\gamma=\F e^2$. The actual coefficient is
\begin{equation}\label{base:new:actualsource}
 a(G)=eGg_0^\T=N_m+\zeta,\qquad
 N_n=\sum_{r\le n}B_rt_rg_0^\T,\qquad
 \zeta=\sum_{\alpha\in\mathrm{out}\cup\mathrm{in}}B_\alpha W_\alpha^Gg_0^\T.
\end{equation}
For the reverse filtration,
\begin{align}
 \E\F{N_m}^2&\le\rho\beta,\label{base:new:source-var}\\
 \E[N_m^\T N_m\mid\Fil_n]&\preceq N_n^\T N_n+
 d^{-1}\sum_{r>n}\F{B_r}^2I.\label{base:new:source-rem}
\end{align}
Conditional on the entire forward ancestor chain $\mathcal C$,
\begin{equation}\label{base:new:offchain}
 \begin{gathered}
 \E[\zeta\mid\mathcal C]=0,\\
 \E[\zeta^\T\zeta\mid\mathcal C]
 \preceq\left(\frac{\gamma_{\rm out}}{q-3k}
       +\frac{\gamma_{\rm in}}{q-2k}\right)I
 \preceq\frac\gamma{q-3k}I.
 \end{gathered}
\end{equation}
Next to the conditioned exterior chain, two retained frames and one discarded frame have already been exposed. The unexposed partner has a $(q-3k)$-dimensional complement. Strict interior directions instead have a fresh $(q-2k)$-complement after conditioning on their parent pair. Integrating the finest unconditioned frame first and using sibling row-sign symmetry removes cross terms, including those between the two classes. This proves \eqref{base:new:offchain} without a new independence assertion about the chain source.

\begin{example}[The omitted interior-source direction]\label{base:con:interior-example}
Take $k=1$, $L=3$, $N=16$ and balanced bases. In leaf $i$ let
$r_i=(e_{2i-1}+e_{2i})/\sqrt2$ and
$w_i=(e_{2i-1}-e_{2i})/\sqrt2$. The residual $E=r_2w_1^\T$ is orthogonal to the fixed space. At the node containing the first four leaves,
\[
 \widetilde u_j=(r_1+r_2-r_3-r_4)/2,
 \qquad e_j=w_1^\T/2.
\]
All exterior coefficients are zero, but
$a_j=W_1g_j^\T/2$ and
\[
 \E|a_j|^2=\frac1{8(q-2)}=\frac1{16}\quad(q=4).
\]
Indeed the available complement for $W_1$ contains squared $g_j$-length $1/2$. An exterior-only decomposition would incorrectly predict zero. The repaired decomposition retains the same coarse $\gamma/(q-3k)$ allowance, but the old exterior-only equality is false. Every assembly below uses the corrected budget.
\end{example}

\section{One-source state, exact signed covariance, and the linear-depth route}\label{base:new:signed}
The following equations specify the actual chain decoder completely. At one step pad the current $k$-row observations by a zero sibling. Draw forward and transpose siblings $t,\tau_H$. Write
\[
 r=A_V^Tg+C_V^Tt,\quad w'=B_V^Tg+D_V^Tt,
 \quad \widehat h=A_U^Th+C_U^T\tau_H,\quad \chi'=B_U^Th+D_U^T\tau_H.
\]
The nonroot coefficient sectors are
\begin{equation}\label{base:new:readout}
 a=B_U^TYr^T,\qquad b=(B_V^TZ\widehat h^T)^T,\qquad d_Y=B_U^TY(w')^T.
\end{equation}
Their total squared energy is $\F a^2+\F b^2+\F{d_Y}^2$. For the stream-swapped analysis, replace the last sector by $d_Z=(B_V^TZ(\chi')^T)^T$. Both choices must be retained when comparing source orientations. The updated observations are
\begin{equation}\label{base:new:chain-update}
 Y'=A_U^TY-bw',\qquad Z'=A_V^TZ-a^T\chi'.
\end{equation}
At the root, pad both observations by a zero sibling and apply the half-sum of forward and transpose regression from \eqref{base:orf:eq:root}.

Freeze all of $G$. Define
\[
 C=Zh^T,\quad D=Z\chi^T,\quad R=ZZ^T-CC^T,
 \quad J_0=R-DD^T,\quad F=Y-C^Tg.
\]
The notation $J_0$ here is a local innovation matrix, not the external John--Nirenberg constant. It is positive semidefinite. Let $K=Z\tau_H^T$ and $\tau=\tr(C_UC_U^T)/d$. Direct substitution into \eqref{base:new:chain-update} gives
\begin{align}
 C'&=A_V^T(CA_U+KC_U),\label{base:new:C-update}\\
 D'&=A_V^TKD_U-rF^TB_U,\\
 F'&=A_U^TF-C_U^TK^Tg.\label{base:new:F-update}
\end{align}
The conditional mean of $K$ is zero and its covariance is $I_k\otimes J_0/d$. Thus
\begin{align}
 \E_H[R'\mid\mathcal F,G]&=A_V^T(R-\tau J_0)A_V+rF^TB_UB_U^TFr^T,\label{base:new:R-update}\\
 \E_H\F{F'}^2&=\F{A_U^TF}^2+\tau\tr J_0.\label{base:new:F-energy}
\end{align}
With $L_U=B_UB_U^T$ and $L_V=B_VB_V^T$, the potential $\mathcal P=\F F^2+\tr R$ obeys
\begin{equation}\label{base:new:P-drop}
 \mathcal P-\E_H\mathcal P'=\F{B_U^TF(I-r^Tr)}^2+
 (1-\tau)\tr(L_VR)+\tau\F{B_V^TD}^2\ge0.
\end{equation}
At a source injection, $Y=0$, $Z=-a^T\chi$, so $F=C=0$, $R=a^Ta$, $J_0=0$. Hence $\E_H\mathcal P_s\le\F a^2$ uniformly, even for $a=a(G)$.

\subsection{Why this first gave only linear depth}
Let $J_h=\sum_{s<h}\E_H[\tau_s\tr J_{0,s}]$. The estimate above gives $J_h\le\rho h\F a^2$. Split the coefficient output into its $F$ part and common-regression part. Orthogonality of emitted sectors and telescoping \eqref{base:new:F-energy} give
\[
 \E_H\F{\mathcal O_F}^2\le J_h,
 \qquad \E_H\F{\mathcal O_C}^2=J_h+\frac\rho4\E_H\tr J_{0,h}.
\]
The squared triangle inequality then gives source gain at most $\rho(4h+1/2)$. Adding the free output, both source orientations, and the orthogonal residual gives
\begin{equation}\label{base:new:linear}
 \E\F{A-\widehat B}^2\le(5+16\rho L+2\rho)d_{\mathcal S}(A)^2;
 \quad q=4k:\quad 8L+6.
\end{equation}
This was a genuine arbitrary-basis, actual-input improvement over exponential depth. Its loss comes from separating two output components that can cancel.

\subsection{Exact cancellation with a linearly growing internal budget}
For $k=1$, repeat $Q_U=\bigl(\begin{smallmatrix}0&-1\\1&0\end{smallmatrix}\bigr)$ and $Q_V=I$. Starting from a unit source, $Y$ remains zero, all nonroot output is zero, and root output is at most $1/4$. Nevertheless, with $d=q-2$,
\[
 u_0=0,\quad v_0=1,\qquad
 u_{s+1}=(1-u_s-v_s)/d,\quad v_{s+1}=u_s
\]
has fixed point $u=v=1/q$ and geometric transient. Hence $J_h=\sum_{s<h}u_{s+1}=h/q+O_q(1)$. Bounding $J_h$ alone cannot finish the proof. The signed equality $Y=F+C^Tg=0$ explains the discrepancy.

\subsection{Exact conditional response matrix}
Retain
\[
 \Sigma=\E_H[\vecop(F,C)\vecop(F,C)^T],\quad\overline R=\E_HR,
 \quad\overline J=\E_HJ_0.
\]
Define linear maps $\mathsf A(F,C)=(A_U^TF,A_V^TCA_U)$ and
$\mathsf B(K)=(-C_U^TK^Tg,A_V^TKC_U)$. Then
\begin{align}
 \Sigma'&=\mathsf A\Sigma\mathsf A^T+
 \mathsf B(I_k\otimes\overline J/d)\mathsf B^T,\label{base:new:sig-update}\\
 \overline R'&=A_V^T(\overline R-\tau\overline J)A_V+
 r\E_H[F^TL_UF]r^T,\\
 \overline J'&=A_V^T(\overline R-\rho\overline J)A_V.\label{base:new:J-update}
\end{align}
In the last equation the common feedback square cancels from $R'$ and $D'D'^T$. Every readout in \eqref{base:new:readout}, including its signed cross term and the root, is a linear functional of this state. Initialization depends on $a$ only through $a^Ta$. Therefore
\begin{equation}\label{base:new:W}
 \E_H[\F{\mathcal O_{G,H}(a(G))}^2\mid G]
 =\tr(W(G)a(G)^Ta(G)),\qquad W(G)=\sum_nW_n(G)\succeq0.
\end{equation}
This finite transfer has $O(k^4)$ real coordinates for $q=O(k)$ and no depth-dependent state dimension. A finite exact transfer is not a uniform gain bound.

\section{Adapted restarts: the full covariance induction and response tails}\label{base:new:carleson}
All expectations in the next argument may be conditional on an actual current history. The two future streams must retain their original conditionally independent Haar laws. In particular, this is not a restart after exposing a complete future forward path.

\begin{lemma}[Full cross tensor and one-sided covariance cone]\label{base:con:restart-cone}
Start from fixed current observations $(Y_0,0)$. Put $P_s=Y_sg_s^\T$, $C_s=Z_sh_s^\T$, $M_0=Y_0g_0^\T$, and $M_{s+1}=A_{U,s}^\T M_sA_{V,s}$. Then
\begin{align}
 \E P_s&=M_s,&\E C_s&=0,\label{base:con:cone-means}\\
 \E[\vecop(P_s)\vecop(C_s)^\T]&=0,\label{base:con:cone-cross}\\
 \operatorname{Cov}(\vecop P_s)&\preceq d^{-1}I_k\otimes\E[Y_sY_s^\T],\label{base:con:cone-P}\\
 \E[\vecop(C_s)\vecop(C_s)^\T]&\preceq d^{-1}I_k\otimes\E[Z_sZ_s^\T].\label{base:con:cone-C}
\end{align}
The cross identity is for the entire tensor, not only its trace.
\end{lemma}
\begin{proof}
Let $K_G=Yt^\T$, $K_H=Z\eta^\T$, $J_Y=Y(I-g^\T g-w^\T w)Y^\T$, and $J_Z=Z(I-h^\T h-\chi^\T\chi)Z^\T$. Their conditional covariances are $I_k\otimes J_Y/d$ and $I_k\otimes J_Z/d$. They are centered and conditionally independent. Frame orthogonality gives
\[
 P'=A_U^\T(PA_V+K_GC_V),\qquad
 C'=A_V^\T(CA_U+K_HC_U).
\]
Centering proves the means. In the cross-tensor recursion every fresh term vanishes; the old cross tensor undergoes deterministic transport and remains zero.

The observation cross terms also vanish in expectation. Indeed
$Y(w')^\T=PB_V+K_GD_V$ and
$b^\T=B_V^\T(CA_U+K_HC_U)$. Their product has one old cross-tensor term, two centered terms, and one conditionally independent fresh product. Therefore, exactly,
\begin{align}
 \E[Y'Y'^\T]&=A_U^\T\E[YY^\T]A_U+\E[bb^\T],\label{base:con:rowY}\\
 \E[Z'Z'^\T]&=A_V^\T\E[ZZ^\T]A_V+\E[a^\T a].\label{base:con:rowZ}
\end{align}
The centered $P'$ covariance is at most
\begin{align*}
 &d^{-1}(A_V^\T A_V)\otimes(A_U^\T\E[YY^\T]A_U)\\
 &\quad+d^{-1}(C_V^\T C_V)\otimes(A_U^\T\E[J_Y]A_U)\\
 &\preceq d^{-1}I_k\otimes(A_U^\T\E[YY^\T]A_U)
 \preceq d^{-1}I_k\otimes\E[Y'Y'^\T].
\end{align*}
Here $J_Y\preceq YY^\T$, $A_V^\T A_V+C_V^\T C_V=I$, and \eqref{base:con:rowY} were used. The transpose cone follows identically. Both centered initial covariances are zero, completing the induction.
\end{proof}

\begin{lemma}[Adapted restart and response tail]\label{base:new:restart}
With $S_s=\F{Y_s}^2+\F{Z_s}^2$ and
$D_s^{\rm obs}=\F{B_{U,s}^\T Y_s}^2+\F{B_{V,s}^\T Z_s}^2$, a one-sided restart satisfies
\begin{equation}\label{base:con:strong-ledger}
 (1-\rho)\sum_{s<h}\E D_s^{\rm obs}+\E S_h
 \le\F{Y_0}^2+\F{Y_0g_0^\T}^2.
\end{equation}
Consequently, for current-history-measurable data,
\begin{align}
 \E[\mathscr W_{\rm future}(Y,0)\mid\mathcal A]
 &\le\frac{\F Y^2+\F{Yg^\T}^2}{1-\rho},\label{base:new:restart-one}\\
 \E[\F{\mathcal R_{\rm future}(Y,Z)}^2\mid\mathcal A]
 &\le\frac4{1-\rho}(\F Y^2+\F Z^2).\label{base:new:restart-two}
\end{align}
For a deterministic injected coefficient, the positive response increments satisfy
\begin{equation}\label{base:new:response-tail}
 \E[\sum_{j\ge n}W_j\mid\Fil_n^G]\preceq\kappa I,
 \qquad \E W\preceq c_0I,
\end{equation}
where $c_0=(1-\rho)^{-1}$ and
$\kappa=(1+4c_0)(1+2\rho+2\sqrt\rho)$.
\end{lemma}
\begin{proof}
Contraction of the covariance tensors in Lemma~\ref{base:con:restart-cone} gives
\[
 \E\F a^2\le\rho\E\F{B_U^\T Y}^2+\F{B_U^\T MA_V}^2,
 \qquad \E\F b^2\le\rho\E\F{B_V^\T Z}^2.
\]
Taking traces in \eqref{base:con:rowY}--\eqref{base:con:rowZ} yields
$\E S_{s+1}=\E S_s-\E D_s^{\rm obs}+\E\F a^2+\E\F b^2$.
The deterministic contribution telescopes, since
\[
 \F{B_U^\T MA_V}^2+\F{A_U^\T MA_V}^2
 =\F{MA_V}^2\le\F M^2.
\]
This proves \eqref{base:con:strong-ledger}. Each nonroot readout costs at most $D_s^{\rm obs}$ and the root costs at most $S_h/2$. Hence the first restart bound follows. Split arbitrary data into $(Y,0)+(0,Z)$ and use the squared triangle inequality for the second.

For the response tail, freeze the complete forward path for a deterministic test source only. Equations~\eqref{base:new:C-update} and \eqref{base:new:P-drop} imply
\[
 \E_H C^\T C\preceq\frac{\F a^2}{d}I,
 \qquad \E_H(\F Y^2+\F Z^2)
 \le(1+2\rho+2\sqrt\rho)\F a^2.
\]
The first inequality is an induction: its old term is at most
$\F a^2 A_U^\T A_U/d$ and its new term at most
$\F a^2 C_U^\T C_U/d$, because $\E_H\tr J_0\le\F a^2$.
For the second use $Y=F+C^\T g$, $\E_H(\F F^2+\tr R)\le\F a^2$, and Cauchy--Schwarz. Charge the current output to current observation energy, then restart before any remaining forward frames are exposed. This costs a further $4c_0$ times that energy. Average the past transpose state and test all deterministic source Grams to obtain the first part of \eqref{base:new:response-tail}. At injection the retained regressions vanish, so the one-sided ledger gives its mean bound.
\end{proof}
\begin{remark}
The retained-mean term in \eqref{base:con:strong-ledger} is essential at an arbitrary adapted restart. For a canonical injected source it is zero. The explicit cross-tensor proof replaces a previously incomplete proof interface. The later quotient-variance theorem will sharpen the zero-regression \emph{output} bound from $c_0$ to a quantity of order $\rho$; the older constants are kept here to identify the historical route.
\end{remark}

\subsection{Adapted embedding and anticipative terminal coefficients}
For adapted $A_n\succeq0$ with conditional tails bounded by $KI$, scalar stopping and Doob's inequality give
\[
 \sum_n\E\langle A_nf_n,f_n\rangle\le b_rK\E\|f_m\|^2
\]
for an $r$-dimensional martingale, with elementary $b_r=4r$. Indeed replace $A_n$ by $\tr A_n$; its scalar conditional tail is at most $rK$, apply the layer-cake stopping argument and $\E\sup_n\|f_n\|^2\le4\E\|f_m\|^2$. The public matrix Carleson embedding theorem gives the alternative $b_r\le C_{\rm CE}(1+\log r)^2$; see \cite{NTV,NPTV}. For finite Haar filtrations this follows first for finite partitions, represented by dyadic refinements with inserted zero weights, then by $L_2$ approximation. The constants refer to the squared embedding norm.

A finite-time alternative has
\begin{equation}\label{base:new:time-embed}
 b_m^{\rm time}=[1+\lceil\log_2m\rceil]^2.
\end{equation}
Decompose each prefix of martingale differences into maximal dyadic intervals. There are at most $1+\lceil\log_2m\rceil$ pieces. The interval sum is measurable at its right endpoint; the conditional-tail condition bounds all weights subsequently assigned to it. At any fixed dyadic scale, orthogonality bounds total interval energy by terminal martingale energy. One logarithm counts prefix pieces, the other counts scales. Zero deterministic source blocks may be removed, so only active stages need be counted. More generally an ordered binary tree gives profile $(\sum_\ell\sqrt{v_\ell})^2$, where $v_\ell$ is source energy carried by prefix intervals at depth $\ell$.

The terminal coefficient $N_m$ is not adapted at time $n$. Use \eqref{base:new:source-rem} before applying the embedding. Together with \eqref{base:new:offchain} this gives
\begin{equation}\label{base:new:Cstar}
 \E\tr(Wa^Ta)\le C_*\F e^2,\quad
 C_*=
 \max\left\{\rho\left(\kappa b+\frac1{1-\rho}\right),
 \frac{k}{(q-3k)(1-\rho)}\right\},
\end{equation}
where $b$ may be any valid rank or time constant above. No factorization of $\E W$ and $\E a^Ta$ is used. Source orthogonality gives the fundamental assembly rule
\begin{equation}\label{base:new:assembly}
 \boxed{\E\F{A-\widehat B}^2\le(5+4C_*)d_{\mathcal S}(A)^2.}
\end{equation}
At width $4k$, the elementary bound is
$[9+(144+72\sqrt2)k]d_{\mathcal S}(A)^2<255k\,d_{\mathcal S}(A)^2$.
Rank/time embedding gives
\begin{equation}\label{base:new:general-log}
 \E\F{A-\widehat B}^2\le C[1+\log\min\{k,L+1\}]^2d_{\mathcal S}(A)^2.
\end{equation}
One may take the minimum with \eqref{base:new:linear} and any geometric certificate below. Alternatively $q=\lceil(4+b_k)k\rceil$ yields an absolute factor at most $132$, at $O(k(1+\log k)^2)$ refit queries. That earlier oversampling route does not have the desired uniform $O(k)$ query cost. Theorem~\ref{base:con:main-eight} improves its dependence on the embedding constant from linear to square-root, without proving a uniform $O(k)$ bound.


\section{A genuine conditional amplifier}\label{base:new:amplifier}
The statement $W(G)\preceq CI$ for every basis family and every complete forward path is false, even at rank one. This is stronger than the growing-internal-budget example of Section~\ref{base:new:signed}: the \emph{combined actual output} can grow conditionally.

Take $k=1,q=4$. For an integer $M\ge1$, let $n$ be the least odd integer at least $4\pi^2M$, and put
\[
 c=2\cos(\pi/n)-1,\qquad s=\sqrt{1-c^2}.
\]
After two warm-up steps, use $M$ identical blocks. Each block has $n$ steps with
$Q_U=I$ and $Q_V=\bigl(\begin{smallmatrix}c&-s\\s&c\end{smallmatrix}\bigr)$, followed by a release step $Q_U=\bigl(\begin{smallmatrix}0&-1\\1&0\end{smallmatrix}\bigr)$, $Q_V=I$. An admissible forward triple can follow
\begin{equation}\label{base:new:amplifier-triple}
 \begin{pmatrix}g'\\w'\\t'\end{pmatrix}
 =\begin{pmatrix}c&0&s\\-s&0&c\\0&1&0\end{pmatrix}
 \begin{pmatrix}g\\w\\t\end{pmatrix}.
\end{equation}
Each new sibling is perpendicular to the current pair. The displayed matrix has eigenvalues $-1,e^{i\pi/n},e^{-i\pi/n}$; since $n$ is odd its $n$th power is $-I$. Hence $g_n=-g_0$, although $c^n$ is close to one.

At block entry $Y=0$. Write $C=Zh^T$. During rotation, $F=-Cg_0$ and $C_j=c^jC$, so $Y_n=-(1+c^n)Cg_0$. The release turns this forward displacement into output. The exact expected block output is
\begin{equation}\label{base:new:amplifier-output}
 2(1+c^n)u,\qquad u=\E_H C^2.
\end{equation}
For $p=\E_H\|Z\|^2$, $v=\E_H(Zw_H^T)^2$, $b=c^{2n}$ and $r_n=(-1/2)^n$, the block recurrence is
\[
 p'=b(p-u)+u,\quad
 u'=b\left[\frac{2+r_n}{6}(p-u)-\frac{r_n}{2}v\right],\quad v'=u.
\]
The warm-up gives $(p,u,v)=(1,1/2,0)$. The choice of $n$ guarantees $p\ge1/2$ and $u\ge1/32$ during the $M$ blocks: substitute the previous bounds in the recurrence and use $1-c^{2n}\le2\pi^2/n$ to bound cumulative attenuation. Each block therefore emits at least $1/16$, proving
\begin{equation}\label{base:new:conditional-unbounded}
 W(G)\ge M/16,\qquad h=2+M(n+1).
\end{equation}
The path is in the support of the Haar law. Continuity gives a positive-probability neighborhood with gain at least $M/32$. No probability lower bound independent of $M$ is asserted.

The construction can be embedded into a fixed-space residual with one discarded row direction and an exterior-sibling retained column direction. Its actual scalar coefficient is nonzero on the specified path, and homogeneity preserves the gain relative to that coefficient. This does not give an unbounded expected reconstruction error relative to $\F E^2$: the coefficient amplitude and the probability of the path matter.

\subsection{Averaging the forward path changes the result}
For the same fixed bases and an independent unit source, a block with $\alpha=c^n$, $b=\alpha^2$ satisfies
\[
 \E_G[g_n\mid\text{block-entry history}]=\alpha g_0,\quad
 Y_n=C(\alpha g_n-g_0),\quad
 \E_G\|Y_n\|^2=(1-b)C^2.
\]
The rotation output is $(1-b)C^2$ and release output is at most $\|Y_n\|^2$. With $\zeta=g_0g_n^T$,
$Z_{\rm new}=\alpha(Z-Ch)+\zeta Ch$, so $\|Z_{\rm new}\|^2\le\|Z\|^2$. Thus $C^2\le1$, total expected nonroot output is at most $2M(1-b)\le1$, and the root costs at most $1/4$. The fully averaged independent-source gain is at most $5/4$.

Later, taking $n\ge4\pi^2M^2$ makes the cumulative column loss at most $1/(2M)$ while preserving \eqref{base:new:conditional-unbounded}. The loss-sensitive theorem below then gives actual source-weighted average at most $14.402\beta/M$ at width four. This is the explicit separation between a diverging conditional gain and a vanishing actual joint average. It invalidates the universal conditioning strategy, not the original randomized target.

\section{Temporal Gram structure and immediate parity}\label{base:new:temporal}
Let $T_r=t_rg_0^T$ and stack them as $X=[T_1;\ldots;T_m]$. For $B=[B_1\ \cdots\ B_m]$, define
\[
 \mathcal K=\E[XWX^T]\succeq0,\qquad
 \mathcal K_{rs}=\E[T_rWT_s^T].
\]
Then the exact chain response is $\tr(B\mathcal K B^T)$. This temporal operator retains the dependence discarded by an absolute path sum.

At step $n$, the two row-discarded readout sectors use $r=A_V^Tg+C_V^Tt$ and $w'=B_V^Tg+D_V^Tt$. Orthogonality gives
\[
 r^Tr+(w')^Tw'=g^Tg+t^Tt.
\]
Consequently the immediate response has the form $W_n=W_n^{(0)}+\Phi_n(t_n^Tt_n)$ for a predictable positive map $\Phi_n$. It is even under $t_n\mapsto-t_n$, whereas the fresh source increment is odd. Their immediate mixed contribution vanishes exactly. Future states need not be even under that reflection. For $r<s$,
\begin{equation}\label{base:new:delayed-only}
 \mathcal K_{rs}=\E\left[T_r\Big(\sum_{n>s}W_n\Big)T_s^T\right].
\end{equation}
Outputs before $s$ vanish by fresh centering; output at $s$ vanishes by parity. The diagonal satisfies
\[
 \mathcal K_{ss}\preceq\rho(\kappa+c_0)I.
\]
Split the response before/after $s$, apply conditional frame isotropy to the former and the response-tail estimate to the latter. If $B_r^TB_s=0$ for distinct times in every source and both orientations, the cross blocks contribute nothing, giving at width $4k$ the restricted factor
$9+18(2+\sqrt2)<71$. This is a condition on residual source blocks, not a theorem for arbitrary input residuals.

\subsection{An initial-anchor positive decomposition}
Let the first column completion have blocks $(A_c,B_c;C_c,D_c)$. Then $g_0=A_cg_1+B_cw_1$. The first source transition emits nothing and has signed state $\Sigma_1=0$, $\overline R_1=\overline J_1=A_c^Ta^TaA_c$. Conditional on $g_1,t_2$ and the later chain, the response does not use $w_1$ separately. The latter is uniform in their $d$-dimensional complement. With
\[
 P=I-g_1^Tg_1-t_2^Tt_2,\qquad
 \lambda=\tr(B_c^TWB_c)/d,
\]
conditional integration gives
\[
 \E[g_0^TWg_0\mid\text{reduced chain}]
 =g_1^TA_c^TWA_cg_1+\lambda P.
\]
For $r,s\ge2$ this gives two positive temporal Grams,
\[
 \mathcal K^{\rm ret}_{rs}=\E[(t_rg_1^T)A_c^TWA_c(g_1t_s^T)],
 \quad\mathcal K^{\rm inn}_{rs}=\E[\lambda t_rPt_s^T].
\]
The innovation part has a scalar weight. Apply the scalar Carleson embedding to $\sum_{r\ge3}H_rt_rP$ and its remaining variance budget to obtain
\begin{equation}\label{base:new:innovation-bound}
 \|\mathcal K^{\rm inn}\|\le\frac{\|B_c\|_F^2}{d}(4\kappa+c_0)
 \le\rho(4\kappa+c_0).
\end{equation}
The retained part remains coupled. Repeating this initial-anchor regression mechanically is invalid because, after a nonzero noise injection, the response retains the previous forward direction explicitly. Also $T_1=T_2=0$, $\mathcal K_{34}=0$, and
$\mathcal K_{33}=d^{-1}\tr(B_c^T(\E W)B_c)I$. These are genuine early cancellations, not all-time diagonality.

\section{A joint fourth-order transfer under the Haar law}\label{base:new:fourth}
The conditional-$H$ transfer can be extended to integrate both streams while retaining two source-time marks. This supplies an exact finite evaluation method for $\mathcal K$, not a uniform bound by itself.

Use $X=g_0$ as a fixed anchor and retain six projection matrices
\[
 P=Yg^T,\ D=Yw^T,\ Q=Zh^T,\ E_Z=Z\chi^T,
 \ A_X=Xg^T,\ B_X=Xw^T,
\]
and three Gram/cross-Gram matrices $S_Y=YY^T$, $S_Z=ZZ^T$, $C_X=YX^T$, together with accumulated output energy. Draw fresh blocks $K=Yt^T$, $L=Xt^T$, $N_H=Z\tau_H^T$. Their conditional residual Grams are
\[
 R_G=\begin{pmatrix}
 S_Y-PP^T-DD^T&C_X-PA_X^T-DB_X^T\\
 *&I-A_XA_X^T-B_XB_X^T
 \end{pmatrix},\qquad
 R_H=S_Z-QQ^T-E_ZE_Z^T.
\]
The two fresh blocks are conditionally independent. All updates follow by substitution in \eqref{base:new:chain-update} and the new frame pair. For example
\[
 P'=A_U^T(PA_V+KC_V),\qquad
 C_X'=A_U^TC_X-b(A_XB_V+LD_V)^T,
\]
where $b=(B_V^T(QA_U+N_HC_U))^T$. For the remaining registers, form $Y',Z',g',w',h',\chi'$ using \eqref{base:new:chain-update} and multiply by their named partners. This is an explicit polynomial specification of all updates, without requiring an external code file.

The second moments are Lemma~\ref{base:new:frame-mom}. The fourth moments are the orthogonal Haar four-entry formula: sum over the three row pairings and three column pairings, with diagonal coefficient
\[
 \frac{d+1}{d(d-1)(d+2)}
\]
and off-diagonal coefficient $-1/[d(d-1)(d+2)]$. These are the classical orthogonal Weingarten coefficients~\cite{CS}. For one spherical row the formula reduces to the three covariance pairings divided by $d(d+2)$. Using the Gaussian denominator $d^2$ would give a different temporal Gram.

Assign degree one to projection registers and degree two to Gram and energy registers. The local substitutions preserve weighted degree, and Haar integration replaces a fresh moment by residual-Gram polynomials of the same degree. Consequently the space of polynomials of weighted degree at most four is invariant under conditional expectation. It has $O(k^8)$ coordinates at fixed width ratio; its dimension is independent of depth. A temporal insertion $t_rg_0^T=L^T$ has degree one, so two marks and output energy fit in degree four. Compose the local operators with these insertions to compute every entry of $\mathcal K$ under the actual joint law.

\paragraph{Finite certificates recorded by this route.}
In the scalar implementation the filtered polynomial space had 332 monomials. For five nonroot steps with both completions $\frac15\bigl(\begin{smallmatrix}3&-4\\4&3\end{smallmatrix}\bigr)$, the recorded exact original-law values were
\[
 \mathcal K_{35}=\frac{2925051264}{95367431640625}>0,\qquad
 \mathcal K_{45}=\frac{4766750208}{476837158203125}>0.
\]
Thus actual delayed terms are nonzero, not merely artifacts of an abstract counterexample. The recorded row-sum bound for this Gram was less than $1/25$. Enumeration of all 126 scalar path/readout cases of length at most five for the same rotation tree gave a row-sum bound below $0.104$ and a factor-$13$ finite-family refit certificate. These rational calculations are historical computational claims specified here by their recurrence and fractions; their original 332-state implementation is not re-executed in this consolidation. They are not used to prove the all-depth matched-weight theorem below. Negative residual terms and negative Weingarten coefficients also mean the monomial transfer is not entrywise positive; finite closure alone cannot prove stability.

\section{Why generic isotropy does not force dimension-free embedding}\label{base:new:abstract-obstruction}
Take $m\ge4$, $r\ge2m$, and let $v_1,\ldots,v_m$ be the first columns of a Haar orthogonal $r\times r$ matrix. Define
\[
 f_n=\sum_{i\le n}v_i,\quad
 u_j=\sum_{i\le j}\frac{v_i}{j-i+1},\quad
 A_1=0,\quad A_n=u_{n-1}u_{n-1}^T\ (n\ge2).
\]
The increments $v_n$ are centered spherical vectors in the remaining subspace, with deterministic variance budgets. The weights are predictable and hence even under reflection of the current increment. Their norms are at most $\pi^2/6$, and their conditional future sums have a universal norm bound. To check the latter, condition on $v_1,\ldots,v_n$. The past-coordinate contribution is the Gram of the Hilbert matrix $1/(i+j+1)$, whose norm is at most $\pi$; the unrevealed part is scalar on the orthogonal complement and has norm at most $\pi^2/6$. The Hilbert bound follows from Schur's test with weights $(j+1/2)^{-1/2}$ and the integral $\int_0^\infty t^{-1/2}/(1+t)dt=\pi$. Cross terms vanish by conditional centering.

Yet orthogonality gives, pathwise,
\[
 \langle f_n,u_{n-1}\rangle=\langle f_m,u_{n-1}\rangle=H_{n-1},
 \qquad \sum_n\langle A_nf_n,f_n\rangle=\sum_{j<m}H_j^2.
\]
With total deterministic variance budget at most $4m$, the normalized weighted response is at least $\log^2(m/2)/16$. Thus spherical increments, deterministic variance budgets, a response-tail bound and immediate parity do not by themselves yield a dimension-free embedding. This is an abstract adaptation of the harmonic matrix-Carleson mechanism~\cite{NTV,NPTV}, not an HSS realization. Its growing available orthogonal source family differs from a fixed-width HSS chain.

The same example has response norm of order $\log^2m$. Taking $m$ exponential in the Schatten moment order shows that response moments growing like $p^2$ cannot be improved to $O(p)$ from the abstract response-tail condition alone. The specialized source--response relation, rather than separate isotropic moment bounds, must provide any further unrestricted improvement.



\endgroup

% ===== GEOMETRY =====
\begingroup
\let\R\relax
\newcommand{\R}{\mathbb R}
\let\E\relax
\newcommand{\E}{\mathbb E}
\let\Prb\relax
\newcommand{\Prb}{\mathbb P}
\let\T\relax
\newcommand{\T}{\mathsf T}
\let\F\relax
\newcommand{\F}[1]{\left\lVert#1\right\rVert_{F}}
\let\N\relax
\newcommand{\N}[1]{\left\lVert#1\right\rVert}
\let\ip\relax
\newcommand{\ip}[2]{\left\langle#1,#2\right\rangle_F}
\let\tr\relax
\newcommand{\tr}{\operatorname{tr}}
\let\rank\relax
\newcommand{\rank}{\operatorname{rank}}
\let\diag\relax
\newcommand{\diag}{\operatorname{blockdiag}}
\let\range\relax
\newcommand{\range}{\operatorname{range}}
\let\OPT\relax
\newcommand{\OPT}{\operatorname{OPT}}
\let\HSS\relax
\newcommand{\HSS}{\operatorname{HSS}}
\let\diver\relax
\newcommand{\diver}{\operatorname{div}}
\let\Var\relax
\newcommand{\Var}{\operatorname{Var}}
\let\UC\relax
\newcommand{\UC}{\mathcal U}
\let\VC\relax
\newcommand{\VC}{\mathcal V}
\let\err\relax
\newcommand{\err}{\mathcal E}
\let\proofstep\relax
\newcommand{\proofstep}[1]{\par\addvspace{.45\baselineskip}\Needspace{3\baselineskip}\textbf{#1}\par\nobreak}
\let\op\relax
\newcommand{\op}[1]{\left\lVert#1\right\rVert_{\mathrm{op}}}
\let\Sn\relax
\newcommand{\Sn}[2]{\left\lVert#1\right\rVert_{S_{#2}}}
\let\Cov\relax
\newcommand{\Cov}{\mathsf C}
\let\Loss\relax
\newcommand{\Loss}{\mathsf L}
\let\Fil\relax
\newcommand{\Fil}{\mathcal F}
\let\Id\relax
\newcommand{\Id}{I}
\let\supp\relax
\newcommand{\supp}{\operatorname{supp}}
\let\vecop\relax
\newcommand{\vecop}{\operatorname{vec}}
\let\Gammafun\relax
\newcommand{\Gammafun}{\operatorname{Gamma}}
\chapter{Storage, refresh, and protected channels}\label{ch:geometry}
\begin{scopebox}
The restrictions in these specialized geometric certificates are kept explicit. They are sufficient routes, not necessary conditions for the uniform fixed-space theorem proved in Chapter~\ref{ch:uniform}.
\end{scopebox}
\section{Scalar storage, balanced leakage, and finite exceptions}\label{base:new:scalar-storage}
This part gives response bounds without the general Carleson interface. Freeze the complete forward sketch and use the signed state of Section~\ref{base:new:signed}. Set $L_U=B_UB_U^T$, $L_V=B_VB_V^T$, $\tau=\tr L_U/d$, and
\[
 \mathsf Q=4\tau I+2(\rho-\tau)L_V,\qquad
 \mathsf D=(1-\tau)L_V.
\]
Expanding the three coefficient sectors in \eqref{base:new:readout}, using $Y=F+C^Tg$, bounding the two common-component sums by twice their squares, and using the centered covariance of $K$, gives for either readout
\begin{equation}\label{base:new:output-storage-core}
 \E_H[e+2\F{F'}^2+2\F{C'}^2]
 \le2\F F^2+2\F C^2+\tr(\mathsf QJ_0).
\end{equation}
For the second readout the additional local common-regression term is
$(\rho-\tau)\tr(L_VJ_0)$; the last term in $\mathsf Q$ includes it. Dropping this term would falsely identify the two orientations.

For $\lambda\ge2\rho$, put
\[
 \mathcal H_\lambda=(\lambda+2)\F F^2+\lambda\tr R+2\F C^2,
\]
\[
 \delta_t(\lambda)=\left[4\tau_t-
 \{\lambda(1-\tau_t)-2(\rho-\tau_t)\}\lambda_{\min}(L_{V,t})\right]_+.
\]
Combining \eqref{base:new:P-drop} and \eqref{base:new:output-storage-core}, using $0\preceq J_0\preceq R$ and positivity of $\mathsf Q$, proves
\begin{equation}\label{base:new:scalar-storage-ineq}
 \E_H[e+\mathcal H_\lambda']\le\mathcal H_\lambda+
 \delta_t(\lambda)\tr R.
\end{equation}
The root is bounded by terminal storage. Since $\E_H\tr R_t\le\F a^2$, this gives the finite conditional certificate
\begin{equation}\label{base:new:defect-profile}
 W(G)\preceq\mathfrak B I,
 \quad\mathfrak B=\inf_{\lambda\ge2\rho}
 \left[\lambda+\sum_t\delta_t(\lambda)\right].
\end{equation}
It is a convex piecewise-linear scalar minimization with explicit breakpoints. Its finiteness on a finite chain is not uniform boundedness in depth.

\subsection{Matched child weights: factor thirteen at width four}
Suppose
\[
 U_I=\begin{pmatrix}\sqrt{p_I}O_{I,1}\\\sqrt{1-p_I}O_{I,2}\end{pmatrix},\qquad
 V_I=\begin{pmatrix}\sqrt{p_I}P_{I,1}\\\sqrt{1-p_I}P_{I,2}\end{pmatrix},
\]
with arbitrary orthogonal blocks and $p_I\in[0,1]$. At either child, $L_U=L_V=xI$, so at $q=4k$,
$\tau=x/2$, $\mathsf Q=x(3-x)I$, $\mathsf D=x(1-x/2)I$. Thus
$4\mathsf D-\mathsf Q=x(1-x)I\succeq0$. Every defect in \eqref{base:new:scalar-storage-ineq} vanishes for $\lambda=4$, giving $W(G)\preceq4I$ uniformly.

The actual chain source has $\E N^TN\preceq\beta I/(2k)$; hence its weighted output is at most $2\beta$. The off-chain term is at most $2\gamma$ by \eqref{base:new:offchain} and the independent-source mean bound. The assembly rule gives
\begin{equation}\label{base:new:matched13}
 \E\F{A-\widehat B}^2\le13d_{\mathcal S}(A)^2.
\end{equation}
The weights may approach zero or one, and the row and column orthogonal directions need not agree or commute. The condition is matching scalar child weights; even $U=V$ does not imply it if child singular values vary by direction.

More generally, if in both orientations
\[
 k^{-1}\tr L_U\le\Lambda\lambda_{\min}(L_V),\qquad
 k^{-1}\tr L_V\le\Lambda\lambda_{\min}(L_U),\qquad\Lambda\ge1,
\]
then $\lambda=4\Lambda$ has zero defect. At most $b$ arbitrary exceptional steps per source path cost at most $2b$, since $\delta_t(4\Lambda)\le4\tau_t\le2$. The resulting factor is
\begin{equation}\label{base:new:exceptions}
 5+8\Lambda+4b.
\end{equation}
This charges the total number of exceptions; the next matrix route instead handles indefinitely many exceptions with controlled spacing.

\subsection{Why no larger scalar weight repairs directional mismatch}
Take $k=2,q=8$, $P=\operatorname{diag}(1,0)$, $Q=I-P$, and the same completion in both orientations:
\[
 \begin{pmatrix}P&-Q\\Q&P\end{pmatrix}.
\]
A valid source reaches $Y=F=C=D=0$, $R=J_0=e_1e_1^T$. Write the next fresh regression $K=e_1\xi^T$, with $\E\xi_2^2=1/4$. The next local coefficient output is zero, but
\[
 C'=\xi_2e_1e_2^T,\quad F'=-\xi_2e_2g_1,
 \quad R'=(1-\xi_2^2)e_1e_1^T.
\]
Consequently $\E\mathcal P'=1$ while $\E\mathcal H_\lambda'-\mathcal H_\lambda=1$ for every $\lambda$. The actual equality $Y'=F'+(C')^Tg'=0$ explains the missing cancellation. This refutes the scalar storage family, not constant-factor refitting.

\section{Backward matrix storage and temporal loss}\label{base:new:matrix-storage}
Choose deterministic matrices $0\preceq\mathsf M_t\preceq\alpha I$, depending only on the fixed future bases, with terminal value $\mathsf M_h=\rho I/2$. Impose
\begin{equation}\label{base:new:Mrec}
 \mathsf M_t\succeq(1-\tau_t)A_{V,t}\mathsf M_{t+1}A_{V,t}^T
 +\alpha\tau_tI+\mathsf Q_t.
\end{equation}
Then
\[
 \mathcal H_t=(\alpha+2)\F{F_t}^2+\tr(\mathsf M_tR_t)+2\F{C_t}^2
\]
satisfies $\E_H[e_t+\mathcal H_{t+1}]\le\mathcal H_t$, and
\begin{equation}\label{base:new:Mbound}
 W(G)\preceq\mathsf M_0.
\end{equation}
To verify this, weight \eqref{base:new:R-update} by $\mathsf M_{t+1}$. Its $F$ term plus $\alpha\F{A_U^TF}^2$ is at most $\alpha\F F^2$, using $\mathsf M_{t+1}\preceq\alpha I$ and $A_UA_U^T+L_U=I$. The remaining coefficient of $J_0$ is positive; replace $J_0$ by $R$ and use \eqref{base:new:Mrec}. Terminal storage dominates either averaged root readout. This proof is deterministic in $G$ and valid for its dependent source coefficient.

Define $\mathcal T_t(X)=(1-\tau_t)A_{V,t}XA_{V,t}^T$ and propagate
\begin{equation}\label{base:new:OB}
 \mathsf O_h=I,\ \mathsf B_h=\rho I/2,\quad
 \mathsf O_t=(1-\tau_t)L_{V,t}+\mathcal T_t(\mathsf O_{t+1}),\quad
 \mathsf B_t=\mathsf Q_t+\mathcal T_t(\mathsf B_{t+1}).
\end{equation}
At finite depth $0\prec\mathsf O_t\preceq I$. The least ceiling within this construction is
\[
 \alpha_*=
 \max_t\lambda_{\max}(\mathsf O_t^{-1/2}\mathsf B_t\mathsf O_t^{-1/2}).
\]
For $\alpha\ge\alpha_*$, set $\mathsf M_t=\alpha(I-\mathsf O_t)+\mathsf B_t$.
It satisfies equality in \eqref{base:new:Mrec} and the ceiling. The calculation costs $O(hk^3)$ per chain, with shared suffixes reusable.

The directional trace improves the actual-source bound:
\begin{equation}\label{base:new:Mtrace}
 \E\tr(Wa^Ta)\le\frac{\tr\mathsf M_0}{d}\beta+
 \frac{c_0k}{q-3k}\gamma.
\end{equation}
The matrix $\mathsf M_0$ is deterministic, so $\E N^TN\preceq\beta I/d$ suffices without independence. Conditional off-chain centering removes mixed terms. The final factor is $5+4\max\{\sup\tr\mathsf M_0/d,c_0k/(q-3k)\}$.

\subsection{Arbitrary levels separated by balanced levels}
At width $4k$, write $\mathsf S_t=\alpha\mathsf O_t-\mathsf B_t$. A balanced step has $L_V=I/2$, $\tau=1/4$, and
\[
 \mathsf S_t=(3\alpha/8-5/4)I+(3/8)O\mathsf S_{t+1}O^T.
\]
An arbitrary step with $\mathsf S_{t+1}\succeq sI$, $0\le s\le\alpha$, gives
$\mathsf S_t\succeq(s/2-2)I$. This follows by checking the endpoint eigenvalues of $0\preceq L_V\preceq I$ and using $\tau\le1/2$. A balanced step supplies $4I$ at $\alpha=14$, enough for one arbitrary preceding step. If no source path has two consecutive arbitrary nodes, $\mathsf M_t\preceq14I$ at all depths and
\begin{equation}\label{base:new:33}
 \E\F{A-\widehat B}^2\le33d_{\mathcal S}(A)^2.
\end{equation}
If the longest run of arbitrary nodes is $b$, including endpoint runs, one may use
$\alpha_b=(32\,2^b-22)/3$, giving factor $\max\{13,5+2\alpha_b\}$. The number of repetitions is unlimited; only run length is charged.

\subsection{The scalar-storage obstruction becomes neutral}
For even rank let $P$ have rank $k/2$ and $Q=I-P$. Alternate the two projector completions that retain $P$ and $Q$. The matrices
\[
 \alpha=34/9,\quad
 \mathsf M_P=(34/9)P+(22/9)Q,\quad
 \mathsf M_Q=(22/9)P+(34/9)Q
\]
satisfy \eqref{base:new:Mrec} exactly. For the rank-two witness above,
\[
 \E\mathcal H_{t+1}=(\alpha+2)/4+(3/4)(22/9)+1/2=34/9=\mathcal H_t.
\]
The direction surviving now receives a smaller later weight because the next step discards it. The trace source factor is $14/9$, less than the off-chain factor two. This calculation concerns alternating ancestor chains; not every branch of an arbitrary selector tree has that pattern.

\section{Noncommuting directional refresh without balanced nodes}\label{base:new:refresh}
For $k=2n$, let
\[
 P^{(0)}=\begin{pmatrix}I_n&0\\0&0\end{pmatrix},\qquad
 P^{(1)}=\frac12\begin{pmatrix}I_n&I_n\\I_n&I_n\end{pmatrix}.
\]
At even levels choose either $P^{(0)}$ or its complement; at odd levels choose either $P^{(1)}$ or its complement, independently for row and column bases. Use the retained basis $[P;I-P]$ and completion $\bigl(\begin{smallmatrix}P&-(I-P)\\I-P&P\end{smallmatrix}\bigr)$. Every child contraction is singular, but every consecutive pair along every path satisfies
\[
 P_tP_{t+1}P_t=P_t/2.
\]
The weights
\[
 \alpha=50/9,\qquad
 \mathsf M_t=(50/9)P_t+(26/9)(I-P_t)
\]
solve \eqref{base:new:Mrec}. Their normalized trace is $38/9$; at $d=2k$ the chain-source factor is $19/9$. Therefore
\begin{equation}\label{base:new:121}
 \E\F{A-\widehat B}^2\le(121/9)d_{\mathcal S}(A)^2.
\end{equation}
For odd rank add a scalar balanced coordinate, with weight $38/9$; the same normalized trace works. Independent fixed global orthogonal conjugations of the row and column families are allowed. Unrelated local coordinate changes cannot be assumed to preserve the parent--child product identity.

More generally, for rank-half projectors with $P_tP_{t+1}P_t\preceq\gamma P_t$, $\gamma<1$, weights
$a=2+16/[9(1-\gamma)]$, $b=2+4/[9(1-\gamma)]$ give factor
$\max\{13,9+20/[9(1-\gamma)]\}$. This is a multilevel geometric condition, not a local least-eigenvalue condition.

\subsection{A uniform perturbation allowance}
Let the actual retained child blocks be within operator norm $\varepsilon\le1/50$ of this reference family, in both orientations; the actual bases must still be isometries. Double the reference weights and use $\widetilde\alpha=100/9$. Doubling creates margin $\mathsf Q^{(0)}\succeq I$. Since child blocks are contractions,
\[
 \|L_V-L_V^{(0)}\|\le2\varepsilon,\quad
 |\tau-1/4|\le\varepsilon,
\]
and the right side of \eqref{base:new:Mrec} changes by at most
$[(7/2)\widetilde\alpha+5]\varepsilon=(395/9)\varepsilon$.
At $1/50$, margin $11I/90$ remains. Every level absorbs its own perturbation, with no sum over depth. The normalized trace doubles and gives
\begin{equation}\label{base:new:197}
 \E\F{A-\widehat B}^2\le(197/9)d_{\mathcal S}(A)^2<22d_{\mathcal S}(A)^2.
\end{equation}
The distance is relative to the actual perturbed fixed space, not the reference space.

\subsection{Transported observability, not the sum of displayed losses}
For \eqref{base:new:OB}, telescoping $I-\mathcal T_t(I)=\tau_tI+(1-\tau_t)L_{V,t}$ gives
\begin{equation}\label{base:new:OBsandwich}
 4(I-\mathsf O_t)\preceq\mathsf B_t
 \preceq4(I-\mathsf O_t)+2\rho\mathsf O_t.
\end{equation}
If $\varepsilon_h=\min_t\lambda_{\min}\mathsf O_t$, then
$4(\varepsilon_h^{-1}-1)\le\alpha_*\le4(\varepsilon_h^{-1}-1)+2\rho$.
Thus boundedness of this storage construction is equivalent to a positive reserve of transported loss.

A sufficient window condition is: every $r$-step product
$F_{t,r}=A_{V,t}\cdots A_{V,t+r-1}$ has squared norm at most $1-\eta$, and the scalar products $\prod(1-\tau)$ over full and terminal truncated windows are at least $\sigma>0$. Then
\[
 \mathsf O_t\succeq\varepsilon I,\qquad
 \varepsilon=\frac{\sigma\eta}{1-\sigma+\sigma\eta}.
\]
Indeed $\sum_{j<r}F_{t,j}L_{V,t+j}F_{t,j}^T=I-F_{t,r}F_{t,r}^T$. A future reserve $\varepsilon I$ produces at least $\sigma[\eta+(1-\eta)\varepsilon]I=\varepsilon I$. Backward induction starts from the root. At width $4k$ the factor is $\max\{13,8/\varepsilon-1\}$.

An untransported sum of losses is insufficient. For $A_0=e_1e_2^T$, $A_1=e_2e_1^T$, one has $L_0+L_1=I$ but $L_0+A_0L_1A_0^T=L_0$ and $\|A_0A_1\|=1$. The surviving direction changes coordinates rather than disappearing.

\section{Protected directions and the obstruction to forced refresh}\label{base:new:protected}
Let $F_h=A_{V,0}\cdots A_{V,h-1}$ and let $\Pi$ project onto $\ker(I-F_hF_h^T)$. For $z_0\in\operatorname{ran}\Pi$ define $z_{s+1}=A_{V,s}^Tz_s$. Equality of norm through a product of contractions forces equality at every factor; hence $A_{V,s}z_{s+1}=z_s$ and $C_{V,s}z_{s+1}=0$. Equation~\eqref{base:new:reverse} gives $z_s^Tg_s=z_0^Tg_0$. Every sibling is orthogonal to that direction, so
\begin{equation}\label{base:new:protected-source}
 t_rg_0^T\Pi=0\text{ for every }r,\qquad N(G)\Pi=0\quad\text{pathwise}.
\end{equation}
A source $a=a\Pi$ supported entirely in a protected direction has $Y=0$ and $B_V^TZ=0$ at every step, hence emits no nonroot coefficients. Root averaging gives
$\Pi W(G)\Pi\preceq\Pi/4$. No vanishing of mixed response blocks is asserted.

These identities prevent a mistaken learning proposal: impose strict contraction in every direction within a fixed number of levels. The exactly HSS rank-one matrix $e_ie_j^T$, with $i,j$ in opposite root children, must retain the localized row direction along the path to $i$ and column direction along the path to $j$. Those directions lose no norm, otherwise exact representation is impossible. Since the optimal error is zero, a universal multiplicative guarantee cannot forcibly discard them. This is an obstruction to a strict-refresh learning rule, not to refitting: exact protected directions are harmless in \eqref{base:new:protected-source}. Nearly protected directions require the joint averaging studied next.


\section{Directional source excitation and near-protection}\label{base:new:near}
For a chain of $h$ nonroot transitions, write
\[
 N(G)=\sum_{s=0}^{h}B_{s+1}t_{s+1}g_0^T,\quad
 F_0=I,\ F_s=A_0\cdots A_{s-1},\quad D_s=I-F_sF_s^T,
\]
where $A_s=A_{V,s}$ and $D_s=0$ at negative indices. Put
\begin{equation}\label{base:new:CLdef}
 \Cov=\E N^TN,\qquad
 \Loss=\frac1d\sum_{s=0}^{h}\F{B_{s+1}}^2D_{s-1}.
\end{equation}
The exact covariance recursion and its sharp lower comparison are proved in Section~\ref{base:new:coercive}. In particular $0\preceq\Cov\preceq\Loss$, with the two first source increments identically zero. For a deterministic projector $\Pi$, the initial loss profile is
$V_\Pi=\tr(\Pi\Loss)$, so $\E\F{N\Pi}^2\le V_\Pi$.

The first fourth-moment route used $\E\F{TG}^4\le3\F T^4$ and obtained
\[
 (\E\F{N\Pi}^4)^{1/2}\le3\sqrt3 V_\Pi.
\]
Expanding the positive response trace square under \eqref{base:new:response-tail} gives
$\E\tr[(\Pi W\Pi)^2]\le2\kappa c_0\rank\Pi$. Hilbert--Schmidt Cauchy--Schwarz, with no independence assumption, then gave
\begin{equation}\label{base:new:first-near}
 \E\tr(W(N\Pi)^T(N\Pi))\le3\sqrt{6\kappa c_0r}\,V_\Pi.
\end{equation}
If $\Pi D_{h-1}\Pi\preceq\varepsilon\Pi$, then $V_\Pi\le\beta\varepsilon r/d$. At width $4k$, \eqref{base:new:first-near} is at most $28.804\beta\varepsilon r^{3/2}/k$. This was the first direct control of the actual dependent product in nearly protected directions. Its extra $\sqrt r$ was later replaced by a logarithm, but the explicit bound remains useful.

For the rescaled amplifier in Section~\ref{base:new:amplifier}, take the odd $n\ge4\pi^2M^2$. Its loss satisfies $1-c^{2nM}\le2\pi^2M/n\le1/(2M)$. Applying \eqref{base:new:first-near} with $k=r=1$ gives actual weighted response at most $14.402\beta/M$, while the conditional unit-source response is at least $M/16$. Source suppression and exceptional response amplification are different phenomena and must be kept in the same expectation.


\endgroup

% ===== MOMENTS =====
\begingroup
\let\R\relax
\newcommand{\R}{\mathbb R}
\let\E\relax
\newcommand{\E}{\mathbb E}
\let\Prb\relax
\newcommand{\Prb}{\mathbb P}
\let\T\relax
\newcommand{\T}{\mathsf T}
\let\F\relax
\newcommand{\F}[1]{\left\lVert#1\right\rVert_{F}}
\let\N\relax
\newcommand{\N}[1]{\left\lVert#1\right\rVert}
\let\ip\relax
\newcommand{\ip}[2]{\left\langle#1,#2\right\rangle_F}
\let\tr\relax
\newcommand{\tr}{\operatorname{tr}}
\let\rank\relax
\newcommand{\rank}{\operatorname{rank}}
\let\diag\relax
\newcommand{\diag}{\operatorname{blockdiag}}
\let\range\relax
\newcommand{\range}{\operatorname{range}}
\let\OPT\relax
\newcommand{\OPT}{\operatorname{OPT}}
\let\HSS\relax
\newcommand{\HSS}{\operatorname{HSS}}
\let\diver\relax
\newcommand{\diver}{\operatorname{div}}
\let\Var\relax
\newcommand{\Var}{\operatorname{Var}}
\let\UC\relax
\newcommand{\UC}{\mathcal U}
\let\VC\relax
\newcommand{\VC}{\mathcal V}
\let\err\relax
\newcommand{\err}{\mathcal E}
\let\proofstep\relax
\newcommand{\proofstep}[1]{\par\addvspace{.45\baselineskip}\Needspace{3\baselineskip}\textbf{#1}\par\nobreak}
\let\op\relax
\newcommand{\op}[1]{\left\lVert#1\right\rVert_{\mathrm{op}}}
\let\Sn\relax
\newcommand{\Sn}[2]{\left\lVert#1\right\rVert_{S_{#2}}}
\let\Cov\relax
\newcommand{\Cov}{\mathsf C}
\let\Loss\relax
\newcommand{\Loss}{\mathsf L}
\let\Fil\relax
\newcommand{\Fil}{\mathcal F}
\let\Id\relax
\newcommand{\Id}{I}
\let\supp\relax
\newcommand{\supp}{\operatorname{supp}}
\let\vecop\relax
\newcommand{\vecop}{\operatorname{vec}}
\let\Gammafun\relax
\newcommand{\Gammafun}{\operatorname{Gamma}}
\chapter{Original-law moments and covariance-sensitive certificates}\label{ch:moments}
\begin{scopebox}
These moment and symmetry certificates are intermediate results. Their logarithmic losses and source-spread hypotheses are not imposed on the final fixed-space refit theorem.
\end{scopebox}
\section{Whole-probe chi-square comparison}\label{base:new:chi}
Define
\[
 \chi_{d,p}=\prod_{j=0}^{p-1}(1+2j/d),\qquad\chi_{d,0}=1.
\]
These are the moments of $\chi_d^2/d$. They improve the scalar Gaussian allowance $(2p-1)!!$ by retaining the complement dimension.

\begin{theorem}[Original-probe moment envelope]\label{base:new:probe-chi}
For any deterministic compatible $T$ and every integer $p\ge1$,
\[
 \E\F{TG}^{2p}\le\chi_{d,p}\F T^{2p},\qquad
 \E\exp\left(\theta\frac{\F{TG}^2}{\F T^2}\right)
 \le(1-2\theta/d)^{-d/2}\quad(0\le\theta<d/2).
\]
At $p=1$ equality holds. The assertion at $T=0$ is interpreted directly.
\end{theorem}
\begin{proof}
In multiresolution coordinates the probe consists of a root frame and jointly generated discarded sibling pairs. A deterministic linear combination of a discarded pair has deterministic radius equal to its coefficient norm and is uniform on a sphere in a $d$-dimensional complement, conditional on previously generated ancestors. Pairs in different branches can be exposed in top-down order. Their coefficient radii are deterministic and their squared sum, with the root contribution, is the squared norm of the original deterministic row.

For a past-measurable vector $x$ and a sphere increment $\xi$ of radius $a$, expand $(\|x\|^2+2\langle x,\xi\rangle+a^2)^p$. Odd spherical moments vanish; all remaining coefficients are nonnegative. Replace the projection of $x$ into the sphere's subspace by a vector of norm $\|x\|$, and then replace the fixed radius by an isotropic Gaussian radius with the same second moment. Higher even radius moments increase. The resulting Gaussian radial polynomial is
\[
 \E[\|x+\xi\|^{2p}\mid\text{past}]
 \le\sum_{j=0}^{p}\binom pj\frac{\chi_{d,p}}{\chi_{d,p-j}}
 a^{2j}\|x\|^{2(p-j)}.
\]
Its coefficients can be checked by the noncentral Gaussian radial expansion or repeated differentiation of $(1-2\theta a^2/d)^{-d/2}\exp(\theta\|x\|^2/(1-2\theta a^2/d))$. If the past budget is $b$, insertion of the induction hypothesis makes the sum exactly $\chi_{d,p}(a^2+b)^p$. The root's deterministic norm starts the induction. Finally apply Minkowski to squared row norms of $TG$. The Laplace bound follows from its nonnegative power series; no stochastic domination is inferred.
\end{proof}
In particular,
\[
 \Var(\F{TG}^2)\le\frac2d\F T^4,
\]
\begin{equation}\label{base:new:probe-tail}
 \Prb\left\{\F{TG}^2>\F T^2(1+2\sqrt{x/d}+2x/d)\right\}\le e^{-x}.
\end{equation}
Only an upper tail is asserted. A one-sided MGF comparison does not justify a lower tail or convex-order statement.

\section{Source tails, fourth moments, and confidence dependence}\label{base:new:source-moments}
For deterministic vectors $u,z$, define
\[
 v(u,z)=\frac1d\sum_{s=0}^{h}\|u^TB_{s+1}\|^2z^TD_{s-1}z.
\]
Conditional spherical regression makes each scalar source increment sub-Gaussian with random variance allowance
$a_s=d^{-1}\|u^TB_{s+1}\|^2\|P_sg_0^Tz\|^2$, where $P_s$ is the available complement. The residual projection is bounded by a deterministic geometric residual applied to the complete probe. Theorem~\ref{base:new:probe-chi}, and convexity for a sum of variance allowances, give
\[
 \E e^{\theta\sum_s a_s}\le(1-2\theta v/d)^{-d/2}.
\]
The exponential supermartingale satisfies
$\E\exp(2\lambda u^TNz-2\lambda^2\sum_sa_s)\le1$.
Cauchy--Schwarz therefore proves
\begin{equation}\label{base:new:mgf}
 \boxed{\E e^{\lambda u^TNz}\le
 (1-4\lambda^2v(u,z)/d)^{-d/4},\qquad4\lambda^2v<d.}
\end{equation}
The random accumulated source and variance are not separated. The envelope is the MGF of $\sqrt{8v\Gamma/d}\,g$, with $\Gamma\sim\operatorname{Gamma}(d/4,1)$ and independent $g\sim N(0,1)$. Its variance is $2v$, whereas the actual variance is at most $v$. Chernoff's method gives
\begin{equation}\label{base:new:source-tail}
 \Prb\{|u^TNz|>2\sqrt v(\sqrt x+x/\sqrt d)\}\le2e^{-x}.
\end{equation}
The sub-Gaussian regime extends through orders comparable to $d$.

For every deterministic matrix $T$ and integer $j\ge1$, put $V_T=\tr(T^T\Loss T)$. One nonnegative coefficient of the cosh series in \eqref{base:new:mgf}, optimized at $4\lambda^2v/d=4j/(d+4j)$, gives
\begin{equation}\label{base:new:source-allmom}
 (\E\F{NT}^{2j})^{1/j}\le a_{j,d}V_T,
\quad
 a_{j,d}=\frac4d[(2j)!]^{1/j}
 (1+d/(4j))(1+4j/d)^{d/(4j)}
 \le2ej(1+4j/d).
\end{equation}
To pass from bilinear forms to row norms, integrate the test vector against an independent Gaussian, cancel its even moment $(2j-1)!!$, and use convexity over the eigenvalues of each deterministic row variance matrix. Minkowski then sums squared rows. This argument bounds one coefficient using a nonnegative series; it does not assume that pointwise ordered MGFs have coefficientwise ordered Taylor series. The earlier dimension-free bound $15j^2$ is recovered coarsely; the new growth is $O(j+j^2/d)$.

\subsection{Centering the random variance improves the fourth-moment constant}
Let $\delta_d=\sqrt{2/d}$,
\[
 u_d=\frac{3\delta_d+\sqrt{9\delta_d^2+8}}2,
 \qquad c_d=\sqrt{1+u_d^2}.
\]
Then
\begin{equation}\label{base:new:source-fourth}
 \boxed{(\E\F{NT}^4)^{1/2}\le c_d\tr(T^T\Loss T).}
\end{equation}
Here $c_2=3.699278\ldots$, $c_4=3$, and $c_d\to\sqrt3$. To verify the induction, let $m_s=\E\F{S_s}^2$ and $x_s=(\E\F{S_s}^4)^{1/2}$ for the accumulated source. If $U_s$ is the residual upper bound on a conditional variance term, Theorem~\ref{base:new:probe-chi} gives $\E U_s=b_s$ and $\Var U_s\le(2/d)b_s^2$. Centering before Cauchy--Schwarz gives
\[
 \E[\F{S_s}^2U_s]\le
 b_s[m_s+\delta_d\sqrt{x_s^2-m_s^2}].
\]
The local Haar fourth moment contributes at most three times its squared deterministic variance budget. The ansatz $m_s\le V_s$, $x_s\le c_dV_s$ closes because
$c_d^2=3[1+\delta_d\sqrt{c_d^2-1}]$; expansion of $(V_s+b_s)^2$ handles the new pure fourth term. This proves \eqref{base:new:source-fourth} with the actual dependent conditional variances.

\section{Response moments via a signed square-root martingale}\label{base:new:response-moments}
Assume the response-tail estimate of Lemma~\ref{base:new:restart}. Introduce independent signs, independent of the entire original experiment, and set
\[
 Z_*=\sum_n\epsilon_nW_n^{1/2}.
\]
Interleave revealing the original data for $W_n$ with revealing $\epsilon_n$. This is a self-adjoint martingale; its only nonzero increments are $\epsilon_nW_n^{1/2}$. Exactly,
\[
 \E_\epsilon Z_*^2=W,
\]
and its row and column BMO norms are at most $\sqrt\kappa$, because future signs remove all cross terms and the remaining conditional square is the response tail.

The public noncommutative John--Nirenberg theorem of Junge--Musat~\cite{JM}, in the multiplier formulation, states for a deterministic $B$ and $s>2$
\[
 \|Z_*B\|_{L_s(S_s)}\le Js\|Z_*\|_{\rm BMO}\|B\|_{S_s}
\]
with a universal constant $J$. This is the external high-moment input; no numerical value is supplied. Normalize the matrix trace on both sides and then undo the normalization to obtain this formula. Weighted trace convexity and the exact sign square give
\begin{equation}\label{base:new:Wmom}
 \left(\E\tr[(R^{1/2}WR^{1/2})^p]\right)^{1/p}
 \le4J^2\kappa p^2\Sn Rp,\qquad p\ge2,
\end{equation}
for every deterministic $R\succeq0$. At $p=2$, direct expansion of the trace square gives the sharper explicit bound
\begin{equation}\label{base:new:W2}
 \left(\E\tr[(R^{1/2}WR^{1/2})^2]\right)^{1/2}
 \le\sqrt{2\kappa c_0}\Sn R2.
\end{equation}
Indeed with $A_n=R^{1/2}W_nR^{1/2}$, condition each cross term at its earlier index, use conditional tail $\kappa R$ and mean $c_0R$, and obtain at most $2\kappa c_0\tr R^2$.

The auxiliary signs do not change either HSS probe. John--Nirenberg is applied to $Z_*$, not directly to the positive sum $W$. The square explains the $p^2$ growth. Section~\ref{base:new:abstract-obstruction} shows why response-tail information alone cannot generally replace it by $p$.

\section{A loss-spectrum certificate without spectral-band triangle loss}\label{base:new:spectral}
Let $R$ be deterministic and positive on the exact support of $\Cov$. Put
\[
 v_R=\tr(\Cov R^{-1}),\quad V_R=\tr(\Loss R^{-1}),\quad
 X=NR^{-1/2},\quad Y=R^{1/2}WR^{1/2}.
\]
The product of interest is $\tr(YX^TX)$. Schatten and probabilistic H\"older with $p'=p/(p-1)$, followed by interpolation of the exact source second moment and \eqref{base:new:source-fourth}, gives
\begin{equation}\label{base:new:metric-joint}
 \boxed{\E\tr(WN^TN)\le4J^2\kappa p^2\Sn Rp
 v_R^{1-2/p}(c_dV_R)^{2/p},\qquad p\ge2.}
\end{equation}
If $v_R=0$, the source is zero and no ratio is formed. At order two, \eqref{base:new:W2} gives
$c_d\sqrt{2\kappa c_0}\Sn R2\tr(\Loss R^{-1})$.

The exact metric identity
\begin{equation}\label{base:new:one-metric}
 \inf_{R\succ0}\Sn Rp\tr(LR^{-1})=\Sn L{p/(p+1)}
\end{equation}
holds for positive $L$, with minimizer $R=L^{1/(p+1)}$ up to scale. Diagonalizing $R$ in an eigenbasis of $L$ cannot increase either factor: use convexity of $x^p$ on its diagonal entries and $(R^{-1})_{ii}\ge1/R_{ii}$. The remaining scalar inequality is ordinary H\"older. Thus
\begin{equation}\label{base:new:loss-profile}
 \E\tr(WN^TN)\le
 \min\left\{c_d\sqrt{2\kappa c_0}\Sn\Loss{2/3},\quad
 4J^2\kappa\inf_{p\ge2}p^2c_d^{2/p}\Sn\Loss{p/(p+1)}\right\}.
\end{equation}
For $0<a<1$, $\Sn La=(\tr L^a)^{1/a}$ is a Schatten quasi-norm. The explicit branch never worsens the earlier unweighted fourth-moment certificate, since $\Sn L{2/3}\le\sqrt{\rank L}\tr L$. One metric treats all loss directions together; no assumption of orthogonal spectral-sector outputs is made.

For a rank-$r$ nearly protected subspace with $\Pi\Loss\Pi\preceq(\beta\varepsilon/d)\Pi$, define
\[
 \psi(x)=\inf_{p\ge2}p^2x^{1/p}
 =\begin{cases}4\sqrt x,&1\le x\le e^4,\\(e^2/4)(\log x)^2,&x\ge e^4.\end{cases}
\]
The second branch of \eqref{base:new:loss-profile} is at most
$4J^2\kappa(\beta\varepsilon r/d)\psi(c_d^2r)$.
At width $4k$ this has scale
\begin{equation}\label{base:new:near-log}
 C\beta\varepsilon\frac rk(1+\log r)^2.
\end{equation}
The extra $\sqrt r$ in \eqref{base:new:first-near} has been removed, but a logarithm remains. For a full-rank nearly protected sector the sufficient small-loss scale improves from $k^{-1/2}$ to $(1+\log k)^{-2}$. This is a source-level sufficient condition, not a claim about every source in every tree.

\subsection{Exact noncommuting covariance--loss optimization}
For $p>2$, let $a=1-2/p$, $b=2/p$, $\theta=p/(p+1)$, and define
\[
 \mathfrak G_p(C,L)=\inf_{R\succ0}\Sn Rp
 [\tr(CR^{-1})]^a[\tr(LR^{-1})]^b.
\]
The identity $u^av^b=\min_{t>0}(at^bu+bt^{-a}v)$ and \eqref{base:new:one-metric} give
\begin{equation}\label{base:new:mixed-metric}
 \boxed{\mathfrak G_p(C,L)=\min_{t>0}\Sn{at^bC+bt^{-a}L}{\theta}.}
\end{equation}
No simultaneous diagonalization is needed. At an optimizing $t_*$,
$R_*=(at_*^bC+bt_*^{-a}L)^{1/(p+1)}$ up to scale, and
\[
 \lambda_{\min}(C^{-1/2}LC^{-1/2})\le t_*\le
 \lambda_{\max}(C^{-1/2}LC^{-1/2}).
\]
The bracket follows because the optimal scalar for fixed $R$ is
$\tr(LR^{-1})/\tr(CR^{-1})$. The logarithm of the objective is convex in $\log t$: write its square-root factor as the vertical stack
$[\sqrt a\,e^{bz/2}C^{1/2};\sqrt b\,e^{-az/2}L^{1/2}]$ and apply Schatten interpolation; imaginary shifts are block-unitary. The norm exponent of that stack is $2\theta>1$, so ordinary norm interpolation applies despite $\theta<1$ for the positive matrix itself.

Replacing the second branch of \eqref{base:new:loss-profile} by
\[
 4J^2\kappa\inf_{p\ge2}p^2c_d^{2/p}\mathfrak G_p(\Cov,\Loss)
\]
retains the exact covariance. At $p=2$, set $\mathfrak G_2(C,L)=\Sn L{2/3}$. If $C=\eta L$, the value is $\eta^{1-2/p}\Sn L\theta$ and $t_*=1/\eta$. The final sharp covariance theorem proves that, for actual HSS sources at width $4k$, $\eta$ cannot be smaller than $1/4$ in Loewner order. This turns a possible asymptotic route into a bounded constant refinement.

\subsection{Conditional-risk tails, not realized-error tails}
For rank-$r$ projector $\Pi$, let $V=\tr(\Pi\Loss\Pi)$,
$p_\delta=\max\{2,\log(2r/\delta)\}$ and
$j_\delta=\max\{1,\lceil\log(2/\delta)\rceil\}$. Markov applied to \eqref{base:new:Wmom} and \eqref{base:new:source-allmom}, and a union bound not requiring independence, show with $G$-probability at least $1-\delta$ that
\begin{equation}\label{base:new:cond-risk}
 \E_H[\F{\mathcal O(N(G)\Pi)}^2\mid G]
 \le4e^2J^2\kappa Vp_\delta^2\widehat a_{j_\delta,d}.
\end{equation}
Here $\widehat a$ uses the exact second-moment coefficient one and the fourth-moment coefficient $c_d$ at orders one and two. Its growth is
$O(\kappa J^2V\log^2(r/\delta)\log(1/\delta)[1+\log(1/\delta)/d])$.
This is a high-probability statement about conditional expected source risk, not an equal-probability bound for the realized two-sketch output. The extra argument needed to control the latter is not supplied by this estimate.


\section{Sharp original-law covariance coercivity}\label{base:new:coercive}
This theorem has a direct proof from the reverse frame law. It does not use the response-tail lemma, Carleson embedding, the signed response moment bound, or any claim about a learned basis.

Fix deterministic orthogonal completions along an arbitrary source path. Let $A_s$ and $C_s^{\rm comp}$ be their upper and lower retained blocks, so
$A_s^TA_s+(C_s^{\rm comp})^TC_s^{\rm comp}=I$. Put
\[
 F_0=I,\quad F_s=A_0\cdots A_{s-1},\quad
 D_s=I-F_sF_s^T,\qquad D_s=0\ (s<0),
\]
\[
 N=\sum_{s=0}^{h}B_{s+1}t_{s+1}g_0^T,\quad
 \Cov=\E N^TN,\quad
 \Loss=\frac1d\sum_{s=0}^{h}\F{B_{s+1}}^2D_{s-1}.
\]
All source blocks are deterministic, and all random frames follow the original reverse construction.

\begin{theorem}[Sharp covariance comparison]\label{base:new:sharp-C}
For $q\ge4k$, $d=q-2k$, $\rho=k/d$, and $\eta=(1-\rho)^2$,
\begin{equation}\label{base:new:sharp-main}
 \boxed{\eta\Loss\preceq\Cov\preceq\Loss.}
\end{equation}
Both endpoints are attained by actual frame chains at every rank and admissible width. In particular, at $q=4k$,
\[
 \frac14\Loss\preceq\Cov\preceq\Loss,\qquad
 \ker\Cov=\ker\Loss.
\]
\end{theorem}

\subsection{The exact second-moment map}
Put $X=g_0$,
$P_s=I-g_s^Tg_s-w_s^Tw_s$, and
\[
 \Xi_s=\E[XP_sX^T],\quad
 \mathcal C_s(H)=\E[(Xg_s^T)H(g_sX^T)].
\]
Then $\mathcal C_0(H)=H$, $\Xi_0=0$, and
\begin{align}
 \mathcal C_{s+1}(H)
 &=\mathcal C_s(A_sHA_s^T)
 +\frac{\tr(C_s^{\rm comp}H(C_s^{\rm comp})^T)}d\Xi_s,
 \label{base:new:Cs}\\
 \Xi_{s+1}&=I-\mathcal C_s(I)-\rho\Xi_s.\label{base:new:Xirec}
\end{align}
The first identity follows by expanding \eqref{base:new:reverse}; the mixed terms vanish by conditional centering and the fresh square is Lemma~\ref{base:new:frame-mom}. The second follows because $[g_{s+1};w_{s+1}]$ has the same row space as $[g_s;t_{s+1}]$.

The expectation $\E[Xg_s^T]=F_s$. Jensen therefore gives
$\mathcal C_s(I)\succeq F_sF_s^T$, and \eqref{base:new:Xirec} implies
\begin{equation}\label{base:new:Xiupper}
 0\preceq\Xi_{s+1}\preceq D_s,
 \qquad\Xi_0=\Xi_1=0.
\end{equation}
The exact covariance is
\begin{equation}\label{base:new:exactC}
 \Cov=\frac1d\sum_{s=0}^{h}\F{B_{s+1}}^2\Xi_s.
\end{equation}
Indeed, different source increments are orthogonal martingale differences, and the conditional right covariance of the $s$th increment is
$d^{-1}\F{B_{s+1}}^2XP_sX^T$.

\subsection{Unroll the covariance memory before bounding it}
Unrolling \eqref{base:new:Cs} at $H=I$ gives
\begin{equation}\label{base:new:memory}
 \mathcal C_s(I)=F_sF_s^T+K_s,\quad
 K_s=\sum_{j<s}a_{s,j}\Xi_j,\quad
 a_{s,j}=\frac1d\|C_j^{\rm comp}A_{j+1}\cdots A_{s-1}\|_F^2.
\end{equation}
Every coefficient is nonnegative. Moreover,
\begin{equation}\label{base:new:memory-total}
 \sum_{j<s}a_{s,j}=\frac{k-\|F_s\|_F^2}{d}\le\rho.
\end{equation}
To verify the equality, let $R_j=A_{j+1}\cdots A_{s-1}$ and use
$\|C_j^{\rm comp}R_j\|_F^2=\|R_j\|_F^2-\|A_jR_j\|_F^2$; the sum telescopes. Extending a product on the right by a contraction decreases its Frobenius norm, so
\[
 a_{s,j}\le a_{s-1,j}\ (j<s-1),\qquad a_{s,s-1}\le\rho.
\]
No covariance matrices need commute.

From \eqref{base:new:Xiupper} and \eqref{base:new:memory-total},
\begin{equation}\label{base:new:memory-bd1}
 K_s\preceq\rho D_{s-2}.
\end{equation}
From the horizon monotonicity and \eqref{base:new:Xirec},
\begin{equation}\label{base:new:memory-bd2}
 K_s\preceq K_{s-1}+\rho\Xi_{s-1}=D_{s-1}-\Xi_s.
\end{equation}
Now \eqref{base:new:Xirec} becomes $\Xi_{s+1}=D_s-K_s-\rho\Xi_s$. The two bounds yield
\[
 \Xi_{s+1}\succeq D_s-\rho D_{s-2}-\rho\Xi_s,
 \qquad
 \Xi_{s+1}\succeq D_s-D_{s-1}+(1-\rho)\Xi_s.
\]
Take weights $1-\rho$ and $\rho$. The intermediate matrix $\Xi_s$ cancels exactly:
\begin{equation}\label{base:new:age-Xi}
 \Xi_{s+1}\succeq D_s-\rho D_{s-1}-\rho(1-\rho)D_{s-2}.
\end{equation}
The right side equals
\[
 (D_s-D_{s-1})+(1-\rho)(D_{s-1}-D_{s-2})+(1-\rho)^2D_{s-2}
 \succeq(1-\rho)^2D_s.
\]
The loss matrices are monotone because $A_s$ is a contraction. Combining with \eqref{base:new:Xiupper} and summing \eqref{base:new:exactC} proves Theorem~\ref{base:new:sharp-C}.

The first naive memory bound alone would have produced $1-2\rho$, zero at width $4k$. The successful step uses both valid bounds on the same memory matrix and cancels the intermediate covariance before scalarization.

\subsection{An age-resolved source lower profile}
Define
\[
 \Loss_{\rm age}=\frac1d\sum_{s=0}^{h}\F{B_{s+1}}^2
 [D_{s-1}-\rho D_{s-2}-\rho(1-\rho)D_{s-3}].
\]
Each bracket is positive semidefinite by the loss-increment decomposition above, and
\begin{equation}\label{base:new:ageC}
 \eta\Loss\preceq\Loss_{\rm age}\preceq\Cov\preceq\Loss.
\end{equation}
Newly lost directions have coefficient one, loss one step older has coefficient $1-\rho$, and older loss has coefficient $(1-\rho)^2$. Keeping the finite age profile can be more informative than collapsing immediately to its smallest coefficient.

\subsection{Sharpness and a realizable residual}
Take every completion to be the full swap
$Q_s=\bigl(\begin{smallmatrix}0&-I\\I&0\end{smallmatrix}\bigr)$.
Then $A_s=0$, $D_s=I$ for $s\ge1$. Writing $\mathcal C_s(I)=u_sI$ and $\Xi_s=x_sI$ gives
\[
 u_0=1,\ x_0=0,\qquad u_{s+1}=\rho x_s,
 \quad x_{s+1}=1-u_s-\rho x_s.
\]
Hence $(x_0,x_1,x_2,x_3,x_4)=(0,0,1,1-\rho,(1-\rho)^2)$. A source with only $B_3\ne0$ attains $\Cov=\Loss$; a source with only $B_5\ne0$ attains $\Cov=\eta\Loss$.

These arrays arise from $E=(\widetilde u_jb)(v_{s_5}z)^T$ for an initial discarded row basis and the relevant exterior-sibling retained column basis. Disjoint exterior support makes the residual orthogonal to discrepancies at and below the source; its discarded row direction is orthogonal to ancestor retained row spaces and the root. Thus $E\perp\mathcal S$. Sharpness concerns the source covariance, not a lower bound on reconstruction error.

\subsection{The full vectorized covariance}
With column-major vectorization define
\[
 \Sigma_N=\E[\vecop N(\vecop N)^T],\qquad
 \Gamma_N=\frac1d\sum_{s=0}^{h}D_{s-1}\otimes(B_{s+1}B_{s+1}^T).
\]
Entrywise conditional frame covariance gives
\[
 \Sigma_N=\frac1d\sum_s\Xi_s\otimes(B_{s+1}B_{s+1}^T),\qquad
 \eta\Gamma_N\preceq\Sigma_N\preceq\Gamma_N.
\]
The same sharp constants and age profile apply. Thus every deterministic linear relation in $N$ is characterized by $\Gamma_N$: $\langle T,N\rangle_F=0$ almost surely exactly when $\vecop(T)^T\Gamma_N\vecop(T)=0$.

\section{Consequences: covariance refinement has a bounded ceiling}\label{base:new:cov-consequences}
On the common exact support of $\Cov,\Loss$,
\[
 \operatorname{spec}(\Cov^{-1/2}\Loss\Cov^{-1/2})\subseteq[1,\eta^{-1}].
\]
The scalar search in \eqref{base:new:mixed-metric} therefore always lies in $[1,\eta^{-1}]$, hence in $[1,4]$ at width $4k$. For $p>2$,
\begin{equation}\label{base:new:metric-ceiling}
 \eta^{1-2/p}\Sn\Loss{p/(p+1)}
 \le\mathfrak G_p(\Cov,\Loss)\le\Sn\Loss{p/(p+1)}.
\end{equation}
This follows by applying the covariance sandwich to $\tr(\Cov R^{-1})$ before taking the infimum. The full-swap source attains the lower endpoint. At $p=2$ there is no exact-covariance improvement. Even after optimizing $p$, the exact-covariance profile can improve the loss-only version by at most four at width $4k$.

This settles the scope of a previously tempting route. The exact covariance may sharpen constants and preserve directions, but cannot be arbitrarily smaller than its geometric allowance. In an order-one-loss subspace the source does not universally vanish with rank. For example, if $\Loss=\beta I/(2k)$, then $\beta/8\le\E\F N^2\le\beta/2$ at width $4k$. Further asymptotic progress must use the source's alignment with $W(G)$, not only its unweighted covariance.

\subsection{Covariance-only moments, tails and a small-ball statement}
The source fourth moment becomes
\[
 (\E\F{NT}^4)^{1/2}\le(c_d/\eta)\tr(T^T\Cov T).
\]
Consequently a covariance-only response certificate is
\[
 \E\tr(WN^TN)\le\min\left\{
 \frac{c_d}{\eta}\sqrt{2\kappa c_0}\Sn\Cov{2/3},\quad
 4J^2\kappa\inf_{p\ge2}p^2(c_d/\eta)^{2/p}\Sn\Cov{p/(p+1)}
 \right\}.
\]
For $S=u^TNz$ with actual variance $\sigma^2$, its geometric allowance satisfies $\eta v\le\sigma^2\le v$. Hence
\[
 \E e^{\lambda S}\le(1-4\lambda^2\sigma^2/(d\eta))^{-d/4},
 \qquad
 \Prb\{|S|>(2\sigma/\sqrt\eta)(\sqrt x+x/\sqrt d)\}\le2e^{-x}.
\]
At $q=4k$ the prefactor is $4\sigma$. If $\sigma=0$, the variable is zero almost surely. Paley--Zygmund gives, for $v_T=\E\F{NT}^2>0$ and $0<\theta<1$,
\[
 \Prb\{\F{NT}^2\ge\theta v_T\}\ge(1-\theta)^2\eta^2/c_d^2.
\]
These are source statements, not approximation-error lower bounds. Their constants are not optimized.


\section{A fourth-Gram comparison and the coherent-source reduction}\label{base:con:gram}
The preceding covariance sandwich concerns $\E N^\T N$. A different statistic records the distribution of source energy across its left coordinates. For
\[
 N=\sum_s B_st_sX^\T,\qquad XX^\T=I_k,
\]
define
\[
 H_B=\sum_sB_sB_s^\T,\quad \beta=\tr H_B,
 \quad h_2=\tr(H_B^2),\quad r_2(H_B)=\beta^2/h_2\quad(\beta>0).
\]
The number of rows of the $B_s$ may be arbitrary; the retained frame rank is still $k$. For an HSS source with $k$ rows, $1\le r_2(H_B)\le k$.

\begin{theorem}[Gaussian comparison for two quartic functionals]\label{base:con:gram-thm}
For every deterministic compatible $T$, put $K=T^\T T$. Under the original reverse frame law,
\begin{align}
 \E\tr\{[(NT)^\T(NT)]^2\}
 &\le d^{-2}\{\beta^2\tr K^2+h_2[(\tr K)^2+\tr K^2]\},\label{base:con:gram-fourth}\\
 \E\F{NT}^4&\le d^{-2}\{\beta^2(\tr K)^2+2h_2\tr K^2\}.\label{base:con:frob-fourth}
\end{align}
The right sides are the corresponding Gaussian moments for
$d^{-1/2}H_B^{1/2}\mathcal G K^{1/2}$, with a standard rectangular Gaussian $\mathcal G$. The comparison is restricted to these two functionals.
\end{theorem}
\begin{proof}
Write $Z_s=B_st_sX^\T T$, $S_s=\sum_{j<s}Z_j$,
$H_s=B_sB_s^\T$, and
$A_s=T^\T XP_{s-1}X^\T T$. Conditional on the past, $Z_s$ is symmetric,
$0\preceq A_s\preceq K$, and
\[
 \E[\vecop Z_s(\vecop Z_s)^\T\mid\Fil_{s-1}]
 =d^{-1}A_s\otimes H_s.
\]
Therefore
$\E[\vecop S_s(\vecop S_s)^\T]\preceq
 d^{-1}K\otimes\sum_{j<s}H_j$.

For one fresh regression $Z=BOR$, with $O$ Haar on $O(d)$, put
$H=BB^\T$, $A=R^\T R$, $b=\tr H$, $h=\tr H^2$, $a=\tr A$, $v=\tr A^2$. Contracting the fourth Haar tensor gives
\[
 \E\tr[(Z^\T Z)^2]
 =\frac{d(ha^2+b^2v)+(d-2)hv-b^2a^2}{d(d-1)(d+2)}
 \le\frac{b^2v+h(a^2+v)}{d^2}.
\]
The analogous contraction for $\E\F Z^4$ is bounded by
$d^{-2}(b^2a^2+2hv)$. These inequalities follow also by subtracting the expressions and using $h\le b^2$, $v\le a^2$. Zero blocks follow by continuity.

Expanding $\tr[((S_s+Z_s)^\T(S_s+Z_s))^2]$ eliminates odd terms by conditional symmetry. The mixed terms are
\[
 \frac{2\tr H_s}{d}\tr(S_s^\T S_sA_s),\quad
 \frac2d\tr(S_s^\T H_sS_sA_s),\quad
 \frac{2\tr A_s}{d}\tr(H_sS_sS_s^\T).
\]
Replace $A_s$ by the deterministic upper bound $K$ inside these positive contractions, then use the covariance bound for $S_s$. This preserves $\tr(H_s\sum_{j<s}H_j)$ instead of replacing it by a product of scalar traces. Finally
\[
 \sum_s[b_s^2+2b_s\sum_{j<s}b_j]=\beta^2,
 \quad
 \sum_s[\tr H_s^2+2\tr(H_s\sum_{j<s}H_j)]=h_2.
\]
The Frobenius fourth-power expansion uses the same conditional covariance and the second local inequality. Both inductions have coefficient one on the previous fourth moment, proving the result without a per-level multiplicative loss.
\end{proof}

The comparison cannot be extended to all nonnegative quartics. A $2\times2$ determinant minor of a Haar $O(d)$ matrix has second moment $2/[d(d-1)]$, whereas independent $N(0,1/d)$ entries give $2/d^2$. Thus a blanket Gaussian replacement already fails at degree four.

\subsection{Source spread closes an all-input-basis, but not all-input, theorem}
The response-tail estimate gives
$\E\tr W^2\le2\kappa c_0k$. Hilbert--Schmidt Cauchy--Schwarz and \eqref{base:con:gram-fourth} with $T=I_k$ imply
\begin{equation}\label{base:con:gram-risk}
 \E\tr(WN^\T N)
 \le C_G\sqrt{\beta^2+(k+1)h_2},
 \qquad C_G=\rho\sqrt{2\kappa c_0}.
\end{equation}
The factors remain dependent. For a rank-$r$ deterministic projector $\Pi$, the corresponding allowance is
\[
 \E\tr(W(N\Pi)^\T(N\Pi))
 \le\frac r d\sqrt{2\kappa c_0}\,
 \beta\sqrt{1+(r+1)/r_2(H_B)}.
\]
At $q=4k$, if every nonzero source of the actual residual in both orientations has $r_2(H_B)\ge\theta k$, then
\begin{equation}\label{base:con:spread-global}
 \E\F{A-\widehat B}^2
 \le[5+4\sqrt{\kappa(1+2/\theta)}]d_{\mathcal S}(A)^2,
 \qquad\kappa=9(2+\sqrt2).
\end{equation}
For isotropic row energy $H_B=(\beta/k)I$, retain the sharper factor
$5+4\sqrt{\kappa(2+1/k)}<44$. At rank one this condition is automatic. At arbitrary rank, $r_2\ge1$ still gives the explicit general certificate
$5+4\sqrt{\kappa(k+2)}$, of order $\sqrt k$. The logarithmic-rank certificate is asymptotically smaller.

A centered variant uses $\Cov=\E N^\T N$ and
$M_4=k[\beta^2+(k+1)h_2]/d^2$:
\begin{equation}\label{base:con:centered-risk}
 \E\tr(WN^\T N)\le c_0\tr\Cov+
 \sqrt{k c_0(2\kappa-c_0)[M_4-\tr(\Cov^2)]}.
\end{equation}
Indeed subtract both means, use Cauchy--Schwarz, and use
$\E\|W-\E W\|_F^2\le k c_0(2\kappa-c_0)$, obtained by maximizing
$2\kappa\tr\E W-\tr(\E W)^2$ over $0\preceq\E W\preceq c_0I$.
One may aggregate these source-specific allowances with their actual energy budgets rather than maximize over sources. This produces residual-dependent certificates, not quantities claimed to be observed by the algorithm.

\subsection{The exact coherent hard case and a noisy auxiliary estimator}
Let $T_s=t_sX^\T$, $\mathsf X=[T_1;\ldots;T_m]$, and
$\mathcal K=\E[\mathsf XW\mathsf X^\T]\succeq0$. For $B=[B_1\ \cdots\ B_m]$,
\[
 R(B)=\tr(B\mathcal K B^\T),\qquad
 \sup_{B\ne0}R(B)/\F B^2=\op{\mathcal K}.
\]
An extremizer $B=uv^\T$ has one left source row and an arbitrary temporal vector $v\in\R^{mk}$. The retained rank remains $k$, so the $k=1$ theorem does not solve this reduction.

Choosing the leading $r$ eigenvectors of $\mathcal K$ as rows of $B$ in \eqref{base:con:gram-risk} proves
\begin{equation}\label{base:con:KyFan}
 \sum_{i\le r}\lambda_i(\mathcal K)\le C_G\sqrt{r(r+k+1)},
 \qquad \lambda_r(\mathcal K)\le C_G\sqrt{1+(k+1)/r}.
\end{equation}
Outside the leading $k$ temporal directions the norm is at most $\sqrt2C_G$. Their removed trace is at most $C_G\sqrt{k(2k+1)}$.

This suggested an auxiliary regression route. If an independent Gaussian sketch $b\Omega$ of the actual residual-source row were available, regress on the leading $k$-dimensional temporal eigenspace using $3k+2$ columns. The coordinate-estimation error has mean zero and covariance
$\|b_\perp\|^2I_k/(2k+1)$. Trace rather than maximum-eigenvalue charging would then give risk at most
\[
 \left(\sqrt2C_G+\frac{C_G\sqrt{k(2k+1)}}{2k+1}\right)\|b_\perp\|^2
 \le\frac3{\sqrt2}C_G\|b\|^2.
\]
This is an ideal-estimation lemma, not a legal new query algorithm: the source row belongs to $A-P_{\mathcal S}A$, and a sketch of $A$ can contain arbitrarily large in-space contamination. An exact-on-$\mathcal S$ simultaneous estimator has not been constructed. Adding rows to a coherent source also fails: preserving $B^\T B$ preserves the nonzero spectrum of $BB^\T$, while making it isotropic by adding energy can cost a factor $k$.

\section{Symmetry multiplicity closes the complete interaction}\label{base:con:symmetry-multiplicity}
Suppose, in compatible tree coordinates, $k=ab$ and
\[
 Q_{U,I}=\widetilde Q_{U,I}\otimes I_b,\qquad
 Q_{V,I}=\widetilde Q_{V,I}\otimes I_b,
 \qquad\widetilde Q_{U,I},\widetilde Q_{V,I}\in O(2a).
\]
The reduced completions may be unrelated between streams and noncommuting between levels. The probes remain full-width original Haar probes.

\begin{theorem}[Repeated-coordinate fixed-space bound]\label{base:con:multiplicity-thm}
At width $4k$, every deterministic input obeys
\begin{equation}\label{base:con:multiplicity-error}
 \E\F{A-\widehat B}^2
 \le[5+4\sqrt{\kappa(a+1+1/b)}]d_{\mathcal S}(A)^2.
\end{equation}
In particular, arbitrary scalar child weights in the two orientations, without matching them, give factor $5+4\sqrt{\kappa(2+1/k)}<44$ at $8k$ refit queries.
\end{theorem}
\begin{proof}
For one forward-origin source, suppose column completion blocks commute with an orthogonal $S$. Transform every forward frame by left multiplication by $S$, the source by $a\mapsto aS^\T$, and keep the same transpose realization. The local equations give
$Y\mapsto Y$, $Z\mapsto SZ$, $r\mapsto Sr$, $w'\mapsto Sw'$, and
$a_+,b_+\mapsto a_+S^\T,b_+S^\T$. Both doubly discarded readouts and the root preserve their norms. Thus
\[
 W(SG)=SW(G)S^\T,\qquad T_s(SG)=ST_s(G)S^\T.
\]
In the invariant initial gauge, simultaneous left rotation preserves the original joint frame law. A common right rotation connects this gauge to any fixed root gauge without changing overlaps. Therefore
\[
 \mathcal K=(I_m\otimes S)\mathcal K(I_m\otimes S^\T).
\]
For $S=I_a\otimes R$ and every signed permutation $R$ of $b$ coordinates, sign changes annihilate off-replication-coordinate blocks and permutations equalize diagonal ones. Hence
$\mathcal K=\widetilde{\mathcal K}\otimes I_b$.
Its leading eigenvalue has multiplicity at least $b$. Equation~\eqref{base:con:KyFan} with $r=b$ proves
$\op{\mathcal K}\le C_G\sqrt{1+(k+1)/b}$.
Use the corrected off-chain allowance \eqref{base:new:offchain}, both orientations, and \eqref{base:new:assembly}. At $q=4k$, $C_G=\sqrt\kappa$ and the coarse off-chain constant is $2$, smaller than the displayed chain allowance.
\end{proof}

The scalar-child case allows
\[
 U_I=\begin{bmatrix}\sqrt{p_I}O_{I,1}\\\sqrt{1-p_I}O_{I,2}\end{bmatrix},\qquad
 V_I=\begin{bmatrix}\sqrt{r_I}P_{I,1}\\\sqrt{1-r_I}P_{I,2}\end{bmatrix},
 \quad 0\le p_I,r_I\le1,
\]
with independent weights $p_I,r_I$. Its orthogonal gauges can be removed recursively: once a parent gauge is chosen, choose each child gauge by multiplying the parent's gauge by its own child factor. A tree has no compatibility cycle. Complement rotations are absorbed by Haar invariance. This does not change child singular values and cannot convert arbitrary anisotropic bases into this family.

For arbitrary bases define
\[
 C_{rs}=k^{-1}\tr\mathcal K_{rs},\qquad
 \overline{\mathcal K}=C\otimes I_k,
 \qquad\mathcal A=\mathcal K-\overline{\mathcal K}.
\]
Testing $B_s=z_sI_k$ in \eqref{base:con:gram-risk} gives
\begin{equation}\label{base:con:anisotropic}
 0\preceq\overline{\mathcal K}\preceq C_G\sqrt{2+1/k}I,
 \qquad \tr\mathcal A_{rs}=0.
\end{equation}
Thus the coordinate-averaged part, including all its temporal cross terms, is uniformly bounded. A uniform upper bound on $\lambda_{\max}(\mathcal A)$ would settle the general chain problem. Analytically averaging coordinates is not permission to replace the actual kernel by its average.

A scalar-child learner cannot cover every exact HSS input. For $k=2$, a cut whose exterior column space is spanned by one coordinate in each of its two children has child Gram spectra $\{0,1\}$ on both sides, not scalar spectra. Embed it as the only nonzero rank-two exterior block. Its best HSS error is zero, but no same-rank scalar-child space represents it. This is an obstruction to restricting the learner, not to the theorem for that family.


\endgroup

% ===== PREDICTABLE =====
\begingroup
\let\R\relax
\newcommand{\R}{\mathbb R}
\let\E\relax
\newcommand{\E}{\mathbb E}
\let\Prb\relax
\newcommand{\Prb}{\mathbb P}
\let\T\relax
\newcommand{\T}{\mathsf T}
\let\F\relax
\newcommand{\F}[1]{\left\lVert#1\right\rVert_{F}}
\let\N\relax
\newcommand{\N}[1]{\left\lVert#1\right\rVert}
\let\ip\relax
\newcommand{\ip}[2]{\left\langle#1,#2\right\rangle_F}
\let\tr\relax
\newcommand{\tr}{\operatorname{tr}}
\let\rank\relax
\newcommand{\rank}{\operatorname{rank}}
\let\diag\relax
\newcommand{\diag}{\operatorname{blockdiag}}
\let\range\relax
\newcommand{\range}{\operatorname{range}}
\let\OPT\relax
\newcommand{\OPT}{\operatorname{OPT}}
\let\HSS\relax
\newcommand{\HSS}{\operatorname{HSS}}
\let\diver\relax
\newcommand{\diver}{\operatorname{div}}
\let\Var\relax
\newcommand{\Var}{\operatorname{Var}}
\let\UC\relax
\newcommand{\UC}{\mathcal U}
\let\VC\relax
\newcommand{\VC}{\mathcal V}
\let\err\relax
\newcommand{\err}{\mathcal E}
\let\proofstep\relax
\newcommand{\proofstep}[1]{\par\addvspace{.45\baselineskip}\Needspace{3\baselineskip}\textbf{#1}\par\nobreak}
\let\op\relax
\newcommand{\op}[1]{\left\lVert#1\right\rVert_{\mathrm{op}}}
\let\Sn\relax
\newcommand{\Sn}[2]{\left\lVert#1\right\rVert_{S_{#2}}}
\let\Cov\relax
\newcommand{\Cov}{\mathsf C}
\let\Loss\relax
\newcommand{\Loss}{\mathsf L}
\let\Fil\relax
\newcommand{\Fil}{\mathcal F}
\let\Id\relax
\newcommand{\Id}{I}
\let\supp\relax
\newcommand{\supp}{\operatorname{supp}}
\let\vecop\relax
\newcommand{\vecop}{\operatorname{vec}}
\let\Gammafun\relax
\newcommand{\Gammafun}{\operatorname{Gamma}}
\chapter{Bellman registers, product decompositions, and quotient variance}\label{ch:predictable}
\begin{scopebox}
The predictable interaction isolated here is paid in Chapters~\ref{ch:feedback}--\ref{ch:drift}. Any open checkpoint concerning only that interaction describes this preliminary reduction, not the present status of fixed-space refitting.
\end{scopebox}
\section{Finite Bellman registers and the marked telescope}\label{base:con:bellman}
Hereafter $P=Yg^\T$, $C=Zh^\T$, $D_Y=Yw^\T$, $D_Z=Z\chi^\T$,
$J_Y=Y(I-g^\T g-w^\T w)Y^\T$, and
$J_Z=Z(I-h^\T h-\chi^\T\chi)Z^\T$.
The symbol $Q_s$ denotes a future-output quadratic, not a completion.

\begin{proposition}[Complete deterministic response registers]\label{base:con:registers}
The conditional expected future output has the form
\begin{equation}\label{base:con:five-total}
 Q_s=\tr(M_YJ_Y)+\tr(M_ZJ_Z)
 +\tr(M_{D_Y}D_YD_Y^\T)+\tr(M_{D_Z}D_ZD_Z^\T)
 +\zeta^\T K_{PC}\zeta,
\end{equation}
where $\zeta=(\vecop P,\vecop C)$ and all coefficient matrices are deterministic, positive semidefinite, and at most $4c_0I$.
\end{proposition}
\begin{proof}
Independent right rotations fixing each current frame pair act as all of $O(d)$ on its innovation complement. Invariance kills cross terms between innovation sectors and all fixed registers, and between the two independent innovation sectors. Each remaining innovation quadratic is a row-matrix trace. Independently rotating the coordinates of $w$ and $\chi$ leaves their spans and all future rules unchanged; it forces the two discarded-regression quadratics to the displayed trace form and kills their mixed terms. Only $P,C$ can couple. Each sector can be prescribed independently as observation data. Positivity and the bound follow by testing the arbitrary-state restart estimate \eqref{base:new:restart-two} on those sectors.
\end{proof}

Exact register updates are
\begin{align}
 P'&=A_U^\T(PA_V+K_GC_V),&
 C'&=A_V^\T(CA_U+K_HC_U),\label{base:con:retained-update}\\
 D_Y'&=A_U^\T(P-C^\T)B_V+A_U^\T K_GD_V-C_U^\T K_H^\T B_V,\label{base:con:DYupdate}\\
 D_Z'&=A_V^\T(C-P^\T)B_U+A_V^\T K_HD_U-C_V^\T K_G^\T B_U.\label{base:con:DZupdate}
\end{align}
In particular, for fixed fresh transpose frame, the one-step output-plus-future quadratic has degree at most two in the fresh forward frame \emph{after} simplification on the Stiefel support. Directly expanding moving innovation projectors creates apparent higher powers that cancel.

More explicitly, let $R_Y=A_UM_{Y,s+1}A_U^\T$, $R_Z=A_VM_{Z,s+1}A_V^\T$. There is a positive port matrix $\mathsf H_s$ on $v=(\vecop P,\vecop C,\vecop K_G,\vecop K_H)$ such that
\begin{align}
 O_s^*O_s+L_s^*Q_{s+1}L_s
 ={}&\tr[R_Y(YY^\T-PP^\T-K_GK_G^\T)]\nonumber\\
 &+\tr[R_Z(ZZ^\T-CC^\T-K_HK_H^\T)]+v^\T\mathsf H_sv.\label{base:con:ports}
\end{align}
Indeed $J_Y'=A_U^\T(YY^\T-PP^\T-K_GK_G^\T)A_U$ and similarly for $J_Z'$. Substitute \eqref{base:con:retained-update}--\eqref{base:con:DZupdate} into \eqref{base:con:five-total}; all remaining terms are positive quadratics in $v$.
Defining the column partial trace by
$\E[(\vecop K)^\T T\vecop K]=d^{-1}\tr[(\operatorname{ptr}T)J]$ when its covariance is $I\otimes J/d$, the exact backward equations are
\[
 K_{PC,s}=\mathsf H_{00,s},\quad M_{D_Y,s}=R_Y,\quad M_{D_Z,s}=R_Z,
\]
\[
 M_{Y,s}=(1-\rho)R_Y+d^{-1}\operatorname{ptr}\mathsf H_{GG,s},\qquad
 M_{Z,s}=(1-\rho)R_Z+d^{-1}\operatorname{ptr}\mathsf H_{HH,s}.
\]
Only $M_Y\succeq(1-\rho)M_{D_Y}$ follows, not $M_Y\succeq M_{D_Y}$. At $k=1$ with identity completions one step before the root,
$M_{D_Y}=1/(4d)$ but $M_Y=(1-1/d)/(4d)$. A proof based on monotone backward differences would therefore be invalid.

\subsection{Past output must remain in the state}
Fix a unit left source $u$. For the $k$ canonical starts $ue_\ell^\T$, let $T_s$ map source weights to current observations and $V_s$ map them to the outputs already emitted. Put
\[
 A_s=V_s^*V_s,\quad R_s=A_s+T_s^*Q_sT_s,
 \quad\mathcal K_s=A_s+T_s^*(O_s^*O_s+L_s^*Q_{s+1}L_s)T_s.
\]
Then $\E[\mathcal K_s\mid\Fil_s]=R_s$ and $R_{s+1}=\mathcal K_s$. For the actual row mark
$n_0=0$, $\delta n_s=b_st_sX^\T$, $n_{s+1}=n_s+\delta n_s$, $\beta=\sum_s\|b_s\|^2$, expansion and conditioning give
\begin{equation}\label{base:con:bellman-identity}
 \E[n_mR_mn_m^\T]
 =\sum_s\E[\delta n_s\mathcal K_s\delta n_s^\T]
 +2\sum_s\E[n_s\mathcal K_s^{\rm odd}\delta n_s^\T].
\end{equation}
The odd part is linear in the fresh forward frame. In coordinates $t=\T heta E_G$, $F=XE_G^\T$ and
$\mathcal K_s^{\rm odd}=\sum_{i,\alpha}\T heta_{i\alpha}\mathcal L_{i\alpha}$, its conditional contraction is exactly
\[
 \frac2d\sum_{i,\alpha,j}b_{s,i}F_{j\alpha}\,
 n_s\mathcal L_{i\alpha}e_j.
\]
Only second moments are needed here. A bound on each coefficient alone does not sum this expression over depth. Omitting $A_s$ would miss the new terminal mark's effect on already emitted coefficients.

\subsection{The sharp output--state tradeoff and self-mark bound}
For every local realization and either readout, write $e_s=\|O_s(Y,Z)\|^2$, $S_s=\F Y^2+\F Z^2$. Then
\begin{equation}\label{base:con:sharp-local}
 e_s+\lambda S_{s+1}\le\mu(\lambda)S_s,
 \quad\mu(\lambda)=\lambda+\tfrac12+\sqrt{\lambda^2+\tfrac14},\quad\lambda\ge0.
\end{equation}
To see this, decompose the observation rows in the orthonormal new frame pairs. The only coupled terms are $\|b\|^2+\lambda\|x-b\|^2$, whose coefficient matrix is
$\bigl(\begin{smallmatrix}\lambda&-\lambda\\-\lambda&\lambda+1\end{smallmatrix}\bigr)$.
Its top eigenvalue is $\mu(\lambda)$; every uncoupled sector has coefficient at most $\max\{1,\lambda\}$. Scalar identity/swap completions and the top eigenvector attain equality. In particular $S_{s+1}\le2S_s$.

Summing \eqref{base:con:strong-ledger} over canonical initial source coordinates gives the joint budget
\[
 (1-\rho)\E\tr A_s+\E\tr(T_s^*T_s)\le k.
\]
Since $\mathcal K_s\preceq A_s+\mu(4c_0)T_s^*T_s$ and
$\E[\delta n_s^\T\delta n_s\mid\Fil_s]\preceq\|b_s\|^2I/d$, it follows that
\begin{equation}\label{base:con:selfmark}
 \sum_s\E[\delta n_s\mathcal K_s\delta n_s^\T]
 \le\rho\mu(4c_0)\beta.
\end{equation}
At width $4k$ this coefficient is $(17+\sqrt{257})/4=8.257804885\ldots$, improving the earlier $35/2$ estimate. The mixed sum in \eqref{base:con:bellman-identity} is not included. Under the multiplicity family, the total-risk bound implies that this signed sum lies between $-\rho\mu(4c_0)\beta$ and the chain allowance in \eqref{base:con:multiplicity-thm}; its separate port terms need not vanish.

\section{The local Haar defect and the mean-forcing route}\label{base:con:marked-local}
A useful local comparison was obtained without changing the entire source law. Let $O$ be Haar on $O(d)$ and $Z$ have independent $N(0,1/d)$ entries. Write a quadratic polynomial as
\[
 f(M)=c+\sum_{ij}A_{ij}M_{ij}+\sum_{ij,\ell r}T_{ij,\ell r}M_{ij}M_{\ell r},
 \qquad T_{ij,\ell r}=T_{\ell r,ij}.
\]
Let $\mathscr R$ interchange the two row indices of $T$, and set $T_-=(T-\mathscr RT)/2$, $T_+=(T+\mathscr RT)/2$. The first is antisymmetric in both row and column pairs. Orthogonally split $T_+=T_{00}+T_{\rm tr}$, where $T_{00}$ has both partial traces zero. If
$a_{jr}=\sum_i(T_+)_{ij,ir}$, $b_{i\ell}=\sum_j(T_+)_{ij,\ell j}$, and $\tau=\sum_{ij}T_{ij,ij}$, then
\[
 (T_{00})_{ij,\ell r}=(T_+)_{ij,\ell r}
 -\frac{\delta_{i\ell}}d a_{jr}-\frac{\delta_{jr}}d b_{i\ell}
 +\frac\tau{d^2}\delta_{i\ell}\delta_{jr}.
\]
With $\mu=\tau/d$, orthogonal fourth moments give
\begin{align}
 \E|f(O)|^2&=|c+\mu|^2+\F A^2/d
 +\frac{2\F{T_-}^2}{d(d-1)}+\frac{2\F{T_{00}}^2}{(d-1)(d+2)},\label{base:con:haar-sectors}\\
 \E|f(Z)|^2&=|c+\mu|^2+\F A^2/d
 +\frac2{d^2}(\F{T_-}^2+\F{T_{00}}^2+\F{T_{\rm tr}}^2).\nonumber
\end{align}
Only the exterior-square sector produces a positive Haar-minus-Gaussian excess. Consequently, for Hilbert-valued polynomials of degree at most two,
\[
 \E\|f(O)\|^2\le\E\|f(Z)\|^2+\frac1{d-1}\E\|f_-(Z)\|^2
 \le\frac d{d-1}\E\|f(Z)\|^2.
\]
The determinant minor attains the last factor. The proof is a contraction of the three-pairing orthogonal tensor in the scalar, symmetric-traceless, and antisymmetric sectors; summing output coordinates gives the Hilbert-valued statement. This is an application of the classical Haar integration calculus, not a new integration formula \cite{CS}.

For a running marked response, the local transition is affine in $t$ after fixing the fresh transpose frame. Its quadratic-in-$t$ term is solely
\[
 \mathcal L_1(t)\mathcal U_s(B_{s+1}tX^\T).
\]
Thus the positive comparison defect is attached to the current source mark, not to every empty propagation step. This avoids automatically raising $d/(d-1)$ to the depth. It does not bound the whole process: old-state--new-source cross terms are second-moment contractions and remain even under Gaussian comparison.

\subsection{Summable mean forcing, and why it was not a disjoint risk bound}
Let $Y_{s,\ell},Z_{s,\ell}$ be the canonical unmarked observations and $x_\ell$ row $\ell$ of $X$. Define
\[
 H_s^\sharp=\sum_\ell Y_{s,\ell}\Pi_{G,s}x_\ell^\T,
 \qquad U_s=d^{-1}H_s^\sharp b_s.
\]
The mean mark contribution to the next forward retained regression is
$A_U^\T U_sC_V$. Its mean current coefficient outputs are
$B_U^\T U_sC_V$ and $B_U^\T U_sD_V$; the other off-diagonal mean is zero. Cauchy--Schwarz and the unmarked state ledger give
\[
 \sum_s\E\F{U_s}^2\le\rho^2\beta.
\]
The complementary deterministic recursion
$M'=A_U^\T(MA_V+U_sC_V)$,
$J=MB_V+U_sD_V$ satisfies
$\F{M'}^2+\F J^2\le\F M^2+\F{U_s}^2$.
This is simply multiplication of the block row $(M,U_s)$ by $Q_V$, followed by a contraction of the first output.

An earlier decomposition propagated the conditional observation mean forcing with separate adapted restarts. Its norm-sum bound was
$\sum_s\|m_s\|_{L_2}\le\rho\sqrt{c_0\beta}$ and its future-output bound was $4c_0^2\rho^2\beta$. Adding the mean current readouts by Minkowski gave
$\rho^2(1+2c_0)^2\beta$. These components overlap with the Bellman self-mark decomposition; their constants cannot be added as if they partitioned terminal risk. Section~\ref{base:con:predictable} gives a different exact product decomposition in which one passive first-moment filter pays the complete predictable mean at cost $\rho^2\beta$.

\section{Symmetric dilation and common-probe restarts}\label{base:con:symmetric}
A separate route changes the joint probe law. Interleave the two copies of every physical node and set
\[
 \mathcal A=\begin{pmatrix}0&A\\A^\T&0\end{pmatrix},\qquad
 W_I=\operatorname{diag}(U_I,V_I).
\]
This gives common nested row/column bases of rank $2k$ and leaf size $4k$. For the associated symmetric fixed space $\mathcal T$,
\[
 P_{\mathcal T}\mathcal A=
 \begin{pmatrix}0&P_{\mathcal S}A\\(P_{\mathcal S}A)^\T&0\end{pmatrix},
 \qquad d_{\mathcal T}(\mathcal A)^2=2d_{\mathcal S}(A)^2.
\]
Each dilation matvec uses one query to each of $A,A^\T$. Width $8k$ therefore costs $16k$ original queries. A factor-$C$ symmetric theorem transfers with the same factor to the off-diagonal block, whose HSS rank is at most $k$. Temporary diagonal-copy coefficients must be propagated before extraction; discarding them earlier changes the recursion.

For the common-probe chain, let
$Q=\bigl(\begin{smallmatrix}A&B\\C&D\end{smallmatrix}\bigr)$ and use one frame pair $(g,w)$ and one fresh sibling $t$. The transition is
\[
 r=A^\T g+C^\T t,\quad w'=B^\T g+D^\T t,\quad
 a=B^\T Yr^\T,\quad Y'=A^\T Y-a^\T w'.
\]
Its local symmetric readout energy is
\[
 e=2\F a^2+\F{\operatorname{sym}(B^\T Y(w')^\T)}^2.
\]
Put $P=Yg^\T$, $K=P-P^\T$, $R=YY^\T-PP^\T$,
$J=Y(I-g^\T g-w^\T w)Y^\T$, $L=BB^\T$, $M=I-L$, and $\tau=\tr L/d$.
The signed potential is
\[
 V=\F Y^2-\tr(P^2)=\tr R+\tfrac12\F K^2\ge0.
\]

\begin{proposition}[Symmetric state cone and its limitation]\label{base:con:sym-cone}
For a deterministic initial state at an adapted history, let $M_0=P_0$ and $M_{s+1}=A_s^\T M_sA_s$. Put $H_s=R_s+K_sK_s^\T$. Then
\[
 \E P_s=M_s,\qquad
 \operatorname{Cov}(\vecop P_s)\preceq d^{-1}I_k\otimes\E H_s,
\]
\[
 \E H_{s+1}=A_s^\T\E H_sA_s+\frac{\E\tr J_s}{d}C_s^\T C_s.
\]
For $q\ge4k$,
\[
 \E\F{Y_s}^2\le(1+2\rho)V_0+\F{P_0}^2.
\]
For an injected start $Y_0=-a^\T w_0$, this is at most $(1+2\rho)\F a^2$.
\end{proposition}
\begin{proof}
The fresh regression $N=Yt^\T$ is centered with covariance $I_k\otimes J/d$. Expand $P'=A^\T(PA+NC)$ and $H'$. The mixed terms vanish, giving the displayed $H$ identity. If the cone holds, the next covariance is bounded by
\[
 d^{-1}(A^\T A)\otimes(A^\T\E H A)
 +d^{-1}(C^\T C)\otimes(A^\T\E J A)
 \preceq d^{-1}I\otimes\E H'.
\]
Initially it is zero. Direct expansion of $V'$ gives a nonnegative conditional decrement for $d\ge k+1$, in particular for $d\ge2k$. Since $\tr H_s\le2V_s$, the centered regression energy is at most $2\rho V_0$. Expanding $\E\tr P_s^2$ about $M_s$ and using $\tr Z^2\le\F Z^2$ proves the state bound.
\end{proof}

The accumulated output is different. Exact local calculations give
\begin{align}
 \E[e+\F{P'}^2]-\F P^2
 ={}&\tr[L(PP^\T-P^\T P)]
 -\tfrac14\F{B^\T KB}^2+\mathcal I_L(J),\label{base:con:sym-output}\\
 V-\E V'={}&\tr(LR)+\tfrac12\F{B^\T KB}^2
 -\tau\tr(LJ)-d^{-1}\tr(LMJ),\label{base:con:sym-V}
\end{align}
where
\[
 \mathcal I_L(J)=\tau\tr J+\tfrac{\tau+\rho}{2}\tr(LJ)
 +\tfrac1{2d}\tr(LMJ).
\]
These formulas follow by substituting the transition, using $\E N=0$ and its covariance, and splitting $P$ into symmetric and skew parts. The remaining signed term in \eqref{base:con:sym-output} is a commutator trace. It can have either sign.

If $L_s=\ell_sI$ at every future step, it vanishes. Then
$\mathcal I_L(J)\le2\rho\ell\tr J$ and
$V-\E V'\ge(1-\rho)\ell\tr J$. Hence
\[
 S_s=\F{P_s}^2+\frac{2\rho}{1-\rho}V_s
\]
pays current output plus next storage. At the root,
$\E e_{\rm root}=\F{\operatorname{sym}P_h}^2+(\rho/2)\tr J_h\le S_h$.
Thus an adapted injected source satisfies
\begin{equation}\label{base:con:sym-adapted}
 \E\|\mathcal O_{\rm sym}(a)\|^2
 \le\frac{2\rho}{1-\rho}\F a^2.
\end{equation}
At width four times the retained rank, this is factor two. The scalar-child hypothesis is automatic only when that retained rank is one; the dilation of an original rank-one problem has retained rank two.

For nonscalar $L$, the state cone pays $\|B^\T KB\|^2$ but not every one-leg term $\|B^\T KA\|^2$. Charging the latter separately loses cancellation. A stronger cone with $\E YY^\T$ in place of $\E H$ was explored numerically but not proved; finite covariance-state screens are not a replacement for invariance. Finally an actual common-probe source can contain future frame information. The adapted theorem \eqref{base:con:sym-adapted} does not allow such anticipative conditioning.

\section{Generic identifiability and nonlinear global fitting}\label{base:con:identifiability}
This route uses the original Gaussian measurements
$\mathcal M(X)=(X\Omega,X^\T\Psi)$, with independent widths $s\ge3k$, rather than the fixed-space orthogonal-refit law.

\begin{theorem}[Saturated fixed-input identifiability]\label{base:con:identifiable}
Fix an HSS matrix $B$ whose exterior row and column cuts all have rank exactly $k$. With probability one, it is the unique HSS-rank-at-most-$k$ matrix $C$ satisfying $\mathcal M(C)=\mathcal M(B)$.
\end{theorem}
\begin{proof}
At a leaf, nullifying its local design leaves at least $k$ independent Gaussian columns acting on its rank-$k$ exterior block. This sample has rank $k$ almost surely and forces the exterior row space of every competing $C$ with the same observations. The transpose sample forces its column space. Equal forced spaces and observations give equal exact discrepancies. Compression preserves the remaining cut ranks and identical observations, so induction forces the parent spaces. At the root, the retained design has full row rank almost surely and fixes the core. There are finitely many rank events. For the fixed true input one may choose deterministic frames of the exact cut subspaces; changes of gauge do not alter their ranks. The competitor may depend on the sketches, since the same successful events force every feasible competitor's spaces.
\end{proof}

For almost every successful realization there exist $K<\infty$ and $\epsilon>0$, depending on $B$ and the sketches, such that every rank-at-most-$k$ HSS $C$ satisfies
\begin{equation}\label{base:con:local-inverse}
 \|\mathcal M(C-B)\|\le\epsilon
 \quad\Longrightarrow\quad
 \F{C-B}\le K\|\mathcal M(C-B)\|.
\end{equation}
Indeed the sample rank gaps and full-row-rank design conditions are open. In local coordinate charts, the exact raw decoder is smooth in a neighborhood of the successful observations. Every nearby model completion is recovered exactly by the same rank induction, and a derivative bound proves \eqref{base:con:local-inverse}. This prevents model completions escaping to infinity while their data converge to $\mathcal M(B)$.

For $y=\mathcal M(B+E)$ and $2\|\mathcal M(E)\|<\epsilon$, a global least-squares model fit exists and every minimizer obeys
$\F{C-B}\le2K\|\mathcal M(E)\|$. To prove existence, take an infimizing sequence. Eventually its observations lie in the controlled neighborhood; \eqref{base:con:local-inverse} bounds the matrices. Closedness of the cut-rank constraints gives a convergent subsequence. Neither $K$ nor $\epsilon$ is uniform, and their dependence on $\mathcal M(E)$ cannot be removed by factoring expectations.

Saturation is necessary for this uniqueness argument. If $u$ is a unit vector in the common null space of $\Omega^\T$ and $\Psi^\T$, then $I$ and $I-uu^\T$ have the same data and HSS exterior ranks at most one. Minimum-Frobenius fitting selects the latter rather than the identity. A mathematical exact selector is still available: lexicographically minimize all exterior cut ranks, ordered from leaves upward, among feasible completions. Observed nullified ranks force the fixed true input's ranks almost surely; use variable-size bases and repeat upward. A minimum exists because the possible rank vectors form a finite set. This gives noiseless selection, not a robust rank-selection rule for inconsistent data. Weak singular values, rank transitions, and random inverse constants remain the quantitative obstruction.


\section{A passive first-moment filter for complete output}\label{base:con:predictable}
This part treats arbitrary fixed bases and the original independent streams. Its output space $\mathscr E$ is the direct sum of every emitted coefficient sector and the root; the norm on each summand is Frobenius. Already emitted coefficients remain in that space throughout.

For retained regressions $P=Yg^\T$, $C=Zh^\T$ and arbitrary fresh-regression inputs $K_G,K_H$, the exact equations are \eqref{base:con:retained-update}. Their nonroot readouts are
\begin{align}
 a&=B_U^\T(PA_V+K_GC_V),&
 b&=\{B_V^\T(CA_U+K_HC_U)\}^\T,\label{base:con:passive-readout}\\
 d_Y&=B_U^\T(PB_V+K_GD_V),&
 d_Z&=\{B_V^\T(CB_U+K_HD_U)\}^\T.\nonumber
\end{align}
Use either $d_Y$ or $d_Z$, not both. At the root the output is the block matrix with $(P+C^\T)/2$ in its occupied diagonal block, $K_G/2$ and $K_H^\T/2$ in the two off-diagonal blocks, and zero in the other diagonal block.

Setting future fresh-regression inputs to zero computes the conditional mean of complete future output, because the retained recursion is linear and each fresh regression is centered. It does not compute the actual observation process.

Let $\mathcal M_0(P_0,C_0)$ be the zero-input complete output, and let $\mathcal I_s(K_G,K_H)$ be the deterministic complete output due to an impulse in the fresh regressions at stage $s$. The following inequality holds for arbitrary input arrays, random or deterministic:
\begin{equation}\label{base:con:passive-filter}
 \left\|\mathcal M_0(P_0,C_0)+\sum_s\mathcal I_s(K_{G,s},K_{H,s})\right\|^2
 \le\F{P_0}^2+\F{C_0}^2+\sum_s(\F{K_{G,s}}^2+\F{K_{H,s}}^2).
\end{equation}
To prove it, multiply $[P\ K_G]$ by $Q_V$ and $[C\ K_H]$ by $Q_U$. Orthogonal row splitting partitions each first output into its retained state and discarded readout; selecting the appropriate second-block readout costs no more than its energy. Summing over stages telescopes. At the root
\[
 \tfrac14\F{P+C^\T}^2+\tfrac14(\F{K_G}^2+\F{K_H}^2)
 \le\F P^2+\F C^2+\F{K_G}^2+\F{K_H}^2.
\]
This proves \eqref{base:con:passive-filter} with constant one.

\section{An exact four-term product decomposition}\label{base:con:product}
Condition on the initial history, with $X=g_0$ a row-isometry and the future joint history $\Fil_s$ ending just before the next sibling pair. Fix a unit left source $u$. Let $\mathcal V:\R^k\to\mathscr E$ be the complete random output map of the canonical starts $ue_\ell^\T$. Its current observations are $Y_{s,\ell},Z_{s,\ell}$. The actual source weights are columns
\[
 x_0=0,\qquad \delta x_s=Xt_s^\T b_s^\T,\qquad
 x_{s+1}=x_s+\delta x_s,\qquad\beta=\sum_s\|b_s\|^2,
\]
where the rows $b_s$ are deterministic. The terminal output is $\mathcal Vx_m$, not a sum of outputs with contemporaneous weights.

Let $\mathcal J_s=\E[\mathcal V\mid\Fil_s]$ and $D_s=\mathcal J_{s+1}-\mathcal J_s$. Initially $\mathcal J_0=0$. The passive first-moment equations prove the exact formula
\begin{equation}\label{base:con:Doob-increment}
 D_se_\ell=\mathcal I_s(Y_{s,\ell}t_s^\T,Z_{s,\ell}\eta_s^\T).
\end{equation}
In particular it is linear in the fresh pair. Set $r_s=\E[D_s\delta x_s\mid\Fil_s]$ and
\begin{align}
 \mathcal A&=\sum_s\mathcal J_s\delta x_s,&
 \mathcal B&=\sum_sD_sx_s,\nonumber\\
 \mathcal C&=\sum_s(D_s\delta x_s-r_s),&
 \mathcal D&=\sum_sr_s.\label{base:con:ABCD}
\end{align}
The discrete product rule gives
\begin{equation}\label{base:con:product-identity}
 \mathcal Vx_m=\mathcal A+\mathcal B+\mathcal C+\mathcal D.
\end{equation}
The first three sums are Hilbert-valued martingale differences. At equal times, $\mathcal A$ and $\mathcal B$ are odd under $(t_s,\eta_s)\mapsto(-t_s,-\eta_s)$ and the centered product in $\mathcal C$ is even. Across distinct times, martingale orthogonality applies. Hence
\begin{equation}\label{base:con:parity-product}
 \E\langle\mathcal A,\mathcal C\rangle=
 \E\langle\mathcal B,\mathcal C\rangle=0.
\end{equation}
No other pair is asserted to be orthogonal.

\subsection{Three terms have dimension-free bounds}
Initially the unmarked canonical sources have zero retained regressions. Equation~\eqref{base:con:strong-ledger}, summed over them, gives
\[
 \sum_\ell\E(\F{Y_{s,\ell}}^2+\F{Z_{s,\ell}}^2)\le k,
 \qquad\E\|\mathcal V\|_{\rm HS}^2\le c_0k.
\]
Since $\E[\delta x_s\delta x_s^\T\mid\Fil_s]\preceq\|b_s\|^2I/d$, conditional Jensen yields
\begin{equation}\label{base:con:A-old}
 \E\|\mathcal A\|^2\le\rho c_0\beta.
\end{equation}
For the predictable mean let
$H_s^\sharp=\sum_\ell Y_{s,\ell}\Pi_{G,s}x_\ell^\T$ and
$U_s=H_s^\sharp b_s/d$. The fresh second moments give
$r_s=\mathcal I_s(U_s,0)$. Thus \eqref{base:con:passive-filter} and the canonical state budget imply
\begin{equation}\label{base:con:D}
 \E\|\mathcal D\|^2\le\sum_s\E\F{U_s}^2\le\rho^2\beta.
\end{equation}
This pays the complete predictable mean, including its current readout, by one deterministic filter.

For the centered product use the following exact Haar identity. If $\T heta$ comprises the first $k$ rows of Haar $O(d)$, $b\in\R^{1\times k}$, and $A\in\R^{d\times d}$, write
\[
 S=\tfrac12(A+A^\T)-d^{-1}\tr(A)I,\qquad K=\tfrac12(A-A^\T).
\]
Then
\begin{equation}\label{base:con:conjugation}
 \E\left\|b\T heta A\T heta^\T-\frac{\tr A}{d}b\right\|^2
 =\|b\|^2\left[
 \frac{k+1-2/d}{(d-1)(d+2)}\F S^2+
 \frac{k-1}{d(d-1)}\F K^2\right].
\end{equation}
This follows by contracting the fourth Haar tensor separately on symmetric traceless and skew matrices. For $d\ge2k$, its first coefficient dominates the second and also $k/d^2$.

In forward innovation coordinates the $i$th row of the same-frame quadratic mark is $b\T heta A_i\T heta^\T$, with
$A_i=\sum_\ell (x_\ell E_G^\T)^\T(Y_{s,\ell}E_G^\T)_{i,:}$. Since $XE_G^\T$ is a contraction,
$\sum_i\F{A_i}^2\le\sum_\ell\F{Y_{s,\ell}\Pi_G}^2$.
The forward--transpose product is centered and, by independence of the fresh pair, costs at most
$(k/d^2)\|b_s\|^2\sum_\ell\F{Z_{s,\ell}\Pi_H}^2$.
Apply \eqref{base:con:conjugation}, then the impulse contraction and the canonical state budget. Martingale orthogonality yields
\begin{equation}\label{base:con:C}
 \E\|\mathcal C\|^2\le\gamma_{k,d}\beta,
 \qquad\gamma_{k,d}=\frac{k(k+1-2/d)}{(d-1)(d+2)}\le2\rho^2.
\end{equation}
At $d=2k$, $\gamma_{k,2k}\le5/18$, with exact difference
$(k-2)^2/[9(4k^2+2k-2)]$.

\subsection{The equivalent nonnegative predictable target}
Let $R(b)=\E\|\mathcal Vx_m\|^2$ and
$\mathcal P(b)=\E\|\mathcal B\|^2$. Equations~\eqref{base:con:parity-product}--\eqref{base:con:C} and the reverse triangle inequality give
\begin{equation}\label{base:con:equivalence}
 |\sqrt{R(b)}-\sqrt{\mathcal P(b)}|
 \le[\rho+\sqrt{\rho c_0+\gamma_{k,d}}]\sqrt\beta.
\end{equation}
At width $4k$ this coefficient is at most
$\nu_*=(1/2)+\sqrt{23/18}=1.630389\ldots$. Therefore a uniform bound on $\mathcal P(b)/\beta$ is equivalent, up to absolute constants, to the general chain-source bound.

The predictable energy is explicit. Let
\[
 (\Lambda_{G,s})_{ab}=\sum_{i=1}^k
 \langle\mathcal I_s(E_{ai},0),\mathcal I_s(E_{bi},0)\rangle,
\]
and define $\Lambda_{H,s}$ analogously. These are deterministic positive matrices, at most $kI$. Put $Y_s(x)=\sum_\ell x_\ell Y_{s,\ell}$ and similarly $Z_s(x)$. Exact conditional second moments give
\begin{equation}\label{base:con:P-explicit}
 \begin{split}
 \mathcal P(b)=d^{-1}\sum_s\E\bigl[&\tr(\Lambda_{G,s}Y_s(x_s)\Pi_{G,s}Y_s(x_s)^\T)\\
 &+\tr(\Lambda_{H,s}Z_s(x_s)\Pi_{H,s}Z_s(x_s)^\T)\bigr].
 \end{split}
\end{equation}
Every summand is nonnegative and predictable. Uniform bounds on the weights alone do not establish summability. The next section uses a sharper variance budget for this particular filter.

\section{The quotient contracts pathwise, but pays variance rather than common output}\label{base:con:quotient}
For an isolated source chain with zero sibling observations define
\begin{equation}\label{base:con:Phi}
 \Phi(Y,Z)=\F{Y-Pg}^2+\F{Z-Ch}^2+\F{P-C^\T}^2
 =\F Y^2+\F Z^2-2\tr(PC).
\end{equation}
For every local realization,
\begin{equation}\label{base:con:Phi-pathwise}
 \begin{split}
 \Phi-\Phi'={}&\F{B_U^\T Y(I-g^\T g-t^\T t)}^2
 +\F{B_V^\T Z(I-h^\T h-\eta^\T\eta)}^2\\
 &+\F{d_Y-d_Z}^2\ge0.
 \end{split}
\end{equation}
Indeed, before compression stack zero sibling observations. Their regression mismatch is
$\bigl(\begin{smallmatrix}P-C^\T&K_G\\-K_H^\T&0\end{smallmatrix}\bigr)$.
The innovations lose $\F{K_G}^2+\F{K_H}^2$, exactly compensated by the new mismatch blocks. Compression then discards the three orthogonal sectors displayed in \eqref{base:con:Phi-pathwise}; this is also a specialization of the compression identity in Section~\ref{base:new:energy}.

Averaging gives the all-state signed identity
\begin{equation}\label{base:con:Phi-mean}
 \begin{split}
 \Phi-\E[\Phi'\mid\Fil]={}&\F{B_U^\T D_Y}^2+\F{B_V^\T D_Z}^2
 +\F{B_U^\T(P-C^\T)B_V}^2\\
 &+(1-\alpha_V)\tr(L_UJ_Y)+(1-\alpha_U)\tr(L_VJ_Z),
 \end{split}
\end{equation}
where $L_U=B_UB_U^\T$, $L_V=B_VB_V^\T$ and
$\alpha_U=\tr L_U/d$, $\alpha_V=\tr L_V/d$. It requires no zero-cross-tensor assumption. Substituting conditional frame second moments in \eqref{base:con:Phi-pathwise}, or expanding the observation and regression cross terms, proves it.

At a canonical injected start, $\Phi_0=\|x\|^2$. Consequently its quadratic form on initial source coordinates satisfies
\begin{equation}\label{base:con:Qmonotone}
 \Phi(T_sx)=x^\T\mathsf Q_sx,
 \qquad0\preceq\mathsf Q_{s+1}\preceq\mathsf Q_s\preceq I.
\end{equation}
This is an isolated-source-chain statement, not a claim about merging arbitrary nonzero sibling observations.

If $Y=Mg$, $Z=M^\T h$, then $\Phi=0$ but current coefficient output need not vanish. Hence \eqref{base:con:Phi-pathwise} cannot be summed and called an output bound. The correct next object is centered future-output variance.

\section{A width-sensitive output bound for zero-regression restarts}\label{base:con:theta}
Set
\begin{equation}\label{base:con:theta-def}
 \vartheta_\rho=\frac{\rho(2+\rho)}{1-\rho^2}.
\end{equation}
For a one-sided adapted start with both retained regressions zero,
\begin{equation}\label{base:con:theta-output}
 \E\|\mathcal O_{\rm future}(Y_0,Z_0)\|^2
 \le\vartheta_\rho(\F{Y_0}^2+\F{Z_0}^2),
\end{equation}
for either readout, including the root. This sharpens $c_0$ to a quantity of order $\rho$.

Here is the full energy argument. The cross tensor in Lemma~\ref{base:con:restart-cone} is zero and both deterministic means vanish. Let $D_y,D_z$ be accumulated expected discarded-observation energies and $S_y,S_z$ terminal expected observation energies. For a $Y$-only start of energy $S_0$,
\[
 D_y+S_y\le S_0+\rho D_z,\qquad
 D_z+S_z\le\rho D_y.
\]
The output sums satisfy, respectively,
$\E\sum e_s\le\rho(2D_y+D_z)$ and
$\E\sum e_s\le\rho(D_y+2D_z)$.
For example the two forward readouts together equal
$\F{B_U^\T P}^2+\F{B_U^\T K_G}^2$; the covariance cone bounds each by $\rho\E\F{B_U^\T Y}^2$. The transpose readout is bounded by $\rho\E\F{B_V^\T Z}^2$. The root costs at most $(\rho/2)(S_y+S_z)$: its two retained-regression means and full cross tensor vanish, and each of its two fresh and retained squares is controlled by the cone.

Put $v=\rho(2+\rho)/(1-\rho^2)$ and $z=\rho(1+2\rho)/(1-\rho^2)$. Then
$v-\rho z=2\rho$, $z-\rho v=\rho$, and $v,z\ge\rho/2$.
Multiplying the first flow inequality by $v$ and the second by $z$ pays all coefficient and root costs by $vS_0$. Exchange the weights for the other readout and the streams for a $Z$-only start. This proves \eqref{base:con:theta-output}. At $q=4k$, $\vartheta_{1/2}=5/3$; as $\rho\downarrow0$, $\vartheta_\rho=2\rho+O(\rho^2)$.

\section{Five positive registers for the future output variance}\label{base:con:variance}
At an actual current joint history, let
\[
 \mathscr V_s(Y,Z)=\E[\|\mathcal O_s(Y,Z)-\mathcal M_s(P,C)\|^2\mid\Fil_s],
\]
where $\mathcal M_s$ is the deterministic conditional mean filter of Section~\ref{base:con:predictable}. The bases and horizon are fixed.

\begin{theorem}[Quotient variance bound]\label{base:con:variance-thm}
For arbitrary current observations and every admissible fixed continuation,
\begin{equation}\label{base:con:variance-main}
 0\le\mathscr V_s(Y,Z)\le\vartheta_\rho\Phi(Y,Z).
\end{equation}
It is valid after conditioning on the current joint history, but not after revealing future forward frames.
\end{theorem}
\begin{proof}
\proofstep{Common regressions add deterministic output.}
Adding $(Mg,M^\T h)$ to the observations produces no fresh regressions and updates to the common retained coefficient $A_U^\T MA_V$. By induction all its output is deterministic. Variance is unchanged by this shift. Its retained-regression dependence is therefore only through $\Delta=P-C^\T$.

\proofstep{The exact invariant quadratic has five positive blocks.}
Apply the independent innovation-complement rotations and discarded-frame coordinate rotations used in Proposition~\ref{base:con:registers}, now to the centered output. This gives deterministic positive matrices such that
\begin{equation}\label{base:con:variance-registers}
 \begin{split}
 \mathscr V_s={}&\tr(M_{Y,s}J_Y)+\tr(M_{Z,s}J_Z)
 +\tr(M_{D_Y,s}D_YD_Y^\T)\\
 &+\tr(M_{D_Z,s}D_ZD_Z^\T)
 +\vecop(\Delta)^\T H_s\vecop(\Delta).
 \end{split}
\end{equation}
The innovation groups kill cross terms with fixed regressions and with the other independent probe space. Rotating only the current discarded coordinates kills mixed terms involving $D_Y,D_Z$ and makes each of their quadratics a row trace. Positivity follows by prescribing one sector at a time as observation data. These arguments do not infer independence of later errors.

\proofstep{Bound four blocks by zero-regression starts.}
Pure $Y$ innovation and pure $Y=D_Yw$ starts are one-sided with zero retained regressions; \eqref{base:con:theta-output} bounds their total output, hence their variance. They show $M_Y,M_{D_Y}\preceq\vartheta_\rho I$. Pure transpose starts give the other two bounds. Their means in fact vanish, but this is not needed for the inequality.

\proofstep{Bound mismatch by backward induction.}
At the root, only the fresh regressions are random, so
$\mathscr V_h=(\rho/4)(\tr J_Y+\tr J_Z)$ and $H_h=0$.
Test $Y=\Delta g$, $Z=0$ at a preceding stage. Its fresh regressions and current variance are zero. The only next quotient registers are
\[
 \Delta'=A_U^\T\Delta A_V,\quad
 D_Y'=A_U^\T\Delta B_V,\quad
 D_Z'=-A_V^\T\Delta^\T B_U,
\]
with both next innovations zero. Their squared norms sum to
$\F\Delta^2-\F{B_U^\T\Delta B_V}^2\le\F\Delta^2$ by orthogonal two-sided block decomposition. The already established bounds on four blocks and the induction bound on $H_{s+1}$ therefore give $H_s\preceq\vartheta_\rho I$.

Finally, the five squared register energies in \eqref{base:con:variance-registers} sum exactly to $\Phi$. Summing their bounds proves \eqref{base:con:variance-main}, without multiplying by the number of registers.
\end{proof}

This theorem reconciles two earlier observations. A common coefficient can emit nonzero output while the quotient is zero; all such output is deterministic. An innovation can create a retained coefficient without an immediate coefficient loss; its \emph{future variance} is nevertheless paid by the quotient. Total output and centered output variance must not be interchanged.

\section{General constant-factor refitting with one logarithm of oversampling}\label{base:con:oversampling}
Let $D_s=\mathcal J_{s+1}-\mathcal J_s$ as in Section~\ref{base:con:product}, and set
$\Gamma_s=\E[D_s^*D_s\mid\Fil_s]$. Conditional martingale orthogonality, Theorem~\ref{base:con:variance-thm}, and \eqref{base:con:Qmonotone} yield the stronger predictable tail estimate
\begin{equation}\label{base:con:predictable-tail}
 \E[\sum_{r\ge s}\Gamma_r\mid\Fil_s]
 =\E[(\mathcal V-\mathcal J_s)^*(\mathcal V-\mathcal J_s)\mid\Fil_s]
 \preceq\vartheta_\rho\mathsf Q_s\preceq\vartheta_\rho I.
\end{equation}
Past output is already known and has no future variance. The equality is an equality of quadratic forms on canonical source weights.

Let $b\ge4$ be any valid squared-energy embedding allowance in Section~\ref{base:new:carleson}, with the harmless enlargement to four when necessary. Since $x_s$ is adapted and $\E\|x_m\|^2\le\rho\beta$,
\begin{equation}\label{base:con:pred-bound}
 \mathcal P(b_{\bullet})=\sum_s\E[x_s^\T\Gamma_sx_s]
 \le\rho\vartheta_\rho b\,\beta.
\end{equation}
Here $b_{\bullet}$ denotes the source rows, whereas the scalar $b$ is the embedding constant. The distinction avoids identifying a source block with a complexity parameter.

The sharper zero-regression output theorem also gives
$\E\|\mathcal V\|_{\rm HS}^2\le\vartheta_\rho k$ and hence
$\E\|\mathcal A\|^2\le\rho\vartheta_\rho\beta$.
The bounds \eqref{base:con:C} and \eqref{base:con:D} are unchanged. Grouping $\mathcal A+\mathcal C$ using their parity orthogonality in \eqref{base:con:product-identity} proves
\begin{equation}\label{base:con:chain-current}
 R(b_{\bullet})\le C_{\rm chain}(k,q,b)\,\beta,
 \quad
 C_{\rm chain}=
 \left[\sqrt{\rho\vartheta_\rho b}+\rho+
 \sqrt{\rho\vartheta_\rho+\gamma_{k,d}}\right]^2.
\end{equation}
This is a bound on the full terminal source--response product. It is not a sum of overlapping bounds from earlier decompositions.

\subsection{Reattach the corrected physical source}
The mean response now satisfies $\E W\preceq\vartheta_\rho I$. Conditioning on the ancestor chain in \eqref{base:new:offchain} therefore gives
\begin{equation}\label{base:con:corrected-source-final}
 \E\tr(Wa^\T a)\le C_{\rm chain}\beta+
 \frac{\vartheta_\rho k}{q-3k}\gamma_{\rm out}
 +\frac{\vartheta_\rho k}{d}\gamma_{\rm in}.
\end{equation}
The mixed terms vanish by conditional centering of the off-chain component. Sum over initial discarded sources, use their conditional output orthogonality and free-output contraction, then combine the two orientations. Fixed-space exactness and $E\perp\mathcal S$ give the general bound
\begin{equation}\label{base:con:general-refit}
 \boxed{\E\F{A-\widehat B}^2\le
 \left[5+4\max\left\{C_{\rm chain},
 \frac{\vartheta_\rho\rho}{1-\rho}\right\}\right]d_{\mathcal S}(A)^2.}
\end{equation}
Both interior and exterior discarded directions are included. No assumption on the residual's row-energy spread or on the geometry of the fixed bases has entered.

\begin{theorem}[General one-logarithm oversampling theorem]\label{base:con:main-eight}
For arbitrary fixed nested bases and every deterministic matrix $A$, choose a valid squared embedding bound $b\ge4$ and
\begin{equation}\label{base:con:width}
 q=\left\lceil(2+4\sqrt b)k\right\rceil.
\end{equation}
The unchanged sibling-orthogonal decoder, using $2q$ refit matrix--vector queries, satisfies
\begin{equation}\label{base:con:eight}
 \boxed{\E\F{A-\widehat B}^2\le8d_{\mathcal S}(A)^2.}
\end{equation}
The proof below gives a factor smaller than $7.234$.
\end{theorem}
\begin{proof}
The chosen width implies $\rho\le1/(4\sqrt b)\le1/8$. Hence
$\vartheta_\rho/\rho\le136/63$ and $\gamma_{k,d}\le2\rho^2$. Substitution into \eqref{base:con:chain-current} gives
\[
 C_{\rm chain}\le\frac1{16}
 \left[\sqrt{\frac{136}{63}}+
 \frac{1+\sqrt{262/63}}2\right]^2
 =0.5583489429\ldots.
\]
The off-chain constant is at most $17/441$. Equation~\eqref{base:con:general-refit} is therefore bounded by
$7.233395772\ldots<8$.
\end{proof}

\begin{corollary}[Completely explicit depth-dependent width]\label{base:con:explicit-width}
For $L\ge1$, put $a_L=1+\lceil\log_2(L+1)\rceil$. The choice
\[
 q=(2+4a_L)k,
 \qquad2q=(4+8a_L)k
\]
satisfies Theorem~\ref{base:con:main-eight}. No unspecified public-theorem constant is used in this choice.
\end{corollary}
\begin{proof}
The source contains at most $L+1$ reverse sibling stages. The binary-prefix proof of \eqref{base:new:time-embed} gives $b=a_L^2\ge4$. Apply \eqref{base:con:width}.
\end{proof}
Combining the explicit time bound with the public logarithmic-rank embedding bound gives
\[
 2q=O\bigl(k[1+\log\min\{k,L+1\}]\bigr)
\]
for constant-factor fixed-space refitting. The rank-based absolute constant is not numerically calibrated. The earlier sufficient constant-factor width was proportional to $kb$; the variance calculation makes it proportional to $k\sqrt b$. At fixed $q/k$, this argument still permits a squared-logarithmic approximation factor.

\subsection{A uniform running-state bound is not an output-sum bound}
There is also a useful fixed-time estimate at the original width. For the running source $x_n$ with budget $\beta_n$, the endpoint retained-register filter has impulse weights satisfying
\[
 \sum_{s<n}\|C_{V,s}A_{V,s+1}\cdots A_{V,n-1}\|_F^2
 =k-\|A_{V,0}\cdots A_{V,n-1}\|_F^2\le k,
\]
and the row-side identity is the same. The quotient inequality bounds each marked innovation by $\|x_s\|^2$, whose mean is at most $\rho\beta_n$. Thus the endpoint version of the predictable term costs at most $2\rho^2\beta_n$. The endpoint predictable-transform term costs at most $\rho^2\beta_n$ by the regression covariance cone; the centered-product and mean terms retain costs $\gamma_{k,d}\beta_n$ and $\rho^2\beta_n$. Apply the product parity and triangle inequalities, then use
$\F Y^2+\F Z^2\le\Phi+\F P^2+\F C^2$, to obtain
\begin{equation}\label{base:con:state-current}
 \E[\F{Y_n(x_n)}^2+\F{Z_n(x_n)}^2]
 \le\left\{\rho+
 \left[\rho+\sqrt{(1+\sqrt2)^2\rho^2+\gamma_{k,d}}\right]^2\right\}\beta_n.
\end{equation}
At $q=4k$ this is below $3.803\beta_n$. It is a bound at every deterministic stage, not on an expected maximum and not on the sum of outputs. Summing it over stages would reintroduce depth; Theorem~\ref{base:con:main-eight} instead uses future variance and embedding.

\subsection{The exact remaining analytic gap}
The current proof uses the refined inequality
$\E[\sum_{r\ge s}\Gamma_r\mid\Fil_s]\preceq\vartheta_\rho\mathsf Q_s$
with $\mathsf Q_{s+1}\preceq\mathsf Q_s$, but then replaces $\mathsf Q_s$ by $I$ before generic embedding. A uniform estimate for the specialized pair $(\Gamma_s,x_s)$ at fixed $\rho>0$ would close the general constant-factor refit. Monotonicity alone does not establish it: even the constant sequence $\mathsf Q_s=I$ allows the generic embedding obstruction of Section~\ref{base:new:abstract-obstruction}. The special coupling between the quotient forms, the passive filter, and the actual source must be used, or a different exact decoder must be analyzed.

This is the endpoint of the preliminary refit chain. Chapter~\ref{ch:uniform} subsequently removes its query logarithm. It is a refit theorem, not the end-to-end $O(k)$-query best-HSS theorem. Section~\ref{base:new:remaining} states the remaining learning requirement explicitly.



\endgroup

\part{Closing the uniform fixed-space refit}

% ===== CENTERED =====
\begingroup
\let\E\relax
\newcommand{\E}{\mathbb E}
\let\R\relax
\newcommand{\R}{\mathbb R}
\let\F\relax
\newcommand{\F}[1]{\lVert#1\rVert_F}
\let\op\relax
\newcommand{\op}[1]{\lVert#1\rVert_{\mathrm{op}}}
\let\tr\relax
\newcommand{\tr}{\operatorname{tr}}
\chapter{Residual centering and cumulative quotient loss}\label{ch:centered}
\begin{scopebox}
The embedding-dependent bounds in this chapter are preliminary certificates. The centered free-output and source-assembly identities remain the inputs to Chapter~\ref{ch:uniform}, where the logarithmic allowance is removed.
\end{scopebox}
\begin{chaptersynopsis}
For arbitrary fixed nested bases, the unchanged sibling-orthogonal HSS refit achieves expected squared Frobenius error at most $(1+\varepsilon)d_{\mathcal S}(A)^2$ with
$O(k/\varepsilon+k\sqrt b/\sqrt\varepsilon)$ queries, where $b\ge4$ is a valid squared-energy matrix martingale embedding constant. An explicit sufficient width is
$\lceil k[2+6(1+\sqrt{1+b\varepsilon})/\varepsilon]\rceil$ for $0<\varepsilon\le1$.
The temporal choice of $b$ is completely explicit. The new step centers the uncorrected coefficient estimates at the fixed-space residual, reducing their expected energy to $O(k/(q-2k))$ instead of an absolute allowance. The correction estimates are proved in Section~\ref{base:con:variance}. At its previous factor-eight width, the new assembly gives a factor below five.

Separately, the pathwise decreasing quotient forms imply a cumulative energy bound for the actual running source: total marked quotient loss plus terminal quotient energy is at most $4\rho\beta$ in expectation. This is stronger than a pointwise-in-time state bound. A reachable rank-one example shows that the current output variance can nevertheless be positive when the current quotient loss is identically zero, so the new loss bound does not alone pay the remaining predictable interaction.

The remaining logarithm at fixed accuracy and the independent basis-learning obligation are not removed. The relative-error theorem uses stated, inherited correction and source-orthogonality interfaces; it is not an independent audit of the complete preceding proof chain. No priority or formal-verification claim is made.
\end{chaptersynopsis}

\section{Scope and endpoint}
Fix the nested bases before drawing either refit probe, and let $\mathcal S$ be their telescoping coefficient space. The local retained bases have size $2k\times k$; the root coefficient is an arbitrary $2k\times2k$ matrix. The lifted discrepancy sectors, with gauge $U_I^{\mathsf T}D_IV_I=0$, and the root sector are mutually Frobenius-orthogonal. Put
\[
 E=A-P_{\mathcal S}A,\qquad d_{\mathcal S}(A)=\F E.
\]
The original decoder is linear in the two observations, exact on $\mathcal S$, and returns an element of $\mathcal S$ (Chapter~\ref{ch:probes}). It draws the original independent forward and transpose sibling-orthogonal probes $G,H$ of width $q\ge4k$ and uses $2q$ matrix--vector queries. No fresh sketch is introduced at individual levels.

Set
\begin{equation}\label{centered:eq:parameters}
 d=q-2k,\qquad \rho=k/d\le\tfrac12,\qquad
 \vartheta=\frac{\rho(2+\rho)}{1-\rho^2},\qquad
 \gamma=\frac{k(k+1-2/d)}{(d-1)(d+2)}.
\end{equation}
Let $b\ge4$ be a valid squared-energy embedding constant for the source martingales in the preceding quotient-variance theorem (Section~\ref{base:con:variance}). Define its correction constants
\begin{align}
 C_{\rm ch}&=\left[\sqrt{\rho\vartheta b}+\rho+\sqrt{\rho\vartheta+\gamma}\right]^2,\label{centered:eq:chain}\\
 C_{\rm off}&=\frac{\vartheta\rho}{1-\rho},\qquad
 C_*=\max\{C_{\rm ch},C_{\rm off}\}.\label{centered:eq:cstar}
\end{align}
These are not new estimates in this chapter. Their exact interface is stated in Section~\ref{centered:sec:assemble}.

\begin{theorem}[Residual-centered general certificate]\label{centered:thm:main}
Let $f_Y=\min\{1,4\rho\}$ and $f_Z=\min\{1,2\rho\}$. For every deterministic matrix and every fixed nested basis family, either choice of the doubly discarded readout satisfies
\begin{equation}\label{centered:eq:main}
 \boxed{\E\F{A-\widehat B}^2
 \le\left[1+\left(\sqrt{f_Y+C_*}+\sqrt{f_Z+C_*}\right)^2\right]
 d_{\mathcal S}(A)^2.}
\end{equation}
In particular,
\begin{equation}\label{centered:eq:simple}
 \E\F{A-\widehat B}^2\le[1+12\rho+4C_*]d_{\mathcal S}(A)^2.
\end{equation}
The exact readout reversal interchanges $f_Y,f_Z$, leaving the displayed bound unchanged.
\end{theorem}

\begin{corollary}[Arbitrary-accuracy fixed-space refitting]\label{centered:cor:eps}
For $0<\varepsilon\le1$, choose
\begin{equation}\label{centered:eq:width}
 \boxed{q=\left\lceil k\left[2+\frac{6(1+\sqrt{1+b\varepsilon})}{\varepsilon}\right]\right\rceil.}
\end{equation}
Then
\begin{equation}\label{centered:eq:eps}
 \boxed{\E\F{A-\widehat B}^2\le(1+\varepsilon)d_{\mathcal S}(A)^2,\qquad
 \text{refit queries}=2q.}
\end{equation}
\end{corollary}

For depth $L\ge1$, take
\[
 a_L=1+\lceil\log_2(L+1)\rceil,\qquad b=a_L^2.
\]
This is the explicit temporal embedding proved in the preceding analysiss. It gives
\[
 O\!\left(\frac{k}{\varepsilon}+\frac{k[1+\log(L+1)]}{\sqrt\varepsilon}\right)
\]
refit queries without a numerically unspecified embedding constant. Using also the inherited classical rank estimate $b\le C_{\rm CE}(1+\log k)^2$ \cite[Theorem 1.1, squared-energy form]{NTV} gives
\begin{equation}\label{centered:eq:asymptotic}
 O\!\left(\frac{k}{\varepsilon}+
 \frac{k[1+\log\min\{k,L+1\}]}{\sqrt\varepsilon}\right).
\end{equation}
The rank-based constant is not assigned a numerical value here. When $b\varepsilon\le1$, \eqref{centered:eq:width} is $O(k/\varepsilon)$ with no multiplicative embedding factor. At fixed accuracy the single logarithm remains.

\section{A centered regression estimate from the actual probe law}\label{centered:sec:regression}
We only need the fourth-order case of the whole-probe chi-square comparison from (Section~\ref{base:new:chi}); here is a proof of that case from the stated construction.

\begin{lemma}[Whole-probe fourth moment]\label{centered:lem:probe}
For every deterministic row $v$,
\begin{equation}\label{centered:eq:probe}
 \E\|vG\|^2=\|v\|^2,\qquad
 \operatorname{Var}(\|vG\|^2)\le\frac2d\|v\|^4.
\end{equation}
The statement remains valid conditional on a fixed root frame.
\end{lemma}
\begin{proof}
Expand $v$ in the deterministic orthonormal multiresolution coordinates of the column bases. The root contribution has deterministic norm. Generate the remaining discarded sibling pairs top-down. A fixed linear combination of the rows of one pair is conditionally uniform on a sphere of deterministic radius $a$ in a $d$-dimensional complement. Its conditional mean is zero and its covariance is $(a^2/d)P$ for a projection $P$.

For the current sum $x$ and the next contribution $\xi$,
\[
 \E[\|x+\xi\|^4\mid\text{past}]
 =\|x\|^4+2a^2\|x\|^2+a^4+4\E[\langle x,\xi\rangle^2\mid\text{past}]
 \le\|x\|^4+(2+4/d)a^2\|x\|^2+a^4.
\]
If the accumulated deterministic squared coefficient budget is $B$, induction gives
$\E\|x\|^2=B$ and $\E\|x\|^4\le(1+2/d)B^2$.
The last display closes this induction at $B+a^2$. Parseval makes the final budget $\|v\|^2$. Subtract its squared mean to obtain \eqref{centered:eq:probe}. No independence of compressed probe rows is assumed.
\end{proof}

\begin{lemma}[Regression onto an exact probe frame]\label{centered:lem:regression}
Let $V\in\R^{N\times r}$ be a deterministic isometry such that
$V^{\mathsf T}GG^{\mathsf T}V=I_r$ pathwise. If a deterministic matrix $F$ satisfies $FV=0$, then
\begin{equation}\label{centered:eq:regression}
 \E[FGG^{\mathsf T}V]=0,\qquad
 \boxed{\E\F{FGG^{\mathsf T}V}^2
 \le\min\{1,2r/d\}\F F^2.}
\end{equation}
The random frame $V^{\mathsf T}G$ is not required to be independent of $FG$.
\end{lemma}
\begin{proof}
First take one row $f$ of $F$ and one unit column $v$ of $V$. Let $a=\|f\|$, $S=\|fG\|^2$, and $Z=\langle fG,v^{\mathsf T}G\rangle$. Isotropy gives $\E Z=fv=0$, while $\|v^{\mathsf T}G\|=1$ for every realization. For $a>0$,
\[
 \tfrac12\operatorname{Var}(\|(f+a v^{\mathsf T})G\|^2)
 +\tfrac12\operatorname{Var}(\|(f-a v^{\mathsf T})G\|^2)
 =\operatorname{Var}S+4a^2\E Z^2.
\]
Both input rows on the left have squared norm $2a^2$. Lemma~\ref{centered:lem:probe} bounds the left side by $8a^4/d$, so $\E Z^2\le2a^2/d$. The case $a=0$ is immediate. Summing over the columns of $V$ and rows of $F$ gives the factor $2r/d$. Also $\F{FGG^{\mathsf T}V}\le\F{FG}$ pathwise, since $V^{\mathsf T}G$ has orthonormal rows; isotropy gives the competing factor one.
\end{proof}

This is a centered residual estimate, not a concentration event or an inverse bound. Its usefulness depends on $FV=0$. Applying the old pathwise contraction alone would lose the factor $r/d$.

\section{The free coefficient outputs cost only \texorpdfstring{$O(\rho)$}{O(rho)}}\label{centered:sec:free}
Use the word \emph{free} for the following analytical decomposition: propagate the original data by the fixed retained isometries without subtracting any estimated discrepancies, and read the usual coefficient sectors. This does not change the actual decoder; its difference from this free output is the sum of the established source corrections.

At node $j$, let $\widetilde U_j$ and $\widetilde V_j$ be the lifted discarded bases, and $U_j,V_j$ the lifted retained bases. Let $V_j^{\rm loc}$ denote the full lifted $2k$-dimensional precompression column space. With forward-only data $EG$, the free nonroot coefficient pair has total squared norm
\[
 \F{\widetilde U_j^{\mathsf T}E\,GG^{\mathsf T}V_j^{\rm loc}}^2.
\]
With transpose-only data $E^{\mathsf T}H$, its free nonroot coefficient has squared norm
\[
 \F{\widetilde V_j^{\mathsf T}E^{\mathsf T}HH^{\mathsf T}U_j}^2.
\]
These are the forward doubly discarded readout. For the other readout the roles of the one- and two-sector estimates are interchanged.

Since $E\perp\mathcal S$,
\begin{equation}\label{centered:eq:zero}
 \widetilde U_j^{\mathsf T}EV_j^{\rm loc}=0,\qquad
 \widetilde V_j^{\mathsf T}E^{\mathsf T}U_j=0.
\end{equation}
The local physical probe frames have exactly orthonormal rows by construction. Lemma~\ref{centered:lem:regression} applies with $r=2k$ and $r=k$, respectively.

Let $U_*,V_*$ be the root bases, each of width $2k$. The forward root contribution is
$\frac12U_*^{\mathsf T}EGG^{\mathsf T}V_*$; its expected squared norm is at most
$\min\{1/4,\rho\}\F{U_*^{\mathsf T}E}^2$ because $U_*^{\mathsf T}EV_*=0$. The transpose root has the corresponding estimate.

The lifted discarded spaces and root form orthogonal decompositions on their respective physical sides:
\begin{align*}
 \sum_j\F{\widetilde U_j^{\mathsf T}E}^2+\F{U_*^{\mathsf T}E}^2&=\F E^2,\\
 \sum_j\F{\widetilde V_j^{\mathsf T}E^{\mathsf T}}^2+\F{V_*^{\mathsf T}E^{\mathsf T}}^2&=\F E^2.
\end{align*}

\begin{proposition}[Residual free-output bound]\label{centered:prop:free}
Writing $\mathcal F_Y(EG)$ and $\mathcal F_Z(E^{\mathsf T}H)$ for the free coefficient stacks,
\begin{equation}\label{centered:eq:free}
 \boxed{\E\|\mathcal F_Y\|^2\le f_Y\F E^2,\qquad
 \E\|\mathcal F_Z\|^2\le f_Z\F E^2,}
\end{equation}
where $f_Y=\min\{1,4\rho\}$, $f_Z=\min\{1,2\rho\}$. Both free stacks have zero mean. The opposite readout interchanges the constants.
\end{proposition}
\begin{proof}
All nonroot and root coefficient sectors are orthogonal. Sum the estimates preceding the statement and use Parseval. The root constant is no larger than either relevant nonroot constant. Every mean is zero by \eqref{centered:eq:zero}, its root counterpart, and $\E GG^{\mathsf T}=\E HH^{\mathsf T}=I$.
\end{proof}

The residual $E$ is a proof device. The algorithm does not query $P_{\mathcal S}A$ or require dense access to the input. Linearity and exactness permit analyzing its actual error through the residual inputs.

\section{Correction interface and final assembly}\label{centered:sec:assemble}
We state the inherited interfaces precisely. The canonical chain correction for an actual source, including its terminal source weights on earlier emitted coefficients, obeys the bound $C_{\rm ch}\beta$ in \eqref{centered:eq:chain}, from (Section~\ref{base:con:variance}). The corrected source decomposition from (Section~\ref{base:new:carleson}; \ref{base:con:interior-example}) is
\[
 a=N+\zeta,\qquad
 \beta+\gamma_{\rm out}+\gamma_{\rm in}=\F e^2.
\]
Here strict descendant discarded column directions inside the source node are included. Conditional on the entire forward ancestor chain,
\[
 \E[\zeta\mid\mathcal C]=0,\qquad
 \E[\zeta^{\mathsf T}\zeta\mid\mathcal C]
 \preceq\left(\frac{\gamma_{\rm out}}{q-3k}
 +\frac{\gamma_{\rm in}}d\right)I.
\]
The opposite-stream mean response obeys $\E W\preceq\vartheta I$. Therefore the expected correction energy for a source is at most
\[
 C_{\rm ch}\beta+\frac{\vartheta k}{q-3k}\gamma_{\rm out}
 +\frac{\vartheta k}{d}\gamma_{\rm in}\le C_*\F e^2.
\]
Distinct source corrections are orthogonal in conditional expectation over the opposite stream. They are also orthogonal in that expectation to the same one-sided free output. Their deterministic row budgets sum to at most $\F E^2$. These source-orthogonality facts are inherited, not consequences of the centered regression lemma.

\begin{proof}[Proof of Theorem~\ref{centered:thm:main}]
Let $\mathcal R_Y=\mathcal R(EG,0)$ and $\mathcal R_Z=\mathcal R(0,E^{\mathsf T}H)$, as coefficient stacks or equivalently as reconstructed fixed-space matrices. Proposition~\ref{centered:prop:free} and the just-stated orthogonality give
\[
 \E\|\mathcal R_Y\|^2\le(f_Y+C_*)\F E^2,\qquad
 \E\|\mathcal R_Z\|^2\le(f_Z+C_*)\F E^2.
\]
The two complete one-sided outputs are not asserted to be orthogonal or independent. Apply the $L_2$ triangle inequality between them. Exactness on $\mathcal S$ and $E\perp\mathcal S$ give
\[
 \F{A-\widehat B}^2=\F E^2+\|\mathcal R_Y+\mathcal R_Z\|^2.
\]
This proves \eqref{centered:eq:main}. The inequality $(\sqrt a+\sqrt b)^2\le2(a+b)$ and $f_Y+f_Z\le6\rho$ prove \eqref{centered:eq:simple}.
\end{proof}

\begin{proof}[Proof of Corollary~\ref{centered:cor:eps}]
For $\rho\le1/8$,
\[
 \vartheta/\rho\le136/63,\qquad \gamma\le2\rho^2.
\]
Since $b\ge4$,
\begin{align*}
 \sqrt{C_{\rm ch}}
 &\le\rho\left[\sqrt{(136/63)b}+1+\sqrt{262/63}\right]\\
 &\le\rho\sqrt b\left[\sqrt{136/63}+\frac{1+\sqrt{262/63}}2\right]
 <3\rho\sqrt b.
\end{align*}
For an explicit rational check of the strict inequality, use
$\sqrt{136/63}<147/100$ and $\sqrt{262/63}<51/25$; the bracket is less than $299/100<3$.
Also $C_{\rm off}\le(1088/441)\rho^2<3\rho^2$. Thus $C_*\le9b\rho^2$ and
\begin{equation}\label{centered:eq:smallrho}
 \E\F{A-\widehat B}^2\le[1+12\rho+36b\rho^2]d_{\mathcal S}(A)^2.
\end{equation}
The width \eqref{centered:eq:width} guarantees
\[
 \rho\le\rho_*:=\frac{\varepsilon}{6(1+\sqrt{1+b\varepsilon})}\le\frac1{12}<\frac18.
\]
Direct substitution gives $12\rho_*+36b\rho_*^2=\varepsilon$. This proves the result, with upward integer rounding only increasing the width.
\end{proof}

At the earlier factor-eight width $q=\lceil(2+4\sqrt b)k\rceil$, $\rho\le1/(4\sqrt b)\le1/8$. Using the preceding analysis's sharper common bound $C_*\le0.5583489429\ldots$ in \eqref{centered:eq:simple} gives
\[
 \E\F{A-\widehat B}^2\le4.733395772\ldots\,d_{\mathcal S}(A)^2<5d_{\mathcal S}(A)^2.
\]
Alternatively, \eqref{centered:eq:smallrho} gives the simpler factor $19/4=4.75<5$. Thus the new assembly also improves the previous constant at the same width, not only in the high-accuracy regime.

\section{A depth-free cumulative bound for the running quotient}\label{centered:sec:dissipation}
The next result is separate from the proof of the relative-error theorem. It advances the paid-energy analysis of the remaining $O(k)$ question.

\begin{lemma}[Decreasing positive forms and a martingale]\label{centered:lem:decreasing}
Let $(x_s)_{s=0}^m$ be a square-integrable real vector martingale with $x_0=0$. Let $Q_s$ be adapted self-adjoint matrices satisfying
$0\preceq Q_{s+1}\preceq Q_s\preceq I$ pathwise. Put $D_s=Q_s-Q_{s+1}$. Then
\begin{equation}\label{centered:eq:decreasing}
 \boxed{\sum_{s<m}\E[x_s^{\mathsf T}D_sx_s]
 +\E[x_m^{\mathsf T}Q_mx_m]\le4\E\|x_m\|^2.}
\end{equation}
The matrices need not commute with one another or be independent of the martingale.
\end{lemma}
\begin{proof}
Use the direct-sum $L_2$ norm of the vectors
$(D_s^{1/2}x_s)_{s<m}$ together with $Q_m^{1/2}x_m$.
Write this vector as the difference of
\[
 U=((D_s^{1/2}x_{s+1})_{s<m},Q_m^{1/2}x_m),\qquad
 V=((D_s^{1/2}(x_{s+1}-x_s))_{s<m},0).
\]
Since $D_s$ is $\mathcal F_{s+1}$-measurable, conditional Jensen gives
\[
 \|U\|_{L_2}^2\le\E\left[x_m^{\mathsf T}\left(\sum_{s<m}D_s+Q_m\right)x_m\right]
 =\E[x_m^{\mathsf T}Q_0x_m]\le\E\|x_m\|^2.
\]
Also $D_s\preceq I$, so martingale orthogonality gives
$\|V\|_{L_2}^2\le\sum_s\E\|x_{s+1}-x_s\|^2=\E\|x_m\|^2$.
The triangle inequality proves \eqref{centered:eq:decreasing}. No terminal-source conditional law is substituted into a fresh transition.
\end{proof}

In the HSS source chain, write $T_s$ for the canonical map from source coordinates to the current observation pair and define
\[
 \Phi(Y,Z)=\F{Y-Pg}^2+\F{Z-Ch}^2+\F{P-C^{\mathsf T}}^2,
 \quad P=Yg^{\mathsf T},\ C=Zh^{\mathsf T}.
\]
The preceding quotient identity gives adapted matrices $\mathsf Q_s$ with
\[
 \Phi(T_sx)=x^{\mathsf T}\mathsf Q_sx,\qquad
 I=\mathsf Q_0\succeq\mathsf Q_1\succeq\cdots\succeq0.
\]
The actual source increments are $\delta x_s=Xt_s^{\mathsf T}b_s^{\mathsf T}$, where $X$ is the initial retained frame and the rows $b_s$ are deterministic. Hence
\[
 \E\|x_m\|^2\le\rho\beta,\qquad \beta=\sum_s\|b_s\|^2.
\]
If the terminal root mark is included, extend the quotient form constantly over that last draw; it has no additional nonroot compression loss.

\begin{corollary}[Actual running-source dissipation]\label{centered:cor:loss}
In the original joint Haar law,
\begin{equation}\label{centered:eq:hssloss}
 \boxed{\sum_{s<m}\E\left[x_s^{\mathsf T}(\mathsf Q_s-\mathsf Q_{s+1})x_s\right]
 +\E[x_m^{\mathsf T}\mathsf Q_mx_m]\le4\rho\beta.}
\end{equation}
Moreover $\E\max_{s\le m}\Phi(T_sx_s)\le4\rho\beta$ by $\mathsf Q_s\preceq I$ and the $L_2$ maximal inequality for a martingale. The same finite-horizon statements apply after stopping both processes at a bounded stopping time.
\end{corollary}

Here $x_s^{\mathsf T}(\mathsf Q_s-\mathsf Q_{s+1})x_s=\Phi(T_sx_s)-\Phi(T_{s+1}x_s)$: the current source is held fixed through compression, before its next increment is added. This is not the difference using $T_{s+1}x_{s+1}$. Each local loss in \eqref{centered:eq:hssloss} is the actual sum of the two discarded innovation energies and the squared difference of the two discarded-coefficient readouts from the quotient identity. This is a cumulative bound, not a sum of separate stagewise upper bounds. It does not, however, bound the complete coefficient-output energy.

\section{Current variance versus current quotient loss}\label{centered:sec:obstruction}
A local witness distinguishes the two quantities. Take $k=1$, $q\ge4$, with current frames $g=e_1^{\mathsf T}$ and $w=e_2^{\mathsf T}$, observations $Y=e_3^{\mathsf T}$ and $Z=0$, and completions
\[
 Q_U=I_2,\qquad
 Q_V=\begin{pmatrix}0&-1\\1&0\end{pmatrix}.
\]
All current coefficients vanish and $Y'=Y,Z'=0$, while $g'=t,w'=-g$. Thus
\[
 \Phi(Y,Z)=\Phi(Y',Z')=1
\]
for every new frame. The current quotient loss is identically zero.

Suppose the root follows this step. The conditional mean of the root's retained block after revealing $t$ is $Yt^{\mathsf T}/2$; before revealing it, that mean is zero. The corresponding output-map Doob increment therefore has variance
\begin{equation}\label{centered:eq:witness}
 \boxed{\E\|D_s\|^2=\frac1{4d}>0,\qquad \Phi_s-\Phi_{s+1}=0.}
\end{equation}
The state is reachable: an identity-completion warmup from a one-sided discarded-frame source $Y=w_{\rm old},Z=0$ retains $g$ and makes the new discarded frame orthogonal to $Y$. A global coordinate rotation gives the displayed state. Interchanging streams gives the canonical opposite orientation.

Consequently there is no all-state local inequality making output quadratic variation at time $s$ at most a constant times the expected quotient loss at the same time. The example is not a counterexample to a global bound: its final output is bounded, and the terminal quotient reserve can pay it. It explains why \eqref{centered:eq:hssloss} cannot simply replace the remaining predictable-energy estimate.

The current predictable interaction is
\[
 \mathcal P(b)=\sum_s\E[x_s^{\mathsf T}\Gamma_sx_s],\qquad
 \Gamma_s=\E[D_s^*D_s\mid\mathcal F_s],
\]
with the stronger inherited tail
\[
 \E[\sum_{r\ge s}\Gamma_r\mid\mathcal F_s]\preceq\vartheta\mathsf Q_s.
\]
Replacing $\mathsf Q_s$ by $I$ and using a generic matrix embedding produces the remaining logarithm. The cumulative bound \eqref{centered:eq:hssloss} pays genuine quotient dissipation, but \eqref{centered:eq:witness} shows that an argument connecting variance to later dissipation or retained-coefficient transport is still needed. No uniform $O(k)$ fixed-space theorem is claimed.

\section{Verification and remaining scope}
The new checker records 120 accuracy/width certificates, with the simplified inequality \eqref{centered:eq:smallrho} rechecked using exact rational arithmetic at the chosen integer widths. Twenty additional certificates verify the factor below five at the preceding temporal width. Twenty-one finite binary-tree martingale models (up to 128 paths each) use rational noncommuting contraction matrices and verify the decreasing-form telescope exactly. Three local conditional second-moment designs verify \eqref{centered:eq:witness}; these designs are not substituted for fourth-order Haar laws.

Eight full-tree original-Haar configurations use arbitrary fixed Haar bases, including a directional-selector perturbation, ranks one through three, depths two and three, and both readouts. They verify the residual-cut orthogonality, the complete decoder's linearity, the free/correction decomposition, the same-side cross expectation as a sampling diagnostic, and the new bounds. Sample sizes range from 48 to 384 probe pairs. Two cases use the new $\varepsilon=1/2$ width. The largest deterministic residual-cut discrepancy is below $1.9\cdot10^{-16}$, and the decoder linearity discrepancy below $1.7\cdot10^{-16}$.

Six running-source original-Haar configurations use ranks one through four, horizons through forty, and 384--1536 paths each. They keep the terminal and running sources correlated with the forward frames and check the pathwise positive quotient-loss matrices and the cumulative ledger. These sample means are diagnostics, not interval-certified expectations or estimates of sharp constants. The source/correction theorem from (Section~\ref{base:con:variance}) is an inherited premise; this chapter does not re-audit every earlier lemma.

The full public objective asks for $O(k)$ queries and an $O(L)$ expected-error factor against the best HSS-rank-$k$ matrix. Here the rank-$k$ space is fixed before the refit. The relative-error result can be conditioned on an independently learned space, but it does not supply
\[
 \E_{\rm learn}d_{\mathcal S}(A)^2\le C_bL\,\mathrm{OPT}(A)
\]
from $O(k)$ queries. At fixed accuracy the new query bound still contains one logarithm. The new gain is a residual-centered $(1+\varepsilon)$ guarantee and a dimension-free cumulative marked quotient bound, not a solution of either remaining all-input problem. The external HSS approximation theorem \cite{AmselHSS} remains a different, best-model benchmark.

\begingroup\footnotesize
\endgroup

\endgroup

% ===== DELAYED =====
\begingroup
\let\E\relax
\newcommand{\E}{\mathbb E}
\let\R\relax
\newcommand{\R}{\mathbb R}
\let\F\relax
\newcommand{\F}[1]{\lVert#1\rVert_F}
\let\T\relax
\newcommand{\T}{\mathsf T}
\let\tr\relax
\newcommand{\tr}{\operatorname{tr}}
\chapter{Direct, delayed, and root predictable weights}\label{ch:delayed}
\begin{scopebox}
The delayed term is isolated without being discarded. Chapters~\ref{ch:feedback} and~\ref{ch:drift} supply its uniform closure; the direct and root calculations here remain part of that proof.
\end{scopebox}
This chapter uses the exact quotient and predictable-product identities in
Sections~\ref{base:con:product} and~\ref{base:con:variance}, and the cumulative running-quotient lemma in Chapter~\ref{ch:centered}. It proves a sharper
localization of the remaining difficulty. It does not establish a
uniform general refit theorem.

\paragraph{Subsequent closure.}
The delayed remainder isolated here is now controlled in
Chapter~\ref{ch:feedback} and
Chapter~\ref{ch:drift}. This chapter supplies
the paid-term reduction used there; its original open-status language
describes this intermediate reduction only.

There are nonroot stages $0,\ldots,h-1$ and one root stage $h$.
Put $d=q-2k$, $\rho=k/d<1$, and write the deterministic orthogonal
completions as
\[
 Q_{U,s}=\begin{pmatrix}A_{U,s}&B_{U,s}\\C_{U,s}&D_{U,s}\end{pmatrix},
 \qquad
 Q_{V,s}=\begin{pmatrix}A_{V,s}&B_{V,s}\\C_{V,s}&D_{V,s}\end{pmatrix}.
\]
Let $L_{U,s}=B_{U,s}B_{U,s}^{\T}$, and similarly $L_{V,s}$.
The current running-source innovation Grams are
\[
 J_{Y,s}=Y_s(x_s)\Pi_{G,s}Y_s(x_s)^{\T},\qquad
 J_{Z,s}=Z_s(x_s)\Pi_{H,s}Z_s(x_s)^{\T}.
\]
In this chapter these symbols include the actual running source. They are
not the unmarked Grams. The source martingale has $x_0=0$ and
$\E\|x_s\|^2\le\rho\beta$ at every stage, with
$\beta=\sum_s\|b_s\|^2$.

\paragraph{Audit of the cumulative lemma.}
For adapted forms $0\preceq Q_{s+1}\preceq Q_s\preceq I$, the
difference $D_s=Q_s-Q_{s+1}$ is measurable at time $s+1$.
Consequently, conditional Jensen applies to
$D_s^{1/2}x_{s+1}=\E[D_s^{1/2}x_m\mid\mathcal F_{s+1}]$.
This validates the direct-sum proof of
\[
 \sum_s\E[x_s^{\T}D_sx_s]+\E[x_m^{\T}Q_mx_m]
 \le4\E\|x_m\|^2.
\]
It does not require $D_s$ to be predictable at time $s$.
Applied to the exact HSS quotient forms, it gives
\begin{equation}\label{delayed:eq:lossbudget}
 \sum_{s<h}\E\bigl[\Phi(T_sx_s)-\Phi(T_{s+1}x_s)\bigr]
 \le4\rho\beta.
\end{equation}
The source is held fixed in each difference. We also use the pathwise
bound $\tr J_{Y,s}+\tr J_{Z,s}\le\Phi(T_sx_s)\le\|x_s\|^2$.

\section{Explicit deterministic impulse weights}
Define, for $s<r\le h$,
\[
 T_U(s,r)=A_{U,s}\cdots A_{U,r-1},\qquad
 F_V(s,r)=C_{V,s}A_{V,s+1}\cdots A_{V,r-1},
 \qquad c_V(s,r)=\F{F_V(s,r)}^2,
\]
with the empty product interpreted as $I$. Define $T_V,F_U,c_U$
by exchanging the two families.

For an upper envelope of either legal decoder readout, include both
discarded diagonal mean-filter readouts $d_Y,d_Z$ as separate orthogonal
outputs. This only enlarges the impulse norm. It is an analytical
envelope, not a change to the decoder. The enlarged deterministic
filter is still passive: the two row splits use separate $(P,K_G)$
and $(C,K_H)$ state-input pairs.

\begin{proposition}[Impulse partial traces]\label{delayed:prop:weights}
The forward impulse partial-trace weight of this enlarged filter is
\begin{align}
 \Lambda^+_{G,s}={}&kL_{U,s}
 +\sum_{s<r<h}c_V(s,r)T_U(s,r)L_{U,r}T_U(s,r)^{\T}\\
 &+\frac14c_V(s,h)T_U(s,h)T_U(s,h)^{\T},\qquad s<h.
 \label{delayed:eq:weights}
\end{align}
The transpose weight is obtained by swapping $U,V$. At the root,
$\Lambda^+_{G,h}=\Lambda^+_{H,h}=(k/4)I$.
For the forward doubly discarded readout,
$\Lambda_{G,s}=\Lambda^+_{G,s}$ and
$\Lambda_{H,s}\preceq\Lambda^+_{H,s}$; the opposite readout
interchanges these statements.
\end{proposition}
\begin{proof}
An impulse $K$ in the forward fresh register at stage $s$ has current
readouts $B_{U,s}^{\T}KC_{V,s}$ and
$B_{U,s}^{\T}KD_{V,s}$. Their squared norms sum to
$\F{B_{U,s}^{\T}K}^2$. Its retained state at $r>s$ is exactly
$T_U(s,r)^{\T}KF_V(s,r)$. The two forward readouts at $r<h$
then have combined squared norm
$\F{B_{U,r}^{\T}T_U(s,r)^{\T}KF_V(s,r)}^2$.
The retained root readout contributes one quarter of the retained-state
norm. Taking the partial trace over the $k$ columns of $K$ gives
\eqref{delayed:eq:weights}. A root fresh impulse occupies its own off-diagonal
block with factor $1/2$, giving $(k/4)I$. Exchanging streams proves
the other formula. The omitted legal readout is a squared norm and
therefore contributes a positive partial-trace matrix.
\end{proof}

\begin{lemma}[Incoming Bessel mass]\label{delayed:lem:bessel}
For every $r\le h$,
\[
 \sum_{s<r}F_V(s,r)^{\T}F_V(s,r)
 =I-T_V(0,r)^{\T}T_V(0,r),
 \qquad \sum_{s<r}c_V(s,r)\le k.
\]
The same statement holds for the other family.
\end{lemma}
\begin{proof}
Insert $C_{V,s}^{\T}C_{V,s}=I-A_{V,s}^{\T}A_{V,s}$ into
each summand. The resulting difference telescopes, including for
noncommuting completion blocks. Take traces for the scalar bound.
\end{proof}

\section{The terms already paid by the running quotient}
Define the nonnegative delayed nonroot term
\begin{equation}\label{delayed:eq:remainder}
\begin{split}
 \mathcal R_{\rm delay}={1\over d}
 \sum_{0\le s<r<h}\Bigl[&c_V(s,r)
 \E\tr\!\bigl(T_U(s,r)L_{U,r}T_U(s,r)^{\T}J_{Y,s}\bigr)\\
 &+c_U(s,r)
 \E\tr\!\bigl(T_V(s,r)L_{V,r}T_V(s,r)^{\T}J_{Z,s}\bigr)\Bigr].
\end{split}
\end{equation}
All weights in this expression are deterministic. All state factors
retain their original source--probe dependence.

\begin{proposition}[Paid direct terms and root reserve]\label{delayed:prop:paid}
For either legal discarded readout, the actual predictable variance
satisfies
\begin{equation}\label{delayed:eq:paid}
 \boxed{\displaystyle
 \mathcal P(b_\bullet)\le
 \left({4\over1-\rho}+{3\over4}\right)\rho^2\beta
 +\mathcal R_{\rm delay}.}
\end{equation}
In particular, the entire root contribution, including impulses created
at earlier stages and transported to the root, costs at most
$(3/4)\rho^2\beta$.
\end{proposition}
\begin{proof}
The exact predictable-variance formula and Proposition~\ref{delayed:prop:weights}
split its upper envelope into the direct nonroot term, the delayed
nonroot term \eqref{delayed:eq:remainder}, and root terms.
The direct nonroot term is
\[
 \rho\sum_{s<h}\E\tr(L_{U,s}J_{Y,s}+L_{V,s}J_{Z,s}).
\]
The exact conditional quotient-loss identity bounds each conditional
expected loss below by
$(1-\rho)\tr(L_{U,s}J_{Y,s}+L_{V,s}J_{Z,s})$.
Since $x_s$ is adapted, this applies to the running state without
changing its law. Equation~\eqref{delayed:eq:lossbudget} gives the contribution
$4\rho^2\beta/(1-\rho)$.

For root-transported forward impulses, $T_U(s,h)$ is a contraction,
so their total contribution is at most
\[
 {1\over4d}\sum_{s<h}c_V(s,h)\E\tr J_{Y,s}
 \le {\rho\beta\over4d}\sum_{s<h}c_V(s,h)
 \le\frac14\rho^2\beta.
\]
The transpose impulses cost the same. The fresh-root term is
$(\rho/4)\E(\tr J_{Y,h}+\tr J_{Z,h})\le\rho^2\beta/4$.
Add the three contributions.
\end{proof}

Thus a uniform bound on \eqref{delayed:eq:remainder}, at any sufficiently small
fixed positive $\rho$, is sufficient for the general fixed-space target.
The incoming Bessel bound cannot be used by replacing $J_{Y,s}$ with
$J_{Y,r}$: the observations undergo random feedback and the running
source changes between these stages. Neither replacement is an identity
or a consequence of the scalar quotient contraction. No such comparison
is used in Proposition~\ref{delayed:prop:paid}.

\endgroup

% ===== CROSSFEEDBACK =====
\begingroup
\let\R\relax
\newcommand{\R}{\mathbb R}
\let\E\relax
\newcommand{\E}{\mathbb E}
\let\Prb\relax
\newcommand{\Prb}{\mathbb P}
\let\T\relax
\newcommand{\T}{\mathsf T}
\let\F\relax
\newcommand{\F}[1]{\left\lVert#1\right\rVert_{F}}
\let\N\relax
\newcommand{\N}[1]{\left\lVert#1\right\rVert}
\let\ip\relax
\newcommand{\ip}[2]{\left\langle#1,#2\right\rangle_F}
\let\tr\relax
\newcommand{\tr}{\operatorname{tr}}
\let\rank\relax
\newcommand{\rank}{\operatorname{rank}}
\let\diag\relax
\newcommand{\diag}{\operatorname{blockdiag}}
\let\range\relax
\newcommand{\range}{\operatorname{range}}
\let\OPT\relax
\newcommand{\OPT}{\operatorname{OPT}}
\let\HSS\relax
\newcommand{\HSS}{\operatorname{HSS}}
\let\diver\relax
\newcommand{\diver}{\operatorname{div}}
\let\Var\relax
\newcommand{\Var}{\operatorname{Var}}
\let\UC\relax
\newcommand{\UC}{\mathcal U}
\let\VC\relax
\newcommand{\VC}{\mathcal V}
\let\err\relax
\newcommand{\err}{\mathcal E}
\let\proofstep\relax
\newcommand{\proofstep}[1]{\par\addvspace{.45\baselineskip}\Needspace{3\baselineskip}\textbf{#1}\par\nobreak}
\let\op\relax
\newcommand{\op}[1]{\left\lVert#1\right\rVert_{\mathrm{op}}}
\let\Sn\relax
\newcommand{\Sn}[2]{\left\lVert#1\right\rVert_{S_{#2}}}
\let\Cov\relax
\newcommand{\Cov}{\mathsf C}
\let\Loss\relax
\newcommand{\Loss}{\mathsf L}
\let\Fil\relax
\newcommand{\Fil}{\mathcal F}
\let\Id\relax
\newcommand{\Id}{I}
\let\supp\relax
\newcommand{\supp}{\operatorname{supp}}
\let\vecop\relax
\newcommand{\vecop}{\operatorname{vec}}
\let\Gammafun\relax
\newcommand{\Gammafun}{\operatorname{Gamma}}
\chapter{Transpose feedback and predictable forcing}\label{ch:crossfeedback}
\begin{scopebox}
This chapter records the valid forcing decomposition and the limitation of isolated-impulse estimates. Its coherent forcing is subsequently controlled by the complete-sequence argument of Chapter~\ref{ch:drift}.
\end{scopebox}
This chapter records the valid transpose product decomposition. The fixed-at-interval-start bilinear estimate must not be applied to a coefficient map that changes during the interval. The coherent term isolated below is closed by Chapter~\ref{ch:drift}.

\section{The actual product split}
Write
\[
 E_i(a)=Y_i(a)-C_i(a)^Tg_i,
 \qquad F_i(a)=Z_i(a)-P_i(a)^Th_i.
\]
The quotient contraction gives, pathwise for every fixed $a$,
\[
 \|E_i(a)\|_F^2\le\Phi_i(a)\le\|a\|^2,
 \qquad \|F_i(a)\|_F^2\le\Phi_i(a)\le\|a\|^2.
\]
Thus the sum over $a=e_\ell$ of either squared norm is at most $k$. Use the normalized canonical source with zero initial retained regressions. Let
\[
 \delta x_i=Xt_i^Tb_i^T,\qquad x_{i+1}=x_i+\delta x_i,
 \qquad\beta=\sum_i\|b_i\|^2.
\]
The exact recurrence is
\begin{align}\label{crossfeedback:eq:product}
 F_{i+1}(x_{i+1})={}&A_{V,i}^TF_i(x_i)
 -C_{V,i}^Tt_iE_i(x_i)^Th_i\notag\\
 &+A_{V,i}^TF_i(\delta x_i)
 -C_{V,i}^Tt_iE_i(\delta x_i)^Th_i.
\end{align}
Freezing the entire $H$ stream leaves the $E_i,F_i$ maps adapted to the forward filtration. The first forcing and source-increment forcing are centered and linear in $t_i$. The last term splits into a centered quadratic forcing and its predictable mean. The linear and centered quadratic terms are orthogonal in conditional second moment by $t_i\mapsto-t_i$ parity.

\section{One impulse and two martingale forcing terms}
An $F$-state impulse $W$ at time $n$, subsequently propagated by $A_V$, has delayed transpose cost
\[
 \frac1d\sum_{n\le s<r<h}c_U(s,r)
 \|B_{V,r}^TT_V(n,r)^TW\Pi_{H,s}\|_F^2
 \le\rho\|W\|_F^2.
\]
Drop the projector, use $\sum_{s<r}c_U(s,r)\le k$, and telescope
\[
 \sum_{r\ge n}T_V(n,r)B_{V,r}B_{V,r}^TT_V(n,r)^T\preceq I.
\]
This is a single-impulse operator norm bound, not an $\ell_2$ gain for coherent impulses.

For the source-increment forcing $U_i=A_{V,i}^TF_i(\delta x_i)$ and the raw quadratic forcing $V_i=-C_{V,i}^Tt_iE_i(\delta x_i)^Th_i$,
\[
 \E_G[\|U_i\|_F^2\mid\text{past }G,\text{ all }H]\le\rho\|b_i\|^2,
 \qquad
 \E_G[\|V_i\|_F^2\mid\text{past }G,\text{ all }H]\le\rho\|b_i\|^2.
\]
The second inequality uses $t_it_i^T=I$ and contraction. Centering reduces the variance. After $H$ is frozen, the output maps are deterministic, so distinct forward martingale times are orthogonal. Each of these two processes therefore costs at most $\rho^2\beta$ in the delayed seminorm. Combining them with adapted feedback still requires a triangle inequality or a cross-covariance calculation; constant-one addition is not asserted.

\section{The predictable mean must be retained}
Let $X_\ell$ be row $\ell$ of $X$ and define
\[
 H_i^\sharp=\sum_{\ell=1}^{k}E_i(e_\ell)\Pi_{G,i}X_\ell^T.
\]
The mean of the quadratic forcing is exactly
\[
 m_i=-\frac1d C_{V,i}^Tb_i^T(H_i^\sharp)^Th_i.
\]
The elementary bound $\|H_i^\sharp\|\le k$ yields an input-energy allowance of order $\rho^2\beta$ before the $C_V$ factors. It does not control the coherent sum after propagation. The complete-sequence estimate of Chapter~\ref{ch:drift} is what supplies that missing implication.

\section{Additional structure for transpose-only initialization}
Assume $E_0(a)=0$, $F_0(a)=av_0$, $\|v_0\|=1$, and freeze all $G$. For $j<i$ put
\[
 A_{ji}=C_{U,j}T_U(j+1,i),\qquad
 W_{ji}=\sum_\ell F_j(e_\ell)^Tg_j\Pi_{G,i}X_\ell^T.
\]
The vector $W_{ji}$ is measurable before the transpose frame $\eta_j$ is drawn. Exact unrolling gives
\[
 H_i^\sharp=-\sum_{j<i}A_{ji}^T\eta_jW_{ji}.
\]
Thus $H_i^\sharp$ is centered under $H$, with
\[
 \E_H[H_i^\sharp(H_i^\sharp)^T\mid G]
 =\frac1d\sum_{j<i}\E_H[\|\Pi_{H,j}W_{ji}\|^2\mid G]A_{ji}^TA_{ji}.
\]
The response contraction gives $\|W_{ji}\|\le k$, and the Bessel bound gives $\sum_{j<i}A_{ji}^TA_{ji}\preceq I$. Consequently
\[
 \E_H[H_i^\sharp(H_i^\sharp)^T\mid G]\preceq\frac{k^2}{d}I,
 \qquad \E_H\|H_i^\sharp\|^2\le\frac{k^3}{d}.
\]
For $N_i^T=b_i^T(H_i^\sharp)^T/d$, this improves the input estimate to
\[
 \sum_i\E\|N_i\|_F^2\le\rho^3\beta.
\]
It is still an input estimate, not by itself a delayed-output estimate.

Multiplication by $h_i$ destroys the naive centering:
\[
 \E_H[\eta_j^TA_{ji}h_i\mid\text{history before }\eta_j,G]
 =\frac{\|A_{ji}\|_F^2}{d}\Pi_{H,j}.
\]
The $\eta_j$ contribution in the retained-frame recursion is $A_{ji}^T\eta_j$; earlier terms are centered against it and later innovations have zero conditional mean. Thus the corresponding contribution to $(H_i^\sharp)^Th_i$ has mean
\[
 -\frac{\|A_{ji}\|_F^2}{d}W_{ji}^T\Pi_{H,j},
\]
which generally does not vanish. Centering $H_i^\sharp$ alone cannot eliminate coherent drift. The subsequent cancellation proof controls the actual retained-frame-weighted mean.

\endgroup

% ===== ROWMAJORANT =====
\begingroup
\let\R\relax
\newcommand{\R}{\mathbb R}
\let\E\relax
\newcommand{\E}{\mathbb E}
\let\Prb\relax
\newcommand{\Prb}{\mathbb P}
\let\T\relax
\newcommand{\T}{\mathsf T}
\let\F\relax
\newcommand{\F}[1]{\left\lVert#1\right\rVert_{F}}
\let\N\relax
\newcommand{\N}[1]{\left\lVert#1\right\rVert}
\let\ip\relax
\newcommand{\ip}[2]{\left\langle#1,#2\right\rangle_F}
\let\tr\relax
\newcommand{\tr}{\operatorname{tr}}
\let\rank\relax
\newcommand{\rank}{\operatorname{rank}}
\let\diag\relax
\newcommand{\diag}{\operatorname{blockdiag}}
\let\range\relax
\newcommand{\range}{\operatorname{range}}
\let\OPT\relax
\newcommand{\OPT}{\operatorname{OPT}}
\let\HSS\relax
\newcommand{\HSS}{\operatorname{HSS}}
\let\diver\relax
\newcommand{\diver}{\operatorname{div}}
\let\Var\relax
\newcommand{\Var}{\operatorname{Var}}
\let\UC\relax
\newcommand{\UC}{\mathcal U}
\let\VC\relax
\newcommand{\VC}{\mathcal V}
\let\err\relax
\newcommand{\err}{\mathcal E}
\let\proofstep\relax
\newcommand{\proofstep}[1]{\par\addvspace{.45\baselineskip}\Needspace{3\baselineskip}\textbf{#1}\par\nobreak}
\let\op\relax
\newcommand{\op}[1]{\left\lVert#1\right\rVert_{\mathrm{op}}}
\let\Sn\relax
\newcommand{\Sn}[2]{\left\lVert#1\right\rVert_{S_{#2}}}
\let\Cov\relax
\newcommand{\Cov}{\mathsf C}
\let\Loss\relax
\newcommand{\Loss}{\mathsf L}
\let\Fil\relax
\newcommand{\Fil}{\mathcal F}
\let\Id\relax
\newcommand{\Id}{I}
\let\supp\relax
\newcommand{\supp}{\operatorname{supp}}
\let\vecop\relax
\newcommand{\vecop}{\operatorname{vec}}
\let\Gammafun\relax
\newcommand{\Gammafun}{\operatorname{Gamma}}
\chapter{A positive-majorant obstruction and its cancellation requirement}\label{ch:rowmajorant}
\begin{scopebox}
The counterexample concerns a positive upper bound, not the decoder: the majorant grows while the actual predictable energy is zero. It explains why the later proof retains signed quotient terms.
\end{scopebox}
This chapter checks the proposed feedback closure based on the positive row quotient \(H_Y\). The fixed-forward unrolling and subsequent martingale timing identity are valid. The comparison \(J_Y\preceq H_Y\), however, can create a linearly growing spurious predictable budget even for an actual source with one nonzero deterministic source block.

Use the notation of Section~\ref{base:con:variance}, and start from the canonical transpose-only source, so \(H_{Y,0}=0\). Freeze all forward frames. For a fixed source coordinate \(x\), set

\[
h_j(x;G)=\mathbb E_H\operatorname{tr}J_{Z,j}(x)
=x^{\top}B_j(G)x.
\]

Here \(B_j\) is measurable with respect to the forward history before step \(j\), and \(0\preceq B_j\preceq I\), because \(\operatorname{tr}J_Z\le\Phi\le\|x\|^2\) pathwise. Define deterministic matrices

\[
K_U(j,s)=F_U(j+1,s)^{\top}C_{U,j}^{\top}C_{U,j}F_U(j+1,s),
\qquad j<s.
\]

The exact opposite-stream average is

\[
\boxed{
\mathbb E_H H_{Y,s}(x)
=\frac1d\sum_{j<s}h_j(x;G)K_U(j,s).}
\]

The proof simply iterates \(\mathbb E_HH_Y'=A_U^{\top}\mathbb E_HH_YA_U+
(\mathbb E_H\operatorname{tr}J_Z/d)C_U^{\top}C_U\). Fresh transpose projections remain centered after all forward frames are fixed. Consequently the identity remains valid upon substituting \(x=x_s(G)\).

There is also a valid favorable-order martingale identity. If \(x_s\) is the actual running forward source and \(j<s\), then

\[
\boxed{
\begin{aligned}
\mathbb E h_j(x_s;G)
={}&\mathbb E h_j(x_j;G)\\
&+\frac1d\sum_{r=j}^{s-1}\|b_r\|^2
\mathbb E\operatorname{tr}\!\left[B_jX\Pi_{G,r}X^{\top}\right].
\end{aligned}}
\]

Indeed \(B_j\) is measurable before every subsequent increment. Conditional centering kills the old-source/future-increment cross terms and all cross terms between distinct future increments. The remaining terms use the exact conditional covariance of \(\delta x_r\). In particular,

\[
\mathbb E h_j(x_s;G)
\le\mathbb E h_j(x_j;G)+\rho\sum_{r=j}^{s-1}\|b_r\|^2.
\]

Thus this part of the proposed argument does not incorrectly condition away future forward randomness.

\section{The deterministic feedback weights cannot be absorbed}\label{rowmajorant:the-deterministic-feedback-weights-cannot-be-absorbed}

Substituting \(J_{Y,s}\preceq H_{Y,s}\) into the forward predictable energy leads to the nonnegative majorant

\[
\overline{\mathcal P}_G
=\frac1{d^2}\sum_{j<s}
\operatorname{tr}[\Lambda_{G,s}K_U(j,s)]\,
\mathbb E h_j(x_s;G).
\]

For rank one, repeat the completions

\[
Q_U=\begin{pmatrix}0&-1\\1&0\end{pmatrix},
\qquad Q_V=I_2.
\]

Then

\[
\Lambda_{G,s}=1,\qquad
K_U(j,s)=\begin{cases}1,&s=j+1,\\0,&s>j+1.\end{cases}
\]

So each interior transpose innovation receives feedback weight one in the majorant. In contrast, \(\Lambda_{H,j}=0\) whenever at least one further nonroot swap follows before the root: its current transpose readout is zero since \(B_V=0\), and its retained impulse is killed by the following \(A_U=0\). Hence no uniform deterministic domination of these feedback weights by \(\Lambda_{H,j}\) is possible.

This is not merely an unreachable positive-matrix configuration. In the actual canonical process, \(Y=0\) remains identically zero and \(Z\) is fixed throughout this repeated-swap segment. Therefore \(J_Y=0\) and the genuine forward predictable energy is zero. But \(H_Y=C^2\) is generally positive.

If \(S=\|Z\|^2\), \(u_s=\mathbb E_H C_s^2\), and \(v_s=\mathbb E_H D_{Z,s}^2\), then

\[
u_{s+1}=\frac{S-u_s-v_s}{d},\qquad v_{s+1}=u_s.
\]

For \(d\ge2\), this stable affine recursion has limit \(u_s,v_s\to S/(d+2)=S/q\). In particular,

\[
\sum_{s<n}u_s=\frac{nS}{q}+O(S).
\]

\section{An actual nonzero running source before the swap segment}\label{rowmajorant:an-actual-nonzero-running-source-before-the-swap-segment}

The preceding loss can be realized with a source budget independent of the repeated-swap horizon. Take \(k=1\), \(X=g_0\), and initially use three row-identity completions \(Q_U=I\). For those three steps use the column rotation

\[
Q_V=\begin{pmatrix}c&-\sigma\\\sigma&c\end{pmatrix},
\qquad c=\sigma=1/\sqrt2.
\]

Let the deterministic source blocks be \(b_0=b_1=0\), \(b_2=1\), and zero thereafter. The first two retained/discarded frame updates imply

\[
\|X\Pi_{G,2}\|^2=\sigma^2=\frac12,
\qquad x_3=t_2X^{\top},\qquad
\mathbb E x_3^2=\frac1{2d}>0.
\]

Indeed the current frame span before \(t_2\) is \(\operatorname{span}(g_1,t_1)\), where \(g_1=cg_0+\sigma t_0\), and \(t_1\perp g_0,t_0\). Throughout these three steps, \(Y=0\), the transpose retained frame remains \(h_0\), and \(Z\) is multiplied by \(c\) each time. Thus after the prefix \(Z=c^3Z_0\), with energy \(S=c^6=1/8\) for a unit canonical source.

Now repeat the swap/identity pair from the preceding section for \(n\) steps. The actual source stays \(x_s=x_3\) because all later blocks vanish, so \(\beta=1\). Forward and transpose randomness are independent. The majorant arising from \(J_Y\preceq H_Y\) therefore satisfies

\[
\boxed{
\overline{\mathcal P}_{G,\mathrm{swap\ segment}}
=\frac{n}{16d^2q}+O(1),
\qquad \mathcal P_{G,\mathrm{swap\ segment}}=0.}
\]

The \(O(1)\) term depends on \(q\) but not on the number of repeated swaps. At \(q=4\), the slope is \(1/256\). The actual transpose predictable budget in this segment receives at most its last nonroot impulse and its root impulse, and stays bounded.

Thus neither the favorable martingale timing nor a larger fixed width ratio repairs this particular majorization: at any fixed \(q\), the spurious upper bound still grows with the horizon. A successful feedback argument must retain enough of \(J_Y=H_Y-D_YD_Y^{\top}-\Delta\Delta^{\top}\) to cancel the common-regression mismatch generated by transpose noise.

\endgroup

% ===== FEEDBACK =====
\begingroup
\let\R\relax
\newcommand{\R}{\mathbb R}
\let\E\relax
\newcommand{\E}{\mathbb E}
\let\Prb\relax
\newcommand{\Prb}{\mathbb P}
\let\T\relax
\newcommand{\T}{\mathsf T}
\let\F\relax
\newcommand{\F}[1]{\left\lVert#1\right\rVert_{F}}
\let\N\relax
\newcommand{\N}[1]{\left\lVert#1\right\rVert}
\let\ip\relax
\newcommand{\ip}[2]{\left\langle#1,#2\right\rangle_F}
\let\tr\relax
\newcommand{\tr}{\operatorname{tr}}
\let\rank\relax
\newcommand{\rank}{\operatorname{rank}}
\let\diag\relax
\newcommand{\diag}{\operatorname{blockdiag}}
\let\range\relax
\newcommand{\range}{\operatorname{range}}
\let\OPT\relax
\newcommand{\OPT}{\operatorname{OPT}}
\let\HSS\relax
\newcommand{\HSS}{\operatorname{HSS}}
\let\diver\relax
\newcommand{\diver}{\operatorname{div}}
\let\Var\relax
\newcommand{\Var}{\operatorname{Var}}
\let\UC\relax
\newcommand{\UC}{\mathcal U}
\let\VC\relax
\newcommand{\VC}{\mathcal V}
\let\err\relax
\newcommand{\err}{\mathcal E}
\let\proofstep\relax
\newcommand{\proofstep}[1]{\par\addvspace{.45\baselineskip}\Needspace{3\baselineskip}\textbf{#1}\par\nobreak}
\let\op\relax
\newcommand{\op}[1]{\left\lVert#1\right\rVert_{\mathrm{op}}}
\let\Sn\relax
\newcommand{\Sn}[2]{\left\lVert#1\right\rVert_{S_{#2}}}
\let\Cov\relax
\newcommand{\Cov}{\mathsf C}
\let\Loss\relax
\newcommand{\Loss}{\mathsf L}
\let\Fil\relax
\newcommand{\Fil}{\mathcal F}
\let\Id\relax
\newcommand{\Id}{I}
\let\supp\relax
\newcommand{\supp}{\operatorname{supp}}
\let\vecop\relax
\newcommand{\vecop}{\operatorname{vec}}
\let\Gammafun\relax
\newcommand{\Gammafun}{\operatorname{Gamma}}
\chapter{Frozen-stream feedback and future source marks}\label{ch:feedback}
This chapter proves a one-sided small-gain inequality and a companion transpose decomposition under the original HSS law. The coherent mean-forcing term isolated here is now bounded in Chapter~\ref{ch:drift}, closing this predictable-variance obstruction for arbitrary fixed nested bases. The separate basis-learning target remains open.

Write \(d=q-2k\), \(\rho=k/d\), and \(T_U(i,r)=A_{U,i}\cdots A_{U,r-1}\), with the analogous definition for \(T_V\). The stages \(0,\ldots,h-1\) are nonroot. Put

\[
c_V(s,r)=\|C_{V,s}T_V(s+1,r)\|_F^2,
\quad L_{U,r}=B_{U,r}B_{U,r}^{\mathsf T},
\quad u_U(i,r)=\|C_{U,i}T_U(i+1,r)B_{U,r}\|_F^2.
\]

The actual source is \(x_s=\sum_{j<s}Xt_j^{\mathsf T}b_j^{\mathsf T}\), where \(XX^{\mathsf T}=I_k\), the rows \(b_j\) are deterministic, and \(\beta=\sum_j\|b_j\|^2\). The formulas below allow either one-sided canonical initialization, with vanishing initial retained regressions. For the displayed free term use \(Y_0(a)=av_0\), \(Z_0(a)=0\), \(\|v_0\|=1\). In the actual forward-mark/opposite-response convention, \(Y_0(a)=0\) and \(Z_0(a)=-av_{H,0}\); the free term in \eqref{feedback:manual:1} is then zero. All subsequent upper bounds remain valid (and the term \(\rho^2\beta\) in \eqref{feedback:manual:2}, \eqref{feedback:manual:8}, and \eqref{feedback:manual:9} may be omitted).

\section{Exact projected feedback decomposition}\label{feedback:exact-projected-feedback-decomposition}

Define

\[
E_s(a)=Y_s(a)-C_s(a)^{\mathsf T}g_s,
\qquad F_s(a)=Z_s(a)-P_s(a)^{\mathsf T}h_s.
\]

The full response equations imply exactly

\[
E_{s+1}=A_{U,s}^{\mathsf T}E_s-C_{U,s}^{\mathsf T}K_{H,s}^{\mathsf T}g_s,
\qquad
F_{s+1}=A_{V,s}^{\mathsf T}F_s-C_{V,s}^{\mathsf T}K_{G,s}^{\mathsf T}h_s,
\]

where \(K_{G,s}=E_st_s^{\mathsf T}\) and \(K_{H,s}=F_s\eta_s^{\mathsf T}\). For example, the first identity follows by subtracting \(C_{s+1}^{\mathsf T}g_{s+1}\) from the update for \(Y_{s+1}\), using \(A_{V,s}g_{s+1}+B_{V,s}w_{s+1}=g_s\).

Fix the entire forward stream \(G\) and any vector \(a=a(G)\). The transpose increments \(K_{H,i}(a)\) are conditionally centered, so their different times are orthogonal in the expansion of \(E_s(a)\). Applying the final forward innovation projector before taking this second moment gives the \textbf{exact} identity

\[
\begin{split}
\mathbb E_H J_{Y,s}(a)
={}&\|v_0\Pi_{G,s}\|^2T_U(0,s)^{\mathsf T}aa^{\mathsf T}T_U(0,s)\\
&+\frac1d\sum_{i<s}T_U(i+1,s)^{\mathsf T}C_{U,i}^{\mathsf T}C_{U,i}T_U(i+1,s)\,
\mathbb E_H\operatorname{tr}\!\left(g_i\Pi_{G,s}g_i^{\mathsf T}J_{Z,i}(a)\right).
\end{split}\tag{\thechapter.1}\label{feedback:manual:1}
\]

Indeed the conditional second moment of \(K_H^{\mathsf T}MK_H\) is \(d^{-1}\operatorname{tr}(MJ_Z)I_k\). The factor \(g_i\Pi_{G,s}g_i^{\mathsf T}\) in \eqref{feedback:manual:1} is essential: replacing it by \(I_k\) would restore the unpaid one-leg mismatch.

Consequently the forward part of the delayed remainder satisfies

\[
\mathcal R_G\le\rho^2\beta+\mathcal B_G,
\qquad
\mathcal B_G=\frac1{d^2}\sum_{i<s<r<h}c_V(s,r)u_U(i,r)
\mathbb E\operatorname{tr}\!\left(g_i\Pi_{G,s}g_i^{\mathsf T}J_{Z,i}(x_s)\right).
\tag{\thechapter.2}\label{feedback:manual:2}
\]

For the free term, use \(\mathbb E[x_sx_s^{\mathsf T}]\preceq(\beta/d)I_k\), then the incoming bound \(\sum_{s<r}c_V(s,r)\le k\), and finally

\[
\sum_{r<h}\operatorname{tr}\!\left(T_U(0,r)L_{U,r}T_U(0,r)^{\mathsf T}\right)\le k.
\]

This proves its cost at most \(\beta k^2/d^2=\rho^2\beta\).

\section{Deterministic complementary-loss routing}\label{feedback:deterministic-complementary-loss-routing}

Set \(D_V(i,s)=I-T_V(i,s)T_V(i,s)^{\mathsf T}\). Let \(\Lambda^+_{H,i}\) be the transpose partial-trace weight of the enlarged mean filter that includes both discarded readouts, as defined in Chapter~\ref{ch:delayed}.

Then

\[
\boxed{\displaystyle
\sum_{i<s<r<h}c_V(s,r)u_U(i,r)D_V(i,s)
\preceq k\Lambda^+_{H,i}.}\tag{\thechapter.3}\label{feedback:manual:3}
\]

Here \(i\) is fixed on the left. To prove this, expand

\[
D_V(i,s)=\sum_{j=i}^{s-1}T_V(i,j)L_{V,j}T_V(i,j)^{\mathsf T}.
\]

For each \(j\), the coefficient of its summand is at most

\[
k\sum_{r>j+1}u_U(i,r)
\le\begin{cases}k^2,&j=i,\\
k\|C_{U,i}T_U(i+1,j)\|_F^2,&j>i.
\end{cases}
\]

The first inequality uses the incoming \(V\)-Bessel bound at each \(r\); the second uses the outgoing \(U\)-Bessel identity. These are exactly the direct and delayed nonroot terms in \(k\Lambda^+_{H,i}\). Its root term is nonnegative.

Also, conditional on the history at \(i\),

\[
\mathbb E[g_i\Pi_{G,s}g_i^{\mathsf T}\mid\mathcal F_i]
\preceq D_V(i,s).\tag{\thechapter.4}\label{feedback:manual:4}
\]

Indeed \(\mathbb E[g_s\mid\mathcal F_i]=T_V(i,s)^{\mathsf T}g_i\), both retained frames have orthonormal rows, and \(g_i\Pi_{G,s}=(g_i-T_V(i,s)g_s)\Pi_{G,s}\). The conditional Gram of the unprojected difference is exactly \(D_V(i,s)\).

Equations \eqref{feedback:manual:3}--\eqref{feedback:manual:4} imply that replacing \(x_s\) by the earlier \(x_i\) in \eqref{feedback:manual:2} gives

\[
\mathcal B_G^{\mathrm{old}}\le\rho\mathcal P_H^+,
\qquad
\mathcal P_H^+=\frac1d\sum_i\mathbb E\operatorname{tr}(\Lambda^+_{H,i}J_{Z,i}(x_i)).\tag{\thechapter.5}\label{feedback:manual:5}
\]

\section{A bilinear martingale lemma}\label{feedback:a-bilinear-martingale-lemma}

Let \(t_j\) be adapted \(k\)-row orthonormal frames with conditional mean zero and conditional entry covariance \(d^{-1}I_k\otimes\Pi_j\), where \(0\preceq\Pi_j\preceq I\). Let \(D_j\) and \(b_j\) be deterministic, \(XX^{\mathsf T}=I\), and define

\[
a=\sum_jXt_j^{\mathsf T}b_j^{\mathsf T},\quad R=\sum_jD_jt_j,
\quad S=\sum_jD_jD_j^{\mathsf T},\quad\beta_*=\sum_j\|b_j\|^2.
\]

For a fixed linear matrix map \(M(a)=\sum_\ell a_\ell M_\ell\), put \(G_M=\sum_\ell M_\ell M_\ell^{\mathsf T}\). Then

\[
\boxed{\displaystyle
\mathbb E\|R^{\mathsf T}M(a)\|_F^2
\le16\frac{\beta_*}{d}\operatorname{tr}(SG_M).}\tag{\thechapter.6}\label{feedback:manual:6}
\]

\textbf{Proof.} Write \(c_j=Xt_j^{\mathsf T}b_j^{\mathsf T}\), \(a_j=\sum_{l<j}c_l\), and \(R_j=\sum_{l<j}D_lt_l\). Expand the product as

\[
R^{\mathsf T}M(a)=\sum_jR_j^{\mathsf T}M(c_j)
+\sum_jt_j^{\mathsf T}D_j^{\mathsf T}M(a_j)
+\sum_jt_j^{\mathsf T}D_j^{\mathsf T}M(c_j).
\]

The first two sums are Hilbert-space martingales. Each has squared \(L^2\) norm at most \(\beta_*d^{-1}\operatorname{tr}(SG_M)\). For the first use \(\mathbb E[c_jc_j^{\mathsf T}\mid\mathcal F_j]\preceq\|b_j\|^2I/d\) and \(\mathbb E R_jR_j^{\mathsf T}=\sum_{l<j}D_lD_l^{\mathsf T}\). For the second use \(t_jt_j^{\mathsf T}=I\) and \(\mathbb E[a_ja_j^{\mathsf T}]\preceq\beta_*I/d\).

Center each summand of the third sum. Its centered part is another martingale with squared \(L^2\) norm at most the same bound: orthonormal rows remove the apparent fourth moment, because

\[
\|t_j^{\mathsf T}D_j^{\mathsf T}M(c_j)\|_F^2
=\|D_j^{\mathsf T}M(c_j)\|_F^2.
\]

Its conditional mean is \(m_j=d^{-1}\Pi_jX^{\mathsf T}\Gamma_j\), where row \(\ell\) of \(\Gamma_j\) is \(b_jD_j^{\mathsf T}M_\ell\). Therefore

\[
\Big\|\sum_jm_j\Big\|_F
\le\frac{\sqrt{\beta_*}}d\sqrt{\operatorname{tr}(SG_M)}.
\]

The \(L^2\) triangle inequality gives constant \((3+d^{-1/2})^2\le16\) for \(d\ge1\), proving \eqref{feedback:manual:6}. No independence between successive projectors is assumed.

\section{The actual future mark}\label{feedback:the-actual-future-mark}

Apply \eqref{feedback:manual:6} conditionally at \(i\), with \(a=x_s-x_i\), \(M(a)=Z_i(a)\Pi_{H,i}\), and the deterministic representation

\[
g_i-T_V(i,s)g_s
=D_V(i,s)g_i-\sum_{j=i}^{s-1}T_V(i,s)T_V(j+1,s)^{\mathsf T}C_{V,j}^{\mathsf T}t_j.
\]

The Gram of the constant term is \(D_V(i,s)^2\); the sum of coefficient Grams is \(D_V(i,s)-D_V(i,s)^2\). Combining \eqref{feedback:manual:6} with the elementary second-moment bound for the constant term gives, conservatively,

\[
\boxed{\displaystyle
\mathbb E\!\left[\operatorname{tr}\!\left(g_i\Pi_{G,s}g_i^{\mathsf T}J_{Z,i}(x_s-x_i)\right)\middle|\mathcal F_i\right]
\le25\frac{\beta_{i:s}}d\operatorname{tr}\!\left(D_V(i,s)G_{Z,i}\right),}\tag{\thechapter.7}\label{feedback:manual:7}
\]

where \(\beta_{i:s}=\sum_{j=i}^{s-1}\|b_j\|^2\) and \(G_{Z,i}=\sum_\ell J_{Z,i}(e_\ell)\). This estimate retains the complementary matrix \(D_V(i,s)\), including its protected directions.

Define the deterministic-source trace budget

\[
\mathfrak T_H^+=\frac1d\sum_i\sum_\ell\mathbb E\operatorname{tr}(\Lambda^+_{H,i}J_{Z,i}(e_\ell)).
\]

Using \(x_s=x_i+(x_s-x_i)\), the squared-norm inequality with factor two, \eqref{feedback:manual:3}, and \eqref{feedback:manual:7}, now proves

\[
\boxed{\displaystyle
\mathcal R_G\le\rho^2\beta+2\rho\mathcal P_H^+
+50\frac{\beta k}{d^2}\mathfrak T_H^+.}\tag{\thechapter.8}\label{feedback:manual:8}
\]

The enlarged filter consists of the union of the two legal discarded readouts. If each legal deterministic-source variance obeys the supplied bound \(\vartheta\|a\|^2\), then \(\mathfrak T_H^+\le2\vartheta k\). Thus \eqref{feedback:manual:8} yields

\[
\mathcal R_G\le(1+100\vartheta)\rho^2\beta+2\rho\mathcal P_H^+.\tag{\thechapter.9}\label{feedback:manual:9}
\]

Swapping streams in the preceding proof would place the source in the stream being averaged, so that swap is not licensed. The following product decomposition supplies a companion estimate with its predictable mean retained explicitly.

\section{Companion transpose estimate}\label{feedback:companion-transpose-estimate}

Define a seminorm on transpose-state sequences by

\[
\|f\|_{\mathscr D_H}^2=\frac1d\sum_{s<r<h}c_U(s,r)
\|B_{V,r}^{\mathsf T}T_V(s,r)^{\mathsf T}f_s\Pi_{H,s}\|_F^2.
\tag{\thechapter.10}\label{feedback:manual:10}
\]

Then \(\mathcal R_H=\mathbb E\|(F_s(x_s))_s\|_{\mathscr D_H}^2\). For a state \(Z\) introduced at stage \(j\), dropping the projectors, summing \(\sum_{j\le s<r}c_U(s,r)\le k\), and applying outgoing \(V\) Bessel gives the deterministic single-impulse estimate

\[
\|(T_V(j,s)^{\mathsf T}Z)_{s\ge j}\|_{\mathscr D_H}^2\le\rho\|Z\|_F^2.
\tag{\thechapter.11}\label{feedback:manual:11}
\]

Let \(\delta x_i=Xt_i^{\mathsf T}b_i^{\mathsf T}\) and \(f_i=F_i(x_i)\). Since \(x_0=0\), \(f_0=0\) for either canonical response orientation. Linearity in the source gives the exact product recurrence

\[
\begin{split}
f_{i+1}={}&A_{V,i}^{\mathsf T}f_i
-C_{V,i}^{\mathsf T}t_iE_i(x_i)^{\mathsf T}h_i\\
&+A_{V,i}^{\mathsf T}F_i(\delta x_i)
-C_{V,i}^{\mathsf T}t_iE_i(\delta x_i)^{\mathsf T}h_i.
\end{split}\tag{\thechapter.12}\label{feedback:manual:12}
\]

With \(X_\ell\) denoting row \(\ell\) of \(X\), put

\[
U_i=\frac1d\left(\sum_{\ell=1}^k E_i(e_\ell)\Pi_{G,i}X_\ell^{\mathsf T}\right)b_i.
\tag{\thechapter.13}\label{feedback:manual:13}
\]

Conditional covariance gives \(\mathbb E[t_iE_i(\delta x_i)^{\mathsf T}\mid\mathcal F_i]=U_i^{\mathsf T}\). Also \(E_i(e_\ell)\Pi_{G,i}=Y_i(e_\ell)\Pi_{G,i}\). Thus \eqref{feedback:manual:13} is the actual source-response mean port.

Freeze the entire \(H\) stream. The seminorm \eqref{feedback:manual:10} and propagation matrices are now fixed, and the future \(G\) conditional law is preserved by independence of the two streams. Split \eqref{feedback:manual:12} into the predictable-linear forcing, the source-increment forcing, the centered quadratic forcing, and the predictable quadratic mean. The first three are martingale differences in this enlarged \(G\) filtration, so different times within each of those sequences are orthogonal after deterministic propagation. No mutual orthogonality of the three sequences is asserted.

The predictable-linear forcing has expected squared seminorm at most \(\rho\mathcal P_G^+\). Its exact second moment leaves \(h_i\Pi_{H,s}h_i^{\mathsf T}\); unfreezing \(H\) and applying \eqref{feedback:manual:3}--\eqref{feedback:manual:4} with the streams interchanged proves the claim. The mark here is \(x_i\), fixed before all future \(H\) innovations.

The quotient identity gives, pathwise, \(\|E_i(a)\|_F^2,\|F_i(a)\|_F^2\le\Phi_i(a)\le\|a\|^2\). Consequently the sums of these squared norms over \(a=e_\ell\) are each at most \(k\). Together with \(\mathbb E[\delta x_i\delta x_i^{\mathsf T}\mid\mathcal F_i]\preceq\|b_i\|^2I/d\) and \eqref{feedback:manual:11}, this pays at most \(\rho^2\beta\) for the source-increment forcing. It also pays at most \(\rho^2\beta\) for the centered quadratic forcing: centering reduces its variance, and \(\|C_{V,i}^{\mathsf T}t_iE_i(\delta x_i)^{\mathsf T}h_i\|_F\le\|E_i(\delta x_i)\|_F\). This uses orthonormal rows and requires no fourth-moment estimate.

Define the remaining coherent state and its cost by

\[
m_0=0,\quad m_{i+1}=A_{V,i}^{\mathsf T}m_i-C_{V,i}^{\mathsf T}U_i^{\mathsf T}h_i,
\qquad \mathcal D_{\rm drift}=\mathbb E\|m\|_{\mathscr D_H}^2.
\tag{\thechapter.14}\label{feedback:manual:14}
\]

The \(L^2\) triangle inequality for the four sequences proves

\[
\boxed{\sqrt{\mathcal R_H}\le\sqrt{\rho\mathcal P_G^+}
+2\rho\sqrt\beta+\sqrt{\mathcal D_{\rm drift}}.}
\tag{\thechapter.15}\label{feedback:manual:15}
\]

\section{Explicit reduction to the mean-forcing term}\label{feedback:explicit-reduction-to-the-mean-forcing-term}

The cumulative quotient estimate pays the direct and root terms:

\[
\mathcal P^+:=\mathcal P_G^++\mathcal P_H^+
\le\left(\frac4{1-\rho}+\frac34\right)\rho^2\beta+\mathcal R_G+\mathcal R_H.
\tag{\thechapter.16}\label{feedback:manual:16}
\]

Use \((a+b+c)^2\le2a^2+4b^2+4c^2\) in \eqref{feedback:manual:15}, followed by \eqref{feedback:manual:9}. For \(0<\rho<1/2\) this gives

\[
\boxed{\mathcal P^+\le
\frac{\left[\frac4{1-\rho}+\frac34+17+100\vartheta\right]\rho^2\beta
+4\mathcal D_{\rm drift}}{1-2\rho}.}
\tag{\thechapter.17}\label{feedback:manual:17}
\]

For \(\vartheta=\rho(2+\rho)/(1-\rho^2)\) and \(\rho=1/4\), the bound is \(\mathcal P^+\le10.385417\beta+8\mathcal D_{\rm drift}\). The optional forward free term is retained; the actual opposite-response source permits deleting it.

There is a pathwise input-energy estimate \(\sum_i\|U_i\|_F^2\le\rho^2\beta\): Cauchy--Schwarz over \(\ell\) and the quotient bound give \(\|\sum_\ell E_i(e_\ell)\Pi_{G,i}X_\ell^{\mathsf T}\|^2\le k^2\). The single-impulse gain \eqref{feedback:manual:11} alone does not bound the coherent sum \eqref{feedback:manual:14}. That missing step is proved in Chapter~\ref{ch:drift} by introducing a corrected retained coefficient and telescoping two passive energy identities. Its conclusion is

\[
\mathcal D_{\rm drift}\le4\sum_i\mathbb E\|U_i\|_F^2\le4\rho^2\beta.
\tag{\thechapter.18}\label{feedback:manual:18}
\]

No independence of \(U_i\) from future readouts or temporal orthogonality of the mean is assumed. Substituting \eqref{feedback:manual:18} in \eqref{feedback:manual:17} proves

\[
\mathcal P^+\le C_P(\rho)\rho^2\beta,\qquad
C_P(\rho)=\frac{4/(1-\rho)+3/4+33+100\vartheta}{1-2\rho},
\quad 0<\rho<1/2.
\tag{\thechapter.19}\label{feedback:manual:19}
\]

This closes the depth- and rank-uniform predictable bound for fixed bases. The source-to-refit assembly and explicit widths are recorded in Chapter~\ref{ch:uniform}. The stronger deterministic whole-array transport bound in Chapter~\ref{ch:frame_transport} supplies an independent proof of drift control; its auxiliary randomness is only a proof device.

\section{Closing the original width endpoint}\label{feedback:closing-the-original-width-endpoint}

The condition \texttt{rho\textless{}1/2} in \eqref{feedback:manual:19} is an artifact of the equal-weight squared-norm bounds. Before those bounds, \eqref{feedback:manual:5} and \eqref{feedback:manual:7} give

\[
\mathcal B_G^{\rm old}\le\rho\mathcal P_H^+,
\qquad \mathcal B_G^{\rm future}\le50\vartheta\rho^2\beta.
\]

Use \(\|a+b\|^2\le(3/2)\|a\|^2+3\|b\|^2\) in this split and in \eqref{feedback:manual:15} after applying \eqref{feedback:manual:18}. The result is

\[
\mathcal R_G\le(1+150\vartheta)\rho^2\beta+\tfrac32\rho\mathcal P_H^+,
\qquad
\mathcal R_H\le48\rho^2\beta+\tfrac32\rho\mathcal P_G^+.
\]

Together with \eqref{feedback:manual:16}, this proves, for \(0<\rho\le1/2\),

\[
\boxed{\mathcal P^+\le
\frac{4/(1-\rho)+3/4+49+150\vartheta}{1-3\rho/2}\rho^2\beta.}
\tag{\thechapter.20}\label{feedback:manual:20}
\]

At the original \texttt{q=4k}, the coefficient before \texttt{rho\^{}2\ beta} equals \texttt{1231}. The resulting source and residual assembly gives expected fixed-space error at most \texttt{1607\ d\_S(A)\^{}2} with the original \texttt{8k} queries. This constant is deliberately conservative; the \texttt{36k}-query certificate is \texttt{15/4}, and the relative-error certificate remains \texttt{O(k/epsilon)}.

\endgroup

% ===== DRIFT =====
\begingroup
\let\E\relax
\newcommand{\E}{\mathbb E}
\let\T\relax
\newcommand{\T}{\mathsf T}
\let\F\relax
\newcommand{\F}[1]{\lVert#1\rVert_F}
\let\tr\relax
\newcommand{\tr}{\operatorname{tr}}
\chapter{Coherent drift: retained-coefficient cancellation}\label{ch:drift}
This chapter proves a uniform expected gain for the coherent forcing isolated
in Chapter~\ref{ch:feedback}. The gain uses the original
conditional Haar law and predictable forcing. It does not replace the
forcing by independent inputs, and it does not assert a pathwise bound
on the delayed quadratic functional. The proof below is self-contained.
The subsequent fixed-space consequence explicitly uses the earlier
feedback reduction and the source-to-refit assembly of Chapter~\ref{ch:interfaces}.

\section{The forced process and the delayed functional}

Fix stages $0,\ldots,H-1$ and deterministic orthogonal completions
\[
 Q_{U,i}=\begin{pmatrix}A_{U,i}&B_{U,i}\\C_{U,i}&D_{U,i}\end{pmatrix},
 \qquad
 Q_{V,i}=\begin{pmatrix}A_{V,i}&B_{V,i}\\C_{V,i}&D_{V,i}\end{pmatrix},
\]
whose blocks are $k\times k$. Write
\[
 T_U(s,r)=A_{U,s}\cdots A_{U,r-1},\qquad
 T_V(s,r)=A_{V,s}\cdots A_{V,r-1},
 \qquad c_U(s,r)=\F{C_{U,s}T_U(s+1,r)}^2.
\]
The products are identity matrices when their endpoints coincide.

Let $(h_i,\chi_i)$ be the retained and discarded row frames of the
transpose probe, with $2k$ orthonormal rows in $\mathbb R^q$.
Set $d=q-2k\ge k$ and
\[
 \Pi_i=I-h_i^\T h_i-\chi_i^\T\chi_i.
\]
Conditional on the history $\mathcal F_i$, draw $\eta_i$, a uniform
$k$-row orthonormal frame in this $d$-dimensional complement. Thus
\[
 \E[\eta_i\mid\mathcal F_i]=0,\qquad
 \E[\eta_i^\T R\eta_i\mid\mathcal F_i]=\frac{\tr R}{d}\Pi_i
 \quad(R\in\mathbb R^{k\times k}\text{ predictable}),
\]
and
\[
 h_{i+1}=A_{U,i}^\T h_i+C_{U,i}^\T\eta_i,\qquad
 \chi_{i+1}=B_{U,i}^\T h_i+D_{U,i}^\T\eta_i.
\]
The filtration may include an independent entire forward probe. The
conditional laws above remain valid with that enlargement.

Let $U_i\in\mathbb R^{k\times k}$ be any predictable square-integrable
inputs. Starting with $f_0=0$, define
\begin{equation}\label{drift:eq:forced}
 f_{i+1}=A_{V,i}^\T f_i-C_{V,i}^\T U_i^\T h_i.
\end{equation}
Its delayed nonroot cost is
\begin{equation}\label{drift:eq:cost}
 \mathfrak D(f)=\frac1d\sum_{0\le s<r<H}c_U(s,r)
 \E\F{B_{V,r}^\T T_V(s,r)^\T f_s\Pi_s}^2.
\end{equation}

\begin{theorem}[Coherent retained-forcing gain]\label{drift:thm:gain}
Under these assumptions,
\begin{equation}\label{drift:eq:gain}
 \boxed{\mathfrak D(f)\le4\sum_{i<H}\E\F{U_i}^2.}
\end{equation}
The constant is independent of the rank, horizon, and completions.
\end{theorem}

\begin{proof}
Introduce the auxiliary matrix filter
\begin{equation}\label{drift:eq:M}
 M_0=0,\qquad M_{i+1}=A_{U,i}^\T(M_iA_{V,i}+U_iC_{V,i}),
\end{equation}
and the corrected retained coefficient
\[
 R_i=f_i h_i^\T+M_i^\T,\qquad K_i=f_i\eta_i^\T.
\]
Because $h_i\eta_i^\T=0$ and $h_ih_i^\T=I$, direct multiplication
of \eqref{drift:eq:forced} by $h_{i+1}^\T$ gives
\begin{align*}
 f_{i+1}h_{i+1}^\T
 &=A_{V,i}^\T f_i h_i^\T A_{U,i}
   +A_{V,i}^\T K_iC_{U,i}-C_{V,i}^\T U_i^\T A_{U,i},\\
 M_{i+1}^\T
 &=A_{V,i}^\T M_i^\T A_{U,i}+C_{V,i}^\T U_i^\T A_{U,i}.
\end{align*}
The two input terms cancel exactly. Consequently
\begin{equation}\label{drift:eq:R}
 R_0=0,\qquad
 R_{i+1}=A_{V,i}^\T R_iA_{U,i}+A_{V,i}^\T K_iC_{U,i}.
\end{equation}
The $K_i$ are martingale differences. All matrices transporting them
in \eqref{drift:eq:R} are deterministic. Therefore, for every $r$,
\[
 R_r=\sum_{s<r}T_V(s,r)^\T K_sC_{U,s}T_U(s+1,r),
\]
and distinct summands are orthogonal in expected Frobenius inner product.
The conditional second-moment identity for $\eta_s$ yields the exact
observability identity
\begin{equation}\label{drift:eq:observability}
 \boxed{\E\sum_{r<H}\F{B_{V,r}^\T R_r}^2=\mathfrak D(f).}
\end{equation}
In particular, the projectors in \eqref{drift:eq:cost} have not been dropped.

It remains to bound these actual coefficient outputs. Define
\[
 e_i=B_{V,i}^\T f_i-D_{V,i}^\T U_i^\T h_i,\qquad
 N_i=M_iB_{V,i}+U_iD_{V,i}.
\]
Orthogonality of $Q_{V,i}$ gives the first exact energy identity
\begin{equation}\label{drift:eq:fenergy}
 \F{f_{i+1}}^2+\F{e_i}^2=\F{f_i}^2+\F{U_i}^2.
\end{equation}
Here $\F{U_i^\T h_i}=\F{U_i}$ because $h_i$ has orthonormal rows.
For the matrix filter, put $S_i=M_iA_{V,i}+U_iC_{V,i}$. Splitting
$S_i$ by $[A_{U,i}\ B_{U,i}]$ gives a second exact identity:
\begin{equation}\label{drift:eq:Menergy}
 \F{M_{i+1}}^2+\F{B_{U,i}^\T S_i}^2+\F{N_i}^2
 =\F{M_i}^2+\F{U_i}^2.
\end{equation}
Both statements hold pathwise, even for history-dependent $U_i$.
Finally, the identity
\begin{equation}\label{drift:eq:output}
 B_{V,i}^\T R_i=e_i h_i^\T+N_i^\T
\end{equation}
follows by canceling $D_{V,i}^\T U_i^\T$ between its two right-hand
terms. Thus, pathwise,
\[
 \sum_{i<H}\F{B_{V,i}^\T R_i}^2
 \le2\sum_{i<H}\F{e_i}^2+2\sum_{i<H}\F{N_i}^2
 \le4\sum_{i<H}\F{U_i}^2.
\]
Take expectations and use \eqref{drift:eq:observability}.
\end{proof}

\section{Closing the predictable-variance reduction}

For the actual HSS source, let $X$ have orthonormal rows and
\[
 x_0=0,\quad x_{i+1}-x_i=Xt_i^\T b_i^\T,\quad
 \beta=\sum_i\|b_i\|^2,\quad \rho=k/d.
\]
Write $E_i(a)=Y_i(a)-C_i(a)^\T g_i$. The exact quotient contraction
implies $\F{E_i(a)}^2\le\|a\|^2$. The predictable mean port in the
transpose product expansion is
\[
 U_i=\frac1d\left(\sum_{\ell=1}^k E_i(e_\ell)\Pi_{G,i}X_\ell^\T\right)b_i,
\]
where $X_\ell$ is row $\ell$ of $X$. Cauchy--Schwarz gives, pathwise,
\begin{equation}\label{drift:eq:inputbudget}
 \sum_i\F{U_i}^2\le\rho^2\beta.
\end{equation}
This input is predictable for the joint filtration and also for the
transpose filtration with the entire independent forward probe revealed.
Theorem~\ref{drift:thm:gain} therefore proves
\begin{equation}\label{drift:eq:drift}
 \boxed{\mathcal D_{\rm drift}\le4\rho^2\beta.}
\end{equation}
No additional estimate on the transpose-only initialization is needed.

For clarity, the earlier feedback reduction has the following three
interfaces, with $\vartheta=\rho(2+\rho)/(1-\rho^2)$:
\begin{align*}
 \mathcal R_G&\le(1+100\vartheta)\rho^2\beta+2\rho\mathcal P_H^+,\\
 \sqrt{\mathcal R_H}&\le\sqrt{\rho\mathcal P_G^+}
                         +2\rho\sqrt\beta+\sqrt{\mathcal D_{\rm drift}},\\
 \mathcal P^+&\le\left(\frac4{1-\rho}+\frac34\right)\rho^2\beta
                         +\mathcal R_G+\mathcal R_H.
\end{align*}
They are proved in Chapter~\ref{ch:feedback} and
Chapter~\ref{ch:delayed}; the first conservatively
retains an optional free term. Substitution of \eqref{drift:eq:drift} gives
$\mathcal R_H\le2\rho\mathcal P_G^++32\rho^2\beta$. Hence
\begin{equation}\label{drift:eq:predictable}
 \boxed{\mathcal P^+\le C_P(\rho)\rho^2\beta,\qquad
 C_P(\rho)=\frac{\frac4{1-\rho}+\frac34+33+100\vartheta}{1-2\rho},
 \quad 0<\rho<\frac12.}
\end{equation}
This is uniform in the depth and rank.

\paragraph{The original width is also covered.}
The restriction $\rho<1/2$ in \eqref{drift:eq:predictable} is only a consequence
of choosing equal weights in Young's inequality. In the forward feedback
proof, the old-mark and future-mark contributions satisfy, respectively,
$\mathcal B_G^{\rm old}\le\rho\mathcal P_H^+$ and
$\mathcal B_G^{\rm future}\le50\vartheta\rho^2\beta$.
Using $\|a+b\|^2\le(3/2)\|a\|^2+3\|b\|^2$ instead gives
\[
 \mathcal R_G\le(1+150\vartheta)\rho^2\beta
                         +\tfrac32\rho\mathcal P_H^+.
\]
The same inequality applied to
$\sqrt{\mathcal R_H}\le\sqrt{\rho\mathcal P_G^+}+4\rho\sqrt\beta$
gives $\mathcal R_H\le\tfrac32\rho\mathcal P_G^++48\rho^2\beta$.
Consequently a second certificate, valid throughout $0<\rho\le1/2$, is
\begin{equation}\label{drift:eq:originalwidth}
 \boxed{\mathcal P^+\le\widetilde C_P(\rho)\rho^2\beta,\qquad
 \widetilde C_P(\rho)=
 \frac{4/(1-\rho)+3/4+49+150\vartheta}{1-3\rho/2}.}
\end{equation}
At the original width $q=4k$, $\rho=1/2$, $\vartheta=5/3$, and
$\widetilde C_P(1/2)=1231$. Thus that width also has a uniform constant;
no increase in the original $8k$ query budget is needed.

\section{What this establishes, and the remaining general target}

The product decomposition of Section~\ref{base:con:product} now bounds each complete chain-source
correction by $C_{\rm ch}\beta$, where
\[
 C_{\rm ch}=\left[\rho\sqrt{C_P(\rho)}+\rho
                       +\sqrt{\rho\vartheta+\gamma_{k,d}}\right]^2,
 \qquad \gamma_{k,d}\le2\rho^2.
\]
Together with the off-chain source and orthogonality interfaces of Chapter~\ref{ch:interfaces},
put $C_* =\max\{C_{\rm ch},\vartheta\rho/(1-\rho)\}$.
The residual-centered assembly in Chapter~\ref{ch:centered} gives
\begin{equation}\label{drift:eq:refit}
 \E\F{A-\widehat B}^2
 \le[1+12\rho+4C_*]\,d_{\mathcal S}(A)^2.
\end{equation}
In particular it yields a constant-factor fixed-space refit with
$O(k)$ queries and no logarithmic factor. The source-decomposition and orthogonality arguments used at this step are in Chapter~\ref{ch:interfaces} and Sections~\ref{base:new:source} and~\ref{base:new:carleson}.

At $q=4k$, substitute \eqref{drift:eq:originalwidth} instead of
\eqref{drift:eq:predictable} in $C_{\rm ch}$. Since
\[
 C_{\rm ch}\le
 \left[\tfrac12\sqrt{1231}+\tfrac12+\sqrt{4/3}\right]^2<400,
 \qquad C_{\rm off}=5/3,
\]
\eqref{drift:eq:refit} gives the explicit, deliberately conservative bound
\[
 \boxed{\E\F{A-\widehat B}^2\le1607\,d_{\mathcal S}(A)^2,
 \qquad\text{queries}=8k.}
\]
Here $\sqrt{1231}<36$, $\sqrt{4/3}<3/2$, and
$1+12\rho=7$. The smaller constants below use wider probes.

\begin{corollary}[Explicit widths without a logarithm]
For arbitrary fixed nested bases, the original two-stream refit with
$q=18k$ obeys
\[
 \E\F{A-\widehat B}^2\le\frac{15}{4}d_{\mathcal S}(A)^2,
 \qquad \text{queries}=36k.
\]
For $0<\varepsilon\le1$, the choice
\[
 \boxed{q=\left\lceil k\left(2+\frac{32}{\varepsilon}\right)\right\rceil}
\]
suffices for
\[
 \boxed{\E\F{A-\widehat B}^2\le(1+\varepsilon)d_{\mathcal S}(A)^2,
 \qquad \text{queries}=2q=O(k/\varepsilon).}
\]
The source and residual-assembly premises are the same as in
\eqref{drift:eq:refit}.
\end{corollary}
\begin{proof}
For $0<\rho\le1/16$, monotonicity of the positive terms in
\eqref{drift:eq:predictable} gives
\[
 C_P(\rho)\le C_P(1/16)<64,\qquad
 \vartheta/\rho\le176/85<9/4.
\]
Since $\gamma_{k,d}\le2\rho^2$,
\[
 \sqrt{C_{\rm ch}}\le\rho\left(8+1+\sqrt{17/4}\right)
 <\frac{45}{4}\rho,
 \qquad C_{\rm ch}<128\rho^2.
\]
Also $\vartheta\rho/(1-\rho)\le(12/5)\rho^2<128\rho^2$.
Thus \eqref{drift:eq:refit} is bounded by
\[
 [1+12\rho+512\rho^2]d_{\mathcal S}(A)^2.
\]
For $q=18k$, substitute $\rho=1/16$. For the accuracy choice,
$\rho\le\varepsilon/32\le1/16$ and
\[
 12\rho+512\rho^2\le\frac38\varepsilon+\frac12\varepsilon^2
 \le\frac78\varepsilon<\varepsilon.
\]
\end{proof}

The full public objective additionally requires constructing a nested
rank-$k$ space $\mathcal S$ from $O(k)$ queries with
\[
 \E_{\rm learn}d_{\mathcal S}(A)^2
 \le C_b L\min_{B\in\mathrm{HSS}(L,k)}\F{A-B}^2.
\]
That basis-learning statement is not proved here. If it is supplied,
an independent refit satisfying \eqref{drift:eq:refit} completes the target by
conditioning on the learned space. Thus the coherent-drift obstruction
in fixed-space refitting is removed; the full best-HSS benchmark still
has a separate unresolved premise.

\section{Checks}
The accompanying \texttt{hss\_coherent\_drift\_check.py} verifies the
change of variables, the centered recurrence \eqref{drift:eq:R}, the output
identity \eqref{drift:eq:output}, and both pathwise energy identities in 24
chains with ranks one through four and horizons one, eight, and forty.
It includes random and nearly protected completions and predictable
history-dependent inputs. The largest reported normalized algebraic
residual is below $7\cdot10^{-15}$. These checks are consistency tests;
the expectation identity is justified by the martingale argument above.

\endgroup

% ===== FRAME_TRANSPORT =====
\begingroup
\let\R\relax
\newcommand{\R}{\mathbb R}
\let\E\relax
\newcommand{\E}{\mathbb E}
\let\Prb\relax
\newcommand{\Prb}{\mathbb P}
\let\T\relax
\newcommand{\T}{\mathsf T}
\let\F\relax
\newcommand{\F}[1]{\left\lVert#1\right\rVert_{F}}
\let\N\relax
\newcommand{\N}[1]{\left\lVert#1\right\rVert}
\let\ip\relax
\newcommand{\ip}[2]{\left\langle#1,#2\right\rangle_F}
\let\tr\relax
\newcommand{\tr}{\operatorname{tr}}
\let\rank\relax
\newcommand{\rank}{\operatorname{rank}}
\let\diag\relax
\newcommand{\diag}{\operatorname{blockdiag}}
\let\range\relax
\newcommand{\range}{\operatorname{range}}
\let\OPT\relax
\newcommand{\OPT}{\operatorname{OPT}}
\let\HSS\relax
\newcommand{\HSS}{\operatorname{HSS}}
\let\diver\relax
\newcommand{\diver}{\operatorname{div}}
\let\Var\relax
\newcommand{\Var}{\operatorname{Var}}
\let\UC\relax
\newcommand{\UC}{\mathcal U}
\let\VC\relax
\newcommand{\VC}{\mathcal V}
\let\err\relax
\newcommand{\err}{\mathcal E}
\let\proofstep\relax
\newcommand{\proofstep}[1]{\par\addvspace{.45\baselineskip}\Needspace{3\baselineskip}\textbf{#1}\par\nobreak}
\let\op\relax
\newcommand{\op}[1]{\left\lVert#1\right\rVert_{\mathrm{op}}}
\let\Sn\relax
\newcommand{\Sn}[2]{\left\lVert#1\right\rVert_{S_{#2}}}
\let\Cov\relax
\newcommand{\Cov}{\mathsf C}
\let\Loss\relax
\newcommand{\Loss}{\mathsf L}
\let\Fil\relax
\newcommand{\Fil}{\mathcal F}
\let\Id\relax
\newcommand{\Id}{I}
\let\supp\relax
\newcommand{\supp}{\operatorname{supp}}
\let\vecop\relax
\newcommand{\vecop}{\operatorname{vec}}
\let\Gammafun\relax
\newcommand{\Gammafun}{\operatorname{Gamma}}
\chapter{A deterministic tensor-transport gain}\label{ch:frame_transport}
This lemma bounds the coherent input sequence, not only an isolated impulse. No HSS-frame independence is used in the statement.

Let both sequences of block matrices \(Q_{U,r}=[A_{U,r},B_{U,r};C_{U,r},D_{U,r}]\) and \(Q_{V,r}\) be orthogonal, with \(k\)-by-\(k\) blocks, for \(0\le r<h\). Write \(T_U(s,r)=A_{U,s}\cdots A_{U,r-1}\) and similarly for \(V\). For any fixed sequence of matrices \(u_r\in\mathbb R^{k\times q}\) let

\[
m_0=0,\qquad m_{r+1}=A_{V,r}^{\mathsf T}m_r-C_{V,r}^{\mathsf T}u_r.
\tag{\thechapter.1}\label{frame_transport:manual:1}
\]

For arbitrary orthogonal projectors \(\Pi_s\) on \(\mathbb R^q\), put

\[
\mathscr A(m)=\frac1d\sum_{s<r<h}
\|C_{U,s}T_U(s+1,r)\|_F^2
\|B_{V,r}^{\mathsf T}T_V(s,r)^{\mathsf T}m_s\Pi_s\|_F^2.
\tag{\thechapter.2}\label{frame_transport:manual:2}
\]

\textbf{Lemma.} Uniformly in the horizon, ranks, orthogonal completions, inputs, and projectors,

\[
\boxed{\mathscr A(m)\le16\frac{k}{d}\sum_{r<h}\|u_r\|_F^2.}
\tag{\thechapter.3}\label{frame_transport:manual:3}
\]

\section{Proof}\label{frame_transport:proof}

Drop \(\Pi_s\) in \eqref{frame_transport:manual:2}. Introduce auxiliary independent centered vectors \(\xi_r\in\mathbb R^k\) with covariance \(I_k\), independent of all the fixed data. Let \(p_0=0\) and

\[
p_{r+1}=A_{U,r}^{\mathsf T}p_r+C_{U,r}^{\mathsf T}\xi_r.
\tag{\thechapter.4}\label{frame_transport:manual:4}
\]

Its covariance is bounded by \(I_k\), because \(A_{U,r}^{\mathsf T}A_{U,r}+C_{U,r}^{\mathsf T}C_{U,r}=I_k\). In particular \(\mathbb E_\xi\|p_r\|^2\le k\).

Work in the Hilbert tensor product \(\mathbb R^k_U\otimes\mathbb R^k_V\otimes\mathbb R^q\). All operators below act as the identity on the final factor. Define

\[
K_r=\sum_{s<r}
\left[T_U(s+1,r)^{\mathsf T}C_{U,s}^{\mathsf T}\xi_s\right]
\otimes\left[T_V(s,r)^{\mathsf T}m_s\right].
\tag{\thechapter.5}\label{frame_transport:manual:5}
\]

Independence and zero means of the auxiliary vectors give exactly

\[
d\,\mathscr A(m)\le
\mathbb E_\xi\sum_{r<h}\|(I\otimes B_{V,r}^{\mathsf T})K_r\|^2.
\tag{\thechapter.6}\label{frame_transport:manual:6}
\]

The tensor recurrence is

\[
K_{r+1}=(A_{U,r}^{\mathsf T}\otimes A_{V,r}^{\mathsf T})K_r
+(C_{U,r}^{\mathsf T}\xi_r)\otimes(A_{V,r}^{\mathsf T}m_r).
\tag{\thechapter.7}\label{frame_transport:manual:7}
\]

Set \(Z_r=p_r\otimes m_r\) and \(J_r=K_r-Z_r\). Expanding \eqref{frame_transport:manual:1} and \eqref{frame_transport:manual:4} gives

\[
J_{r+1}=(A_{U,r}^{\mathsf T}\otimes A_{V,r}^{\mathsf T})J_r
+p_{r+1}\otimes(C_{V,r}^{\mathsf T}u_r),\qquad J_0=0.
\tag{\thechapter.8}\label{frame_transport:manual:8}
\]

First control \(Z\). Put \(e_r=B_{V,r}^{\mathsf T}m_r-D_{V,r}^{\mathsf T}u_r\). Orthogonality of \(Q_{V,r}\) gives

\[
\|m_{r+1}\|_F^2+\|e_r\|_F^2
=\|m_r\|_F^2+\|u_r\|_F^2.
\]

Hence \(\sum_r\|B_{V,r}^{\mathsf T}m_r\|_F^2\le4\sum_r\|u_r\|_F^2\), and therefore

\[
\mathbb E_\xi\sum_r\|(I\otimes B_{V,r}^{\mathsf T})Z_r\|^2
\le4k\sum_r\|u_r\|_F^2.
\tag{\thechapter.9}\label{frame_transport:manual:9}
\]

Next control \(J\). Let

\[
\bar J_r=(A_{U,r}^{\mathsf T}\otimes I)J_r,\quad
L_r=(B_{U,r}^{\mathsf T}\otimes I)J_r,\quad
v_r=p_{r+1}\otimes u_r,
\] \[
y_r=(I\otimes B_{V,r}^{\mathsf T})\bar J_r
+(I\otimes D_{V,r}^{\mathsf T})v_r.
\]

The \(U\) row identity and the \(V\) orthogonal splitting yield the exact energy relation

\[
\|J_{r+1}\|^2+\|y_r\|^2+\|L_r\|^2
=\|J_r\|^2+\|v_r\|^2.
\tag{\thechapter.10}\label{frame_transport:manual:10}
\]

Moreover,

\[
\begin{split}
\|(I\otimes B_{V,r}^{\mathsf T})J_r\|^2
={}&\|(A_{U,r}^{\mathsf T}\otimes B_{V,r}^{\mathsf T})J_r\|^2
+\|(B_{U,r}^{\mathsf T}\otimes B_{V,r}^{\mathsf T})J_r\|^2\\
\le{}&2\|y_r\|^2+2\|v_r\|^2+\|L_r\|^2.
\end{split}
\]

Summing, using \eqref{frame_transport:manual:10}, then taking auxiliary expectation gives

\[
\mathbb E_\xi\sum_r\|(I\otimes B_{V,r}^{\mathsf T})J_r\|^2
\le4\mathbb E_\xi\sum_r\|v_r\|^2
\le4k\sum_r\|u_r\|_F^2.
\tag{\thechapter.11}\label{frame_transport:manual:11}
\]

Since \(K=Z+J\), the \(L^2\) triangle inequality, \eqref{frame_transport:manual:9}, and \eqref{frame_transport:manual:11} bound the right side of \eqref{frame_transport:manual:6} by \(16k\sum_r\|u_r\|_F^2\). This proves \eqref{frame_transport:manual:3}.

The same statement holds pathwise for random original data: condition on that data first and use fresh auxiliary vectors. The auxiliary randomness is a proof device and changes no HSS sketch law.

\section{Application to the HSS mean forcing}\label{frame_transport:application-to-the-hss-mean-forcing}

In equation~\eqref{feedback:manual:14}, use \(u_i=U_i^{\mathsf T}h_i\) and \(\Pi_i=\Pi_{H,i}\). Since \(h_i h_i^{\mathsf T}=I_k\), \(\|u_i\|_F=\|U_i\|_F\). The pathwise input bound in Chapter~\ref{ch:feedback} and \eqref{frame_transport:manual:3} give

\[
\boxed{\mathcal D_{\rm drift}\le16\rho^3\beta,\qquad \rho=k/d.}
\tag{\thechapter.12}\label{frame_transport:manual:12}
\]

Together with equation (17) there, this yields

\[
\mathcal P^+\le
\frac{\left[\frac4{1-\rho}+\frac34+17+100\vartheta+64\rho\right]\rho^2}
{1-2\rho}\,\beta,\qquad 0<\rho<\tfrac12.
\tag{\thechapter.13}\label{frame_transport:manual:13}
\]

At \(\rho=1/4\) and \(\vartheta=3/5\), the coefficient is \(12.385416667\). This establishes a uniform predictable-variance bound for the fixed-basis chain. Translating that bound into the final fixed-basis refit theorem still uses the preceding product and source-injection reductions; no basis-learning guarantee is asserted by this lemma alone.

\endgroup

% ===== INTERFACES =====
\begingroup
\let\R\relax
\newcommand{\R}{\mathbb R}
\let\E\relax
\newcommand{\E}{\mathbb E}
\let\Prb\relax
\newcommand{\Prb}{\mathbb P}
\let\T\relax
\newcommand{\T}{\mathsf T}
\let\F\relax
\newcommand{\F}[1]{\left\lVert#1\right\rVert_{F}}
\let\N\relax
\newcommand{\N}[1]{\left\lVert#1\right\rVert}
\let\ip\relax
\newcommand{\ip}[2]{\left\langle#1,#2\right\rangle_F}
\let\tr\relax
\newcommand{\tr}{\operatorname{tr}}
\let\rank\relax
\newcommand{\rank}{\operatorname{rank}}
\let\diag\relax
\newcommand{\diag}{\operatorname{blockdiag}}
\let\range\relax
\newcommand{\range}{\operatorname{range}}
\let\OPT\relax
\newcommand{\OPT}{\operatorname{OPT}}
\let\HSS\relax
\newcommand{\HSS}{\operatorname{HSS}}
\let\diver\relax
\newcommand{\diver}{\operatorname{div}}
\let\Var\relax
\newcommand{\Var}{\operatorname{Var}}
\let\UC\relax
\newcommand{\UC}{\mathcal U}
\let\VC\relax
\newcommand{\VC}{\mathcal V}
\let\err\relax
\newcommand{\err}{\mathcal E}
\let\proofstep\relax
\newcommand{\proofstep}[1]{\par\addvspace{.45\baselineskip}\Needspace{3\baselineskip}\textbf{#1}\par\nobreak}
\let\op\relax
\newcommand{\op}[1]{\left\lVert#1\right\rVert_{\mathrm{op}}}
\let\Sn\relax
\newcommand{\Sn}[2]{\left\lVert#1\right\rVert_{S_{#2}}}
\let\Cov\relax
\newcommand{\Cov}{\mathsf C}
\let\Loss\relax
\newcommand{\Loss}{\mathsf L}
\let\Fil\relax
\newcommand{\Fil}{\mathcal F}
\let\Id\relax
\newcommand{\Id}{I}
\let\supp\relax
\newcommand{\supp}{\operatorname{supp}}
\let\vecop\relax
\newcommand{\vecop}{\operatorname{vec}}
\let\Gammafun\relax
\newcommand{\Gammafun}{\operatorname{Gamma}}
\chapter{Source-product normalization and residual assembly}\label{ch:interfaces}
This chapter proves the interfaces that connect source-product estimates, residual centering, and the uniform fixed-space refit theorem.

\textbf{Conclusion.} The substitution of the new predictable bound into the source-product interface is consistent. No further logarithmic or depth-dependent allowance is present in the audited source or residual assembly. The result applies to the original fixed-basis decoder, for either legal doubly discarded readout. It remains conditional on bases fixed independently of the refit probes; a basis-learning theorem is a separate premise.

\section{Exact normalization of the predictable term}\label{interfaces:exact-normalization-of-the-predictable-term}

In Section~\ref{base:con:product}, the complete output Hilbert space contains every coefficient already emitted and the root coefficient. It is not just the current observation state. The canonical output map is \(\mathcal V\), its Doob martingale is \(\mathcal J_s\), and \(D_s=\mathcal J_{s+1}-\mathcal J_s\). The source-product term is

\[
\mathcal B=\sum_s D_sx_s,\qquad
\mathcal P=\mathbb E\|\mathcal B\|^2.
\]

The deterministic impulse partial trace is explicitly defined by

\[
(\Lambda_{G,s})_{ab}
=\sum_{i=1}^k
\langle\mathcal I_s(E_{ai},0),\mathcal I_s(E_{bi},0)\rangle.
\]

Thus the conditional Haar covariance gives exactly

\[
\mathcal P=\frac1d\sum_s\mathbb E\left[
\operatorname{tr}(\Lambda_{G,s}J_{Y,s}(x_s))
+\operatorname{tr}(\Lambda_{H,s}J_{Z,s}(x_s))\right].
\]

This is the normalization of Chapters~\ref{ch:delayed} and~\ref{ch:feedback}. In particular, the \(1/d\) is outside the partial trace; the latter traces over \(k\) fresh-frame rows. At the final root, a fresh register appears with coefficient \(1/2\), so its partial trace is \(kI/4\), as in Chapter~\ref{ch:delayed}.

The enlarged filter adds the missing doubly discarded sector as an additional orthogonal analytical output. It does not modify the decoder. Its partial traces satisfy \(\Lambda_G\preceq\Lambda_G^+\) and \(\Lambda_H\preceq\Lambda_H^+\), hence

\[
\mathcal P\le\mathcal P^+.
\]

There is no extra factor two in this inequality or in substituting the new bound into the product theorem. The factor two used in the feedback proof has a different role: the squared norm of the enlarged unmarked output is at most the sum of the squared norms of the two legal outputs. Therefore its unmarked trace budget is at most \(2\vartheta k\). This does not redefine \(\mathcal P\).

The direct, delayed, and root decomposition of Chapter~\ref{ch:delayed} is precisely a decomposition of this positive partial-trace sum. The proof pays its enlarged envelope even though the older proposition statement only mentions the legal \(\mathcal P\). Thus using the paid estimate for \(\mathcal P^+\) is licensed by that proof.

\section{Terminal source weights are retained}\label{interfaces:terminal-source-weights-are-retained}

The actual source has increments

\[
\delta x_s=Xt_s^{\mathsf T}b_s^{\mathsf T},\qquad
\beta=\sum_s\|b_s\|^2,
\]

where \(X\) is the initial retained frame and the rows \(b_s\) are deterministic after the initial conditioning. The quantity being bounded is the full terminal product \(\mathcal Vx_m\). The exact product rule of Section~\ref{base:con:product} is

\[
\mathcal Vx_m=\mathcal A+\mathcal B+\mathcal C+\mathcal D.
\]

It does not replace terminal weights on earlier coefficients by the weights current when those coefficients were emitted. The only use of the embedding constant \(b\) in the final chain estimate is the old bound

\[
\mathbb E\|\mathcal B\|^2\le\rho\vartheta b\,\beta.
\]

The other inputs to that estimate remain

\[
\mathbb E\|\mathcal A\|^2\le\rho\vartheta\beta,\quad
\mathbb E\|\mathcal C\|^2\le\gamma_{k,d}\beta,\quad
\mathbb E\|\mathcal D\|^2\le\rho^2\beta,
\]

with \(\mathcal A\) and \(\mathcal C\) orthogonal by the stated fresh-pair parity and martingale argument. Consequently the new bound \(\mathcal P^+\le C_P(\rho)\rho^2\beta\) replaces the chain constant by exactly

\[
\left[\rho\sqrt{C_P(\rho)}
+\rho+\sqrt{\rho\vartheta+\gamma_{k,d}}\right]^2.
\]

No embedding theorem or condition \(b\ge4\) is needed after this replacement. That condition belonged to the old choice of a sufficient bound on \(\mathcal P\), not to the product identity, the other three terms, or the decoder.

The final root sibling is also covered. Its Doob increment is included in the product decomposition, and its predictable contribution uses the source value before that final increment. This matches the root term of the partial-trace formula.

The canonical forward-marked source starts in the opposite response stream, with zero retained regressions. The local feedback proof explicitly permits that orientation; its optional forward free term is only a conservative allowance.

\section{Reattaching a physical residual source}\label{interfaces:reattaching-a-physical-residual-source}

For each discarded physical row source \(e\), the corrected spatial decomposition is

\[
e=\text{exterior retained parts}
+\text{exterior discarded parts}
+\text{strict interior discarded parts}.
\]

The squared budgets satisfy \(\beta+\gamma_{\rm out}+\gamma_{\rm in}=\|e\|^2\) by an orthogonal spatial decomposition. The strict interior term is present in the preliminary propagation analysis and in the near-optimal assembly; the new theorem does not revert to the older exterior-only decomposition.

Conditional on the entire forward ancestor chain, the off-chain coefficient \(\zeta\) is centered, and

\[
\mathbb E[\zeta^{\mathsf T}\zeta\mid\mathcal C]
\preceq\left(\frac{\gamma_{\rm out}}{q-3k}
+\frac{\gamma_{\rm in}}{q-2k}\right)I.
\]

These two denominators are distinct. The stated bound on the opposite-stream mean response is an unconditional mean bound \(\mathbb EW\preceq\vartheta I\) for a canonical zero-regression start. It is used correctly: first condition on the ancestor chain to integrate \(\zeta\), then average \(\operatorname{tr}W\). No bound on \(W\) conditional on future forward frames is substituted. The mixed chain/off-chain terms vanish by that conditional centering.

The resulting source allowance is

\[
C_{\rm ch}\beta+
\frac{\vartheta k}{q-3k}\gamma_{\rm out}
+\frac{\vartheta k}{d}\gamma_{\rm in}.
\]

Since \(q-3k=d-k\), its common coefficient is \(C_*=\max\{C_{\rm ch},\vartheta\rho/(1-\rho)\}\), with no hidden rank factor.

\section{Why summing physical sources introduces no depth factor}\label{interfaces:why-summing-physical-sources-introduces-no-depth-factor}

Source orthogonality follows from the following conditional argument. With all of \(G\) fixed, a source correction is linear in its initial discarded \(H\) frame when coarser \(H\) frames are fixed. Its subsequent operations occur only at ancestors and do not use that initial frame a second time. For two comparable sources, integrate the finer initial frame first. The coarser correction does not depend on it. For incomparable nodes or rows in a paired discarded draw, the appropriate independent row-sign symmetry removes the cross term.

This argument applies to the complete coefficient stack, including earlier emitted coefficients. It also makes every correction mean-zero conditionally on \(G\), so it is orthogonal in expectation to the same one-sided free output, which depends only on \(G\). Individual rows of a matrix source can be separated by the row-sign symmetry of the initial frame. Their squared budgets add, instead of incurring a triangle-inequality factor \(k\).

The discarded physical row bases together with the root basis form an orthonormal multiresolution decomposition. Hence the source row budgets sum to at most \(\|E\|_F^2\) on each orientation. Neither a sum of pointwise-in-time response bounds nor an absolute sum along a path is used here.

This proves the source-orthogonality interface using its actual conditioning order. It is not inferred from fresh-level independence.

\section{Residual centering and the actual decoder}\label{interfaces:residual-centering-and-the-actual-decoder}

The centered-regression bound in Chapter~\ref{ch:centered} is obtained from the original whole-probe law, not a Gaussian replacement or a favorable event. Its fourth-moment induction uses deterministic Parseval budgets, so its constant has no horizon dependence. The regression lemma only needs a deterministic isometry \(V\) with \(V^{\mathsf T}GG^{\mathsf T}V=I\) pathwise and a residual row \(F\) satisfying \(FV=0\); it does not need independence of those probe quantities.

For \(E=A-P_{\mathcal S}A\), the stated orthogonal coefficient-sector decomposition gives exactly the required local zero conditions. Summing the centered regressions over orthogonal discarded row sectors and the root yields

\[
\mathbb E\|\mathcal F_Y\|^2\le\min\{1,4\rho\}\|E\|_F^2,\qquad
\mathbb E\|\mathcal F_Z\|^2\le\min\{1,2\rho\}\|E\|_F^2.
\]

The root has its own factor \(1/2\) in each orientation and therefore a squared factor \(1/4\); the root bound in Chapter~\ref{ch:centered} already includes it. Reversing the legal readout exchanges the two free constants.

The decoder remains linear in the observations, exact on \(\mathcal S\), and valued in \(\mathcal S\). Thus analyzing \(E\) is only a proof reduction. It does not require querying \(P_{\mathcal S}A\) or accessing dense residual entries. The two complete one-sided outputs are combined by an \(L^2\) triangle inequality, not by an unsupported independence or orthogonality assertion. This gives the stated residual-centered certificate with \(C_*\).

\section{Scope of the audit}\label{interfaces:scope-of-the-audit}

No hidden logarithmic or depth-dependent assumption was found in the interfaces audited above. The old logarithm occurs only in obtaining a generic embedding allowance for \(\mathcal P\), which the new specialized bound replaces.

The hypotheses retained are substantive: the bases and their completions are fixed before the refit probes; the two streams are independent and follow the sibling-orthogonal law; the source row budgets are deterministic after the stated conditioning; and the decoder uses one legal discarded readout. The corollaries' widths satisfy the stronger required sampling condition \(q\ge4k\).

Conditioning on an independently learned space preserves all these hypotheses. The construction of that space with the desired best-HSS error and query bounds is not supplied by this audit or by the fixed-basis closure.

\section{\texorpdfstring{7. Endpoint audit: the original width \(q=4k\)}{7. Endpoint audit: the original width q=4k}}\label{interfaces:endpoint-audit-the-original-width-q4k}

The restriction \(\rho<1/2\) in the first small-gain certificate came only from its symmetric use of Young's inequality. It is not an assumption of the inherited source or residual interfaces.

Before splitting the old and future source marks, the forward delayed-feedback estimates are

\[
B_{G,\mathrm{old}}\le \rho\mathcal P_H^+,\qquad
B_{G,\mathrm{future}}\le 50\vartheta\rho^2\beta .
\]

Using \(\|u+v\|^2\le(3/2)\|u\|^2+3\|v\|^2\) therefore gives

\[
R_G\le (1+150\vartheta)\rho^2\beta
       +\frac32\rho\mathcal P_H^+.
\]

The optional free term \(\rho^2\beta\) is conservative for the actual opposite-stream initialization. On the transpose side, the coherent-drift closure and the already proved product decomposition give

\[
\sqrt{R_H}\le\sqrt{\rho\mathcal P_G^+}+4\rho\sqrt\beta,
\qquad
R_H\le\frac32\rho\mathcal P_G^++48\rho^2\beta .
\]

Combining these with the paid direct/root allowance yields

\[
\mathcal P^+\le\widetilde C_P(\rho)\rho^2\beta,\qquad
\widetilde C_P(\rho)=
\frac{4/(1-\rho)+3/4+49+150\vartheta}{1-3\rho/2},
\qquad 0<\rho\le\frac12 .
\]

At \(\rho=1/2\), one has \(\vartheta=5/3\), numerator \(1231/4\), denominator \(1/4\), and hence \(\widetilde C_P=1231\) exactly.

The endpoint is allowed by every other interface used above. In particular, \(q=4k\) gives \(d=2k\), which satisfies the fourth-moment estimates including their equality case; the quotient-loss denominator \(1-\rho\) and the unmarked-response denominator \(1-\rho^2\) remain positive; and \(q-3k=k>0\) for the exterior off-chain source. The bilinear estimate needs only \(d\ge1\), and the root increment remains included. Thus no hidden strict inequality excludes this width.

The source-product constant becomes

\[
C_{\rm ch}=
\left[\rho\sqrt{\widetilde C_P}
      +\rho+\sqrt{\rho\vartheta+\gamma}\right]^2.
\]

At the endpoint, \(\sqrt{1231}<36\), \(\gamma\le2\rho^2=1/2\), and therefore \(\sqrt{\rho\vartheta+\gamma}<3/2\). These elementary bounds give \(C_{\rm ch}<400\), whereas \(C_{\rm off}=\vartheta\rho/(1-\rho)=5/3\). The residual certificate \(1+12\rho+4C_*\) is consequently strictly less than \(1607\), using \(2q=8k\) queries.

This verifies a depth-independent constant-factor refit at the original width. It does not add a basis-learning guarantee.

\endgroup

% ===== UNIFORM =====
\begingroup
\let\E\relax
\newcommand{\E}{\mathbb E}
\let\F\relax
\newcommand{\F}[1]{\lVert#1\rVert_F}
\chapter{Uniform and near-optimal fixed-space refitting}\label{ch:uniform}
\begin{scopebox}
\textbf{Update, 20 September 2026.} This fixed-space refit remains an established component of the later route. It applies after conditioning on training bases only when the refit probes are new and independent, with the stated physical-leaf convention. For the synchronized learner the substitution is retained width $s=4k$ and physical leaves at most $8k=2s$, so the refit costs $32k$ additional products. It supplies neither the training-space error bound nor an efficient final rank-$k$ reduction.
\end{scopebox}

\paragraph{Scope and dependencies.}
Fix the nested rank-$k$ bases before drawing the two refit probes,
and let $\mathcal S$ be their telescoping coefficient space.
The decoder is the original linear, exact-on-$\mathcal S$,
sibling-orthogonal refit; either doubly discarded readout is allowed.
The new uniform predictable-variance estimate is proved by
Chapter~\ref{ch:feedback} and
Chapter~\ref{ch:drift}.
The residual-centered free-output, source decomposition, off-chain
covariance, source orthogonality, and source-product interfaces are
proved in Chapters~\ref{ch:centered} and~\ref{ch:interfaces}, using the propagation identities in Part~I. The calculation below substitutes these estimates and evaluates the constants.

Put $d=q-2k$, $\rho=k/d$, and
\[
 \vartheta=\frac{\rho(2+\rho)}{1-\rho^2},\qquad
 \gamma_{k,d}=\frac{k(k+1-2/d)}{(d-1)(d+2)}.
\]
The new predictable bound is
\[
 \mathcal P^+\le C_P(\rho)\rho^2\beta,\qquad
 C_P(\rho)=
 \frac{4/(1-\rho)+3/4+33+100\vartheta}{1-2\rho},
 \qquad 0<\rho<\tfrac12.
\]
Since the enlarged filter contains either legal readout, the same bound
applies to its predictable interaction. The source-product decomposition of Section~\ref{base:con:product} therefore has constants
\begin{align*}
 C_{\rm ch}
 &=\left[\rho\sqrt{C_P(\rho)}+\rho+
                  \sqrt{\rho\vartheta+\gamma_{k,d}}\right]^2,\\
 C_{\rm off}&=\frac{\vartheta\rho}{1-\rho},\qquad
 C_*=\max\{C_{\rm ch},C_{\rm off}\}.
\end{align*}
Here the predictable term has simply replaced
$\sqrt{\rho\vartheta b}$ in the earlier logarithmic certificate by
$\rho\sqrt{C_P(\rho)}$.

\begin{theorem}[The original eight-$k$ query budget]
For every fixed nested basis family and every deterministic input $A$,
the original refit with $q=4k$ satisfies
\[
 \boxed{\E\F{A-\widehat B}^2\le1607\,d_{\mathcal S}(A)^2,
 \qquad\text{queries}=8k.}
\]
\end{theorem}
\begin{proof}
The Young-weight refinement in
Chapter~\ref{ch:drift} gives, for $0<\rho\le1/2$,
\[
 \mathcal P^+\le\widetilde C_P(\rho)\rho^2\beta,\qquad
 \widetilde C_P(\rho)=
 \frac{4/(1-\rho)+3/4+49+150\vartheta}{1-3\rho/2}.
\]
At $q=4k$, $\rho=1/2$, $\vartheta=5/3$, and
$\widetilde C_P=1231$. Replacing $C_P$ by $\widetilde C_P$ in
the chain formula gives
\[
 C_{\rm ch}\le
 [\tfrac12\sqrt{1231}+\tfrac12+\sqrt{4/3}]^2<400,
 \qquad C_{\rm off}=5/3.
\]
The estimate $\gamma_{k,d}\le2\rho^2$ and the residual assembly
proved below are valid also at $\rho=1/2$. Their bound
$1+12\rho+4C_*$ is at most $7+1600=1607$.
This constant is not optimized.
\end{proof}

\begin{theorem}[Explicit uniform fixed-basis certificate]
For every $k\ge1$, every fixed nested basis family, and every deterministic
input matrix $A$, if $q\ge18k$, then
\[
 \boxed{\E\F{A-\widehat B}^2
 \le[1+12\rho+512\rho^2]\,d_{\mathcal S}(A)^2.}
\]
The decoder uses $2q$ matrix--vector queries and no new sketches at
individual levels.
\end{theorem}
\begin{proof}
For $0<\rho\le1/16$,
\[
 \frac{\vartheta}{\rho}\le\frac{176}{85}<\frac94,\qquad
 C_P(\rho)\le C_P(1/16)=\frac{103954}{1785}<64.
\]
The latter follows from monotonicity of every positive term in its
numerator and of $(1-2\rho)^{-1}$.
Also $\gamma_{k,d}\le2\rho^2$: for $d\ge2$,
\[
 \frac{\gamma_{k,d}}{\rho^2}
 =\frac{d^2}{(d-1)(d+2)}
       \left(1+\frac{1-2/d}{k}\right)
 \le\frac{2d(d-1)}{(d-1)(d+2)}<2.
\]
Consequently
\[
 \frac{C_{\rm ch}}{\rho^2}
 <\left(9+\frac{\sqrt{17}}2\right)^2
 <\left(\frac{45}4\right)^2<128,\qquad
 \frac{C_{\rm off}}{\rho^2}
 \le\frac{2816}{1275}<128.
\]

For completeness, the residual-centered assembly is as follows.
Let $E=A-P_{\mathcal S}A$. The two one-sided free coefficient stacks
have expected energies at most $f_Y\F E^2$ and $f_Z\F E^2$, where
$f_Y=\min\{1,4\rho\}$ and $f_Z=\min\{1,2\rho\}$.
For each source row the chain, exterior, and interior budgets
$\beta,\gamma_{\rm out},\gamma_{\rm in}$ sum to its squared norm.
Its correction energy is bounded by
\[
 C_{\rm ch}\beta+
 \frac{\vartheta k}{q-3k}\gamma_{\rm out}+
 \frac{\vartheta k}{d}\gamma_{\rm in}
 \le C_*(\beta+\gamma_{\rm out}+\gamma_{\rm in}).
\]
Here $k/(q-3k)=\rho/(1-\rho)$.
The inherited opposite-stream orthogonality adds distinct source
energies and their same-sided free output; their row budgets total
at most $\F E^2$ on each side.
The two complete one-sided outputs need not be independent or orthogonal.
Their $L^2$ triangle inequality, exactness on $\mathcal S$, and
$E\perp\mathcal S$ give
\[
 \E\F{A-\widehat B}^2
 \le\left[1+
   \left(\sqrt{f_Y+C_*}+\sqrt{f_Z+C_*}\right)^2\right]\F E^2
 \le[1+12\rho+4C_*]\F E^2.
\]
Use $C_*\le128\rho^2$.
\end{proof}

\begin{corollary}[A concrete linear-query refit]
At $q=18k$, the original refit uses $36k$ queries and satisfies
\[
 \boxed{\E\F{A-\widehat B}^2\le\frac{15}{4}d_{\mathcal S}(A)^2.}
\]
\end{corollary}
\begin{proof}
Here $\rho=1/16$, so
$1+12\rho+512\rho^2=1+3/4+2=15/4$.
\end{proof}

\begin{corollary}[Relative error with no logarithmic width factor]
For $0<\varepsilon\le1$, choose
\[
 \boxed{q=\left\lceil k\left(2+\frac{32}{\varepsilon}\right)\right\rceil.}
\]
Then the unchanged refit uses $2q=O(k/\varepsilon)$ queries and satisfies
\[
 \boxed{\E\F{A-\widehat B}^2\le(1+\varepsilon)d_{\mathcal S}(A)^2.}
\]
\end{corollary}
\begin{proof}
The ceiling only increases $d$, so $\rho\le\varepsilon/32\le1/32$.
Therefore
\[
 12\rho+512\rho^2
 \le\frac{3\varepsilon}{8}+\frac{\varepsilon^2}{2}
 \le\frac{7\varepsilon}{8}<\varepsilon.
\]
\end{proof}

\paragraph{Remaining basis-learning premise.}
These are fixed-space guarantees. They also hold conditionally on a
random nested space if its learning probes are independent of the refit
probes. An $O(k)$-query learner satisfying
\[
 \E_{\rm learn}d_{\mathcal S}(A)^2
 \le C_bL\min_{B\in{\rm HSS}(L,k)}\F{A-B}^2
\]
would combine with the $36k$-query corollary by conditional expectation.
No such learner is proved in this chapter, so the full best-HSS benchmark
retains that separate unresolved premise.

\endgroup

\part{Selection reductions and limitations of basis-learning arguments}

% ===== BICRITERIA =====
\begingroup
\let\R\relax
\newcommand{\R}{\mathbb R}
\let\E\relax
\newcommand{\E}{\mathbb E}
\let\Prb\relax
\newcommand{\Prb}{\mathbb P}
\let\T\relax
\newcommand{\T}{\mathsf T}
\let\F\relax
\newcommand{\F}[1]{\left\lVert#1\right\rVert_{F}}
\let\N\relax
\newcommand{\N}[1]{\left\lVert#1\right\rVert}
\let\ip\relax
\newcommand{\ip}[2]{\left\langle#1,#2\right\rangle_F}
\let\tr\relax
\newcommand{\tr}{\operatorname{tr}}
\let\rank\relax
\newcommand{\rank}{\operatorname{rank}}
\let\diag\relax
\newcommand{\diag}{\operatorname{blockdiag}}
\let\range\relax
\newcommand{\range}{\operatorname{range}}
\let\OPT\relax
\newcommand{\OPT}{\operatorname{OPT}}
\let\HSS\relax
\newcommand{\HSS}{\operatorname{HSS}}
\let\diver\relax
\newcommand{\diver}{\operatorname{div}}
\let\Var\relax
\newcommand{\Var}{\operatorname{Var}}
\let\UC\relax
\newcommand{\UC}{\mathcal U}
\let\VC\relax
\newcommand{\VC}{\mathcal V}
\let\err\relax
\newcommand{\err}{\mathcal E}
\let\proofstep\relax
\newcommand{\proofstep}[1]{\par\addvspace{.45\baselineskip}\Needspace{3\baselineskip}\textbf{#1}\par\nobreak}
\let\op\relax
\newcommand{\op}[1]{\left\lVert#1\right\rVert_{\mathrm{op}}}
\let\Sn\relax
\newcommand{\Sn}[2]{\left\lVert#1\right\rVert_{S_{#2}}}
\let\Cov\relax
\newcommand{\Cov}{\mathsf C}
\let\Loss\relax
\newcommand{\Loss}{\mathsf L}
\let\Fil\relax
\newcommand{\Fil}{\mathcal F}
\let\Id\relax
\newcommand{\Id}{I}
\let\supp\relax
\newcommand{\supp}{\operatorname{supp}}
\let\vecop\relax
\newcommand{\vecop}{\operatorname{vec}}
\let\Gammafun\relax
\newcommand{\Gammafun}{\operatorname{Gamma}}
\chapter{Bicriteria reduction and the final-rank interface}\label{ch:bicriteria}
\begin{scopebox}
\textbf{Update, 20 September 2026.} The exact best-HSS projection used here is a query-only mathematical reduction: it is computed from a known intermediate matrix and requires no further products with the unknown input. Existence and its approximation inequality remain valid. The target also requires polynomial additional work; neither this chapter nor the later consolidation provides an efficient same-rank projection with the required error factor. A recompression factor depending on depth cannot be discarded when composing the guarantees. The final reduction must be onto the original target tree, not merely the learner's coarsened physical-leaf tree.
\end{scopebox}

The rank of the final answer need not be imposed on every intermediate randomized object. This chapter records the query-only reduction and the obstruction to one tempting interpolation shortcut.

\section{Offline return to the original HSS rank}
Let $\mathcal H_k$ denote the fixed-tree HSS rank-at-most-$k$ class, and suppose a procedure returns a known matrix $C$ with
\[
 \E\|A-C\|_F^2\le\alpha L\,\operatorname{dist}_F(A,\mathcal H_k)^2.
\]
The known $C$ can have HSS rank $ck$, or have no asserted structure. Let $\widehat B$ be an exact Frobenius-best approximation to $C$ in $\mathcal H_k$. The minimum exists: the class is closed, contains zero, and its minimizing sequences can be confined to a bounded ball. This optimization uses only the known matrix $C$, not additional products with $A$.

For a best comparator $B_*$ for $A$,
\[
 \|A-\widehat B\|_F
 \le\|A-C\|_F+\|C-\widehat B\|_F
 \le2\|A-C\|_F+\|A-B_*\|_F.
\]
Minkowski's inequality in probability gives
\begin{equation}\label{bicriteria:eq:projection}
 \boxed{\E\|A-\widehat B\|_F^2
 \le(2\sqrt{\alpha L}+1)^2\OPT_k(A)
 \le(2\sqrt\alpha+1)^2L\OPT_k(A),\qquad L\ge1.}
\end{equation}
Thus an $O(k)$-query bicriteria learner with expected $O(L)$ error is sufficient in the computation-free query model. Exact nonlinear projection introduces neither additional queries nor another depth factor. No polynomial-time algorithm for this projection is asserted. A generic deterministic recompression with an $O(L)$ approximation factor can instead introduce another factor of depth.

\section{Capturing both nullified samples does not interpolate both observations}
Consider local compressed size $2r$ and width $q=3r$, with full-row-rank designs. Each design nullspace has dimension $r$, so spaces $U,V$ can contain every column of the nullified innovations $E_Y,E_Z$. Write $P=UU^T$, $Q=VV^T$, and let
\[
 B_Y=Y\Omega^\dagger,\qquad
 B_Z=(\Psi^\dagger)^TZ^T,\qquad
 D=(I-P)B_Y+PB_Z(I-Q).
\]
Then
\[
 (I-P)(Y-D\Omega)=(I-P)E_Y=0,
\]
but
\begin{equation}\label{bicriteria:eq:mismatch}
 (I-Q)(Z-D^T\Psi)
 =(I-Q)(B_Z^T-B_Y^T)(I-P)\Psi.
\end{equation}
The discarded/discarded regression mismatch can therefore survive and enter later data. Exact innovation capture is neither simultaneous two-sided interpolation nor a physical Frobenius-error estimate.

There is a matching dimension check. With balanced leaves of size $2r$, the saturated HSS-rank-$r$ stratum has dimension
\[
 5Nr-6r^2.
\]
Leaf diagonal blocks contribute $2Nr$ parameters; sibling couplings contribute $Nr-2r^2$; row and column Grassmann parameters contribute $2Nr-4r^2$. The two-sided observation map at generic full-column-rank width-$q$ probes has rank $2Nq-q^2$. At $q=3r$, its rank is $6Nr-9r^2$, exceeding the HSS dimension by $r(N-3r)>0$ when $N>3r$. Generic two-sided data therefore do not have an exact HSS-rank-$r$ interpolant.

This only rules out the interpolation shortcut. It does not rule out stable bicriteria approximation or the general $O(k)$-query target.

\endgroup

% ===== FINITE =====
\begingroup
\let\E\relax
\newcommand{\E}{\mathbb E}
\let\R\relax
\newcommand{\R}{\mathbb R}
\let\T\relax
\newcommand{\T}{\mathsf T}
\let\F\relax
\newcommand{\F}[1]{\lVert#1\rVert_F}
\let\op\relax
\newcommand{\op}[1]{\lVert#1\rVert_2}
\let\tr\relax
\newcommand{\tr}{\operatorname{tr}}
\let\rank\relax
\newcommand{\rank}{\operatorname{rank}}
\let\orth\relax
\newcommand{\orth}{\operatorname{orth}}
\let\argmin\relax
\newcommand{\argmin}{\operatorname*{arg\,min}}
\chapter{High-probability finite-family selection}\label{ch:finite}
\begin{chaptersynopsis}
The public finite-family construction~\cite{Note} gives a Gaussian two-sided selector from a known constant-factor pivot. Its main argument is valid in its stated exact-real query model. A different Frobenius-product estimate, together with rank-zero-compatible rounding, improves its accuracy dependence: the sufficient query count is
$O(\gamma(1+\gamma)\varepsilon^{-1}\sqrt{\log(M/\delta)}+\log(1/\delta))$, rather than $O(\gamma^4\varepsilon^{-2}\sqrt{\log(M/\delta)}+\log(1/\delta))$.
We give a complete proof with explicit constants. The result concerns selection from known matrices; it does not construct an HSS candidate list, remove dependence in the recycled HSS refit, or supply a uniform expected-error theorem for the continuous HSS class. The restricted-query discussion in the source has additional qualifications. No priority, optimal-constant, or independent-verification claim is made.
\end{chaptersynopsis}

\section{Source, model, and audit verdict}
The external finite-family starting point is \emph{Constant factor to $(1+\epsilon)$-approximation for structured matrix learning}, the public finite-family approximation argument in the Structured-Matrix-Learning repository \cite{Note}. \par
\noindent\textbf{Model.} An unknown $A\in\R^{n\times n}$ is accessible through exact products $Ax,A^\T x$. A finite family $\mathcal F=\{A_1,\ldots,A_M\}$ and a pivot $A_0$ are known as matrices. All non-query computation is free. With
\[
\Delta=\min_i\F{A-A_i},\qquad B=A-A_0,\qquad D_i=A_i-A_0,
\]
a known $\gamma\ge1$ satisfies $\F B\le\gamma\Delta$. Probabilities concern fresh selector randomness; $A$, the family, and the pivot are fixed before that randomness. A random family from an independent preceding transcript is covered by conditioning on that transcript.

The source's Gaussian main theorem follows from its rangefinder event, deterministic head bias, conditional quadratic-form concentration, and deterministic score selection. Its product $D_i^{H\T}R$ need not be symmetric: symmetrization preserves its trace and only decreases both relevant norms. The score randomness is independent of the rangefinder. The two query batches really use $2m$ right products and at most $m$ left products, not $m$ products per candidate. Exact candidate truncations and their products are offline work.

The improvement below keeps that method, but changes the Frobenius bound and uses a floor for the truncation rank. A positive rank must not be forced when the accuracy is small. The source's original ceiling is harmless at its larger width after the constants are chosen in the stated order, but it is not appropriate for the sharper width.

\section{Two concentration inputs with explicit constants}
\begin{lemma}[Gaussian deflation]\label{finite:lem:range}
Let $8\le m<n$, draw a standard Gaussian $\Omega\in\R^{n\times m}$, and let $Q=\orth(B\Omega)$, retaining only its nonzero-range columns. Put
\[
\widehat B=QQ^\T B,\qquad R=(I-QQ^\T)B.
\]
With probability at least $1-6e^{-m/2}$,
\begin{equation}\label{finite:eq:range}
\op R\le C_R\F B/\sqrt m,\qquad C_R=52.
\end{equation}
Always $\F R\le\F B$.
\end{lemma}
\begin{proof}
Use Halko--Martinsson--Tropp, Corollary 10.9 \cite{HMT}, with target rank $r=\lfloor m/2\rfloor$ and oversampling $p=m-r\ge4$. Its failure probability is at most $6e^{-p}\le6e^{-m/2}$. The bound is
\[
\op R\le
(1+17\sqrt{1+r/p})\sigma_{r+1}(B)
+\frac{8\sqrt m}{p+1}\left(\sum_{j>r}\sigma_j(B)^2\right)^{1/2}.
\]
Use $r/p\le1$, $\sigma_{r+1}(B)\le\sqrt{2/m}\F B$, and $p+1\ge m/2$. The resulting coefficient is less than $\sqrt2+34+16<52$. If $B$ has rank smaller than $r$, the Gaussian rangefinder spans its range almost surely, so $R=0$. The Frobenius assertion follows from orthogonal projection.
\end{proof}

\begin{lemma}[Averaged Gaussian quadratic form]\label{finite:lem:quad}
For a fixed real $F\in\R^{n\times n}$ and independent standard Gaussian vectors $y_1,\ldots,y_m$, with probability at least $1-2e^{-t}$,
\begin{equation}\label{finite:eq:quad}
\left|\frac1m\sum_{j=1}^m y_j^\T Fy_j-\tr F\right|
\le 2\F F\sqrt{t/m}+2\op F\,t/m.
\end{equation}
\end{lemma}
\begin{proof}
Replace $F$ by $S=(F+F^\T)/2$. The random quadratic form and trace are unchanged, while both norms decrease. Diagonalizing $S$ and using the Gaussian square moment-generating function gives, for $0<\theta<m/(2\op S)$,
\[
\log\E\exp\!\left[\theta\left(\frac1m\sum y_j^\T Sy_j-\tr S\right)\right]
\le \frac{\theta^2\F S^2/m}{1-2\theta\op S/m}.
\]
The same bound holds for the negative variable. Chernoff optimization yields \eqref{finite:eq:quad}. This is the Gaussian quadratic-form case of the Hanson--Wright method; see \cite{RV} for the general sub-Gaussian theorem.
\end{proof}

\section{The sharper simultaneous score estimator}
\begin{theorem}[Simultaneous inner products]\label{finite:thm:score}
Let $0<\eta\le1$, $0<\delta<1/2$, and $M\ge1$. Define
\begin{equation}\label{finite:eq:m}
\begin{split}
t&=\log(4M/\delta),\\
m&=\max\!\left\{8,\ \left\lceil2\log(12/\delta)\right\rceil,
\ \left\lceil8C_R\eta^{-1}\sqrt t\right\rceil\right\}.
\end{split}
\end{equation}
There is an algorithm using at most $3m$ matrix--vector queries such that, with probability at least $1-\delta$,
\begin{equation}\label{finite:eq:score}
|\widehat\varphi_i-\langle B,D_i\rangle|
\le\eta\F B\F{D_i}\qquad(1\le i\le M).
\end{equation}
If $m\ge n$, $n$ standard-basis right queries instead recover $A$ exactly.
\end{theorem}
\begin{proof}
Consider $m<n$. Draw independent standard Gaussian matrices $\Omega,Y\in\R^{n\times m}$ and query $A\Omega,AY$. Construct $Q=\orth(B\Omega)$. Query $A^\T Q$ to compute $Q^\T B$ and hence $\widehat B=QQ^\T B$ exactly. This uses $2m+\rank Q\le3m$ queries. In particular, $RY=BY-\widehat BY$ is known even though $R$ is not.

Set
\begin{equation}\label{finite:eq:r}
r_0=\min\!\left\{n,\left\lfloor\frac{\eta^2m}{64C_R^2}\right\rfloor\right\},\qquad
D_i^L=[D_i]_{r_0},\quad D_i^H=D_i-D_i^L.
\end{equation}
Here $[D_i]_0=0$; rank zero is allowed. The exact offline SVD is legitimate because every $D_i$ is known. Output
\begin{equation}\label{finite:eq:estimator}
\widehat\varphi_i=\langle D_i,\widehat B\rangle
+\frac1m\langle D_i^HY,RY\rangle.
\end{equation}
Its error is a centered quadratic-sketch error minus $\langle D_i^L,R\rangle$.

On the event of Lemma~\ref{finite:lem:range}, the head bias satisfies
\begin{equation}\label{finite:eq:bias}
|\langle D_i^L,R\rangle|
\le\sqrt{r_0}\F{D_i}\op R
\le\frac\eta8\F{D_i}\F B.
\end{equation}
The floor in \eqref{finite:eq:r} gives the last inequality without an additional condition $m\gtrsim\eta^{-2}$.

Condition on $\Omega$. Then $R$ is fixed and $Y$ remains independent Gaussian. Write $F_i=(D_i^H)^\T R$ and $P_i=\F{D_i}\F B$. The decisive estimate is
\begin{equation}\label{finite:eq:betterF}
\boxed{\F{F_i}\le\F{D_i^H}\op R\le C_RP_i/\sqrt m.}
\end{equation}
This replaces the source's valid but larger estimate $\op{D_i^H}\F R=O(P_i/(\eta\sqrt m))$.

If the tail is nonzero, $r_0<n$ and
\[
r_0+1\ge\frac{\eta^2m}{64C_R^2},\qquad
\op{D_i^H}\le\frac{\F{D_i}}{\sqrt{r_0+1}}.
\]
Therefore
\begin{equation}\label{finite:eq:betterop}
\op{F_i}\le\frac{8C_R^2P_i}{\eta m}.
\end{equation}
If $r_0=n$, the tail vanishes and its error is zero.

Lemma~\ref{finite:lem:quad} now bounds the stochastic error, divided by $P_i$, by
\[
\frac{2C_R\sqrt t}{m}+\frac{16C_R^2t}{\eta m^2}
\le\frac\eta4+\frac\eta4=\frac\eta2
\]
except with probability $2e^{-t}$. Union bounding over the fixed family costs at most $2Me^{-t}=\delta/2$. The deflation event fails with probability at most $\delta/2$ by \eqref{finite:eq:m}. Adding \eqref{finite:eq:bias} gives error at most $5\eta P_i/8<\eta P_i$. If $B=0$ or $D_i=0$, the corresponding error is identically zero. No norm of an unknown residual is needed to choose $m$.
\end{proof}

\clearpage
\section{Selection and the warm start}
\begin{theorem}[Sharper warm-start selector]\label{finite:thm:main}
Under the model of Section 1, for $0<\varepsilon\le1$ and $0<\delta<1/2$, one can return $A_{i_*}\in\mathcal F$ with
\[
\F{A-A_{i_*}}\le(1+\varepsilon)\Delta
\]
with probability at least $1-\delta$, using
\begin{equation}\label{finite:eq:main}
\boxed{O\!\left(\frac{\gamma(1+\gamma)}{\varepsilon}
\sqrt{\log(M/\delta)}+\log(1/\delta)\right)}
\end{equation}
queries. Conditional on the pivot and family, there are two query batches. A high-probability norm guarantee is asserted, not an expected squared-error guarantee.
\end{theorem}
\begin{proof}
Take $\eta=\varepsilon/[2\gamma(1+\gamma)]$ in Theorem~\ref{finite:thm:score} and minimize
$s_i=\F{D_i}^2-2\widehat\varphi_i$.
Let $d_* =\F{A-A_{i_*}}$ and choose $i_o$ with error $\Delta$. The simultaneous score bound gives
\[
d_*^2-\Delta^2
\le2\eta\F B\bigl(\F{D_{i_*}}+\F{D_{i_o}}\bigr)
\le2\eta\F B(d_*+\Delta+2\F B).
\]
If $\Delta>0$, then $d_*\ge\Delta$; divide by $d_*+\Delta$ to obtain
\[
d_*-\Delta\le2\eta\F B\left(1+\frac{\F B}{\Delta}\right)
\le2\eta\gamma(1+\gamma)\Delta=\varepsilon\Delta.
\]
If $\Delta=0$, the pivot equals $A$ and selection is exact. The query count follows from \eqref{finite:eq:m}. A pivot/family fixed by an earlier independent transcript is covered by conditioning. Their generation queries must still be added to the budget.
\end{proof}

\paragraph{Interpretation of the improvement.}
At target inner-product accuracy $\eta P_i$, \eqref{finite:eq:betterF} gives Gaussian variance exponent of order $\eta^2m^2$, rather than $\eta^4m^2$. Equation~\eqref{finite:eq:betterop} gives the same $\eta^2m^2$ scale in the linear branch. This proves the $\varepsilon^{-1}$ improvement suggested in the source's open-question list, in its warm-start exact-real Gaussian model. It does not prove optimality in accuracy or a sign-restricted analogue.

\paragraph{Obtaining a pivot.}
Theorem 1 of \cite{QESML} already returns a $(3+\zeta)$ norm approximation from a known finite family, without a supplied pivot. At fixed $\zeta$, its query cost is $\widetilde O(\sqrt{\log M})$, with displayed log-log and confidence factors in that paper. Use failure probability $\delta/2$ for this stage, freeze its output as a pivot with a known constant $\gamma$, and use fresh Gaussian randomness and failure probability $\delta/2$ in Theorem~\ref{finite:thm:main}. A union bound gives a scale-free high-probability finite-family selector; at fixed confidence its total cost is $\widetilde O(\varepsilon^{-1}\sqrt{\log M})$. The bootstrap has more than two query rounds and retains the first stage's log-log factors. It does not settle the source's sharper finite-precision/no-overhead warm-start question.

\section{The precise HSS interface and its limitations}
The linked theorem selects \emph{known matrices}; an HSS coefficient space or basis family is not itself a known candidate matrix. The local HSS refit instead constructs coefficients in a continuous fixed space. The outstanding project objective is to learn bases and return rank-at-most-$k$ HSS output with expected squared error $O(L)$ times the best HSS squared error in $O(k)$ total queries. The new selector is most naturally an independent validation stage, not a substitute for that training theorem.

A sufficient list-generation route would produce, using $O(k)$ queries, a fixed-after-training list of fully known HSS matrices of cardinality $M$ with $\log M=O(k^2)$, a known constant-factor pivot relative to the best list error, and a best-list approximation guarantee against the HSS optimum. At constant accuracy and confidence, Theorem~\ref{finite:thm:main} then adds $O(k)$ validation queries. All candidate-construction queries must be counted. One still needs failure-event control to convert the high-probability norm conclusion into the project's expected squared-error statement. Using the general pivot theorem of \cite{QESML} gives soft-$O(k)$, not automatically the literal $O(k)$ target.

The fresh validation pair must be independent of the training transcript. Choosing the family after examining those same validation probes is not covered by the union bound. The theorem therefore does not remove the anticipative dependence in the existing recycled-sketch calculation.

\paragraph{Why a direct net is insufficient.}
For a bounded family, a Frobenius net of radius $\tau$ only ensures
\[
\min_{C\in\text{net}}\F{A-C}
\le\inf_{B\in\mathrm{HSS}}\F{A-B}+\tau.
\]
A fixed positive $\tau$ cannot give a pure multiplicative guarantee at zero or arbitrarily small optimum. The corresponding finite-to-infinite passage in \cite{QESML} explicitly has an additive error term.

There is also an entropy cost. With leaf size $2k$ and $2k\mid n$, arbitrary block-diagonal matrices form an HSS subspace of dimension $2nk$. A constant-resolution net of its Frobenius unit ball has $\log M=\Omega(nk)$. The generic selector's $\sqrt{\log M}$ allowance is then of order $\sqrt{nk}$, not $k$. This is \emph{not} an HSS query lower bound: the block-diagonal subclass itself can be recovered with $2k$ right queries by repeating the same identity block over the leaves. It shows only that generic entropy discards useful structure.

\paragraph{Expected error and accuracy convention.}
The source and Theorem~\ref{finite:thm:main} bound the Frobenius norm with high probability. For a high-probability squared factor $1+\epsilon$, substitute $\varepsilon=\sqrt{1+\epsilon}-1$. An expectation bound needs an additional control on failure outputs. Finite cardinality alone does not supply a uniform relative bound on their errors.

\section{Qualifications to the source's supplementary discussion}
The Gaussian main theorem and the preceding sharpening should be separated from the restricted-query appendix. That appendix explicitly leaves the Rademacher rangefinder guarantee as a hypothesis. Its binary simulation of real-valued left queries also carries a scale- and dimension-dependent logarithmic overhead, rather than changing only absolute constants.

One definite correction is its zero-sketch statement after switching to signs. If $B=u(e_1+e_2)^\T\ne0$ and $\Omega$ has $m$ independent Rademacher columns, then
\[
\Pr\{B\Omega=0\}=2^{-m}>0.
\]
This can be included in a failure budget; it is not a probability-zero event. The same event has probability zero for Gaussian probes.

The codeword counting lower bound applies to its explicit finite-information models: sign/zero query entries, or only $O(\log\log M)$ informative response bits per query on that construction. It does not establish a lower bound in every unspecified ``bounded-precision'' model or for arbitrary exact-real queries. No noise-robust multiplicative theorem at zero optimum is established by that discussion. These qualifications do not affect the Gaussian theorem checked above.

\section{Reproducible checks}
The accompanying script checks 112 width/confidence parameter combinations, including very small accuracies, plus 20 additional rank-rounding regimes (five with positive truncation rank). It verifies 60 matrix error decompositions and norm inequalities; the maximum normalized identity discrepancy is below $2.8\times10^{-17}$. Twelve product three-point Gaussian quadratures check quadratic-form means and variances exactly up to floating point, with normalized discrepancies below $8.3\times10^{-16}$. Twenty selection-parameter identities and complete sign-sketch enumerations through width six are also included.

The conservative constant $52$ makes the truncation rank zero in all 112 primary parameter examples; this is reported rather than hidden. Separate rounding and arbitrary-rank algebra tests cover the nonzero-head regime. These checks verify algebra and sufficient-parameter bookkeeping, not the probability theorems. No Monte Carlo accuracy claim, formal verification, or optimal-constant claim is made. These finite-family checks do not validate the separate HSS approximation theorems.

\endgroup

% ===== EXPECTED =====
\begingroup
\let\E\relax
\newcommand{\E}{\mathbb E}
\let\Prb\relax
\newcommand{\Prb}{\mathbb P}
\let\R\relax
\newcommand{\R}{\mathbb R}
\let\F\relax
\newcommand{\F}[1]{\lVert#1\rVert_F}
\let\argmin\relax
\newcommand{\argmin}{\operatorname*{arg\,min}}
\chapter{Expected finite-family selection}\label{ch:expected}
This chapter closes the failure-event expectation gap in the
Chapter~\ref{ch:finite}, conditional on its proved
simultaneous inner-product estimator. It does not construct an HSS
training list or a warm pivot. In particular it does not prove the
general HSS approximation target.

\paragraph{Simultaneous score estimate.}
From Chapter~\ref{ch:finite}, for fixed known matrices $A_0,A_1,\ldots,A_M$, put
\[
 B=A-A_0,\qquad D_i=A_i-A_0,\qquad
 b=\F B,\qquad \Delta=\min_i\F{A-A_i}.
\]
For $0<\eta\le1$ and $0<\delta<1/2$, the audit constructs scores
\[
 s_i=\F{D_i}^2-2\widehat\varphi_i
\]
using
\[
 O\!\left(\eta^{-1}\sqrt{\log(M/\delta)}+
                   \log(1/\delta)\right)
\]
queries. With probability at least $1-\delta$, simultaneously
\begin{equation}\label{expected:eq:score}
 |\widehat\varphi_i-\langle B,D_i\rangle|
       \le\eta b\F{D_i}.
\end{equation}
This estimator needs no warm-pivot hypothesis; that hypothesis is used
only in the audit's subsequent relative-error comparison. Its
Frobenius-product improvement, truncation floor, and width arithmetic
are consistent. The Gaussian rangefinder input is precisely the first
bound in Halko--Martinsson--Tropp, Corollary 10.9; its substitution
$r=\lfloor m/2\rfloor$, $p=m-r$ gives the audit's constant $52$.

\begin{lemma}[An independent residual radius]\label{expected:lem:radius}
Let $G\in\R^{n\times s}$ be standard Gaussian and define
\[
 r^2=s^{-1}\F{BG}^2.
\]
Then $\E r^2=b^2$. If $b>0$, then
\[
 \Prb\{r<b/2\}\le e^{-s/5}.
\]
Thus $s=\lceil5\log(1/\delta)\rceil$ right queries suffice
to obtain this radius with lower-tail failure at most $\delta$.
\end{lemma}
\begin{proof}
Normalize the squared singular values of $B$ as $w_a\ge0$ with
$\sum_a w_a=1$. Then
\[
 sr^2/b^2\ \overset{d}=\ \sum_a w_a Z_a,
       \qquad Z_a\sim\chi_s^2\text{ independently}.
\]
The mean assertion follows. For $t>0$,
\[
 \E e^{-tsr^2/b^2}
 =\prod_a(1+2tw_a)^{-s/2}
 \le(1+2t)^{-s/2},
\]
since $\prod_a(1+2tw_a)\ge1+2t\sum_a w_a=1+2t$.
Use $t=1/2$ and Markov's inequality to obtain
\[
 \Prb\{r^2<b^2/4\}
 \le\exp[-s((\log2)/2-1/8)]\le e^{-s/5}.
\]
For $b=0$, the radius is identically zero and no lower-tail
failure is declared.
\end{proof}

\paragraph{The modified selector.}
Use independent randomness for the radius $r$ and the scores for
the \emph{entire fixed family}. The radius batch also makes
$(A-A_i)G$ available for every known candidate. If any such matrix
is exactly zero, return that candidate immediately. For a fixed finite
family, a nonzero residual has a nonzero Gaussian sketch almost surely.
Thus this check detects $\Delta=0$ exactly, even without a warm pivot,
and has no false positive almost surely in the exact-real model.

If this check does not return, compute the known quantity
\[
 d_0=\min_i\F{D_i},
\]
retain indices in
\begin{equation}\label{expected:eq:filter}
 \mathcal I(r)=\{i:\F{D_i}\le d_0+4r\},
\end{equation}
and return any score minimizer over $\mathcal I(r)$. The set is
nonempty because an index attaining $d_0$ belongs to it. Computing
the scores before restricting the family avoids any appeal to a
union bound over a family chosen from the score randomness.

\begin{theorem}[Expected oracle inequality without a warm hypothesis]
\label{expected:thm:oracle}
Let $0<\eta\le\alpha\le1$ and $0<\delta<1/2$.
The modified selector satisfies
\begin{equation}\label{expected:eq:oracle}
 \boxed{\quad
 \E\F{A-A_{\widehat i}}^2
 \le(1+\alpha+5\delta)\Delta^2
       +(15\eta+92\delta)b^2.
 \quad}
\end{equation}
Its query count is
\[
 O\!\left(\eta^{-1}\sqrt{\log(M/\delta)}+
                   \log(1/\delta)\right).
\]
Whenever the proposed count exceeds $n$, exact recovery using $n$
standard-basis right queries and offline selection is an alternative.
\end{theorem}
\begin{proof}
The zero-residual return has error zero, so it satisfies all bounds
below. For the other branch, proceed as follows.
Let $i_o$ attain $\Delta$, let $\mathcal E$ be the simultaneous score
event \eqref{expected:eq:score}, and let $\mathcal H=\{r\ge b/2\}$, with
$\mathcal H$ the whole space if $b=0$.
The triangle inequalities give
\[
 \Delta\le b+d_0,\qquad
 \F{D_{i_o}}\le b+\Delta\le d_0+2b.
\]
Consequently $i_o\in\mathcal I(r)$ on $\mathcal H$.

Write $d=\F{A-A_{\widehat i}}$. On $\mathcal E\cap\mathcal H$,
score minimality and \eqref{expected:eq:score} imply
\begin{equation}\label{expected:eq:comparison}
 d^2-\Delta^2
 \le2\eta b(\F{D_{\widehat i}}+\F{D_{i_o}})
 \le2\eta b(d+\Delta+2b).
\end{equation}
Set $\beta=\alpha/(2+\alpha)$. The two inequalities
\[
 2\eta bd\le\beta d^2+\eta^2b^2/\beta,
 \qquad
 2\eta b\Delta\le\beta\Delta^2+\eta^2b^2/\beta
\]
give
\[
 d^2\le(1+\alpha)\Delta^2
 +\left(2(2+\alpha)\eta+
             \frac{(2+\alpha)^2\eta^2}{\alpha}\right)b^2
 \le(1+\alpha)\Delta^2+15\eta b^2.
\]
The last step uses $\alpha\le1$ and $\eta\le\alpha$.

No score event is needed for the pathwise cutoff bound
\begin{equation}\label{expected:eq:pathwise}
 d\le b+d_0+4r\le\Delta+2b+4r,
\end{equation}
where $d_0\le\F{D_{i_o}}\le b+\Delta$ was used.
Since $\mathcal E$ is independent of $r$, Lemma~\ref{expected:lem:radius}
and $(x+y+z)^2\le3(x^2+y^2+z^2)$ imply
\[
 \E[d^2\mathbf1_{\mathcal E^c}]
 \le\delta(3\Delta^2+12b^2+48\E r^2)
 =\delta(3\Delta^2+60b^2).
\]
On $\mathcal H^c$, \eqref{expected:eq:pathwise} gives
$d\le\Delta+4b$ and hence
\[
 \E[d^2\mathbf1_{\mathcal H^c}]
 \le\delta(2\Delta^2+32b^2).
\]
Add these two failure allowances to the good-event bound.
Overcounting their intersection only increases the bound and yields
\eqref{expected:eq:oracle}. The radius costs $O(\log(1/\delta))$
additional queries.
\end{proof}

\begin{corollary}[Pure expected squared relative error]\label{expected:cor:warm}
Suppose a known $\gamma\ge1$ satisfies $b\le\gamma\Delta$.
For $0<\epsilon\le1$, the modified selector can guarantee
\[
 \E\F{A-A_{\widehat i}}^2\le(1+\epsilon)\Delta^2
\]
with
\begin{equation}\label{expected:eq:warm-cost}
 O\!\left(\frac{\gamma^2}{\epsilon}
       \sqrt{\log(M\gamma^2/\epsilon)}
       +\log(\gamma^2/\epsilon)\right)
\end{equation}
queries. In particular, fixed accuracy and a constant warm factor
cost $O(\sqrt{\log(M+1)})$ queries, including expectation control.
\end{corollary}
\begin{proof}
In Theorem~\ref{expected:thm:oracle} take
\[
 \alpha=\epsilon/2,\qquad
 \eta=\epsilon/(60\gamma^2),\qquad
 \delta=\epsilon/(400\gamma^2).
\]
The excess coefficient over $\Delta^2$ is at most
\[
 \epsilon/2+5\epsilon/(400\gamma^2)
             +15\epsilon/60+92\epsilon/400
 \le397\epsilon/400<\epsilon.
\]
If $\Delta=0$, the preliminary residual check returns $A$ almost
surely. This endpoint check does not require the warm hypothesis.
\end{proof}

\begin{corollary}[Random training output and a mean warm bound]
\label{expected:cor:training}
Let a preceding training transcript $\mathcal T$ determine a pivot
and at most $M$ fully known candidates. Draw all selector randomness
independently of $\mathcal T$. Then
\[
 \E\F{A-A_{\widehat i}}^2
 \le(1+\alpha+5\delta)\E\Delta(\mathcal T)^2
      +(15\eta+92\delta)\E b(\mathcal T)^2.
\]
In particular, if a known $\Gamma\ge1$ satisfies
\[
 \E b(\mathcal T)^2\le\Gamma^2\E\Delta(\mathcal T)^2<\infty,
\]
then the parameter choice of Corollary~\ref{expected:cor:warm}, with
$\gamma=\Gamma$, yields
\[
 \E\F{A-A_{\widehat i}}^2
    \le(1+\epsilon)\E\Delta(\mathcal T)^2.
\]
The warm condition need not hold separately for every transcript.
\end{corollary}
\begin{proof}
Apply Theorem~\ref{expected:thm:oracle} conditional on $\mathcal T$ and take
expectations. Its parameters are fixed from $M,\Gamma,\epsilon$,
so no transcript-dependent confidence or accuracy is used.
\end{proof}

\begin{theorem}[A pivot-free expected validator for smaller lists]
\label{expected:thm:ordinary}
For any fixed finite family of $M$ known matrices, let $q\ge8$ be a
multiple of four. Draw a fresh standard Gaussian $G\in\R^{n\times q}$
and minimize $\F{(A-A_i)G}^2$. Then
\begin{equation}\label{expected:eq:ordinary}
 \E\F{A-A_{\widehat i}}^2
 \le2\sqrt{1+2/q}\,M^{4/q}\Delta^2.
\end{equation}
In particular $q=8\lceil\max\{1,\log M\}\rceil$ gives a factor
strictly less than four in expected squared error, using only $q$
right queries and no warm pivot. The same conclusion holds
conditional on an independently generated random family.
\end{theorem}
\begin{proof}
The zero-optimum case follows from the exact residual-sketch check.
Otherwise put $F_i=\F{A-A_i}^2$ and
\[
 T_i=\frac{\F{(A-A_i)G}^2}{qF_i}.
\]
For each fixed residual, $T_i$ is a convex combination of independent
$\chi_q^2/q$ variables. Consequently
\[
 \E T_i^2\le1+2/q.
\]
For the integer $p=q/4\ge2$, Jensen's inequality and the inverse
chi-square moment give
\[
 \E T_i^{-p}
 \le\prod_{j=1}^{p}\frac{q}{q-2j}\le2^p.
\]
Let $i_o$ minimize $F_i$. Score minimality yields
\[
 F_{\widehat i}\le F_{i_o} T_{i_o}\max_iT_i^{-1}.
\]
No independence among the scores is asserted. Cauchy--Schwarz,
monotonicity of $L_p$ means, and a sum bound give
\begin{align*}
 \E[F_{\widehat i}]/F_{i_o}
 &\le\sqrt{\E T_{i_o}^2}
              \sqrt{\E(\max_iT_i^{-1})^2}\\
 &\le\sqrt{1+2/q}\,
                \bigl(\E\max_i T_i^{-p}\bigr)^{1/p}\\
 &\le\sqrt{1+2/q}\,
                \left(\sum_i\E T_i^{-p}\right)^{1/p}
 \le2\sqrt{1+2/q}\,M^{4/q}.
\end{align*}
For the stated choice of $q$, the factor is at most
$2\sqrt{5/4}\,e^{1/2}<4$.
\end{proof}

\paragraph{What this resolves for the HSS interface.}
An $O(k)$-query training list with $\log M=O(k^2)$ and an appropriate
constant warm pivot can now be followed by an $O(k)$-query selector
with a \emph{proved expected squared-error} conclusion. The
failure-event allowance no longer depends on the largest candidate
error or an ambient matrix dimension. Selection preserves any rank
restriction satisfied by every candidate.

There is a cleaner pivot-free sufficient interface if a training
procedure produces at most $\exp(C_0 k)$ candidates using $C_1k$
queries and satisfies
\[
 \E_{\rm train}\min_i\F{A-A_i}^2
       \le C_2L\operatorname{OPT}(A).
\]
Theorem~\ref{expected:thm:ordinary} then proves the full \emph{query and error}
target with an additional $O(k)$ queries and an absolute factor at
most four, without any moment assumption on a pivot or on each
individual candidate. Constructing such a list remains unproved.
For example, a useful list polynomial in $k$ would meet the query
allowance and avoid an exponential list-enumeration cost. A merely
polynomial-in-$n$ list is not automatically enough for $O(k)$
validation when $\log n$ can greatly exceed $k$.

\paragraph{Query complexity versus additional arithmetic.}
This is a pure exact-real query theorem, with computation on known
candidates treated as free. Enumerating $M=\exp(O(k^2))$ candidates,
computing their singular-value truncations, and scoring them can require
exponential additional arithmetic. If polynomial additional arithmetic
is also required, this conditional selection implication does not
establish that requirement. Such a computational result would also
need an implicit representation and an efficient selection procedure,
or a polynomial-size useful list.

\paragraph{What it does not resolve.}
The candidates must already be known matrices. Candidate-generation
and pivot-generation queries remain part of the budget. A pivot that
is warm only with fixed high probability does not automatically meet
the mean hypothesis of Corollary~\ref{expected:cor:training}; its bad-transcript
second moment is still relevant. Nor does a supplied coefficient space
count as a finite list of known matrices.

For a constant-size list, Section~\ref{base:sec:validation} gives a
simpler expected guarantee using ordinary Gaussian residual validation.
The present result matters when the list grows as $\exp(O(k^2))$,
where that constant-size argument has an unacceptable cardinality
factor. It does not establish that such a useful list can be generated
with $O(k)$ queries. Indeed, a warm HSS pivot with an expected
$O(L)\operatorname{OPT}$ error would already meet the approximation
target before selection. The new lemma removes the validation
expectation gap; the training and stability problem remains.

\endgroup

% ===== FALLBACK =====
\begingroup
\let\R\relax
\newcommand{\R}{\mathbb R}
\let\E\relax
\newcommand{\E}{\mathbb E}
\let\Prb\relax
\newcommand{\Prb}{\mathbb P}
\let\T\relax
\newcommand{\T}{\mathsf T}
\let\F\relax
\newcommand{\F}[1]{\left\lVert#1\right\rVert_{F}}
\let\N\relax
\newcommand{\N}[1]{\left\lVert#1\right\rVert}
\let\ip\relax
\newcommand{\ip}[2]{\left\langle#1,#2\right\rangle_F}
\let\tr\relax
\newcommand{\tr}{\operatorname{tr}}
\let\rank\relax
\newcommand{\rank}{\operatorname{rank}}
\let\diag\relax
\newcommand{\diag}{\operatorname{blockdiag}}
\let\range\relax
\newcommand{\range}{\operatorname{range}}
\let\OPT\relax
\newcommand{\OPT}{\operatorname{OPT}}
\let\HSS\relax
\newcommand{\HSS}{\operatorname{HSS}}
\let\diver\relax
\newcommand{\diver}{\operatorname{div}}
\let\Var\relax
\newcommand{\Var}{\operatorname{Var}}
\let\UC\relax
\newcommand{\UC}{\mathcal U}
\let\VC\relax
\newcommand{\VC}{\mathcal V}
\let\err\relax
\newcommand{\err}{\mathcal E}
\let\proofstep\relax
\newcommand{\proofstep}[1]{\par\addvspace{.45\baselineskip}\Needspace{3\baselineskip}\textbf{#1}\par\nobreak}
\let\op\relax
\newcommand{\op}[1]{\left\lVert#1\right\rVert_{\mathrm{op}}}
\let\Sn\relax
\newcommand{\Sn}[2]{\left\lVert#1\right\rVert_{S_{#2}}}
\let\Cov\relax
\newcommand{\Cov}{\mathsf C}
\let\Loss\relax
\newcommand{\Loss}{\mathsf L}
\let\Fil\relax
\newcommand{\Fil}{\mathcal F}
\let\Id\relax
\newcommand{\Id}{I}
\let\supp\relax
\newcommand{\supp}{\operatorname{supp}}
\let\vecop\relax
\newcommand{\vecop}{\operatorname{vec}}
\let\Gammafun\relax
\newcommand{\Gammafun}{\operatorname{Gamma}}
\chapter{Leaf-diagonal safeguards and translation invariance}\label{ch:fallback}
The general HSS-relative target remains unproved. This chapter records an unconditional all-input safeguard for the raw decoder, followed by a precisely qualified invariance statement. The safeguard does not assume a warm HSS pivot, raw inverse integrability, or a basis-learning theorem.

Let the fixed original leaves have size \texttt{m=2k}. Let \texttt{D} be the linear space of matrices supported on these leaf diagonal blocks, and write \texttt{off\_D(A)=A-P\_D\ A}. These leaf-diagonal matrices have zero exterior ranks and belong to every fixed nested-basis HSS coefficient space.

\section{A stable candidate available from the existing training sketch}\label{fallback:a-stable-candidate-available-from-the-existing-training-sketch}

The raw algorithm has already queried \texttt{Y=A\ Omega}, with \texttt{Omega} standard Gaussian of width \texttt{s=8k}. For each original leaf \texttt{I}, form

\[
 C_I=Y_I\Omega_I^\dagger,
 \qquad C=\operatorname{blockdiag}_I C_I.
\]

This costs no further queries. The local designs have full row rank almost surely. For every fixed input \texttt{A},

\[
 \boxed{\quad
 \mathbb E\|A-C\|_F^2
 =\left(1+\frac{2k}{6k-1}\right)
       \|\operatorname{off}_{\mathcal D}(A)\|_F^2.
 \quad}
\]

Indeed,

\[
 C_I-A_{II}=A_{I,I^c}\Omega_{I^c}\Omega_I^\dagger.
\]

Conditional on \texttt{Omega\_I}, the other row blocks remain independent standard Gaussian, so the mean squared coefficient error is

\[
 \|A_{I,I^c}\|_F^2\,\mathbb E\|\Omega_I^\dagger\|_F^2
 =\frac{2k}{6k-1}\|A_{I,I^c}\|_F^2.
\]

Different leaf coefficient errors occupy disjoint diagonal blocks, and the original off-diagonal residual is orthogonal to them. Summing proves the identity. No independence between different leaf estimates is required.

\section{An unconditional all-input wrapper}\label{fallback:an-unconditional-all-input-wrapper}

Construct the canonical raw \texttt{18k}-query HSS output \texttt{B\_raw} and the candidate \texttt{C} above from the same training data. Draw a fresh Gaussian matrix \texttt{G} with eight columns, query \texttt{AG}, and return whichever of \texttt{C,B\_raw} minimizes the squared residual sketch norm.

For any two fixed candidates independent of \texttt{G}, the Gaussian validation lemma in the preliminary propagation analysis gives

\[
 \mathbb E_G\|A-B_{\rm sel}\|_F^2
 \le c\min\{\|A-C\|_F^2,\|A-B_{\rm raw}\|_F^2\},
 \qquad c=1+\sqrt{10/3}<3.
\]

Conditioning first on the complete training transcript yields the following all-input conclusion:

\[
 \boxed{\quad
 \mathbb E\|A-B_{\rm sel}\|_F^2
 \le c\,\mathbb E\min\{\|A-C\|_F^2,\|A-B_{\rm raw}\|_F^2\}
 <4\|\operatorname{off}_{\mathcal D}(A)\|_F^2.
 \quad}
\]

The strict factor four follows from

\[
 c\left(1+\frac{2k}{6k-1}\right)
 \le\frac75\left(1+\sqrt{10/3}\right)<4.
\]

The query cost is \texttt{18k+8\ \textless{}=\ 26k}. Both candidates are HSS rank at most \texttt{k}, and selection preserves this restriction. All additional arithmetic is polynomial. The displayed bound proves integrability of the selected error even if the raw decoder has no finite second moment. If \texttt{A} is exactly HSS rank at most \texttt{k}, the raw candidate equals \texttt{A} almost surely; fresh validation selects an exact candidate almost surely. Thus the wrapper retains the exact zero-optimum endpoint.

This is not an HSS-relative approximation theorem: \texttt{\textbar{}\textbar{}off\_D(A)\textbar{}\textbar{}\^{}2/OPT\_HSS(A)} can be arbitrarily large or infinite. It is a rigorously stable fallback, not a replacement for the missing learner estimate.

\section{Translation and scaling on full-rank raw trajectories}\label{fallback:translation-and-scaling-on-full-rank-raw-trajectories}

Fix a realization of the two raw sketches. Suppose every local design along the raw trajectory for \texttt{E} has full row rank. Let \texttt{D} be any original-leaf-diagonal matrix, and let \texttt{epsilon\ !=\ 0}. With the same sketches and the same Gram-based canonical basis rule, the raw trajectory on \texttt{D+epsilon\ E} has exactly the same row and column bases as that on \texttt{E}, and

\[
 \boxed{\quad
 B_{\rm raw}(D+\epsilon E)=D+\epsilon B_{\rm raw}(E).
 \quad}
\]

At a local block, adding common-coefficient observations \texttt{(C\ Omega,C\^{}T\ Psi)} adds \texttt{C} to both local regressions, leaves both innovations unchanged, and changes the discrepancy by \texttt{C-PCQ}. The retained observations gain \texttt{(U\^{}T\ C\ V)(V\^{}T\ Omega)} and its transpose counterpart. An original-leaf-diagonal matrix remains block diagonal after every compression and sibling merge, so the argument iterates. Multiplication by \texttt{epsilon} scales the innovation Grams by \texttt{epsilon\^{}2}; the canonical bases therefore remain the same. The added discrepancy terms and the added exact root core telescope to \texttt{D}.

The full-row-rank hypothesis is material. For the specified rank-deficient pseudoinverse extension, adding \texttt{C\ Omega} changes the regression by \texttt{C\ Omega\ Omega\^{}dagger}, which need not equal \texttt{C}. This chapter does not prove full row rank almost surely for every arbitrary non-HSS input, and does not silently extend the translation statement to that locus.

Independently of any algorithm, the HSS class is invariant under leaf-diagonal translations and nonzero scalar multiplication. Consequently,

\[
 \operatorname{OPT}_{\rm HSS}(D+\epsilon E)
 =\epsilon^2\operatorname{OPT}_{\rm HSS}(E).
\]

On the full-rank trajectory locus, the learned coefficient space is identical for the two inputs and contains \texttt{D}, so its approximation error also scales by \texttt{epsilon\^{}2}.

Thus an arbitrary raw-decoder difficulty can persist along inputs converging to a leaf-diagonal exact HSS matrix. Merely being close to an exact input, or small compared with \texttt{\textbar{}\textbar{}D\textbar{}\textbar{}}, does not supply the uniform residual-relative estimate. The safeguard handles neighborhoods of leaf-diagonal matrices through an actual known candidate, but extending this to varying nonzero cut ranks still requires a new argument.

\endgroup

% ===== HALFSOLVE =====
\begingroup
\let\R\relax
\newcommand{\R}{\mathbb R}
\let\E\relax
\newcommand{\E}{\mathbb E}
\let\Prb\relax
\newcommand{\Prb}{\mathbb P}
\let\T\relax
\newcommand{\T}{\mathsf T}
\let\F\relax
\newcommand{\F}[1]{\left\lVert#1\right\rVert_{F}}
\let\N\relax
\newcommand{\N}[1]{\left\lVert#1\right\rVert}
\let\ip\relax
\newcommand{\ip}[2]{\left\langle#1,#2\right\rangle_F}
\let\tr\relax
\newcommand{\tr}{\operatorname{tr}}
\let\rank\relax
\newcommand{\rank}{\operatorname{rank}}
\let\diag\relax
\newcommand{\diag}{\operatorname{blockdiag}}
\let\range\relax
\newcommand{\range}{\operatorname{range}}
\let\OPT\relax
\newcommand{\OPT}{\operatorname{OPT}}
\let\HSS\relax
\newcommand{\HSS}{\operatorname{HSS}}
\let\diver\relax
\newcommand{\diver}{\operatorname{div}}
\let\Var\relax
\newcommand{\Var}{\operatorname{Var}}
\let\UC\relax
\newcommand{\UC}{\mathcal U}
\let\VC\relax
\newcommand{\VC}{\mathcal V}
\let\err\relax
\newcommand{\err}{\mathcal E}
\let\proofstep\relax
\newcommand{\proofstep}[1]{\par\addvspace{.45\baselineskip}\Needspace{3\baselineskip}\textbf{#1}\par\nobreak}
\let\op\relax
\newcommand{\op}[1]{\left\lVert#1\right\rVert_{\mathrm{op}}}
\let\Sn\relax
\newcommand{\Sn}[2]{\left\lVert#1\right\rVert_{S_{#2}}}
\let\Cov\relax
\newcommand{\Cov}{\mathsf C}
\let\Loss\relax
\newcommand{\Loss}{\mathsf L}
\let\Fil\relax
\newcommand{\Fil}{\mathcal F}
\let\Id\relax
\newcommand{\Id}{I}
\let\supp\relax
\newcommand{\supp}{\operatorname{supp}}
\let\vecop\relax
\newcommand{\vecop}{\operatorname{vec}}
\let\Gammafun\relax
\newcommand{\Gammafun}{\operatorname{Gamma}}
\chapter{Why two exact half-refits do not imply a small-gain theorem}\label{ch:halfsolve}
This is an abstract countermodel for a proposed proof, not an HSS counterexample. It does not refute global least squares.

\section{Independent small-gain half-maps can accumulate a shared exceptional event}
Suppose a fixed coefficient space splits orthogonally as $\mathcal S=\mathcal L\oplus\mathcal R$. Two unbiased one-sided regressions have independent cross-maps $C:\mathcal R\to\mathcal L$ and $D:\mathcal L\to\mathcal R$. Even if $C$ strictly raises level, $D$ preserves or raises level, both maps are pathwise contractions, and their expected squared gains are small, these properties do not bound $(I-CD)^{-1}$ uniformly in depth.

Take $\mathcal L=\mathcal R=\R^h$, $h\ge2$. Define $Se_i=e_{i+1}$ for $i<h$, $Se_h=0$, and $P=\operatorname{diag}(1,\ldots,1,0)$. Fix $0<a,b<1$. Draw independent $\xi,\eta,\zeta$, where $\xi$ equals $+1$ and $-1$ with probability $a/2$ each and is otherwise zero, $\eta$ has the analogous law with $b$, and $\zeta$ is a uniform sign. Set
\[
 C=\xi S,\qquad D=\eta P.
\]
Then $C,D$ are independent and centered, $\|C\|,\|D\|\le1$, and
\[
 \E C^*C\preceq aI,\qquad \E D^*D\preceq bI.
\]
Also $CD=\xi\eta S$ is nilpotent of index $h$. Nevertheless,
\begin{equation}\label{halfsolve:eq:resolvent}
 \boxed{\E\|(I-CD)^{-1}e_1\|^2=1+ab(h-1).}
\end{equation}
Indeed,
\[
 (I-CD)^{-1}e_1=\sum_{j=0}^{h-1}(\xi\eta)^je_{j+1}.
\]
The vectors are orthogonal, and each nonconstant term has expected squared norm $ab$, not $(ab)^j$: the same two exceptional events are reused at every step.

\section{The same example includes residual-relative half-refit bounds}
Add a one-dimensional residual coordinate $t$, and define
\[
 T_L(l,r,t)=l+Cr+\sqrt a\,\zeta t e_1,\qquad
 T_R(l,r,t)=r+Dl.
\]
These maps are independent, exact on their own model halves, and unbiased. For every deterministic input,
\[
 \E\|T_L(l,r,t)-l\|^2\le a(\|r\|^2+t^2),\qquad
 \E\|T_R(l,r,t)-r\|^2\le b(\|l\|^2+t^2).
\]
The first inequality uses $\E\xi\zeta=0$. On the pure residual $(0,0,1)$, the coupled exact solve gives
\[
 \widehat l=(I-CD)^{-1}\sqrt a\,\zeta e_1,\qquad
 \widehat r=-D\widehat l,
\]
so
\begin{equation}\label{halfsolve:eq:error}
 \boxed{\E\|\widehat l\|^2=a[1+ab(h-1)].}
\end{equation}
The additional $\mathcal R$-output is nonnegative and cannot repair this divergence. For $t=0$, nilpotence makes the coupled solve exact.

Thus exactness, stream independence, nilpotence, and small expected one-step gains are insufficient. A valid fixed-space closure must use additional original-law geometry, as the source--response proof in Part~II does. A global least-squares decoder uses a different positive normal operator and is not ruled out by this countermodel.

\endgroup

% ===== BASISOBS =====
\begingroup
\let\R\relax
\newcommand{\R}{\mathbb R}
\let\E\relax
\newcommand{\E}{\mathbb E}
\let\Prb\relax
\newcommand{\Prb}{\mathbb P}
\let\T\relax
\newcommand{\T}{\mathsf T}
\let\F\relax
\newcommand{\F}[1]{\left\lVert#1\right\rVert_{F}}
\let\N\relax
\newcommand{\N}[1]{\left\lVert#1\right\rVert}
\let\ip\relax
\newcommand{\ip}[2]{\left\langle#1,#2\right\rangle_F}
\let\tr\relax
\newcommand{\tr}{\operatorname{tr}}
\let\rank\relax
\newcommand{\rank}{\operatorname{rank}}
\let\diag\relax
\newcommand{\diag}{\operatorname{blockdiag}}
\let\range\relax
\newcommand{\range}{\operatorname{range}}
\let\OPT\relax
\newcommand{\OPT}{\operatorname{OPT}}
\let\HSS\relax
\newcommand{\HSS}{\operatorname{HSS}}
\let\diver\relax
\newcommand{\diver}{\operatorname{div}}
\let\Var\relax
\newcommand{\Var}{\operatorname{Var}}
\let\UC\relax
\newcommand{\UC}{\mathcal U}
\let\VC\relax
\newcommand{\VC}{\mathcal V}
\let\err\relax
\newcommand{\err}{\mathcal E}
\let\proofstep\relax
\newcommand{\proofstep}[1]{\par\addvspace{.45\baselineskip}\Needspace{3\baselineskip}\textbf{#1}\par\nobreak}
\let\op\relax
\newcommand{\op}[1]{\left\lVert#1\right\rVert_{\mathrm{op}}}
\let\Sn\relax
\newcommand{\Sn}[2]{\left\lVert#1\right\rVert_{S_{#2}}}
\let\Cov\relax
\newcommand{\Cov}{\mathsf C}
\let\Loss\relax
\newcommand{\Loss}{\mathsf L}
\let\Fil\relax
\newcommand{\Fil}{\mathcal F}
\let\Id\relax
\newcommand{\Id}{I}
\let\supp\relax
\newcommand{\supp}{\operatorname{supp}}
\let\vecop\relax
\newcommand{\vecop}{\operatorname{vec}}
\let\Gammafun\relax
\newcommand{\Gammafun}{\operatorname{Gamma}}
\chapter{Adaptive erasure and pathwise rank novelty}\label{ch:basisobs}
These constructions delimit prospective learning arguments. Neither is an information-theoretic lower bound for the general query target, and the first does not prescribe the choices of a canonical learner.

\section{Exact containment and conditioning do not justify adaptive nullification}
Let $I=a\cup b$ have children of size at least $s+1$, with sketch width $s\ge2$. Fix unit vectors $u_a,u_b$ on the children, choose $j\notin I$, and set
\[
 u=2^{-1/2}(u_a+u_b),\qquad A=ue_j^T.
\]
This fixed matrix has HSS rank at most one. On the subtree at $I$, its column exterior spaces are zero; its row exterior spaces are restrictions of $u$. Unused column spaces therefore can be arbitrary without violating exact representation, provided the row spaces contain the restrictions of $u$.

Write $g=e_j^T\Omega$ and choose $a_0=g^T/\|g\|$, with a unit $b_0\perp a_0$. The event
\[
 \sigma_{\min}(\Omega_a)>1,\qquad\sigma_{\min}(\Omega_b)>1
\]
has positive probability. On this event there exist unit $v_a,v_b$ such that
\[
 \Omega_a^Tv_a=a_0,\qquad\Omega_b^Tv_b=b_0.
\]
The minimum-norm solution of each equation has norm below one; a vector from the nonzero transpose nullspace supplies the missing norm. The choices can be measurable. Every $v_c$ has a nested rank-one representation by its normalized nonzero restrictions, with completion on zero restrictions.

Choose row templates $u_a,u_b$ and column templates $v_a,v_b$. The compressed design is
\[
 G=\begin{bmatrix}v_a^T\Omega_a\\v_b^T\Omega_b\end{bmatrix}
 =\begin{bmatrix}a_0^T\\b_0^T\end{bmatrix},\qquad GG^T=I_2.
\]
Its singular values are exactly one. However the compressed row observation is
\[
 T=2^{-1/2}\begin{bmatrix}1\\1\end{bmatrix}g,
\]
and, because $g\in\operatorname{row}(G)$,
\begin{equation}\label{basisobs:eq:erasure}
 \boxed{T(I-G^\dagger G)=0.}
\end{equation}
The true exterior block is nonzero and every lower-level representation loss can be zero. A threshold that retains singular values equal to one does not repair this example.

Therefore no local rangefinder lemma can be uniform over arbitrary sketch-dependent, exactly containing lower-level spaces under conditioning alone. This construction deliberately chooses unused column directions from the forward probe; it does not show that a specified canonical zero-space completion does so. Chapters~\ref{ch:actual} and~\ref{ch:anchor} address later, actual-rule obstructions separately.

\section{The total newly appearing rank along a path can be proportional to depth}
Choose a root-to-leaf chain $I_0\supset I_1\supset\cdots\supset I_h$, coordinates $j_t\in I_t\setminus I_{t+1}$ for $0\le t<h$, and $j_h\in I_h$. Set
\[
 A=\sum_{t=1}^{h}\left(e_{j_{t-1}}e_{j_t}^T+e_{j_t}e_{j_{t-1}}^T\right).
\]
At the cut $I_t$, the only crossing edge is $(j_{t-1},j_t)$, so
\[
 \operatorname{range}A[I_t,I_t^c]
 =\operatorname{span}\{e_{j_t}|_{I_t}\}.
\]
For $t<h$, this range restricts to zero on $I_{t+1}$, whereas the next exterior range is the nonzero direction $e_{j_{t+1}}$.

Every node outside the chain lies in a sibling branch containing at most one selected coordinate. All its exterior edges have the same row factor; its exterior rank is at most one. The final descendants have the same property, and symmetry supplies the column condition. Thus the matrix has HSS rank at most one everywhere while a wholly new direction appears at each of $h$ levels.

Replacing selected coordinates by disjoint rank-$k$ frames and edges by identity couplings gives total new-direction dimension $kh$, when the branch sizes allow it. This refutes charging all pathwise novelty to $k$. It does not imply that $kh$ queries are necessary: different nodes can be learned in parallel.

\endgroup

% ===== NUCLEAR =====
\begingroup
\let\HH\relax
\newcommand{\HH}{\mathcal H}
\let\MM\relax
\newcommand{\MM}{\mathcal M}
\let\JJ\relax
\newcommand{\JJ}{\mathcal J}
\let\E\relax
\newcommand{\E}{\mathbb E}
\let\PP\relax
\newcommand{\PP}{\mathbb P}
\chapter{Failure of depth-weighted tree nuclear regularization}\label{ch:nuclear}
This chapter rules out a specific global regularizer. It is not an
impossibility theorem for nonlinear HSS recovery, and it does not concern
the fixed-basis refit whose source--response estimate has been closed.

Let a balanced binary hierarchy partition $N=2n$ coordinates, with root
children $I,J$ of size $n$. Within each root half, let depth $\ell$ have
$m_\ell=2^\ell$ equal-size nodes, for $\ell=0,\ldots,H$.
Write $\HH_k$ for matrices whose row and column exterior cuts at every
nonroot node have rank at most $k$.
Let $\Omega,\Psi\in\mathbb R^{2n\times q}$ be independent standard Gaussian
matrices, and set
\[
 \MM(X)=(X\Omega,X^{\mathsf T}\Psi).
\]
For arbitrary nonnegative depth weights $w_\ell$, not all zero, define
\begin{equation}\label{nuclear:eq:penalty}
 \JJ(X)=
 \sum_{\ell=0}^H w_\ell
 \sum_{\substack{S\text{ a node at depth }\ell\\
                 \text{within }I\text{ or }J}}
 \left(\|X(S,S^c)\|_*+\|X(S^c,S)\|_*\right).
\end{equation}
The weights may depend on the hierarchy and dimensions; they are the
same for all nodes at a given depth.

\begin{theorem}\label{nuclear:thm:failure}
Suppose $k\ge2$ and $n>4q$. There is a deterministic rank-one matrix
$B\in\HH_1$ with $\|B\|_F=1$ such that, with probability at least
$1-4q/n$, a matrix $C\in\HH_2\subseteq\HH_k$ satisfies
\[
 \MM(C)=\MM(B),\qquad
 \JJ(C)<\JJ(B),\qquad
 \|C\|_F<1/\sqrt2.
 \tag{1}
\]
Consequently, neither minimizing $\JJ$ over exact HSS-rank-$k$
completions nor minimizing
\[
 \|\MM(X)-\MM(B)\|^2+\lambda\JJ(X),\qquad X\in\HH_k,
 \qquad\lambda>0,
 \tag{2}
\]
can return $B$ on this event. In particular these selectors cannot
satisfy an all-input expected residual-relative approximation bound.
\end{theorem}

\begin{proof}
Let $u,v\in\mathbb R^n$ be the constant unit vectors
$u=v=n^{-1/2}(1,\ldots,1)^{\mathsf T}$, and let $B$ vanish except for
\[
 B(I,J)=uv^{\mathsf T}.
\]
Every exterior cut has rank at most one.

Let $P$ be the orthogonal projector onto
$\operatorname{range}(\Psi(I,:))$, and let $Q$ project onto
$\operatorname{range}(\Omega(J,:))$. Both have rank $q$ almost surely.
Define $C$ to vanish outside $I\times J$, and put
\begin{align}
 C(I,J)
 &=Puv^{\mathsf T}+uv^{\mathsf T}Q-Puv^{\mathsf T}Q \notag\\
 &=(Pu)v^{\mathsf T}+(u-Pu)(Qv)^{\mathsf T}.
 \label{nuclear:eq:competitor}
\end{align}
Thus $C$ has rank at most two, and in particular belongs to $\HH_2$.
The difference in the only nonzero block is
\[
 B(I,J)-C(I,J)=(I-P)uv^{\mathsf T}(I-Q).
 \tag{3}
\]
Since $(I-Q)\Omega(J,:)=0$ and $(I-P)\Psi(I,:)=0$, (3) proves
$\MM(C)=\MM(B)$.

Set $a^2=\|Pu\|^2$ and $b^2=\|Qv\|^2$. The observed three blocks in the
orthogonal row/column decompositions determined by $P,Q$ give
\begin{equation}\label{nuclear:eq:norm}
 \|C\|_F^2=a^2+b^2-a^2b^2\le a^2+b^2.
\end{equation}
Rotational invariance of a Gaussian column space implies
$\E P=\E Q=(q/n)I$, so
\[
 \E(a^2+b^2)=2q/n,\qquad
 \PP\{a^2+b^2<1/2\}\ge1-4q/n.
 \tag{4}
\]
Only this elementary first-moment estimate is needed.

It remains to compare all tree cuts, rather than only the root cut.
Fix depth $\ell$ and write $m=2^\ell$. For a node $S\subseteq I$,
the only nonzero cut is $B(S,S^c)=B(S,J)$, whose nuclear norm is
$\|u_S\|\|v\|=m^{-1/2}$. A node $S\subseteq J$ contributes
the same amount through $B(S^c,S)$. Hence the unweighted depth
contribution of $B$ is exactly $2\sqrt m$.

Every restriction of $C$ has rank at most two. The row restrictions
$C(S,J)$ over the $m$ nodes $S\subseteq I$ have disjoint rows. Therefore
\begin{align*}
 \sum_{\substack{S\subseteq I\\S\text{ at depth }\ell}}
 \|C(S,J)\|_*
 &\le\sqrt2
 \sum_{\substack{S\subseteq I\\S\text{ at depth }\ell}}
 \|C(S,J)\|_F\\
 &\le\sqrt{2m}\,\|C\|_F.
\end{align*}
The analogous column-restriction bound over nodes in $J$ is identical.
Thus the depth contribution of $C$ is at most
$2\sqrt{2m}\|C\|_F$. Summing with the nonnegative weights proves
\begin{equation}\label{nuclear:eq:tree-comparison}
 \JJ(C)\le\sqrt2\,\|C\|_F\,\JJ(B).
\end{equation}
On the event in (4), (\ref{nuclear:eq:norm}) and
(\ref{nuclear:eq:tree-comparison}) give the strict inequality in (1).
Because $\JJ(B)>0$, $B$ is not a minimizer of either stated objective.

Finally, for the input $A=B$ the best HSS-rank-$k$ residual is zero.
An expected residual-relative bound would require the decoder to return
$B$ almost surely. Its failure on an event of positive probability
contradicts that requirement.
\end{proof}

\paragraph{Scope.}
The result also rules out unconstrained exact-completion minimization
of (\ref{nuclear:eq:penalty}), since the same competitor is feasible there.
Adding nonnegative multiples of the global nuclear norm or Frobenius
norm does not repair this example:
$\|C\|_*\le\sqrt2\|C\|_F<1=\|B\|_*$ on the event above.
Penalties on diagonal leaf blocks do not distinguish $B$ and $C$,
because those blocks vanish.

The obstruction uses a rank transition: the true matrix has rank one
while the permitted competitor has rank two. It therefore does not
contradict generic uniqueness for fully saturated rank-$k$ inputs.
It also does not rule out a selector that first minimizes the vector
of exterior cut ranks, a nonconvex rank penalty, or a robust
rank-selection algorithm. Establishing polynomial-time, quantitative
stability for such a method remains a separate question.

\endgroup

% ===== SLACK =====
\begingroup
\let\F\relax
\newcommand{\F}[1]{\lVert#1\rVert_F}
\let\op\relax
\newcommand{\op}[1]{\lVert#1\rVert_{\mathrm{op}}}
\chapter{Deterministic rank slack and the two-sided mismatch}\label{ch:slack}
\noindent
\textbf{Status.} This chapter proves a local estimate for the actual leading-singular-subspace operation. It does not prove the all-input basis-learning theorem, close the two-sided feedback, or supply a new query-complexity bound. The fixed-basis refit is a separate result.

\begin{lemma}[A low-rank view of the discarded sample]
Let $S=C+F\in\mathbb R^{m\times q}$, with $\operatorname{rank}C\le k$, and let $k\le r<\min(m,q)$. Let $P_r$ be the orthogonal projector onto any leading $r$ left singular vectors of $S$. Then
\begin{equation}\label{slack:eq:spectral}
 \op{(I-P_r)S}^2
 \le \frac{\F F^2}{r-k+1}.
\end{equation}
For every matrix $Z$ with $q$ rows, rank at most $k$, and $\op Z\le1$,
\begin{equation}\label{slack:eq:lowrank}
 \F{(I-P_r)SZ}^2
 \le \frac{k}{r-k+1}\F F^2.
\end{equation}
All matrices may depend on the same randomness. No probabilistic independence or spectral gap is required.
\end{lemma}
\begin{proof}
The variational characterization of singular values and the best rank-$k$ approximation identity give
\[
 (r-k+1)\sigma_{r+1}(S)^2
 \le\sum_{j=k+1}^{r+1}\sigma_j(S)^2
 \le\sum_{j>k}\sigma_j(S)^2
 \le\F{S-C}^2.
\]
Also $\op{(I-P_r)S}=\sigma_{r+1}(S)$, which proves~\eqref{slack:eq:spectral}. The product $(I-P_r)SZ$ has rank at most $k$ and operator norm at most $\op{(I-P_r)S}\op Z$. Bounding its squared Frobenius norm by $k$ times its squared operator norm proves~\eqref{slack:eq:lowrank}. For general $Z$ the right side is multiplied by $\op Z^2$.
\end{proof}

\paragraph{Connection to the actual local learner.}
With the notation of the original reused-sketch construction, put
\[
 B_Y=Y\Omega^\dagger,\quad
 B_Z=(\Psi^\dagger)^T Z_{\rm obs}^T,\quad
 E_Y=Y-B_Y\Omega,\quad E_Z=Z_{\rm obs}-B_Z^T\Psi.
\]
Let $P,Q$ be the actual leading-$r$ projectors of $E_Y,E_Z$, and use the raw discrepancy
\[
 D=(I-P)B_Y+PB_Z(I-Q).
\]
Direct multiplication gives the exact identities
\begin{align}
 (I-P)(Y-D\Omega)&=(I-P)E_Y,\label{slack:eq:forward}\\
 (I-Q)(Z_{\rm obs}-D^T\Psi)
 &=(I-Q)E_Z
 +(I-Q)(B_Z^T-B_Y^T)(I-P)\Psi.\label{slack:eq:transpose}
\end{align}
Consequently the lemma applies directly to rank-$k$ contractions of the forward discarded observation whenever a decomposition $E_Y=C_Y+F_Y$, $\operatorname{rank}C_Y\le k$, is available. It also applies to the first term on the right of~\eqref{slack:eq:transpose}. The regression-mismatch term in that display is not covered by the lemma.

If local size is $2r$, width is $3r$, both designs have full row rank, and the learner spans every innovation column, the two innovation residuals vanish. The mismatch can remain. Exact innovation capture therefore does not prove a stable two-sided decoder.

\paragraph{Why the estimate cannot be moved to the unknown true signal.}
A proposed inequality
\[
 \F{(I-P_r)C}^2\le\gamma\F F^2,\qquad \gamma<1,
\]
cannot hold for all $S=C+F$. Take $S=0$, choose a nonzero rank-one $C$ with range orthogonal to the selected $P_r$, and set $F=-C$. Both sides before multiplication by $\gamma$ are equal. Oversampling gives a small factor on a low-rank view of the \emph{discarded observed sample}; it does not by itself give a small factor on an arbitrary comparator defect.

\paragraph{Remaining implication that has not been proved.}
For $r=ck$, the coefficient in~\eqref{slack:eq:lowrank} is at most $1/(c-1)$ when $c>1$. A useful global argument would have to express inherited learning error through products of the controlled form in~\eqref{slack:eq:lowrank}, bound their multipliers, and control the remaining term in~\eqref{slack:eq:transpose} relative to the fixed best-HSS residual. None of those global claims follows from the local lemma. A rank-adaptive, unpadded variant is another possible direction, but no threshold rule preserving the exact-input endpoint and the desired all-input error bound has been established here.

\paragraph{Relation to the query-only bicriteria reduction.}
If an $O(k)$-query learner did produce a known $C_{\rm out}$ satisfying
\[
 \mathbb E\F{A-C_{\rm out}}^2\le\alpha L\,\mathrm{OPT}_k(A),
\]
exact offline projection of $C_{\rm out}$ onto the original HSS rank-$k$ class would give
\[
 \mathbb E\F{A-\widehat B}^2
 \le(2\sqrt{\alpha L}+1)^2\mathrm{OPT}_k(A)
 \le(2\sqrt\alpha+1)^2L\,\mathrm{OPT}_k(A).
\]
This established reduction allows intermediate rank $r=O(k)$ without more queries, but does not construct the required learner or assert efficient nonlinear optimization.

\endgroup

\part{Commutator learning, actual-rule obstructions, and the union}

% ===== COMM =====
\begingroup
\let\E\relax
\newcommand{\E}{\mathbb E}
\let\R\relax
\newcommand{\R}{\mathbb R}
\let\F\relax
\newcommand{\F}[1]{\lVert #1\rVert_F}
\let\op\relax
\newcommand{\op}[1]{\lVert #1\rVert_{\mathrm{op}}}
\let\ran\relax
\newcommand{\ran}{\operatorname{range}}
\let\rank\relax
\newcommand{\rank}{\operatorname{rank}}
\let\diag\relax
\newcommand{\diag}{\operatorname{blockdiag}}
\let\OPT\relax
\newcommand{\OPT}{\operatorname{OPT}}
\let\dist\relax
\newcommand{\dist}{\operatorname{dist}}
\let\tr\relax
\newcommand{\tr}{\operatorname{tr}}
\let\HH\relax
\newcommand{\HH}{\mathcal H}
\let\SSS\relax
\newcommand{\SSS}{\mathcal S}
\let\J\relax
\newcommand{\J}{\mathcal J}
\let\Id\relax
\newcommand{\Id}{\mathrm I}
\chapter{Original-data commutators and adaptive lifting}\label{ch:comm}
\begin{scopebox}
\textbf{Update, 20 September 2026.} The inverse-first commutator learner defined here is not the later full-range, inverse-optimally completed candidate, and neither is the completed synchronized learner. The later full-range candidate has its own actual noisy-input failure, recorded in Chapter~\ref{upd:chap:later-obstructions}. That result is not a retraction of the commutator identities or an automatic counterexample to every raw or annihilated learner.
\end{scopebox}

\begin{scopebox}
This inverse-first original-data learner is not the later sketch-metric un-nullified companion. The all-depth basis estimate here is a sufficient, unproved condition. Chapters~\ref{ch:actual} and~\ref{ch:ports} rule out a depth-uniform unweighted inverse argument for this rule, not its general approximation risk.
\end{scopebox}
\thispagestyle{plain}
\begin{chaptersynopsis}
The fixed-space refit is taken as a completed interface; the unresolved task is constructing good nested bases from a fixed number of sketch columns. We derive a basis-only observable, the two-sided sketch of a coordinate commutator, that can be formed at every tree node from the original two matrix products. Its comparator rank is at most $2k$, and its total expected residual energy is at most $2q^2 L\OPT_k(A)$. A cross-sample exact-capture theorem avoids the particular adaptive-nullification obstruction, on full-row-rank compressed designs. An exact local defect identity, valid also on singular designs, reduces approximate learning to lifted residuals and the return of previously lost comparator directions. For coarse physical leaves of size $8k$, width $16k$, and retained rank $4k$, the first layer has total expected comparator loss at most $45$ times the squared off-leaf residual. The corresponding all-depth adaptive estimate is not proved. Thus this is a concrete continuation toward, not a proof of, the general $O(k)$-query HSS-relative theorem.
\end{chaptersynopsis}


\section{Target, refitting interface, and scope}
Let $A\in\R^{N\times N}$ and fix a balanced binary coordinate tree with original leaves of size $2k$. Write $\HH_k$ for matrices whose exterior blocks $B[I,I^c]$ and $B[I^c,I]$ have rank at most $k$ at every nonroot node. Put
\[
 \OPT_k(A)=\min_{B\in\HH_k}\F{A-B}^2.
\]
The desired conclusion is an algorithm with at most $C_qk$ queries to $A$ and $A^T$, returning $\widehat B\in\HH_k$ such that
\begin{equation}\label{comm:eq:target}
 \E\F{A-\widehat B}^2\le C L\OPT_k(A).
\end{equation}
Here $L\ge1$ bounds the number of nonroot tree partitions. Enlarging this convention by an absolute constant does not change the target.

\paragraph{Completed interface, proved in Chapter~\ref{ch:uniform}.}
For arbitrary nested rank-$r$ bases fixed independently of the refit probes, the fixed-space theorem gives a matrix $C\in\SSS$ with
\begin{equation}\label{comm:eq:refit}
 \E[\F{A-C}^2\mid\SSS]\le1607\,\dist_F(A,\SSS)^2
\end{equation}
using $8r$ queries. A $36r$-query version has factor $15/4$. This chapter uses this interface only in its final conditional implication; none of its new local proofs depends on the previous refit calculations.

The public fresh-level result of Amsel et al. uses $O(k\log(N/k))$ queries with an $O(\log(N/k))$ expected squared-error factor; reuse is explicitly raised as a possible way to remove the query logarithm \cite{AmselHSS}. Levitt--Martinsson supply the related black-box reused-sketch construction \cite{LM}. No priority claim for the commutator viewpoint or for the elementary identities below is made here.

\paragraph{What changes in this chapter.}
The new learner does not recycle coefficient estimates into later learning observations. It combines the two original products before learning each node. It is a modified learner, not an analysis that silently changes the original one-sided innovation rule. Intermediate rank $4k$ is allowed; final rank $k$ is imposed only after a known approximation is available.

\section{Original-data commutators and their all-depth budget}
Draw independent standard Gaussian matrices $\Omega,\Psi\in\R^{N\times q}$ and query
\[
 Y=A\Omega,\qquad Z=A^T\Psi.
\]
For a coordinate node $I$, let $P_I$ denote its coordinate projector. Define the available $q\times q$ matrix
\begin{equation}\label{comm:eq:J}
 \boxed{\J_I(A)=\Psi_I^T Y_I-Z_I^T\Omega_I
       =\Psi^T(P_IA-AP_I)\Omega.}
\end{equation}
The entries of $A[I,I]$ cancel exactly. If $I=I_1\dot\cup I_2$, then
\begin{equation}\label{comm:eq:add}
 \J_I(A)=\J_{I_1}(A)+\J_{I_2}(A).
\end{equation}
Thus all node commutators can be formed by summing leaf commutators; no additional products with $A$ or $A^T$ are required. The root commutator is zero.

\begin{proposition}[Fixed-comparator rank and noise budget]\label{comm:prop:budget}
Fix $B\in\HH_k$ before drawing the probes, and set $E=A-B$. Then
\begin{align}
 \rank\J_I(B)&\le 2k,\label{comm:eq:rankJ}\\
 \E\F{\J_I(E)}^2
 &=q^2\bigl(\F{E[I,I^c]}^2+\F{E[I^c,I]}^2\bigr).\label{comm:eq:noiseJ}
\end{align}
For the nonroot nodes in at most $L$ partitions,
\begin{equation}\label{comm:eq:totalJ}
 \boxed{\sum_I\E\F{\J_I(E)}^2\le2q^2L\F E^2.}
\end{equation}
No independence between node commutators is asserted or needed.
\end{proposition}
\begin{proof}
Relative to $I\cup I^c$, the commutator has block form
\[
 [P_I,B]=\begin{pmatrix}0&B[I,I^c]\\-B[I^c,I]&0\end{pmatrix}.
\]
Its rank is the sum of the two exterior ranks, at most $2k$. Left and right multiplication do not increase rank. For any deterministic matrix $M$, two successive Gaussian second-moment calculations give
$\E\F{\Psi^T M\Omega}^2=q^2\F M^2$.
The off-diagonal commutator blocks have disjoint supports, proving \eqref{comm:eq:noiseJ}. On one partition, each off-partition entry is counted twice by the sum of row- and column-exterior squared norms. Summing over the partitions proves \eqref{comm:eq:totalJ}.
\end{proof}

A rank-$2k$ truncated SVD $\widehat J_I$ of $\J_I(A)$ also satisfies
\[
 \F{\widehat J_I-\J_I(B)}\le2\F{\J_I(E)},\qquad
 \sum_I\E\F{\widehat J_I-\J_I(B)}^2\le8q^2L\F E^2.
\]
This is a guarantee in sketch coordinates, not yet a guarantee for physical nested bases. The learner below uses $\J_I(A)$ itself, not this optional truncation.

\section{Why the regression mismatch should be observed}
Consider local designs $\Omega_c\in\R^{n\times q}$ and $\Psi_c\in\R^{m\times q}$, both of full row rank, and local observations $Y_c\in\R^{m\times q}$, $Z_c\in\R^{n\times q}$. Write
\[
 B_Y=Y_c\Omega_c^\dagger,\quad B_Z=(\Psi_c^\dagger)^T Z_c^T,
 \quad E_Y=Y_c-B_Y\Omega_c,\quad E_Z=Z_c-B_Z^T\Psi_c,
 \quad\Delta=B_Y-B_Z.
\]
The cross-samples are
\begin{align}
 S_U&=Y_c-B_Z\Omega_c=E_Y+\Delta\Omega_c,\label{comm:eq:crossU}\\
 S_V&=Z_c-B_Y^T\Psi_c=E_Z-\Delta^T\Psi_c.\label{comm:eq:crossV}
\end{align}
For $J_c=\Psi_c^TY_c-Z_c^T\Omega_c$, direct multiplication gives
\begin{equation}\label{comm:eq:crossJ}
 S_U=(\Psi_c^\dagger)^TJ_c,\qquad
 S_V=-(\Omega_c^\dagger)^TJ_c^T.
\end{equation}
Moreover,
\begin{equation}\label{comm:eq:crossrange}
 \ran S_U=\ran[E_Y,\Delta],\qquad
 \ran S_V=\ran[E_Z,\Delta^T].
\end{equation}
For example, $S_U(\Id-\Omega_c^\dagger\Omega_c)=E_Y$ and $S_U\Omega_c^\dagger=\Delta$, while the reverse range inclusion follows from \eqref{comm:eq:crossU}.

In the earlier erasure construction, $Y_c=u g$, $Z_c=0$, and $g$ belongs to the row space of $\Omega_c$. Then $E_Y=0$, but $B_Z=0$ and $S_U=u g$. The cross-sample retains exactly the exterior direction that one-sided nullification erased. This fact does not require a different draw of the original probes.

\begin{remark}[Local versus original commutator]
Equation \eqref{comm:eq:crossJ} uses a local compressed commutator. After approximate compression, this need not equal the original $\J_I(A)$ from \eqref{comm:eq:J}. The new algorithm below explicitly uses the \emph{original} commutator. Its local analysis is derived afresh, including the resulting comparator-return term.
\end{remark}

\section{Fixed Gaussian anchors and adaptive exact capture}
Fix a node $I$ and a deterministic comparator $B$. Factor its exterior blocks as
\begin{equation}\label{comm:eq:factor}
 B[I,I^c]=C R^T,\qquad B[I^c,I]^T=V H,
\end{equation}
where $R$ and $V$ have orthonormal columns, $C\in\R^{|I|\times a}$, $H\in\R^{b\times|I^c|}$, and $a,b\le k$. Zero-dimensional factors are allowed. Define
\begin{equation}\label{comm:eq:anchors}
 \Gamma=R^T\Omega_{I^c},\qquad X=V^T\Omega_I,
 \qquad\Theta=H\Psi_{I^c}.
\end{equation}
The matrices $\Gamma$ and $X$ are independent standard Gaussian matrices of sizes $a\times q$ and $b\times q$. They are tied to the fixed comparator, not to the learned templates. Exactly,
\begin{equation}\label{comm:eq:Jfactor}
 \J_I(B)=\Psi_I^TC\Gamma-\Theta^T X.
\end{equation}

\begin{theorem}[Exact capture despite adaptive child spaces]\label{comm:thm:capture}
Assume $q\ge a+b$. Let $W\in\R^{|I|\times m}$ be any isometry, possibly a measurable function of all probes, such that $\ran C\subseteq\ran W$. Suppose $G=W^T\Psi_I$ has full row rank. Put
\[
 S=(G^\dagger)^T\J_I(B),\qquad C_c=W^TC.
\]
On a single probability-one event depending only on the fixed comparator anchors, simultaneously for every such $W$,
\begin{equation}\label{comm:eq:capture}
 \ran C_c\subseteq\ran S,\qquad\rank S\le a+b.
\end{equation}
Consequently any retained space of dimension $r\ge a+b$ containing the sample range contains the true compressed exterior row space. The analogous assertion holds for column spaces after transposition and exchanging the probes.
\end{theorem}
\begin{proof}
Full row rank gives $(G^\dagger)^TG^T=\Id_m$, and exact containment gives $\Psi_I^TC=G^TC_c$. Hence, with $K=(G^\dagger)^T\Theta^T$,
\[
 S=C_c\Gamma-KX=[C_c,-K]\begin{bmatrix}\Gamma\\X\end{bmatrix}.
\]
The fixed Gaussian stack on the right has full row rank almost surely. Right multiplication by a full-row-rank matrix preserves the column range of its left factor. Thus $\ran S=\ran[C_c,K]$, proving both claims. The same full-rank event works for every adaptive left factor; no conditional Gaussian law for $W$ is used.
\end{proof}

This avoids the specific well-conditioned-design erasure mechanism. It does not prove that every compressed design produced by an all-depth learner has full row rank or a uniformly integrable inverse. In particular, the theorem is not a uniform approximation statement over arbitrary adaptive templates.

\section{A concrete original-commutator learner}
Use $r=4k$ and $q=16k$. For $N\le16k$, exact coordinate queries are an allowed $O(k)$ fallback. Otherwise group four original leaves into coarse leaves of size $8k=2r$.

At a coarse leaf $I$, set $W_{U,I}=W_{V,I}=\Id_{|I|}$. At an internal nonroot node with children $I_1,I_2$, set
\[
 W_{U,I}=\diag(\mathcal U_{I_1},\mathcal U_{I_2}),\qquad
 W_{V,I}=\diag(\mathcal V_{I_1},\mathcal V_{I_2}),
\]
where $\mathcal U,\mathcal V$ are the physical child bases already retained. Form
\begin{align}
 G_I&=W_{U,I}^T\Psi_I,& H_I&=W_{V,I}^T\Omega_I,\\
 S_{U,I}&=(G_I^\dagger)^T\J_I(A),&
 S_{V,I}&=-(H_I^\dagger)^T\J_I(A)^T.\label{comm:eq:learner}
\end{align}
Choose $U_I,V_I$ to span leading $r$ left singular spaces of these samples and set
\[
 \mathcal U_I=W_{U,I}U_I,\qquad\mathcal V_I=W_{V,I}V_I.
\]
Use a fixed measurable singular-value tie rule and coordinate-based zero-space completion. Moore--Penrose inverses define the procedure on singular designs as well. At the root, no exterior space is learned.

\paragraph{Access and computation.}
Only the original $2q=32k$ matvec queries are needed for training. Equation \eqref{comm:eq:add} forms parent data by addition, and the compressed designs propagate as $U_I^TG_I$ and $V_I^TH_I$. Storing local transfers, rather than dense physical bases, gives $O(Nk^2)$ training arithmetic under standard dense local linear algebra. This is an algorithmic cost statement, not an error theorem. The row and column learners use the same $\J_I(A)$ and are not independent.

The learned space is coarse-tree HSS rank at most $r$. Every proper descendant of a coarse leaf has size at most $r$, so its exterior rank is automatically at most $r$. Therefore this is also an original-tree HSS rank-$r$ space.

\paragraph{Invariance.}
For any matrix supported on the original leaf diagonal blocks, all relevant commutators vanish. The learner is pathwise invariant under adding such a matrix. It is also invariant under nonzero scalar multiplication of $A$, with the stated canonical projector rule. These facts remain valid for singular designs; they follow by induction from \eqref{comm:eq:learner}, not from an inverse-integrability assertion.

\section{The all-trajectory local defect identity}
Keep the fixed factors \eqref{comm:eq:factor}. Let $W$ be the actual row template at the node, whether or not it contains $C$. Define
\begin{align}
 G&=W^T\Psi_I,& \mathcal L&=(G^\dagger)^T,& Q_G&=\mathcal L G^T,\\
 C_c&=W^TC,& D&=\mathcal L\Psi_I^T(\Id-WW^T)C,&
 K&=\mathcal L\Theta^T,\\
 F&=\mathcal L\J_I(E),&
 \widetilde D&=D-(\Id-Q_G)C_c.\label{comm:eq:defs}
\end{align}
Here $Q_G$ is the orthogonal projector onto $\ran G\subseteq\R^m$. The term $D$ returns already lost physical comparator directions through a local oblique regression. The additional term in $\widetilde D$ records comparator directions invisible to a singular compressed design.

\begin{lemma}[Exact sample decomposition]\label{comm:lem:sample}
For every realization, including singular $G$,
\begin{equation}\label{comm:eq:sample}
 \boxed{S:=\mathcal L\J_I(A)=(Q_GC_c+D)\Gamma-KX+F.}
\end{equation}
The part before $F$ has rank at most $a+b$.
\end{lemma}
\begin{proof}
Split $C=WW^TC+(\Id-WW^T)C$ in \eqref{comm:eq:Jfactor}. Applying $\mathcal L$ proves the identity, and its two comparator terms have ranks at most $a$ and $b$.
\end{proof}

For $a>0$, let $N_X$ be any orthonormal basis for $\ker X$ and define
\begin{equation}\label{comm:eq:T}
 T=N_X(\Gamma N_X)^\dagger.
\end{equation}
When $b=0$, take $N_X=\Id_q$. For $q\ge a+b$, almost surely
\begin{equation}\label{comm:eq:ann}
 \Gamma T=\Id_a,\qquad XT=0.
\end{equation}
These equalities involve fixed-comparator Gaussian anchors only, and hold regardless of the adaptivity or rank of $G$.

\begin{lemma}[Rank slack]\label{comm:lem:slack}
If $S=M+F$, $\rank M\le h$, and $P$ projects onto leading $r$ left singular vectors of $S$, with $h\le r<\min(m,q)$, then
\begin{equation}\label{comm:eq:slack}
 \op{(\Id-P)S}^2\le\frac{\F F^2}{r-h+1}.
\end{equation}
\end{lemma}
\begin{proof}
The squared singular values numbered $h+1$ through $r+1$ are all at least $\sigma_{r+1}(S)^2$. Their sum is at most the squared best rank-$h$ approximation error, which is at most $\F F^2$.
\end{proof}

\begin{theorem}[Pathwise comparator loss, with no hidden rank exception]\label{comm:thm:defect}
Let $P$ be the actual leading-$r$ projector of $S$ in \eqref{comm:eq:sample}, with $r\ge a+b$. On the anchor event \eqref{comm:eq:ann},
\begin{equation}\label{comm:eq:defect}
 \boxed{(\Id-P)C_c=(\Id-P)ST-(\Id-P)FT-(\Id-P)\widetilde D.}
\end{equation}
Consequently,
\begin{equation}\label{comm:eq:defectbound}
 \begin{split}
 \F{(\Id-P)C_c}\le{}&
 \sqrt{\frac{a}{r-a-b+1}}\,\F F\,\op T\\
 &+\F{(\Id-P)FT}+\F{(\Id-P)\widetilde D}.
 \end{split}
\end{equation}
For $r=4k$ and $a,b\le k$, the coefficient squared in the first term is at most $1/2$. If $G$ has full row rank, $\widetilde D=D$. If also $\ran C\subseteq\ran W$, this term vanishes. If $a=0$, the row comparator loss is identically zero and no right inverse is needed.
\end{theorem}
\begin{proof}
Multiply \eqref{comm:eq:sample} by $T$ and use \eqref{comm:eq:ann}. Then
$ST=Q_GC_c+D+FT=C_c+\widetilde D+FT$.
This proves \eqref{comm:eq:defect}. The matrix $(\Id-P)ST$ has rank at most $a$, so its Frobenius norm is at most $\sqrt a\,\op{(\Id-P)S}\op T$. Apply Lemma~\ref{comm:lem:slack} with $h=a+b$ and the triangle inequality. The case $r\ge m$ has zero loss directly.
\end{proof}

\paragraph{What has disappeared and what has not.}
No recycled coefficient discrepancy enters $F$: it is the lift of the original fixed residual commutator. The matrices $F,T,W,G,P$ are nevertheless dependent. The fixed-anchor inverse moment
\begin{equation}\label{comm:eq:TF}
 \E\F T^2=\frac{a}{q-a-b-1},\qquad q>a+b+1,
\end{equation}
therefore cannot be multiplied by a marginal moment of $F$. Equation \eqref{comm:eq:TF} is the usual inverse-Wishart Frobenius identity, applied conditionally on $X$ \cite{HMT}. Singular designs are not discarded by the analysis: they contribute explicitly through $\widetilde D$.

\section{An unconditional first-layer approximation bound}
This section closes the joint Gaussian calculation when $I$ is a physical coarse leaf, so $W=\Id_m$ and $G=\Psi_I$. Write
\[
 E_o=E[I,I^c],\qquad E_i=E[I^c,I],\qquad
 \nu=\frac{m}{q-m-1},\quad\tau=\frac{a}{q-a-b-1}.
\]
Assume $q>m+1$ and $q>a+b+1$. Then $G$ has full row rank almost surely, $\widetilde D=0$, and
\begin{equation}\label{comm:eq:Fleaf}
 F=E_o\Omega_{I^c}-(\Psi_I^\dagger)^T\Psi_{I^c}^T E_i\Omega_I.
\end{equation}

\begin{lemma}[A joint product moment, not a decoupled bound]\label{comm:lem:FTmoment}
For $a>0$,
\begin{equation}\label{comm:eq:FTmoment}
 \boxed{\E\F{FT}^2
 =\F{E_oR}^2+\tau\F{E_o(\Id-RR^T)}^2
 +\tau\nu\F{E_i(\Id-VV^T)}^2.}
\end{equation}
\end{lemma}
\begin{proof}
Decompose the forward probes into their components along $R,V$ and their orthogonal complements. Since $\Gamma T=\Id_a$ and $XT=0$,
\[
 E_o\Omega_{I^c}T=E_oR+E_o(\Id-RR^T)\Omega_{I^c}T,
 \quad
 E_i\Omega_IT=E_i(\Id-VV^T)\Omega_IT.
\]
The Gaussian complement coordinates are independent of $\Gamma,X$, hence of $T$. Their conditional means are zero and their conditional squared Frobenius moments equal the relevant residual energy times $\F T^2$. Also, for every fixed matrix $M$, averaging the independent transpose probes gives
\[
 \E_\Psi\F{(\Psi_I^\dagger)^T\Psi_{I^c}^TM}^2=\nu\F M^2.
\]
The cross term with the first term of \eqref{comm:eq:Fleaf} vanishes by centering $\Psi_{I^c}$ conditional on $\Psi_I,\Omega$. Use \eqref{comm:eq:TF} to finish.
\end{proof}

Let $\omega=\E\op{H_0^\dagger}^2$, where $H_0$ is an $a\times(q-b)$ standard Gaussian matrix. Thus $\omega=\E\op T^2$.
\begin{lemma}[The shared-anchor weighted energy]\label{comm:lem:weighted}
For $a(q-b)>2$,
\begin{equation}\label{comm:eq:weighted}
 \begin{split}
 \E[\F F^2\op T^2]
 &=\omega\left[q\bigl(\F{E_o}^2+\nu\F{E_i}^2\bigr)
                   -\frac2a\F{E_oR}^2\right]\\
 &\le q\omega\bigl(\F{E_o}^2+\nu\F{E_i}^2\bigr).
 \end{split}
\end{equation}
\end{lemma}
\begin{proof}
After averaging $\Psi$, \eqref{comm:eq:Fleaf} has conditional squared norm
$\F{E_o\Omega_{I^c}}^2+\nu\F{E_i\Omega_I}^2$.
Conditional on $X$, split $\Gamma$ into its $b$ row-space coordinates and its $d=q-b$ null-space coordinates. The latter matrix $H_0$ is standard Gaussian and independent of $X$; the former and all orthogonal physical probe coordinates are independent of $H_0$. Left rotational invariance implies, for fixed $M$,
\[
 \E[\F{MH_0}^2\op{H_0^\dagger}^2]
 =\frac{\F M^2}{a}\E[\F{H_0}^2\op{H_0^\dagger}^2].
\]
Write $H_0=\mathcal R U$, with radius $\mathcal R^2\sim\chi^2_{ad}$ independent of its direction $U$. Homogeneity and $\E\mathcal R^{-2}=1/(ad-2)$ give
\[
 \E[\F{H_0}^2\op{H_0^\dagger}^2]=(ad-2)\omega.
\]
Thus the $E_oR$ term contributes $(q-2/a)\omega\F{E_oR}^2$, and its orthogonal complement contributes $q\omega\F{E_o(\Id-RR^T)}^2$. The conditional distribution of $\op T$ given $X$ is constant, so the entire $E_i$ term contributes $q\omega\F{E_i}^2$. This proves the exact expression.
\end{proof}

\paragraph{External Gaussian input and its explicit integration.}
For an $a\times(a+p)$ Gaussian matrix with $p\ge4$, the spectral pseudoinverse tail in \cite[Sec.~10.1, Eq.~(71) in the arXiv HTML rendering]{HMT} is
\[
 \Pr\left\{\op{H_0^\dagger}\ge\frac{e\sqrt{a+p}}{p+1}t\right\}
 \le t^{-(p+1)},\qquad t\ge1.
\]
Integrating $2t$ times this tail above its threshold gives
\begin{equation}\label{comm:eq:omega}
 \omega\le\frac{e^2(a+p)}{p^2-1},\qquad p=q-a-b.
\end{equation}
The one-row case can alternatively be obtained directly from a chi-square inverse moment. In particular, for $m=8k$, $q=16k$, and $a,b\le k$,
\begin{equation}\label{comm:eq:constants}
 \nu\le\frac87,\qquad\tau\le1,\qquad q\omega<10,
 \qquad\frac{a}{4k-a-b+1}\le\frac12.
\end{equation}
For the spectral estimate, $e^2<7.4$ and
$7.4\cdot256k^2/(196k^2-1)<10$ suffice.

\begin{theorem}[The first layer is residual-relative]\label{comm:thm:leaf}
At the coarse physical leaves, use $m=8k$, $q=16k$, and $r=4k$. For every fixed comparator with exterior ranks at most $k$, the row projector satisfies
\begin{equation}\label{comm:eq:leafrow}
 \E\F{(\Id-P)C}^2
 \le(1+\sqrt5)^2\bigl(\F{E_o}^2+\nu\F{E_i}^2\bigr).
\end{equation}
Summing row and column comparator losses over that entire leaf partition gives
\begin{equation}\label{comm:eq:leafglobal}
 \boxed{\sum_{I\text{ coarse leaf}}\E\left[
 \F{(\Id-P_{U,I})B[I,I^c]}^2+
 \F{(\Id-P_{V,I})B[I^c,I]^T}^2\right]
 \le45\F{\operatorname{off}_{\rm leaf}(E)}^2.}
\end{equation}
There is no singular-value gap assumption on $B$ and no conditioning on a favorable event.
\end{theorem}
\begin{proof}
At leaves $\widetilde D=0$. The first term in \eqref{comm:eq:defectbound} has $L^2$ norm at most $\sqrt5(\F{E_o}^2+\nu\F{E_i}^2)^{1/2}$ by Lemma~\ref{comm:lem:weighted}. The second has $L^2$ norm at most $(\F{E_o}^2+\nu\F{E_i}^2)^{1/2}$ by Lemma~\ref{comm:lem:FTmoment} and contraction by $\Id-P$. Minkowski proves \eqref{comm:eq:leafrow}; $a=0$ is trivial. Transpose the argument for columns. Each exterior-energy sum over a partition equals the squared off-partition residual norm. The resulting constant is
\[
 2(1+\nu)(1+\sqrt5)^2\le\frac{30}{7}(6+2\sqrt5)<45.
\]
\end{proof}

The transpose/row dependence between distinct leaves is irrelevant to this sum of expectations. In contrast, conditioning on already learned child bases at later levels does \emph{not} preserve the Gaussian independence used in Lemmas~\ref{comm:lem:FTmoment}--\ref{comm:lem:weighted}. That extension is not asserted.

\section{Exact assembly into the remaining basis-learning condition}
For every learned node and orientation, use the quantities in Theorem~\ref{comm:thm:defect}. For the column orientation apply the same definitions to $A^T$, exchanging $\Omega,\Psi$ and $a,b$. Let $e_{U,I}$ and $e_{V,I}$ denote the physical local comparator losses.

\begin{proposition}[Orthogonal multilevel assembly]\label{comm:prop:assembly}
For any nested bases, with fixed-space coefficient space $\SSS$,
\begin{equation}\label{comm:eq:assembly}
 \dist_F(B,\SSS)^2\le\sum_I(e_{U,I}^2+e_{V,I}^2),
 \quad e_{U,I}=\F{(\Id-P_{U,I})W_{U,I}^TB[I,I^c]}.
\end{equation}
No additional depth factor is introduced by this inequality.
\end{proposition}
\begin{proof}
At one compressed partition, set the diagonal discrepancy equal to the diagonal blocks minus their retained-retained projections. The remaining off-diagonal error splits orthogonally as the discarded-row sector and the retained-row/discarded-column sector. Their squared norms are bounded by the full row and column exterior losses. The retained-retained sector is isometrically passed to the next level. Thus the discrepancy sectors emitted at successive levels and the final root coefficient are mutually orthogonal. Summing the Pythagorean identities gives the multilevel error of this particular coefficient approximation. At a given node, the opposite-side compression of its exterior block is an isometry contraction, so dropping that compression increases the bounding local loss to the expression in \eqref{comm:eq:assembly}. The best coefficient approximation is no worse.
\end{proof}

Retain three separate, dependent-error ledgers:
\begin{align}
 \mathfrak A&=\sum_{I,o}\frac{a_{I,o}}{r-a_{I,o}-b_{I,o}+1}
             \E[\F{F_{I,o}}^2\op{T_{I,o}}^2],\label{comm:eq:Aledger}\\
 \mathfrak B&=\sum_{I,o}\E\F{(\Id-P_{I,o})F_{I,o}T_{I,o}}^2,\label{comm:eq:Bledger}\\
 \mathfrak D&=\sum_{I,o}\E\F{(\Id-P_{I,o})\widetilde D_{I,o}}^2.\label{comm:eq:Dledger}
\end{align}
Here $o$ ranges over the two orientations, and a zero comparator rank contributes zero. Applying Minkowski in the direct sum of the node/output spaces and then in probability yields
\begin{equation}\label{comm:eq:three}
 \boxed{\bigl(\E\dist_F(B,\SSS)^2\bigr)^{1/2}
 \le\sqrt{\mathfrak A}+\sqrt{\mathfrak B}+\sqrt{\mathfrak D}.}
\end{equation}
This statement is valid with extended nonnegative expectations; it does not presume integrability of the learner's inverses.

\paragraph{A sufficient next theorem.}
For the \emph{specified original-commutator learner}, a bound
\begin{equation}\label{comm:eq:remaining}
 \mathfrak A+\mathfrak B+\mathfrak D\le C_\ell L\OPT_k(A)
\end{equation}
would complete the missing learning step. Equation \eqref{comm:eq:totalJ} pays the unlifted residual data at all levels, and Theorem~\ref{comm:thm:leaf} pays the first physical layer. The unproved step is \eqref{comm:eq:remaining} for later adaptive layers, including singular-design losses. Bounds on separate marginal inverse moments do not prove it.

The three-ledger sufficient condition may itself be stronger than necessary. The signed identity \eqref{comm:eq:defect} is the exact object; if cancellation between $FT$ and $\widetilde D$ is essential, one should retain it instead of forcing separate positive majorants. Failure of one proposed ledger estimate would not refute the learner or the general target.

\section{Conditional closure of the query target}
The following implication identifies the exact benefit of a successful learning bound, without assuming that bound has been proved.

\begin{proposition}[Query ledger and final-rank reduction]\label{comm:prop:conditional}
Suppose \eqref{comm:eq:remaining} holds for every fixed input, with a best comparator $B\in\HH_k$, and the independent refit interface \eqref{comm:eq:refit} is available. Then the original target follows using at most $64k$ queries, with an absolute error constant depending only on $C_\ell$. The final offline rank-$k$ projection need not be computationally efficient.
\end{proposition}
\begin{proof}
Equation \eqref{comm:eq:three} gives $\E\dist_F(B,\SSS)^2\le3C_\ell L\OPT_k(A)$. Since $\dist_F(A,\SSS)\le\F{A-B}+\dist_F(B,\SSS)$,
\[
 \E\dist_F(A,\SSS)^2\le(1+\sqrt{3C_\ell})^2L\OPT_k(A).
\]
The $32k$ training queries followed by $8r=32k$ independent refit queries produce a known matrix $C\in\HH_{4k}$ with
\[
 \E\F{A-C}^2\le\alpha L\OPT_k(A),\qquad
 \alpha=1607(1+\sqrt{3C_\ell})^2.
\]
The class $\HH_k$ is closed, since it is specified by finitely many matrix rank inequalities, and contains zero. An exact Frobenius-best rank-$k$ HSS approximation $\widehat B$ to the known $C$ therefore exists: a minimizing sequence can be restricted to a bounded ball. Its construction requires no further access to $A$. For a best comparator $B$,
\[
 \F{A-\widehat B}\le\F{A-C}+\F{C-\widehat B}
 \le2\F{A-C}+\F{A-B}.
\]
Minkowski proves
\[
 \E\F{A-\widehat B}^2\le(2\sqrt\alpha+1)^2L\OPT_k(A).
\]
This is the query-only bicriteria reduction, not a polynomial-time algorithm for the final nonlinear projection.
\end{proof}

\section{Checks, interpretation, and next mathematical task}
The accompanying script uses NumPy/SciPy only; it does not query external services. It checks the commutator identity, the fixed-anchor annihilator, the exact local lift and defect identities, and the rank-slack estimate in 24 full-design local configurations, including adaptive templates that do and do not contain the comparator signal. The largest normalized identity residual in those local checks is below $3.5\times10^{-15}$; exact-capture errors are below $1.1\times10^{-14}$ in absolute Frobenius norm. Six additional deliberately singular-design tests verify the $\widetilde D$ correction, with normalized residuals below $6.3\times10^{-17}$.

Eight full-tree runs cover $N=64,128,256$ and $k=1,2$, with both exact synthetic HSS inputs and perturbed inputs. On the four exact runs the ideal coefficient-space distance squared is below $1.9\times10^{-27}$. These are finite consistency checks, not a theorem about the design distribution, inverse moments, or worst-case expected risk. Two separate 1,600-draw Monte Carlo diagnostics test the first-leaf moment formulas; their sampled means are not used in the proofs.

\paragraph{What this chapter establishes.}
There is now an all-depth, fixed-residual sketch-energy budget; a cross-sample rule that escapes the specific adaptive one-sided erasure; a concrete fixed-query basis learner; an all-trajectory local physical-error identity; and an unconditional first-layer approximation theorem. The first-layer proof explicitly handles a shared Gaussian inverse and its residual numerator jointly, rather than multiplying their expectations.

\paragraph{What remains.}
The next mathematical task is a rule-specific, all-depth estimate for the adaptive lifting operation in \eqref{comm:eq:learner}. It must pay both the original residual lift and the return/invisibility term of \eqref{comm:eq:defs}. A general theorem uniform over arbitrary adaptive templates is neither needed nor claimed. The standard first-layer Gaussian conditioning cannot simply be iterated after the spaces have been learned from the same probes. There is also no all-depth inverse-integrability claim hidden in the pseudoinverse definition.

The present derivation uses exact bilinear Gaussian second moments, comparator-annihilating right inverses, and deterministic singular-value inequalities. It does not yet constitute a Wick--Hermite or graded-moment proof of adaptive learner stability; the pseudoinverse and selected singular spaces are the nonlinear interface that still needs control.


\section{A paired local refit that controls both discarded observations}
This optional algebraic repair is independent of the original-commutator learner. It explains how the cross-sample viewpoint also removes the old one-sided mismatch from a local residual formula.

Use the full-row-rank local notation of Section~3. Let $P,Q$ be any row and column projectors, set $\bar B=(B_Y+B_Z)/2$, and choose
\[
 D_{\rm pair}=\bar B-P\bar BQ.
\]
Then exactly
\begin{align}
 (\Id-P)(Y_c-D_{\rm pair}\Omega_c)
  &=(\Id-P)(E_Y+\tfrac12\Delta\Omega_c),\label{comm:eq:pairedU}\\
 (\Id-Q)(Z_c-D_{\rm pair}^T\Psi_c)
  &=(\Id-Q)(E_Z-\tfrac12\Delta^T\Psi_c).\label{comm:eq:pairedV}
\end{align}
Because $E_Y\Omega_c^T=0$ and $E_Z\Psi_c^T=0$, the Gram matrix of each right side is bounded above by the Gram matrix of the corresponding discarded cross-sample $(\Id-P)S_U$ or $(\Id-Q)S_V$. Both the Frobenius and operator norms therefore contract. If $P,Q$ are leading-$r$ cross-sample projectors, the rank-slack lemma applies to both discarded observations, not just one.

This does not provide a generic low-rank interpolant: cross-samples can have full local row rank on noisy inputs, unlike one-sided nullified samples whose column counts are reduced. It also does not prove that the paired refit has a stable all-depth feedback law. The original-data learner avoids introducing that feedback in the learning stage in the first place.

\endgroup

% ===== ACTUAL =====
\begingroup
\let\E\relax
\newcommand{\E}{\mathbb E}
\let\PP\relax
\newcommand{\PP}{\mathbb P}
\let\R\relax
\newcommand{\R}{\mathbb R}
\let\F\relax
\newcommand{\F}[1]{\lVert #1\rVert_F}
\let\op\relax
\newcommand{\op}[1]{\lVert #1\rVert_{\rm op}}
\let\ran\relax
\newcommand{\ran}{\operatorname{range}}
\let\rank\relax
\newcommand{\rank}{\operatorname{rank}}
\let\diag\relax
\newcommand{\diag}{\operatorname{blockdiag}}
\let\dist\relax
\newcommand{\dist}{\operatorname{dist}}
\let\OPT\relax
\newcommand{\OPT}{\operatorname{OPT}}
\let\HH\relax
\newcommand{\HH}{\mathcal H}
\let\SSS\relax
\newcommand{\SSS}{\mathcal S}
\let\Id\relax
\newcommand{\Id}{\mathrm I}
\let\tr\relax
\newcommand{\tr}{\operatorname{tr}}
\chapter{Actual design collapse and annihilate-first learning}\label{ch:actual}
\begin{scopebox}
\textbf{Update, 20 September 2026.} The actual inverse-growth theorem below concerns its specified canonical un-nullified rule; exact representation on that family remains valid. The later full-range inverse-optimal learner is a different rule and has both a separate full-inverse obstruction and a genuine physical-error counterexample. The annihilate-first basis failure proved in Chapter~\ref{ch:anchor} is also unchanged. These results are separated by algorithm and conclusion in Chapter~\ref{upd:chap:later-obstructions}; none alone is an impossibility theorem for the synchronized candidate.
\end{scopebox}

\begin{scopebox}
The annihilate-first learner in this chapter has the exact-input, first-layer, and sketch-tail guarantees proved below. Its prospective all-input basis-quality bound is false: Chapter~\ref{ch:anchor} gives a fixed-input counterexample for this actual rule. The local theorems are not retracted.
\end{scopebox}
\thispagestyle{plain}
\begin{chaptersynopsis}
For the original-data, un-nullified commutator learner, we prove a quantitative obstruction involving its actual canonical selection rule, rather than arbitrary adaptive spaces. On a deterministic rank-one family, its unused row spaces repeatedly learn an incoming sketch direction. A rotationally invariant two-child calculation contracts the squared tangent of the sketch-space angle by a factor strictly smaller than one half. Consequently the smallest singular value of the actual compressed design tends to zero in probability, and its inverse second moment grows at least exponentially with depth. This does not refute an approximation theorem: the true outgoing signal on that subtree is zero. It does rule out a depth-uniform unweighted inverse-moment argument.

We then specify an annihilate-first learner with retained rank $2k$ and width $16k$. It uses ordinary nullified SVDs at physical leaves and truncation in sketch coordinates at internal nodes. It has an all-depth, all-rank-strata exact-input theorem; standard Gaussian marginal design laws on exact inputs; a first-layer comparator-loss constant $32$; and an all-input expected bound $6q^2L\OPT_k(A)$ on its discarded sketch data after the actual adaptive truncations. A physical-basis approximation theorem on arbitrary approximate inputs is not proved. The independent fixed-space refit is not reopened.
\end{chaptersynopsis}


\section{Target and separation of claims}
Fix a balanced binary coordinate hierarchy with original leaves of size $2k$. Write $\HH_k$ for the matrices whose outgoing and incoming exterior blocks at every nonroot node have ranks at most $k$. Set
\[
 \OPT_k(A)=\min_{B\in\HH_k}\F{A-B}^2.
\]
The general target is an algorithm using at most $C_q k$ products with $A$ and $A^T$, returning $\widehat B\in\HH_k$ with
\begin{equation}\label{actual:eq:target}
 \E\F{A-\widehat B}^2\le C L\OPT_k(A).
\end{equation}
Here $L$ is at least the number of nonroot coarse partitions used by the learner. All input matrices are fixed before the probes. Constants must not depend on input conditioning, matrix dimension, or depth.

\paragraph{Fixed-space refitting interface.}
The independent fixed-space refit is used only in the final conditional implication: for a known nested rank-$r$ coefficient space $\SSS$ independent of the refit probes, it uses $8r$ queries and returns $C\in\SSS$ with
\begin{equation}\label{actual:eq:refit}
 \E[\F{A-C}^2\mid\SSS]\le1607\dist_F(A,\SSS)^2.
\end{equation}
This is an explicitly stated input, used here through the complete proof in Chapters~\ref{ch:interfaces} and~\ref{ch:uniform}.

\paragraph{External background and attribution.}
Gaussian rangefinding, pseudoinverse moments, and Gaussian operator-norm tails are classical inputs; we use the formulations in Halko--Martinsson--Tropp \cite{HMT}. The fresh-level HSS analysis and the distinction between fresh sketches and reuse are discussed by Amsel et al. \cite{AmselHSS}. The alignment calculation and the modified learner below are derived here. No literature-priority or optimal-constant claim is made.

\section{The original commutator learner}
Draw independent standard Gaussian $\Omega,\Psi\in\R^{N\times q}$, and obtain
\[
 Y=A\Omega,\qquad Z=A^T\Psi.
\]
For a coordinate node $I$, write $P_I$ for its coordinate projector and define
\begin{equation}\label{actual:eq:J}
 J_I=\Psi_I^TY_I-Z_I^T\Omega_I
     =\Psi^T[P_I,A]\Omega.
\end{equation}
The parent commutator is the sum of its two child commutators. The root commutator is zero. Thus the original two products determine all these matrices without more queries.

A row template $W_I$ is the block-diagonal physical row basis of the children; at a physical coarse leaf it is the identity. The un-nullified row learner forms
\begin{equation}\label{actual:eq:old}
 G_I=W_I^T\Psi_I,\qquad S_I=(G_I^\dagger)^T J_I,
\end{equation}
retains a leading rank-$r$ left space, and lifts it by $W_I$. On samples of rank at most $r$, use a fixed canonical extension depending only on the physical column space: project coordinate vectors in order, orthogonalize, skip zero vectors, and complete to dimension $r$. This removes arbitrary dependence on rotations of nonzero singular vectors or zero singular vectors. The previous concrete parameters were $r=4k$, $q=16k$, coarse leaf size $m=8k=2r$.

\section{An actual-rule obstruction: learning a direction at a zero cut}
\subsection{A fixed rank-one family}
Let $J$ be one root half and choose a coordinate $j$ in the other half. Let $v$ be a deterministic unit vector supported on $J$, with every coarse-leaf restriction nonzero; a constant vector on $J$ suffices. Put
\begin{equation}\label{actual:eq:oneway}
 B=e_jv^T.
\end{equation}
This matrix has global rank one and hence belongs to $\HH_1$ for every hierarchy. It is nonsymmetric. For every node $I\subseteq J$,
\[
 B[I,I^c]=0,\qquad B[I^c,I]^T=v_Ie_j^T.
\]
Write $\zeta=\Psi^Te_j$, and condition on its nonzero value. Its direction
$z=\zeta/\|\zeta\|$ is independent of all physical probe rows inside $J$. For each such node,
\begin{equation}\label{actual:eq:phantom}
 J_I=-\zeta(v_I^T\Omega_I),\qquad
 \ran S_I=\operatorname{span}\{(G_I^\dagger)^Tz\}
\end{equation}
whenever this vector is nonzero. The outgoing row space is zero, but the learner still retains this incoming-regression direction. Its row choices within $J$ depend on the local $\Psi$ rows and $z$, not on $\Omega$: the latter contributes only a nonzero row factor to a rank-one sample.

Call the retained physical row basis $\mathcal U_I$ and its sketch image
\[
 \mathscr U_I=\ran(\Psi_I^T\mathcal U_I)\subset\R^q.
\]
For any subspace $S$ having a nonzero projection of $z$, define
\begin{equation}\label{actual:eq:tangent}
 t(S)=\frac{\|(\Id-\Pi_S)z\|^2}{\|\Pi_Sz\|^2}
     =\tan^2\angle(z,S).
\end{equation}
This is an angle to a subspace, not to an arbitrarily chosen basis vector.

\subsection{A two-child geometric calculation}
\begin{lemma}[Independent axial subspaces]\label{actual:lem:axial}
Let $r>1$, $q>2r$, and let $S_1,S_2$ be independent random $r$-dimensional subspaces, conditionally on a fixed unit vector $z\in\R^q$. Suppose their laws are invariant under every orthogonal transformation fixing $z$, and almost surely $0<\|\Pi_{S_i}z\|<1$. Let $S$ be any subspace of $S_1+S_2$ that contains $\Pi_{S_1+S_2}z$. Then
\begin{equation}\label{actual:eq:axial}
 \E t(S)\le\frac{q-2r+1}{4(q-r)}
                   \bigl(\E t(S_1)+\E t(S_2)\bigr).
\end{equation}
Both child subspaces are in direct sum almost surely, and their sum does not contain $z$.
\end{lemma}
\begin{proof}
The projections of the two subspaces onto $z^\perp$ have dimension $r$ and, by axial invariance, are independent uniform $r$-subspaces of the $(q-1)$-dimensional space $z^\perp$. Since $2r\le q-1$, they are in direct sum almost surely. This implies both claims about $S_1+S_2$.

Put
\[
 D_i=S_i\cap z^\perp,\qquad
 a_i=\frac{\Pi_{S_i}z}{\|\Pi_{S_i}z\|^2}=z+\eta_i.
\]
Then $\dim D_i=r-1$, $\eta_i\perp z,D_i$, and
$\|\eta_i\|^2=t(S_i)$. Conditional on $D_i$ and $t(S_i)$, axial invariance makes $\eta_i$ uniform on the corresponding sphere in
$D_i^\perp\cap z^\perp$, a space of dimension $q-r$. In particular its conditional mean is zero. The two pairs remain independent.

Let $D=D_1+D_2$. The vector
\[
 w=\frac{a_1+a_2}{2}-\Pi_D\frac{\eta_1+\eta_2}{2}
   =z+(\Id-\Pi_D)\frac{\eta_1+\eta_2}{2}
\]
belongs to $S_1+S_2$ and satisfies $z^Tw=1$. For any subspace $T$,
\[
 t(T)=\min_{w\in T,\ z^Tw=1}\|w-z\|^2.
\]
Also $t(S)=t(S_1+S_2)$ because $S$ contains the projection of $z$ onto the larger space. Thus
\[
 t(S)\le\tfrac14\|(\Id-\Pi_D)(\eta_1+\eta_2)\|^2.
\]
Conditional on $D_1,D_2$ and the two squared lengths, the mixed term has mean zero. For the first square, the projection of the independent uniform $(r-1)$-space $D_2$ onto $D_1^\perp\cap z^\perp$ is a uniform $(r-1)$-space there. Consequently
\[
 \E\|(\Id-\Pi_D)\eta_1\|^2
 =\left(1-\frac{r-1}{q-r}\right)\E\|\eta_1\|^2.
\]
The same holds with the indices exchanged. This proves \eqref{actual:eq:axial}. Truncation and monotone convergence cover extended expectations; in the application below the leaf moment is finite.
\end{proof}

\begin{theorem}[Actual phantom-space alignment]\label{actual:thm:alignment}
Apply the un-nullified learner \eqref{actual:eq:old} with fixed $r>1$, width $q>2r$, and physical coarse leaves of size $m=2r$ to the fixed input \eqref{actual:eq:oneway}. Define
\begin{equation}\label{actual:eq:parameters}
 \mu=\frac{q-m}{m-2},\qquad
 \beta=\frac{q-2r+1}{q-r}<1,\qquad \alpha=\frac\beta2.
\end{equation}
At every finite depth, every row design in the subtree $J$ has full row rank almost surely. A node at height $h$ above the coarse leaves satisfies
\begin{equation}\label{actual:eq:angle-bound}
 \boxed{\E\,t(\mathscr U_I)\le\mu\alpha^h.}
\end{equation}
The expectation may also be taken conditionally on $\zeta$; the bound is unchanged.
\end{theorem}
\begin{proof}
At a leaf, the physical design is an $m\times q$ standard Gaussian matrix and has full row rank almost surely. Its image is a uniform $m$-subspace. The retained image contains the projection of $z$ onto that image by \eqref{actual:eq:phantom}, since
\[
 G^T(G^\dagger)^Tz=\Pi_{\ran G^T}z.
\]
It is contained in the same image, so the angle is unchanged by retention. The squared projection length has beta law
$\operatorname{Beta}(m/2,(q-m)/2)$. Equivalently, its squared tangent is
$X/Y$ for independent $X\sim\chi^2_{q-m}$ and $Y\sim\chi^2_m$. Hence
$\E t=(q-m)/(m-2)=\mu$.

Conditional on $z$, disjoint subtrees use independent physical probe rows. A right orthogonal transformation fixing $z$ leaves the physical vector $(G^\dagger)^Tz$, and thus its canonical physical completion, unchanged, while rotating its sketch image. Induction therefore gives precisely the axial invariance and subtree independence of Lemma~\ref{actual:lem:axial}.

At a parent the design image is the sum of the two retained child images. The lemma proves it has dimension $2r$ and does not contain $z$. Thus the design is full row rank, the sample vector is nonzero, and the retained image contains the projection of $z$ onto the sum. Apply \eqref{actual:eq:axial}; the two equal-height children have the same distribution. This gives a factor $\beta/2$ at each merge and proves \eqref{actual:eq:angle-bound}.
\end{proof}

\subsection{The actual design inverse cannot have a uniform moment}
\begin{theorem}[Design collapse and exponential inverse growth]\label{actual:thm:inverse}
In Theorem~\ref{actual:thm:alignment}, consider the design $G^{(h+1)}$ at a parent whose two children each have height $h$ and physical size $n_h=m2^h$. If $n_h\ge\max\{q,8\}$, then
\begin{equation}\label{actual:eq:median-collapse}
 \PP\left\{\sigma_{\min}(G^{(h+1)})^2
                     \le36m\mu\beta^h\right\}\ge\tfrac12,
\end{equation}
and therefore
\begin{equation}\label{actual:eq:inverse-growth}
 \boxed{\E\op{(G^{(h+1)})^\dagger}^2
             \ge\frac{\beta^{-h}}{72m\mu}.}
\end{equation}
In fact $\sigma_{\min}(G^{(h+1)})\to0$ in probability as the family depth grows.
For the previous parameters $r=4$, $m=8$, $q=16$,
\begin{equation}\label{actual:eq:kone}
 \E\op{(G^{(h+1)})^\dagger}^2\ge\frac1{768}
          \left(\frac43\right)^h,\qquad h\ge1.
\end{equation}
These are statements about the learner's actual designs on deterministic, nonsymmetric rank-one inputs.
\end{theorem}
\begin{proof}
Normalize $z$ and use the two child vectors $a_i=z+\eta_i$ from the preceding proof. Write $M_i=\Psi_i^T\mathcal U_i$, so
$G^{(h+1)T}=[M_1,M_2]$. Since $a_i\in\ran M_i$, let $M_ix_i=a_i$. Then
\[
 \|x_i\|\ge\frac1{\op{\Psi_i}},\qquad
 G^{(h+1)T}\begin{bmatrix}x_1\\-x_2\end{bmatrix}
       =\eta_1-\eta_2.
\]
Putting $M_* = \max_i\op{\Psi_i}$ gives
\begin{equation}\label{actual:eq:sigma-pointwise}
 \sigma_{\min}(G^{(h+1)})^2
       \le\frac{M_*^2}{2}\|\eta_1-\eta_2\|^2.
\end{equation}
Conditional on $z$, independence and centering give
\[
 \E\|\eta_1-\eta_2\|^2\le2\mu\alpha^h.
\]
Markov's inequality puts this quantity below $8\mu\alpha^h$ with probability at least $3/4$.

Each physical $\Psi_i$ is an $n_h\times q$ standard Gaussian matrix. The usual Gaussian norm estimate \cite{HMT} gives
\[
 \PP\{\op{\Psi_i}>\sqrt{n_h}+\sqrt q+t\}\le e^{-t^2/2}.
\]
With $t=\sqrt{n_h}$ and $n_h\ge q$, both norms are at most $3\sqrt{n_h}$ except on an event of probability at most $2e^{-n_h/2}\le1/4$. No independence between these norm events and the angles is used. On their intersection with the Markov event, \eqref{actual:eq:sigma-pointwise} is at most
$36n_h\mu\alpha^h=36m\mu\beta^h$. This proves \eqref{actual:eq:median-collapse} and \eqref{actual:eq:inverse-growth}.

More generally, for every $\varepsilon>0$, the same argument gives
\[
 \PP\{\sigma_{\min}(G^{(h+1)})^2>\varepsilon\}
 \le2e^{-n_h/2}+\frac{9m\mu}{\varepsilon}\beta^h\longrightarrow0.
\]
The specialization uses $\mu=4/3$ and $\beta=3/4$.
\end{proof}

\paragraph{What is and is not ruled out.}
A fixed larger width ratio only makes $\beta$ closer to one: it does not restore a depth-uniform unweighted inverse moment. This is stronger than an example selecting arbitrary adversarial nested spaces. It concerns the actual canonical commutator learner.

It is \emph{not} an approximation lower bound. On the whole subtree used in the proof, the true outgoing comparator is zero. An error term weighted by that comparator can vanish even while the inverse moment diverges. The example therefore mandates a signal-weighted, cancellation-preserving analysis or a learner modification. It does not change the independent fixed-space refit theorem, and it does not disprove \eqref{actual:eq:target}. In fact $B$ still belongs to the learned coefficient space on this family: the nonzero outgoing ranges in the source half are the deterministic restrictions of $e_j$, and the nonzero incoming ranges in $J$ are the restrictions of $v$. Each of those two orientations has a pure rank-one physical sample and a deterministic canonical comparison trajectory. The other two orientations, including the collapsing row designs in $J$, are unused by the true matrix.

\section{An annihilate-first, sketch-metric learner}
We now change the learner explicitly. The original commutators are still used; no coefficient estimates are recycled into them.

Choose
\begin{equation}\label{actual:eq:newparams}
 r=2k,\qquad q=16k,\qquad m=4k=2r.
\end{equation}
Group pairs of original leaves into coarse leaves of size $4k$. For $N\le16k$, exact coordinate queries are an allowed fallback. Otherwise all local template sizes are $m=4k$. As before,
\[
 W_{U,I}=\diag(\mathcal U_a,\mathcal U_b),\quad
 W_{V,I}=\diag(\mathcal V_a,\mathcal V_b),\quad
 G_I=W_{U,I}^T\Psi_I,\quad H_I=W_{V,I}^T\Omega_I.
\]
At physical leaves both templates are identities. Put
\[
 \Pi_G=G^\dagger G,\quad\Pi_H=H^\dagger H,
 \qquad N_G=\Id-\Pi_G,\quad N_H=\Id-\Pi_H.
\]
These are projectors on the $q$-dimensional sketch coordinate space, including on singular designs.

\paragraph{Physical leaves.}
Choose ordinary leading rank-$r$ left spaces from
\begin{equation}\label{actual:eq:leafalgo}
 X_U=Y_I N_H=A[I,I^c]\Omega_{I^c}N_H,\qquad
 X_V=Z_I N_G=A[I^c,I]^T\Psi_{I^c}N_G.
\end{equation}
Use canonical physical completion whenever the sample rank is at most $r$. The row leaf rule depends only on $\Omega$ and the column rule only on $\Psi$.

\paragraph{Internal nodes.}
Form the inverse-free projected matrices
\begin{equation}\label{actual:eq:metricdata}
 M_U=\Pi_G J_I N_H,\qquad M_V=-\Pi_HJ_I^T N_G.
\end{equation}
If $\rank M_U>r$, take a leading $r$-dimensional left space $Z_U\subset\ran G^T$, and choose the physical local row space
\begin{equation}\label{actual:eq:lift}
 \ran U_I=\ran\bigl((G^T)^\dagger Z_U\bigr).
\end{equation}
If $\rank M_U\le r$, canonically extend the full physical range
$\ran((G^T)^\dagger M_U)$ to dimension $r$. Use the analogous column rule. Fix measurable tie rules and canonical frames of the resulting physical spaces. Lift by $W_{U,I},W_{V,I}$ as usual.

For full-row-rank $G$, write $G^T=QR$ with $Q$ column-orthonormal. The internal rule truncates $Q^T J_I N_H$ \emph{before} applying $R^{-1}$ to the selected vectors. It differs from truncating the physically lifted matrix $R^{-1}Q^T J_I N_H$. It is an SVD in the sketch metric, not a claim that the physical metric and sketch metric are uniformly equivalent.

The Moore--Penrose formulation and the full-range branch make the procedure well-defined on singular designs as well. No singular-design probability assertion is required for the all-input data-fit theorem below.

\paragraph{Cost and rank.}
Training uses exactly the two width-$16k$ original streams, hence $32k$ matvec queries. Parent commutators add, and retained designs propagate by $U_I^TG_I,V_I^TH_I$. A local-transfer implementation uses $O(Nk^2)$ arithmetic. The resulting coarse-tree coefficient space has rank at most $2k$. A proper original descendant of a coarse leaf has size at most $2k$, so the space also lies in the original-tree rank-$2k$ class. These access and representation claims do not assert \eqref{actual:eq:target}.

\section{The exact-input endpoint is now unconditional}
\begin{theorem}[All-depth exact containing spaces]\label{actual:thm:exact}
For every fixed $B\in\HH_k$, the annihilate-first learner, applied to $A=B$, produces a coefficient space $\SSS$ containing $B$ almost surely. Its physical bases coincide with a deterministic canonical hierarchy depending only on $B$, including at zero cuts and rank-deficient cuts. Every reached row and column design has full row rank almost surely, and each marginal design is standard Gaussian. In particular, at a size-$4k$ node,
\begin{equation}\label{actual:eq:exact-inverse}
 \E\F{G_I^\dagger}^2=\E\F{H_I^\dagger}^2
                  =\frac{4k}{12k-1}.
\end{equation}
No independence between designs at different nodes or depths is claimed.
\end{theorem}
\begin{proof}
Construct a deterministic comparison hierarchy directly from $B$ for the proof only. At each node canonically extend its true outgoing and incoming exterior ranges, in the coordinates of its comparison child templates, to dimension $r=2k$. Each true range has dimension at most $k$. Restrictions of parent exterior columns to a child belong to that child's exterior range, so induction proves containment by these templates. Thus all comparison physical frames are deterministic isometries.

At any fixed comparison node factor
\[
 B[I,I^c]=CR^T,\qquad B[I^c,I]^T=VK,
\]
with the exterior factors $R,V$ column-orthonormal. Write
$\Gamma=R^T\Omega_{I^c}$, $X=V^T\Omega_I$, and
$\Theta=K\Psi_{I^c}$. Then
\[
 J_I(B)=\Psi_I^TC\Gamma-\Theta^TX.
\]
Containment by the comparison column template gives $XN_H=0$. Row containment gives $\Psi_I^TC=G^TC_c$, where $C_c=W_U^TC$. Hence
\begin{equation}\label{actual:eq:exact-annihilation}
 J_I(B)N_H=G^TC_c\Gamma N_H.
\end{equation}
All comparison designs are full-row-rank Gaussian matrices almost surely. Conditional on the physical probe rows inside $I$, $\Gamma$ is an independent standard Gaussian matrix. A basis of $\ker H$ has $q-m=12k$ columns, so $\Gamma N_H$ has full row rank almost surely for every exterior rank at most $k$. Consequently the physical lift in \eqref{actual:eq:exact-annihilation} has range exactly $\ran C_c$. The same argument applies to columns. Physical leaves satisfy the same assertion directly from \eqref{actual:eq:leafalgo}.

There are finitely many comparison nodes. Intersect these probability-one events, and induct from the leaves: the actual physical full-range branch returns precisely the deterministic comparison spaces at each node. This proves the asserted actual trajectory and containment of $B$. The Gaussian inverse-Frobenius identity \cite{HMT} gives \eqref{actual:eq:exact-inverse}.
\end{proof}

\begin{remark}[Removal of the specific phantom mechanism]
On \eqref{actual:eq:oneway}, every outgoing row space inside $J$ is zero. Equation \eqref{actual:eq:exact-annihilation} makes every actual row sample there zero, and the canonical zero-space completion is deterministic. The learner no longer repeatedly retains $(G_I^\dagger)^Tz$. In particular the exponential inverse phenomenon of Theorem~\ref{actual:thm:inverse} is absent on this family. This does not extend \eqref{actual:eq:exact-inverse} to arbitrary perturbed inputs.
\end{remark}

\section{All-depth control after the actual adaptive truncations}
Fix any $B\in\HH_k$ and put $E=A-B$. Direct Gaussian second moments give
\begin{equation}\label{actual:eq:Jbudget}
 \E\F{J_I(E)}^2=q^2\left(\F{E[I,I^c]}^2+\F{E[I^c,I]}^2\right),\qquad
 \sum_I\E\F{J_I(E)}^2\le2q^2 L\F E^2.
\end{equation}
Also $\rank J_I(B)\le2k$, since the two off-diagonal blocks of $[P_I,B]$ each have rank at most $k$.

Let $\Pi_{GU}$ be the projector onto $\ran(G^TU_I)$ and $\Pi_{HV}$ the analogous column projector. Define actual discarded sketch matrices
\begin{equation}\label{actual:eq:residuals}
 \mathcal R_{U,I}=(\Pi_G-\Pi_{GU})J_I(A)N_H,
 \qquad
 \mathcal R_{V,I}=(\Pi_H-\Pi_{HV})J_I(A)^TN_G.
\end{equation}
\begin{theorem}[Projected data-fit budget]\label{actual:thm:datafit}
At every internal node, on every trajectory, including singular designs,
\begin{equation}\label{actual:eq:local-tail}
 \F{\mathcal R_{U,I}}^2\le\F{J_I(E)}^2,
 \qquad
 \F{\mathcal R_{V,I}}^2\le\F{J_I(E)}^2.
\end{equation}
For the whole hybrid learner, including its physical leaves,
\begin{equation}\label{actual:eq:global-tail}
 \boxed{\sum_I\E\left[\F{\mathcal R_{U,I}}^2+
                            \F{\mathcal R_{V,I}}^2\right]
           \le6q^2L\F E^2.}
\end{equation}
In particular a best comparator gives $6q^2L\OPT_k(A)$. This is a theorem about the actual selected sketch spaces, not only the unprocessed commutators.
\end{theorem}
\begin{proof}
At an internal node, the matrix $\Pi_GJ_I(B)N_H$ has rank at most $2k=r$ regardless of any dependence of the projectors on the probes. The leading-$r$ choice in sketch coordinates therefore has squared tail at most
$\F{\Pi_GJ_I(E)N_H}^2\le\F{J_I(E)}^2$.
The image of the lifted physical space contains the selected sketch space. In the full-range branch it contains the entire projected sample range and the tail is zero. This proves \eqref{actual:eq:local-tail}, also when $G$ is singular. Repeat for columns.

At a physical leaf, almost surely
$J_I(A)N_H=\Psi_I^TX_U$. The image $\ran(G^TU_I)$ contains $\Psi_I^TU_IU_I^TX_U$, whence
\[
 \F{\mathcal R_{U,I}}
 \le\op{\Psi_I}\F{(\Id-U_IU_I^T)X_U}
 \le\op{\Psi_I}\F{E[I,I^c]\Omega_{I^c}N_H}.
\]
The last inequality is the best-rank approximation property, using the rank-at-most-$k$ comparator sample. All other factors in this bound depend only on $\Omega$, so averaging $\Psi_I$ separately is legitimate. Put $d=q-m=12k$. Conditional on $\Omega_I$, the exterior projected Gaussian matrix has $d$ effective columns. Thus
\[
 \E\F{\mathcal R_{U,I}}^2
 \le d\E\op{\Psi_I}^2\F{E[I,I^c]}^2.
\]
The Gaussian upper-tail bound gives
\[
 (\E\op{\Psi_I}^2)^{1/2}\le\sqrt q+\sqrt m+\sqrt2,
 \qquad d\E\op{\Psi_I}^2<660k^2<3q^2.
\]
The transpose argument is identical. Exterior squared energies sum to the off-leaf residual energy in each orientation, so the leaf layer costs at most $6q^2\F E^2$. Each subsequent partition costs at most $4q^2\F E^2$ by \eqref{actual:eq:Jbudget} and \eqref{actual:eq:local-tail}. The total is at most
$(6+4(L-1))q^2\F E^2\le6q^2L\F E^2$.
\end{proof}

\section{A first-layer physical error bound}
\begin{theorem}[First physical layer]\label{actual:thm:leaf}
For the parameters \eqref{actual:eq:newparams}, every fixed comparator $B$ with exterior ranks at most $k$ satisfies
\begin{equation}\label{actual:eq:leaf-bound}
 \boxed{\sum_{I\text{ coarse leaf}}\E\left[
 \F{(\Id-U_IU_I^T)B[I,I^c]}^2+
 \F{(\Id-V_IV_I^T)B[I^c,I]^T}^2\right]
 \le32\F{\operatorname{off}_{\rm leaf}(A-B)}^2.}
\end{equation}
No singular-value gap or favorable-event assumption is used.
\end{theorem}
\begin{proof}
Consider one outgoing block $A_o=B_o+E_o$. An orthonormal basis of $\ker\Omega_I$ reduces \eqref{actual:eq:leafalgo} to
$X=A_o\Xi$, where $\Xi$ is an exterior-by-$d$ standard Gaussian matrix and $d=12k$. This has the same left Gram matrix as the projected $q$-column sample. Factor $B_o=CR^T$, with $R$ orthonormal and $a=\rank B_o\le k$. The case $a=0$ is trivial. Put
\[
 \Gamma=R^T\Xi,\qquad T=\Gamma^\dagger,\qquad F=E_o\Xi.
\]
Then $XT=C+FT$. If $P$ is the retained leading-$2k$ projector, the singular-value variational inequality gives
\[
 \F{(\Id-P)C}
 \le\sqrt{\frac{a}{2k-a+1}}\F F\op T+\F{FT}.
\]
Indeed $(\Id-P)XT$ has rank at most $a$, and
$\sigma_{2k+1}(X)^2\le\F F^2/(2k-a+1)$.

The Gaussian coordinates orthogonal to $R$ are independent of $\Gamma$, so
\begin{equation}\label{actual:eq:joint1}
 \E\F{FT}^2=\F{E_oR}^2+
       \frac{a}{d-a-1}\F{E_o(\Id-RR^T)}^2\le\F{E_o}^2.
\end{equation}
The other moment must also be taken jointly. Write
$\omega=\E\op{\Gamma^\dagger}^2$. Rotational invariance in the $a$ rows and a Gaussian radial decomposition imply
\begin{equation}\label{actual:eq:joint2}
 \E[\F F^2\op T^2]
 =\omega\left[d\F{E_o}^2-\frac2a\F{E_oR}^2\right]
 \le d\omega\F{E_o}^2.
\end{equation}
For completeness, if $\Gamma=\mathcal R U$, then
$\mathcal R^2\sim\chi^2_{ad}$ is independent of $U$, and homogeneity gives
$\E[\F\Gamma^2\op{\Gamma^\dagger}^2]=(ad-2)\omega$.
Row rotational invariance gives the $-2/a$ term. The orthogonal physical coordinates contribute $d\omega$ times their squared residual norm.

For $p=d-a$, the standard spectral pseudoinverse tail \cite{HMT}, integrated against $2t\,dt$, gives
\[
 \omega\le\frac{e^2d}{p^2-1},\qquad
 d\omega\le\frac{7.4\cdot144k^2}{121k^2-1}<9.
\]
For $a=1$, the inverse chi-square moment gives the same spectral bound directly. Also $a/(2k-a+1)<1$. Minkowski's inequality in $L^2$ now bounds the local comparator loss by
$(3+1)^2\F{E_o}^2=16\F{E_o}^2$. Apply the transpose version, then sum over the partition. Each orientation counts the off-leaf energy once, proving \eqref{actual:eq:leaf-bound}.
\end{proof}

\section{Why the remaining implication is still substantive}
The new results do not turn the all-depth sketch-data budget into a physical-basis budget. Right annihilation can erase a direction when the learned opposite template is correlated with an exterior Gaussian anchor. For example, with full-row-rank designs, take
\[
 J=G^Tc\,g,\qquad g\in\operatorname{row}(H).
\]
Then both $J N_H$ and $J^T N_G$ vanish even though $c\ne0$. Thus both projected learning samples and both tails in \eqref{actual:eq:residuals} can vanish while a true physical exterior direction is absent from a canonically completed space. Appropriate adaptive templates can realize this algebraic obstruction; it is not a claim that the specified learner chooses those templates.

The exact-input theorem defeats that issue on its actual deterministic comparison trajectory. On approximate inputs, the bases become genuinely sketch-dependent and can change across rank strata. Exact recovery, Gaussian marginal designs on exact inputs, and finite-dimensional numerical continuity do not furnish a uniform residual-relative bound in those neighborhoods.

A sufficient remaining theorem is
\begin{equation}\label{actual:eq:remaining}
 \E\dist_F(A,\SSS_{\rm learned})^2
                 \le C_bL\OPT_k(A)
\end{equation}
for the specified learner, or an equivalent bound on the sum of its local physical comparator losses. The fixed-space orthogonal multilevel assembly introduces no additional depth factor after such a sum is bounded. A proof must exploit the actual selection rule and handle loss inherited from both child orientations. It cannot follow from the projected-tail theorem alone, nor from an unweighted inverse bound for the old learner.

\begin{proposition}[Conditional query consequence]\label{actual:prop:queries}
If \eqref{actual:eq:remaining} is proved, the general target follows with at most $48k$ queries in the query-only model, using the inherited refit \eqref{actual:eq:refit} and exact offline projection onto $\HH_k$.
\end{proposition}
\begin{proof}
Training costs $32k$. The independent rank-$2k$ refit costs $16k$, giving a known $C\in\HH_{2k}$ with
$\E\F{A-C}^2\le1607C_bL\OPT_k(A)$.
The original class $\HH_k$ is closed, since it is defined by finitely many rank inequalities, and contains zero. Hence an exact best approximation $\widehat B$ to the known $C$ exists and requires no additional queries. For a best comparator $B$ to $A$,
\[
 \F{A-\widehat B}\le\F{A-C}+\F{C-\widehat B}
                  \le2\F{A-C}+\F{A-B}.
\]
Minkowski gives \eqref{actual:eq:target} with constant
$(2\sqrt{1607C_b}+1)^2$. No efficient algorithm for the final nonlinear optimization is asserted.
\end{proof}

\section{Numerical diagnostics and their exact interpretation}
\subsection{Actual coefficient-space distance, not a surrogate inverse bound}
For fixed learned transfers and a rank-one comparator $B=uv^T$, its exact coefficient-space distance can be computed without a dense $N\times N$ matrix. At one compressed partition, write its row and column factors as $u_i,v_i$, and let their retained factors be $u_i^+,v_i^+$. Set
\[
 x_i=\|u_i\|^2,\quad y_i=\|v_i\|^2,\quad
 d_i=\|u_i-U_iu_i^+\|^2,\quad e_i=\|v_i-V_iv_i^+\|^2,
 \quad x_i^+=\|u_i^+\|^2.
\]
The orthogonal off-diagonal loss at that level is exactly
\begin{equation}\label{actual:eq:rankone-distance}
 \sum_i d_i\sum_{j\ne i}y_j+
 \sum_i x_i^+\sum_{j\ne i}e_j.
\end{equation}
The first term is the discarded-row sector, the second the retained-row/discarded-column sector. All same-block entries are absorbed by the diagonal coefficient sector. The retained-retained sector passes isometrically to the next level. Consequently the sum of \eqref{actual:eq:rankone-distance} over levels equals $\dist_F(B,\SSS)^2$. The diagnostic uses this positive formula, avoiding subtraction of nearly equal unit-size quantities.

For each $N$, fix independent unit vectors $u,v$ and a permutation matrix $P$ using input seed $20260919$, before any probe draw. Use
\[
 B=uv^T,\qquad A=B+\epsilon E,\qquad
 E=P/\sqrt N,\quad\epsilon=10^{-4}.
\]
Then $B\in\HH_1$, $\F E=1$, and $\OPT_1(A)\le\epsilon^2$. Three paired probe seeds, $1101,1102,1103$, give the following medians of
$\dist_F(B,\SSS(A))^2/\epsilon^2$:
\begin{center}
\begin{tabular}{rrrr}
\toprule
$N$ & Old, $r=4$ & Old, $r=2$ & Annihilate-first hybrid, $r=2$\\
\midrule
512   & 2.816   & 0.490 & 0.416\\
2048  & 11.211  & 0.700 & 0.477\\
8192  & 43.234  & 0.925 & 0.498\\
32768 & 212.979 & 0.933 & 0.446\\
\bottomrule
\end{tabular}
\end{center}
All three columns use width $q=16$ per stream. Coarse leaf size is $2r$, so rank and leaf-size changes are explicit. The middle column isolates much of the effect of reducing retained rank; the improvement must not be attributed solely to the metric modification.

For each individual run, not merely in expectation,
\begin{equation}\label{actual:eq:actual-lower}
 \dist_F(A,\SSS)^2\ge
 \left(\dist_F(B,\SSS)-\epsilon\right)_+^2.
\end{equation}
Thus a large reported value is a genuine representational obstruction for that realized basis: refitting inside the space cannot remove it. Three paired seeds do not estimate worst-case expected risk, and growth over four dimensions is not an asymptotic impossibility proof. The rank-one family in the analytic inverse theorem is different from this perturbed diagnostic family.

\subsection{Consistency checks delivered in the computational archive}
Twenty small exact-input cases cover full HSS inputs, one-way rank-one matrices, rank-one inputs below the allowed rank, leaf-diagonal matrices, and zero matrices, with $k=1,2$. The largest normalized squared projection error was below $1.6\cdot10^{-27}$. Four additional checks verify idempotence and self-adjointness of the fixed-space projection used to measure the diagnostic error. Twenty-four local projected-tail checks include deliberately singular row and column designs; the largest observed ratio to the permitted residual energy was below $0.102$.

The phantom-space checker uses 32 trees of height eight, with $r=4,q=16$. The measured mean squared tangent decreased from approximately $1.308$ at the leaves to $7.35\cdot10^{-5}$ at height eight. The proved bound there is $\mu\alpha^8=(4/3)(3/8)^8$, about $5.21\cdot10^{-4}$. Nodes within a partition are separate samples conditional on the common direction, but levels are not independent. These checks support the algebraic implementation; no numerical value is used in the proofs.

\section{Limits of this candidate}
There is now a proved, quantitative obstruction to a depth-uniform \emph{unweighted} inverse analysis of the actual un-nullified commutator learner. It persists for every fixed oversampling ratio with $r>1$ and $q>2r$. The old general approximation claim remains unproved, not refuted.

The annihilate-first hybrid is a precisely defined alternative with the same $O(k)$ training access. Its all-depth exact-input theorem covers zero and deficient cut ranks and gives genuinely Gaussian marginal designs on that locus. Its all-depth projected-tail bound and first-layer physical bound hold for every fixed input. Their conversion into \eqref{actual:eq:remaining} is impossible for this exact rule: Chapter~\ref{ch:anchor} disproves that all-input basis-distance bound. The exact-input, projected-tail, and first-layer results remain valid, and the independent fixed-space refit remains a completed component.

\endgroup

% ===== PORTS =====
\begingroup
\let\E\relax
\newcommand{\E}{\mathbb E}
\let\R\relax
\newcommand{\R}{\mathbb R}
\let\Id\relax
\newcommand{\Id}{\mathrm I}
\let\ran\relax
\newcommand{\ran}{\operatorname{range}}
\let\rank\relax
\newcommand{\rank}{\operatorname{rank}}
\let\tr\relax
\newcommand{\tr}{\operatorname{tr}}
\let\dist\relax
\newcommand{\dist}{\operatorname{dist}}
\let\OPT\relax
\newcommand{\OPT}{\operatorname{OPT}_k}
\let\HH\relax
\newcommand{\HH}{\mathcal H}
\let\SSS\relax
\newcommand{\SSS}{\mathcal S}
\let\F\relax
\newcommand{\F}[1]{\left\|#1\right\|_{\mathrm F}}
\let\op\relax
\newcommand{\op}[1]{\left\|#1\right\|_{\mathrm{op}}}
\let\ip\relax
\newcommand{\ip}[2]{\left\langle #1,#2\right\rangle}
\let\transpose\relax
\newcommand{\transpose}{\mathsf T}
\chapter{Port geometry and exact-intersection recovery}\label{ch:ports}
\begin{scopebox}
The exact-intersection branch is an exact-real recovery statement on its recognized class. It is not a perturbation-stability theorem for the maintained noisy union.
\end{scopebox}
\begin{chaptersynopsis}
The all-input, $O(k)$-query, expected $O(L)$ best-HSS approximation target is not established here. We isolate a new obstruction for the specified commutator basis learner itself, not for arbitrarily chosen adaptive spaces. On a fixed symmetric HSS-rank-$k$ family with both exterior ranks saturated at every nonroot node, the expected squared operator norm of its compressed-design pseudoinverse is bounded below by a quantity growing geometrically with height. Nevertheless, the learner represents this exact family without basis loss. The proof supplies a positive all-depth result: a port graph-angle energy contracts by a factor at most $3/8$ per merge at retained rank $4k$ and width $16k$. Thus full inverse control is the wrong intermediate objective for this route.

We also give a query-neutral global rank test and an exact-intersection branch. For every fixed input whose two exterior ranks sum to at most $4k$ at each node, this branch learns deterministic, exactly containing nested bases almost surely. It includes all HSS-rank-$2k$ inputs and handles rank-deficient cuts. Combining with the independent fixed-space refit recovers such inputs using at most $64k$ queries. This exact-real result does not prove perturbation stability. For the remaining noisy branch, we retain the actual discarded sample tail and combine the two signed error terms before squaring.
\end{chaptersynopsis}

\section{Target, refitting interface, and scope}
Fix a binary coordinate tree. The HSS-rank-at-most-$k$ class $\HH_k$ consists of matrices whose two exterior blocks at every nonroot node have ranks at most $k$. Let $L\ge1$ count nonroot partitions and write
\[
 \OPT(A)=\min_{B\in\HH_k}\F{A-B}^2.
\]
The target is a rank-$k$ HSS output satisfying
\begin{equation}\label{ports:eq:target}
 \E\F{A-\widehat B}^2\le C L\OPT(A)
 \quad\text{using at most }C_q k\text{ products with }A,A^\transpose.
\end{equation}
The external fresh-level benchmark in \cite{AmselHSS} uses a logarithmic number of batches; its sketch-reuse question motivates this chapter. This chapter does not claim an exhaustive literature or priority assessment.

\paragraph{Inherited fixed-space theorem, proved in Chapter~\ref{ch:uniform}.}
For nested bases of rank $r$ fixed independently of the refit randomness, the sibling-orthogonal refit is linear, exact on the resulting coefficient space $\SSS$, and obeys
\begin{equation}\label{ports:eq:inherited}
 \E[\F{A-C}^2\mid\SSS]\le1607\dist_F(A,\SSS)^2
 \quad\text{using }8r\text{ queries}.
\end{equation}
The complete proof of this interface is in Chapters~\ref{ch:interfaces} and~\ref{ch:uniform}. The port-geometry and exact-intersection arguments below do not use it; only the final query consequence uses \eqref{ports:eq:inherited}.

A known intermediate approximation may have rank $r=O(k)$. Exact offline projection onto $\HH_k$ uses no further queries. Indeed, if $\widehat B$ minimizes $\F{C-\widehat B}$ over $\HH_k$, then
\begin{equation}\label{ports:eq:offline}
 \F{A-\widehat B}\le2\F{A-C}+\sqrt{\OPT(A)}.
\end{equation}
The minimum exists because the class is closed, contains zero, and a minimizing sequence can be restricted to a bounded ball. This is a query-only statement, not a polynomial-time claim.

\section{Original observations and the two-sided local sample}
Draw independent standard Gaussian $\Omega,\Psi\in\R^{N\times q}$ and query
\[
 Y=A\Omega,\qquad Z=A^\transpose\Psi.
\]
For a node $I$ with coordinate projector $P_I$, define
\begin{equation}\label{ports:eq:J}
 J_I=\Psi_I^\transpose Y_I-Z_I^\transpose\Omega_I
     =\Psi^\transpose[P_I,A]\Omega.
\end{equation}
Internal diagonal blocks cancel. For children $I_1,I_2$,
$J_{I_1\cup I_2}=J_{I_1}+J_{I_2}$.

If $A=B+E$ for fixed $B\in\HH_k$, then
\begin{equation}\label{ports:eq:unlifted}
 \rank J_I(B)\le2k,\qquad
 \sum_I\E\F{J_I(E)}^2\le2q^2L\F E^2.
\end{equation}
For completeness, $[P_I,E]$ has blocks $E[I,I^c]$ and $-E[I^c,I]$ on disjoint supports. For every fixed $M$, independent Gaussian second moments give
$\E\F{\Psi^\transpose M\Omega}^2=q^2\F M^2$.
At each partition the row exterior energies sum to at most $\F E^2$, and the column energies have the same bound. This proves \eqref{ports:eq:unlifted}. It is an unlifted sketch-energy estimate.

Let $W_I$ be an orthonormal row template formed from the two retained child bases. The original-data row learner computes
\begin{equation}\label{ports:eq:learner}
 G_I=W_I^\transpose\Psi_I,\qquad
 S_I=(G_I^\dagger)^\transpose J_I(A)
\end{equation}
and retains a leading-$r$ left singular space of $S_I$. The column learner exchanges the streams and uses $-J_I^\transpose$. Use a deterministic measurable SVD/tie/completion rule depending only on the local sample and the fixed coordinates. If $\rank S_I\le r$, the rule must retain its full range. Coarse leaves have size $2r$; the parameters of interest are
\begin{equation}\label{ports:eq:parameters}
 r=4k,\qquad q=16k.
\end{equation}
Small matrices may instead be recovered by coordinate queries. Coarsening two levels of an original size-$2k$-leaf balanced hierarchy gives the size-$8k$ leaves used here; the final offline projection can use the original hierarchy.

\section{Graph geometry relative to a fixed query-space port}
Let $H_0\in\R^{q\times b}$ have orthonormal columns. We call its range a \emph{port}: it is a distinguished subspace of the query-coordinate space. This is a definition, not an extra query. Let $\mathscr S\subset\R^q$ be an $r$-plane with
\[
 0<K_{\mathscr S}:=H_0^\transpose P_{\mathscr S}H_0<\Id_b.
\]
Define
\begin{equation}\label{ports:eq:graph}
 X_{\mathscr S}=P_{\mathscr S}H_0K_{\mathscr S}^{-1}
               =H_0+Z_{\mathscr S},\qquad
 t(\mathscr S)=\F{Z_{\mathscr S}}^2
             =\tr K_{\mathscr S}^{-1}-b.
\end{equation}
Here $H_0^\transpose Z_{\mathscr S}=0$. The matrix $X_{\mathscr S}$ is the minimum-Frobenius-norm matrix with columns in $\mathscr S$ satisfying $H_0^\transpose X=\Id_b$. The space
\[
 \mathscr T_{\mathscr S}=\mathscr S\cap\ran(H_0)^\perp
\]
has dimension $r-b$, and $Z_{\mathscr S}$ is orthogonal to it. Thus $t$ is a trace of squared tangent-angle parameters, not a physical comparator error.

\begin{lemma}[Two-child graph contraction]\label{ports:lem:graph}
Assume $r>b$, $q\ge2r+b$. Let $\mathscr S_1,\mathscr S_2$ be independent random $r$-planes, with distributions invariant under every orthogonal transformation fixing $H_0$ pointwise. Assume their port matrices lie strictly between zero and identity and $\E t(\mathscr S_i)<\infty$. Put $\mathscr M=\mathscr S_1+\mathscr S_2$ and
\[
 \lambda=\frac{q-2r+b}{q-r}.
\]
Then $\dim\mathscr M=2r$ and $\mathscr M\cap\ran H_0=\{0\}$ almost surely, and
\begin{equation}\label{ports:eq:contraction}
 \boxed{\E t(\mathscr M)\le\frac\lambda4
      \bigl(\E t(\mathscr S_1)+\E t(\mathscr S_2)\bigr).}
\end{equation}
\end{lemma}
\begin{proof}
Write $Z_i=Z_{\mathscr S_i}$ and $\mathscr T_i=\mathscr T_{\mathscr S_i}$. The nuisance planes $\mathscr T_i$ are independent uniform $(r-b)$-planes in the $(q-b)$-dimensional complement of $H_0$. Conditional on $\mathscr T_i$, $Z_i$ is centered: reflect the complement of $\ran H_0+\mathscr T_i$, leaving $\mathscr T_i$ and $H_0$ fixed. Conditional distributions are invariant under this stabilizer, so $Z_i$ and $-Z_i$ have the same conditional law.

Let $Q$ be orthogonal projection, within $\ran(H_0)^\perp$, onto $(\mathscr T_1+\mathscr T_2)^\perp$. Subtracting a vector in $\mathscr T_1+\mathscr T_2$ shows that both $H_0+QZ_i$ belong columnwise to $\mathscr M$. Therefore
\begin{equation}\label{ports:eq:candidate}
 t(\mathscr M)\le\frac14\F{QZ_1+QZ_2}^2.
\end{equation}
Conditionally on the two nuisance planes, the cross term has expectation zero by independence and centering. In $\mathscr T_1^\perp\cap\ran(H_0)^\perp$, which has dimension $q-r$, the projection of $\mathscr T_2$ is a uniform $(r-b)$-plane. Consequently
\[
 \E\F{QZ_1}^2=\left(1-\frac{r-b}{q-r}\right)\E\F{Z_1}^2
 =\lambda\E t(\mathscr S_1).
\]
The same statement holds for the second child, proving \eqref{ports:eq:contraction}.

We justify the genericity assertions used above and at the next level. One may introduce independent Haar rotations fixing $H_0$ as auxiliary proof randomness without changing either law. For any fixed pair of nondegenerate graph spaces, choose their two nuisance planes and their two $b$-dimensional graph ranges mutually orthogonal inside $\ran(H_0)^\perp$. This is possible because these four dimensions sum to $2r\le q-b$. This configuration has zero intersection between the two $r$-planes, and their sum has zero intersection with $\ran H_0$. Failure of either property is the vanishing of appropriate minors in the rotation entries. Since those minors are not identically zero on the relevant rotation orbits, the failure set has Haar measure zero. Applying the auxiliary rotations proves the assertions for the original independent invariant laws. The port projection of the sum is positive definite since either child already has that property.
\end{proof}

\begin{lemma}[Initial graph moment]\label{ports:lem:initial}
If $\mathscr M$ is a uniform $m$-plane in $\R^q$, $m>b+1$ and $m+b\le q$, then
\begin{equation}\label{ports:eq:t0}
 \E t(\mathscr M)=\frac{b(q-m)}{m-b-1}.
\end{equation}
\end{lemma}
\begin{proof}
By rotational invariance the principal-angle matrix has the same law as the Gram of the first $m$ columns of a Haar $b$-row orthonormal frame in $\R^q$. Generate this frame by normalizing a standard Gaussian matrix $[G_1,G_2]$, with $m$ and $q-m$ columns respectively. Put $W_1=G_1G_1^\transpose$, $W_2=G_2G_2^\transpose$. Cyclicity of trace gives
\[
 \tr K^{-1}=b+\tr(W_1^{-1}W_2).
\]
The independent Gaussian Gram identities are $\E W_2=(q-m)\Id_b$ and $\E W_1^{-1}=\Id_b/(m-b-1)$; these are the standard inverse-Wishart identities used in Gaussian rangefinder analysis \cite{HMT}. Subtract $b$.
\end{proof}

\section{An actual-rule example with contracting angles and growing inverses}
Take root halves of size $n=2r\,2^h$. Choose a fixed $n\times k$ orthonormal-column matrix $U$ whose restriction to every coarse leaf has rank $k$. For example, at the parameters \eqref{ports:eq:parameters}, use a normalized repeated identity frame. Define the fixed symmetric matrix
\begin{equation}\label{ports:eq:B}
 B=\begin{pmatrix}0&UU^\transpose\\UU^\transpose&0\end{pmatrix}.
\end{equation}
Its global rank is $2k$, and both exterior ranks at every nonroot tree node are exactly $k$. This is not an inactive-cut or rank-transition example.

Inside the left half, put
\[
 H=\Psi_R^\transpose U,\qquad \Gamma=U^\transpose\Omega_R,
 \qquad X_I=U_I^\transpose\Omega_I.
\]
Then the original observable is
\begin{equation}\label{ports:eq:familyJ}
 J_I(B)=\Psi_I^\transpose U_I\Gamma-HX_I.
\end{equation}
The exterior port $H$ and outgoing anchor $\Gamma$ use only right-half probes. Let $H_0$ be an orthonormal frame for $\ran H$.

\begin{lemma}[Retention, conditional locality, and invariance]\label{ports:lem:actual}
Suppose $r\ge2k$ and $q\ge2r+k$. For the original learner \eqref{ports:eq:learner} applied to \eqref{ports:eq:B}, all finite-depth local designs have full row rank almost surely; all true exterior row and column spaces are retained. If
\[
 \mathscr M_I=\ran G_I^\transpose,\qquad
 \mathscr S_I=\ran(G_I^\transpose U_I^{\rm loc}),
\]
where $U_I^{\rm loc}$ is the retained local row basis, then
\begin{equation}\label{ports:eq:protected}
 P_{\mathscr S_I}H_0=P_{\mathscr M_I}H_0.
\end{equation}
Conditionally on $(H,\Gamma)$, the image spaces from disjoint left-half subtrees are independent and invariant under query-space rotations fixing $H_0$ pointwise.
\end{lemma}
\begin{proof}
The fixed Gaussian anchor matrix $[\Gamma;X_I]$ has full row rank $2k$ almost surely, simultaneously at all nodes of a finite tree. Assume the template contains $U_I$ and $G_I$ has full row rank. Then
\[
 S_I=C_I\Gamma-K_I X_I,\qquad
 C_I=W_I^\transpose U_I,\qquad K_I=(G_I^\dagger)^\transpose H.
\]
Full row rank of the anchor matrix gives
\begin{equation}\label{ports:eq:span}
 \ran S_I=\ran[C_I,K_I].
\end{equation}
This range has dimension at most $2k\le r$ and is retained in full. It contains the comparator factor. Also
$G_I^\transpose K_I=P_{\mathscr M_I}H$, proving \eqref{ports:eq:protected}.

Conditional on $(H,\Gamma)$, every operation in a left-half subtree uses only its own $\Psi,\Omega$ rows and fixed matrix coefficients. Hence disjoint subtrees are conditionally independent. If an orthogonal query-space matrix $Q$ fixes $H$ pointwise, replace all $\Psi$ rows in one subtree by $\Psi Q$. Inductively the physical child bases are unchanged. At its next node,
\[
 G_I\longmapsto G_IQ,\qquad J_I\longmapsto Q^\transpose J_I,
 \qquad (G_I^\dagger)^\transpose J_I\longmapsto
 (G_I^\dagger)^\transpose J_I.
\]
Thus the actual local sample and its deterministic selection are unchanged, while its image space rotates by $Q^\transpose$. Gaussian invariance proves the asserted law; no independence of learned bases and the probes is claimed.

At physical leaves, $G_I=\Psi_I$ has full row rank and its image is a uniform $2r$-plane, with nondegenerate port angles. After retention the same port projection remains, by \eqref{ports:eq:protected}. At a parent, the input image is the sum of the two child images. Apply the genericity part of Lemma~\ref{ports:lem:graph}, conditionally on $(H,\Gamma)$, to obtain dimension $2r$ and nondegenerate port angles. This proves the full-rank and exact-containment induction. The right half and the column orientation follow from the same argument with the streams and halves exchanged.
\end{proof}

\begin{theorem}[All-depth graph-angle contraction]\label{ports:thm:angles}
At a specified node of height $j$ above the coarse leaves, let $t_j=t(\mathscr S_I)$, with $b=k$. Put
\[
 t_0^{\rm av}=\frac{k(q-2r)}{2r-k-1},\qquad
 \lambda=\frac{q-2r+k}{q-r}.
\]
Then
\begin{equation}\label{ports:eq:angleheight}
 \boxed{\E t_j\le t_0^{\rm av}(\lambda/2)^j.}
\end{equation}
In particular, along any prescribed nested chain,
\begin{equation}\label{ports:eq:anglesum}
 \sum_{j=0}^h\E t_j\le\frac{t_0^{\rm av}}{1-\lambda/2}.
\end{equation}
At $r=4k,q=16k$, these become
\begin{equation}\label{ports:eq:explicitangle}
 \E t_j\le\frac{8k^2}{7k-1}(3/8)^j,
 \qquad
 \sum_j\E t_j\le\frac{64k^2}{5(7k-1)}.
\end{equation}
\end{theorem}
\begin{proof}
At leaves, Lemma~\ref{ports:lem:initial} applies to the uniform input image and \eqref{ports:eq:protected} preserves $t$ after retention. At each parent, apply Lemma~\ref{ports:lem:graph} conditionally on $(H,\Gamma)$, then preserve $t$ again under retention. The conditional leaf expectation equals the deterministic value $t_0^{\rm av}$, so induction gives the bound for every node. Summing a geometric series proves the chain estimate. No assertion of independence between heights is used.
\end{proof}

\begin{lemma}[A deterministic inverse lower-bound mechanism]\label{ports:lem:singular}
At a parent with two child image graphs $H_0+Z_a,H_0+Z_b$, set $R_I^2=\F{\Psi_I}^2$. Then
\begin{equation}\label{ports:eq:smin}
 \sigma_{\min}(G_I)^2\le R_I^2(t_a+t_b).
\end{equation}
\end{lemma}
\begin{proof}
For any unit $e\in\R^k$, choose physical retained child vectors $x_a,x_b$ with query images $(H_0+Z_a)e$ and $(H_0+Z_b)e$. Their norms are at least $1/R_I$, since each image has a unit component in $\ran H_0$. The signed stacked vector belongs to the parent template, has squared norm at least $2/R_I^2$, and has query image $(Z_a-Z_b)e$. Hence
\[
 \sigma_{\min}(G_I)^2
 \le\tfrac12R_I^2\op{Z_a-Z_b}^2
 \le R_I^2(\F{Z_a}^2+\F{Z_b}^2).
\]
\end{proof}

\begin{theorem}[No depth-uniform full inverse second moment]\label{ports:thm:inverse}
For the actual learner and fixed inputs \eqref{ports:eq:B}, at a specified node of height $j\ge1$,
\begin{equation}\label{ports:eq:inverselower}
 \boxed{\E\op{G_I^\dagger}^2
 \ge\frac{\lambda^{-(j-1)}}{8rq\,t_0^{\rm av}}.}
\end{equation}
At $r=4k,q=16k$,
\begin{equation}\label{ports:eq:inverseexplicit}
 \boxed{\E\op{G_I^\dagger}^2
 \ge\frac{7k-1}{4096k^4}(4/3)^{j-1}.}
\end{equation}
The expectations may be interpreted as extended nonnegative expectations. This theorem does not claim finite inverse moments at every height.
\end{theorem}
\begin{proof}
Put $u=t_a+t_b$, which is positive almost surely by the nondegenerate graph property. Lemma~\ref{ports:lem:singular} and full row rank yield
\[
 \E\op{G_I^\dagger}^2\ge\E\frac1{R_I^2u}.
\]
Crucially, no independent-factorization step is needed. Cauchy--Schwarz gives
\[
 \E\frac1{R_I^2u}\ge
 \frac{(\E R_I^{-1})^2}{\E u}
 \ge\frac1{\E R_I^2\,\E u}.
\]
The last inequality is Jensen applied to $x^{-1/2}$. The node has $n_j=2r\,2^j$ physical rows, so $\E R_I^2=n_jq$. Theorem~\ref{ports:thm:angles} gives
$\E u\le2t_0^{\rm av}(\lambda/2)^{j-1}$.
Substitution and simplification give \eqref{ports:eq:inverselower}. The explicit values give \eqref{ports:eq:inverseexplicit}.
\end{proof}

\paragraph{What this obstruction rules out.}
For each fixed width ratio with $r>k$, $\lambda<1$. Therefore a height-independent bound on the full inverse second moment cannot hold for this actual rule on all exact inputs. Increasing the fixed width ratio changes the geometric rate but does not remove this obstruction. The theorem concerns a prescribed node as its physical subtree grows, not just the worst node among an increasing number of nodes.

\paragraph{What it does not rule out.}
The same inputs are represented exactly at all depths by Lemma~\ref{ports:lem:actual}. In particular,
\begin{equation}\label{ports:eq:portkilled}
 (\Id-U_I^{\rm loc}(U_I^{\rm loc})^\transpose)
 (G_I^\dagger)^\transpose H=0.
\end{equation}
The large inverse behavior can be confined to directions that the error readout annihilates. Thus Theorem~\ref{ports:thm:inverse} is not a counterexample to the learner's residual-relative risk, to its prior sufficient joint-error conditions, or to \eqref{ports:eq:target}. It specifically excludes an attractive full-inverse intermediate bound.

\section{A query-neutral exact-intersection branch}
This section changes the learner on an exactly recognizable class. It does not silently substitute a new rule in the inverse theorem above.

\begin{lemma}[The original cut sketch recognizes capped total cut rank]\label{ports:lem:rankflag}
For every fixed $A$, almost surely simultaneously at all nonroot nodes,
\begin{equation}\label{ports:eq:rankflag}
 \rank J_I(A)=\min\{q,a_I+b_I\},\qquad
 a_I=\rank A[I,I^c],\quad b_I=\rank A[I^c,I].
\end{equation}
\end{lemma}
\begin{proof}
The fixed commutator $[P_I,A]$ has rank $a_I+b_I$, since its two nonzero blocks have disjoint row and column supports. Right multiplication by a standard Gaussian matrix preserves rank up to the width almost surely. Conditional on this right product, independent left Gaussian multiplication preserves its rank up to the same width. A finite intersection of probability-one events proves the simultaneous statement.
\end{proof}

Define
\[
 \mathcal C_r=\{A:a_I+b_I\le r\text{ at every nonroot node}\}.
\]
Here the nodes are those of the coarse learning tree; the final rank-$k$ projection still uses the original tree. Because $q>r$, the global flag
\begin{equation}\label{ports:eq:flag}
 \max_{I\ne\mathrm{root}}\rank J_I(A)\le r
\end{equation}
recognizes membership in $\mathcal C_r$ almost surely for each fixed input. The probability-one statement is pointwise in $A$; it is not uniform over inputs chosen after seeing the sketches. At $r=4k$, $\HH_{2k}\subseteq\mathcal C_{4k}$.

\paragraph{Algorithm on the global flag.}
At every node compute the subspace intersection
\begin{equation}\label{ports:eq:intersection}
 \mathscr C_I=\ran J_I(A)\cap\ran G_I^\transpose.
\end{equation}
Lift this subspace by $(G_I^\dagger)^\transpose$, and choose a canonical orthonormal extension of its range of dimension $\min\{r,m_I\}$ in the local physical coordinates, where $m_I$ is the template dimension. A specific rule is Gram--Schmidt applied to the coordinate projections onto the lifted subspace, followed by coordinate completion. It is a function of the subspace projector, not of the arbitrary basis used to represent the intersection. The column side applies the same rule to $J_I^\transpose$ and $W_{V,I}^\transpose\Omega_I$.

Use this branch throughout the tree only when the \emph{global} flag holds. On its complement keep the original leading-$r$ commutator learner. The flag and all intersections require no extra matvecs.

\begin{theorem}[All-depth exact recovery of the exterior spaces]\label{ports:thm:exactbranch}
Let coarse leaves have size at most $2r$, use $q\ge3r$, and assume $q>2r+1$ for the inverse-moment statement. For every fixed $A\in\mathcal C_r$, the intersection branch learns exactly containing row and column bases at every node almost surely. With the prescribed canonical completion, these bases equal deterministic comparison bases depending only on $A$ and the tree. Every compressed design is therefore standard Gaussian in a deterministic coordinate system and has full row rank almost surely. If its row count is $m_I\le2r$, then
\begin{equation}\label{ports:eq:branchinverse}
 \E\F{G_I^\dagger}^2=\frac{m_I}{q-m_I-1}
 \le\frac{2r}{q-2r-1}.
\end{equation}
This is a per-node expectation, not a bound for the maximum over all nodes.
\end{theorem}
\begin{proof}
Construct a deterministic comparison tree from the true row and column exterior spaces of $A$, using the same canonical extension rule. These spaces have dimensions at most $r$. Their restrictions to a child lie in that child's exterior space, so the comparison spaces are nested.

At a node of this deterministic comparison tree, factor the fixed exterior blocks as
\[
 A[I,I^c]=CR^\transpose,\qquad A[I^c,I]=DV^\transpose,
\]
where $C,V$ have $a,b$ columns respectively, $a+b\le r$, and $R,D$ have full column rank. Write
\[
 \Gamma=R^\transpose\Omega_{I^c},\quad
 X=V^\transpose\Omega_I,\quad\Theta=\Psi_{I^c}^\transpose D.
\]
Then
\[
 J_I=\Psi_I^\transpose C\Gamma-\Theta X.
\]
The fixed-factor Gaussian anchor stack $[\Gamma;X]$ has full row rank $a+b$ almost surely. Thus
\begin{equation}\label{ports:eq:Jrange}
 \ran J_I=\ran[\Psi_I^\transpose C,\Theta].
\end{equation}
Let $W$ be the deterministic row comparison template and $G=W^\transpose\Psi_I$. Since $W$ contains $C$ and $G$ has full row rank,
$\ran(\Psi_I^\transpose C)\subseteq\ran G^\transpose$.
The random $b$-plane $\ran\Theta$ is independent of $G$. Its projection onto $(\ran G^\transpose)^\perp$ is injective almost surely because
$q-m_I\ge q-2r\ge r\ge b$.
Consequently \eqref{ports:eq:Jrange} implies
\[
 \ran J_I\cap\ran G^\transpose=\ran(\Psi_I^\transpose C).
\]
Its lifted range is exactly $\ran(W^\transpose C)$, because
$(G^\dagger)^\transpose G^\transpose=\Id_{m_I}$.
Canonical extension therefore reproduces the deterministic comparison basis. Starting with identity templates at physical leaves proves that the actual and comparison trajectories agree everywhere, for both orientations. Take a finite intersection of their probability-one events. Zero ranks cause no difficulty: the corresponding intersection is zero and canonical coordinate completion is deterministic.

Each actual design agrees almost surely with the design formed from its deterministic comparison template. Orthogonal invariance makes that design standard Gaussian. The inverse-Wishart Frobenius identity gives \eqref{ports:eq:branchinverse}. The global flag is true with probability one by Lemma~\ref{ports:lem:rankflag}; no conditioning on a positive-probability rank event changes the Gaussian law.
\end{proof}

\begin{corollary}[An unconditional exact-data branch at $64k$ queries]\label{ports:cor:exactquery}
At $r=4k,q=16k$, the branch uses $32k$ training queries. Combined with the independent rank-$4k$ refit interface \eqref{ports:eq:inherited}, an additional $32k$ queries recover every fixed $A\in\mathcal C_{4k}$ exactly almost surely. In particular, every fixed HSS-rank-at-most-$2k$ input is recovered. Exact offline projection of this known matrix onto $\HH_k$ attains $\OPT(A)$, without extra queries, on this class. At each full-size local design,
\[
 \E\F{G_I^\dagger}^2=\frac{8k}{8k-1}.
\]
\end{corollary}
\begin{proof}
The learned coefficient space contains $A$ exactly, and training and refit probes are independent. Its conditional expected refit error is therefore zero by \eqref{ports:eq:inherited}. A nonnegative random squared error with zero expectation is zero almost surely. The offline assertion follows because the known recovered matrix equals $A$.
\end{proof}

\paragraph{Scope and limitations.}
This is an exact-real branch, not a robust numerical rank selector. An arbitrarily small unstructured perturbation can change the global flag. Likewise a literal nodewise intersection is not stable when lower-level templates only approximately contain a true exterior space. The global flag is essential to the proof because it supplies an exactly containing deterministic comparison tree at every level. No all-input perturbation theorem is obtained by continuity from this exact branch. Exact HSS recovery is an established topic; the contribution here is a specified original-commutator branch with a transparent deterministic-comparison proof, not a priority claim for $O(k)$ exact HSS recovery in general.

\section{The remaining noisy estimate: tails and signed cancellation}
We restate the local error identity to make the remaining condition explicit. Fix $A=B+E$ with $B\in\HH_k$. At a row node factor
\[
 B[I,I^c]=CR^\transpose,\qquad B[I^c,I]^\transpose=VH_B,
\]
with orthonormal $R,V$ and ranks $a,b\le k$. Put
\[
 \Gamma=R^\transpose\Omega_{I^c},\quad X=V^\transpose\Omega_I,
 \quad T=N_X(\Gamma N_X)^\dagger,
\]
where $N_X$ is an orthonormal basis for $\ker X$. Almost surely $\Gamma T=\Id_a$ and $XT=0$. For the actual possibly adaptive, possibly singular design, define
\begin{align*}
 G&=W^\transpose\Psi_I,&\mathcal L&=(G^\dagger)^\transpose,&Q_G&=\mathcal L G^\transpose,\\
 C_c&=W^\transpose C,&F&=\mathcal L J_I(E),&
 D&=\mathcal L\Psi_I^\transpose(\Id-WW^\transpose)C,\\
 \widetilde D&=D-(\Id-Q_G)C_c.&
\end{align*}
Let $P$ be the actual leading-$r$ left projector of $S=\mathcal L J_I(A)$. Direct multiplication gives
\begin{equation}\label{ports:eq:signed}
 \boxed{(\Id-P)C_c=(\Id-P)ST-(\Id-P)(FT+\widetilde D).}
\end{equation}
Indeed $ST=C_c+\widetilde D+FT$. No integrability or full-rank design event is needed for this identity.

Define two quantities over the row and column orientations of all nodes in the noisy branch:
\begin{align}
 \mathfrak A_{\rm tail}
 &=\sum_{I,o}a_{I,o}\E\!\left[
 \sigma_{r+1}(S_{I,o})^2\op{T_{I,o}}^2\right],\label{ports:eq:Atail}\\
 \mathfrak C_{\rm joint}
 &=\sum_{I,o}\E\F{(\Id-P_{I,o})
       (F_{I,o}T_{I,o}+\widetilde D_{I,o})}^2.\label{ports:eq:Cjoint}
\end{align}
Zero comparator ranks contribute zero.

\begin{proposition}[Tail-and-signed-source sufficient condition]\label{ports:prop:remaining}
For the nested coefficient space $\SSS$ produced by the noisy branch,
\begin{equation}\label{ports:eq:twoledger}
 \left(\E\dist_F(B,\SSS)^2\right)^{1/2}
 \le\sqrt{\mathfrak A_{\rm tail}}+\sqrt{\mathfrak C_{\rm joint}}.
\end{equation}
It is therefore sufficient to prove, uniformly over fixed inputs outside $\mathcal C_{4k}$,
\begin{equation}\label{ports:eq:missing}
 \boxed{\mathfrak A_{\rm tail}+\mathfrak C_{\rm joint}
        \le C_{\rm learn}L\OPT(A).}
\end{equation}
No proof of \eqref{ports:eq:missing} is supplied here.
\end{proposition}
\begin{proof}
The first term of \eqref{ports:eq:signed} has at most $a$ columns. Since $P$ is a leading-$r$ projector,
\[
 \F{(\Id-P)ST}\le\sqrt a\,\sigma_{r+1}(S)\op T.
\]
The usual orthogonal multilevel coefficient construction gives
\[
 \dist_F(B,\SSS)^2\le\sum_{I,o}\F{(\Id-P_{I,o})C_{c,I,o}}^2.
\]
To see that no depth factor is hidden here, split each compressed exterior block into discarded-row and retained-row/discarded-column sectors. These sectors are orthogonal to one another and to the retained-retained sector passed upwards. The emitted sectors at different heights are therefore orthogonal; diagonal coefficient sectors and the root core are unrestricted. Their squared errors telescope. Removing the opposite-side isometric compression only increases each bounding local exterior loss. Apply Minkowski to \eqref{ports:eq:signed} in this direct-sum Hilbert space and then in probability. This proves \eqref{ports:eq:twoledger}.
\end{proof}

\paragraph{Why this condition is better targeted than full inverse control.}
The comparator part of $S$ has rank at most $a+b$. Consequently, the deterministic singular-value inequality gives
\begin{equation}\label{ports:eq:tailweyl}
 \sigma_{r+1}(S)\le\sigma_{r-a-b+1}(F)
 \le\frac{\F F}{\sqrt{r-a-b+1}}.
\end{equation}
Thus the actual observed tail can be substantially smaller than a full Frobenius noise norm. Likewise \eqref{ports:eq:Cjoint} keeps the signed cancellation between residual lift and comparator return; it does not square the two terms separately. Theorem~\ref{ports:thm:inverse} explains why blindly replacing these quantities by full inverse norms is not viable. It does not prove that \eqref{ports:eq:missing} holds, nor does it prove it fails.

If \eqref{ports:eq:missing} holds, \eqref{ports:eq:twoledger} gives comparator-space loss at most $2C_{\rm learn}L\OPT(A)$. The ordinary triangle inequality adds the fixed residual $E$, the independent refit supplies \eqref{ports:eq:inherited}, and \eqref{ports:eq:offline} supplies the final rank-$k$ output. The exact branch already attains the optimum on $\mathcal C_{4k}$. The total query count remains $64k$.

\section{Checks, findings, and next task}
The companion checker runs the actual row learner on five symmetric bidirectional-root test trees, with $k=1,2$ and root-half sizes from $512$ to $4096$. It checks exact comparator capture, protected-port annihilation, preservation of the port projection, the deterministic graph candidate inequality \eqref{ports:eq:candidate}, and the singular-value bound \eqref{ports:eq:smin}. The largest recorded projection/capture identity residual is below $1.5\cdot10^{-13}$; neither deterministic inequality has a positive numerical violation.

Six full two-sided intersection-branch tests include nested HSS inputs of ranks $k$ and $2k$, and one-way root interactions with zero exterior orientations and arbitrary leaf diagonals. Sizes range from $64$ to $256$. All capped-rank identities match. Three additional generic or perturbed inputs fail the global rank flag, as predicted, with no capped-rank mismatch. The largest relative exterior loss is below $2.1\cdot10^{-14}$, and the maximum difference from the deterministic canonical comparison basis is below $1.2\cdot10^{-12}$. A separate $2500$-draw leaf diagnostic has sample mean $1.3841$ against the exact value $4/3$, with estimated standard error $0.0285$. This is a diagnostic, not a premise of the proof.

The checker uses numerical thresholds for ranks and intersections, whereas the theorems concern exact-real arithmetic. No Monte Carlo output establishes the inverse expectation lower bound, its asymptotics, or an all-input approximation guarantee. 

\paragraph{Status.}
The fixed-space refit remains an inherited closed component. This chapter proves an actual-rule inverse obstruction and a positive all-depth angle theorem on a saturated exact HSS family. It also closes the all-depth exact-input endpoint for an explicitly modified commutator branch. The unresolved part of the general target is still robust basis learning outside the exactly recognizable class, with joint control of the discarded tail and the signed return term. A larger fixed sketch width cannot fix the full-inverse obstruction; a successful proof must exploit the projected or quotient geometry, or a genuinely different noisy learner.

\endgroup

% ===== RETURN =====
\begingroup
\let\E\relax
\newcommand{\E}{\mathbb E}
\let\PP\relax
\newcommand{\PP}{\mathbb P}
\let\R\relax
\newcommand{\R}{\mathbb R}
\let\Id\relax
\newcommand{\Id}{I}
\let\HH\relax
\newcommand{\HH}{\mathcal H}
\let\SSS\relax
\newcommand{\SSS}{\mathcal S}
\let\F\relax
\newcommand{\F}[1]{\left\|#1\right\|_F}
\let\op\relax
\newcommand{\op}[1]{\left\|#1\right\|_{\mathrm{op}}}
\let\ran\relax
\newcommand{\ran}{\operatorname{range}}
\let\rank\relax
\newcommand{\rank}{\operatorname{rank}}
\let\diag\relax
\newcommand{\diag}{\operatorname{diag}}
\let\tr\relax
\newcommand{\tr}{\operatorname{tr}}
\let\dist\relax
\newcommand{\dist}{\operatorname{dist}}
\let\OPT\relax
\newcommand{\OPT}{\operatorname{OPT}}
\chapter{Inherited-error return and the first adaptive merge}\label{ch:return}
\begin{scopebox}
The first-adaptive-merge and zero-tail statements remain valid. They do not iterate into an all-depth guarantee for the annihilate-first rule; Chapter~\ref{ch:anchor} disproves that guarantee. The inherited-error identity is retained as a valid local interface.
\end{scopebox}
\begin{chaptersynopsis}
The all-input, $O(k)$-query, expected $O(L)$-relative-error HSS theorem remains unproved. This chapter gives an unconditional physical-error estimate at the first adaptive merge of the specified annihilate-first learner. The newly discarded comparator loss is bounded by $8/11$ times inherited loss plus an absolute multiple of the original residual crossing the current cut. With no crossing residual the coefficient improves to $4/11$. The proof uses two different valid conditioning orders and sixth-moment H\"older bounds; it never asserts independence of the dependent inverse factors. An explicit actual-rule example shows that current residual data and current sketch tails can both vanish while new physical loss is positive, even with well-conditioned current designs and a nonvanishing exterior anchor. We also prove a fixed-input, all-depth zero-tail characterization of HSS rank $2k$, and record an optional query-only safeguard against global low-rank and leaf-diagonal benchmarks. None of these statements converts small all-depth sketch tails into the missing best-HSS physical approximation bound.
\end{chaptersynopsis}

\section{Target, scope, and established interfaces}
A balanced binary tree on $N$ coordinates has original leaf size $2k$. Write $\HH_k$ for the matrices whose outgoing and incoming exterior blocks at every nonroot node have rank at most $k$, and set
\[
 \OPT_k(A)=\min_{B\in\HH_k}\F{A-B}^2.
\]
The target is an algorithm using at most $C_qk$ products with $A,A^T$, returning $\widehat B\in\HH_k$ with
\begin{equation}\label{return:eq:target}
 \E\F{A-\widehat B}^2\le CL\OPT_k(A).
\end{equation}
Here $L$ bounds the number of nonroot partitions. Computation other than oracle products is free; an exact nonlinear projection used below is not claimed to be efficient.

The following already established interfaces are retained, not reproved. For a nested rank-$s$ coefficient space $\SSS$ chosen independently of fresh refit probes, the fixed-space refit uses $8s$ queries and obeys
\begin{equation}\label{return:eq:refit}
 \E[\F{A-C_h}^2\mid\SSS]\le1607\dist_F(A,\SSS)^2.
\end{equation}
Chapters~\ref{ch:interfaces} and~\ref{ch:uniform} prove this theorem without a depth logarithm. Also, if a known matrix $C$ has a useful approximation guarantee, exact offline projection onto $\HH_k$ uses no further queries and satisfies
\begin{equation}\label{return:eq:projection}
 \F{A-\widehat B}\le2\F{A-C}+\sqrt{\OPT_k(A)}.
\end{equation}
Existence follows because $\HH_k$ is closed, contains zero, and a minimizing sequence for distance from $C$ may be restricted to a bounded ball.

The public fresh-level HSS analysis provides the surrounding approximation setting \cite{AmselHSS}. The results here concern a particular reused-sketch learner, not a replacement of that paper's independence hypotheses. The Gaussian inverse and rangefinder estimates used below are classical \cite{HMT}. No priority or optimal-constant claim is made.

\section{The actual learner and its error registers}
Use
\[
 r=2k,\qquad m=4k=2r,\qquad q=16k.
\]
Group pairs of original leaves into coarse leaves of size $m$. Unless otherwise stated, consider the hierarchical branch $N>q$; small $N$ can be handled separately by exact coordinate queries. Draw independent standard Gaussian matrices $\Omega,\Psi\in\R^{N\times q}$ and query
\[
 Y=A\Omega,\qquad Z=A^T\Psi.
\]
For a node $I$, with physical coordinate projector $P_I$, define the original-data commutator
\begin{equation}\label{return:eq:J}
 J_I(A)=\Psi_I^TY_I-Z_I^T\Omega_I
       =\Psi^T[P_I,A]\Omega.
\end{equation}
Parent commutators are sums of their child commutators. No recursively fitted coefficient is substituted into these observations.

At a physical leaf the row and column samples are
\[
 X_U=Y_I(\Id-\Omega_I^\dagger\Omega_I),\qquad
 X_V=Z_I(\Id-\Psi_I^\dagger\Psi_I).
\]
Take their leading-$r$ physical left spaces. If the sample rank is at most $r$, canonically extend its full physical range. Fix the following canonical convention: frame a range by projecting coordinate vectors in order and orthonormalizing; extend with coordinate vectors in order, and frame the resulting space the same way. In exact mathematics rank is algebraic rank.

At an internal node use the child templates
\[
 W_U=\diag(\mathcal U_a,\mathcal U_b),\quad
 W_V=\diag(\mathcal V_a,\mathcal V_b),\quad
 G=W_U^T\Psi_I,\quad H=W_V^T\Omega_I.
\]
They have $m$ columns. Set
\[
 \Pi_G=G^\dagger G,\quad N_G=\Id-\Pi_G,\qquad
 \Pi_H=H^\dagger H,\quad N_H=\Id-\Pi_H.
\]
Form
\begin{equation}\label{return:eq:metric}
 M_U=\Pi_GJ_I(A)N_H,\qquad
 M_V=-\Pi_HJ_I(A)^TN_G.
\end{equation}
When $\rank M_U>r$, retain a leading-$r$ left sketch space and lift it through $(G^T)^\dagger$. When $\rank M_U\le r$, canonically extend the full physical range of $(G^T)^\dagger M_U$. Apply the analogous column rule. This rule is defined with Moore--Penrose inverses also on singular designs.

Write $U,V$ for the local physical bases, and $P=UU^T$. Let $\Pi_{GU}$ project onto $\ran(G^TU)$ and define
\begin{equation}\label{return:eq:tail}
 \mathcal R_U=(\Pi_G-\Pi_{GU})J_I(A)N_H,
 \qquad
 \mathcal R_V=(\Pi_H-\Pi_{HV})J_I(A)^TN_G.
\end{equation}
The sign convention for the column sample does not affect its squared tail.

For fixed $B\in\HH_k$ and $E=A-B$, inherited results for this rule are
\begin{align}
 \F{\mathcal R_U}^2,\F{\mathcal R_V}^2
       &\le\F{J_I(E)}^2 &&\text{at internal nodes, pathwise},\label{return:eq:oldtail}\\
 \sum_{I,o}\E\F{\mathcal R_{o,I}}^2
       &\le6q^2L\F E^2,\label{return:eq:oldglobal}\\
 \sum_{I\text{ coarse leaf}}\E(e_{U,I}^2+e_{V,I}^2)
       &\le32\F{\operatorname{off}_{\rm leaf}(E)}^2.\label{return:eq:oldleaf}
\end{align}
The first follows directly from $\rank J_I(B)\le2k=r$ and best-rank approximation in sketch coordinates. The last is the earlier first-physical-layer theorem. On each fixed exact HSS-$k$ input the hierarchy is deterministic almost surely; that statement must not be extended to approximate inputs by conditioning on their learned spaces.

\section{A stream-separated local theorem}
The useful measurability property at the first adaptive merge is
\begin{equation}\label{return:eq:separate}
 W_U\text{ is measurable with respect to all }\Omega,
 \qquad
 W_V\text{ is measurable with respect to all }\Psi.
\end{equation}
Indeed the physical-leaf row rule uses only $A\Omega,\Omega$, whereas the column rule uses only $A^T\Psi,\Psi$. The result below permits \emph{arbitrary} isometries satisfying \eqref{return:eq:separate}, not only the actual first-merge templates. Both may depend on the fixed input $A$.

Fix a node and factorizations
\[
 B[I,I^c]=CR^T,\qquad B[I^c,I]^T=VK,
\]
where $R,V$ have orthonormal columns, with ranks $a,b\le k$. Put
\begin{align*}
 C_c&=W_U^TC,& C_\perp&=(\Id-W_UW_U^T)C,\\
 \Gamma&=R^T\Omega_{I^c},& X&=V^T\Omega_I,\qquad
 \Theta=K\Psi_{I^c},\\
 \mathcal L&=(G^\dagger)^T,&D&=\mathcal L\Psi_I^TC_\perp.
\end{align*}
Define the local active and inherited comparator losses
\[
 e_U=\F{(\Id-P)C_c},\qquad h_U=\F{C_\perp}.
\]
These equal the corresponding outgoing-block losses because $R$ is an isometry. Define $e_V,h_V$ by transposition. Finally set
\begin{equation}\label{return:eq:cutenergy}
 \sigma_I^2=\F{E[I,I^c]}^2+\F{E[I^c,I]}^2
            =\F{[P_I,E]}^2.
\end{equation}

\begin{lemma}[A separating right inverse]\label{return:lem:separation}
Assume \eqref{return:eq:separate} and $a>0$. Almost surely $G$ has full row rank and there exists $T\in\R^{q\times a}$ such that
\begin{equation}\label{return:eq:T}
 HT=0,\qquad XT=0,\qquad\Gamma T=\Id_a.
\end{equation}
It can be chosen as $T=N(\Gamma N)^\dagger$, where the columns of $N$ are an orthonormal basis of $\ker[H;X]$. Conditional on all $\Psi$ and on $\Omega_I$, $\Gamma N$ is an $a\times d$ standard Gaussian matrix with $d\ge q-m-b\ge11k$.
\end{lemma}
\begin{proof}
Conditional on all $\Omega$, $W_U$ is fixed, so $G=W_U^T\Psi_I$ is standard Gaussian and full row rank almost surely. This conditioning is used only for that conclusion and later for the return moment.

For the right inverse, use a \emph{different} conditioning: fix all $\Psi$ and $\Omega_I$. Now $W_V,H,X,N$ are fixed and $\Gamma=R^T\Omega_{I^c}$ is independent standard Gaussian. The nullspace dimension obeys $d\ge q-m-b\ge a$, so $\Gamma N$ has full row rank almost surely. This proves \eqref{return:eq:T}. In particular, it is not necessary to condition on $G$ or on the $\Omega$-learned row template when proving this anchor rank.
\end{proof}

\begin{lemma}[Signed local physical-loss identity]\label{return:lem:signed}
Whenever $G$ has full row rank and \eqref{return:eq:T} holds, the actual selected projector satisfies
\begin{equation}\label{return:eq:signed}
 \boxed{(\Id-P)C_c=(\Id-P)\mathcal L
 \big[(\mathcal R_U-J_I(E))T-\Psi_I^TC_\perp\big].}
\end{equation}
Consequently, under \eqref{return:eq:separate}, this identity holds almost surely, and
\begin{equation}\label{return:eq:pathbound}
 e_U\le2\op{\mathcal L}\op T\F{J_I(E)}+\F D.
\end{equation}
The identity itself may be used at later levels on its stated full-rank and separation locus; no probability-one assertion for that locus at all later levels is made here.
\end{lemma}
\begin{proof}
Direct block multiplication gives
\[
 J_I(B)=\Psi_I^TC\Gamma-\Theta^TX.
\]
Multiplication by $T$ removes the incoming comparator term and gives
\[
 \mathcal L J_I(B)T=C_c+D.
\]
Since $\mathcal L\Pi_G=\mathcal L$ and $\mathcal L$ maps $\ran(G^TU)$ into $\ran U$, we have
\[
 (\Id-P)\mathcal L\mathcal R_UT
  =(\Id-P)\mathcal LJ_I(A)T.
\]
Here $N_HT=T$ is essential. Substituting $A=B+E$ proves \eqref{return:eq:signed}. The inherited sketch-metric tail inequality \eqref{return:eq:oldtail} and norm contraction then give \eqref{return:eq:pathbound}. If $a=0$, the outgoing loss is zero and no right inverse is needed.
\end{proof}

\begin{lemma}[Exact return moment]\label{return:lem:D}
Under \eqref{return:eq:separate},
\begin{equation}\label{return:eq:Dmoment}
 \E[\F D^2\mid\Omega]=\rho\F{C_\perp}^2,
 \qquad \rho=\frac{m}{q-m-1}=\frac{4k}{12k-1}\le\frac4{11}.
\end{equation}
\end{lemma}
\begin{proof}
Conditional on $\Omega$, both $W_U$ and $C_\perp$ are fixed. The Gaussian physical projections $G=W_U^T\Psi_I$ and $\Psi_I^TC_\perp$ are independent, since $W_U^TC_\perp=0$. Averaging the latter gives
\[
 \E[\F{\mathcal L\Psi_I^TC_\perp}^2\mid\Omega,G]
     =\F{\mathcal L}^2\F{C_\perp}^2.
\]
The Gaussian inverse-Frobenius identity \cite{HMT} gives
$\E[\F{\mathcal L}^2\mid\Omega]=m/(q-m-1)$. No independence of $P$ and $D$ is asserted.
\end{proof}

\section{Paying a noisy cut without false independence}
The residual term in \eqref{return:eq:pathbound} involves dependent factors. Sixth moments are sufficient to pay it by H\"older's inequality.

\begin{lemma}[Uniform sixth-moment bounds]\label{return:lem:six}
Under the hypotheses of Lemma~\ref{return:lem:separation},
\begin{align}
 (\E\op{\mathcal L}^{6})^{1/3}&<\frac{14}{q},&
 (\E\op T^{6})^{1/3}&<\frac{14}{q},\label{return:eq:inverse6}\\
 (\E\F{J_I(E)}^{6})^{1/3}&<1.26q^2\sigma_I^2.\label{return:eq:J6}
\end{align}
In particular,
\begin{equation}\label{return:eq:holder}
 \E[\op{\mathcal L}^2\op T^2\F{J_I(E)}^2]<250\sigma_I^2.
\end{equation}
\end{lemma}
\begin{proof}
For an $s\times d$ Gaussian matrix $F$, put $p=d-s>5$. The standard spectral pseudoinverse tail \cite[Sec.~10.1, Eq.~(71) in the arXiv HTML rendering]{HMT} is
\[
 \PP\{\op{F^\dagger}\ge e\sqrt d\,t/(p+1)\}\le t^{-(p+1)},\qquad t\ge1.
\]
Integrating $6t^5$ times the tail gives
\begin{equation}\label{return:eq:integrated6}
 (\E\op{F^\dagger}^6)^{1/3}
 \le\frac{e^2d}{(p+1)^2}\left(\frac{p+1}{p-5}\right)^{1/3}.
\end{equation}
For $G$, use $(s,d)=(4k,16k)$. For $T$, conditional on the variables in Lemma~\ref{return:lem:separation}, its nonzero singular values are those of $(\Gamma N)^\dagger$. Its inverse norm is stochastically bounded by the inverse norm of a $k\times11k$ Gaussian matrix: append rows to reach $k$, and remove columns to reach $11k$, using eigenvalue interlacing and positive-semidefinite order. Apply \eqref{return:eq:integrated6} with $(s,d)=(k,11k)$.

For either bound the expression after multiplication by $q$ decreases with $k$. To verify this, the varying factor has the form
\[
 f(x)=\frac{x^2}{(x+1)^{5/3}(x-5)^{1/3}},\qquad
 \frac{d}{dx}\log f(x)=\frac{-10}{x(x+1)(x-5)}<0.
\]
At $k=1$, the two factors are bounded using $e^2<7.4$ by
\[
 \frac{7.4\cdot256}{13^2}(13/7)^{1/3}<14,
 \qquad
 \frac{7.4\cdot176}{11^2}(11/5)^{1/3}<14.
\]
The one-row case can also be obtained directly from an inverse chi-square moment.

For \eqref{return:eq:J6}, a centered Gaussian vector with covariance $\Sigma$ obeys
\[
 \E\|z\|^6=(\tr\Sigma)^3+6\tr\Sigma\tr(\Sigma^2)
                         +8\tr(\Sigma^3).
\]
This follows by expansion of independent squared Gaussian coordinates. For a Gaussian matrix with $q$ independent identically distributed columns and covariance $K K^T$ per column, it implies
\[
 \E\F{K\Omega}^6\le q(q+2)(q+4)\F K^6.
\]
Apply this first conditionally to $\Psi^TK\Omega$, then to $K\Omega$, with $K=[P_I,E]$. Thus
\[
 \E\F{J_I(E)}^6\le[q(q+2)(q+4)]^2\sigma_I^6.
\]
Since $q\ge16$, the normalized cube-root factor is at most
$((1+2/16)(1+4/16))^{2/3}=(45/32)^{2/3}<1.26$.
Finally H\"older, without any independence assumption between its three factors, gives \eqref{return:eq:holder}, because $14^2\cdot1.26<250$.
\end{proof}

\begin{theorem}[Noisy first adaptive return]\label{return:thm:first}
For any stream-separated templates satisfying \eqref{return:eq:separate}, and therefore at every actual parent of two coarse physical leaves, the outgoing local physical loss satisfies
\begin{equation}\label{return:eq:first}
 \boxed{\E e_U^2\le2\rho\E h_U^2+2000\sigma_I^2
                 \le\frac8{11}\E h_U^2+2000\sigma_I^2.}
\end{equation}
The same assertion holds for the incoming loss. No singular-value gap, exact lower-level containment, or favorable-event assumption is required.

If no residual crosses the current cut, $\sigma_I=0$, the sharper estimate is
\begin{equation}\label{return:eq:clean}
 \boxed{\E e_U^2\le\rho\E h_U^2\le\frac4{11}\E h_U^2.}
\end{equation}
The residual inside the node can still be arbitrary, and the children can have nonzero representation error.
\end{theorem}
\begin{proof}
Squaring \eqref{return:eq:pathbound} with $(x+y)^2\le2x^2+2y^2$ and using Lemmas~\ref{return:lem:D}--\ref{return:lem:six} proves \eqref{return:eq:first}. For a clean cut, $J_I(A)=J_I(B)$ has rank at most $2k=r$, so the actual full-range branch has $\mathcal R_U=0$. The exact identity reduces to
\begin{equation}\label{return:eq:cleanidentity}
 (\Id-P)C_c=-(\Id-P)D.
\end{equation}
Orthogonal projection cannot increase norm; Lemma~\ref{return:lem:D} proves \eqref{return:eq:clean}. Transpose the construction and exchange the two independent streams for the column assertion.
\end{proof}

\begin{remark}[What contracts]
$e_U^2$ is the \emph{new} loss generated at the parent, not the full loss of the parent exterior block. The full physical loss equals $h_U^2+e_U^2$ by Pythagoras. Thus \eqref{return:eq:clean} gives at most $(1+\rho)\E h_U^2$ for that full loss; it is not a contraction of total error.
\end{remark}

\begin{corollary}[Two learned partitions]\label{return:cor:two}
Let $\mathcal P_0$ be the coarse-leaf partition and $\mathcal P_1$ its parent partition, omitting the root if necessary. For the actual learner,
\begin{align}
 \sum_{I\in\mathcal P_1}\E(e_{U,I}^2+e_{V,I}^2)
 &\le64\rho\F{\operatorname{off}_{\mathcal P_0}(E)}^2
      +8000\F{\operatorname{off}_{\mathcal P_1}(E)}^2,\label{return:eq:secondsum}\\
 \sum_{I\in\mathcal P_0\cup\mathcal P_1}\E(e_{U,I}^2+e_{V,I}^2)
 &\le8056\F E^2.\label{return:eq:twosum}
\end{align}
\end{corollary}
\begin{proof}
The parent exterior block restricted to a child is a restriction of that child's exterior block. Hence summed inherited parent losses are at most the summed leaf losses in \eqref{return:eq:oldleaf}. Also
$\sum_{I\in\mathcal P_1}\sigma_I^2=2\F{\operatorname{off}_{\mathcal P_1}(E)}^2$.
Sum both orientations of \eqref{return:eq:first}, and then add \eqref{return:eq:oldleaf}. The final scalar bound is
$32+64(4/11)+8000<8056$.
\end{proof}

\section{An actual-rule obstruction with well-conditioned current designs}
The next example shows why an all-depth argument must retain inherited physical loss. It is not an approximation lower bound.

\begin{theorem}[Zero current tails need not mean zero new loss]\label{return:thm:example}
For $k=1$, $r=2$, $m=4$, $q=16$, the actual hybrid rule admits fixed inputs $A_\epsilon=B+\epsilon E$ and positive-probability Gaussian probe events at a first adaptive parent $I$ on which
\[
 J_I(E)=0,\qquad \mathcal R_{U,I}=\mathcal R_{V,I}=0,
 \qquad e_{U,I}>0.
\]
The current row design can have every singular value in $[3/4,5/4]$, the column design can be arbitrarily close to a row isometry, and the surviving exterior anchor can be bounded away from zero. Moreover no finite deterministic gain $e_U^2\le\gamma h_U^2$ holds uniformly over all realizations, even with these current-design and anchor safeguards.
\end{theorem}
\begin{proof}
Take $N=32$, coarse leaves $a,b$ of size four, parent $I=a\cup b$, and a coordinate $j\notin I$. Let $e_{a,t},e_{b,t}$ denote physical coordinate vectors and set
\[
 u=e_{a,1}+e_{b,1},\quad B=ue_j^T,\quad
 E[a,b]=e_3e_1^T+2e_2e_2^T+\tfrac32e_4e_3^T,
\]
with all other entries of $E$ zero. Then $B$ is HSS rank one, $\F E^2=29/4$, and $E$ does not cross the parent cut. Fix $M\ge1$ and
$0<\epsilon\le\min\{0.1,0.1/M\}$.

Let $f_1,\ldots,f_{16}$ denote standard row vectors in sketch coordinates. Consider the following reference probe rows:
\begin{center}
\begin{tabular}{@{}cll@{}}\toprule
Physical rows & $\Omega$ & $\Psi$\\\midrule
$a$ & $f_1,f_2,f_3,f_4$ & $f_1,f_2,f_1+f_3+Mf_6,f_4$\\
$b$ & $f_5,f_6,f_7,f_8$ & $f_5,f_6,f_7,f_8$\\
$j$ & $f_1+f_5$ & $f_9$\\\bottomrule
\end{tabular}
\end{center}
Other rows may be filled arbitrarily; they do not affect this local calculation.

At the reference probes the actual leaf row spaces are
\[
 U_a=\left[\frac{e_1+\epsilon e_3}{s},e_2\right],\qquad
 U_b=[e_1,e_2],\qquad s=\sqrt{1+\epsilon^2},
\]
and the leaf column spaces are
\[
 V_a=[e_1,e_2],\qquad V_b=[e_2,e_3].
\]
Indeed the outgoing sample at $a$ has orthogonal sketch components
$(e_1+\epsilon e_3)f_5+2\epsilon e_2f_6+(3\epsilon/2)e_4f_7$, so its two leading directions are as stated. At $b$ the outgoing sample is $e_1f_1$. There is no incoming sample at $a$. The incoming sample at $b$, after nullifying $\Psi_b$, has row singular magnitudes $\epsilon\sqrt2,2\epsilon,3\epsilon/2$; its leading two directions are $e_2,e_3$. All relevant truncation gaps are strict.

In the resulting canonical local coordinates the parent designs have rows
\begin{align*}
 G&=\left[
 \frac{(1+\epsilon)f_1+\epsilon f_3+\epsilon Mf_6}{s};\
 f_2;\ f_5;\ f_6\right],\\
 H&=[f_1;f_2;f_6;f_7].
\end{align*}
The surviving right anchor is $(f_1+f_5)N_H=f_5$, of norm one. The difference between $G$ and the row isometry with rows $f_1,f_2,f_5,f_6$ has operator norm below $1/4$, since its only changed row has squared norm at most $2\epsilon^2+\epsilon^2M^2\le0.03$. Hence its singular values are in $[3/4,5/4]$.

Put $d_\epsilon=(1+\epsilon)^2+\epsilon^2$. The lifted rank-one parent sample has physical direction
\[
 c=\left(\frac{s(1+\epsilon)}{d_\epsilon},\ 0,\ 1,\
          -\frac{\epsilon M(1+\epsilon)}{d_\epsilon}\right)^T.
\]
Canonical full-range extension selects the space $\ran[c,e_1]$. The true comparator coordinates are $(1/s,0,1,0)^T$. Therefore direct projection gives
\begin{equation}\label{return:eq:exampleformulas}
 \boxed{h_U^2=\frac{\epsilon^2}{1+\epsilon^2},\qquad
 e_U^2=\frac{\epsilon^2M^2(1+\epsilon)^2}
 {d_\epsilon^2+\epsilon^2M^2(1+\epsilon)^2}>0.}
\end{equation}
In particular $e_U^2/h_U^2\longrightarrow M^2$ as $\epsilon\downarrow0$.

For this parent, $J_I(E)=0$ for every probe realization, and $J_I(A)=J_I(B)$ has rank at most one. Both projected samples have rank at most one and hence both actual reported tails vanish in the full-range branch. The losses above therefore cannot be charged solely to current residual or current tail energy.

The reference is converted to a positive-probability event as follows. For fixed $\epsilon>0,M<\infty$, the leaf spectral gaps, design ranks, nonzero anchor, and the row-frame and parent-extension pivots just used persist in a sufficiently small open neighborhood of the reference rows. The relevant spaces and the two losses vary continuously there. The column frame itself need not be continuous: the parent annihilator depends only on its subspace, so rotations of that frame do not affect this conclusion. The structural identities $J_I(E)=0$ and the rank-one bound on $J_I(B)$ persist exactly, so both current tails remain zero. Independent Gaussian probes have positive density on every such open set. Given any finite $\gamma$, choose $M^2>\gamma$ and then sufficiently small $\epsilon$; the strict gain violation also persists on a positive-probability neighborhood.
\end{proof}

\begin{remark}[Compatibility with the mean contraction]
The probability of the event in Theorem~\ref{return:thm:example} is not bounded uniformly below as $M$ grows or $\epsilon$ shrinks. The theorem is fully consistent with the unconditional clean-cut mean gain $\rho\le4/11$. It rules out a realization-wise current-data converse and a realization-wise inherited-error contraction. It does not rule out the learner, the expected all-depth estimate, or the general query target.
\end{remark}

\section{A qualitative all-depth zero-tail certificate}
Let
\[
 \tau(A;\Omega,\Psi)=\sum_{I,o}\F{\mathcal R_{o,I}}^2
\]
include every nonroot coarse-tree node, with the physical-leaf selection rule unchanged. This is computable from the sketches and local factors.

\begin{theorem}[Fixed-input zero-tail characterization]\label{return:thm:zero}
For every fixed input matrix $A$, on the hierarchical branch $N>q$,
\begin{equation}\label{return:eq:zero}
 \boxed{\tau(A;\Omega,\Psi)=0
       \quad\Longleftrightarrow\quad A\in\HH_{2k},\qquad\text{almost surely}.}
\end{equation}
On this zero-tail endpoint, the actual learned coefficient space contains $A$. The probability-one event is allowed to depend on the fixed input $A$; the statement is not simultaneous over all matrices for one realization of the probes.
\end{theorem}
\begin{proof}
First suppose $A\in\HH_r$, $r=2k$. The deterministic canonical comparison induction used at rank $k$ extends to rank $r$: at a leaf the exterior rank is at most $r$ and the nullified exterior Gaussian sample has $q-m=12k\ge r$ effective columns. It therefore recovers the full true exterior range almost surely. At a comparison internal node, the templates are deterministic, contain the true outgoing and incoming ranges, and have full-row-rank Gaussian designs. Incoming containment kills the incoming commutator term. The remaining sample is
\[
 J_I(A)N_H=G^TC_c\Gamma N_H,
\]
with $\Gamma$ independent of the inside probes and $\Gamma N_H$ full row rank. The actual full-range branch returns exactly the canonical comparison space. Inducting over the finitely many nodes proves exact containment and $\tau=0$.

Conversely suppose $A\notin\HH_r$. Choose a finest coarse-tree node $I$ at which either exterior rank exceeds $r$. Every strict descendant has both ranks at most $r$, so the preceding deterministic comparison argument applies throughout its descendant subtrees, regardless of failures elsewhere in the tree. This locality uses the original-data commutators, rather than recursively modified coefficient observations.

If $I$ is a coarse leaf and its outgoing rank exceeds $r$, its nullified exterior sample has rank $\min\{\rank A[I,I^c],q-m\}>r$ almost surely. The retained leading-$r$ physical space leaves a nonzero residual; injectivity of $\Psi_I^T$ prevents its image from lying in the retained sketch image. Hence $\mathcal R_U\ne0$. The column case is the same.

If $I$ is internal, its deterministic child templates contain the full outgoing and incoming exterior ranges. The outgoing rank $a$ is at most $m=2r$ because each child restriction has rank at most $r$. Compression by $W_U^T$ preserves this rank. With $q-m=12k\ge m$, the Gaussian anchor after nullifying $H$ still has row rank $a$. Thus $M_U$ has rank exactly $a$. If $a>r$, leading-$r$ truncation has a nonzero tail; because $G^T$ is injective, lifting the selected $r$-space cannot increase its sketch dimension and cannot remove that tail. The incoming case is identical. Therefore $\tau>0$ almost surely.

Finally coarse-tree membership in $\HH_r$ is equivalent to membership on the original tree: every proper original descendant of a coarse leaf has physical size at most $r$, so its two exterior ranks are automatically at most $r$.
\end{proof}

\begin{remark}[What this certificate does not say]
This is an exact-real zero test. It gives no numerical threshold, no stability modulus near zero, and no inequality bounding physical distance by a small value of $\tau$. The local example in Section~5 is compatible with the theorem: that example has nonzero lower-level tails. A uniform statement over sketch-dependent candidate inputs would also be a different assertion and is not proved.
\end{remark}

\section{Diagnostics beyond a globally rank-one comparator}
The scripts accompanying this chapter use matrix entries only for synthetic generation and oracle diagnostics; these uses are not hidden learner queries. None of the sampled values proves a uniform expected-risk claim.

The local suite checks 24 signed identities and eight explicit reference constructions. The largest normalized signed-identity discrepancy is below $6.6\times10^{-15}$. A further 240 actual noisy first-merge tests cover $k=1,2,3$ and have normalized identity discrepancies below $5.6\times10^{-15}$. Four independent positive-sector versus full-projection checks agree within $4.5\times10^{-16}$ after normalization. Six zero-tail tests cover exact HSS rank $2k$, dense matrices, and perturbed HSS rank $2k$ matrices.

A separate clean-cut Monte Carlo diagnostic uses 400 probe draws for each of $k=1,2$, with both orientations. Its sampled active/inherited energy ratios are approximately $0.161$ and $0.171$. The corresponding sampled regression/inherited ratios are $0.365$ and $0.366$, compared with exact theoretical factors $4/11\approx0.364$ and $8/23\approx0.348$. These finite-sample ratios are not used in proving the exact conditional calculation.

For deeper tests, generate a fixed general HSS-rank-one matrix $B$ from recursive leaf blocks and sibling couplings, normalize $\F B=1$, then fix a permutation matrix $P_0$ independently of the probes. Set
\[
 A=B+\epsilon P_0/\sqrt N,\qquad\epsilon=10^{-4}.
\]
This comparator is not restricted to global rank one. The data below measure the actual coefficient-space distance
$\dist_F(B,\SSS(A))^2/\epsilon^2$, using a positive orthogonal multilevel decomposition, not an upper-bound ledger. Input generation uses seed $20260919+N$, and each dimension uses five probe seeds, $11,12,13,14,15$.
\begin{center}
\begin{tabular}{@{}rrrr@{}}\toprule
$N$ & Minimum & Median & Maximum\\\midrule
128  &0.2880&0.4382&0.7224\\
256  &0.3494&0.3911&0.5391\\
512  &0.4411&0.4655&0.6241\\
1024 &0.4389&0.5758&0.6090\\
2048 &0.5304&0.6088&0.7853\\\bottomrule
\end{tabular}
\end{center}
There are 25 runs, not enough to estimate a worst-case expectation or infer a depth-uniform theorem. The exact deterministic comparison
\[
 \dist_F(A,\SSS)\ge\big(\dist_F(B,\SSS)-\epsilon\big)_+
\]
continues to explain why comparator-space loss is a relevant diagnostic.

\section{What remains for the general target}
The first two learned partitions are now paid in physical Frobenius norm. The next templates, however, are produced by the two-sided internal rule and generally depend on both $\Omega$ and $\Psi$. Two hypotheses used above are then absent: $G$ need not be Gaussian after fixing the row template, and the surviving exterior anchor need not remain independent after fixing the opposite template.

Consequently the proofs of Lemmas~\ref{return:lem:separation}, \ref{return:lem:D}, and \ref{return:lem:six} cannot simply be iterated. The new strict first-return coefficient does not itself prevent inherited error from being presented again at several ancestors. Even a levelwise inequality with coefficient less than one would require a valid accounting of which inherited sources it reuses.

A sufficient all-depth result remains
\begin{equation}\label{return:eq:remaining}
 \E\dist_F(A,\SSS_{\rm learned})^2\le C_bL\OPT_k(A).
\end{equation}
Orthogonal coefficient assembly introduces no extra depth factor if the sum of local physical losses is already controlled. The exact signed register in \eqref{return:eq:signed}, rather than a separate bound on an inverse alone, is the natural next object. A general learner needs an actual-law, multilevel control of inherited return, including right-anchor separation or an alternative to dividing by that anchor. The local theorem here does not provide it for the hybrid: Chapter~\ref{ch:anchor} disproves that learner's global basis bound. Theorem~\ref{return:thm:example} shows that the inherited term cannot be discarded just because present cut data fit exactly.

Training still costs $32k$ products. If \eqref{return:eq:remaining} is established, an independent rank-$2k$ refit costs $16k$ more, so \eqref{return:eq:refit} and \eqref{return:eq:projection} give the original target at $48k$ queries. This implication is conditional; the all-depth expected best-HSS theorem is not claimed here.


\section{An optional all-input safeguard from the existing sketches}
This appendix records a practical query-only consequence, using classical low-rank sampling and a fresh validation batch. It does not remove the basis-learning gap. No new approximation or selection method is claimed as novel.

\subsection{Two additional known candidates at no training-query cost}
Let $Q$ be an orthonormal basis of the nonzero range of the first $4k$ columns of $Y=A\Omega$, and put $s=\rank Q\le4k$. Define
\[
 C_g=Q(\Psi^TQ)^\dagger Z^T.
\]
Because $Q$ depends only on $\Omega$, it is independent of $\Psi$. With $F=(\Id-QQ^T)A$, orthogonality and conditional Gaussian regression give
\[
 \E_\Psi\F{A-C_g}^2
   =\left(1+\frac{s}{16k-s-1}\right)\F F^2.
\]
The standard Gaussian rangefinder calculation \cite{HMT}, with target rank $k$ and $4k$ columns, gives the squared-error bound
\[
 \E\F F^2\le\left(1+\frac{k}{3k-1}\right)
            \min_{\rank X\le k}\F{A-X}^2.
\]
For $k=1$ the same calculation uses the inverse chi-square identity. Thus
\begin{equation}\label{return:eq:globalguard}
 \E\F{A-C_g}^2\le\frac{45}{22}\OPT_k^{\rm global}(A).
\end{equation}
This rank-at-most-$4k$ candidate is known from the same two original streams.

For each \emph{original} leaf $I$ of size $2k$, set
\[
 C_d[I,I]=Y_I\Omega_I^\dagger,\qquad C_d[I,J]=0\quad(I\ne J).
\]
The same local Gaussian regression calculation gives
\begin{equation}\label{return:eq:diagguard}
 \E\F{A-C_d}^2
  =\left(1+\frac{2k}{14k-1}\right)
       \F{\operatorname{off}_{\rm original}(A)}^2
  \le\frac{15}{13}\F{\operatorname{off}_{\rm original}(A)}^2.
\end{equation}
Both constructions are analysis-independent algorithms: they do not require a comparator or its residual.

\subsection{Fresh validation with an expected-error bound}
\begin{lemma}[A finite-candidate validation estimate]\label{return:lem:validation}
Let $C_1,\ldots,C_M$ be known matrices fixed independently of a standard Gaussian $V\in\R^{N\times t}$, $t>4$. Select a candidate minimizing $\F{(A-C_i)V}^2$. Then
\begin{equation}\label{return:eq:validation}
 \E_V\F{A-C_{\rm sel}}^2
 \le\left(\frac{Mt(t+2)}{(t-2)(t-4)}\right)^{1/2}
      \min_i\F{A-C_i}^2.
\end{equation}
The candidates may be dependent on each other. A random list from an earlier independent transcript is covered by conditioning on that transcript.
\end{lemma}
\begin{proof}
Write $x_i=\F{A-C_i}^2$ and $z_i=t^{-1}\F{(A-C_i)V}^2$. If some $x_i=0$, validation selects an exact candidate almost surely, so assume every $x_i>0$. Spectral decomposition gives $z_i/x_i$ as a convex combination of independent $\chi_t^2/t$ variables. Therefore
\[
 \E(z_i/x_i)^2\le1+2/t,\qquad
 \E(x_i/z_i)^2\le\frac{t^2}{(t-2)(t-4)},
\]
where the second inequality is Jensen for the convex function $u^{-2}$. For a best candidate $*$,
\[
 x_{\rm sel}\le z_*\max_i(x_i/z_i).
\]
Cauchy--Schwarz and $\max_i b_i^2\le\sum_i b_i^2$ prove \eqref{return:eq:validation}. No independence of the different validation scores is used.
\end{proof}

Apply this with $M=3$, $t=8$, and candidates $C_h,C_g,C_d$, where $C_h$ is the independent fixed-space refit in \eqref{return:eq:refit}. The factor is $\sqrt{10}$. Define
\[
 D(A)=\min\left\{1607\E\dist_F(A,\SSS_{\rm learned})^2,
 \frac{45}{22}\OPT_k^{\rm global}(A),
 \frac{15}{13}\F{\operatorname{off}_{\rm original}(A)}^2\right\}.
\]
Conditioning and $\E\min\le\min\E$ give
\[
 \E\F{A-C_{\rm sel}}^2\le\sqrt{10}\,D(A).
\]
Projecting the known selected matrix exactly onto $\HH_k$ yields
\begin{equation}\label{return:eq:wrapper}
 \E\F{A-\widehat B}^2
 \le\left(2\sqrt{\sqrt{10}\,D(A)}+\sqrt{\OPT_k(A)}\right)^2.
\end{equation}
Since global rank-$k$ matrices and original-leaf-diagonal matrices both lie in $\HH_k$, conservative consequences are
\begin{equation}\label{return:eq:wrapperbenchmarks}
 \E\F{A-\widehat B}^2\le38\OPT_k^{\rm global}(A),\qquad
 \E\F{A-\widehat B}^2\le24\F{\operatorname{off}_{\rm original}(A)}^2.
\end{equation}
The query cost is $32k+16k+8=48k+8\le56k$. The final rank is at most $k$, and exact HSS-$k$ inputs are still returned exactly almost surely. Both guarantees hold for arbitrary fixed inputs, conditional only on the already established refit interface. They are \emph{not} best-HSS-relative guarantees: either displayed benchmark can be arbitrarily larger than $\OPT_k(A)$. The exact final nonlinear projection is a query-model operation, not an efficient recompression claim.

\section{Proof dependencies and computational scope}
The computational archive retains the scripts
\texttt{learner\_core.py}, \texttt{return\_checks.py},
\texttt{noisy\_return\_checks.py}, and \texttt{full\_hss\_stress.py}, with recorded JSON outputs. Run with \texttt{OPENBLAS\_NUM\_THREADS=1}. Floating-point rank thresholds in the scripts are diagnostic conventions; the theorems use exact rank and Moore--Penrose inverses. The core implements the inherited hybrid selection rule; the new scripts check the added identities and bounds. The optional nonlinear final HSS projection and the independent refit are not newly implemented by these scripts.

\endgroup

% ===== ANCHOR =====
\begingroup
\let\E\relax
\newcommand{\E}{\mathbb E}
\let\PP\relax
\newcommand{\PP}{\mathbb P}
\let\R\relax
\newcommand{\R}{\mathbb R}
\let\F\relax
\newcommand{\F}[1]{\lVert #1\rVert_F}
\let\op\relax
\newcommand{\op}[1]{\lVert #1\rVert_{\rm op}}
\let\ran\relax
\newcommand{\ran}{\operatorname{range}}
\let\rank\relax
\newcommand{\rank}{\operatorname{rank}}
\let\diag\relax
\newcommand{\diag}{\operatorname{blockdiag}}
\let\dist\relax
\newcommand{\dist}{\operatorname{dist}}
\let\OPT\relax
\newcommand{\OPT}{\operatorname{OPT}}
\let\HH\relax
\newcommand{\HH}{\mathcal H}
\let\SSS\relax
\newcommand{\SSS}{\mathcal S}
\let\Id\relax
\newcommand{\Id}{\mathrm I}
\let\tr\relax
\newcommand{\tr}{\operatorname{tr}}
\let\supp\relax
\newcommand{\supp}{\operatorname{supp}}
\chapter{Anchor trapping and the shared-query union}\label{ch:anchor}
\begin{scopebox}
\textbf{Update, 20 September 2026.} The annihilate-first counterexample and its companion rescue remain valid as stated. The union introduced here propagates two complete hierarchies separately before taking their physical unions. Later work reaches an additional failure of this separate-view propagation: a clean parent can lose a comparator direction already contained in its child union. That is a positive-probability fixed-input witness, not by itself an asymptotic expected-error lower bound. The current synchronized learner instead feeds a common child union into both parent views, with a full-raw branch and revised completion; see Chapters~\ref{upd:chap:later-obstructions} and~\ref{upd:chap:current-route}.
\end{scopebox}

\begin{scopebox}
This chapter proves a failure of the annihilate-first learned coefficient space, not an information-theoretic lower bound for HSS recovery. The union is a different candidate, run from the same observations; its all-input basis theorem remains open.
\end{scopebox}
\thispagestyle{plain}
\begin{chaptersynopsis}
The previously unproved all-depth basis-quality bound is false for the specified annihilate-first hybrid. For retained rank $r=2k$, width $q=16k$, and canonical full-range completion, we construct deterministic inputs $A_h=B_h+\delta_h E_{0,h}+\epsilon_h E_{1,h}$ with $B_h$ globally rank one, unit Frobenius norm, and $L_h\OPT_k(A_h)\to0$, while the expected squared distance to the actual learned coefficient space tends to one. The bottom perturbation $\delta_h$ is positive and fixed before the probes; its existence follows from a one-sided continuity argument, not from sketch-dependent input selection.

A positive bottom perturbation activates otherwise unused column directions. Above its support, their actual evolution retains projections of a fixed exterior Gaussian anchor. A two-child geometric contraction makes the surviving anchor small. A weak rank-$r$ signal in a sibling then occupies the retained sketch singular space and displaces the rank-one comparator. A conditional Gaussian residual lemma proves that the activating samples remain nonzero at every finite depth. This is a basis-error counterexample, not merely an inverse-moment obstruction.

The independent fixed-space refit remains valid. We give an un-nullified companion that avoids the identified displacement at the problematic cut and a nested union of the two hierarchies using the same training probes. This union has rank at most $4k$ and a pointwise space-domination guarantee. No all-input approximation theorem for the union, or impossibility theorem for the general $O(k)$-query problem, is asserted.
\end{chaptersynopsis}



\section{Target, algorithm, and scope}
Fix a balanced binary coordinate tree. Its original leaves have size $2k$. Let $\HH_k$ be the matrices whose outgoing and incoming exterior blocks at every nonroot node have ranks at most $k$, and put
\[
 \OPT_k(A)=\min_{B\in\HH_k}\F{A-B}^2.
\]
The general research target is an $O(k)$-query algorithm returning $\widehat B\in\HH_k$ with expected squared error at most $CL\OPT_k(A)$. The published fresh-level benchmark uses $O(kL)$ products and has an $O(L)$ expected error factor \cite{AmselHSS}. Nothing proved here contradicts that result or rules out the general constant-query target.

The independent fixed-space interface is used only in Section~\ref{anchor:sec:union}: for rank-$s$ nested bases chosen before its independent probes, the refit of Chapter~\ref{ch:uniform} has
\begin{equation}\label{anchor:eq:refit}
 \E[\F{A-\widehat C}^2\mid\SSS]\le1607\dist_F(A,\SSS)^2
 \quad\text{using }8s\text{ queries}.
\end{equation}
This is an input to this chapter, not proved anew here.

\subsection{The unchanged annihilate-first hybrid}
Write $r=2k$, use coarse physical leaves of size $m=2r=4k$, and take independent standard Gaussian $N\times q$ probes $\Omega,\Psi$, with $q=16k=8r$. From $Y=A\Omega$ and $Z=A^T\Psi$, form the original-data commutators
\begin{equation}\label{anchor:eq:J}
 J_S(A)=\Psi_S^TY_S-Z_S^T\Omega_S=\Psi^T[P_S,A]\Omega.
\end{equation}
At a physical leaf, the row space is the leading-$r$ left space of $Y_SN_{\Omega_S}$ and the column space is the analogous left space of $Z_SN_{\Psi_S}$. Here $N_D=\Id-D^\dagger D$.

At an internal node, let $W_U,W_V$ be the block-diagonal child templates and set
\[
 G=W_U^T\Psi_S,\quad H=W_V^T\Omega_S,\quad
 \Pi_G=G^\dagger G,\quad\Pi_H=H^\dagger H.
\]
The row rule takes a leading-$r$ left sketch space of
\begin{equation}\label{anchor:eq:hybrid}
 M_U=\Pi_GJ_S(A)N_H
\end{equation}
and lifts it through $(G^\dagger)^T$. The column rule uses
$M_V=-\Pi_HJ_S(A)^TN_G$ and $(H^\dagger)^T$. When a sample has rank at most $r$, its full lifted physical range is extended canonically to dimension $r$. The resulting nested physical bases define a linear coefficient space $\SSS_{\rm hyb}(A;\Omega,\Psi)$.

The canonical rule is defined as follows: obtain a coordinate-projector frame for the given range, append the first available coordinate directions, and recanonicalize the resulting projector. In particular, for a rank-one local vector $v\in\R^{2r}$ with a nonzero coordinate after position $r-1$, its completed space is
\begin{equation}\label{anchor:eq:canonical}
 \operatorname{span}\{v,e_1,\ldots,e_{r-1}\}.
\end{equation}
Its first $r-1$ canonical frame vectors are exactly $e_1,\ldots,e_{r-1}$.

For every coarse node $S$, every $C\in\SSS_{\rm hyb}$ satisfies
\begin{equation}\label{anchor:eq:cutrange}
 \ran C[S,S^c]\subseteq\ran\mathcal U_S,
\end{equation}
where $\mathcal U_S$ is the learned physical row basis. Descendant coefficient blocks are internal to $S$ and ancestor blocks use this retained space; this proves \eqref{anchor:eq:cutrange} directly.

\subsection{What the counterexample will and will not imply}
We disprove the proposed learner inequality
\begin{equation}\label{anchor:eq:false}
 \E\dist_F(A,\SSS_{\rm hyb})^2\le C_bL\OPT_k(A).
\end{equation}
It follows that no coefficient refit confined to this learned space can repair the exhibited errors. An additional nonlinear HSS projection that changes the bases is a different decoder: a lower bound on its final error does not follow solely from \eqref{anchor:eq:false}. In particular, this chapter is not an information-theoretic lower bound on all constant-query algorithms.

\section{A deterministic family of inputs}\label{anchor:sec:input}
It is useful to state the construction for any fixed $r\ge2$ and integer width $q>4r$; the project parameters are $r=2k,q=8r$. Let
\[
 n=4r\,2^h,\quad N=4n,\quad
 I=[0,n),\quad K=[n,2n),\quad P=I\cup K,\quad j=2n.
\]
Indices are zero-based. Let $u=n^{-1/2}\mathbf1_I$ and define
\begin{equation}\label{anchor:eq:B}
 B_h=ue_j^T.
\end{equation}
Thus $\F{B_h}=1$ and $B_h\in\HH_1\subseteq\HH_k$.

Let $\mathcal P_h$ be the permutation on $I$ that swaps each pair of adjacent coarse leaves of size $2r$. On each $4r$-block it has the form
\[
 \begin{bmatrix}0&\Id_{2r}\\ \Id_{2r}&0\end{bmatrix}.
\]
Extend it by zero outside $I$ and set
\begin{equation}\label{anchor:eq:noises}
 E_{0,h}=\frac{\mathcal P_h}{\sqrt n},\qquad
 E_{1,h}=\frac1{\sqrt r}\sum_{a=1}^r e_{n+a-1}e_{j+a}^T.
\end{equation}
Both have unit Frobenius norm. Their supports and that of $B_h$ are pairwise disjoint. For fixed positive $\delta,\epsilon$, put
\begin{equation}\label{anchor:eq:A}
 A_h(\delta,\epsilon)=B_h+\delta E_{0,h}+\epsilon E_{1,h}.
\end{equation}
In particular,
\begin{equation}\label{anchor:eq:optupper}
 \OPT_k(A_h)\le\delta^2+\epsilon^2,\qquad
 \F{A_h}^2=1+\delta^2+\epsilon^2.
\end{equation}
Every node of size at least $4r$ inside $I$ has a clean cut:
\begin{equation}\label{anchor:eq:cleanJ}
 J_S(A_h)=a_Sg,\qquad a_S=\Psi_S^Tu_S,\qquad g=e_j^T\Omega.
\end{equation}
The noise $E_0$ has become internal, and $E_1$ touches no row or column in $I$. Nevertheless, the child spaces on those nodes remember $E_0$.

Inside $K$, the outgoing space is the coordinate span of $e_n,\ldots,e_{n+r-1}$ (restricted to the node); the incoming space is zero. The actual row hierarchy is exactly containing and deterministic there, and the column hierarchy is its canonical zero-space completion. This follows by induction: with exact row containment, $J_S(A_h)^TN_G=0$; its column rule therefore stays on the zero branch. The outgoing Gaussian anchor has $q-2r>r$ available columns, so the row rule captures its full rank. At $K$ itself, the retained row and column spaces are both the span of its first $r$ coordinate vectors.

\section{Why a positive perturbation activates every clean column sample}
A central point is that \eqref{anchor:eq:cleanJ} does not imply a zero column sample. It gives
\begin{equation}\label{anchor:eq:pure}
 M_U=(\Pi_Ga_S)(gN_H),\qquad
 M_V=-(\Pi_Hg^T)(a_S^TN_G).
\end{equation}
Whenever both scalar-rank factors are nonzero, the physical full-range rules become
\begin{equation}\label{anchor:eq:surrogate}
 \mathcal U_S=W_U\operatorname{can}_r((G^\dagger)^Ta_S),\qquad
 \mathcal V_S=W_V\operatorname{can}_r((H^\dagger)^Tg^T).
\end{equation}
The second direction is retained regardless of how small $a_S^TN_G$ is.

We first define the two recursions in \eqref{anchor:eq:surrogate} as auxiliary, always-active rules above the physical leaves. We then prove that they equal the actual rules almost surely for each fixed $\delta>0$. This avoids circular conditioning on the other rule's nonvanishing.

\begin{lemma}[Positive leaf loss]\label{anchor:lem:leafloss}
For every fixed $\delta>0$, at each physical leaf $S\subset I$, the actual row space does not contain $u_S$, almost surely. Its column space depends only on the transpose probes in its $4r$ sibling pair and is independent of $\delta>0$.
\end{lemma}
\begin{proof}
For the row rule, up to the common $1/\sqrt n$ factor, the nullified sample is
\[
 (\mathbf1_{2r}g+\delta\Omega_{S'})N_{\Omega_S},
\]
where $S'$ is its sibling. Conditional on $\Omega_S,g$, the coordinates of $\Omega_{S'}N_{\Omega_S}$ in the nullspace are a standard $2r\times(q-2r)$ Gaussian matrix. Its shifted law has a density on all such matrices. Since $r<2r$, the event that its leading-$r$ left space contains the fixed nonzero vector $\mathbf1_{2r}$ has measure zero. Singular-value ties also have measure zero.

The incoming leaf sample is exactly
$\delta\Psi_{S'}N_{\Psi_S}/\sqrt n$. Its leading space is independent of the positive scalar $\delta$, and uses no forward probes. All assertions hold simultaneously over the finitely many leaves.
\end{proof}

\begin{lemma}[A conditional Gaussian residual through all clean merges]\label{anchor:lem:residual}
Fix the entire forward stream on a realization satisfying Lemma~\ref{anchor:lem:leafloss}. In the auxiliary row recursion of \eqref{anchor:eq:surrogate}, all clean predesigns have full row rank $2r$ almost surely. At every such node,
\begin{equation}\label{anchor:eq:nonzero-return}
 a_S^TN_G\ne0\quad\text{almost surely}.
\end{equation}
\end{lemma}
\begin{proof}
At a leaf, write its retained design and full signal image as
\[
 D_S=\mathcal U_S^T\Psi_S,\qquad x_S=\Psi_S^Tu_S.
\]
The leaf template is fixed under the present conditioning. Orthogonal Gaussian projections give
\[
 x_S=D_S^Tc_S+\eta_S,\qquad
 \eta_S\sim N(0,h_S^2\Id_q),\quad h_S>0,
\]
with $\eta_S$ independent of $D_S$. Different leaf pairs $(D_S,x_S)$ are independent because their transpose-probe rows are disjoint; dependence of the fixed templates on forward rows causes no problem.

Maintain a record containing all leaf $D$ matrices and, at each completed clean merge, the predesign $D$ and the regression observation $y=Dx$. The selected physical range depends only on $D,y$, since $(D^\dagger)^Tx=(DD^T)^{-1}y$. Thus the transmitted design is measurable with respect to this record. Conditional on the record, the untransmitted vector $x$ is Gaussian, with some mean and covariance $Q$.

At a parent, before revealing $y$, the two conditional Gaussian child vectors are independent and $x=x_1+x_2$ has covariance $Q_1+Q_2$. This is positive definite at the first merge. At every subsequent merge it is also positive definite almost surely: the inductive conditional-covariance formula below shows
\[
 \ker Q_i=\ran D_{i,\rm pre}^T.
\]
These are independent, rotationally invariant $2r$-subspaces in $\R^q$. They have zero intersection almost surely because $q>4r$. Hence $\ker(Q_1+Q_2)=\{0\}$.

For completeness, the rotational assertion follows from an actual symmetry of this auxiliary row algorithm: right-rotating the local transpose probes rotates $D$ and $x$ together, leaves $(DD^T)^{-1}Dx$ unchanged, and therefore leaves every canonical physical choice unchanged. Subtrees use disjoint transpose rows after the forward stream has been fixed. The same argument makes the two child retained $r$-dimensional query images independent uniform subspaces. Their stack $D$ has rank $2r$ almost surely.

Write $Q_*=Q_1+Q_2\succ0$. Upon observing $y=Dx$, Gaussian conditioning gives
\begin{equation}\label{anchor:eq:conditionalQ}
 Q=Q_*-Q_*D^T(DQ_*D^T)^{-1}DQ_*.
\end{equation}
Its rank is $q-2r$ and its kernel is exactly $\ran D^T$. The conditional covariance of $N_Dx$ is also $Q$, so $N_Dx$ is not the zero vector almost surely. This proves \eqref{anchor:eq:nonzero-return}. The regression vector is nonzero almost surely as well, since $DQ_*D^T\succ0$. Adding the transmitted design to the record reveals nothing further: it is a function of the recorded $D,y$.

This induction conditions on regression observations, not on the full original probe rows or on a fictitious fresh law for learned templates.
\end{proof}

\section{Actual exterior-anchor trapping}
Put $z=g^T/\|g\|$. For a subspace $S$ not orthogonal to $z$, define
\begin{equation}\label{anchor:eq:t}
 t(S)=\frac{1-\|\Pi_Sz\|^2}{\|\Pi_Sz\|^2}.
\end{equation}

\begin{lemma}[Two-child axial contraction]\label{anchor:lem:axial}
Let $S_1,S_2$ be independent $r$-subspaces whose laws are invariant under rotations fixing a unit vector $z$, with neither orthogonal to nor containing $z$. Suppose $q>2r$ and $r>1$. If $S\subseteq S_1+S_2$ contains $\Pi_{S_1+S_2}z$, then
\begin{equation}\label{anchor:eq:axial}
 \E t(S)\le\frac{q-2r+1}{4(q-r)}
                  (\E t(S_1)+\E t(S_2)).
\end{equation}
The child sum has dimension $2r$ and does not contain $z$ almost surely.
\end{lemma}
\begin{proof}
The projections of the child subspaces onto $z^\perp$ are independent uniform $r$-subspaces in dimension $q-1$, so their sum is direct. This proves the dimension and noncontainment assertions.

Let $D_i=S_i\cap z^\perp$ and normalize the closest lifts as
\[
 a_i=\frac{\Pi_{S_i}z}{\|\Pi_{S_i}z\|^2}=z+\eta_i.
\]
Then $\dim D_i=r-1$, $\eta_i\perp D_i,z$, and $\|\eta_i\|^2=t(S_i)$. Conditional on $D_i$ and this length, $\eta_i$ is uniform on the appropriate sphere in dimension $q-r$, hence centered. With $D=D_1+D_2$, the vector
\[
 w=z+\tfrac12(\Id-\Pi_D)(\eta_1+\eta_2)
\]
belongs to $S_1+S_2$ and has $z^Tw=1$. The variational identity
\[
 t(T)=\min_{w\in T,\ z^Tw=1}\|w-z\|^2
\]
therefore bounds the parent tangent by $\tfrac14\|(\Id-\Pi_D)(\eta_1+\eta_2)\|^2$. The mixed term has conditional mean zero. Projecting away the independent $(r-1)$ transverse child directions retains the fraction $(q-2r+1)/(q-r)$ of either squared length in expectation. This proves \eqref{anchor:eq:axial}. Since $S$ contains the projection onto the full sum, its angle is the same as that of the sum. Truncation handles extended expectations; the application below has finite initial mean.
\end{proof}

\begin{theorem}[Actual-rule anchor contraction]\label{anchor:thm:anchor}
For each fixed $h,\delta>0,\epsilon>0$, the actual hybrid agrees with \eqref{anchor:eq:surrogate} at all clean nodes inside $I$, almost surely. Its column query images are independent of $\delta>0$ and $\epsilon$. At a node of size $4r\,2^s$ in $I$, let
$S_s=\ran(\Omega_S^T\mathcal V_S)$. Set
\begin{equation}\label{anchor:eq:constants}
 \mu=\frac{q-2r}{2r-2},\qquad
 \alpha=\frac{q-2r+1}{2(q-r)}<\frac12.
\end{equation}
Then
\begin{equation}\label{anchor:eq:angle}
 \E[t(S_s)\mid\Psi,g]\le\mu\alpha^s.
\end{equation}
For the actual parent design $H_P$ at $P=I\cup K$,
\begin{equation}\label{anchor:eq:anchor-energy}
 \boxed{\E\|gN_{H_P}\|^2\le q\mu\alpha^h.}
\end{equation}
In particular, when $k=1,r=2,q=16$, $\mu=6$, $\alpha=13/28$, and the bound is $96(13/28)^h$.
\end{theorem}
\begin{proof}
Analyze the auxiliary column recursion after fixing all transpose probes and $g$. Its physical-leaf templates are fixed by Lemma~\ref{anchor:lem:leafloss}. The predesign at each first clean $4r$-node is thus a standard $2r\times q$ Gaussian matrix, independently over those nodes. Its query image is uniform of dimension $2r$. Retention contains the projection of $z$ onto that image and preserves its angle. The initial squared tangent is a ratio of independent chi-squares with degrees $q-2r$ and $2r$, giving mean $\mu$.

Right rotations fixing $z$ leave $(H^\dagger)^Tz$ and its canonical physical completion unchanged, while rotating the query image. The two child image laws therefore have exactly the symmetry and independence in Lemma~\ref{anchor:lem:axial}. Induction proves full row rank of every predesign, nonzero $gN_H$ at each clean node, and \eqref{anchor:eq:angle}.

Separately, Lemma~\ref{anchor:lem:residual} proves that the auxiliary row recursion has nonzero return $a_S^TN_G$ at every clean node. Intersect these two probability-one events. Induction in \eqref{anchor:eq:pure} now identifies both auxiliary rules with the actual hybrid. This step does not condition the row lemma on the column spaces, or vice versa.

The top annihilator removes at least the image $S_h$ retained by $I$, so
\[
 \|gN_{H_P}\|^2\le\|g\|^2\frac{t(S_h)}{1+t(S_h)}
                        \le\|g\|^2t(S_h).
\]
Apply \eqref{anchor:eq:angle}, then average $\E\|g\|^2=q$. The extra $K$ image is independent Gaussian and has dimension $r$, so $H_P$ has rank $2r$ and $gN_{H_P}\ne0$ almost surely. The same conclusions follow by generic subspace transversality conditional on the $I$ image and $g$.
\end{proof}

\paragraph{Discontinuity at zero noise.}
The theorem concerns every fixed \emph{positive} $\delta$. At $\delta=0$, exact row containment makes $a_S^TN_G=0$ and the column rule takes its canonical zero-sample branch. The activated column hierarchy need not converge to that zero-noise hierarchy. Exact recovery on HSS inputs and \eqref{anchor:eq:anchor-energy} are therefore consistent.

\section{The one-sided row limit is exactly containing}\label{anchor:sec:limit}
We next let $\delta\downarrow0$ at fixed $h,\epsilon$. This is a proof device; the final inputs will have positive deterministic $\delta_h$.

\begin{lemma}[One-sided physical row limit]\label{anchor:lem:rowlimit}
For each fixed finite tree, the physical row spaces inside $I$ have limits as $\delta\downarrow0$ through positive values, almost surely. The limits depend only on $\Omega$, contain the restrictions of $u$, and have dimension $r$. Their compressed transpose designs have full row rank almost surely. The column hierarchy stays the activated hierarchy of Theorem~\ref{anchor:thm:anchor}, not the zero-noise hierarchy.
\end{lemma}
\begin{proof}
At a physical leaf, the row sample has the form $ab^T+\delta W$ with fixed nonzero $a=\mathbf1_{2r}$ and $b=gN_{\Omega_S}$ under conditioning on $\Omega_S,g$. In nullspace coordinates, $W$ is standard Gaussian. The leading singular direction converges to $a$. The remaining $r-1$ selected directions converge to the leading $r-1$ left directions of
\[
 (\Id-P_a)W(\Id-P_b).
\]
To see the last statement, expand the sample Gram matrix in the splitting $\ran a\oplus a^\perp$ and eliminate its separated top eigenvalue. Its lower Schur complement, divided by $\delta^2$, converges to the Gram matrix of the displayed projected noise. In orthonormal coordinates this is a $(2r-1)\times(q-2r-1)$ standard Gaussian matrix. Its relevant singular values are positive and distinct almost surely. The leaf projector and its canonical frame consequently have limits, all measurable in $\Omega$ alone. The limiting projector contains $a$.

At a clean node, suppose the child limiting spaces have this property. Their block-diagonal template $W_0$ is independent of $\Psi$ and contains $u_S$. Thus $W_0^T\Psi_S$ has full row rank almost surely, and the physical lifted vector converges to
\[
 W_0(W_0^T\Psi_S)^{\dagger T}\Psi_S^Tu_S=u_S.
\]
The completion in \eqref{anchor:eq:canonical} appends the first $r-1$ child coordinate directions. Their physical versions are supported in the left child; $u_S$ has a nonzero right-child restriction, so their union with $u_S$ has dimension $r$. The completed space therefore converges continuously.

More explicitly, the first $r-1$ physical frame vectors at every clean node are inherited from its leftmost physical leaf. They have limits depending only on $\Omega$. The limiting completed space is their span together with $u_S$. This observation removes any possible orientation ambiguity at coordinate-frame pivots. It proves the claimed independence of $\Psi$ and closes the induction. Full-rank and continuity events are intersected over finitely many nodes.
\end{proof}

At $P$, let $C_K=[e_n,\ldots,e_{n+r-1}]$ and define
\[
 F=[a,b],\qquad a=\Psi_I^Tu,\quad b=\Psi_K^TC_K,
 \qquad \Gamma=\begin{bmatrix}\Omega_{j+1,:}\\\vdots\\\Omega_{j+r,:}\end{bmatrix}.
\]
Here $F$ is a standard $q\times(r+1)$ Gaussian matrix marginally. Let $N$ be an orthonormal basis for $\ker H_P$, of dimension $d=q-2r$, and set
\begin{equation}\label{anchor:eq:limit-sample}
 w=gN,\quad Z=\Gamma N,\qquad
 M=a w+\frac\epsilon{\sqrt r}\,bZ.
\end{equation}
Conditional on $\Psi,\Omega_I,\Omega_K,g$, the matrix $Z$ is standard $r\times d$ Gaussian. This remains true although $N$ depends on both probe streams. The unused rows $\Gamma$ are independent of all those variables.

In the one-sided limit, the row template contains the orthonormal physical frame
\[
 Q_*=[u,C_K].
\]
In its compressed coordinates, call this isometry $Q_c$. It obeys
\begin{equation}\label{anchor:eq:lift-F}
 G_P^TQ_c=F,\qquad (G_P^\dagger)^TF=Q_c.
\end{equation}
Hence $\Pi_{G_P}J_P=J_P$, and \eqref{anchor:eq:limit-sample} is the actual sketch-space sample (up to the harmless right isometry $N^T$). Every selected leading-$r$ left space lies in $\ran F$; its physical lift is obtained by applying $Q_cF^\dagger$. The conditioning of the remaining directions of $G_P$ does not enter this restricted lift.

\section{Displacement of the comparator and failure of the learner bound}
\begin{lemma}[Rank-$r$ displacement by a weak sibling]\label{anchor:lem:displacement}
Let $\ell_h^+(\epsilon)$ be the squared loss of the physical unit vector $u$ at $P$ in the one-sided limit of Lemma~\ref{anchor:lem:rowlimit}. Set
\begin{equation}\label{anchor:eq:eps}
 \epsilon_h=\alpha^{h/4}.
\end{equation}
Then $\ell_h^+(\epsilon_h)\to1$ in probability and in expectation.
\end{lemma}
\begin{proof}
Let $M_0=(\epsilon_h/\sqrt r)bZ$ and $E_M=aw$ in \eqref{anchor:eq:limit-sample}. Put $P_0=\Pi_{\ran b}$ and let $P_M$ be a leading-$r$ left projector of $M=M_0+E_M$. Since $\rank M_0=r$,
\[
 \op{(\Id-P_M)M_0}\le\op{(\Id-P_M)M}+\op{E_M}
 \le2\op{E_M}.
\]
Consequently the largest principal-angle sine satisfies
\begin{equation}\label{anchor:eq:sin}
 \op{P_M-P_0}\le
 \min\left\{1,\frac{2\sqrt r\,\|a\|\|w\|}
 {\epsilon_h\sigma_{\min}(b)\sigma_{\min}(Z)}\right\}.
\end{equation}
The nonzero noise singular value bound used here is
$\sigma_r(bZ)\ge\sigma_{\min}(b)\sigma_{\min}(Z)$.

By Theorem~\ref{anchor:thm:anchor},
\[
 \E\frac{\|w\|^2}{\epsilon_h^2}\le q\mu\alpha^{h/2}\longrightarrow0.
\]
The random variables $\|a\|$, $\sigma_{\min}(b)^{-1}$, and $\sigma_{\min}(Z)^{-1}$ form tight families: their marginal laws have fixed dimensions and the matrices are full rank almost surely. The same is true of $\kappa(F)$. No independence of these quantities and $w$ is required to conclude that the right side of \eqref{anchor:eq:sin}, multiplied by $\kappa(F)$, tends to zero in probability.

For any unit coefficient vector $c$ whose image $Fc$ lies in $\ran P_M$,
\[
 \dist(c,\operatorname{span}\{e_2,\ldots,e_{r+1}\})
 \le\op{F^\dagger}\dist(Fc,\ran b)
 \le\kappa(F)\op{P_M-P_0}.
\]
The physical frame $Q_c$ in \eqref{anchor:eq:lift-F} is an isometry. Thus the retained physical plane becomes orthogonal to its first direction, which is $u$, and $\ell_h^+\to1$ in probability. Since $0\le\ell_h^+\le1$, the convergence also holds in expectation.
\end{proof}

\begin{theorem}[Actual noisy-input basis-risk counterexample]\label{anchor:thm:main}
Fix $k\ge1$, $r=2k$, and $q=16k$. Use the unchanged annihilate-first hybrid specified in Section~1. There are positive deterministic numbers
\[
 0<\delta_h\le\epsilon_h^2,\qquad
 \epsilon_h=\left(\frac{12k+1}{28k}\right)^{h/4},
\]
such that the deterministic inputs \eqref{anchor:eq:A} satisfy
\begin{equation}\label{anchor:eq:main}
 \boxed{\E\dist_F(A_h,\SSS_{\rm hyb}(A_h))^2\longrightarrow1,\qquad
        L_h\OPT_k(A_h)\longrightarrow0.}
\end{equation}
One may take $L_h=h+4$ for the original tree with leaf size $2k$. In particular, no finite constant $C_b$ makes \eqref{anchor:eq:false} true for all fixed inputs and depths. The same construction works at any fixed $r\ge2,q>4r$, with \eqref{anchor:eq:constants}; merely increasing the fixed width does not cure this failure.
\end{theorem}
\begin{proof}
For each fixed $h,\epsilon>0$, let $\ell_h(\delta,\epsilon)$ be the actual squared physical row loss of $u$ at $P$. Theorems~\ref{anchor:thm:anchor} and Lemma~\ref{anchor:lem:rowlimit} imply
\[
 \ell_h(\delta,\epsilon)\longrightarrow\ell_h^+(\epsilon)
 \quad\text{almost surely as }\delta\downarrow0.
\]
The leading-$r$ singular subspace at $P$ is continuous on the full-rank, strict-gap event, which has probability one: conditional on the other variables, $\Gamma N$ has a Gaussian density, $w\ne0$, and $F$ has full column rank. The restricted lift is continuous there as well. A countable sequence of positive $\delta$ suffices throughout this argument.

Bounded convergence permits a deterministic choice $\delta_h\in(0,\epsilon_h^2]$ such that
\[
 \left|\E\ell_h(\delta_h,\epsilon_h)-\E\ell_h^+(\epsilon_h)\right|
 \le\frac1{h+1}.
\]
For example, choose a sufficiently far member of a predetermined positive rational sequence. The choice uses the expectation as a mathematical existence criterion; it is not made after the realized probes are seen. No closed-form rate for $\delta_h$ is claimed.

The exterior column $j$ of $A_h$ on $P$ is exactly $u$. Thus \eqref{anchor:eq:cutrange} gives, for every realized learned space,
\begin{equation}\label{anchor:eq:rowlower}
 \dist_F(A_h,\SSS_{\rm hyb})^2\ge\ell_h(\delta_h,\epsilon_h).
\end{equation}
There is no residual subtraction in this bound: neither perturbation changes that exterior column. The lower expectation tends to one by Lemma~\ref{anchor:lem:displacement}; the upper bound $\dist_F(A_h,\SSS_{\rm hyb})^2\le\F{A_h}^2\to1$ proves the first limit.

The second follows from
\[
 L_h\OPT_k(A_h)\le(h+4)(\alpha^h+\alpha^{h/2})\longrightarrow0.
\]
The optimum is positive for every $\delta_h>0$: at a coarse-leaf cut inside $I$, the perturbation contributes an invertible $2r$-row sibling block, so the exterior rank is $2r>k$. Since $\HH_k$ is closed, its distance is nonzero. This also justifies interpreting \eqref{anchor:eq:main} as a diverging approximation ratio.
\end{proof}

\paragraph{No contradiction with the first adaptive-merge theorem.}
The earlier physical first-merge estimate controls newly discarded comparator loss by inherited loss and the current crossing residual. It is valid when the child templates are stream-separated. In this example the column comparator is zero on the clean subtree, so retaining a harmful column direction has zero immediate comparator cost. At later nodes the same direction affects the outgoing annihilator. The stream-separation premise no longer holds there. A small first-step comparator loss does not prevent this latent information from accumulating in unused directions.

\section{A direct-sum version and the scope of fallback protection}
\begin{corollary}[Large global rank does not remove the obstruction]\label{anchor:cor:directsum}
Let $T$ be a fixed power of two and take the normalized direct sums
\[
 \overline A_h=T^{-1/2}\diag(A_h,\ldots,A_h),\qquad
 \overline B_h=T^{-1/2}\diag(B_h,\ldots,B_h).
\]
On the aligned larger tree, $\overline B_h\in\HH_1$, has global rank $T$, and has $T$ singular values equal to $T^{-1/2}$. The expected distance of $\overline A_h$ to its actual hybrid space still tends to one, while $(L_h+\log_2T)\OPT_k(\overline A_h)\to0$.
\end{corollary}
\begin{proof}
All computations inside a diagonal block use only its own probe rows; homogeneity removes the normalization from the basis choices. Each copy has the same marginal law as before, and its distinguished-cut loss is divided by $T$. The disjoint cut-column errors add. Averaging proves the same lower bound; the Frobenius upper bound and comparator estimate are unchanged. At nodes above the copies the comparator has zero exterior blocks, proving its HSS rank claim. Its singular values follow from the disjoint rank-one summands.
\end{proof}

For $T>4k$, no globally rank-at-most-$4k$ candidate has small residual on this family. The triangle inequality for distance to the rank-$4k$ set gives
\begin{equation}\label{anchor:eq:globalfallback}
 \dist_F(\overline A_h,\{\rank\le4k\})
 \ge\sqrt{1-4k/T}-\sqrt{\delta_h^2+\epsilon_h^2}.
\end{equation}
Every original-leaf-diagonal candidate misses all the unit-energy comparator entries. In fact all entries of this input are off those diagonal leaf blocks. Therefore a selector among a hybrid-space candidate, a rank-$4k$ global candidate, and a leaf-diagonal candidate, \emph{before any further nonlinear HSS projection}, has limiting expected squared error at least $1-4k/T$. To verify the expectation statement, lower-bound the hybrid error by the average of the bounded cut losses in the proof above; that average converges to one in probability and expectation. Then take the minimum of it and the two deterministic lower bounds.

This is a statement about candidate residuals and the fallback benchmarks. It does not by itself prove failure of an additional best-HSS projection of the selected matrix, which may change the bases. The direct-sum construction nevertheless shows that the bad basis mechanism is not restricted to matrices well approximated at small \emph{global} rank.

\section{An un-nullified companion preserves the displaced signal}\label{anchor:sec:backup}
Use the same physical-leaf rules and the same original sketches, but at internal nodes take a leading-$r$ left sketch space of
\begin{equation}\label{anchor:eq:raw}
 \widetilde M_U=\Pi_GJ_S(A),\qquad
 \widetilde M_V=-\Pi_HJ_S(A)^T,
\end{equation}
then lift and use the same canonical full-range completion. This is an explicitly specified companion; it is not silently substituted into Theorem~\ref{anchor:thm:main}.

On the clean subtree $I$ with positive $\delta$, its row and column spaces coincide with the hybrid's: both rank-one full-range rules give \eqref{anchor:eq:surrogate}. Inside $K$ both row hierarchies contain the coordinate noise space exactly. The column hierarchies in $K$ can differ, but the row sample \eqref{anchor:eq:raw} does not use their annihilator.

In the one-sided limit, the companion's sample at $P$ is
\[
 \widetilde M=ag+\frac\epsilon{\sqrt r}b\Gamma,
\]
not \eqref{anchor:eq:limit-sample}. The genuine anchor has not been removed.

\begin{proposition}[A physical rescue at the failing cut]\label{anchor:prop:rescue}
For $r=2,q=16$, let $\widetilde\ell_h^+(\epsilon)$ be the one-sided-limit physical row loss of $u$ at $P$ for the companion. Uniformly in $h$,
\begin{equation}\label{anchor:eq:rescue}
 \boxed{\E\widetilde\ell_h^+(\epsilon)\le\frac{92}{9}\epsilon^2.}
\end{equation}
This is a cut-specific comparison, not an all-depth noisy-input theorem for the companion.
\end{proposition}
\begin{proof}
Let $T=N_\Gamma(gN_\Gamma)^\dagger$, with an orthonormal nullspace basis understood, so $gT=1$ and $\Gamma T=0$. Thus $\widetilde M T=a$. Let $P_M$ be its leading-two sketch projector. By \eqref{anchor:eq:lift-F}, projection of the lifted $P_Ma$ lies in the retained physical plane. Therefore
\[
 \widetilde\ell_h^+
 \le\op{F^\dagger}^2\sigma_3(\widetilde M)^2\|T\|^2
 \le\frac{\epsilon^2}{2}\F{F^\dagger}^2\F b^2\F\Gamma^2\|T\|^2.
\]
Here $F=[a,b]$ is standard $q\times3$ Gaussian and is independent of $(g,\Gamma)$. Conditional on $\Gamma$, $\|T\|^2$ is inverse chi-square with $q-2$ degrees of freedom. Hence
\[
 \E[\F\Gamma^2\|T\|^2]=\frac{2q}{q-4}.
\]
The Gaussian inverse-Frobenius identity \cite{HMT}, right-column rotational invariance, and radius-direction independence give
\[
 \E[\F{F^\dagger}^2\F b^2]
 =\frac23\E[\F{F^\dagger}^2\F F^2]
 =\frac{2(3q-2)}{q-4}.
\]
For the last equality, write $F=RU$ with $R^2\sim\chi^2_{3q}$. Since $\E R^{-2}=1/(3q-2)$ and $\E\F{F^\dagger}^2=3/(q-4)$, homogeneity yields the displayed product moment. Multiplication gives
\[
 \E\widetilde\ell_h^+\le
 \frac{2q(3q-2)}{(q-4)^2}\epsilon^2
 =\frac{92}{9}\epsilon^2\quad(q=16).
\]
The restricted lift through $F$ is essential; no unweighted estimate on the full adaptive inverse was used.
\end{proof}

\section{A nested union at the same training-query count}\label{anchor:sec:union}
Run the hybrid and companion on the same $\Omega,\Psi,Y,Z$. At each node take the union of their physical row spaces and separately the union of their physical column spaces. Orthogonalize each union, retaining all its directions. The dimensions are at most $2r=4k$.

\begin{proposition}[Nested union and pointwise domination]\label{anchor:prop:union}
The unions are nested and admit a coefficient space $\SSS_\cup$ satisfying
\begin{equation}\label{anchor:eq:union}
 \SSS_{\rm hyb}+\SSS_{\rm raw}\subseteq\SSS_\cup,\qquad
 \boxed{\dist_F(A,\SSS_\cup)^2
 \le\min\{\dist_F(A,\SSS_{\rm hyb})^2,
          \dist_F(A,\SSS_{\rm raw})^2\}.}
\end{equation}
Training uses the same $32k$ products. With the independent rank-$4k$ refit, total access is at most $64k$ products. The union statement itself has no missing probabilistic premise; an all-input bound on its distance to $A$ remains unproved.
\end{proposition}
\begin{proof}
Each parent's original row space is contained in the direct sum of its own child spaces, hence in the direct sum of the child unions. Taking the two parent spaces together preserves this containment. The column proof is identical. Every coefficient representation from either original hierarchy remains valid after embedding its bases into the union bases. Linearity then proves the space inclusion and the distance inequality.

If the uniform-rank refit convention needs leaves of size $8k$, coarsen to those leaves and allow their full diagonal blocks; this only enlarges the union space. All omitted finer cuts have physical side size at most $4k$, so the resulting matrices still have HSS rank at most $4k$ on the original tree. Apply \eqref{anchor:eq:refit} with $s=4k$, using $32k$ fresh products in addition to the $32k$ common training products. Small dimensions can be read exactly within the same budget.
\end{proof}

If a future theorem proves $\E\dist_F(A,\SSS_\cup)^2\le C_bL\OPT_k(A)$, the refit of Chapter~\ref{ch:uniform} and the query-free exact offline projection of Chapter~\ref{ch:bicriteria} onto $\HH_k$ complete the general target. Indeed, for a known output $C$ and its best HSS-$k$ approximation $\widehat B$,
\[
 \F{A-\widehat B}\le2\F{A-C}+\sqrt{\OPT_k(A)}.
\]
This final minimization is a computation-free query reduction, not a polynomial-time claim. Neither the union construction nor the local rescue estimate proves the premise. The companion may have other instabilities; the purpose of retaining both hierarchies is to avoid replacing one unproved universal claim by another.

\section{Numerical checks and a zero-sample implementation issue}\label{anchor:sec:checks}
The included script evaluates the deterministic construction in factored form. It never forms a large dense matrix. It evaluates mathematically clean rank-one samples and zero orientations directly, instead of obtaining them by subtracting nearly equal floating-point arrays. This specialization is an algebraic diagnostic on a known test input, not an additional query available to the learner.

There are two separate run types. A positive-$\delta$ run is the actual input \eqref{anchor:eq:A}. A run with the script parameter \texttt{delta=None} evaluates the proved one-sided limit of Section~\ref{anchor:sec:limit}; it is \emph{not} a run on the exact input $\delta=0$.

For $k=1$, $\delta=10^{-3}$, and $\epsilon=(13/28)^{h/4}$, the following are medians over three predetermined probe seeds. The reported quantity is the squared physical loss of the unit comparator column at $P$, so it directly lower-bounds the full coefficient-space error for the hybrid. The companion column reports the cut-specific un-nullified comparison, not its full-tree error.
\begin{center}
\begin{tabular}{rrrrr}
\toprule
$h$&$N$&$\epsilon$&Hybrid cut loss&Un-nullified cut loss\\
\midrule
3&256&0.562457&0.162790&0.038082\\
5&1,024&0.383250&0.034542&0.000858\\
7&4,096&0.261141&0.990766&0.000090\\
9&16,384&0.177937&0.983456&0.000501\\
11&65,536&0.121244&0.995294&0.000082\\
\bottomrule
\end{tabular}
\end{center}
The finite sample is not monotone at small depths and is not an expectation estimate. The proof of Theorem~\ref{anchor:thm:main}, including the deterministic selection of $\delta_h$, does not rely on this table. Additional runs compare four positive $\delta$ values with the one-sided limit on shared probes. Twelve dense original-product checks, at $N=32,64,128,256$, agree with the factored cut-loss calculations to within $6.7\times10^{-16}$ in absolute squared loss. The local union has dimension three in all fifteen positive-noise grid runs; its median cut losses range from $9.6\times10^{-8}$ to $2.2\times10^{-7}$. These are numerical checks of this family, not a general risk theorem for the union.

\paragraph{Why earlier generic floating-point checks can mask this mechanism.}
The inherited checker classified rank relative to the largest singular value of the \emph{computed sample itself}. A mathematically zero sample can contain cancellation roundoff; a purely relative test can treat that roundoff as genuine singular directions and replace the intended canonical zero-space completion. In the sibling $K$ of this construction, that changes the zero incoming space and can also suppress the competing noise anchors. It therefore changes the very mechanism under test. This is a limitation of the numerical checker, not a correction to the exact-arithmetic first-merge theorem. The factored checks avoid that ambiguity rather than interpreting a tolerance experiment as a proof.

\section*{Research status after this chapter}
The status of the unchanged hybrid has changed: the desired all-depth basis-distance inequality is false, not merely missing a conditioning argument. The retained conclusions from the earlier work remain compatible with this fact: exact fixed-input recovery, first-layer and first-adaptive-merge physical estimates, and all-depth sketch-tail budgets do not control the unused directions that drive the trap.

The new conditional Gaussian residual calculation is an all-depth actual-law tool on this clean-cut family. It proves persistence of a nonzero innovation without assuming that a learned template is independent of the probes. In this chapter it supports a negative result, rather than a universal stability theorem.

The explicit next candidate is the shared-query nested union in Proposition~\ref{anchor:prop:union}. It retains the good directions of both rules and removes the need to choose between them during training. Establishing its all-input residual-relative quality remains the mathematical task. The general $O(k)$-query approximation target has neither been proved nor disproved here.

\endgroup

\part{Regression-tree stability and source-coupled inverse budgets}

% ===== TANGENT =====
\begingroup
\let\E\relax
\newcommand{\E}{\mathbb E}
\let\PP\relax
\newcommand{\PP}{\mathbb P}
\let\R\relax
\newcommand{\R}{\mathbb R}
\let\F\relax
\newcommand{\F}[1]{\lVert #1\rVert_F}
\let\op\relax
\newcommand{\op}[1]{\lVert #1\rVert_{\mathrm{op}}}
\let\ran\relax
\newcommand{\ran}{\operatorname{range}}
\let\rank\relax
\newcommand{\rank}{\operatorname{rank}}
\let\diag\relax
\newcommand{\diag}{\operatorname{blockdiag}}
\let\dist\relax
\newcommand{\dist}{\operatorname{dist}}
\let\OPT\relax
\newcommand{\OPT}{\operatorname{OPT}}
\let\HH\relax
\newcommand{\HH}{\mathcal H}
\let\SSS\relax
\newcommand{\SSS}{\mathcal S}
\let\Id\relax
\newcommand{\Id}{\mathrm I}
\let\tr\relax
\newcommand{\tr}{\operatorname{tr}}
\chapter{Exact two-view capture and the tangent regression tree}\label{ch:tangent}
\begin{scopebox}
\textbf{Update, 20 September 2026.} The exact two-view range theorem and fixed-height tangent result below retain their original hypotheses. They concern the specified trapping family and separately trained views, not every noisy parent with unequal templates. The later reached clean-parent witness shows why a union taken only after separate propagation cannot be assigned a general no-new-loss principle. The synchronized replacement has a different exact-capture argument and a clean-parent conclusion with an explicit visibility hypothesis. Neither its completion nor the present tangent theorem establishes general noisy stability or a height-uniform perturbation remainder.
\end{scopebox}

\begin{scopebox}
The exact positive-noise two-view range theorem and the fixed-reference tangent calculation are distinct. The expectation--limit gap explicitly left at the tangent stage is closed in Chapter~\ref{ch:nonlinear}; no height-uniform nonlinear remainder is inferred.
\end{scopebox}
\thispagestyle{plain}
\begin{chaptersynopsis}
We advance the nested union of an annihilate-first HSS learner and its un-nullified companion. On the positive-noise exterior-anchor-trapping family, the two actual rank-$r$ row truncations at the distinguished cut jointly span the full $(r+1)$-dimensional lifted sample range almost surely. This is an exact finite-noise statement, stronger than the preceding cut-specific $O(\epsilon^2)$ comparison for the companion alone. The union removes the weak sibling's displacement effect at that cut, but does not remove errors inherited from the children.

A separate reusable theorem controls a fixed-reference Gaussian regression tree. For deterministic nested rank-$r$ reference spaces, arbitrary deterministic orthogonal leaf errors of total squared norm $\beta$, and width $q>2r+1$, its accumulated physical error has a depth-independent expected bound $C_\infty\beta$, where
\[
 C_\infty=\frac{(q-1)(2q-3r-1)}{(q-r-1)(q-2r-1)}.
\]
The proof retains a regression register and a query-space innovation. Conditional innovation covariance decays geometrically; an exact orthogonal energy identity sums the resulting gains without separating dependent expectations.

This regression tree is the actual first variation of a clean-subtree range-containing learner about an independently determined exactly containing reference hierarchy. For the trapping family, $r=2,q=16$, the union's distinguished-cut loss $\ell_h(\delta,\epsilon)$ has an almost-sure one-sided quadratic coefficient $Q_h$, with $\E Q_h\le75/572$ uniformly in depth and in fixed $\epsilon>0$. This is a mean-square tangent theorem, not a finite-noise expected-risk theorem: interchange of the small-noise limit and an untruncated expectation is not asserted. The all-input best-HSS approximation theorem for the union remains unresolved.
\end{chaptersynopsis}


\section{Target, maintained algorithm, and scope}
Fix a balanced binary coordinate tree of depth $L$, with original leaf size $2k$. Let $\HH_k$ denote matrices whose outgoing and incoming exterior blocks at every nonroot node have ranks at most $k$, and set
\[
 \OPT_k(A)=\min_{B\in\HH_k}\F{A-B}^2.
\]
The target is $O(k)$ products with $A,A^T$, followed by an HSS-rank-$k$ output with expected squared error at most $CL\OPT_k(A)$. The public fresh-level benchmark uses $O(kL)$ products and has an $O(L)$ expected approximation factor \cite{AmselHSS}. This chapter concerns a reused-sketch route, not an extension of that paper's independence assumptions.

Use $r=2k$, coarse physical leaf size $2r=4k$, and independent standard Gaussian $N\times q$ probes $\Omega,\Psi$, with $q=16k=8r$. Query $Y=A\Omega$, $Z=A^T\Psi$, and form
\begin{equation}\label{tangent:eq:J}
 J_I(A)=\Psi_I^TY_I-Z_I^T\Omega_I=\Psi^T[P_I,A]\Omega.
\end{equation}
For a full-row-rank matrix $D$, write $\Pi_D=D^\dagger D$ and $N_D=\Id-\Pi_D$; the Moore--Penrose definitions extend these expressions to singular matrices.

Both hierarchies have physical-leaf row and column spaces given by leading-$r$ left spaces of $Y_IN_{\Omega_I}$ and $Z_IN_{\Psi_I}$. At an internal node their own child templates give $G=W_U^T\Psi_I$ and $H=W_V^T\Omega_I$. The hybrid uses
\[
 M_{\rm h}=\Pi_GJ_I(A)N_H,
\]
whereas the companion uses
\[
 M_{\rm c}=\Pi_GJ_I(A).
\]
Each selects a leading-$r$ left space in query coordinates, lifts through $(G^\dagger)^T$, and orthonormalizes. The transpose counterparts define column spaces. When a sample has rank at most $r$, its full lifted physical range is extended canonically: project coordinate vectors in order, append coordinate directions as needed, and canonicalize the resulting projector. In particular, a rank-one local vector not supported on the first $r-1$ coordinates is completed to $\operatorname{span}\{v,e_1,\ldots,e_{r-1}\}$.

The algorithm under consideration takes physical row and column unions of the two \emph{complete, separately propagated} hierarchies. Their union is not silently fed back into either original learner. The union ranks are at most $2r=4k$.

\paragraph{Inherited interfaces, proved in Chapter~\ref{ch:uniform}.}
The fixed-space theorem of Chapter~\ref{ch:uniform} gives, for a rank-$s$ nested coefficient space chosen independently of fresh refit probes,
\begin{equation}\label{tangent:eq:refit}
 \E[\F{A-\widehat C}^2\mid\SSS]\le1607\dist_F(A,\SSS)^2
 \quad\text{using }8s\text{ products}.
\end{equation}
It has no basis-learning conclusion. The nested union contains both original coefficient spaces, so its distance is no larger than either individual distance. These facts yield $32k$ shared training products plus $32k$ independent rank-$4k$ refit products, but the all-input error guarantee still needs
\begin{equation}\label{tangent:eq:missing}
 \E\dist_F(A,\SSS_\cup(A))^2\le C_bL\OPT_k(A).
\end{equation}
Nothing below claims \eqref{tangent:eq:missing}. All Gaussian subspace assertions are pointwise for an input fixed before the probes.

\section{Two truncations can recover the full lifted range}
\begin{theorem}[Two-view hyperplane completion]\label{tangent:thm:twoview}
Let $r\ge2$ and $d=r+1$. Fix $F=[a,b]\in\R^{q\times d}$ with full column rank, where $b$ has $r$ columns. Let $\Pi$ be an orthogonal projector, $N=\Id-\Pi$, with
\[
 \rank N\ge r+1,\qquad \rank\Pi\ge1.
\]
Fix a row $g$ with $gN\ne0$ and $g\Pi\ne0$, and a number $\epsilon>0$. Let $\Gamma\in\R^{r\times q}$ be standard Gaussian and independent of these quantities. Put
\[
 M_0=ag+\epsilon b\Gamma,\qquad M_1=M_0N.
\]
Almost surely, their leading-$r$ left spaces $\mathscr T_0,\mathscr T_1$ are uniquely defined and satisfy
\begin{equation}\label{tangent:eq:span}
 \boxed{\mathscr T_0+\mathscr T_1=\ran F.}
\end{equation}
The same assertion holds conditionally when $F,g,\Pi$ are random and an independent $\Gamma$ remains after conditioning, provided the displayed hypotheses hold almost surely. This is not a statement for arbitrary adaptive projectors without that independence.
\end{theorem}
\begin{proof}
Work inside the $d$-dimensional Euclidean space $\ran F$. Both matrices have rank $d$ almost surely. For $M_1$, rotate query coordinates so that $N$ is a coordinate projector. Its first coefficient row is the fixed nonzero row $gN$ and its other $r$ coefficient rows have independent Gaussian entries. Since $\rank N\ge d$, full row rank follows. The same argument applies to $M_0$.

Their nonzero singular values are distinct almost surely. One detailed justification is to consider the coefficient Gram matrix. With the first row fixed nonzero and at least $d$ columns, the remaining Gaussian rows generate a density on positive definite Gram matrices whose first diagonal entry is fixed. Every such Gram matrix is attainable: resolve the remaining rows along the first row and in its orthogonal complement, and realize the Schur complement there. After congruence by any fixed invertible coordinate matrix for $F$, this set contains a matrix with simple spectrum. For example, scale a simple-spectrum positive definite matrix to satisfy the one positive linear constraint. The discriminant polynomial is therefore nonzero, so its zero set has probability zero.

Condition on $\Gamma N$ and define
\[
 C_N=M_1M_1^T,\qquad C_\Pi=(M_0\Pi)(M_0\Pi)^T.
\]
The remaining coordinates $\Gamma\Pi$ are independent Gaussian. For any fixed nonzero $v\in\ran F$, the event $C_\Pi v\in\operatorname{span}\{v\}$ is the zero set of a nonzero polynomial in those remaining coordinates. To verify nontriviality, first suppose $b^Tv\ne0$. Since $r\ge2$, choose $w\in\ran b$ with $w^Tv\ne0$ and $w$ not parallel to $v$. Vary $b\Gamma\Pi$ along $twh^T$ for some nonzero $h\in\ran\Pi$. The quadratic coefficient of $C_\Pi v$ is a nonzero multiple of $w(w^Tv)$, not parallel to $v$. If instead $b^Tv=0$, then $a^Tv\ne0$. The term linear in $\Gamma\Pi$ is
\[
 \epsilon(a^Tv)b\Gamma\Pi g^T,
\]
which varies in $\ran b$ and is nonzero in a direction orthogonal to $v$, because $g\Pi\ne0$.

There are only $d$ eigenlines of the now fixed simple-spectrum $C_N$. None is an eigenline of $C_\Pi$ almost surely. If the two leading-$r=d-1$ spaces were equal, their one-dimensional complements inside $\ran F$ would coincide. That common line would be an eigenline of $C_N$ and $C_N+C_\Pi$, hence of $C_\Pi$, a contradiction. Two distinct hyperplanes span $\ran F$, proving \eqref{tangent:eq:span}.
\end{proof}

\begin{remark}[No conditioning estimate is hidden in the span statement]
The two spaces may be arbitrarily close. The theorem retains their exact union, not an angle-thresholded numerical union. It gives neither a lower bound on the angle nor a bound on coefficients used to express a vector as a sum of the two spaces. For the exact-arithmetic query model, orthonormalizing the full physical union and applying an independent fixed-space refit does not require an angle-dependent refit constant.
\end{remark}

\section{An exact finite-noise rescue at the trapping cut}\label{tangent:sec:family}
Here $r\ge2,q>4r$ may be general; the project parameters are $r=2k,q=8r$. For $h\ge0$, let
\[
 n=4r\,2^h,\quad N=4n,\quad I=[0,n),\quad K=[n,2n),\quad P=I\cup K,\quad j=2n.
\]
Indices are zero-based. Set $u=n^{-1/2}\mathbf1_I$ and
\begin{equation}\label{tangent:eq:family}
 \begin{split}
 B&=ue_j^T,\\
 A(\delta,\epsilon)&=B+\delta E_0+\epsilon E_1,\qquad \delta,\epsilon>0,\\
 E_0&=\mathcal P_I/\sqrt n,\qquad
 E_1=\frac1{\sqrt r}\sum_{a=1}^r e_{n+a-1}e_{j+a}^T.
 \end{split}
\end{equation}
The permutation $\mathcal P_I$ swaps each pair of adjacent coarse leaves of size $2r$ and is zero outside $I$. The comparator and the two unit-Frobenius perturbations have disjoint supports. Thus $\OPT_k(A)\le\delta^2+\epsilon^2$. At $K$, the $r$ nonzero singular values of its outgoing cut are $\epsilon/\sqrt r$. Consequently, for the project choice $r=2k$,
\begin{equation}\label{tangent:eq:optlower}
 \OPT_k(A)\ge\epsilon^2/2.
\end{equation}
This lower bound follows by restricting any HSS-$k$ approximation to that cut.

Every node $S\subset I$ of size at least $4r$ has
\[
 J_S(A)=a_Sg,\qquad a_S=\Psi_S^Tu_S,\quad g=e_j^T\Omega.
\]
The following actual-rule properties are the parts of the preceding construction needed here.

\begin{lemma}[Shared row templates and an unused Gaussian block]\label{tangent:lem:shared}
For each fixed $h,\delta,\epsilon>0$, almost surely the hybrid and companion have identical row spaces throughout $I$ and at $K$. Their common row template at $P$ is $W=\diag(\mathcal U_I,C_K)$, where $C_K=[e_n,\ldots,e_{n+r-1}]$ in its physical block. Its design $G=W^T\Psi_P$ has full row rank. The hybrid column design $H$ at $P$ has full row rank, $gN_H\ne0$, and $g\Pi_H\ne0$.

Let $a=\Psi_I^Tu$, $b=\Psi_K^TC_K$, and $\Gamma$ consist of the forward-probe rows $j+1,\ldots,j+r$. Conditional on all other variables used here, $\Gamma$ is standard Gaussian. Neither $G,H,a,b,g$ nor the clean row hierarchy depends on $\epsilon$ or on $\Gamma$.
\end{lemma}
\begin{proof}
At an $I$ leaf, the outgoing nullified sample, after omitting a common scale, is $\mathbf1_{2r}gN_{\Omega_S}+\delta\Omega_{S'}N_{\Omega_S}$. Conditional on $\Omega_S,g$, the shifted second term has full Gaussian density in $2r\times(q-2r)$ coordinates. Its leading-$r$ space misses the fixed signal almost surely for every $\delta>0$. The incoming leaf space uses only $\Psi$ and is independent of the positive value of $\delta$.

At a clean node, define auxiliary full-range row and column recursions by retaining $(G^\dagger)^Ta_S$ and $(H^\dagger)^Tg^T$. The hybrid samples are nonzero rank-one multiples of these vectors exactly when $a_S^TN_G\ne0$ and $gN_H\ne0$; the companion retains the same ranges whenever the vectors are nonzero.

For the row recursion, fix all $\Omega$. At each leaf, the signal image is $x=D^Tc+\eta$, where $D=\mathcal U_S^T\Psi_S$ and $\eta$ is an independent nondegenerate isotropic Gaussian vector. Keep a record of leaf designs and, at each merge, the predesign $D$ and regression observation $Dx$, but not the unobserved residual. Conditional on the record, $x$ is Gaussian. If its pre-observation covariance is $Q_*\succ0$, observing $Dx$ leaves
\[
 Q=Q_*-Q_*D^T(DQ_*D^T)^{-1}DQ_*,\qquad
 \ker Q=\ran D^T,\quad \rank Q=q-2r.
\]
Sibling records use disjoint transpose rows. Right rotations of those rows leave physical choices unchanged, so their predesign row spaces are independent uniform $2r$-subspaces. Since $q>4r$, their intersection is zero almost surely. Their covariance sum is therefore positive definite at the next merge. Induction proves full row rank and $a_S^TN_G\ne0$ at every finite clean level.

Separately fix $\Psi,g$ for the auxiliary column recursion. First clean-node designs are independent Gaussian $2r\times q$ matrices. Their retained query images contain the projection of $g^T$ onto the corresponding predesign image. The image laws are invariant under rotations fixing $g$. At every merge, two independent $r$-subspaces of this form have direct sum not containing $g^T$ almost surely: their projections onto $g^\perp$ are independent uniform $r$-subspaces and $q-1\ge2r$. This proves $gN_H\ne0$ at all finite clean levels; the retained projection itself is nonzero. Intersecting the two probability-one events identifies the auxiliary recursions with both actual clean rules.

Inside $K$, the outgoing physical range is exactly the deterministic coordinate range $C_K$ and is retained by both row learners. The hybrid incoming sample is zero when the row range is exactly containing, so its column basis at $K$ is canonical and deterministic. Its $\Omega_K$ design is independent Gaussian relative to the $I$ image. This proves the stated full-rank and anchor assertions at $P$. The rows $j+1,\ldots,j+r$ of $\Omega$ occur in none of these choices, which proves the last assertions. All induction statements are on a fixed finite tree.
\end{proof}

Let $L=(G^\dagger)^T$, let $E_K\in\R^{2r\times r}$ embed the second child coordinates, and put
\[
 a'=\Pi_Ga,\qquad v=La.
\]
\begin{theorem}[Exact union range for every positive noise level]\label{tangent:thm:exactunion}
For each fixed $h,\delta,\epsilon>0$, the two actual row spaces at $P$ satisfy, almost surely,
\begin{equation}\label{tangent:eq:unionrange}
 \boxed{\mathcal U_{\cup,P}
   =\ran\bigl[Wv,\ C_K\bigr].}
\end{equation}
The right side has dimension $r+1$. Therefore the squared cut loss
\begin{equation}\label{tangent:eq:ell}
 \ell_h(\delta,\epsilon)
   =\|(\Id-P_{\mathcal U_{\cup,P}})u\|^2
\end{equation}
is almost surely equal to a quantity depending on the clean subtree and its designs, but not on $\epsilon$ or the competing anchor $\Gamma$. The statement is for each fixed parameter choice; no simultaneous probability-one claim for all uncountably many parameters is needed.
\end{theorem}
\begin{proof}
The row template is shared by Lemma~\ref{tangent:lem:shared}. The samples at $P$ are
\[
 M_0=a'g+\frac\epsilon{\sqrt r}b\Gamma,\qquad M_1=M_0N_H.
\]
Here $b=G^TE_K$ and $Lb=E_K$. The matrix $F'=[a',b]$ has full column rank almost surely. To see the only nontrivial assertion, condition on the clean $I$ data. Its transmitted design $D_I$ satisfies $D_Ia\ne0$, since it retains a nonzero regression vector for $a$. The extra Gaussian $K$ design is independent. The coefficient of $La$ in the first child is a rational function of that extra design and is not identically zero: choose its row space orthogonal to $\operatorname{span}(\ran D_I^T,a)$, which is possible since $q>4r$. Then that first-child coefficient is $(D_ID_I^T)^{-1}D_Ia\ne0$. Gaussian absolute continuity proves the assertion.

Now condition on $F',g,H$ and the other variables in Lemma~\ref{tangent:lem:shared}. Theorem~\ref{tangent:thm:twoview} applies because $\rank N_H=q-2r\ge r+1$ and both anchor components are nonzero. The two selected query spaces span $\ran F'$. The lift $L$ is injective on $\ran G^T$, so their physical union is exactly $W\ran[La,Lb]$, which is \eqref{tangent:eq:unionrange}.
\end{proof}

\paragraph{What the identity removes.}
The hybrid alone can lose the main signal when $gN_H$ becomes tiny relative to $\epsilon\Gamma N_H$. Equation~\eqref{tangent:eq:unionrange} removes that displacement mechanism completely at this cut. It does \emph{not} say $u\in\mathcal U_{\cup,P}$ for positive $\delta$: the inherited regression vector $WLa$ can still differ from $u$. The next sections control its first variation, rather than discarding this inherited error.

\section{A fixed-reference Gaussian regression tree}\label{tangent:sec:tree}
This section is independent of the HSS construction. Let a balanced binary tree have height $H$ above its physical leaves. At every leaf $\ell$, fix an isometry $U_\ell$ with $r$ columns. At an internal node $I=a\cup b$, fix an orthogonal matrix $[T_I,S_I]\in\R^{2r\times2r}$ and set
\[
 W_I=\diag(U_a,U_b),\qquad U_I=W_IT_I,\qquad D_I^{\rm phys}=W_IS_I.
\]
Ordinarily $T_I$ has $r$ columns; at the root only, it may have any number $s\le2r$. All these spaces are deterministic before a standard Gaussian physical probe $X\in\R^{N\times q}$ is drawn.

At each leaf choose a deterministic error matrix $E_\ell\in\R^{|\ell|\times p}$ with $U_\ell^TE_\ell=0$. Let
\[
 \beta_\ell=\F{E_\ell}^2,\qquad \beta_I=\sum_{\ell\subset I}\beta_\ell,
 \qquad\beta=\beta_{\rm root}.
\]
The number of response columns $p$ is arbitrary. Initialize the physical residual $E_\ell^{\rm out}=E_\ell$. At a parent, let $F_I$ be the stacked child residuals, and define
\begin{equation}\label{tangent:eq:regtree}
 \begin{split}
 G_I&=W_I^TX_I,\qquad
 B_I=S_I^T(G_I^\dagger)^TX_I^TF_I,\\
 E_I^{\rm out}&=F_I-D_I^{\rm phys}B_I.
 \end{split}
\end{equation}
All displayed designs have full row rank almost surely when $q\ge2r$. These are fixed-reference Gaussian designs, not arbitrary adaptive-template designs.

\begin{theorem}[Exact all-depth energy identity]\label{tangent:thm:energy}
Assume $q>2r+1$. Put
\begin{equation}\label{tangent:eq:consts}
 \lambda=\frac{q-2r}{q-r},\quad
 \rho_0=\frac r{q-r-1},\quad
 \tau=\frac{r(q-1)}{(q-r-1)(q-2r-1)},
\end{equation}
and
\[
 C_H=1+\rho_0+\tau\frac{1-\lambda^H}{1-\lambda}.
\]
Let $R_I=U_I^TX_I$ and define the retained regression at the root by
\[
 c_{\rm root}=(R_{\rm root}^\dagger)^TX_{\rm root}^TE_{\rm root}^{\rm out}.
\]
Then
\begin{equation}\label{tangent:eq:energyexact}
 \boxed{\E\F{E_{\rm root}^{\rm out}}^2+\E\F{c_{\rm root}}^2=C_H\beta.}
\end{equation}
In particular,
\begin{equation}\label{tangent:eq:cinf}
 \boxed{\E\F{E_{\rm root}^{\rm out}}^2
 \le C_\infty\beta,\qquad
 C_\infty=\frac{(q-1)(2q-3r-1)}{(q-r-1)(q-2r-1)}.}
\end{equation}
The root-rank variation described above does not alter this identity. The result is uniform in height, reference rotations, leaf error allocation, and the number $p$ of response columns.
\end{theorem}

\subsection{Regression and innovation are different registers}
At a node let $Y_I=X_I^TE_I^{\rm out}$ and decompose it as
\[
 Y_I=R_I^Tc_I+w_I,\qquad
 c_I=(R_I^\dagger)^TY_I,\quad w_I=N_{R_I}Y_I.
\]
At an ordinary parent set $d_I=[c_a;c_b]$ and $w_*=w_a+w_b$, and define
\[
 z_I=(G_I^\dagger)^Tw_*.
\]
Direct multiplication in \eqref{tangent:eq:regtree} gives
\begin{equation}\label{tangent:eq:registers}
 \boxed{\quad
 c_I=T_I^T(d_I+z_I),\qquad
 B_I=S_I^T(d_I+z_I),\qquad
 w_I=N_{G_I}(w_a+w_b).
 \quad}
\end{equation}
The innovation recurrence has no pseudoinverse. The retained regression cannot be dropped: it contributes to later physical readouts.

\begin{lemma}[Conditional covariance and a vanishing cross register]\label{tangent:lem:cov}
At an ordinary rank-$r$ node of height $t$,
\begin{equation}\label{tangent:eq:cov}
 \E[w_Iw_I^T\mid R_I]=\beta_I\lambda^t N_{R_I},
 \qquad \E[c_Iw_I^T\mid R_I]=0.
\end{equation}
Both registers have conditional mean zero. All second moments used here are finite.
\end{lemma}
\begin{proof}
At leaves, the source $X_\ell^TE_\ell$ is independent of $R_\ell$ because $U_\ell^TE_\ell=0$. Its row covariance is $\beta_\ell\Id_q$. Orthogonal Gaussian projections prove \eqref{tangent:eq:cov} at $t=0$ and give finite regression moments.

The physical reference choices are fixed. Under a simultaneous right rotation of all probe rows in a subtree, the physical recursion \eqref{tangent:eq:regtree} is unchanged, $c_I$ is unchanged, and $w_I$ rotates in query coordinates. Conditional on $R_I$, reflection on its orthogonal complement fixes $c_I$ and changes the sign of $w_I$. This proves the second assertion of \eqref{tangent:eq:cov}, once second-moment finiteness is known. Conditional centering follows alternatively by conditioning on all leaf reference designs: the recursion is linear in independent centered Gaussian leaf sources. Those sources can have correlated response columns; the trace covariance is all that is needed.

Suppose the covariance assertion holds for both children. Conditional on $G_I=[R_a;R_b]$, the two subtree states are independent and centered. The last identity in \eqref{tangent:eq:registers} gives
\[
 \E[w_Iw_I^T\mid G_I]
 =\lambda^{t-1}N_{G_I}(\beta_aN_{R_a}+\beta_bN_{R_b})N_{G_I}
 =\beta_I\lambda^{t-1}N_{G_I}.
\]
Under the fixed physical orthogonal split, $R_I=T_I^TG_I$ and $S_I^TG_I$ are independent Gaussian matrices. Conditional on $R_I$, the additional discarded rows span a uniform $r$-subspace of its $(q-r)$-dimensional complement. Thus
\[
 \E[N_{G_I}\mid R_I]=\frac{q-2r}{q-r}N_{R_I}.
\]
This proves the covariance induction. Finiteness at each stage follows from the inverse calculation in the next lemma, together with \eqref{tangent:eq:registers}. Hence the reflection argument is legitimate, closing the joint induction. No conditional Gaussian law for $w_I$ after arbitrary adaptive basis selection is asserted.
\end{proof}

\begin{lemma}[A joint regression cost]\label{tangent:lem:tau}
For independent standard Gaussian $A,B\in\R^{r\times q}$, $G=[A;B]$, and $q>2r+1$,
\begin{equation}\label{tangent:eq:tau}
 \E\F{(G^\dagger)^TN_A}^2
 =\left(1+\frac r{q-r-1}\right)\frac r{q-2r-1}=\tau.
\end{equation}
At an ordinary merge of height $t$,
\begin{equation}\label{tangent:eq:z}
 \E\F{z_I}^2=\tau\lambda^{t-1}\beta_I,
 \qquad \E\langle d_I,z_I\rangle_F=0.
\end{equation}
\end{lemma}
\begin{proof}
For \eqref{tangent:eq:tau}, introduce $\eta=N_A\xi$, where $\xi$ is a standard Gaussian vector independent of $A,B$. Its conditional covariance is $N_A$. Regress it on $G^T$. Resolve the rows of $B$ into the row space of $A$ and its complement. If the fitted coefficient blocks are $x_A,x_B$, the normal equations give
\[
 x_B=(BN_AB^T)^{-1}B\eta,\qquad
 x_A=-(AA^T)^{-1}AB^Tx_B.
\]
Conditional on $A$, the perpendicular part of $B$ is an $r\times(q-r)$ standard Gaussian matrix. The expected squared norm of $x_B$ is $r/(q-2r-1)$. The parallel part of $B$ is independent of that perpendicular part. Integrating it multiplies this value by $\tr((AA^T)^{-1})$ for the $x_A$ contribution. Averaging $A$ gives $r/(q-r-1)$. The Gaussian inverse-Frobenius identity used twice here is classical \cite{HMT}; this calculation keeps the numerator and inverse jointly in the regression.

Conditional on $G_I$, Lemma~\ref{tangent:lem:cov} gives the covariance of $w_a+w_b$ as $\lambda^{t-1}(\beta_aN_{R_a}+\beta_bN_{R_b})$. Apply \eqref{tangent:eq:tau} and row-block symmetry to obtain the first identity in \eqref{tangent:eq:z}. The conditional $c$--$w$ cross registers vanish within each child by Lemma~\ref{tangent:lem:cov}; cross-child terms vanish by conditional independence and centering. This proves the second identity. It is not inferred by separating a marginal inverse moment from a marginal residual moment.
\end{proof}

\begin{proof}[Proof of Theorem~\ref{tangent:thm:energy}]
The reference discarded physical sectors are mutually orthogonal and orthogonal to all leaf error matrices. Equation~\eqref{tangent:eq:regtree} therefore gives the pathwise identity
\begin{equation}\label{tangent:eq:physical}
 \F{E_{\rm root}^{\rm out}}^2=\beta+\sum_{I\text{ internal}}\F{B_I}^2.
\end{equation}
The first two identities in \eqref{tangent:eq:registers} are an orthogonal split, so
\[
 \F{c_I}^2+\F{B_I}^2=\F{d_I+z_I}^2.
\]
Take expectations, use \eqref{tangent:eq:z}, and telescope the retained coefficients through the tree. The leaf regression energy is $\rho_0\beta$. At height $t$, the node budgets $\beta_I$ sum to $\beta$, and the total increment energy is $\tau\lambda^{t-1}\beta$. Adding \eqref{tangent:eq:physical} gives \eqref{tangent:eq:energyexact}. At the last split, any root rank $s\le2r$ still gives exactly the same orthogonal energy identity; no subsequent covariance induction is needed. Finally, sum the geometric series and simplify to obtain \eqref{tangent:eq:cinf}.
\end{proof}

\paragraph{Concrete width.}
At $r=2,q=16$,
\begin{equation}\label{tangent:eq:values}
 \lambda=\frac67,\qquad\rho_0=\frac2{13},\qquad\tau=\frac{30}{143},
 \qquad C_\infty=\frac{375}{143}<2.623.
\end{equation}
For $q=8r$ and $r\ge2$, $C_\infty\le375/143$. This is a bound on accumulated physical error, not merely on current innovation energy.

\section{Why this is the tangent of an actual learner}\label{tangent:sec:tangent}
The fixed-reference theorem has a precise nonlinear interface. Consider a family of leaf and parent spaces $U_I(\delta)$ for $\delta>0$ and fixed target matrices $C_I$ obtained by restricting one global physical target. Assume, on a fixed finite tree, that the one-sided limits $U_I(0^+)$ exist, contain $C_I$, and are fixed independently of the regression probe $X$. Assume the leaf projection residuals have derivatives
\[
 (\Id-P_{U_\ell(\delta)})C_\ell=\delta E_\ell+o(\delta),
 \qquad U_\ell(0^+)^TE_\ell=0.
\]
At a parent suppose the actual selected space contains
\begin{equation}\label{tangent:eq:containfit}
 W_I(\delta)\bigl((W_I(\delta)^TX_I)^\dagger\bigr)^TX_I^TC_I
\end{equation}
for all sufficiently small positive $\delta$. Full-range selection on a clean outgoing sample has precisely this property. Fixed additional nuisance directions, such as $C_K$ at the union root, are allowed.

\begin{proposition}[Differentiation gives the regression tree]\label{tangent:prop:derivative}
On the full-rank, one-sided differentiability locus, the derivatives
\[
 Z_I=\left.\frac{d}{d\delta}(\Id-P_{U_I(\delta)})C_I\right|_{0^+}
\]
obey \eqref{tangent:eq:regtree} with the limiting reference spaces and leaf sources $E_\ell$. Therefore, conditionally on any randomness determining those reference spaces and leaf sources independently of $X$,
\begin{equation}\label{tangent:eq:tangentbound}
 \E_X\F{Z_{\rm root}}^2\le C_\infty\sum_\ell\F{E_\ell}^2.
\end{equation}
The statement is about the derivative. It does not assert an expectation bound on a finite difference quotient.
\end{proposition}
\begin{proof}
Write $W=W_I(\delta)$, $L=((W^TX_I)^\dagger)^T$, and $F_\delta=(\Id-WW^T)C_I$. The fitted vector in \eqref{tangent:eq:containfit} equals
\[
 WLX_I^TC_I=C_I-F_\delta+WLX_I^TF_\delta.
\]
Because it belongs to the selected parent space,
\[
 (\Id-P_{U_I(\delta)})C_I
 =(\Id-P_{U_I(\delta)})(F_\delta-WLX_I^TF_\delta).
\]
At zero, $F_0=0$, and its derivative is the stacked child derivative $F_I$. It is orthogonal to the limiting $W_I$. Differentiating gives
\[
 Z_I=F_I-W_IS_IS_I^T(G_I^\dagger)^TX_I^TF_I,
\]
which is \eqref{tangent:eq:regtree}. Derivatives of the inverse and selected projector multiply $F_0=0$ and disappear from this first-order identity. No independence is assigned to the derivatives themselves. The limiting reference hierarchy, not the derivative trajectory, is the object independent of $X$ in Theorem~\ref{tangent:thm:energy}.
\end{proof}

\section{A uniform mean-square tangent bound on the trapping family}\label{tangent:sec:application}
Return to \eqref{tangent:eq:family}. Fix $h$ and $\epsilon>0$, and let $\delta\downarrow0$ through positive values. On each $I$ leaf, set $a=\mathbf1_{2r}$, $b=gN_{\Omega_\ell}$, and $W=\Omega_{\ell'}N_{\Omega_\ell}$ in nullspace coordinates. The local sample is $ab+\delta W$, with a harmless common $1/\sqrt n$ scale.

\begin{lemma}[Leaf tangent budget]\label{tangent:lem:leafderiv}
Almost surely the leaf projector has a one-sided limit $P_\ell^0$ of rank $r$. It contains $a$ and the leading $r-1$ left directions of $(\Id-P_a)W(\Id-P_b)$. Its signal residual derivative is
\begin{equation}\label{tangent:eq:leafderiv}
 E_\ell=-\frac{(\Id-P_\ell^0)Wb^T}{\sqrt n\,\|b\|^2}.
\end{equation}
These limiting leaf spaces and derivatives depend only on $\Omega$. For
$\beta_h=\sum_{\ell\subset I}\|E_\ell\|^2$,
\begin{equation}\label{tangent:eq:betamean}
 \boxed{\E_\Omega\beta_h=\frac1{2(q-2r-2)}.}
\end{equation}
\end{lemma}
\begin{proof}
Expand $(ab+\delta W)(ab+\delta W)^T$ in $\operatorname{span}(a)\oplus a^\perp$. Eliminating its separated nonzero eigenvalue leaves, after division by $\delta^2$, the Gram matrix of $(\Id-P_a)W(\Id-P_b)$. In orthonormal coordinates this is an $(2r-1)\times(q-2r-1)$ Gaussian matrix. Its relevant singular values are positive and distinct almost surely. This proves the one-sided selected-space limit and differentiability of the projected signal. The off-diagonal first-order block of the Gram matrix gives
\[
 \left.\frac d{d\delta}P_\ell(\delta)a\right|_{0^+}
 =\frac{(\Id-P_\ell^0)Wb^T}{\|b\|^2},
\]
which proves \eqref{tangent:eq:leafderiv} after scaling and changing sign.

Conditional on $\Omega_\ell,g$, $Wb^T/\|b\|$ is standard Gaussian and independent of $W(\Id-P_b)$, hence independent of $P_\ell^0$. Its projection onto the $r$-dimensional complement has expected squared norm $r$. Therefore
\[
 \E[\|E_\ell\|^2\mid\Omega_\ell,g]=\frac r{n\|b\|^2}.
\]
The squared norm of $gN_{\Omega_\ell}$ has a chi-square law with $q-2r$ degrees of freedom. Its inverse mean is $1/(q-2r-2)$. There are $n/(2r)$ leaves inside $I$. Summing expectations gives \eqref{tangent:eq:betamean}; no independence between different leaf choices is needed.
\end{proof}

At a clean internal node, the row rule contains $W(G^\dagger)^Ta_S$ and completes its rank-one range. The limiting row template contains $u_S$, so this fitted vector tends to $u_S$. The first $r-1$ canonical physical completion vectors are inherited from the leftmost coarse leaf. They have limits depending only on $\Omega$, and $u_S$ has a nonzero component in the other child. Thus the completion is regular at the limit. Induction proves an exactly containing limiting reference hierarchy independent of $\Psi$. This remains a statement about the positive-$\delta$ branch, not a license to replace its column hierarchy by the branch at $\delta=0$.

At $P$, Theorem~\ref{tangent:thm:exactunion} gives the union space $\ran[WLa,C_K]$. Its limit is $\operatorname{span}\{u,C_K\}$, of rank $r+1$. It contains the fitted target required in \eqref{tangent:eq:containfit}. The root-rank variant of Theorem~\ref{tangent:thm:energy} therefore applies.

There is no uncountable-intersection inference here. Apply Theorem~\ref{tangent:thm:twoview} also to the limiting samples on this positive-$\delta$ branch. The unused Gaussian block $\Gamma$ is still independent of the limiting $F'$, $g$, and $H$. Almost surely the limiting selected spectra are simple and their two hyperplanes are distinct. These are open conditions and persist for every sufficiently small positive $\delta$. On that single event the exact range identity holds on an entire one-sided interval, licensing the derivative below. This argument does not use the algorithm run at $\delta=0$.

\begin{theorem}[All-depth mean-square tangent stability at the distinguished cut]\label{tangent:thm:applied}
For each fixed finite $h$ and $\epsilon>0$, the union cut loss \eqref{tangent:eq:ell} satisfies
\begin{equation}\label{tangent:eq:Qdef}
 \frac{\ell_h(\delta,\epsilon)}{\delta^2}\longrightarrow Q_h
 \quad\text{almost surely as }\delta\downarrow0,
\end{equation}
where $Q_h$ is independent of the weak-sibling amplitude and its Gaussian anchor. Uniformly in $h$,
\begin{equation}\label{tangent:eq:Qbound}
 \boxed{\E Q_h\le\frac{C_\infty}{2(q-2r-2)}.}
\end{equation}
In particular, at $r=2,q=16$,
\begin{equation}\label{tangent:eq:number}
 \boxed{\E Q_h\le\frac{75}{572}<0.132.}
\end{equation}
\end{theorem}
\begin{proof}
Use the one-sided hierarchy described above and Proposition~\ref{tangent:prop:derivative}, with $X=\Psi$ after conditioning on $\Omega$. The $K$ subtree contributes zero leaf derivatives and has deterministic row spaces. It can be included as a balanced zero-source subtree. The limiting union root has $r+1$ retained columns, which is allowed by Theorem~\ref{tangent:thm:energy}. Thus $Q_h$ is the squared norm of the resulting physical root derivative, and
\[
 \E_\Psi[Q_h\mid\Omega]\le C_\infty\beta_h.
\]
Average \eqref{tangent:eq:betamean}. The parameter independence follows already from the exact finite-$\delta$ range identity. The scalar value is $(375/143)/20=75/572$.
\end{proof}

\paragraph{Consequences that do not interchange an unbounded limit.}
For every finite cap $M>0$, bounded convergence gives
\[
 \lim_{\delta\downarrow0}\E\min\{\ell_h(\delta,\epsilon)/\delta^2,M\}
 =\E\min\{Q_h,M\}\le\frac{C_\infty}{2(q-2r-2)}.
\]
For $t>0$, convergence in probability and Markov's inequality give the valid weak statement
\[
 \limsup_{\delta\downarrow0}\PP\{\ell_h(\delta,\epsilon)>t\delta^2\}
 \le\PP\{Q_h\ge t\}\le\frac{C_\infty}{2(q-2r-2)t}.
\]
The constants are uniform in $h$; the limit is taken at each fixed finite $h$. Neither formula controls a simultaneous sequence $h\to\infty$, $\delta\to0$ without another argument.

Since $0\le\ell_h\le1$ and its limit is zero, ordinary bounded convergence also permits a deterministic positive choice of $\delta_h$ with $\E\ell_h(\delta_h,\epsilon_h)\le\epsilon_h^2$ for any prescribed positive sequence $\epsilon_h$. With \eqref{tangent:eq:optlower}, the distinguished-cut loss can then be made at most $2\OPT_k(A_h)$. This can be required in addition to any finite list of one-sided limiting conditions from the earlier trapping construction by choosing $\delta_h$ smaller. It is not a statement about the full row-and-column coefficient-space error.

\paragraph{What is not concluded.}
It would be invalid to infer
$\limsup_{\delta\downarrow0}\E[\ell_h(\delta,\epsilon)/\delta^2]\le75/572$
from \eqref{tangent:eq:Qdef}--\eqref{tangent:eq:number} alone. Uniform integrability of these quotients, and especially a depth-uniform nonlinear remainder bound, have not been proved. The all-input union-space theorem \eqref{tangent:eq:missing} remains separate.

\section{Reproducible checks and numerical scope}
The accompanying programs use fixed seeds and NumPy. The results are diagnostics, not numerical proofs.

\paragraph{Two-view algebra.}
Eighty generic tests use $r=2,3,4,8$, $q=8r$, independent Gaussian anchors, and amplitudes spanning four orders of magnitude. The largest projector discrepancy between the paired union and the full $\ran F$ was $2.5\times10^{-13}$.

\paragraph{Fixed-reference regression tree.}
Six Monte Carlo configurations cover ranks $1,2,4$ and heights $0,1,3,4,5,6$, with 21,000 total independent probe realizations. Their combined physical-plus-retained energies agree with the exact identity within 1.44 reported standard errors. For $r=2,q=16$:
\begin{center}
\begin{tabular}{rrrr}
\toprule Height&Predicted combined energy&Sample mean&Standard error\\\midrule
0&1.153846&1.149896&0.002744\\
1&1.363636&1.365028&0.005288\\
3&1.697588&1.694770&0.006100\\
6&2.040003&2.035391&0.004630\\\bottomrule
\end{tabular}
\end{center}
The leaf source budget is one. The remaining two configurations are recorded in the JSON file. Innovations are checked against $\lambda^H(q-r)\beta$ as well.

\paragraph{Actual paired learner and its tangent.}
There are nine finite-difference checks, five fixed-probe sibling-amplitude checks, and 2,000 additional paired-learner/tangent realizations. The latter give:
\begin{center}
\begin{tabular}{rrrr}
\toprule Clean height $h$&Realizations&Mean $Q_h$&Standard error\\\midrule
0&1,000&0.063744&0.002193\\
2&600&0.080800&0.002114\\
4&300&0.089696&0.002589\\
6&100&0.102343&0.004258\\\bottomrule
\end{tabular}
\end{center}
These are tangent coefficients, not full approximation errors. The common theorem bound is $75/572\simeq0.131119$. At $h=5$ and a fixed probe seed, changing $\epsilon$ from $10^{-3}$ to $30$ left the finite-$\delta$ cut loss equal to $7.3786148\times10^{-8}$ up to rounding, at $\delta=10^{-3}$. This tests the exact amplitude-independent range identity, not an expectation estimate.

The known clean rank-one and zero blocks are evaluated structurally. A numerical rank test relative only to cancellation roundoff is not used to decide their exact branch. These programs specialize the known test input for verification; they do not supply extra information to the general learning algorithm. Exact orthonormalization of a nearly dependent union is a mathematical operation; finite-precision safeguards would require a separate analysis.

\section{What remains for the general target}
The unchanged hybrid-only all-depth inequality remains ruled out by the previous construction. This chapter does not retract that obstruction. It establishes three positive facts for the union route: exact recovery of the entire lifted rank-$(r+1)$ sample range at its failing cut; a reusable all-depth regression-tree identity; and an actual clean-subtree tangent bound with an absolute constant.

The obstacle has not disappeared. At arbitrary noisy inputs, the reference spaces need not be independent of the regression probe, comparator cuts need not be clean, and both hierarchies may incur new losses at many levels. The Gaussian stabilizer symmetry in Lemma~\ref{tangent:lem:cov} cannot be assigned to such adaptive spaces without proof. The separate hierarchies can also develop different child templates away from the trapping family, so Theorem~\ref{tangent:thm:twoview} cannot simply be applied to every union node.

A sufficient next result would bound the all-depth physical union loss by the original comparator residual, retaining a nonlinear counterpart of both registers in \eqref{tangent:eq:registers}. A finite-noise small-step inequality alone is not enough unless the innovation source is shown to decay or be charged once. The linearized theorem supplies exactly such a decay-and-telescope mechanism in its stated model; its extension is the substantive remaining work.

If \eqref{tangent:eq:missing} is proved, the independent refit \eqref{tangent:eq:refit} yields a known rank-$4k$ approximation with the same depth order using at most $64k$ products. Exact offline best approximation of that known matrix in $\HH_k$ costs no further oracle products and obeys
\[
 \F{A-\widehat B}\le2\F{A-\widehat C}+\sqrt{\OPT_k(A)}.
\]
This is the existing query-only final-rank reduction, not a polynomial-time optimization claim. The present work has not proved its missing basis-quality premise.

\endgroup

% ===== NONLINEAR =====
\begingroup
\let\E\relax
\newcommand{\E}{\mathbb E}
\let\R\relax
\newcommand{\R}{\mathbb R}
\let\Id\relax
\newcommand{\Id}{\mathrm I}
\let\F\relax
\newcommand{\F}[1]{\lVert #1\rVert_F}
\let\op\relax
\newcommand{\op}[1]{\lVert #1\rVert_{\mathrm{op}}}
\let\tr\relax
\newcommand{\tr}{\operatorname{tr}}
\let\ran\relax
\newcommand{\ran}{\operatorname{range}}
\let\diag\relax
\newcommand{\diag}{\operatorname{blockdiag}}
\let\rank\relax
\newcommand{\rank}{\operatorname{rank}}
\let\OPT\relax
\newcommand{\OPT}{\operatorname{OPT}}
\let\dist\relax
\newcommand{\dist}{\operatorname{dist}}
\chapter{Nonlinear regression and mean-square stability}\label{ch:nonlinear}
\begin{scopebox}
\textbf{Update, 20 September 2026.} The nonlinear regression and fixed-height mean-square limit results are theorems for the local model and application stated here. Later independent-source results strengthen that local theory, but do not identify a general HSS residual history with independent Gaussian source blocks. Shared exterior probes, history-dependent targets, and the opposite hierarchy remain genuine interface issues. The current synchronized route keeps the signed physical source, inherited return, and possible invisible component together rather than importing subtree independence from resemblance of local formulas.
\end{scopebox}

\begin{scopebox}
The local-equivariant model permits nonlinear, source-dependent selections. Its hypotheses are not automatically satisfied by general HSS cuts with shared exterior probes. The subsequent inverse-budget theorems control this model under additional selection-specific conditions.
\end{scopebox}
\thispagestyle{plain}
\begin{chaptersynopsis}
We extend the preceding fixed-reference Gaussian regression-tree calculation to nonlinear, subtree-local, query-rotation-equivariant choices of retained spaces. Learned designs need not be Gaussian. Nevertheless their normalized row frames are uniform, and the query-space innovation has an exact geometric decay in expected energy. We derive the exact physical-energy cost in terms of inverse-weighted innovation moments, including the cross-sibling terms. A general bound is finite but grows with depth; no depth-independent physical-risk bound for arbitrary adaptive splits is claimed.

A fourth-moment argument, also valid for these nonlinear trees, supplies a uniform finite-height bound over all admissible selection rules and source amplitudes. Applied to the actual paired-union learner on the positive-noise anchor-trapping family, this proves uniform integrability of the squared error quotients at each fixed depth. The one-sided derivative of Chapter~\ref{ch:tangent} is consequently a genuine mean-square derivative. In particular, for $r=2,q=16$, the expected cut-error quotient converges to a number at most $75/572$, with that limiting bound uniform over depths. Interchanging the depth limit with the small-noise limit is not justified. The all-input $O(k)$-query, expected $O(L)$ best-HSS theorem remains unproved.
\end{chaptersynopsis}


\section{Target, maintained learner, and inherited inputs}
Fix a balanced coordinate tree and let $\mathcal H_k$ be its HSS-rank-at-most-$k$ class. The target remains
\[
 \E\F{A-\widehat B}^2\le CL\OPT_k(A),\qquad
 \OPT_k(A)=\min_{B\in\mathcal H_k}\F{A-B}^2,
\]
with $O(k)$ products with $A,A^T$. The current learner uses $r=2k$, coarse leaf size $2r$, and independent Gaussian probes $\Omega,\Psi$ of width $q=16k=8r$. Its original-data cut observable is
\[
 J_I=\Psi^T[P_I,A]\Omega.
\]
With each hierarchy's own child templates, write $G=W_U^T\Psi_I$, $H=W_V^T\Omega_I$, $\Pi_D=D^\dagger D$, and $N_D=\Id-\Pi_D$. The hybrid truncates $\Pi_GJ_IN_H$ to rank $r$ in query coordinates before lifting through $(G^\dagger)^T$. The companion truncates $\Pi_GJ_I$ before the same lift. Their transpose counterparts give column spaces. Their complete hierarchies are built separately; their physical unions are taken afterwards, and have ranks at most $2r=4k$. This is not a recursively union-fed algorithm.

At physical leaves both hierarchies use the leading-$r$ spaces of $Y_IN_{\Omega_I}$ and $Z_IN_{\Psi_I}$. Low-rank samples are completed by a deterministic rule depending only on the recovered physical subspace. This makes physical frame choices unchanged by rotations of query coordinates wherever the clean regression representation applies.

\paragraph{Inherited results and scope of this chapter.}
The preceding analysis (Chapter~\ref{ch:tangent}) establishes an exact paired-union range identity on the trapping family and an almost-sure tangent limit with $\E Q_h\le75/572$ at $r=2,q=16$. It explicitly does not interchange the untruncated expectation and the small-noise limit. Those two interfaces are restated in Section~\ref{nonlinear:sec:application}; the new moment proof closes precisely that interchange. The fixed-space refit is a separate theorem of Chapter~\ref{ch:uniform}: independent rank-$s$ refitting costs $8s$ products and has factor $1607$ relative to the chosen coefficient space. It supplies no basis-quality conclusion; the final-rank reduction is proved separately in Chapter~\ref{ch:bicriteria}.

Sections~\ref{nonlinear:sec:model}--\ref{nonlinear:sec:fourth} prove the nonlinear regression results from their stated assumptions. The application retains the preceding exact clean-trajectory and two-view identities; it does not assert that arbitrary noisy HSS cuts meet the nonlinear tree's locality hypothesis. Earlier manuscripts are unchanged.

\section{A genuinely adaptive regression tree}\label{nonlinear:sec:model}
Take a balanced binary tree of height $H$. At every physical leaf $\ell$, fix an isometry $U_\ell$ with $r$ columns and a source matrix $E_\ell$ with $p$ response columns such that $U_\ell^TE_\ell=0$. Put
\[
 \beta_\ell=\F{E_\ell}^2,\qquad
 \beta_I=\sum_{\ell\subset I}\beta_\ell,\qquad \beta=\beta_{\rm root}.
\]
The leaf data are fixed before drawing a standard Gaussian physical probe $X$ of width $q$. Equivalently, all statements may be applied conditionally on side information independent of $X$; every expectation in the tree theorems is taken with that side information fixed.

At a merge $I=a\cup b$, set $W_I=\diag(U_a,U_b)$ and $G_I=W_I^TX_I$. Choose an orthogonal split $[T_I,S_I]$ of $\R^{2r}$. Ordinarily $T_I$ has $r$ columns; at the root it may have any number $1\le s\le2r$. Set $U_I=W_IT_I$ and define
\begin{equation}\label{nonlinear:eq:physicalrec}
 \begin{split}
 F_I&=\begin{bmatrix}E_a^{\rm out}\\E_b^{\rm out}\end{bmatrix},\\
 B_I&=S_I^T(G_I^\dagger)^TX_I^TF_I,\\
 E_I^{\rm out}&=F_I-W_IS_IB_I.
 \end{split}
\end{equation}
At leaves, $E_\ell^{\rm out}=E_\ell$. The split is permitted to depend nonlinearly on the probe and the states. It is subject to two requirements:

\paragraph{Subtree locality.}
The physical choices and states in $I$ use only $X_I$ and the fixed side information. Disjoint subtrees therefore depend on independent Gaussian row blocks. A shared exterior block of $X$ is not allowed as hidden side information.

\paragraph{Query-rotation equivariance.}
Replacing every probe row in a subtree by $X_IQ$, for an orthogonal $q\times q$ matrix $Q$, leaves all its physical bases and physical source/output matrices unchanged. Thus designs transform as $R\mapsto RQ$, regression coefficients are unchanged, and query-space residuals transform as $w\mapsto Q^Tw$. Canonical physical frames can be used to make this a frame-level statement. No differentiability or Gaussian distribution of selected designs is assumed.

At an ordinary node write
\begin{equation}\label{nonlinear:eq:states}
 R_I=U_I^TX_I,\quad Y_I=X_I^TE_I^{\rm out},\quad
 c_I=(R_I^\dagger)^TY_I,\quad w_I=N_{R_I}Y_I.
\end{equation}
In particular $R_Iw_I=0$. For a merge, let $d_I=[c_a;c_b]$ and
$z_I=(G_I^\dagger)^T(w_a+w_b)$. Direct substitution gives, even for adaptive splits,
\begin{equation}\label{nonlinear:eq:register}
 \boxed{c_I=T_I^T(d_I+z_I),\qquad
 B_I=S_I^T(d_I+z_I),\qquad
 w_I=N_{G_I}(w_a+w_b).}
\end{equation}
Also $U_I^TE_I^{\rm out}=0$. These are finite-noise identities, not differential identities.

\section{The nonlinear innovation law}
\begin{lemma}[Uniform row frames, not Gaussian retained designs]\label{nonlinear:lem:haar}
For $q\ge2r$, all designs have full row rank almost surely on each fixed finite tree. At an ordinary node, the row frame
\[
 Q_I=(R_IR_I^T)^{-1/2}R_I
\]
is uniform on the row-Stiefel manifold and independent of every jointly observed query-rotation-invariant state variable at that node, including $R_IR_I^T$. This does not imply that $R_I$ is Gaussian. Child row frames are independent conditional on the fixed side information.
\end{lemma}
\begin{proof}
Leaves have Gaussian full-row-rank designs. Given full-row-rank child designs, right-rotation invariance implies uniform row frames by disintegration under the orthogonal group; the same argument gives independence from any invariant tuple. Locality gives independence between child tuples. Two independent uniform $r$-subspaces in $\R^q$ have direct sum almost surely for $q\ge2r$, so $G_I$ has full row rank. Compression by $T_I^T$ preserves full row rank, because $G_IG_I^T$ is positive definite. This closes the induction. The right-invariance statement follows from the stated equivariance of physical choices, not from independence of a learned template and $X_I$.
\end{proof}

\begin{lemma}[Conditional isotropy with a random amplitude]\label{nonlinear:lem:isotropy}
At every ordinary node there is a nonnegative function $\eta_I(R_IR_I^T)$ such that
\begin{equation}\label{nonlinear:eq:isotropy}
 \E[w_Iw_I^T\mid R_I]=\eta_I(R_IR_I^T)N_{R_I},\qquad
 \E[c_Iw_I^T\mid R_I]=0,\qquad \E[w_I\mid R_I]=0.
\end{equation}
The second assertion is used after second-moment finiteness has been established below. The conditional mean of $c_I$ is not asserted to be zero.
\end{lemma}
\begin{proof}
Conditional on $R_I$, the stabilizer of its row frame acts as the whole orthogonal group on $\ker R_I$. Equivariance makes the conditional law invariant under that action. Since $w_I\in\ker R_I$, its covariance is scalar on that complement, and its mean is zero. Reflection of the complement changes $w_I$ to $-w_I$ without changing $c_I$, proving the cross identity when integrable. The scalar coefficient is a function of the Gram matrix by full right invariance. These statements can be formalized using an equivariant regular conditional distribution obtained by compact-group averaging.
\end{proof}

\begin{theorem}[Exact nonlinear innovation decay]\label{nonlinear:thm:innovation}
Assume $q\ge2r$ and put $\lambda=(q-2r)/(q-r)$. At an ordinary node of height $t$,
\begin{equation}\label{nonlinear:eq:innovation}
 \boxed{\E\F{w_I}^2=(q-r)\lambda^t\beta_I.}
\end{equation}
The rule may be nonlinear and may choose its retained subspace from the same probe and source. The result requires the two hypotheses of Section~\ref{nonlinear:sec:model}.
\end{theorem}
\begin{proof}
At leaves, $X_\ell^TE_\ell$ is independent of $R_\ell$, with trace covariance $\beta_\ell\Id_q$. The base case follows. At a merge condition on $G_I=[R_a;R_b]$. Locality and Lemma~\ref{nonlinear:lem:isotropy} give
\[
 \E[w_Iw_I^T\mid G_I]
   =(\eta_a+\eta_b)N_{G_I}.
\]
Taking traces, using $\rank N_{G_I}=q-2r$, and averaging gives
\[
 \E\F{w_I}^2=\lambda\bigl(\E\F{w_a}^2+\E\F{w_b}^2\bigr).
\]
This proves the formula. Finiteness of these innovations also follows directly from the last identity of \eqref{nonlinear:eq:register}, since projections are contractions. No deterministic conditional amplitude is claimed: unlike the fixed-reference calculation, $\eta_I$ generally depends on the selected Gram matrix.
\end{proof}

\section{An exact physical-energy budget with weighted inverses}\label{nonlinear:sec:energy}
For an ordinary node define
\begin{equation}\label{nonlinear:eq:scalars}
 V_I=\tr((R_IR_I^T)^{-1}),\qquad v_I=\E V_I,\qquad
 m_I=\frac{\E\F{w_I}^2}{q-r},\qquad
 a_I=\frac{\E[V_I\F{w_I}^2]}{q-r}.
\end{equation}
Thus $m_I=\lambda^t\beta_I$ and $a_I=\E[\eta_I V_I]$. In general $a_I\ne m_Iv_I$.
Set
\[
 \kappa=\frac{q-r-1}{q-2r-1},\qquad \rho_0=\frac{r}{q-r-1}.
\]

\begin{lemma}[A two-frame inverse identity]\label{nonlinear:lem:stiefel}
Let $A,B$ have independent uniform row frames conditional on their positive definite Gram matrices, and put $G=[A;B]$. For $q>2r+1$, averaging their row-frame orientations gives
\begin{align}
 \E_{\rm frames}\F{(G^\dagger)^TN_A}^2
   &=(\kappa-1)\tr((AA^T)^{-1})+\kappa\tr((BB^T)^{-1}),\label{nonlinear:eq:stiefelcost}\\
 \E_{\rm frames}\tr((GG^T)^{-1})
   &=\kappa\bigl(\tr((AA^T)^{-1})+\tr((BB^T)^{-1})\bigr).\label{nonlinear:eq:stiefeltrace}
\end{align}
The analogous formula with $N_B$ exchanges $A,B$.
\end{lemma}
\begin{proof}
Write $A=K_A^{1/2}Q_A$, $B=K_B^{1/2}Q_B$, and $C=Q_AQ_B^T$. Regress a query vector with covariance $N_A$ on $G^T$. In normalized coefficient coordinates the two coefficient covariances are
\[
 C(\Id-C^TC)^{-1}C^T,\qquad (\Id-C^TC)^{-1}.
\]
Independent rotations of the two row frames make their expectations scalar. To compute the second scalar, fix $Q_A$ and obtain $Q_B$ by normalizing a Gaussian matrix $Z=[Z_1,Z_2]$, where $Z_1$ has $r$ columns and $Z_2$ has $q-r$ columns. Then
\[
 \tr((Q_BN_AQ_B^T)^{-1})
 =r+\tr((Z_2Z_2^T)^{-1}Z_1Z_1^T).
\]
Independence and the Gaussian inverse-Wishart mean yield expectation
$r+r^2/(q-2r-1)=r\kappa$. The first scalar is $\kappa-1$, by
$C(\Id-C^TC)^{-1}C^T=(\Id-CC^T)^{-1}-\Id$. Reinsert $K_A^{-1/2},K_B^{-1/2}$ to obtain \eqref{nonlinear:eq:stiefelcost}. The diagonal blocks of the block inverse of $GG^T$ give \eqref{nonlinear:eq:stiefeltrace}. The Gaussian inverse mean used here is a standard input \cite{HMT}; the averaging is over actual uniform row frames, not a Gaussian replacement of learned designs.
\end{proof}

\begin{theorem}[Exact adaptive physical-energy identity]\label{nonlinear:thm:adaptiveenergy}
For $q>2r+1$, all second moments below are finite at each fixed finite height. At every merge,
\begin{equation}\label{nonlinear:eq:cost}
 \boxed{\E\F{z_I}^2
   =(\kappa-1)(a_a+a_b)+\kappa(m_av_b+m_bv_a),\qquad
 \E\langle d_I,z_I\rangle_F=0.}
\end{equation}
Consequently,
\begin{equation}\label{nonlinear:eq:energy}
 \boxed{\E\F{E_{\rm root}^{\rm out}}^2+\E\F{c_{\rm root}}^2
  =(1+\rho_0)\beta+\sum_{I\ {\rm internal}}\E\F{z_I}^2.}
\end{equation}
The root can retain any rank $s\le2r$. All products of separate expectations in \eqref{nonlinear:eq:cost} use independence of different subtrees with the side information fixed; no within-node product $a_I$ is factorized.
\end{theorem}
\begin{proof}
Conditional on $G_I$, the covariance of $w_a+w_b$ is
$\eta_aN_{R_a}+\eta_bN_{R_b}$. Apply Lemma~\ref{nonlinear:lem:stiefel} after conditioning on the two Grams. Each $\eta$ depends only on its own Gram. Same-child terms give $a_a,a_b$; opposite-child products factor by locality. Lemma~\ref{nonlinear:lem:isotropy} removes the same-child $c$--$w$ cross terms; cross-child terms vanish because the conditional means of the $w$ registers are zero. This proves \eqref{nonlinear:eq:cost} once finiteness is established.

For that purpose, inverse compression gives pathwise
\[
 V_I\le\tr((G_IG_I^T)^{-1}).
\]
For an ordinary parent, \eqref{nonlinear:eq:stiefeltrace} and the conditional covariance before the parent selection imply
\begin{equation}\label{nonlinear:eq:upperrec}
 v_I\le\kappa(v_a+v_b),\qquad
 a_I\le\lambda\kappa(a_a+a_b+m_av_b+m_bv_a).
\end{equation}
To prove the second bound, first replace the parent inverse trace by the full predesign inverse trace, then condition on $G_I$ and use the covariance $(\eta_a+\eta_b)N_{G_I}$. This step is valid even when the parent split depends on the new innovation. Leaves satisfy $v_\ell=\rho_0$ and $a_\ell=\rho_0\beta_\ell$. Positive-integrand induction through \eqref{nonlinear:eq:upperrec} proves their finiteness, then \eqref{nonlinear:eq:cost} proves finiteness of $z$, and \eqref{nonlinear:eq:register} proves finiteness of $c$. The reflection argument is now integrable.

Finally, all physical discarded sectors are mutually orthogonal, pathwise even for adaptive splits. Thus
\[
 \F{E_{\rm root}^{\rm out}}^2=\beta+\sum_I\F{B_I}^2.
\]
The orthogonal coefficient split in \eqref{nonlinear:eq:register} gives
$\F{c_I}^2+\F{B_I}^2=\F{d_I+z_I}^2$. Take expectations, use the zero cross term, and telescope. The leaf retained coefficient energy is $\rho_0\beta$, proving \eqref{nonlinear:eq:energy}.
\end{proof}

\begin{corollary}[A finite-noise bound, with a depth-dependent constant]\label{nonlinear:cor:finiteheight}
Let $\theta=2\lambda\kappa$ and $\tau=(2\kappa-1)\rho_0$. Then
\begin{equation}\label{nonlinear:eq:badconstant}
 \E\F{E_{\rm root}^{\rm out}}^2\le C_H^{\rm ad}\beta,\qquad
 C_H^{\rm ad}=1+\rho_0+\tau\sum_{t=0}^{H-1}\theta^t.
\end{equation}
At $r=2,q=16$,
\[
 \lambda=\frac67,\quad\kappa=\frac{13}{11},\quad\rho_0=\frac2{13},\quad
 \tau=\frac{30}{143},\quad\theta=\frac{156}{77}>2.
\]
This is not a depth-independent HSS bound.
\end{corollary}
\begin{proof}
Equation~\eqref{nonlinear:eq:upperrec} gives at height $t$
\[
 v_I\le\rho_0(2\kappa)^t,\qquad
 a_I\le\rho_0\theta^t\beta_I.
\]
Substitute these and $m_I=\lambda^t\beta_I$ into \eqref{nonlinear:eq:cost}. The level-$t$ total is at most $\tau\theta^{t-1}\beta$. Sum \eqref{nonlinear:eq:energy} and drop the nonnegative root coefficient.
\end{proof}

\paragraph{What has been isolated.}
Innovation decay survives arbitrary admissible nonlinear selection, but its covariance amplitude can correlate with the selected Gram matrix. The remaining physical cost is exactly the positive sum in \eqref{nonlinear:eq:cost}. A bound of that sum by $C\beta$ would give a depth-independent nonlinear gain in this model. Equation~\eqref{nonlinear:eq:badconstant} is only a coarse universal certificate; its growth is neither a lower bound for the actual HSS union nor a claim that such a certificate is sharp.

\section{Fourth moments without fixed reference spaces}\label{nonlinear:sec:fourth}
\begin{lemma}[An angle factor with finite fourth moment]\label{nonlinear:lem:angle}
For independent uniform row-Stiefel matrices $Q_a,Q_b\in\R^{r\times q}$, put
\[
 \gamma=\sigma_{\min}([Q_a;Q_b])^{-1}.
\]
If $q>2r+3$, then $D:=(\E\gamma^4)^{1/4}<\infty$. At a tree merge, this angle factor is independent of the child tuple of rotation-invariant scalar states.
\end{lemma}
\begin{proof}
Writing $C=Q_aQ_b^T$, we have
\[
 \gamma^2=\frac1{1-\op C}
 \le2\op{(\Id-C^TC)^{-1}}.
\]
Fix $Q_a$ and obtain $Q_b$ from a Gaussian matrix $Z=[Z_1,Z_2]$ as in Lemma~\ref{nonlinear:lem:stiefel}. Then
\[
 \gamma^2\le2\op Z^2\op{Z_2^\dagger}^2.
\]
The Gaussian pseudoinverse tail has exponent $q-2r+1>4$; choose a H\"older exponent slightly larger than one so that the corresponding inverse moment remains below that exponent. Gaussian norm moments of every finite order are finite. This proves finiteness of $\E\gamma^4$ \cite{HMT}. Independence from the invariant child states follows from Lemma~\ref{nonlinear:lem:haar}; it does not say that $\gamma$ is independent of arbitrary non-invariant child vectors.
\end{proof}

\begin{theorem}[Rule-uniform fourth moments at each finite height]\label{nonlinear:thm:fourth}
Assume $q>2r+3$. There is a finite constant $K_{H,r,q}$, independent of the physical leaf dimensions, number of response columns, source sizes, and admissible selection rule, such that
\begin{equation}\label{nonlinear:eq:fourth}
 \boxed{\E\F{E_{\rm root}^{\rm out}}^4\le K_{H,r,q}^4\beta^2.}
\end{equation}
One explicit finite choice is
\begin{equation}\label{nonlinear:eq:K}
 K_{H,r,q}=1+2^{H/2}AB\sum_{t=1}^H\sum_{j=0}^t(2D)^j,
\end{equation}
where $A=(\E\op{D_0^\dagger}^4)^{1/4}$ for a standard $r\times q$ Gaussian $D_0$, $B=(q(q+2))^{1/4}$, and $D$ is defined in Lemma~\ref{nonlinear:lem:angle}. No uniform-in-$H$ bound is claimed.
\end{theorem}
\begin{proof}
For a node define the invariant scalars $t_I=\op{R_I^\dagger}$ and $b_I=\F{w_I}$, and put $s_I=\sum_{\ell\subset I}\F{E_\ell}$. At leaves, Gaussian orthogonality gives
\[
 \|t_\ell\|_{L^4}\le A,\quad
 \|b_\ell\|_{L^4}\le B\F{E_\ell},\quad
 \|t_\ell b_\ell\|_{L^4},\ \|\F{c_\ell}\|_{L^4}\le AB\F{E_\ell}.
\]
Here the norm of the Gaussian complementary projection has a distribution independent of the design. For arbitrary response columns, the elementary bound
$\E\F{X^TE}^4\le q(q+2)\F E^4$ follows by diagonalizing $E^TE$ and applying convexity to its normalized nonnegative eigenvalue weights.

At a merge,
\[
 t_I\le\gamma\max(t_a,t_b),\qquad b_I\le b_a+b_b,
\]
and both $\F{z_I}$ and $t_Ib_I$ are bounded by
\[
 \gamma\bigl(t_ab_a+t_bb_b+t_ab_b+t_bb_a\bigr).
\]
The angle is independent of the invariant scalar tuple. Opposite-child products factor in $L^4$ by locality. Minkowski induction, using $\|\max(t_a,t_b)\|_{L^4}\le\|t_a\|_{L^4}+\|t_b\|_{L^4}$, therefore gives at height $t$
\begin{align*}
 \|t_I\|_{L^4}&\le A(2D)^t,\\
 \|b_I\|_{L^4}&\le Bs_I,\\
 \|t_Ib_I\|_{L^4},\ \|\F{z_I}\|_{L^4}&\le AB(2D)^t s_I,\\
 \|\F{c_I}\|_{L^4},\ \|\F{B_I}\|_{L^4}
   &\le AB\left(\sum_{j=0}^t(2D)^j\right)s_I.
\end{align*}
The $B_I$ assertion is for internal nodes. Although these constants are loose, they never require repeatedly applying H\"older to dependent inverse factors along a path; the needed same-child product is retained as its own scalar state.

The physical orthogonal-sector identity and Minkowski give
\[
 \bigl(\E\F{E_{\rm root}^{\rm out}}^4\bigr)^{1/4}
 \le\sqrt\beta+\sum_{I\ {\rm internal}}\|\F{B_I}\|_{L^4}.
\]
At each level, $\sum_Is_I=s_{\rm root}\le2^{H/2}\sqrt\beta$. Substitute the bounds above to obtain \eqref{nonlinear:eq:K}. The final root split can have arbitrary permitted rank, since its coefficient norm is still bounded by $\F{d_I+z_I}$.
\end{proof}

\section{Application to the actual finite-noise union}\label{nonlinear:sec:application}
We spell out the input and interfaces to distinguish this application from an oracle regression model. Let $r=2,q=16$ for the numerical constants. At clean height $h$, put $n=8\,2^h$, $N=4n$, $I=[0,n)$, $K=[n,2n)$, $P=I\cup K$, and $j=2n$, using zero-based indices. Set
\[
 u=\frac{\mathbf1_I}{\sqrt n},\qquad
 A(\delta,\epsilon)=ue_j^T+\delta E_0+\epsilon E_1,\qquad \delta,\epsilon>0,
\]
where $E_0$ is the normalized permutation swapping adjacent four-coordinate leaves inside $I$, and
\[
 E_1=\frac{e_ne_{j+1}^T+e_{n+1}e_{j+2}^T}{\sqrt2}.
\]
All data and parameters are fixed before the probes.

\paragraph{Exact range interface from Chapter~\ref{ch:tangent}.}
On every clean node inside $I$, the two row hierarchies coincide and contain the fitted target
\[
 W\bigl((W^T\Psi_S)^\dagger\bigr)^T\Psi_S^Tu_S.
\]
At $K$, their common row range is the deterministic two-coordinate range $C_K=[e_n,e_{n+1}]$. At $P$, their actual union has the exact positive-noise range
\begin{equation}\label{nonlinear:eq:union}
 \mathcal U_{\cup,P}=\ran[\,WLa,C_K\,],\qquad
 L=((W^T\Psi_P)^\dagger)^T,\quad a=\Psi_I^Tu.
\end{equation}
This is the inherited two-view identity. Its hypotheses are verified in Chapter~\ref{ch:tangent} using the clean rank-one recursions, a nonzero-return argument, and the unused independent exterior Gaussian rows. In particular it holds for every fixed positive $\delta,\epsilon$ and finite $h$. We do not replace it by an assertion for arbitrary adaptive projectors or unequal templates. The two-view mechanism is that the leading-two hyperplanes of $ag+\epsilon b\Gamma$ and its right-projected version are distinct almost surely, hence span their common three-dimensional lifted range.

Let $e_h(\delta)=(\Id-P_{\mathcal U_{\cup,P}})u$ and $\ell_h(\delta,\epsilon)=\|e_h(\delta)\|^2$. By \eqref{nonlinear:eq:union}, its distribution is independent of $\epsilon>0$; inherited error from $\delta$ is not removed.

\begin{proposition}[The finite-noise physical residual is the adaptive tree]\label{nonlinear:prop:actual}
Condition on the entire forward stream $\Omega$. The leaf row spaces and their projection residuals of $u$ are fixed independently of $X=\Psi$. At every clean merge, and at the union root using \eqref{nonlinear:eq:union}, these actual projection residuals obey \eqref{nonlinear:eq:physicalrec} exactly. The choices are subtree-local and query-rotation-equivariant in $\Psi$.
\end{proposition}
\begin{proof}
For a target restriction $C$, put $F=(\Id-WW^T)C$ and $L=((W^TX)^\dagger)^T$. The fitted vector is
\[
 WLX^TC=C-F+WLX^TF.
\]
Since it lies in the selected parent space $WT$, and $F\perp\ran W$,
\[
 (\Id-P_{WT})C=F-WSS^TLX^TF.
\]
This is precisely \eqref{nonlinear:eq:physicalrec}; $F$ is the stacked child residual. It is not a linearization. At the union root the extra directions $C_K$ are included and $T$ has three columns, which the tree theorems permit. The target is zero in $K$, so that subtree has zero source.

After fixing $\Omega$, the clean row choices use only the appropriate rows of $\Psi$ through the fitted target and deterministic physical completions. Right multiplication of all those rows by an orthogonal matrix leaves each fitted physical vector unchanged. Thus locality and equivariance hold. The original two-view samples at the top may be described using exterior forward rows, but their exact reduced physical union \eqref{nonlinear:eq:union} no longer uses them or any exterior transpose rows. This is where the established span identity is essential.
\end{proof}

\begin{lemma}[Uniform leaf bounds at every positive noise level]\label{nonlinear:lem:leaf}
Let $m=2r$, $d=q-m>4$, and let
$\beta(\delta)=\sum_{\ell\subset I}\|(\Id-P_{U_\ell(\delta)})u_\ell\|^2$.
For the same normalized sibling-swap family,
\begin{equation}\label{nonlinear:eq:leafmom}
 \E_\Omega\beta(\delta)\le\frac{4d}{d-2}\delta^2,\qquad
 \E_\Omega\beta(\delta)^2
 \le\frac{16d(md+2)}{m(d-2)(d-4)}\delta^4.
\end{equation}
At $m=4,q=16$, the constants are $24/5$ and $30$.
\end{lemma}
\begin{proof}
In nullspace coordinates a leaf sample, apart from the common factor $1/\sqrt n$, is $ab+\delta Z$, with $a=\mathbf1_m$, $b=gN_{\Omega_\ell}$, and $Z=\Omega_{\ell'}N_{\Omega_\ell}$. Its leading-$r$ projector $P$ satisfies
\[
 \|(\Id-P)a\|\,\|b\|
 \le\F{(\Id-P)(ab+\delta Z)}+\delta\F Z\le2\delta\F Z,
\]
because $ab$ has rank one. Conditional on the local probe, $Z$ is standard $m\times d$ Gaussian and $b$ is an independent standard Gaussian row of length $d$. Hence
\[
 \E\F Z^2=md,\quad \E\F Z^4=md(md+2),\quad
 \E\|b\|^{-2}=\frac1{d-2},\quad
 \E\|b\|^{-4}=\frac1{(d-2)(d-4)}.
\]
There are $n/m$ leaves. Sum the second-moment bounds, and use
$(\sum x_\ell)^2\le(n/m)\sum x_\ell^2$ for the fourth-moment bound. No independence between different leaf choices is used; their sibling probe blocks are shared.
\end{proof}

\begin{theorem}[Finite-noise moments and mean-square differentiability]\label{nonlinear:thm:main}
For every fixed finite $h$, the actual union cut residual on this family satisfies
\begin{align}
 \E\ell_h(\delta,\epsilon)
 &\le\min\left\{1,\frac{24}{5}C_{h+2}^{\rm ad}\delta^2\right\},\label{nonlinear:eq:finite}\\
 \sup_{\delta>0,\,\epsilon>0}\E\left(\frac{\ell_h(\delta,\epsilon)}{\delta^2}\right)^2
 &\le30K_{h+2,2,16}^4<\infty.\label{nonlinear:eq:UI}
\end{align}
Let $Z_h$ denote the one-sided projection-residual derivative in the preceding tangent theorem. Then
\begin{equation}\label{nonlinear:eq:L2}
 \frac{e_h(\delta)}\delta\longrightarrow Z_h\quad\text{in }L^2,
 \qquad
 \frac{\ell_h(\delta,\epsilon)}{\delta^2}\longrightarrow\|Z_h\|^2\quad\text{in }L^1.
\end{equation}
In particular,
\begin{equation}\label{nonlinear:eq:sharp}
 \boxed{\lim_{\delta\downarrow0}\frac{\E\ell_h(\delta,\epsilon)}{\delta^2}
   =\E\|Z_h\|^2\le\frac{75}{572}.}
\end{equation}
The upper bound in \eqref{nonlinear:eq:sharp} is uniform in $h$ and $\epsilon>0$, but the limit is taken at each fixed finite $h$.
\end{theorem}
\begin{proof}
Condition on $\Omega$, apply Proposition~\ref{nonlinear:prop:actual}, and then apply Corollary~\ref{nonlinear:cor:finiteheight} and Theorem~\ref{nonlinear:thm:fourth} to the tree on $P$. It has $H=h+2$ levels above its physical leaves and zero sources on $K$. Average Lemma~\ref{nonlinear:lem:leaf}. This proves \eqref{nonlinear:eq:finite}--\eqref{nonlinear:eq:UI}. These arguments are uniform over the nonlinear selection rule at each fixed height, so they apply to all positive values of $\delta$ rather than only a favorable neighborhood.

The inherited one-sided tangent theorem gives almost-sure convergence $e_h(\delta)/\delta\to Z_h$ on a single sufficiently small positive-parameter interval. It establishes an exactly containing limiting hierarchy determined by $\Omega$ alone; differentiating the exact identity in Proposition~\ref{nonlinear:prop:actual} gives the fixed-reference regression tree. Its leaf derivative budget has expectation $1/20$, and its fixed-reference physical gain is at most $375/143$. Thus $\E\|Z_h\|^2\le(375/143)/20=75/572$. These are the restated tangent interfaces; no expectation-limit interchange is included in that inherited input.

The uniform fourth moment \eqref{nonlinear:eq:UI} makes the squared norms of the vector quotients uniformly integrable. Fatou also gives a finite fourth moment for $Z_h$. Almost-sure convergence therefore upgrades to $L^2$ convergence of the vectors. Cauchy--Schwarz then gives $L^1$ convergence of their squared norms, proving \eqref{nonlinear:eq:L2}--\eqref{nonlinear:eq:sharp}.
\end{proof}

\paragraph{The precise improvement.}
The preceding result bounded $\E$ of the almost-sure quadratic coefficient. Equation~\eqref{nonlinear:eq:sharp} now identifies it with the limit of the actual expected error quotients. For every fixed $h$ and every $\eta>0$, there is a deterministic $\delta_0(h,\eta)>0$ such that
\[
 \E\ell_h(\delta,\epsilon)\le(75/572+\eta)\delta^2
 \quad(0<\delta<\delta_0(h,\eta),\ \epsilon>0).
\]
Amplitude independence permits the same threshold for every positive $\epsilon$. No rate for $\delta_0$ uniform over heights is established. In particular we have proved
\[
 \sup_h\lim_{\delta\downarrow0}\frac{\E\ell_h(\delta,\epsilon)}{\delta^2}\le75/572,
\]
not the reversed-order assertion $\limsup_{\delta\downarrow0}\sup_h\E\ell_h/\delta^2\le75/572$. A distinguished-cut estimate is not a full coefficient-space estimate.

\section{Checks and reproducibility}
The accompanying scripts are diagnostics rather than proofs.

\paragraph{Nonlinear regression trees.}
The adaptive-tree checker uses fixed splits, splits chosen from the smallest eigenvalues of the current Gram matrix, and source-dependent splits. There are 7,600 realizations in six configurations, with $r=1,2,3$, $q=8r$, and heights up to seven. It checks the exact register and energy identities, geometric innovation decay, and the inverse-weighted cost formula. The largest register identity residual is below $1.4\times10^{-15}$, and the largest energy identity residual is below $1.1\times10^{-14}$. Across the tested levels, innovation means differ from their exact predictions by at most 1.68 reported standard errors; the paired weighted-cost differences are within 2.05. No fixed-reference covariance amplitude is imposed in these adaptive tests.

\paragraph{Actual positive-noise family.}
The finite-noise checker uses the proved reduced union range \eqref{nonlinear:eq:union} and accumulates positive physical discarded-sector losses. It avoids computing a tiny residual as $1-\|Pu\|^2$. There are 12 pointwise finite-difference checks and 10,350 positive-noise realizations over four heights and three positive noise levels. At $\delta=10^{-3}$ the means are:
\begin{center}
\begin{tabular}{rrrr}
\toprule Clean height $h$&Realizations&Mean $\ell_h/\delta^2$&Mean $\|Z_h\|^2$\\\midrule
0&1,600&0.06750350&0.06750499\\
2&1,000&0.08339545&0.08339600\\
4&600&0.09366495&0.09366299\\
6&250&0.10270954&0.10271019\\\bottomrule
\end{tabular}
\end{center}
The corresponding standard errors of the finite-noise means are approximately $0.00179$, $0.00182$, $0.00195$, and $0.00281$. The same probes are used for the paired derivative comparison. These are cut-error diagnostics, not full HSS approximation errors or evidence for exchanging the height and noise limits.

\section{The remaining general-target theorem}
This chapter establishes a nonlinear regression law under explicit locality and equivariance assumptions, a finite-height physical bound at positive noise, and the previously missing uniform-integrability step for the actual paired-union trapping cut. It does not change either hierarchy or its query count.

Two difficulties remain separate. Even in the local-equivariant model, a depth-independent bound needs control of the exact weighted sum \eqref{nonlinear:eq:cost}; the generic inverse estimate yields the growing constant \eqref{nonlinear:eq:badconstant}. On general noisy HSS inputs, incoming cut terms and the two distinct child templates can introduce dependence on probe rows outside a subtree. That dependence is excluded by the model proved here. Its independence and stabilizer arguments cannot be transferred without verifying an appropriate conditional law.

The desired missing premise is still
\[
 \E\dist_F(A,\mathcal S_\cup(A))^2\le C_bL\OPT_k(A).
\]
If it is proved, the independent fixed-space refit yields a known rank-$4k$ approximation using $32k$ shared training products and $32k$ refit products. Exact offline best HSS-rank-$k$ projection costs no further oracle products and preserves an $O(L)$ error factor through
\[
 \F{A-\widehat B}\le2\F{A-\widehat C}+\sqrt{\OPT_k(A)}.
\]
This remains a query-only consequence, not a claim of efficient nonlinear optimization or a completed all-input guarantee.

\endgroup

% ===== BUDGET =====
\begingroup
\let\E\relax
\newcommand{\E}{\mathbb E}
\let\R\relax
\newcommand{\R}{\mathbb R}
\let\Id\relax
\newcommand{\Id}{\mathrm I}
\let\F\relax
\newcommand{\F}[1]{\lVert #1\rVert_F}
\let\op\relax
\newcommand{\op}[1]{\lVert #1\rVert_{\rm op}}
\let\tr\relax
\newcommand{\tr}{\operatorname{tr}}
\let\ran\relax
\newcommand{\ran}{\operatorname{range}}
\let\rank\relax
\newcommand{\rank}{\operatorname{rank}}
\let\diag\relax
\newcommand{\diag}{\operatorname{blockdiag}}
\let\dist\relax
\newcommand{\dist}{\operatorname{dist}}
\let\OPT\relax
\newcommand{\OPT}{\operatorname{OPT}}
\chapter{Relative inverse budgets and signal-preserving completion}\label{ch:budget}
\begin{scopebox}
\textbf{Update, 20 September 2026.} The local-equivariant counterexample, sufficient relative guard, and constrained completion optimizer are unchanged. Locality and symmetry alone still do not imply stability, and a mandatory direction can make every fixed relative guard infeasible. Later fixed-source and guard-completion theorems supply additional valid restricted cases; they do not make an arbitrary selected HSS core fixed or independent. References here to the maintained union and its $64k$ accounting refer to the older separately trained algorithm, not the synchronized learner of Chapter~\ref{upd:chap:current-route}.
\end{scopebox}

\begin{scopebox}
The relative guard is a sufficient condition, not a property established for the maintained union. The broad-model counterexample is not an actual-union counterexample. Chapter~\ref{ch:core} provides an additive alternative for optimal signal-preserving completion.
\end{scopebox}
\thispagestyle{plain}
\begin{chaptersynopsis}
The preceding nonlinear regression calculation isolates an exact inverse-weighted innovation sum. We show that locality and query-rotation equivariance alone cannot bound this sum uniformly: an admissible selection on a zero-source sibling can force exponential growth of physical error elsewhere. This is a counterexample for the regression model, not for the maintained HSS union learner.

We then prove a depth-independent nonlinear physical-error theorem under a computable relative inverse-budget condition. The condition permits inverse moments to grow, provided that their growth is slower than innovation decay. At retained rank two and width sixteen, nonincreasing sibling inverse budgets give the same bound $375/143$ as the fixed-reference theorem, while a ten-percent local inflation gives $690/143$. At the critical inflation, the bound is linear in tree height.

To preserve a selected signal rather than replace it by well-conditioned unrelated directions, we solve the local inverse-trace completion problem exactly by a Schur complement and a generalized eigenvalue calculation. A conditional all-depth theorem separates the forced-signal inverse cost from optional completion cost. A local example proves that some forced cost is unavoidable. The general all-input $O(k)$-query HSS theorem remains unproved; no unchanged-learner risk claim is inferred from the new stability criterion.
\end{chaptersynopsis}


\section{Status, target, and the scope of the new results}
The target is an HSS-rank-at-most-$k$ output satisfying
\[
 \E\F{A-\widehat B}^2\le CL\OPT_k(A),\qquad
 \OPT_k(A)=\min_{B\in\mathcal H_k}\F{A-B}^2,
\]
using $O(k)$ products with $A,A^T$. The published fresh-level analysis provides the corresponding logarithmic-query benchmark \cite{AmselHSS}; this chapter concerns the shared-query learning obstruction.

The maintained candidate still runs two complete hierarchies on the same original observations: a right-annihilated hybrid and an un-nullified companion, then takes their physical nested unions. With $r=2k$, width $q=16k$, and physical coarse leaves of size $4k$, the union rank is at most $4k$. The $32k$ training products and independent $32k$ fixed-space refit products total $64k$. The fixed-space refit is an inherited independent-probe theorem, proved in Chapters~\ref{ch:uniform} and~\ref{ch:bicriteria}. Exact offline HSS-rank-$k$ projection is a query-only reduction, not a polynomial-time optimization claim.

Chapter~\ref{ch:nonlinear} proves the local-equivariant regression identities, geometric innovation decay, and finite-height moments. It left an inverse-weighted sum unpaid. Section~\ref{budget:sec:model} restates and proves the identities needed here, so the new mathematical arguments do not depend on unstated local notation. The present counterexample does not implement the HSS union selection rule. The positive inverse-budget theorem applies to a specified class of nonlinear regression selections; that class is not silently identified with the actual union. The completion routine is a new local construction and diagnostic, not an announced replacement with a proved HSS guarantee.

\section{The nonlinear regression model and exact identities}\label{budget:sec:model}
Consider a balanced binary tree of height $H$, with fixed orthonormal leaf frames $U_\ell$ of rank $r$ and fixed source matrices $E_\ell$ satisfying $U_\ell^TE_\ell=0$. Set
\[
 \beta_\ell=\F{E_\ell}^2,\qquad
 \beta_I=\sum_{\ell\subset I}\beta_\ell,\qquad \beta=\beta_{\rm root}.
\]
Draw one standard Gaussian physical probe $X$ with $q$ columns. Fixed side information independent of $X$ may be conditioned on throughout.

At an internal node $I=a\cup b$, let $W_I=\diag(U_a,U_b)$, $G_I=W_I^TX_I$, and choose an orthogonal split $[T_I,S_I]$ in $\R^{2r}$. Ordinary nodes retain $r$ columns in $T_I$; the final root can retain any rank at most $2r$. Put
\begin{equation}\label{budget:eq:physical}
 \begin{split}
 U_I&=W_IT_I,\\
 F_I&=[E_a^{\rm out};E_b^{\rm out}],\\
 B_I&=S_I^T(G_I^\dagger)^TX_I^TF_I,\\
 E_I^{\rm out}&=F_I-W_IS_IB_I.
 \end{split}
\end{equation}
At a leaf $E_\ell^{\rm out}=E_\ell$. Choices are \emph{subtree-local}: they use only $X_I$ and the fixed side information. They are also \emph{query-rotation equivariant}: replacing $X_I$ by $X_IQ$, $Q$ orthogonal, leaves all physical choices and physical residuals unchanged. Dependence on the probe and source is otherwise permitted. These are the same hypotheses as in the preceding nonlinear model.

At an ordinary node define
\[
 R_I=U_I^TX_I,\qquad
 c_I=(R_I^\dagger)^TX_I^TE_I^{\rm out},\qquad
 w_I=N_{R_I}X_I^TE_I^{\rm out},
 \quad N_R=\Id-R^\dagger R.
\]
With $d_I=[c_a;c_b]$ and $z_I=(G_I^\dagger)^T(w_a+w_b)$, direct algebra gives
\begin{equation}\label{budget:eq:registers}
 \boxed{c_I=T_I^T(d_I+z_I),\quad B_I=S_I^T(d_I+z_I),
 \quad w_I=N_{G_I}(w_a+w_b).}
\end{equation}
Also $U_I^TE_I^{\rm out}=0$. The physical discarded sectors are mutually orthogonal, hence
\begin{equation}\label{budget:eq:physicalsum}
 \F{E_{\rm root}^{\rm out}}^2=\beta+\sum_{I\text{ internal}}\F{B_I}^2.
\end{equation}

Assume $q>2r+1$ and abbreviate
\begin{equation}\label{budget:eq:constants}
 \lambda=\frac{q-2r}{q-r},\qquad
 \kappa=\frac{q-r-1}{q-2r-1},\qquad
 \rho_0=\frac r{q-r-1},\qquad \tau=(2\kappa-1)\rho_0.
\end{equation}
Set
\begin{equation}\label{budget:eq:moments}
 V_I=\tr((R_IR_I^T)^{-1}),\quad v_I=\E V_I,\quad
 m_I=\frac{\E\F{w_I}^2}{q-r},\quad
 a_I=\frac{\E[V_I\F{w_I}^2]}{q-r}.
\end{equation}
The symbol $V_I$ is a scalar inverse trace, not a column basis.

\begin{lemma}[Conditional laws and exact energy]\label{budget:lem:inherited}
At each finite height all these second moments are finite. All designs have full row rank almost surely. Their normalized row frames are uniform and independent of their rotation-invariant state tuples, but the selected designs themselves need not be Gaussian. There is a nonnegative function $\eta_I$ of $R_IR_I^T$ such that
\begin{equation}\label{budget:eq:conditional}
 \E[w_Iw_I^T\mid R_I]=\eta_I N_{R_I},\qquad
 \E[w_I\mid R_I]=0,\qquad \E[c_Iw_I^T\mid R_I]=0.
\end{equation}
For an ordinary node of height $t$,
\begin{equation}\label{budget:eq:innovation}
 m_I=\lambda^t\beta_I.
\end{equation}
At a merge,
\begin{equation}\label{budget:eq:cost}
 \boxed{\E\F{z_I}^2=(\kappa-1)(a_a+a_b)+\kappa(m_av_b+m_bv_a),
 \qquad \E\langle d_I,z_I\rangle_F=0.}
\end{equation}
Consequently
\begin{equation}\label{budget:eq:totalenergy}
 \boxed{\E\F{E_{\rm root}^{\rm out}}^2+\E\F{c_{\rm root}}^2
 =(1+\rho_0)\beta+\sum_{I\text{ internal}}\E\F{z_I}^2.}
\end{equation}
\end{lemma}
\begin{proof}
Right rotational invariance and compact-group disintegration make each normalized row frame uniform, independently of all invariant scalars. The two children are independent by locality. Their row spaces have direct sum almost surely for $q\ge2r$; thus the predesign is full row rank, as is every orthogonal row compression. This proves full rank by induction. Conditional on $R_I$, rotations fixing its row frame act arbitrarily on $\ker R_I$. The covariance of $w_I$ is scalar on that complement. Complement reflection changes $w_I$ to $-w_I$ without changing $c_I$, yielding \eqref{budget:eq:conditional} once integrability is established.

At leaves, independent Gaussian projections onto $U_\ell$ and its orthogonal complement give $m_\ell=\beta_\ell$, $v_\ell=\rho_0$, and $a_\ell=\rho_0\beta_\ell$. The inverse-Frobenius identity is standard \cite[Sec.~10.1]{HMT}; rank one follows directly from a chi-square inverse moment. At a merge, conditioning on $G_I=[R_a;R_b]$ gives
\[
 \E[w_Iw_I^T\mid G_I]=(\eta_a+\eta_b)N_{G_I}.
\]
Taking traces proves \eqref{budget:eq:innovation}.

For completeness, the needed two-frame integral is
\begin{align}
 \E_{\rm frames}\F{(G^\dagger)^TN_{R_a}}^2
  &=(\kappa-1)V_a+\kappa V_b,\label{budget:eq:framecost}\\
 \E_{\rm frames}\tr((GG^T)^{-1})&=\kappa(V_a+V_b).\label{budget:eq:frametrace}
\end{align}
The averaging fixes the two positive definite child Grams. Write $R_j=K_j^{1/2}Q_j$ and $C=Q_aQ_b^T$. The diagonal blocks of the inverse of $[Q_a;Q_b][Q_a;Q_b]^T$ are $(\Id-CC^T)^{-1}$ and $(\Id-C^TC)^{-1}$. Each has mean $\kappa\Id$. To evaluate the constant, normalize an $r\times q$ Gaussian $Z=[Z_1,Z_2]$ into $Q_b$, after fixing $Q_a$. Then
\[
 \E\tr((Q_bN_{Q_a}Q_b^T)^{-1})
 =r+\E\tr((Z_2Z_2^T)^{-1}Z_1Z_1^T)
 =r+\frac{r^2}{q-2r-1}=r\kappa.
\]
Subtracting the identity gives the coefficient $\kappa-1$ in \eqref{budget:eq:framecost}. This is an average over the actual uniform frames, not replacement of the learned designs by Gaussian matrices.

Since $\eta_j$ is a function of the child Gram, conditional covariance and \eqref{budget:eq:framecost} give \eqref{budget:eq:cost}. Only opposite-child expectations are multiplied. In particular $a_j=\E[\eta_jV_j]$ is not replaced by $m_jv_j$. Reflection removes same-child coefficient--innovation cross terms, and zero conditional innovation means remove opposite-child terms.

Finiteness follows first, using positive integrands, from
\[
 v_I\le\kappa(v_a+v_b),\qquad
 a_I\le\lambda\kappa(a_a+a_b+m_av_b+m_bv_a).
\]
Indeed $V_I\le\tr((G_IG_I^T)^{-1})$ by inverse compression. For the weighted inequality, replace $V_I$ by this preselection trace before conditioning on $G_I$. This makes the covariance calculation valid even for source-dependent choices. The induction then licenses the cross-term integrations above. Finally \eqref{budget:eq:registers}, \eqref{budget:eq:physicalsum}, and the zero expected cross term telescope the retained coefficient energies. The sum of leaf coefficient energies has mean $\rho_0\beta$, giving \eqref{budget:eq:totalenergy}.
\end{proof}

\section{Why locality and symmetry do not give a rule-uniform bound}\label{budget:sec:negative}
\begin{theorem}[A zero-source sibling can cause exponential physical growth]\label{budget:thm:counter}
For every $r\ge1$ and $q>2r+1$, there are fixed leaf data with one response column and $\beta=1$, and admissible subtree-local, query-rotation-equivariant selections on trees of height $h+1$, such that
\begin{equation}\label{budget:eq:lower}
 \boxed{\E\F{E_{\rm root}^{\rm out}}^2
 \ge1+\rho_0+\kappa\rho_0(\lambda\kappa)^h,\qquad
 \lambda\kappa=1+\frac{r}{(q-r)(q-2r-1)}>1.}
\end{equation}
All designs have full row rank almost surely and all second moments are finite at each fixed finite height. Thus neither a depth-independent nor an $O(H)$ bound holds for all selections allowed by the nonlinear model.
\end{theorem}
\begin{proof}
Use two root subtrees $A,B$ of equal height $h$. Leaf frames are fixed standard coordinate frames. Place a unit source orthogonal to its leaf frame in one fixed leaf of $A$; all other sources, including every source in $B$, are zero. These data are chosen before $X$.

In $A$, use fixed coordinate-retaining splits, for instance $T=[\Id_r;0]$. The innovation identity gives $m_A=\lambda^h$. In $B$, at every merge retain the eigenspace of $GG^T$ corresponding to its $r$ \emph{smallest} eigenvalues. Eigenvalue ties can be completed canonically and have no effect on the argument. This is a local, rotation-equivariant rule depending only on the design Gram. Since it retains the $r$ largest inverse eigenvalues,
\[
 V_I\ge\frac12\tr((G_IG_I^T)^{-1}).
\]
The two bad-subtree children have the same distribution. Equation~\eqref{budget:eq:frametrace} yields $v_I\ge\kappa v_a$, so $v_B\ge\rho_0\kappa^h$. Its physical errors, coefficients, and innovations are nevertheless all identically zero.

At the root, choose $T$ of rank $r$ perpendicular to the single coefficient vector $d+z\in\R^{2r}$. Its orthogonal complement has dimension at least $2r-1\ge r$; a canonical choice on the zero-vector event is harmless. This root rule is also local and rotation-equivariant. It gives $c_{\rm root}=0$. At that merge, \eqref{budget:eq:cost} and the zero source in $B$ imply
\[
 \E\F{z_{\rm root}}^2=(\kappa-1)a_A+\kappa m_Av_B
 \ge\kappa\rho_0(\lambda\kappa)^h.
\]
Insert this in \eqref{budget:eq:totalenergy}; all other merge costs are nonnegative. The displayed formula for $\lambda\kappa$ is elementary algebra.
\end{proof}

At $r=2,q=16$ the conclusion is
\begin{equation}\label{budget:eq:counter16}
 \E\F{E_{\rm root}^{\rm out}}^2
 \ge\frac{15}{13}+\frac2{11}\left(\frac{78}{77}\right)^h.
\end{equation}
The growth rate is deliberately conservative. Appendix~\ref{budget:sec:scalar} gives a stronger scalar version.

\paragraph{Scope.}
This theorem is \emph{not} a fixed-matrix counterexample for the HSS union learner. The zero-source branch deliberately chooses weak Gram directions rather than the current canonical completion, and the root deliberately discards the entire coefficient vector. It need not satisfy the extra fitted-target containment present in an actual clean HSS cut. What is ruled out is a proposed theorem uniform over the entire local-equivariant regression class. The two symmetry hypotheses alone leave enough freedom for physical error to grow. In particular the next proof must preserve information about the actual selection, not just its symmetry class.

\section{A depth-independent nonlinear theorem with an inverse budget}\label{budget:sec:positive}
For a retained space $T_I$ define $V_I$ as above, with child quantities $V_a,V_b$ evaluated \emph{on the same realization}.

\begin{theorem}[Inverse growth slower than innovation decay]\label{budget:thm:guard}
Suppose every ordinary internal node obeys
\begin{equation}\label{budget:eq:guard}
 \boxed{V_I\le\gamma\min\{V_a,V_b\}}
\end{equation}
for a fixed $\gamma\ge1$. All other choices may be nonlinear and source-dependent, subject to the two model hypotheses. The final root need not satisfy \eqref{budget:eq:guard}. Then
\begin{equation}\label{budget:eq:guardbound}
 \boxed{\E\F{E_{\rm root}^{\rm out}}^2+\E\F{c_{\rm root}}^2
 \le\left[1+\rho_0+\tau\sum_{t=0}^{H-1}(\gamma\lambda)^t\right]\beta.}
\end{equation}
In particular $\gamma\lambda<1$ gives a depth-independent bound. At the critical value $\gamma\lambda=1$, the bound is $(1+\rho_0+\tau H)\beta$.
\end{theorem}
\begin{proof}
At ordinary height $t$ we prove
\begin{equation}\label{budget:eq:guardmom}
 v_I\le\rho_0\gamma^t,\qquad
 a_I\le\rho_0(\gamma\lambda)^t\beta_I.
\end{equation}
The leaf case is Lemma~\ref{budget:lem:inherited}. The first induction uses \eqref{budget:eq:guard} and $\E\min(V_a,V_b)\le\min(v_a,v_b)$. For the second, the minimum in the guard is measurable before parent selection. Thus
\begin{align*}
 a_I&\le\frac{\gamma}{q-r}\E[\min(V_a,V_b)\F{w_I}^2]\\
 &=\gamma\lambda\E[\min(V_a,V_b)(\eta_a+\eta_b)]\\
 &\le\gamma\lambda(m_av_b+m_bv_a).
\end{align*}
The last step deliberately charges each source amplitude against the \emph{opposite} child's inverse. This is precisely where locality permits factorization. Substituting the first induction and $m_j=\lambda^{t-1}\beta_j$ proves the second induction without imposing $a_I=m_Iv_I$.

At a merge whose children have height $t$, \eqref{budget:eq:cost} is therefore at most $\tau(\gamma\lambda)^t\beta_I$. Summing over the partition at each height, and then using \eqref{budget:eq:totalenergy}, proves \eqref{budget:eq:guardbound}.
\end{proof}

At $r=2,q=16$, $\lambda=6/7$, so the stable range is $1\le\gamma<7/6$. Two explicit bounds are
\begin{equation}\label{budget:eq:explicitpositive}
 \begin{array}{c|c}
 \gamma&\text{uniform factor for physical plus root-coefficient energy}\\\hline
 1&375/143<2.623\\
 11/10&690/143<4.826.
 \end{array}
\end{equation}
At $\gamma=7/6$, the factor is $15/13+30H/143$. Notice that the theorem \emph{allows} $v_I$ to grow like $\gamma^t$. It is not a return to an unnecessary uniform bound on the full inverse. The product with the decaying innovation is the decisive quantity.

\begin{proposition}[A realizable stable selection]\label{budget:prop:strong}
Selecting the leading-$r$ eigenspace of $G_IG_I^T$ satisfies \eqref{budget:eq:guard} with $\gamma=1$ pathwise. A source-dependent proposed space can also be accepted whenever it satisfies the guard, with the leading-Gram space as fallback.
\end{proposition}
\begin{proof}
The largest $r$ eigenvalues of the $2r\times2r$ predesign Gram dominate, term by term, all eigenvalues of each $r\times r$ principal child Gram. Summing their inverse eigenvalues proves $V_I\le V_a$ and $V_I\le V_b$. A guarded choice preserves the same inequality by construction. Both choices use invariant local quantities.
\end{proof}

\paragraph{A stability theorem is not yet a basis-learning theorem.}
The leading-Gram rule can discard a required signal direction. The guarded fallback also need not preserve the fitted target. Thus Proposition~\ref{budget:prop:strong} makes the nonlinear stability class nonempty and implementable, but it is not a replacement HSS algorithm with the desired approximation guarantee. The next section keeps the signal constraint explicit.

\section{Optimal completion while preserving selected signal directions}\label{budget:sec:completion}
At one node let $K=GG^T\succ0$ be $2r\times2r$, and suppose an already selected physical signal subspace in local coordinates has orthonormal frame $Q\in\R^{2r\times s}$, $s\le r$. Consider
\begin{equation}\label{budget:eq:optimization}
 \min_{T^TT=\Id_r,\ \ran Q\subseteq\ran T}
 \tr((T^TKT)^{-1}).
\end{equation}
All quantities are available from the local design and selected signal; this problem requires no further products with the unknown matrix.

Choose an orthonormal complement $Z$ of $Q$ and write
\[
 [Q,Z]^TK[Q,Z]=\begin{pmatrix}A&B\\B^T&D\end{pmatrix},\qquad
 \mathcal S=D-B^TA^{-1}B,\qquad
 \mathcal M=\Id+B^TA^{-2}B.
\]
Both $\mathcal S$ and $\mathcal M$ are positive definite.

\begin{theorem}[Exact signal-preserving inverse-trace optimizer]\label{budget:thm:completion}
For $0<s\le r$, let
\[
 \nu_1\le\cdots\le\nu_{2r-s}
\]
be the eigenvalues of $\mathcal S^{-1/2}\mathcal M\mathcal S^{-1/2}$. The optimum in \eqref{budget:eq:optimization} equals
\begin{equation}\label{budget:eq:optvalue}
 \boxed{V_{\rm opt}=\tr(A^{-1})+\sum_{j=1}^{r-s}\nu_j.}
\end{equation}
It is attained by $\ran T=\ran[Q,ZV]$, where $V$ spans the generalized eigenvectors of $(\mathcal M,\mathcal S)$ corresponding to the smallest $r-s$ eigenvalues. Any orthonormal basis of this span can be used. For $s=0$, the optimizer is the leading-$r$ eigenspace of $K$.
\end{theorem}
\begin{proof}
Every feasible space has an orthonormal basis $[Q,ZV]$ with $V^TV=\Id_{r-s}$. Block inversion gives
\begin{equation}\label{budget:eq:schurtrace}
 \tr((T^TKT)^{-1})=	r(A^{-1})+
 \tr\left[(V^T\mathcal S V)^{-1}V^T\mathcal M V\right].
\end{equation}
For $H=\mathcal S^{1/2}V(V^T\mathcal S V)^{-1/2}$, $H^TH=\Id$, and the second term is
\[
 \tr\left[H^T\mathcal S^{-1/2}\mathcal M\mathcal S^{-1/2}H\right].
\]
Conversely each $(r-s)$-dimensional $H$-space arises from a $V$-space by this transformation. Diagonalizing the positive definite middle matrix shows that the minimum trace is the sum of its smallest $r-s$ eigenvalues. This proves both the optimum and the generalized-eigenvector construction. Empty sums handle $s=r$; when $s=0$, ordinary spectral truncation gives the result.
\end{proof}

\begin{corollary}[Separating mandatory and optional inverse cost]\label{budget:cor:envelope}
Write
\[
 \phi=\tr((Q^TKQ)^{-1}),\qquad
 \theta=\frac{r-s}{2r-s},\qquad \omega=1-\theta.
\]
Then
\begin{equation}\label{budget:eq:envelope}
 \boxed{\phi\le V_{\rm opt}\le\theta\tr(K^{-1})+\omega\phi.}
\end{equation}
In particular, for $s=r/2$,
\begin{equation}\label{budget:eq:third}
 V_{\rm opt}\le\tfrac13\tr(K^{-1})+\tfrac23\phi.
\end{equation}
\end{corollary}
\begin{proof}
The full block-inverse trace is
\[
 \tr(K^{-1})=\tr(A^{-1})+\tr(\mathcal S^{-1}\mathcal M).
\]
The sum of the smallest $r-s$ among $2r-s$ nonnegative eigenvalues is at most the fraction $\theta$ of their total. Substitute in \eqref{budget:eq:optvalue}. The lower bound follows from the same formula.
\end{proof}

The computation uses a Schur complement, a symmetric definite generalized eigenproblem, and orthonormalization, all at dimension $O(r)$. Its dense arithmetic cost is $O(r^3)$. It changes only the optional completion when the prescribed sample space has rank at most $r$. On full-rank truncated samples there may be no optional dimension to optimize. It should not be described as automatically repairing every noisy basis decision.

\begin{proposition}[A mandatory direction can make the inverse guard infeasible]\label{budget:prop:infeasible}
For every $r\ge2$ and $\varepsilon\in(0,1)$, there are full-row-rank predesigns with both child Grams equal to $\Id_r$ and a fixed one-dimensional mandatory space such that
\begin{equation}\label{budget:eq:impossible}
 V_{\rm opt}=\frac1\varepsilon+\frac1{2-\varepsilon}+r-2,
 \qquad \min(V_a,V_b)=r.
\end{equation}
Thus no finite fixed $\gamma$ can guarantee both preservation of every prescribed signal direction and \eqref{budget:eq:guard}.
\end{proposition}
\begin{proof}
Take
\[
 K=\begin{pmatrix}\Id_r&(1-\varepsilon)e_1e_1^T\\
 (1-\varepsilon)e_1e_1^T&\Id_r\end{pmatrix},\qquad
 Q=\frac1{\sqrt2}\begin{bmatrix}e_1\\-e_1\end{bmatrix}.
\]
The required direction is an eigenvector of $K$ with eigenvalue $\varepsilon$. The other eigenvalues are $2-\varepsilon$ and $1$ with multiplicity $2r-2$. Any feasible compression pays at least $1/\varepsilon$. The best remaining $r-1$ directions pay $1/(2-\varepsilon)+r-2$. Realize $K$ as $GG^T$, padding a square root with zero query columns if necessary. Both child inverse traces equal $r$.
\end{proof}

This is a local algebraic obstruction to a universal deterministic guard with arbitrary forced signal. It is not a probabilistic lower bound for the union learner. In a statistical proof the frequency and source-weight of such configurations must be estimated; deleting the required signal is not an acceptable substitute.

\section{A conditional closure using only the forced-signal inverse}\label{budget:sec:forced}
The optimizer gives another route besides a pointwise guard. Suppose every ordinary internal node has a prescribed rank-$s$ frame $Q_I$ and uses Theorem~\ref{budget:thm:completion} for its completion. Keep $s$ fixed in this theorem. Set
\[
 \phi_I=\tr((Q_I^TG_IG_I^TQ_I)^{-1}),\quad f_I=\E\phi_I,\quad
 b_I=\frac{\E[\phi_I\F{w_I}^2]}{q-r}.
\]
No independence between $\phi_I$ and $w_I$ is asserted. In particular $b_I$ is not automatically $f_Im_I$.

\begin{theorem}[Forced-core moment interface]\label{budget:thm:forced}
Suppose, conditional on the fixed side information, there are constants $f_*,b_*$ independent of height and node such that
\begin{equation}\label{budget:eq:forcedhyp}
 f_I\le f_*,\qquad b_I\le b_*m_I
\end{equation}
at all ordinary internal nodes. Assume $2\theta\kappa<1$, with $\theta$ from Corollary~\ref{budget:cor:envelope}. Define
\begin{align*}
 \overline v&=\max\left\{\rho_0,\frac{\omega f_*}{1-2\theta\kappa}\right\},\\
 \overline a&=\max\left\{\rho_0,
 \frac{\theta\kappa\overline v+\omega b_*}{1-\theta\kappa}\right\}.
\end{align*}
Then
\[
 v_I\le\overline v,\qquad a_I\le\overline a\,m_I
\]
at every ordinary node, and
\begin{equation}\label{budget:eq:forcedconclusion}
 \boxed{\E\F{E_{\rm root}^{\rm out}}^2+\E\F{c_{\rm root}}^2
 \le\left[1+\rho_0+
 \frac{(\kappa-1)\overline a+\kappa\overline v}{1-\lambda}\right]\beta.}
\end{equation}
The final root need not use the optimized completion.
\end{theorem}
\begin{proof}
Apply \eqref{budget:eq:envelope} pathwise and then \eqref{budget:eq:frametrace} to obtain
\[
 v_I\le\theta\kappa(v_a+v_b)+\omega f_I.
\]
For the weighted moment, replace the full-predesign term before conditioning on $G_I$, as in Lemma~\ref{budget:lem:inherited}. Retain the forced term as its joint moment. This gives
\[
 a_I\le\theta\lambda\kappa(a_a+a_b+m_av_b+m_bv_a)+\omega b_I.
\]
The definitions of $\overline v$ and $\overline a$ close these two inductions, using $m_I=\lambda(m_a+m_b)$. Substitute their conclusions into \eqref{budget:eq:cost}. The sum of $m_a+m_b$ over merges at child height $t$ is $\lambda^t\beta$. A geometric sum and \eqref{budget:eq:totalenergy} prove \eqref{budget:eq:forcedconclusion}.
\end{proof}

When $s=r/2$, $\theta=1/3$, and $2\theta\kappa<1$ holds whenever $q>4r+1$. The maintained numerical ratio $q=8r$ satisfies this condition. At $r=2,q=16$ one obtains
\[
 \overline v=\max\{2/13,22f_*/7\},\qquad
 \overline a=\max\{2/13,(13\overline v+22b_*)/20\}.
\]
This is a conditional theorem, not a proof of \eqref{budget:eq:forcedhyp} for the union's selected signal. Its significance is the separation: optional completion is now paid explicitly, leaving two moments of a smaller \emph{mandatory signal} Gram. These forced directions may still be adaptive. Treating their Gram as Gaussian without checking its law would reopen the earlier dependence error. Variable prescribed ranks require their own $\theta$ accounting; the $s=r/2$ coefficient cannot be silently assigned to a zero-rank or full-rank sample.

\section{Diagnostics and the unchanged union}\label{budget:sec:checks}
The scripts are reproducible finite-dimensional diagnostics, not proofs of HSS risk.

\paragraph{Stable nonlinear selections.}
There are 5,040 regression-tree realizations in six configurations, with $r=1,2,3$, $q=8r$, heights up to eight, and one to three response columns. The choices are leading-Gram selection and a source-dependent proposal guarded at $\gamma=1.1$. The largest register residual is below $9\times10^{-16}$ and the largest energy-identity residual below $9\times10^{-16}$. The largest normalized deviation of the innovation mean from its prediction is 2.60 reported standard errors; the largest weighted-cost mean deviation is 2.38. Representative energies, with $\beta=1$, are
\begin{center}
\begin{tabular}{lrrr}
\toprule Rule & $(r,H)$ & Mean physical plus retained & Uniform bound\\\midrule
Leading Gram & $(2,4)$ & $1.58280\ \pm0.00276$ & $375/143$\\
Leading Gram & $(2,8)$ & $1.64968\ \pm0.00308$ & $375/143$\\
Source proposal, guarded & $(2,4)$ & $1.62740\ \pm0.00551$ & $690/143$\\
Source proposal, guarded & $(2,7)$ & $1.70982\ \pm0.00516$ & $690/143$\\\bottomrule
\end{tabular}
\end{center}
The displayed uncertainties are standard errors, not confidence bounds. A guarded proposal may discard its requested direction on fallback, so these tests are not a signal-preserving HSS implementation.

\paragraph{The model counterexample.}
There are 13,250 independent realizations of the two-subtree construction for $r=1,2$. For $r=1,q=8$, the means of actual physical energy at $h=0,2,4,6,8$ are respectively $1.3903$, $2.0994$, $3.0339$, $5.4789$, and $8.2658$. The lower bound is analytical, not fitted to these means.

\paragraph{Completion optimizer.}
One hundred positive-definite problems test the optimum formula, containment, and envelope, with 2,500 additional random feasible completions. The largest normalized formula error is below $5.0\times10^{-14}$ and containment error below $1.1\times10^{-15}$. All random competitor costs exceed the computed optimum. Six explicit ill-conditioned-signal examples verify \eqref{budget:eq:impossible}. Sampling competitors is only a check; optimality is proved by Theorem~\ref{budget:thm:completion}.

\paragraph{Local tests along the actual maintained trajectory.}
A separate diagnostic follows the \emph{unchanged} finite-noise clean row hierarchy from the trapping family. At each current node it computes the best rank-two completion containing that node's fitted rank-one target, without feeding this hypothetical completion into later nodes. There are 290 trees and 6,470 local comparisons, using $\delta=10^{-3}$ and $0.2$. A representative subset at $\delta=10^{-3}$ is
\begin{center}
\begin{tabular}{rrrrr}
\toprule $h$ & Local nodes & Current: $1.1$ & Optimal: $1.1$ & Optimal: $7/6$\\\midrule
0 & 50 & 25 & 7 & 2\\
2 & 350 & 210 & 67 & 46\\
4 & 930 & 525 & 154 & 102\\
6 & 1,905 & 1,092 & 349 & 246\\\bottomrule
\end{tabular}
\end{center}
A ``failure'' here means that the measured inverse ratio exceeds the stated sufficient guard, not failure of approximation. These nodes are dependent and the pooled counts are not probability estimates. The optimized completion improves the local inverse objective but does not always meet even the critical guard. Thus we do not apply Theorem~\ref{budget:thm:guard} to the maintained learner on the basis of these tests. Nor do these comparisons report the risk of a recursively optimized hierarchy.

\section{Consequences for the general target}
The unpaid local-equivariant inverse sum has now been divided into distinct mathematical cases. A uniform theorem for \emph{every} local-equivariant selection is false. An explicit relative inverse-budget condition does give a uniform nonlinear physical theorem, even with growing inverse moments. The exact completion optimizer removes avoidable padding cost while preserving a specified signal. Its unavoidable remainder is the mandatory-signal inverse term, including its joint moment with the innovation.

Two HSS interfaces are still missing. First, the actual sample-selected signal must be shown to satisfy a suitable expected inverse budget or forced-core bound without discarding important directions. The pointwise guard is sufficient, not necessary, and its local violation is not a counterexample to the union's expected risk. Second, the general coupled-cut trajectory must satisfy an appropriate conditional analogue of subtree locality. Shared exterior probe rows and unequal templates remain outside the model proved here.

The desired premise remains
\[
 \E\dist_F(A,\mathcal S_\cup(A))^2\le C_bL\OPT_k(A).
\]
No all-input proof or counterexample for this union is established in this chapter. The earlier exact two-view and mean-square small-noise results are unchanged. A distinguished-cut or abstract regression theorem is not silently substituted for a whole coefficient-space bound.

There is also a depth-accounting distinction. At $\gamma\lambda=1$, the local regression theorem is $O(H)$, which can already be useful if its source budget is $O(\OPT_k(A))$. If introducing general cut sources has already cost a factor $L$, another $O(H)$ gain could lead to $L^2$. The depth-independent version or a sharper source summation would then still be needed.

The maintained query count is unchanged: $32k$ shared training products plus $32k$ fresh refit products. A local completion can be computed without new products, but using it recursively changes the basis learner and must be analyzed as such. The refit independence and exact offline rank-reduction premises remain separate from these new selection results.


\section{A stronger scalar version of the counterexample}\label{budget:sec:scalar}
For $r=1$, let the two bad-subtree child row lengths be $a,b$, and let $c$ be the inner product of their independent unit row frames. Testing the $2\times2$ Gram against $[b,-\operatorname{sign}(c)a]^T$ gives
\[
 \sigma_{\min}(G)^2\le\frac{2a^2b^2(1-|c|)}{a^2+b^2}.
\]
The inverse of the retained weaker eigenvalue therefore satisfies
\[
 V_I\ge\frac{V_a+V_b}{2(1-|c|)}.
\]
The angle is independent of the child lengths. Its density on $(-1,1)$ is
\[
 C_q(1-c^2)^{(q-3)/2},\qquad
 C_q=\frac{\Gamma(q/2)}{\sqrt\pi\,\Gamma((q-1)/2)}.
\]
For $q>3$, direct beta integration gives
\[
 \chi_q:=\E\frac1{1-|c|}
 =\frac{q-2+2C_q}{q-3}.
\]
Indeed $1/(1-|c|)=(1+|c|)/(1-c^2)$; the two expectations are $(q-2)/(q-3)$ and $2C_q/(q-3)$. Consequently $v_B\ge\rho_0\chi_q^h$, and the proof of Theorem~\ref{budget:thm:counter} gives
\[
 \E\F{E_{\rm root}^{\rm out}}^2
 \ge1+\rho_0+\kappa\rho_0(\lambda\chi_q)^h.
\]
At $q=8$,
\[
 \chi_8=\frac65+\frac{32}{25\pi},\qquad
 \lambda\chi_8=\frac{36}{35}+\frac{192}{175\pi}>1.3778.
\]
This sharpening is not needed for the general-rank impossibility, but makes the mechanism visible at smaller test heights.

\endgroup

% ===== CORE =====
\begingroup
\let\E\relax
\newcommand{\E}{\mathbb E}
\let\R\relax
\newcommand{\R}{\mathbb R}
\let\Id\relax
\newcommand{\Id}{\mathrm I}
\let\F\relax
\newcommand{\F}[1]{\lVert #1\rVert_F}
\let\op\relax
\newcommand{\op}[1]{\lVert #1\rVert_{\rm op}}
\let\tr\relax
\newcommand{\tr}{\operatorname{tr}}
\let\ran\relax
\newcommand{\ran}{\operatorname{range}}
\let\rank\relax
\newcommand{\rank}{\operatorname{rank}}
\let\diag\relax
\newcommand{\diag}{\operatorname{blockdiag}}
\let\dist\relax
\newcommand{\dist}{\operatorname{dist}}
\let\OPT\relax
\newcommand{\OPT}{\operatorname{OPT}}
\chapter{Fourth-moment innovation decay and additive inverse budgets}\label{ch:core}
\begin{scopebox}
\textbf{Update, 20 September 2026.} The additive inverse budget and fixed-physical-core constants below remain valid under their displayed assumptions. Later source notes prove further capture, conditional-moment, and sufficiently large-rank fixed-source closure results, as well as fixed-guard transfer and fresh-reset cases. These advances supersede the suggestion that the only available local closure is the fixed-core marginal estimate, but they do not establish arbitrary HSS applicability. For the current synchronized learner, the unresolved requirement is a joint physical source/return/invisible budget, not a blanket assertion that every learned design must have a Gaussian law or a uniform unweighted inverse moment.
\end{scopebox}

\begin{scopebox}
The fixed-physical-core theorems below are unconditional within the stated local regression model. The actual noisy selected-core inverse law and the general-cut dependence remain unproved; the uniform constants do not certify the end-to-end union learner.
\end{scopebox}
\thispagestyle{plain}
\begin{chaptersynopsis}
The relative-guard and mandatory-signal estimates of Chapter~\ref{ch:budget} required either a pointwise inverse guard or two separate moment hypotheses. We prove a different closure. First, the innovation in every admissible local, query-rotation-equivariant regression tree has a geometrically decaying fourth moment, even when its physical error is unstable. Second, inverse-optimal completion of a prescribed signal space obeys an additive inverse budget: its retained inverse trace is at most the smaller child inverse trace plus the forced-signal inverse trace. These two facts convert a marginal second-moment bound for the forced inverse into a depth-independent bound for physical regression error, without factoring a dependent product.

As an unconditional application, consider fixed nested physical signal spaces of ranks at most $k$, surrounded by rank-$2k$ completions selected adaptively from the same width-$16k$ Gaussian probe. Inverse-optimal completion gives physical-plus-retained error at most $(6694/1001+2\sqrt3/77)\beta<6.733\beta$ at every depth. When every ordinary mandatory rank equals $k$, a sharper certificate is $(455+4\sqrt3)\beta/153<3.020\beta$. The completions are genuinely probe-adaptive; they are not treated as fixed Gaussian coordinate maps. The mandatory physical spaces, however, are fixed independently of that probe. The actual noisy HSS union does not yet have the needed mandatory-core moment law, and general cut coupling remains outside the local model. No all-input HSS approximation theorem is asserted.
\end{chaptersynopsis}


\section{Scope and relation to the inverse-budget theorems}
The target remains an $O(k)$-query algorithm returning an HSS-rank-at-most-$k$ matrix with
\[
 \E\F{A-\widehat B}^2\le CL\OPT_k(A),\qquad
 \OPT_k(A)=\min_{B\in\mathcal H_k}\F{A-B}^2.
\]
The published fresh-level result provides the logarithmic-query benchmark \cite{AmselHSS}. The maintained shared-query candidate runs an annihilate-first hierarchy and an un-nullified hierarchy separately on the same observations, then takes physical nested unions. With per-hierarchy rank $2k$ and width $16k$ per stream, its training cost is $32k$ products. The independent rank-$4k$ refit costs another $32k$ products. Their full proofs are in Chapters~\ref{ch:uniform} and~\ref{ch:bicriteria}.

Chapter~\ref{ch:budget} proves a uniform nonlinear theorem under $V_I\le\gamma\min(V_a,V_b)$, and a signal-preserving completion theorem conditional on both $\E\phi_I$ and $\E[\phi_I\|w_I\|_F^2]$. The present work does not claim those hypotheses for the actual union. It replaces the pointwise guard by an automatically valid additive inequality for the optimal completion, and replaces the difficult joint moment hypothesis by a marginal $L^2$ hypothesis. A marginal $L^2$ hypothesis is stronger than a marginal $L^1$ hypothesis; we do not claim logical implication from all the old premises to the new ones. Its advantage is that it is directly verifiable for fixed physical signal cores.

The fourth-moment innovation bound and the additive spectral inequality below are new derivations. The second-moment regression identities of Chapters~\ref{ch:nonlinear} and~\ref{ch:budget} are reproduced with their conditioning argument. All numerical checks are diagnostics, not proofs of expected HSS risk.

\section{The local regression model and its exact second moments}\label{core:sec:model}
Let a balanced binary tree have height $H$. At each physical leaf $\ell$, fix an orthonormal rank-$r$ frame $U_\ell$ and a source matrix $E_\ell$ with any fixed number of response columns, satisfying $U_\ell^TE_\ell=0$. Set
\[
 \beta_\ell=\F{E_\ell}^2,\qquad
 \beta_I=\sum_{\ell\subset I}\beta_\ell,\qquad \beta=\beta_{\rm root}.
\]
Draw a standard Gaussian matrix $X$ with $q$ columns. Independent side information can be fixed throughout, so all source budgets in a conditional argument are deterministic.

At a merge $I=a\cup b$, set
\begin{equation}\label{core:eq:model}
 \begin{split}
 W_I&=\diag(U_a,U_b),\qquad G_I=W_I^TX_I,\\
 F_I&=[E_a^{\rm out};E_b^{\rm out}],\qquad U_I=W_IT_I,\\
 B_I&=S_I^T(G_I^\dagger)^TX_I^TF_I,\\
 E_I^{\rm out}&=F_I-W_IS_IB_I.
 \end{split}
\end{equation}
where $[T_I,S_I]$ is an orthogonal split of $\R^{2r}$. Ordinary nodes retain $r$ dimensions. The final root may retain any rank at most $2r$.

All choices must be subtree-local and unchanged in physical coordinates by right rotations $X_I\mapsto X_IQ$. They can otherwise be nonlinear and source-dependent. A fixed, measurable, rotation-invariant tie convention is understood. These hypotheses are substantive: arbitrary general HSS cut data may not satisfy subtree locality.

At an ordinary node define
\[
 R_I=U_I^TX_I,\quad
 c_I=(R_I^\dagger)^TX_I^TE_I^{\rm out},\quad
 w_I=N_{R_I}X_I^TE_I^{\rm out},\quad
 N_R=\Id-R^\dagger R.
\]
With $d_I=[c_a;c_b]$ and $z_I=(G_I^\dagger)^T(w_a+w_b)$, direct multiplication gives
\begin{equation}\label{core:eq:register}
 \boxed{c_I=T_I^T(d_I+z_I),\quad B_I=S_I^T(d_I+z_I),\quad
 w_I=N_{G_I}(w_a+w_b).}
\end{equation}
For example, $X_I^TF_I=G_I^Td_I+w_a+w_b$. Subtracting $G_I^TS_IB_I$ leaves $G_I^TT_Ic_I+w_I$, with $w_I$ orthogonal to the whole row space of $G_I$.

Put
\begin{equation}\label{core:eq:params}
 d_0=q-r,\quad d=q-2r,\quad
 \lambda=\frac d{d_0},\quad
 \kappa=\frac{d_0-1}{d-1},\quad
 \rho_0=\frac r{d_0-1}.
\end{equation}
Assume $q>2r+1$. Define
\[
 V_I=\tr((R_IR_I^T)^{-1}),\quad v_I=\E V_I,\quad
 m_I=\frac{\E\F{w_I}^2}{d_0},\quad
 a_I=\frac{\E[V_I\F{w_I}^2]}{d_0}.
\]

\begin{lemma}[Exact conditional and second-moment interfaces]\label{core:lem:second}
All finite-tree designs have full row rank almost surely. At ordinary height $t$,
\begin{equation}\label{core:eq:innovationsecond}
 m_I=\lambda^t\beta_I.
\end{equation}
At a merge,
\begin{equation}\label{core:eq:cost}
 \boxed{\E\F{z_I}^2=(\kappa-1)(a_a+a_b)+\kappa(m_av_b+m_bv_a),
 \qquad \E\langle d_I,z_I\rangle_F=0.}
\end{equation}
The accumulated physical energy satisfies
\begin{equation}\label{core:eq:total}
 \boxed{\E\F{E_{\rm root}^{\rm out}}^2+\E\F{c_{\rm root}}^2
 =(1+\rho_0)\beta+\sum_{I\ {\rm internal}}\E\F{z_I}^2.}
\end{equation}
\end{lemma}
\begin{proof}
Each normalized row frame $(R_IR_I^T)^{-1/2}R_I$ is uniform and independent of the invariant state under right rotations. Independent subtrees have independent frames. Their two row spaces have direct sum almost surely, and every orthogonal row compression of the resulting full-row-rank predesign is full row rank. This gives an induction from the Gaussian leaf designs. It does not make the selected designs Gaussian.

Conditional on $R_I$, its stabilizer acts on $\ker R_I$, giving
\[
 \E[w_Iw_I^T\mid R_I]=\eta_I N_{R_I},\qquad
 \E[w_I\mid R_I]=0,\qquad
 \E[c_Iw_I^T\mid R_I]=0,
\]
where $\eta_I$ is a function of the retained Gram $R_IR_I^T$. Complement reflection gives the two zero moments. At a merge,
\[
 \E[w_Iw_I^T\mid G_I]=(\eta_a+\eta_b)N_{G_I}.
\]
Taking traces proves \eqref{core:eq:innovationsecond}.

The two independent frame integrals, with child Grams fixed, are
\begin{align*}
 \E_{\rm frames}\F{(G^\dagger)^TN_{R_a}}^2&=(\kappa-1)V_a+\kappa V_b,\\
 \E_{\rm frames}\tr((GG^T)^{-1})&=\kappa(V_a+V_b).
\end{align*}
To verify the coefficient, write the child normalized frames as $Q_a,Q_b$. The diagonal inverse blocks of $[Q_a;Q_b][Q_a;Q_b]^T$ are $(\Id-Q_aQ_b^TQ_bQ_a^T)^{-1}$ and its counterpart. Their expectation is $\kappa\Id$. After fixing $Q_a$, represent $Q_b$ by normalizing an $r\times q$ Gaussian $Z=[Z_1,Z_2]$, where $Z_2$ has $q-r$ columns. The mean inverse trace is
\[
 r+\E\tr((Z_2Z_2^T)^{-1}Z_1Z_1^T)=r+\frac{r^2}{q-2r-1}=r\kappa.
\]
The Gaussian inverse mean used here is standard \cite{HMT}. Subtracting the identity gives $\kappa-1$ for the same-child coefficient. The conditional innovation covariance then proves \eqref{core:eq:cost}; same-child weighted products stay inside $a_I$. Cross-child factors are legitimate only with independent side information fixed.

Finite-height integrability follows from positive-integrand bounds
\[
 v_I\le\kappa(v_a+v_b),\qquad
 a_I\le\lambda\kappa(a_a+a_b+m_av_b+m_bv_a),
\]
starting at $v_\ell=\rho_0$ and $a_\ell=\rho_0\beta_\ell$. The second inequality replaces $V_I$ by the preselection trace before conditioning. Finally, physical discarded sectors are orthogonal, so $\|E_{\rm root}^{\rm out}\|_F^2=\beta+\sum_I\|B_I\|_F^2$. The orthogonal split in \eqref{core:eq:register}, its zero expected cross term, and the leaf coefficient energy $\rho_0\beta$ prove \eqref{core:eq:total}.
\end{proof}

\section{Fourth-moment innovation decay for every admissible selection}\label{core:sec:fourth}
Define
\begin{equation}\label{core:eq:chi}
 \chi=\frac{d(d+2)}{d_0(d_0+2)},\qquad \nu=\sqrt\chi.
\end{equation}
Then $0<\lambda^2\le\chi<1$.

\begin{theorem}[A source fourth moment that does not accumulate inverses]\label{core:thm:fourth}
At every ordinary node of height $t$,
\begin{equation}\label{core:eq:fourth}
 \boxed{\E\F{w_I}^4\le d_0(d_0+2)\chi^t\beta_I^2.}
\end{equation}
The conclusion holds for any fixed number of response columns and for all the local-equivariant choices in Section~\ref{core:sec:model}. It does not require a stable inverse budget. For one response column and a single nonzero physical leaf source, equality holds.
\end{theorem}
\begin{proof}
Conditional on $R_j$ and the response Gram $C_j=w_j^Tw_j$, the orientation of $w_j$ is uniform in the $d_0$-dimensional complement of $R_j$. This follows by disintegration under its stabilizer; it does not condition away the orientation by retaining the whole vector $w_j$.

Fix both child designs and both response Grams. Let $P=N_G$, of rank $d$, and put $A=Pw_a$, $B=Pw_b$, $s_j=\tr C_j$. The two matrices are conditionally independent and centered. The fourth moment of the projection of a uniform unit vector in $\R^{d_0}$ onto a fixed $d$-plane is $d(d+2)/(d_0(d_0+2))$. Applying this columnwise and using Cauchy--Schwarz for products of column energies gives
\[
 \E[\F A^4\mid G,C_a,C_b]\le\chi s_a^2,
 \qquad \E[\F B^4\mid G,C_a,C_b]\le\chi s_b^2.
\]
Moreover,
\[
 \E\F A^2=\lambda s_a,\quad \E\F B^2=\lambda s_b,\quad
 \E\langle A,B\rangle_F^2=\frac d{d_0^2}\tr(C_aC_b)
 \le\frac d{d_0^2}s_as_b
\]
under this conditioning. Odd mixed terms vanish by reflection. Therefore
\begin{equation}\label{core:eq:fourthmerge}
 \E[\F{A+B}^4\mid G,C_a,C_b]
 \le\chi(s_a^2+s_b^2)+\frac{2d(d+2)}{d_0^2}s_as_b.
\end{equation}
Subtree independence gives $\E(s_as_b)=d_0^2\lambda^{2t}\beta_a\beta_b$ for children of height $t$.

At a leaf, conditional on its Gaussian retained design, the innovation is a Gaussian matrix in its $d_0$-dimensional complement. Its fourth moment is
\[
 d_0^2\beta_\ell^2+2d_0\tr((E_\ell^TE_\ell)^2)
 \le d_0(d_0+2)\beta_\ell^2.
\]
Induct in \eqref{core:eq:fourthmerge}. The mixed term is bounded by the desired cross term because $\lambda^{2t}\le\chi^t$. This proves \eqref{core:eq:fourth}. With one response and one active leaf, the columnwise bound is equality and each merge has just one nonzero source, so equality propagates.
\end{proof}

\begin{corollary}[Paying a dependent multiplier using one marginal moment]\label{core:cor:holder}
If $D_I\ge0$ is any random scalar with $\|D_I\|_{L^2}\le\Phi$, it may depend on every quantity at the node. Then
\begin{equation}\label{core:eq:holder}
 \boxed{\frac{\E[D_I\F{w_I}^2]}{d_0}
 \le B\nu^t\beta_I,\qquad
 B=\Phi\sqrt{\frac{d_0+2}{d_0}}.}
\end{equation}
\end{corollary}
\begin{proof}
Cauchy--Schwarz and \eqref{core:eq:fourth} give the result. No independence is used.
\end{proof}
At $r=2,q=16$, $\lambda=6/7$ and $\nu=\sqrt3/2$. The loss from second to fourth moments is small but not zero. One must retain $\nu$, not replace it by $\lambda$. Innovation decay alone still does not rule out the unstable physical selection in (Chapter~\ref{ch:budget}); that selection has uncontrolled inverse multipliers.

\section{An additive bound for optimal signal-preserving completion}\label{core:sec:additive}
Let $K=GG^T\succ0$ have size $2r$, and let $Q$ be an orthonormal mandatory frame of rank $s\le r$. Denote
\[
 \phi=\tr((Q^TKQ)^{-1}),
\]
with $\phi=0$ for $s=0$. Minimize $\tr((T^TKT)^{-1})$ over orthonormal rank-$r$ frames containing $Q$.

For $0<s<r$, complete $Q$ to $[Q,Z]$ and write
\[
 [Q,Z]^TK[Q,Z]=\begin{pmatrix}A&B_0\\B_0^T&D\end{pmatrix},\quad
 \mathcal S=D-B_0^TA^{-1}B_0,\quad
 \mathcal M=\Id+B_0^TA^{-2}B_0.
\]
The exact optimizer of Chapter~\ref{ch:budget} has value
\begin{equation}\label{core:eq:opt}
 V_{\rm opt}=\phi+\sum_{j=1}^{r-s}\mu_j,
\end{equation}
where $\mu_j$ are the increasing eigenvalues of $\mathcal S^{-1/2}\mathcal M\mathcal S^{-1/2}$. To see this directly, write a feasible frame as $[Q,ZV]$. Block inversion gives the objective
\[
 \phi+\tr\bigl((V^T\mathcal SV)^{-1}V^T\mathcal MV\bigr).
\]
The generalized Rayleigh--Ritz principle yields \eqref{core:eq:opt}. The minimizing completion is computed in $O(r^3)$ arithmetic with no new matrix-vector products.

\begin{theorem}[An automatic additive budget]\label{core:thm:additive}
Let $\lambda_1(K)\ge\cdots\ge\lambda_{2r}(K)>0$. For every $0\le s\le r$,
\begin{equation}\label{core:eq:spectral}
 \boxed{V_{\rm opt}\le\phi+\sum_{i=s+1}^{r}\lambda_i(K)^{-1}.}
\end{equation}
In particular, if the principal child Grams are $K_a,K_b$,
\begin{equation}\label{core:eq:additive}
 \boxed{V_{\rm opt}\le\min\{\tr(K_a^{-1}),\tr(K_b^{-1})\}+\phi.}
\end{equation}
This holds on every positive-definite realization, including those where no fixed relative guard is feasible while preserving the signal.
\end{theorem}
\begin{proof}
In the $[Q,Z]$ coordinates, block inversion shows that
\[
 Y=K^{-1}-\begin{pmatrix}A^{-1}&0\\0&0\end{pmatrix}\succeq0.
\]
It has $s$ zero eigenvalues, and its other eigenvalues are exactly the $\mu_j$ in \eqref{core:eq:opt}: factor it as $[-A^{-1}B_0;\Id]\mathcal S^{-1}[-B_0^TA^{-1},\Id]$. Since $Y\preceq K^{-1}$, eigenvalue monotonicity gives $\mu_j\le\lambda_{s+j}(K)^{-1}$, with the eigenvalues of $K$ ordered decreasingly. Summing proves \eqref{core:eq:spectral}. Principal-submatrix interlacing gives $\lambda_i(K)\ge\lambda_i(K_a),\lambda_i(K_b)$ for $i\le r$, proving \eqref{core:eq:additive}. At $s=0$, the optimum is the leading-Gram space; at $s=r$, the optional sum is empty and $V_{\rm opt}=\phi$.
\end{proof}

The additive inequality does not assert that forced cost disappears. The nearly cancelling mandatory direction in Chapter~\ref{ch:budget} still has $\phi=1/\varepsilon$. It is now an explicit additive term, rather than a failed feasibility constraint. If a noisy truncated sample already has rank $r$, there are no optional directions; the new inequality does not manufacture rank slack.

\section{Depth-independent stability with a marginal core moment}\label{core:sec:stability}
\begin{theorem}[Additive inverse-budget stability]\label{core:thm:stability}
Suppose each ordinary internal node satisfies
\begin{equation}\label{core:eq:deltaguard}
 V_I\le\min(V_a,V_b)+\Delta_I,\qquad \Delta_I\ge0,
\end{equation}
with uniform conditional moment bounds, independent side information fixed,
\[
 \E\Delta_I\le F_*,\qquad \|\Delta_I\|_{L^2}\le\Phi_*.
\]
No guard is required at the final root. Set
\begin{equation}\label{core:eq:closedconstants}
 B_*=\Phi_*\sqrt{\frac{d_0+2}{d_0}},\qquad
 \mathcal U=\frac{\rho_0}{1-\lambda}+\frac{F_*\lambda}{(1-\lambda)^2},\qquad
 \mathcal A=\rho_0+\lambda\mathcal U+\frac{B_*\nu}{1-\nu}.
\end{equation}
Then at every finite height,
\begin{equation}\label{core:eq:additiveuniform}
 \boxed{\E\F{E_{\rm root}^{\rm out}}^2+\E\F{c_{\rm root}}^2
 \le[1+\rho_0+(\kappa-1)\mathcal A+\kappa\mathcal U]\beta.}
\end{equation}
The full retained inverse means may grow linearly with height.
\end{theorem}
\begin{proof}
The unweighted induction is $v_I\le\rho_0+tF_*$ at height $t$. For the weighted term, retain the dependent additive cost and condition only the child-measurable minimum:
\begin{align*}
 a_I&\le\frac{\E[\min(V_a,V_b)\F{w_I}^2]}{d_0}
       +\frac{\E[\Delta_I\F{w_I}^2]}{d_0}\\
 &\le\lambda(m_av_b+m_bv_a)+B_*\nu^t\beta_I\\
 &\le\left[\lambda^t(\rho_0+(t-1)F_*)+B_*\nu^t\right]\beta_I,
 \qquad t\ge1.
\end{align*}
The opposite-child charge in the second line follows from
$\min(V_a,V_b)\eta_a\le V_b\eta_a$ and its counterpart, then subtree independence. The additive product is controlled by Corollary~\ref{core:cor:holder}, not by factorization. At leaves, $a_\ell=\rho_0\beta_\ell$.

At each child height $t$, source budgets sum to $\beta$. The sum of the upper bounds for $\lambda^t v_I$ is $\mathcal U$, and the sum of the upper bounds for $a_I/\beta_I$ is $\mathcal A$. Insert these bounds in \eqref{core:eq:cost}, sum over levels, and use \eqref{core:eq:total}. Finite geometric sums are bounded by the displayed infinite sums of nonnegative terms.
\end{proof}

\begin{corollary}[Signal-preserving completion with one mandatory-core premise]\label{core:cor:core}
If every ordinary node uses an inverse-optimal completion and
\[
 \E\phi_I\le F_*,\qquad (\E\phi_I^2)^{1/2}\le\Phi_*
\]
uniformly, then \eqref{core:eq:additiveuniform} holds with $\Delta_I=\phi_I$. The mandatory ranks may vary arbitrarily between zero and $r$. There is no premise on $\E[\phi_I\|w_I\|_F^2]$ and no half-rank requirement.
\end{corollary}

The result also permits moderately growing mandatory costs. If the two moment bounds are $F_*\zeta^t$ and $\Phi_*\zeta^t$, respectively, at height $t$, with $\zeta\nu<1$, replace
\[
 \mathcal U\ \text{by}\ 
 \frac{\rho_0}{1-\lambda}+\frac{F_*\lambda\zeta}{(1-\lambda)(1-\lambda\zeta)},
 \qquad
 \frac{B_*\nu}{1-\nu}\ \text{by}\ \frac{B_*\nu\zeta}{1-\nu\zeta}.
\]
The same proof applies. Thus even a uniform marginal inverse moment is sufficient rather than necessary; the innovation supplies a discount in depth.

\section{Unconditional theorem for fixed physical signal cores}\label{core:sec:anchored}
Prescribe physical orthonormal frames $Q_I^{\rm phys}$ before $X$, with ranks $s_I\le k$, nested in the sense that the parent range lies in the block sum of its child ranges. They may be selected from independent side information. Require the fixed physical leaf frames $U_\ell$ to contain $Q_\ell^{\rm phys}$. At each ordinary merge use the inverse-optimal rank-$2k$ completion preserving
\[
 Q_I=W_I^TQ_I^{\rm phys}.
\]
Nestedness ensures feasibility and makes $Q_I$ orthonormal. The optional directions are chosen from the current, adaptively compressed Gaussian probe, not fixed beforehand.

\begin{lemma}[The mandatory Gram, not the whole selected Gram, is Gaussian]\label{core:lem:coregauss}
For these selections,
\[
 Q_I^TG_IG_I^TQ_I=(Q_I^{\rm phys})^TX_IX_I^TQ_I^{\rm phys}.
\]
If $s=s_I>0$ and $q>s+3$, then
\begin{equation}\label{core:eq:coregaussian}
 \E\phi_I=\frac{s}{q-s-1},\qquad
 \E\phi_I^2\le\frac{s^2}{(q-s-1)(q-s-3)}.
\end{equation}
For $s=0$, both costs vanish.
\end{lemma}
\begin{proof}
The algebra uses $W_IW_I^TQ_I^{\rm phys}=Q_I^{\rm phys}$; it is valid for the adaptive $W_I$. Since $Q_I^{\rm phys}$ is fixed before $X$, its $s\times q$ projection is standard Gaussian. A diagonal entry of the inverse Gram is the reciprocal of the squared norm of a row after projection off the other $s-1$ rows. That squared norm is chi-square with $q-s+1$ degrees of freedom. Its inverse first and second moments are $1/(q-s-1)$ and $1/[(q-s-1)(q-s-3)]$. Summing gives the mean, and $(\sum_{i=1}^s x_i)^2\le s\sum_i x_i^2$ gives the stated second moment. No independence of different inverse diagonal entries is used.
\end{proof}

\begin{theorem}[Adaptive completion around fixed nested cores]\label{core:thm:anchored}
Use retained rank $r=2k$ and width $q=16k$. For arbitrary nested fixed physical core ranks $0\le s_I\le k$, the procedure above satisfies, at every depth and for arbitrary fixed leaf sources,
\begin{equation}\label{core:eq:variableconstant}
 \boxed{\E\F{E_{\rm root}^{\rm out}}^2+\E\F{c_{\rm root}}^2
 \le\left(\frac{6694}{1001}+\frac{2\sqrt3}{77}\right)\beta
 <6.733\beta.}
\end{equation}
The final root need not use the optimizer or preserve a prescribed core.
\end{theorem}
\begin{proof}
All the constants are bounded by their $k=1$ values:
\[
 \begin{aligned}
 &\lambda=6/7,\quad \kappa\le13/11,\quad \rho_0\le2/13,\quad \nu\le\sqrt3/2,\\
 &F_*\le1/14,\quad \Phi_*\le1/\sqrt{168},\quad B_*\le1/\sqrt{147}.
 \end{aligned}
\]
The rank cap and Lemma~\ref{core:lem:coregauss} justify the last three bounds conditionally on independent side information. The right side of Theorem~\ref{core:thm:stability} is increasing in these nonnegative constants. At their upper endpoints,
\[
 \mathcal U=53/13,\qquad
 \mathcal A=358/91+\sqrt3/7.
\]
Substitution proves \eqref{core:eq:variableconstant}.
\end{proof}

This theorem is genuinely nonlinear: the padding at all internal nodes depends on the same probe and changes the future designs. It is also a finite-source theorem, not a derivative statement. But its mandatory physical spaces are fixed independently of that probe. On a general noisy HSS trajectory, the selected mandatory signal is not such a fixed space. Theorem~\ref{core:thm:anchored} cannot be substituted for that missing identification.

\section{Sharper certificate when the core occupies half the retained rank}\label{core:sec:sharp}
The additive bound handles rank variation. When every ordinary mandatory rank equals $k$ and $r=2k$, retain in addition the earlier averaging envelope
\begin{equation}\label{core:eq:average}
 V_I\le\frac13\tr((G_IG_I^T)^{-1})+\frac23\phi_I.
\end{equation}
Indeed the sum of the smallest $r-s$ generalized eigenvalues is at most their fraction $(r-s)/(2r-s)$ of the full sum. At $s=r/2$ this fraction is $1/3$.

Write $\alpha=\kappa/3$, $\omega=2/3$. With the marginal bounds $F_*,\Phi_*$ from Section~\ref{core:sec:stability}, set $B_*=\Phi_*\sqrt{(d_0+2)/d_0}$. Define scalar bounds
\begin{equation}\label{core:eq:recursions}
 \begin{split}
 v_0^*&=\rho_0,\quad v_t^*=2\alpha v_{t-1}^*+\omega F_*,\\
 a_0^*&=\rho_0,\quad
 a_t^*=\alpha\lambda a_{t-1}^*+\alpha\lambda^t v_{t-1}^*
            +\omega B_*\nu^t.
 \end{split}
\end{equation}
Then $v_I\le v_t^*$ and $a_I\le a_t^*\beta_I$. For the second inequality, the preselection inverse trace can be conditioned before selection, while its forced term is bounded by Corollary~\ref{core:cor:holder}. This gives
\[
 a_I\le\alpha\lambda(a_a+a_b+m_av_b+m_bv_a)+\omega B_*\nu^t\beta_I.
\]
It is not the old unproved joint-core moment premise.

If $2\alpha\lambda<1$, summing \eqref{core:eq:recursions} yields
\begin{equation}\label{core:eq:sharpconstants}
 \mathcal U_*:=\sum_{t\ge0}\lambda^tv_t^*
 =\frac{\rho_0+\omega F_*\lambda/(1-\lambda)}{1-2\alpha\lambda},\qquad
 \mathcal A_*:=\sum_{t\ge0}a_t^*
 =\frac{\rho_0+\alpha\lambda\mathcal U_*+\omega B_*\nu/(1-\nu)}{1-\alpha\lambda}.
\end{equation}
The physical certificate is $1+\rho_0+(\kappa-1)\mathcal A_*+\kappa\mathcal U_*$. Notice that the summation only requires $2\alpha\lambda<1$, not a uniform bound on every $v_t^*$.

\begin{corollary}[Half-rank fixed cores]\label{core:cor:sharp}
Under the hypotheses of Theorem~\ref{core:thm:anchored}, if every ordinary mandatory rank is exactly $k$, then
\begin{equation}\label{core:eq:sharp}
 \boxed{\E\F{E_{\rm root}^{\rm out}}^2+\E\F{c_{\rm root}}^2
 \le\frac{455+4\sqrt3}{153}\,\beta<3.020\beta.}
\end{equation}
\end{corollary}
\begin{proof}
Use the $k=1$ upper endpoints from Theorem~\ref{core:thm:anchored}. The recursions have nonnegative coefficients, so those endpoints dominate every $k$. They give
\[
 \mathcal U_*=88/65,\qquad
 \mathcal A_*=12034/9945+22\sqrt3/153.
\]
Equation~\eqref{core:eq:total} gives the displayed constant after simplification.
\end{proof}
The two numerical constants apply to different permitted core-rank patterns. We do not assign the half-rank coefficient to a variable-rank sample.

\section{Connection to the actual learner and the modified clean pilot}\label{core:sec:pilot}
The maintained union itself is unchanged by the preceding mathematical theorems. Its noisy selected core need not be a fixed physical comparator space. The remaining all-input premise is still
\[
 \E\dist_F(A,\mathcal S_\cup(A))^2\le C_bL\OPT_k(A).
\]
For an inverse-optimally completed variant, a useful sufficient local-model premise is now the marginal bound
\[
 \sup_I\E\left[\tr((Q_I^TG_IG_I^TQ_I)^{-1})^2\right]<\infty
\]
with the correct width/rank scaling. The theorem does not prove this premise for the actual adaptive noisy cores. The first-moment bound then follows by Cauchy--Schwarz, though a sharper separately known first moment improves constants. A full-rank truncated sample still leaves no padding freedom.

\paragraph{A rigorous reference-interface observation.}
On a clean fitted-target recursion, let $C$ denote a fixed physical target and suppose a parent retains $WLX_I^TC$, where $L=(G^\dagger)^T$. With $F=(\Id-WW^T)C$, direct algebra gives
\begin{equation}\label{core:eq:fitinterface}
 (\Id-P_{WT})C=F-WSS^TLX_I^TF.
\end{equation}
Suppose leaf residuals divided by a parameter $\delta$ converge, and the realized spaces and full-rank designs have limits as $\delta\downarrow0$. If the limiting hierarchy contains $C$, then $F_0=0$ and division of \eqref{core:eq:fitinterface} by $\delta$ gives exactly the regression model at the limiting splits. Derivatives of the inverse or split are not required for this step. If those limiting splits are inverse-optimal completions of fixed physical target cores, Corollary~\ref{core:cor:sharp} controls that limiting model even though its padding depends on $X$.

This is an implication with its limit hypotheses stated. We do not use numerical finite differences to assert those hypotheses almost surely for every noisy HSS trajectory. The earlier mean-square theorem for the unchanged canonical learner (Chapter~\ref{ch:nonlinear}) remains separate, with its existing constant and order of limits.

\paragraph{Recursive rather than one-node diagnostic.}
The clean-family implementation was modified only at internal full-range row steps: canonical rank-two completion was replaced by inverse-optimal completion of the fitted rank-one direction. Those new choices are fed into all later templates. Physical leaves are unchanged. The terminal calculation uses the reduced fitted-target-plus-sibling span. Thus the experiment is a modified clean recurrence, not a benchmark for the unchanged HSS union on arbitrary matrices. The reference calculation uses inverse-optimal completion around its fixed rank-one target, making it a probe-adaptive reference rather than the previous deterministic-reference hierarchy.

At $r=2,q=16$, the inherited leaf derivative budget in this clean-family setup is $1/20$ (Chapter~\ref{ch:tangent}). Under the reference-interface hypotheses just stated, the corresponding reference error certificate is $(455+4\sqrt3)/3060<0.151$. This is not an improved constant over the unchanged canonical learner's $75/572$; it covers a different, probe-adaptive reference rule. No finite-noise, depth-uniform bound for the modified recurrence is inferred from it.

\section{Checks, mathematical limitations, and the next estimate}\label{core:sec:checks}
The accompanying scripts use NumPy only for the new model and optimizer calculations. All seeds and configuration counts are included.

\paragraph{Additive completion and general innovation checks.}
Six hundred positive-definite completion problems cover $r=2,4,6$ and every $s=0,\ldots,r$. The largest normalized optimizer-formula residual is below $7.7\times10^{-15}$, mandatory containment error below $1.3\times10^{-15}$, and normalized spectral-bound violation below $1.1\times10^{-15}$. Some optimal completions have an inverse ratio exceeding twenty relative to the smaller child inverse; they still satisfy the additive inequality. These tests distinguish the additive condition from the old relative guard.

There are 4,830 tree realizations in seven configurations. They include fixed-core inverse-optimal completion, changing core ranks, weakest-Gram choices, and source-dependent choices. Innovation second moments agree with their predictions within 2.56 reported standard errors. In the one-source scalar tests, the exact fourth moments agree within 1.21 reported standard errors. The largest weighted-cost difference is 2.92 reported standard errors. The largest register identity residual is below $1.6\times10^{-15}$ and the physical-energy identity residual below $1.6\times10^{-14}$.

Representative fixed-core energies with $\beta=1$ are
\begin{center}
\begin{tabular}{rrrrr}
\toprule $r$ & Height & Realizations & Mean physical plus retained & Standard error\\\midrule
2&4&1200&1.697836&0.007832\\
2&7&350&1.878215&0.008939\\
4&4&650&1.671312&0.003884\\
4&6&180&1.748766&0.007543\\\bottomrule
\end{tabular}
\end{center}
The last row permits decreasing core ranks; the first three have fixed half-rank cores. Response counts and source distributions are recorded in the JSON output. These are independently sampled trees; nodes within a tree are not treated as independent observations.

\paragraph{Modified clean recurrence.}
There are twelve fixed-seed finite differences and 1,060 positive-noise realizations. At $\delta=10^{-3}$:
\begin{center}
\begin{tabular}{rrrr}
\toprule Clean height & Realizations & Mean cut quotient & Mean reference coefficient\\\midrule
0&500&0.0675151&0.0675100\\
2&350&0.0794922&0.0794910\\
4&150&0.0829242&0.0829237\\
6&60&0.0881378&0.0881405\\\bottomrule
\end{tabular}
\end{center}
The cut-quotient standard errors are approximately $0.00366,0.00283,0.00276,0.00350$. These values are not final refitted HSS errors. Finite-noise differences, uniform integrability for this modified rule, and a height-uniform remainder are not certified by this table.

\paragraph{What is now paid.}
The fourth moment of the actual local-model innovation decays under arbitrary admissible nonlinear selections. The inverse-optimal completion automatically supplies an additive budget at every mandatory rank. With one verified marginal core $L^2$ bound, the dependent physical cost is summable in depth. For fixed physical cores that marginal law is exactly a Gaussian projection, yielding the unconditional constants above. No additional depth factor appears in these stability constants.

\paragraph{What remains.}
On noisy HSS inputs, the selected mandatory subspace is itself learned from the probes, so its Gram cannot be replaced by the fixed-core Gaussian law. In addition, shared exterior rows may couple two subtrees and the two complete hierarchies can have different templates. The new local theorem does not remove those dependencies. Establishing a sufficiently controlled marginal core law and the general-cut source interface is still needed before applying the $64k$ query accounting to an all-input risk theorem. If a future verified source decomposition costs $O(L\OPT_k(A))$, the depth-independent stability result would not add another factor $L$; no such general decomposition for the noisy union is asserted here.

\endgroup

\part{Synthesis and the remaining end-to-end implication}

% ===== SYNTHESIS =====
\begingroup
\chapter{The complete dependency chain and its remaining premise}\label{ch:synthesis}
\begin{scopebox}
\textbf{Update, 20 September 2026.} This chapter synthesizes the original separately trained union route. Its conditional composition theorem and $64k$ query count remain correct for that specified algorithm and exact offline projection. They are historical rather than the current algorithm declaration. Chapter~\ref{upd:chap:current-route} defines the synchronized candidate, whose width-$32k$ training streams cost $64k$ products and whose independent rank-$4k$ refit costs $32k$, for a conditional total of $96k$. The current open mathematical interface is the joint physical source/return/invisible estimate for its actual reached hierarchy. Efficient same-rank reduction on the original target tree is a separate unresolved computational requirement. The older selected-core and general-cut interfaces below describe the original proposed route; they must not be presented as an exhaustive statement of what remains after the later developments.
\end{scopebox}

\section{The established composition rule}
Let $\mathcal H_k$ be the original fixed-tree HSS class and let a training procedure return nested row and column spaces of rank at most $R=ck$. Suppose its coefficient space satisfies
\begin{equation}\label{eq:learner-premise}
 \E_{\rm train}\operatorname{dist}_F(A,\mathcal S(A))^2
 \le C_bL\OPT_k(A).
\end{equation}
This is a premise about actual training bases, not a consequence of the refit theorem.

Condition on the complete training transcript. By Chapter~\ref{ch:uniform}, an independent rank-$R$ refit gives a known $C\in\mathcal S(A)$ with
\[
 \E_{\rm refit}[\|A-C\|_F^2\mid\mathcal S(A)]
 \le C_r\operatorname{dist}_F(A,\mathcal S(A))^2.
\]
Averaging the training transcript and using \eqref{eq:learner-premise} gives
\[
 \E\|A-C\|_F^2\le C_rC_bL\OPT_k(A).
\]
An exact offline best approximation $\widehat B$ to the known $C$ in $\mathcal H_k$ exists. The reduction of Chapter~\ref{ch:bicriteria} yields
\begin{equation}\label{eq:final}
 \boxed{\E\|A-\widehat B\|_F^2
 \le(2\sqrt{C_rC_b}+1)^2L\OPT_k(A).}
\end{equation}
All computation in the last step uses known data; polynomial-time nonlinear optimization is not asserted.

For the maintained shared-query union, training uses two width-$16k$ Gaussian streams, hence $32k$ products. The union rank is at most $4k$; the original-width independent refit costs $8(4k)=32k$ products and has $C_r=1607$. Thus the conditional count is $64k$ products. The larger-width and relative-error refit alternatives are available with their separately stated counts. Nothing in this arithmetic proves \eqref{eq:learner-premise}.

\section{What the completed refit proof uses}
The refit proof begins with exact coefficient and observation updates and an orthogonal multiresolution decomposition. The source decomposition retains exterior retained, exterior discarded, and strict interior discarded components. The complete terminal-source product includes all earlier emitted outputs and the root, rather than replacing terminal weights by their values at emission.

The delayed predictable interaction is decomposed after freezing the appropriate independent stream. The coherent mean is not discarded by centering its unweighted coefficient. A corrected retained coefficient cancels the current forcing, and two passive energy telescopes control its complete sequence. This removes the logarithmic allowance in the preliminary embedding certificate. Opposite-stream source orthogonality and centered free-output estimates then add physical budgets without another depth factor. Chapters~\ref{ch:feedback}--\ref{ch:uniform} contain this closed chain.

The independence of the nested bases from the refit probes is present at the beginning and remains necessary at the final conditioning step. The refit theorem therefore does not solve basis learning by being applied to bases selected from those same probes.

\section{What the learning obstructions establish}
The general approximation target has no impossibility theorem in this manuscript. The negative results have narrower and different conclusions.

The arbitrary-template erasure and rank-novelty constructions rule out uniform nullification and pathwise novelty accounting principles. The tree nuclear example rules out the specified depth-weighted regularizers. Actual commutator inverse-collapse examples rule out uniform unweighted inverse estimates, even on cases with zero approximation loss. The anchor-trapping construction is stronger: it disproves the actual annihilate-first learned-space bound. Its final nonlinear postprocessing scope is kept explicit in Chapter~\ref{ch:anchor}.

The union changes the algorithm. Its two-view rescue is an exact positive-noise cut theorem, not merely a comparison of favorable simulations. But neither the union's containment of two coefficient spaces nor its recovery of that one cut establishes its all-input best-HSS error.

\section{The source-coupled regression chain}
In the local-equivariant regression model, the innovation and retained coefficient satisfy
\[
 w_I=N_{G_I}(w_a+w_b),\qquad
 c_I=T_I^T(d_I+z_I),\qquad
 B_I=S_I^T(d_I+z_I).
\]
The innovation has exact second-moment decay and a geometrically decaying fourth moment. Physical discarded sectors are orthogonal, so the retained coefficient and the generated physical costs telescope rather than producing a product of local error factors.

Inverse-optimal completion preserves the mandatory core and obeys
\[
 V_I\le\min(V_a,V_b)+\phi_I,
 \qquad
 \phi_I=\operatorname{tr}\bigl((Q_I^TG_IG_I^TQ_I)^{-1}\bigr).
\]
The fourth-moment innovation estimate converts a marginal second moment of $\phi_I$ into a bound on its dependent source-weighted contribution. For fixed nested physical cores, the core Gram is exactly a fixed Gaussian projection despite adaptive padding. This verifies the required moment and produces the depth-independent constants in Chapter~\ref{ch:core}.

This is a completed propagation theorem under its stated source and core hypotheses. It does not supply the unknown cores, replace selected noisy cores by comparator spaces, or declare the full learned design Gaussian.

\section{The remaining mathematical interfaces}
The all-input route still requires a law or bound for the mandatory signal actually selected from noisy observations. A sufficient local-model input is a controlled marginal second moment of
\[
 \operatorname{tr}\bigl((Q_I^TG_IG_I^TQ_I)^{-1}\bigr),
\]
with the correct rank and width scaling and the required conditioning on independent side information. A selected full-rank truncated sample may leave no optional completion directions, so the padding optimizer alone cannot enforce this law.

A second interface is a general-cut source decomposition. Shared exterior probe rows can couple sibling subtrees; the two separately learned hierarchies can have different templates. Neither subtree independence nor the common-template two-view theorem can be assumed on those trajectories. A valid decomposition costing $O(L\OPT_k(A))$, together with the depth-independent propagation theorem, would avoid a second depth factor. Such a decomposition is not established here.

For the clean trapping family, expected small-noise limits are proved at every fixed height. A height-uniform nonlinear remainder and reversal of the depth/noise limit order are not consequences of those statements. Exact-rank detection and intersection also do not provide finite-precision thresholds or perturbation stability.

\begin{statusbox}
\textbf{Historical endpoint of the original route.} Uniform independent-basis refitting and fixed-physical-core adaptive completion are established; the original annihilate-first basis bound is false. The $64k$ implication above concerns the separately trained union and remains conditional. The subsequent clean-parent obstruction and the completed synchronized candidate are presented in Part~VII. Its current $96k$ ledger still requires a joint physical learning bound and a computationally justified rank-$k$ output reduction; see Chapters~\ref{upd:chap:current-route}--\ref{upd:chap:route-evidence}.
\end{statusbox}

\endgroup


% ===== UPDATE 20 SEPTEMBER 2026: FROZEN FOLLOW-UP RECORD =====
\part{Subsequent obstructions and the current synchronized route}
\label{upd:part:current}
\chapter{Obstructions and limits established after the original manuscript}
\label{upd:chap:later-obstructions}

This chapter records what the subsequent investigations ruled out, what
they repaired, and what remained unresolved when research stopped.  Its
scope is deliberately narrower than an impossibility theorem for HSS
approximation.  A counterexample to a selector, an adaptive moment
estimate, or one prescribed hierarchy is not a lower bound against every
algorithm using a depth-independent number of queries.  The general
conjecture---an $O(k)$-query algorithm with expected squared Frobenius
error $O(L\operatorname{OPT}_k(A))$ and final HSS rank at most $k$---was
not proved or disproved.

All matrices in the probabilistic counterexamples below are fixed before
the Gaussian probes.  Some competing interpolants, conditional states,
or event windows are subsequently selected from the observations; those
quantifiers are stated explicitly.  The original manuscript's adaptive
nullification and annihilate-first counterexamples are recalled only
where needed to distinguish later learners.  The detailed follow-up
catalogue is based on the source notes and their recorded audits, not
on the numerical summaries alone.  \sourcefile{hss_audit_report.md}
provides the chronological overview.

\section{Global interpolation and minimum-norm selectors}
\label{upd:sec:later-global-selectors}

\subsection{A fixed rank-one input with smaller rank-one HSS interpolants}

The global-selector extension removes the qualification that a
regularizer obstruction requires nominal rank at least two.  Let the
two root children $I,J$ each have size $n$, let $q<n$, and draw
independent standard Gaussian probes $\Omega,\Psi\in\mathbb R^{2n\times q}$.
For fixed unit vectors $u,v\in\mathbb R^n$, set
\[
 B=\begin{bmatrix}0&uv^T\\0&0\end{bmatrix}.
\]
This fixed input is HSS rank one on the complete coordinate tree and has
unit Frobenius norm.  Almost surely the systems
\[
 \Psi_J^Ta=\Psi_I^Tu,\qquad
 \Omega_I^Tb=\Omega_J^Tv
\]
have solutions; choose their minimum-Euclidean-norm solutions.  For every
real $t$, the observed-data-dependent competitor
\[
 C(t)=\begin{bmatrix}
 (1-t)ub^T&tuv^T\\
 -(1-t)ab^T&(1-t)av^T
 \end{bmatrix}
\]
satisfies
\[
 C(t)\Omega=B\Omega,\qquad C(t)^T\Psi=B^T\Psi.
\]
It remains HSS rank one: at any node contained in $I$, outgoing and
incoming exterior blocks have the respective common factors $u_S$ and
$b_S^T$; at a node contained in $J$, the corresponding factors are $a_S$
and $v_S^T$.  Thus the assertion concerns every cut, not merely the root
split.

Writing
\[
 \eta=\|a\|^2+\|b\|^2+\|a\|^2\|b\|^2,
\]
the disjoint supports of the four blocks give
\[
 \|C(t)\|_F^2=t^2+(1-t)^2\eta.
\]
Its minimum occurs at $t_*={\eta}/({1+\eta})$, and
\[
 \|C(t_*)\|_F^2=\frac{\eta}{1+\eta}<1,
 \qquad
 \|B-C(t_*)\|_F^2=\frac1{1+\eta}.
\]
Consequently minimum-Frobenius-norm HSS interpolation, including minimum
norm as a tie rule for exact empirical least squares, fails almost
surely on this fixed exact input.  A universal relative guarantee would
require zero error because $\operatorname{OPT}_k(B)=0$ for every $k\ge1$.

The exact displayed error belongs to the constructed competitor, not
to an arbitrary optimizer.  The latter nevertheless has a quantitative
risk lower bound.  If $n>q+1$, put
\[
 \rho=\frac{q}{n-q-1},\qquad \mu=2\rho+\rho^2.
\]
Inverse-Gram averaging gives $\mathbb E\|a\|^2=
\mathbb E\|b\|^2=\rho$; the two vectors use independent probe streams,
so $\mathbb E\eta=\mu$.  Concavity of $x/(1+x)$ and the reverse triangle
inequality in the space of square-integrable matrices imply, for every
measurable minimum-norm selection $\widehat C$,
\[
 \mathbb E\|B-\widehat C\|_F^2
 \ge\left(1-\sqrt{\frac{\mu}{1+\mu}}\right)^2.
\]
This tends to one as $n/q\to\infty$.  Existence of the minimizer follows
from closed rank constraints and compact norm sublevel sets.  These
statements are established in \sourcefile{hss_global_selector_obstruction.tex}.

\subsection{Tree penalties and exact leaf-block calibration}

For constant $u,v$ on balanced root halves, the outgoing-plus-incoming
cut nuclear norm of $B$ at a depth with $m$ nodes per half is $2\sqrt m$.
All cuts of $C(t)$ have rank at most one.  The partition bound
\[
 \sum_S\bigl(\|C[S,S^c]\|_F^2+
                    \|C[S^c,S]\|_F^2\bigr)\le2\|C\|_F^2
\]
and Cauchy--Schwarz bound its nuclear sum by
$2\sqrt{2m}\|C\|_F$.  On $\eta<1$, the smaller interpolant has
$\|C(t_*)\|_F<1/\sqrt2$ and improves every depth strictly.  It therefore
improves every nonzero nonnegative deterministic depth-weighted sum of
these cut nuclear norms.  Its ordinary matrix rank is at most two, so
its ordinary nuclear norm is also smaller than that of $B$ on this
event.  Markov's inequality gives probability at least
$\max(0,1-\mu)$ for the event.  The same strict comparison defeats
positive penalization because both empirical residuals vanish.

Supplying all physical-leaf diagonal blocks exactly does not repair
these selectors.  Mark one leaf in each sibling pair active and the
other inactive.  Support $u,v$ on the active leaves of their respective
halves, and restrict $b,a$ to inactive leaves.  The free dimension in
each interpolation system is then $d=n/2$; assume $d\ge q$.  On every
physical leaf $K$, at least one factor of the corresponding diagonal
outer product vanishes, so
\[
 C(t)[K,K]=B[K,K]=0
\]
exactly.  The smaller-norm construction is otherwise unchanged.

For $d>q+1$, replace $\rho,\mu$ by
$\rho_d=q/(d-q-1)$ and $\mu_d=2\rho_d+\rho_d^2$.  The minimum-norm
lower bound remains valid.  With $u,v$ spread uniformly over active
coordinates, the true leaf-depth nuclear sum is $2\sqrt{m/2}$, while
the competitor still obeys $2\sqrt{2m}\|C\|_F$.  The stronger event
$\eta<1/3$, of probability at least $\max(0,1-3\mu_d)$, gives
$\|C(t_*)\|_F<1/2$ and improves every depth, including the leaves.
Any finite penalty depending only on the supplied leaf diagonal blocks
takes the same value on both candidates and cannot reverse this
comparison.  A single-active-leaf version also gives calibrated
minimum-norm failure, with free dimension $n-s$ for leaf size $s$.
See \sourcefile{selector_with_leaf_data.md} and the final theorem in
\sourcefile{hss_global_selector_obstruction.tex}.

The alternatives above depend on the probes.  They disprove specified
global selection principles; they do not say that two fixed inputs are
indistinguishable with positive probability, or that a different
structured decoder cannot recover each fixed HSS input exactly.

\subsection{Certificates restore exact uniqueness, not robust selection}

The trajectory-certified selector imposes injectivity of each complete
physical template together with its exterior cut anchor.  Under these
certificates, equal transcripts determine equal outgoing and incoming
intersections, hence equal canonically completed child frames.  Induction
then gives a common coefficient space; successive two-sided compression,
ending with an injective root core, proves exact uniqueness.  The false
interpolants above violate these full-template certificates, including
the endpoint where a minimal cut anchor alone would not detect them.

For a fixed exact input the comparison templates are deterministic, so
each $q\times p_I$ certificate matrix is Gaussian and
$\mathbb E\sigma_{\min}^{-2}\le p_I/(q-p_I-1)$.  This marginal
bound cannot be replaced by a tree-size-independent bound on the
maximum certificate cost: disjoint leaves provide arbitrarily many
independent Gaussian direction images.  Hard thresholds also exclude
some exact inputs with positive probability.  The quantitative
compression inequality gives products of levelwise factors, rather
than the desired robust residual-relative bound for changing candidate
spaces.  The exact measurable selector is therefore a valid positive
endpoint, not a completed approximation algorithm.
See \sourcefile{certified_selector_route.md} and
\sourcefile{certified_selector_audit.md}.

\section{The one-sided Gaussian pilot obstruction}

Removing transpose data from a pilot would make its row spaces
independent of a later transpose stream, but exact bounded-rank row
containment already fails for that restricted acquisition model.
Fix query width $q$ and output rank $r$.  Take powers of two
$d\ge2r$ and $M\ge2q+2$.  Partition one root half $I$ into $d$ aligned
subtrees $K_i$ of size $M$, choose fixed unit coordinates $e_i\in K_i$,
and choose a fixed unit coordinate $v$ in the other half.  For arbitrary
$t_i$ and $b_i$ supported in $K_i$, define
\[
 u=\sum_{i=1}^d t_i e_i,
 \qquad A=\sum_{i=1}^d e_i b_i^T+uv^T.
\]
Every such matrix is HSS rank one, including at cuts inside the $K_i$.
The all-cut verification uses the single row factor $e_i$ inside a
block and the single exterior factor $v$ above unions of complete
blocks; no large unrestricted leaf is hidden in the construction.

A pilot observing only $(\Omega,A\Omega)$ obtains, apart from zero
rows, exactly
\[
 y_i=X_i^Tb_i+t_i\gamma,
 \quad X_i=\Omega_{K_i},\quad \gamma=\Omega_J^Tv.
\]
To obtain a fixed-input impossibility, use an input prior only in the
proof: independently take $t_i\sim N(0,1)$ and
$b_i\sim N(0,\tau^2 I_M)$ before drawing $\Omega$.  Conditional on
the complete observations, the posterior covariance of $t$ is diagonal
with strictly positive entries
\[
 w_i=\left[1+\gamma^T(\tau^2X_i^TX_i)^{-1}\gamma\right]^{-1}.
\]
Thus $u$ has a posterior density on the entire fixed $d$-dimensional
span of the $e_i$.  After also fixing any independent algorithmic seed,
a returned space of dimension at most $r<d$ contains $u$ with posterior
probability zero.  Fubini converts this statement into failure for
almost every fixed input under the proof prior, with the bad-input set
allowed to depend on the pilot.

There is also a quantitative Bayes loss bound.  A rank-$r$ projector
captures at most $r$ from a covariance bounded by the identity, giving
\[
 \mathbb E\|(I-P_S)u\|^2
 \ge\frac{d}{1+q/[\tau^2(M-q-1)]}-r.
\]
For $\tau=2$ and the displayed dimensions this is at least $3d/10$.
The result permits unlimited computation and disclosure of the entire
parametric family.  It does not cover structured forward queries: on
this very family the single query $Av=u$ reveals the desired vector.
Transpose products and independently trained two-sided pilots are also
outside its scope.  The theorem and independent review are
\sourcefile{one_sided_gaussian_pilot_obstruction.md} and
\sourcefile{one_sided_gaussian_pilot_audit.md}.

\section{Why the screened-source moment proposed by Claude is false}

\subsection{The valid deterministic cancellation}

For a full-row-rank design $G=U\Sigma V^T$, lift
$L=(G^\dagger)^T$, response $J$, and retained projector $P$ with
$PLJ=LJ$, split the singular directions into the weakest $d$ and the
remaining directions.  If $J_1=V_1^TJ$ has full row rank, then
\[
 (I-P)L=(I-P)U_2\Sigma_2^{-1}
       \bigl[V_2^T-J_2J_1^\dagger V_1^T\bigr].
\]
This algebraic identity is correct.  The failed proposal was to discard
the remaining left projector and assert a uniformly controlled second
moment for the resulting screened source.  Weak singular values must
indeed be the removed block; sorting in decreasing order would not even
give the claimed deterministic inverse bound.  Full rank of $G$ also
does not imply the required response rank of $J_1$.

\subsection{An infinite moment at an ordinary Gaussian leaf}

The counterexample uses the actual full-range learner's leaf dimensions
$m=16,q=64,s=4,r=8$.  Fix orthonormal physical vectors $a,e$ inside a
leaf and a unit $v$ outside it, and take
\[
 A=av^T,\quad B=(a-e)v^T,\quad E=ev^T.
\]
The comparator is legal but need not be best.  Here $G$ is an ordinary
Gaussian leaf design and, for an independent Gaussian response row
$\gamma=v^T\Omega$,
$J=G^Ta\gamma$ and $J(E)=G^Te\gamma$.
Sort the singular values increasingly, and write
$c=U^Ta,d=U^Te$.  Screening only the weakest coordinate gives
\[
 Q(E)_i=\sigma_i(d_i-c_i d_1/c_1)\gamma,
 \qquad i\ge2.
\]
Consequently
\[
 \sigma_2^{-2}\|Q(E)\|_F^2
 \ge\|\gamma\|^2\left(1+\frac{d_1^2}{c_1^2}\right).
\]
Haar orientation gives a density for $c_1$ positive at zero and
$\mathbb E[d_1^2\mid c_1]=(1-c_1^2)/(m-1)$.  The expectation is
infinite, even conditional on any nonzero $\gamma$.  All stated rank
assumptions hold almost surely.  In contrast, the actual fit is
$a\gamma$, so the retained space contains $a$ and
\[
 \|(I-P)LJ(E)\|_F^2\le\|\gamma\|^2.
\]
Its expectation is at most four.  The discarded left projection can
cancel a nonintegrable response ratio.

\subsection{Linear scaling even for an exact best comparator}

Fix orthonormal $a,e$ inside the leaf and orthonormal $v,w$ outside it.
For $0<\varepsilon<1$, let
\[
 A=av^T+\varepsilon ew^T,\quad B=av^T.
\]
The outgoing cut has singular values $1,\varepsilon$, and $B$ attains
the cut lower bound globally.  Thus
$\operatorname{OPT}_1(A)=\varepsilon^2$.  Put
\[
 \gamma=v^T\Omega,\quad \eta=w^T\Omega,\quad
 g=\|\gamma\|^2,\quad \kappa=\|\eta\|^2,\quad
 h=\gamma\eta^T,\quad \Delta=g\kappa-h^2.
\]
The same one-coordinate screening has a lower-bound integrand whose
Haar average is exactly
\[
 \frac{\varepsilon g\kappa}{2\sqrt\Delta}.
\]
One obtains this by setting $t=c_1/d_1$, using its standard Cauchy
density, and integrating the resulting rational function.  For
independent Gaussian response rows in $\mathbb R^4$, the radial and
angular moments give
\[
 \mathbb E[\sigma_2^{-2}\|Q(E)\|_F^2]
 \ge\frac94\varepsilon.
\]
Appending the natural comparator separator $T=\gamma^\dagger$ still
gives
\[
 \mathbb E[\sigma_2^{-2}\|Q(E)T\|_F^2]
 \ge\frac14\varepsilon.
\]
Neither expression is uniformly $O(\varepsilon^2)$.  Yet the full fit
$a\gamma+\varepsilon e\eta$ has range $\operatorname{span}(a,e)$
almost surely, so the actual filtered source, with or without $T$, is
identically zero.  This defeats the proposed proof lemma more strongly
than a badly chosen comparator would, while leaving the leaf learner
successful.  See \sourcefile{claude_screened_residual_audit.md} and
\sourcefile{screened_source_counterexample_audit.md} for the exact
integrals, corrected dimensions, perturbation expansion, and review.

The same audit also rejects two unrelated shortcuts: learned templates
need not be measurable in local probe rows alone, and a per-node bound
by the full residual norm cannot be summed over all tree nodes as though
there were only $L$ of them.  The actual comparator source includes the
signed inherited return as well as the current residual commutator.

Several corrections are needed even before those counterexamples.
The internal predesign is $16k\times64k$, not the retained
$8k\times64k$ design.  Weak directions must be indexed by the smallest
singular values, and $d\le s$ does not imply response rank $d$.
A random inverse singular scale remains inside its joint expectation;
conditioning on a probability-one full-rank event does not make it
constant.  Coefficients formed from $J(B)$ may themselves depend on
exterior $\Psi$ rows.  Finally, with
$B_i=V_i^TJ(B)$, $E_i=V_i^TJ(E)$,
$H=(B_1+E_1)^\dagger$, and $H_B=B_1^\dagger$, the exact expansion is
\[
 Q(E)=E_2-B_2H_BE_1+B_2(H_B-H)E_1-E_2HE_1.
\]
The last two terms cannot be omitted, and the adaptive $V_i$ need not
be comparator-only objects.  These are additional algebraic and
conditioning errors in the relayed proposal, distinct from its false
screened-moment claim.

\section{The old full-range learner: inverse growth and actual failure}
\label{upd:sec:later-fullrange}

Two different obstructions were established for the full-range
commutator candidate with wide-query width $q=64k$, response width
$4k$, retained rank $r=8k$, and physical leaves of size $16k$.  This
candidate retains the entire fitted response and chooses the remaining
directions by inverse-optimal completion.  It is not the maintained
annihilated learner or the later synchronized union learner.  Its
all-depth exact-HSS rank and capture theorem remains valid; the noisy
failure below does not contradict exact recovery.

\subsection{An actual full inverse diverges while approximation is exact}

For $k=1$, let $a,b$ be siblings inside a nonroot node $p$, and take the
fixed one-way rank-one input $A=uv^T$, with $u$ supported in $a$ and $v$
in $b$.  Put $h=\Psi_a^Tu$.  On the $b$ branch every nonzero response
has left factor $h$, despite every outgoing physical block being zero.
For a retained query space $S_j$, full-fit retention implies
\[
 P_{S_j}h=P_{\operatorname{row}(G_j)}h.
\]
Write $\delta_j=\|(I-P_{S_j})\widehat h\|^2$.  Conditional on $h$,
the two phantom children use independent physical Gaussian rows and
have axial symmetry about $\widehat h$.  Their mean residuals equal
$d_\ell\widehat h$, where $d_\ell=\mathbb E\delta_j$.  Comparing the
parent projection with the average of the two child projections gives
\[
 d_{\ell+1}\le\frac{d_\ell+d_\ell^2}{2},\qquad
 d_0=\frac34,\qquad d_\ell\le12\,2^{-\ell}.
\]
The initial value concerns the full sixteen-dimensional Gaussian
predesign: the adaptive eight-dimensional output preserves its fit.

The signal branch represents $u$ exactly.  If $D_a^Tc_a=h$ and
$D_b^Tc_b=P_{S_b}h$, then $\|c_a\|=1$ and
\[
 \sigma_{\min}\!\begin{pmatrix}D_a\\D_b\end{pmatrix}^{\!2}
 \le\|h\|^2\delta_b.
\]
For a range-only tie and canonical-frame rule, scaling $h$ does not
change the selected spaces.  Its Gaussian radius is therefore
independent of $\delta_b$, and
\[
 \mathbb E\|G_p^\dagger\|_{\mathrm{op}}^2
 \ge\frac{1}{(q-2)d_\ell}
 \ge\frac{2^\ell}{12(q-2)}.
 \tag{\thechapter.1}
\]
All the designs in this particular finite-depth family are nonsingular:
the child spaces projected onto $h^\perp$ have independent Haar
orientations and dimensions whose sums are below $q-1$.  This argument
does not declare the selected Grams Wishart.

At $p$, however, both the outgoing block and response are zero.
Inverse-optimal completion may discard the almost-canceling direction;
its retained inverse trace satisfies $V_p\le\min(V_a,V_b)$.  Every
true cut range is captured, and the fixed learned coefficient space
contains $A$ exactly.  Thus (\thechapter.1) refutes a uniform bound for
the \emph{full predesign inverse}, not approximation stability.
The proof, exact scope of the range-only rule, and independent review
are in \sourcefile{phantom_stability.md} and
\sourcefile{phantom_stability_audit.md}.

An associated scalar shortcut also fails.  For a normalized shared
target, define $R=\beta/\|f\|^2$, where $\beta$ is the query residual
and $f$ the minimum-norm physical fit.  At a Gaussian leaf,
\[
 R\ \stackrel{d}{=}
 \frac{X_{q-m}X_{q-m+1}}{X_m},\qquad \mathbb ER=168
 \quad(m=16,q=64),
\]
with independent chi-square variables.  The actual parent law also
depends on child metric matrices and their relative orientations.
Reachable leaf states with an increasingly strong fitted direction and
large optional directions give a limiting conditional mean exceeding
$6/7$ times the child ratio.  Hence the proposed uniform conditional
contraction by $6/7$ is false; this does not rule out the weaker
unconditional estimates subsequently used below.
See \sourcefile{phantom_ratio_route.md} and
\sourcefile{phantom_ratio_audit.md}.

\subsection{A fixed-input lower bound for the entire retained projector}

The later construction controls all seven optional directions, not
only the mandatory fitted line.  For small $\varepsilon>0$, set
\[
 t=\left\lceil\frac72\log(1/\varepsilon)\right\rceil,
 \quad\ell=t+1,\quad N=128\,2^t,\quad L=t+3.
\]
Two physical leaves of size sixteen form a source $a_0$.  Extend it
through $t$ active/blank merges to $a$; let $b$ be its sibling and
$p=a\cup b$ a root child.  With $u=e_1$, $v_b$ a fixed unit vector
uniform on $b$, and $v_{\mathrm{ext}}$ an exterior coordinate, define
\[
 B=\frac{u(v_b+v_{\mathrm{ext}})^T}{\sqrt2},\qquad
 E_0[a_1,a_2]=E_0[a_2,a_1]=I_{16},\qquad
 A_\varepsilon=B+\varepsilon E_0.
\]
All remaining entries of $E_0$ vanish.  The matrix is fixed before
sampling, $B$ is globally rank one, and
$\operatorname{OPT}_1(A_\varepsilon)\le32\varepsilon^2$.
For the full actual retained rank-eight space $U_p$, the reviewed result
is
\[
 \mathbb E\operatorname{dist}(u,U_p)^2
 \ge c\varepsilon^{1.87}.
 \tag{\thechapter.2}
\]
Since $A_\varepsilon[p,p^c]=uv_{\mathrm{ext}}^T/\sqrt2$, every
approximation confined to these learned nested row spaces incurs at
least half this expected squared loss.  Consequently its ratio to
$L\operatorname{OPT}_1(A_\varepsilon)$ diverges.  Refitting or reducing
rank inside the same spaces cannot repair this failure.

The proof has three substantive layers.  First, a fixed compact source
event of probability bounded below gives an initial query gap
$\alpha_0\asymp\varepsilon^2$ and a physical fit within
$O(\varepsilon)$ of the represented signal.  The active/blank physical
martingale has $O(\varepsilon^2)$ squared error uniformly in its length.
Its gap is $\alpha_0$ times $t$ independent
$\operatorname{Beta}(24,4)$ factors.  Since their mean negative
logarithm is less than $4/25$, a high-probability state event satisfies
\[
 c_\alpha\varepsilon^2e^{-(4/25)t}
 \le\alpha\le C_\alpha\varepsilon^2.
\]
For the phantom sibling, the separately proved moment estimates give
on a high-probability event
\[
 \beta\le .439^\ell,\qquad \phi_b\le .699^\ell,\qquad
 V_b\le(9/8)^\ell,\qquad R_b\le(22/25)^\ell.
\]
Here $\phi_b$ is the mandatory inverse cost, $V_b$ the full retained
inverse trace, and $R_b$ the innovation-to-fit ratio.  The fit-decay
bound uses an exactly certified ninth-moment angular estimate, not a
Monte Carlo estimate.  The state event also bounds the total physical
probe Frobenius norm by a constant times $N$.

Second, condition on the complete active state and phantom data
invariant under query rotations fixing $h$.  Eliminate the optional
query coordinates and expose their planes and the remaining radii.
On a further event of conditional probability at least $3/8$, a single
relative angle $c$ remains, with sphere-coordinate density
$g_{49}(c)\propto(1-c^2)^{23}$.  Physical backsubstitution locates a
zero $c_*$ of the \emph{actual physical overlap}; a zero of a projected
query overlap alone would not suffice.  On the interval
\[
 |c-c_*|\le\kappa\frac{\sqrt\alpha}{\sqrt{T_\ell}},
 \qquad T_\ell=C_T(2\cdot .699)^\ell,
\]
the mandatory physical overlap is small and the parent mandatory
inverse cost is at most a constant times $\phi_b$.

Third, the Schur-energy estimate controls every inverse-optimal
padding direction.  Its vanishing terms include
\[
 \varepsilon^2\bigl[(2\cdot .699)(9/8)\bigr]^\ell,
 \qquad [2\cdot .699^2]^\ell,
\]
whose exponents are favorable.  With $\kappa$ fixed sufficiently small,
the full retained space captures at most one quarter of the squared
norm of $u$ throughout the interval.  Its conditional probability is
at least a constant times $\sqrt\alpha/\sqrt{T_\ell}$, hence at least
\[
 c\varepsilon^p,\qquad
 p=1+\frac74\left(\frac4{25}+\log(1.398)\right)
   =1.86632462667\ldots<1.87.
\]
All physical-fidelity and state events are imposed \emph{before} the
final shrinking angle interval.  Subtracting an unrelated
unconditional failure probability from that interval would not prove
the claimed lower bound.

The complete proof and two independent reviews are recorded in
\sourcefile{full_projector_bridge_obstruction.md} and
\sourcefile{full_projector_bridge_audit.md}; its ingredients are
\sourcefile{active_blank_martingale.md},
\sourcefile{phantom_negative_moment_route.md},
\sourcefile{phantom_negative_kernel_bound.md},
\sourcefile{physical_bridge_zero.md}, and
\sourcefile{optional_cost_exclusion.md}.  The theorem concerns exact
full-range selection with the specified completion.  An absolute
singular-value cutoff can change that algorithm at small noise.  The
earlier $\varepsilon^{7/5}$ result concerned only its mandatory line;
it must not be substituted for the full-projector proof.

The maintained annihilated learner survives this particular source
mechanism.  Its phantom sample vanishes and its blank frames are
deterministic.  For the stated leftmost two-leaf source family, with
zero-outgoing siblings and vanishing perturbation commutator above the
source, the reviewed bound is
\[
 \mathbb E\operatorname{dist}_F(A,\mathcal S_\cup)^2
 \le\frac{191}{11}\|E\|_F^2
 \le\frac{764}{33}\operatorname{OPT}_k(A).
\]
These support conditions are essential; this is neither a general
stability theorem nor evidence against (\thechapter.2).
See \sourcefile{maintained_internal_source_bound.md} and
\sourcefile{canonical_blank_path_audit.md}.

\section{Separate views can lose a clean parent signal}
\label{upd:sec:later-separateviews}

Post-training union of independently propagated views does not ensure
that a signal contained in the child union survives the parent union.
The initial oblique example established this only for prescribed local
states.  The later construction reaches it through the actual
maintained and raw two-view rules, including physical leaves, shared
probes, column nullifiers, and canonical completion.

Take $N=32$, $k=1$, retained rank $r=2$, and $q=16$, with physical
leaves of size four.  Let a sixteen-coordinate root child $I$ have
children $a,b$ of size eight.  On the two leaves of $a$ choose
orthonormal coordinate pairs $(u_1,p_1),(u_2,p_2)$ and put
$u=(u_1+u_2)/\sqrt2$.  On $b$ choose coordinates $b_1,\ldots,b_4$;
let $f$ be exterior to $I$.  The fixed input is
\[
 B=uf^T,\qquad E=p_1b_1^T+p_2b_2^T,\qquad
 A_\varepsilon=B+\varepsilon E.
\]
The perturbation is entirely internal to $I$, so $J_I(E)=0$.
For $0<\varepsilon<1$, the singular values of the cut at $a$ are
$1,\varepsilon,\varepsilon$, and
\[
 \operatorname{OPT}_1(A_\varepsilon)=2\varepsilon^2.
\]
Thus the comparator is best, rather than merely a convenient rank-one
matrix.

The explicit full-rank probe witness in
\sourcefile{two_view_union_reachability_route.md} produces child row
planes
\[
 A_a=\operatorname{span}\{(u+p_1)/\sqrt2,p_2\},\quad
 R_a=\operatorname{span}\{(u+p_2)/\sqrt2,p_1\},
\]
and $A_b=\operatorname{span}(b_1,b_2)$,
$R_b=\operatorname{span}(b_3,b_4)$.  Consequently the child union
contains $u$ exactly.  The child sample covariances on
$(u,p_1,p_2)$ can be taken as
\[
 S_{\rm ann}=\begin{pmatrix}2&1&0\\1&2&0\\0&0&3\end{pmatrix},
 \qquad
 S_{\rm raw}=\begin{pmatrix}6&0&2\\0&8&0\\2&0&6\end{pmatrix}.
\]
Their selected projectors have strict retained/discarded gaps.  The
single probe construction realizes both covariances and the required
physical-leaf and column frames simultaneously; they are not prescribed
as independent fictitious responses.

At $I$ both actual row samples have structural rank one.  Their lifted
fits have the forms
$A_1(u+p_1)/\sqrt2+B_1b_1$ and
$A_2(u+p_2)/\sqrt2+B_2b_4$, with all displayed coefficients nonzero.
In each canonical child-template coordinate order, the first $a$
column is therefore independent of the mixed fit.  It is used for
padding before the second $a$ column.  The resulting spaces are
\[
 A_I=\operatorname{span}(u+p_1,b_1),\qquad
 R_I=\operatorname{span}(u+p_2,b_4).
\]
Their union projects $u$ to $(2u+p_1+p_2)/3$, and hence
\[
 \operatorname{dist}(u,A_I+R_I)^2=\frac13.
 \tag{\thechapter.3}
\]
The current residual commutator is zero, the parent outgoing comparator
is rank one, and neither view truncates a nonzero parent sample tail.
The loss is inherited through the separately propagated templates.

This is an actual-recursion reachability statement.  For each fixed
$\varepsilon>0$, full design ranks, spectral gaps, nonzero mixing
coefficients, and the relevant first canonical pivot persist on a
positive-probability Gaussian neighborhood.  Only the needed physical
subspaces and first $a$ pivot must be continuous.  Some unused canonical
columns can change gauge discontinuously, and the proof does not assume
otherwise.  The parent rank-one sample is structural under probe
perturbations, rather than a nongeneric zero singular value imposed by
the witness.

The witness uses narrow-probe coordinates of size $1/\varepsilon$.
Its neighborhood probability may depend on $\varepsilon$; (\thechapter.3)
does not by itself yield an asymptotic expected-error lower bound.
Independent deterministic evaluations of the actual dense learners at
$\varepsilon=1,.1,.01$ reproduce $1/3$ and full rank-four predesigns at
all fifteen nodes.  The algebra and continuity argument, not those
evaluations, prove reachability.  See
\sourcefile{two_view_union_reachability_audit.md} and, for the earlier
unreached local example, \sourcefile{two_view_union_oblique_route.md}.

\paragraph{The synchronized repair and its limits.}
Using the same child-union templates for both parent views removes
this exact mechanism.  The specified synchronized learner retains at
most $4k$ directions with width $32k$, retains the full raw fit on its
low-rank branch, and otherwise combines the lifted raw leading $2k$
directions with the annihilated fit.  Its completion first fills query
visible directions and then physical kernel directions if necessary.
For exact HSS-$k$ and the stated rank-slack inputs, rank and true-cut
containment have proofs.  A clean parent incurs no new loss when its
common child template contains the comparator and its predesign is
injective on the relevant coefficient space; on singular branches an
invisible comparator component must be retained explicitly in the
identity.  Full rank of a retained design is weaker than full rank of
the next stacked predesign.  None of these statements provides the
joint inverse-weighted budget for arbitrary noisy signed sources.
See \sourcefile{synchronized_union_exact_route.md},
\sourcefile{synchronized_union_route.md}, and their audits.

\section{Conditioning, passive memory, and changing targets}
\label{upd:sec:later-conditioning}

\subsection{What the renewed latent law permits one to condition on}

For the independent-source fitted-target tree, the exact innovation
mixture is a theorem for a restricted recursively generated summary.
It may retain physical bases and fits, allowed local and retained
Grams, and deterministic functions of these data.  It may not retain
forgotten query-coordinate designs, old residual vectors or radii, or
arbitrary response cross-Grams.  Selection must be invariant under
right rotations of query coordinates and must contain the entire fit.

For example, after temporarily exposing two current child designs and
their scalar latent variances $\tau_a,\tau_b$, the parent residual
before compression has covariance
\[
 (\tau_a+\tau_b)N_G,
\]
where $N_G$ projects onto the kernel of the full $2r$-row predesign.
The fitted increment and this residual are conditionally independent
Gaussians.  The residual is not isotropic in the larger kernel of the
retained $r$-row design under that same conditioning.  Only after
forgetting the child query frames and using the retained-space
stabilizer can one replace its direction by a uniform direction in
that larger kernel.  The beta--chi-square identity then supplies
\[
 \tau_{\rm new}=(\tau_a+\tau_b)\,Z,\qquad
 Z\sim\operatorname{Beta}\!\left(\frac{q-2r}{2},\frac r2\right)
\]
in the reduced realization.  This is not a license to restore all the
discarded records afterwards.

Two explicit pitfalls illustrate the restriction.  A selector that
inspects a fixed query coordinate can orient an optional direction
toward that coordinate, contradicting the alleged isotropic residual
variance.  Even an invariant record can be excessive: retaining the
old residual radius, the new fit, and its Gram determines the new
radius by subtracting the fitted energy.  A fresh nondegenerate beta
radius is then impossible.  Right-rotation invariance alone is
therefore insufficient.  The corrected theorem and its examples are
in \sourcefile{fitted_target_exact_audit.md}; the matrix version uses
the same restricted-summary obligation in
\sourcefile{matrix_latent_budget.md}.

The source posterior does not renew as a Wishart posterior at every
merge.  With block pre-Gram
$K=\left[\begin{smallmatrix}A&C\\C^T&B\end{smallmatrix}\right]$
and Schur complements $S_a=A-CB^{-1}C^T$,
$S_b=B-C^TA^{-1}C$, the fitted increment covariance is
\[
 \Omega=K^{-1}\operatorname{diag}(\tau_bS_a,\tau_aS_b)K^{-1}.
\]
Its likelihood involves
\[
 \frac{(Kz)_a^TS_a^{-1}(Kz)_a}{\tau_b}
 +\frac{(Kz)_b^TS_b^{-1}(Kz)_b}{\tau_a}.
\]
Already on a scalar cross-block slice this is proportional to
$(\tau_b^{-1}+c^2\tau_a^{-1})/(1-c^2)$, not a linear precision term
in $K$.  Source conjugacy thus does not establish learned-Gram
conjugacy.  Likewise an inverse-weighted residual cost contains the
actual product $\mathbb E[(\tau_a+\tau_b)V_p]$; replacing it by a
product of expectations drops its covariance term.
See \sourcefile{binary_posterior_closure.md}.

These obstacles did not prevent every later local theorem.  The finite
moment ladder bypasses conjugacy, retaining the dependent products and
conditional Haar orientations.  With its stated independent-source
assumptions, $r\ge4096$, $q=8r$, and target width at most $r/2$, it
establishes uniform directional error moments and inverse means with
large absolute constants.  It is not a theorem for $r=2$, nor an
automatic interface to actual exterior HSS commutators.
See \sourcefile{fitted_energy_bootstrap_route.md} and
\sourcefile{fitted_energy_bootstrap_audit.md}.

\subsection{Projection loss, fitted energy, and passive records differ}

A physical projection error can be bounded by the target norm on a bad
event.  A fitted coefficient can be arbitrarily large on that same
event and affect ancestor selections.  Thus paying only the local
projection loss does not dispose of the future fitted-energy cost.
For the fitted-target martingale $S_t=F_t-C_0$, a valid global stopping
argument with exit threshold $M\|C_0\|$ gives
\[
 \mathbb E\|E_H\|^2\le
 \beta+(1+M^{-2})\,
 \mathbb E\|S_{H\wedge\tau}\|^2,
\]
where the stopped quadratic variation includes the exit increment.
Discarding that overshoot changes the process.  A global fit bound does
not normalize every small local target; conditioning on survival can
also couple siblings.  The reduction in
\sourcefile{fitted_target_stopping.md} was valid but did not close its
remaining bound.  Later large-rank source theorems should not be read
back into that scalar stopping argument.

Merely carrying an old target or error vector as a passive record
does not make it an actively retained fit.  The next compression may
discard it, after which its update has an additional projection loss
and is not the centered full-fit recurrence.  Even retaining each
current joint and separate child fit need not preserve older witnesses.
\sourcefile{augmented_memory_route.md} gives a full-rank two-merge
local construction in which current fits and inverse-optimal padding
eventually exclude a nearly represented signal.  That construction
was checked as an admissible local state, not reached from the actual
Gaussian leaf process; it is not an expected-error counterexample.
The joint-plus-split learner nevertheless has a valid exact-input
endpoint at $r=16k$, $q=64k$, response width $4k$, and physical leaves
of size $32k$.  Its three mandatory fits have combined dimension at
most $12k$, and the exact-input boundary-anchor induction proves rank
and cut containment.  A clean rank-one bridge retains a good separate
fit from either child.  That one-bridge protection does not automatically
protect the same vector at the next ancestor.
See \sourcefile{split_fit_learner_audit.md}.

Nor does physical closeness to a fixed envelope suffice.  If
$Q=DC+E$, $D^TE=0$, and $\|E\|_{\rm op}\le\delta<1$, then
\[
 \sigma_{\min}(Q^TX)\ge
 \sqrt{1-\delta^2}\,\sigma_{\min}(D^TX)
   -\delta\|(I-DD^T)X\|_{\rm op}.
\]
This gives good-event inverse moments under a \emph{relative
probe-space} perturbation condition.  Without it, for $n>q$ an adaptive
unit vector arbitrarily close to a fixed $e$ can approach
$\ker X^T$ on a positive-probability event while remaining nonsingular.
Its inverse cost grows like an arbitrary inverse squared parameter.
An accompanying fixed-target example has zero query innovation but
arbitrarily large fitted coefficient.  It can even converge eventually
to the exact core for every realization while its expectations diverge.
These are right-rotation-invariant adaptive constructions, not reached
nested-tree examples.  They refute a bridge based only on geometric
closeness and explain the need for uniform integrability.
See \sourcefile{near_envelope_route.md}.

A related historical capture obstruction is also limited in scope.
The marginal facts $D=T-B$, $T/\sigma^2\sim\chi^2_{15}$, small
$L^2$ residual $B$, and $D\ge_{\rm st}\chi^2_2/2$ alone permit abstract
couplings with an infinite first inverse moment.  Such couplings need
not satisfy the actual fresh latent-radius law.  Once that law is kept,
the later Laplace argument gives
$\mathbb ED^{-p}\le\sigma^{-2p}2^{-p}
\Gamma(3/2-p)/\Gamma(3/2)$ for $p<3/2$, including
$\mathbb ED^{-1}\le\sigma^{-2}$.  Thus the marginal construction is
not an actual-tree counterexample.  Compare
\sourcefile{actual_full_response_capture.md} with
\sourcefile{latent_capture_laplace.md} and its independent audit.

\subsection{Target drops require a retained deterministic guard}

Deterministic compatible transfers
$U_I|_i=U_iA_i$ preserve the additive physical target and transform
omission covariances by congruence,
$C_I=\sum_iA_i^TC_iA_i$.  Invertible changes of output coordinates
preserve the relevant fit ranges.  A strict drop is different: a fixed
physical prefix that would complete the \emph{current} smaller target
may already have been discarded while an earlier larger target was
being fitted.

The explicit two-target source example has a fixed prefix of length
$r-2$ and a random fitted two-plane.  A prescribed next coordinate
almost surely lies outside the retained space.  Dropping one output
at the parent does not recreate that coordinate.  Treating its learned
replacement as a fixed Gaussian direction is invalid.  Without a
history-compatible prefix condition, the inverse decomposition has an
extra nonnegative padding cost
\[
 \mathcal P=\|B_{\rm phys}(I-P_{\operatorname{range}Z})\|_F^2,
\]
which persists when there are more unexcited available directions than
active target columns.  The moment recurrence must then charge
inverse-weighted products involving $\mathcal P$.
See \sourcefile{nested_target_bootstrap_route.md}.

The fixed-guard modification repairs this particular problem.  With
$s_{\max}\le r/2$, put $c=r-s_{\max}$.  Each source has a deterministic
pool of $2c$ directions and retains its first $c$ as a guard.  At a
parent, the block sum of its two child guards is an available
deterministic pool of $2c\ge r$ directions; select the full fit and
complete by scanning this pool.  Since the fit rank is at most
$s_{\max}$, the first $c$ pool directions remain retained.  This is an
implementable induction, not an oracle reopening a discarded physical
coordinate.  The enlarged-core argument then restores the original
local moment constants under arbitrary compatible deterministic drops.
See \sourcefile{fixed_guard_completion_route.md} and
\sourcefile{fixed_guard_completion_audit.md}.  It repairs a local source
tree; independent innovations for arbitrary HSS assembly remain a
separate requirement.

\section{Exterior information and pilot independence}
\label{upd:sec:later-exterior}

\subsection{Comparator anchors do not describe actual residual records}

At a fixed cut $I$, all exterior \emph{comparator} anchors restrict
inside $I$ to its rank-$k$ outgoing space.  Conditioning on those
anchors removes at most $k$ local Gaussian directions.  Descendant
anchors are not included in this statement.  An orthogonal comparator
wavelet decomposition also does not make children conditionally
independent: if
\[
 \begin{pmatrix}X_a\\X_b\end{pmatrix}\Bigm|X_I
  =\begin{pmatrix}R_a\\R_b\end{pmatrix}X_I+QG,
\]
then $\operatorname{Cov}(X_a,X_b\mid X_I)=-R_aR_b^T$ in normalized
coordinates.  Fixing a parent sum creates a contrast coupling.

Actual exterior sketches can expose arbitrarily more directions, even
with a best comparator.  At a root cut $K,J$, take fixed rank-$k$
frames $U,V$ and orthonormal leaf-supported directions
$P=(p_1,\ldots,p_M)$, $Q=(q_1,\ldots,q_M)$ orthogonal to the respective
comparator frames.  Set
\[
 B=\sigma(UV^T+VU^T),\qquad E=\varepsilon P Q^T,
 \qquad 0<\varepsilon<\sigma.
\]
On the specified physical-leaf/coarse tree, $B$ is HSS-$k$, and the
root-cut singular values give
$\operatorname{OPT}_k(B+E)=M\varepsilon^2$ exactly.  Fix the narrow
probe, exterior wide rows, and comparator anchor.  Each actual exterior
leaf commutator then reveals
$\varepsilon h_jx_j$, where $h_j=\Psi_K^Tp_j$ and the observed
row $x_j=q_j^T\Omega_{J_j}$ is nonzero almost surely.  It therefore
reveals every $h_j$.  The remaining local physical Gaussian covariance
is $I-UU^T-PP^T$, losing $M$ residual directions, not just $O(k)$.
This is a conditioning obstruction, not an approximation lower bound.
See \sourcefile{correlated_source_interface_route.md} and its audit.

The same distinction applies to internal injections.  Deterministic
additions in the currently retained guard pool are reproduced exactly
and preserve the local theorem.  A genuinely omitted addition is fresh
only if its physical query directions were reserved from the entire
previous computation, including target responses and conditioning
records.  Orthogonality to the final selected template alone does not
give that independence.  Ordinary binary parent rows already belong to
the children; comparator contrasts of dimension at most $2k$ do not
automatically supply unused Gaussian coordinates.  Even a valid
reserved-source model must charge injections by their ages, since
recent injections have not undergone many residual contractions.
See \sourcefile{guarded_internal_injection_route.md}.

\subsection{Independent pilots leave exterior dependence and width costs}

Freeze an independently learned rank-$r$ column pilot $V_I$, then use
fresh probes and its annihilator $N$.  With
$R=A[I^c,I](I-V_IV_I^T)$ and annihilator width $s$,
\[
 Z=\Psi_{I^c}^TR\Omega_I N,\qquad
 \mathbb E(\|Z\|_F^2\mid V_I)=qs\|R\|_F^2.
\]
For a fixed $m$-dimensional local template the corresponding lifted
mean is $ms\|R\|_F^2/(q-m-1)$.  These are useful exact one-step costs.
But a fixed incoming cut of rank $r+1$ defeats every rank-$r$ pilot.
Conditioning on that pilot and the local design, a nonzero surviving
incoming response produces a Gaussian fit with full physical covariance.
A retained space containing it cannot be measurable in local fresh
rows alone.  Exposing a common exterior anchor subsequently correlates
siblings.  Independence of the pilot from the fresh probes does not
prove local independence of the \emph{new learned} hierarchy.
See \sourcefile{crossfit_hss_interface_route.md} and its audit.

An oracle collection of independent full-row exterior pilot ranges is
also not automatically a bounded-width nested envelope.  Along a path,
fixed shell perturbations can add $p$ mutually orthogonal directions
at each of $h$ levels.  Differences of ancestor sketches expose each
new block through an invertible $p\times p$ Gaussian factor, forcing
dimension at least $hp$ in the nested closure at the terminal node.
Conversely, intersecting a clean parent signal envelope with a noisy
child range can yield the zero space.  Projecting that signal into the
child range restores feasibility but reintroduces adaptive dependence.
These facts in \sourcefile{pilot_route.md} do not exclude all pilots;
they exclude those proposed exact-envelope constructions.

\subsection{Clean full-row linear extraction costs depth}

Let $W\in\mathbb R^{N\times M}$ collect every forward query.  An
input-independent linear extraction of a clean exterior sketch at $I$
has the form $A[I,:]WK_I$ with $W[I,:]K_I=0$.  On a path
$I_0\supset\cdots\supset I_L$, choose a fixed rank-$k$ test block
between $I_d$ and $k$ coordinates of its sibling.  Exact capture of
that test forces
\[
 \operatorname{rank}W[I_{d-1},:]
  -\operatorname{rank}W[I_d,:]\ge k.
\]
Summation gives $M\ge kL$.  Only finitely many fixed test inputs are
needed for the simultaneous probability-one assertion.  Adding
transpose queries of width $M_H$ does not evade it on nodes larger
than $M_H$: an invisible perturbation
$(I-P_H)T(I-P_W)$ shows that exact extraction either has its right
factor in $\operatorname{range}W$, or requires the entire coordinate
space of $I$ in $\operatorname{range}H$.

Finite sign averaging at one partition has the additional problem
that its variance is comparator-sized, even when the approximation
residual is zero; finite Rademacher masks can erase a fixed rank-one
interaction with positive probability.  The lower bound concerns clean
\emph{full-row}, input-independent linear extraction.  The two-sided
compressed commutator already available with $O(k)$ probes lies outside
that model, as do nonlinear and adaptive decoders.
See \sourcefile{tree_mask_route.md}.

\section{Frozen compression and shared external noise}
\label{upd:sec:later-frozen}

For independent pilot frames $P,Q$ and fresh Gaussian probes, frozen
compression has the exact conditional law
\[
 Y_c=A_c\Omega_c+N_y,\qquad X_c=A_c^T\Psi_c+N_x,
 \quad A_c=P^TAQ,
\]
where the four arrays $\Psi_c,\Omega_c,N_y,N_x$ are initially
independent.  The noise columns have fixed, possibly dense physical
covariances $C_y,C_x$, and
\[
 b=\operatorname{tr}C_y+\operatorname{tr}C_x
 \le\|A-PP^TAQQ^T\|_F^2.
\]
If the pilot contains a comparator, $b\le\|E\|_F^2$.  At each complete
level the expected commutator-noise energy is exactly
$q_\psi q_\omega b$.  These are orthogonal energy identities, not
independence between nodes; reusing the probes in learned subsequent
compressions does not renew the four-array law.
See \sourcefile{frozen_compression_noise_route.md} and its audit.

Factor $C_y=L_yL_y^T$, $C_x=L_xL_x^T$.  Every actual core computation
can be realized simultaneously as an exact commutator computation for
\[
 A_{\rm ext}=\begin{pmatrix}A_c&L_y\\L_x^T&0\end{pmatrix},
 \quad\Psi_{\rm ext}=\begin{pmatrix}\Psi_c\\G_x\end{pmatrix},
 \quad\Omega_{\rm ext}=\begin{pmatrix}\Omega_c\\G_y\end{pmatrix},
\]
with independent standard Gaussian exterior arrays.  The external
perturbation has squared Frobenius norm $b$.  This pathwise realization
preserves the actual core recursion; its root is a proper node of the
augmented problem, not a new zero-commutator global root.

Rank-one incoming noise already defeats the independent-source product
law.  Set $A_c=N_y=0$ and $N_x=vh^T$.  The physical-leaf row frames
are deterministic, so two first internal sources have independent
Gaussian $M$-row designs $G_a,G_b$ but the same $h$.  Their residuals
$e_i=(I-P_{G_i})h$ obey
\[
 \begin{aligned}
 B_i&=\|e_i\|^2=\|h\|^2Z_i,\\
 Z_a,Z_b&\ \hbox{independent }
     \operatorname{Beta}((q-M)/2,M/2),\\
 \operatorname{Cov}(B_a,B_b)&=\frac{2(q-M)^2}{q}>0.
 \end{aligned}
\]
Each marginal is the expected chi-square law.  At $q=8r,M=2r$ the
covariance is $9r$.  Exposing the complete source designs and fits gives
$h=\mu+\eta$, where, conditionally on the designs,
$\eta\sim N(0,I-P_{[G_a;G_b]})$ is independent of the fits.  The next
retained design $R$ lies in the old stacked row space, so
\[
 \|(I-P_R)h\|^2
 =\|(I-P_R)\mu\|^2+\chi^2_{q-2M}
\]
under this complete record.  Ignoring the coherent mean is incorrect.
This stronger conditioning is not a contradiction to the restricted
latent theorem; the independent-source hypothesis already fails by the
positive unconditional covariance.  Nor does the example prove physical
instability, since its core matrix is zero.
See \sourcefile{independent_noise_tree_route.md} and its audit.

There is a positive rank-slack exception.  If augmented outgoing and
incoming cut ranks fit within the maintained retained rank, its own
hybrid column hierarchy captures the relevant external incoming factors,
and its nullifiers remove them.  In particular a rank-one incoming
external perturbation of an exact HSS-$k$ core can leave the maintained
row trajectory identical to the noiseless trajectory.  The hybrid's
own column frames are essential; substituting raw companion frames
does not prove the assertion.  Low covariance trace alone does not
imply these rank inequalities.  See
\sourcefile{two_view_external_noise_route.md} and its audit.

The noncentral projector theorem supplies another one-step estimate,
without pretending inherited noise is fresh.  For fixed $M$, standard
Gaussian $Z\in\mathbb R^{t\times s}$, $s<t$, and $U=\operatorname{range}M$,
\[
 \mathbb E P_{M+Z}\preceq P_U+\frac st P_{U^\perp}.
\]
The coefficient $s/t$ is sharp; singular covariance is handled by
quotienting out deterministic mean-only columns before whitening.
There is no general opposite inequality
$\mathbb E P_{M+Z}|_U\succeq(s/t)I_U$: for
$t=3,s=2,M=[a e_1,\delta e_2]$, first large $a$ then small positive
$\delta$ gives limiting leverage $1/2$ on $e_2$, below $2/3$.
In a guarded injection this upper bound interpolates an inherited-fit
metric cost with a full pretemplate cost.  It does not bound the
inherited cost or close arbitrary mixed histories.
See \sourcefile{noncentral_projector_route.md} and its audit.

\section{Consistency correction and global correction barriers}
\label{upd:sec:later-compression}

\subsection{Consistent observations can require large matrix noise}

Noisy compressed products generally violate
$\Psi^TY=X^T\Omega$.  Orthogonal least-squares repair is explicit.
With $S=\Psi\Psi^T$, $T=\Omega\Omega^T$, its minimum-norm matrix
representative solves
\[
 SM_*+M_*T=\Psi X^T+Y\Omega^T.
\]
If the noise-free core is $A$, then
\[
 M_*=A-(I-P_\Psi)A(I-P_\Omega)+Z,
 \quad SZ+ZT=\Psi N_x^T+N_y\Omega^T.
\]
The repaired observations are products of $A+Z$.  The accurate witness
$A+Z$ need not be computable from underdetermined observations; the
computable $M_*$ loses the invisible block.  These two matrix objects
must not be identified.

In eigenbases of $S,T$, conditional noise energy equals
\[
 \mathbb E(\|Z\|_F^2\mid\Psi,\Omega)
 =\sum_{\lambda_i+\mu_j>0}
   \frac{\lambda_i c_{x,j}+\mu_j c_{y,i}}
        {(\lambda_i+\mu_j)^2}.
\]
Away from square dimensions, inverse-Wishart moments give a bound by
$b$ when each larger dimension is at least twice the smaller plus one.
Near-square behavior is different.  With $n_r=2q$, $n_c=q+2$, both
query widths $q$, and $C_x=0$, the mean correction energy is at least
$qb/2$.  A physical nullspace on the opposite side and dimension gap
at most one gives an infinite mean.  In the doubly square case the
mean is also infinite: the smallest two eigenvalues are bounded above
by independent $\chi^2_1$ variables, whose sum has an infinite inverse
mean.  This defeats an unqualified stable orthogonal consistency
repair, not every regularized estimator.  Even where finite, $Z$ depends
on the same designs and cannot be fixed while those designs are
subsequently treated as fresh.
See \sourcefile{compression_consistency_repair_route.md} and its audit.

\subsection{A few global low-rank stages need a global-tail premise}

For a known pilot $M_0$, fresh range queries on $A-M_0$ and adapted
transpose queries give an implementable correction
\[
 M_1=M_0+QQ^T(A-M_0),\qquad
 \mathbb E(\|A-M_1\|_F^2\mid M_0)
 \le\left(1+\frac{k}{p-1}\right)\tau_k(A-M_0),
\]
where $Q$ comes from a Gaussian sample of width $k+p$, $p\ge2$, and
$\tau_k$ is the squared global rank-$k$ tail.  For $p=k+1$ the total
new query cost is at most $4k+2$.  The missing premise is control of
that \emph{global} tail by $L\operatorname{OPT}_k$.

HSS rank does not supply it: $m$ disjoint unit rank-one sibling blocks
are HSS rank one but have global tail $m-k$.  A fixed number of
rank-$O(k)$ corrections from zero cannot recover them when $m$ is
large.  Similarly, with frozen $P,Q$,
\[
 \|A-PCQ^T\|_F^2
 =\|A-PP^TAQQ^T\|_F^2+\|P^TAQ-C\|_F^2;
\]
exactly solving the compressed core cannot improve the fixed geometric
floor.  Requerying full products at each level costs the rank of the
concatenated lifted query spaces, which can be proportional to the
number of levels.  Preserving free leaf-diagonal blocks is also
necessary: an identity matrix is HSS rank zero but has large global
rank.

Even orthogonal projection of an HSS-$k$ comparator onto an arbitrary
fixed nested coefficient space need not retain HSS rank $k$: different
levelwise projections can make formerly common row factors independent
on one cut.  Exact coefficient projection is not automatically
available from small fixed transcripts either.  An invisible-matrix
test shows that exact projection onto a space containing every
diagonal matrix requires
$\dim\operatorname{range}\Psi+\dim\operatorname{range}\Omega\ge N$.
These statements do not invalidate the separately proved approximate
independent refit, nor an offline best-HSS rank reduction once a good
matrix approximation has actually been obtained.  They distinguish
query complexity from a fast or directly computable optimization rule.
See \sourcefile{bounded_stage_compression_route.md} and its audit.

\paragraph{Unreviewed draft extension at the research stop.}
The final \sourcefile{replicated_global_tail_obstruction.md} was saved
as a proved draft but had not received an independent audit.  It
should not be assigned the review status of the preceding results.
Its pointwise compression argument claims that, for disjoint tested
cuts with losses $e_j$ and amplitudes $a_j$, every pilot obeying their
learned row constraints satisfies
\[
 \tau_d(A-M_0)\ge
 \sum_j|a_j|^2e_j-\sum_{\text{largest }d}|a_j|^2e_j.
\]
It proposes applying this to block-diagonal copies of the reviewed old
full-range counterexample, obtaining a failure of any fixed number of
global correction directions to restore an $O(L\operatorname{OPT})$
tail premise for that learner.  It does not target the synchronized
learner or arbitrary nonlinear structured corrections.  No new claim
of independent verification is made here.

\section{Numerical evidence and the final status map}
\label{upd:sec:later-status}

The saved experiments test different objects and should not be pooled
into a single claim.  The old full-range proof-chain batch contained
1064 recursive cases up to $N=65536$, with 32--160 independent draws
per noisy cell.  It identified the bridge mechanism and checked
identities, but the later analytic rare-angle argument proves the
$\varepsilon^{1.87}$ lower bound.  A conditional rotation diagnostic
on one child state is likewise not an unconditional expectation bound.
The guarded full-range batch contained 800 cases from 160 independent
probe pairs, reusing pairs across noise amplitudes.  At $N=131072$ and
$\varepsilon=.1$, mean signal loss rose from $.08175$ to
$.62016\pm.05154$ across the bridge.  This finite amplification neither
refutes an $O(L)$ estimate nor tests the large-rank independent-source
theorem's hypotheses.  See \sourcefile{proof_chain_numerical_summary.md},
\sourcefile{bridge_conditional_rotation.md}, and
\sourcefile{guard_fullrange_bridge_check.md}.

The synchronized learner's separate batch used 624 cases from eighty
independent probe pairs.  Its largest exact squared coefficient-space
error was $3.162\times10^{-23}$; its largest noisy-cell mean error
divided by the \emph{specified comparator residual} was
$1.00792492\pm.00335957$.  The high-rank internal union branch was
exercised 10752 times.  The experiment used an oracle projection onto
the fixed learned coefficient space, not the final random refit or a
rank-$k$ output, and did not compute best-HSS OPT.  Reused seeds across
families and amplitudes are paired cases, not extra independent draws.
The follow-up signed diagnostic reused three probe pairs and checked
400 reached oriented nodes.  Its maximum vector-identity discrepancy
was $2.721\times10^{-15}$, with both substantial cancellation and
amplification between the source and inherited-return terms.  No
expectation factorization follows from those checks.
See \sourcefile{synchronized_union_numerical_check.md},
\sourcefile{synchronized_union_numerical_audit.md}, and
\sourcefile{synchronized_signed_identity_check.md}.

The noncentral projector experiment is narrower still: fixed means,
covariances, and positive semidefinite metrics with fresh Gaussian
arrays.  Its 1150000 evaluations include 20000 repeated noiseless
deterministic evaluations, so only 1130000 are nonempty Gaussian draws.
Scalar and central equality checks, singular quotient identities, and
the anisotropic negative control agree with the proved formulas within
their recorded numerical and sampling errors.  It does not test
freshness of inherited HSS noise or closure of the mixed-history cost.
See \sourcefile{noncentral_projector_numerical_check.md} and its audit.

The endpoint can be stated without conflating these scopes:
\begin{enumerate}[leftmargin=*,label=\arabic*.]
 \item \textbf{Defeated selectors and acquisition models.}
 Minimum-norm and the stated nuclear/tree penalties fail even with
 exact leaf blocks.  The one-sided Gaussian pilot and clean full-row
 linear extraction have their stated lower bounds.  None is an
 impossibility theorem for all nonlinear two-sided decoders.
 \item \textbf{A defeated actual learner.}
 The old inverse-optimal full-range candidate fails the desired
 expected physical bound.  Its earlier mandatory-line padding gap was
 closed by the full-projector proof.  The maintained learner's positive
 theorem on that source family remains compatible with this result.
 \item \textbf{A reached local defect repaired by synchronization.}
 Separate-view propagation can incur a clean-parent loss despite zero
 child-union loss.  A common-template synchronized merge removes that
 specific defect under its visibility hypothesis.  Its exact and
 rank-slack endpoints do not prove noisy physical stability.
 \item \textbf{Local source issues with precise repairs.}
 Restricted conditioning restores the latent law; fixed guards repair
 historical target-drop padding; the finite moment ladder closes the
 sufficiently large-rank independent-source model.  These repairs do
 not manufacture independent external commutator sources.
 \item \textbf{Remaining coupled quantities.}
 Shared external noise, residual-dependent anchors, and the adaptive
 left lift remain coupled to inherited physical errors.  The proved
 unlifted signed-source and tail budgets must still be converted into
 a global expected physical budget without multiplying by an
 independently averaged inverse or double-counting cumulative losses.
 \item \textbf{Reconstruction and the unresolved conjecture.}
 Independent fixed-space refitting and query-only rank reduction remain
 available under their proved premises.  No documented result supplies
 the required all-input synchronized basis-quality premise.  Numerical
 checks do not supply it, and the final replicated-tail note remains
 an explicitly unreviewed draft extension.  The original general
 $O(k)$-query, $O(L\operatorname{OPT}_k)$ conjecture is unresolved.
\end{enumerate}

% Current proof route.  Body only; common macros and environments are supplied
% by the master manuscript.  Documentation snapshot: 20 September 2026.

\chapter{The current route and the theorem still required}
\label{upd:chap:current-route}

\section{Target, conventions, and status}

Let $\mathcal T$ be a fixed balanced binary partition tree of
$\{1,\ldots,N\}$.  The target class uses physical leaves of size at most
$2k$.  For a nonroot node $I$, write $I^c$ for its complement and $P_I$ for
the diagonal coordinate projector onto $I$.  An HSS matrix of rank at most
$k$ has outgoing and incoming exterior-cut ranks at most $k$ at every
nonroot node, with nested row and column bases on this tree.  Equivalently,
these cut ranges admit the usual nested HSS representation.  Physical-leaf
diagonal blocks are unrestricted.  Throughout this chapter,
\[
 e_\star=\OPT_k(A):=
 \inf_{B\in\HSS_k(\mathcal T)}\lVert A-B\rVert_F^2
\]
is a \emph{squared} Frobenius error.  Let $L\geq1$ bound the number of
nonroot levels.  Its order is $\log(N/k)$; changing a bounded number of
bottom levels changes this convention only by an absolute additive
constant.

The desired theorem is an algorithm which, for every fixed real matrix
$A$, returns $\widehat B\in\HSS_k(\mathcal T)$ and satisfies
\begin{equation}
 \E\lVert A-\widehat B\rVert_F^2\leq C L\OPT_k(A),
 \label{upd:eq:route-target}
\end{equation}
using $O(k)$ total products with $A$ and $A^T$ and polynomial additional
work.  The constant must be independent of dimension, depth, input
conditioning, and the size of the best-approximation error.  A method that
is exact only on HSS inputs, a method returning rank $4k$, and a
query-efficient method with an unimplemented global optimization do not
individually establish this target.

The established comparison point is the $O(kL)$-query, rank-$k$,
$O(L)$ expected squared-error guarantee of Amsel et al.\ \cite{AmselHSS}.
The reused-sketch algorithm of Levitt and Martinsson\ \cite{LM} has a
linear query count, but its stated analysis does not supply the required
arbitrary-input approximation guarantee for the adaptive coarser-level
spaces.  These are the specific primary statements relevant here; this
manuscript makes no claim that a bounded literature check proves absence
of every possible later result.  The versions and scope checks are
recorded in \sourcefile{latest_primary_sources.md}.

The current route separates three issues.  A learned nested space must
have small \emph{physical coefficient-space approximation error}; an
independent query batch must recover its coefficients stably; and the
result must be reduced to rank $k$ on the original target tree with the
required computational cost.  The coefficient-recovery step is proved.
Exact capture and several local identities for a new synchronized learner
are proved.  Its general noisy physical-stability estimate and an
efficient same-rank final reduction are not proved.

\Needspace{6\baselineskip}
\subsection{Three learners which must not be identified}

The following distinctions are part of the mathematical statements.
\begin{enumerate}
 \item The \emph{separately trained maintained union} constructs an
 annihilated hierarchy and a raw hierarchy independently, each with
 retained width $2k$, then takes their physical unions.  Its usual
 parameters are query width $16k$ per stream and leaves of size $4k$.
 The final union has width at most $4k$.  Unequal intermediate templates
 are intrinsic to that construction.
 \item The \emph{completed synchronized learner}, defined below, feeds
 each current child union into both views of the next merge.  Its
 per-view truncation width is $r=2k$, retained width is $R=4k$, query
 width is $q=32k$ per stream, and physical leaves have size at most $8k$.
 It also keeps an entire raw fit whenever that fit has rank at most $R$.
 These are substantive algorithmic changes.
 \item The \emph{fixed-source full-fit learners} studied later retain
 an entire matrix-target fit and use either specified prefix or fixed
 guard completion.  Their assumptions include independent Gaussian
 source blocks on fixed disjoint supports.  They are restricted local
 models, rather than alternative names for either HSS learner.
\end{enumerate}
An obstruction for one of these rules is not automatically an
obstruction for the others.  Conversely, a theorem for a restricted
source model cannot be imported merely because its local regression
formula resembles an HSS merge.

The synchronized learner uses a coarsening of the original target tree:
its $8k$ physical leaves combine a bounded number of $2k$ target leaves.
Every comparator on the original tree remains a valid comparator on
this coarsened tree.  The converse is unnecessary and generally false.
In particular, arbitrary dense blocks on the coarser physical leaves
must not be silently declared rank-compatible on a finer target tree.

\section{The independent fixed-space refit}
\label{upd:sec:route-refit}

Fix nested orthonormal row and column bases of width $s$, allowing
endpoint dimensions to be capped by their physical sizes.  Use physical
leaves of size at most $2s$.  Let $\mathcal S$ be their linear coefficient
space: sibling interactions have the prescribed row and column spaces,
physical-leaf diagonal blocks are arbitrary, and the root retains its
full child coefficient block.  For fixed bases this is a linear space
of matrices, not the nonlinear class of all HSS matrices of rank $s$.

For clarity, if $a,b$ are siblings, the orthogonal projection of a
matrix $A$ onto the corresponding coefficient block is
\[
 U_aU_a^T A[a,b]V_bV_b^T.
\]
Together with the reversed sibling blocks and unchanged leaf diagonal
blocks, these disjoint blocks specify the orthogonal projection onto
$\mathcal S$.  Its error is denoted
$\operatorname{dist}_F(A,\mathcal S)^2$.  Computing this projection from
the dense matrix is a useful diagnostic; it is not the query-only refit
algorithm.

\begin{theorem}[Independent fixed-space refit]
\label{upd:thm:route-refit}
Suppose the nested bases defining $\mathcal S$ are fixed independently
of two new standard Gaussian query streams.  There is a
sibling-orthogonal linear coefficient refit using width $4s$ in each
stream, hence $8s$ forward and transpose products in total, whose output
$C\in\mathcal S$ satisfies
\begin{equation}
 \E\bigl[\lVert A-C\rVert_F^2\mid\mathcal S\bigr]
 \leq1607\operatorname{dist}_F(A,\mathcal S)^2.
 \label{upd:eq:route-refit}
\end{equation}
The constant is independent of tree depth and of the geometry of the
fixed bases.
\end{theorem}

The refit uses the known nested bases and the new products to perform
linear sibling-orthogonal coefficient regressions and their prescribed
top-down corrections.  It does not learn new bases.  Its detailed
sourcewise proof and independent mathematical review are recorded in
\sourcefile{audit_refit.md} and
\sourcefile{hss_complete_proof_chain.tex}.

Here is the proof structure relevant to subsequent applications.  After
conditioning on the bases, retained and discarded query coordinates
admit the fixed-frame Gaussian quotient and Haar decompositions.  The
physical residual splits into orthogonal source sectors: interior
discarded components, sibling interaction components, and their
source-to-root terminal outputs.  Centering the corrections in the
appropriate top-down order gives orthogonality between same-sided
source contributions.  The outputs of the two query streams are not
declared independent; their combination is bounded in $L^2$.

The retained coefficient recurrence cancels the intermediate feedback
terms before any norm bound is taken.  The remaining terms are the
transported retained coefficient error and the direct discarded-source
contribution; the detailed notation and cancellation are given in the
fixed-space refit proof in
\sourcefile{hss_complete_proof_chain.tex}, as checked by
\sourcefile{audit_refit.md}.  This cancellation, rather than a sum of
separately bounded feedback terms, yields the retained-output gain.  At
query width $4s$, the quotient ratio is
$s/(q-2s)=1/2$.  The sourcewise predictable-energy bound, the terminal
source-product estimate, and the frozen-stream feedback bound are
summable.  The audited assembly, with its specified slack parameters,
gives the constant $1607$.  This explanation identifies the proof's
independence interface: conditioning on the learned bases is permitted
only because the refit streams are genuinely new.

For the synchronized learner, take $s=R=4k$.  Its physical leaves have
size $8k=2s$, so this theorem applies after training.  Training costs
$64k$ products and this refit costs $32k$, for a conditional total of
$96k$.  If a basis theorem established
\begin{equation}
 \E\operatorname{dist}_F(A,\mathcal S)^2
 \leq C_b L\OPT_k(A),
 \label{upd:eq:route-needed-space}
\end{equation}
the refit would give the same order of error in the learned rank-$4k$
coefficient space.  Equation \eqref{upd:eq:route-needed-space} is the
unproved general-input interface, not a consequence of
Theorem~\ref{upd:thm:route-refit}.

\subsection{What an exact final projection would establish}

Suppose, as a mathematical operation, one chooses a best HSS-rank-$k$
approximation $\widehat B$ to $C$ on the original target tree.  For a
best comparator $B_\star$, or by a limiting argument with approximate
minimizers,
\begin{align}
 \lVert A-\widehat B\rVert_F
 &\leq \lVert A-C\rVert_F+\lVert C-\widehat B\rVert_F\\
 &\leq2\lVert A-C\rVert_F+\sqrt{\OPT_k(A)}.
 \label{upd:eq:route-offline-reduction}
\end{align}
Consequently \eqref{upd:eq:route-needed-space} and the refit imply
\[
 \E\lVert A-\widehat B\rVert_F^2
 \leq \bigl(2\sqrt{1607C_bL}+1\bigr)^2\OPT_k(A).
\]
This is a \emph{query-only reduction}.  It does not provide a
polynomial-time implementation of best HSS projection or a verified
constant-factor same-rank recompression routine.  An additional
depth-dependent squared-error factor in a recompression theorem would
generally produce an extra factor of $L$ here.  Thus even completion of
the basis theorem would leave the polynomial-work requirement to be
settled by an appropriate efficient reduction.

\chapter{A completed synchronized learner and exact capture}
\label{upd:chap:route-synchronized}

\section{A query-only rule with all rank branches specified}

Set
\[
 r=2k,\qquad R=4k,\qquad q=32k.
\]
Draw independent standard Gaussian matrices
$\Psi,\Omega\in\R^{N\times q}$ once, and obtain
$Y=A\Omega$ and $Z=A^T\Psi$.  Every node uses these same original
arrays.  Its accessible response is
\begin{equation}
 J_I=\Psi_I^TY_I-Z_I^T\Omega_I
     =\Psi^T[P_I,A]\Omega\in\R^{q\times q}.
 \label{upd:eq:route-commutator}
\end{equation}
No Gaussian design is resampled at an internal node.

Write $N_D$ for a matrix with orthonormal columns spanning the right
kernel of $D$.  Using this basis instead of the kernel projector changes
only the right coordinates of a nullified sample, not its positive left
singular spaces or the selected physical ranges.  At a physical leaf
$I$, form the row and column samples
\[
 Y_I N_{\Omega_I},\qquad Z_I N_{\Psi_I}.
\]
If a sample has rank at most $R$, keep its full positive left singular
range.  Otherwise keep its leading $r$ positive left singular
directions.  Complete in physical coordinate order to dimension
$\min(R,|I|)$.  A zero sample therefore receives the specified fixed
coordinate prefix.  This full-range branch is part of the algorithm;
it is not a proof convention.

At an internal nonroot node with children $a,b$, both views use the
same completed child unions:
\[
 W=\diag(U_a,U_b),\qquad W_V=\diag(V_a,V_b),
\]
where disjoint child supports are understood.  Their widths are at
most $2R=8k$.  In row coordinates put
\begin{equation}
 G=W^T\Psi_I,\quad H=W_V^T\Omega_I,\quad
 L_G=(G^\dagger)^T,\quad
 \Pi_G=G^\dagger G,\quad Q_G=L_GG^T.
 \label{upd:eq:route-local-design}
\end{equation}
Here $\Pi_G$ acts in query space and $Q_G$ acts in physical-template
coordinates.  Both are orthogonal projectors; these definitions remain
valid for singular $G$.  Form
\[
 M_0=\Pi_GJ_I,\qquad M_1=\Pi_GJ_I N_H.
\]
If $\rank M_0\leq R$, the mandatory local space is
$\range(L_GM_0)$.  Otherwise, take the leading $r$ positive left
singular directions of each $M_0,M_1$, lift them by $L_G$, and take
their union.  This mandatory union has dimension at most $R$.

There is one common completion, performed \emph{after} this union.
Starting with the mandatory space, scan
\[
 Q_Ge_1,\ldots,Q_Ge_m,
 \quad (I-Q_G)e_1,\ldots,(I-Q_G)e_m,
\]
adjoining each nonzero orthogonalized candidate until dimension
$\min(R,m)$ is reached.  A fixed projector-based coordinate convention
chooses the final orthonormal frame $U$ of this space, and the physical
row basis is $WU$.  The two views are not padded separately.
The local projector $P=UU^T$ obeys
\begin{equation}
 PQ_G=Q_GP.
 \label{upd:eq:route-commuting-completion}
\end{equation}
Indeed the mandatory directions are visible, and every additional
direction is taken within either the visible subspace or its orthogonal
kernel.  When $G$ has full row rank, this is ordinary coordinate
completion.

The column rule is the transpose version, using $-J_I^T$ and exchanging
$G,H$.  At the global root there is no exterior-cut selection: the full
child coefficient block is retained.  All operations are local linear
algebra on the original products and stored bases.  This construction,
its exact theorem, and an independent proof audit are recorded in
\sourcefile{synchronized_union_exact_route.md} and
\sourcefile{synchronized_union_exact_audit.md}.

\section{Exact and rank-slack inputs}

\begin{theorem}[Exact capture with rank slack]
\label{upd:thm:route-exact}
Fix an input $A$ for which every outgoing and incoming nonroot exterior
cut on the learner tree has rank at most $2k$.  Almost surely, all reached
internal predesigns have full row rank, and the synchronized row and
column spaces contain the corresponding true cut ranges.  This includes
zero cuts, deficient ranks, and physical endpoint caps.  In particular
$A$ belongs to the learned coefficient space.  Every exact HSS-rank-$k$
input satisfies these hypotheses.
\end{theorem}

The full-row-rank conclusion is a theorem on this exact rank-slack
family.  It does not say that learned designs are Gaussian, and it is
not an arbitrary-noisy-input smallest-singular-value estimate.

\begin{proof}[Proof outline with the conditioning interface]
Let $a_I,b_I\leq2k$ be the outgoing and incoming cut ranks.  For every
probe realization,
\[
 \rank J_I\leq a_I+b_I\leq4k=R.
\]
Thus the rule collapses to the full-raw branch before any conditioning.
The row recursion no longer depends on the opposite learned hierarchy.

Factor the outgoing block as $C_IF_I^T$, with $F_I$ an orthonormal
exterior factor, and write the incoming transpose block as $V_IK_I$,
with $V_I$ orthonormal on $I$.  Fix $\Omega$ on the probability-one
event on which its finitely many true-cut anchors and leaf nullified
samples have their required ranks.  The disjoint-support anchors
\[
 \Gamma_I=F_I^T\Omega_{I^c},\qquad X_I=V_I^T\Omega_I
\]
admit a separator $T_I$ with $\Gamma_IT_I=I$ and $X_IT_I=0$.
At a leaf, $q-|I|\geq24k$ is more than sufficient for the outgoing
rank.  The leaf row space therefore consists of the true outgoing cut
range and a fixed canonical completion under this conditioning.

For the induction, choose deterministic orthonormal true frames $C_I$
on the outgoing range and $Z_I$ on the exterior incoming range, and put
\[
 h_I=\Psi_I^TC_I,\quad e_I=\Psi_{I^c}^TZ_I,
 \quad \mathcal H_I=\range(h_I),\quad
 \mathcal K_I=\range[h_I,e_I].
\]
If $\mathcal S_I$ is the query image of the retained row space and its
dimension is $d_I\leq R$, the induction maintains injectivity,
true-cut containment, and
\[
 \mathcal S_I\cap\mathcal K_I=\mathcal H_I.
\]
The exterior incoming range for child $a$ is contained in the sum of
the sibling outgoing range and the parent's exterior incoming range.
Complete these fixed physical frames orthogonally.  The newly exposed
Gaussian anchor is independent of the child's old physical block and
old exterior anchor.  The rank inequalities
$d_a+a_b+b_I\leq8k$ and their symmetric versions preserve the required
child intersections after these anchors are exposed.

Condition now on $\Omega,h_a,h_b,e_I$, rather than on the entire old
query designs.  The remaining child Gaussian blocks are independent.
Rotations fixing $\range[h_a,h_b,e_I]$ preserve the collapsed full-fit
rule and its canonical completion.  The two residual child query
subspaces are therefore independent Haar Grassmann subspaces in the
remaining query coordinates.  Their available dimensions, and
\[
 d_a+d_b+b_I\leq10k<q,
\]
give disjointness and injectivity of the stacked parent design almost
surely.  The corresponding intersection invariant is preserved.

The child containment implies $C_I\in\range W$.  Multiplication by the
fixed separator gives
\[
 L_GJ_IT_I=L_GG^TW^TC_I=W^TC_I,
\]
where injectivity was just established.  The full-raw branch retains
this range, proving parent capture.  The transpose assertion follows
by reversing the roles of the two streams and repeating the argument.
One does not condition on both learned hierarchies and then declare
the resulting matrices independent Gaussian.
\end{proof}

\section{The exact clean-parent statement}

Fix $A=B+E$ with a fixed HSS-rank-$k$ comparator $B$.  Suppose at a
particular parent that the common row template $W$ contains the
comparator outgoing cut $C$, and that $J_I(E)=0$.  A fixed-comparator
separator as constructed in the next chapter gives
\[
 L_GJ_I(B)T=Q_GW^TC.
\]
Moreover $\rank J_I(B)\leq2k$, so the entire raw lift is retained.
Consequently
\begin{equation}
 (I-WPW^T)C
   =(I-WPW^T)W(I-Q_G)W^TC.
 \label{upd:eq:route-clean-invisible}
\end{equation}
Only comparator directions invisible to the common design can be lost.
If that design is injective, the parent captures the comparator cut
exactly, without a hypothesis on containment in the learned opposite
frame.  This conclusion is a local deterministic statement.  It does
not identify the actual feedback trajectory with the trajectory of a
separately trained two-view hierarchy.

\chapter{The signed source identity and the remaining physical budget}
\label{upd:chap:route-signed}

\section{A separator determined by the fixed comparator}

Fix $A=B+E$, where $B\in\HSS_k(\mathcal T)$ is deterministic.  At a
node $I$, factor
\[
 B[I,I^c]=CF^T,\qquad B[I^c,I]^T=VK,
\]
with $F,V$ orthonormal, of widths $a,b\leq k$ respectively.  Extend
$F$ by zero on $I$ and $V$ by zero outside $I$.  Set
\[
 \Gamma=F^T\Omega,\quad X=V^T\Omega,\quad
 \Theta=K\Psi_{I^c}.
\]
The comparator commutator is
\begin{equation}
 J_I(B)=\Psi_I^TC\Gamma-\Theta^TX.
 \label{upd:eq:route-comparator-J}
\end{equation}
The two rows of notation in this formula concern fixed true cut
factors, not learned bases.  If $N_X$ spans $\ker X$, define
\begin{equation}
 T=N_X(\Gamma N_X)^\dagger.
 \label{upd:eq:route-separator}
\end{equation}
Almost surely $\Gamma T=I_a$ and $XT=0$.  The separator depends on the
fixed comparator and $\Omega$, and is independent of $\Psi$.  It does
not depend on the learned annihilator $H$.  This distinction is what
makes its fixed-source moments available.

For the local template, decompose
\[
 C=WC_c+C_\perp,\qquad C_c=W^TC,
 \quad C_\perp=(I-WW^T)C.
\]
Then, even for singular $G$,
\begin{equation}
 L_GJ_I(B)T=Q_GC_c+L_G\Psi_I^TC_\perp.
 \label{upd:eq:route-separated-signal}
\end{equation}
The inherited defect $C_\perp$ is an adaptive physical matrix; it must
not be treated as fixed when applying Gaussian moment formulas.

\section{A raw tail with an unlifted budget}

Let $U$ be the retained local coordinate frame and $P=UU^T$.  Define
\begin{equation}
 R_0=\bigl(\Pi_G-\operatorname{Proj}(G^TU)\bigr)J_I,
 \label{upd:eq:route-raw-tail}
\end{equation}
where $\operatorname{Proj}(K)=KK^\dagger$ projects onto a column range.
The expression $G^TU$, rather than $GU$, is essential.  The two query
projectors are nested, so their difference is an orthogonal projector.

In the high-rank branch, the retained union contains the leading $2k$
raw directions after lifting, and mapping them back by $G^T$ recovers
those query directions.  Since $\rank\Pi_GJ_I(B)\leq2k$, the
Eckart--Young comparison gives
\begin{equation}
 \lVert R_0\rVert_F\leq\lVert J_I(E)\rVert_F.
 \label{upd:eq:route-unweighted-tail}
\end{equation}
In the full-raw branch $R_0=0$, so the same inequality holds.  No
Gaussian independence is used in this pathwise statement.

For a fixed matrix $M$ and two independent Gaussian streams,
\[
 \E\lVert\Psi^TM\Omega\rVert_F^2=q^2\lVert M\rVert_F^2.
\]
At any one partition level,
$\sum_I\lVert[P_I,E]\rVert_F^2\leq2\lVert E\rVert_F^2$.
Summing over the nonroot levels yields
\begin{equation}
 \sum_I\E\lVert R_{0,I}\rVert_F^2
 \leq2q^2L\lVert E\rVert_F^2.
 \label{upd:eq:route-unweighted-global}
\end{equation}
This is an unlifted query-space budget.  It has not yet paid for the
adaptive dual lift.

\subsection{The stronger separator-weighted source budget}

Let $P_\perp=I-FF^T-VV^T$, with the globally embedded factors above.
For every fixed $M$,
\begin{equation}
 \E\lVert M\Omega T\rVert_F^2
 =\lVert MF\rVert_F^2+
   \frac{a}{q-a-b-1}\lVert MP_\perp\rVert_F^2.
 \label{upd:eq:route-fixed-separator-moment}
\end{equation}
To see this, decompose $\Omega$ along the mutually orthogonal physical
ranges $F,V,P_\perp$.  The first term is fixed by $\Gamma T=I$, the
second disappears by $XT=0$, and the remaining Gaussian block is
independent of the two anchors.  The inverse-Wishart trace identity for
$\Gamma N_X$ gives the displayed coefficient.  Independence of $\Psi$
then multiplies this expression by $q$ when $\Psi^T$ is put on the left.
Since $a,b\leq k$ and $q=32k$, taking $M=[P_I,E]$ gives
\begin{equation}
 \E\lVert J_I(E)T\rVert_F^2
 \leq q\lVert[P_I,E]\rVert_F^2.
 \label{upd:eq:route-source-weighted}
\end{equation}

It is false in general that the raw-tail inequality can simply be
right-weighted pathwise.  For example, a leading-one truncation of
$\diag(2,1)$, with comparator $\diag(0,1)$, residual $\diag(2,0)$, and
$T=e_2$, has weighted tail norm one and weighted residual norm zero.
The required tail estimate instead uses moments.  A fixed net and the
directional inverse-chi-square law give
\[
 \bigl(\E\lVert T\rVert_{\mathrm{op}}^4\bigr)^{1/2}
 \leq\frac{2\,9^{1/3}}{24k+1}.
\]
One way to track the dimensions is to set
$d=q-b$, $\nu=d-a+1$, and use order $3a$ for a net of size at most
$9^a$.  The inverse-chi-square moment is bounded by
$1/(\nu-6a)$ at this order, and the net comparison contributes its
factor two.  The case $a=0$ is vacuous.

For the other factor, direct Gaussian fourth moments give
\[
 \E\lVert\Psi^TM\Omega\rVert_F^4
 \leq[q(q+2)]^2\lVert M\rVert_F^4.
\]
Combining this with \eqref{upd:eq:route-unweighted-tail} by
Cauchy--Schwarz, without assuming independence of $R_0$ and $T$, gives
\[
 \E\lVert R_0T\rVert_F^2
 \leq6q\lVert[P_I,E]\rVert_F^2.
\]
Minkowski and \eqref{upd:eq:route-source-weighted} therefore imply the
useful signed-source estimate
\begin{equation}
 \boxed{\quad
 \sum_I\E\lVert(R_{0,I}-J_I(E))T_I\rVert_F^2
 \leq24qL\lVert E\rVert_F^2.
 \quad}
 \label{upd:eq:route-24q}
\end{equation}
The local constant used is
$(\sqrt6+1)^2<12$, followed by the commutator sum bound.  This result
is recorded with the local identities in
\sourcefile{synchronized_union_route.md}.  It prices the right
separator and the source tail, but not the adaptive left lift.

\section{The exact projected signed identity}

For the completed synchronized rule define
\begin{align*}
 Z_I^{\mathrm{src}}&=(I-P)L_G(R_0-J_I(E))T,\\
 D_I^{\mathrm{ret}}&=(I-P)L_G\Psi_I^TC_\perp,\\
 H_I^{\mathrm{inv}}&=(I-P)(I-Q_G)C_c.
\end{align*}
Then
\begin{equation}
 (I-P)C_c=Z_I^{\mathrm{src}}-D_I^{\mathrm{ret}}
                  +H_I^{\mathrm{inv}}.
 \label{upd:eq:route-signed}
\end{equation}
Indeed, left projection by $I-P$ annihilates the lifted retained query
range, so $(I-P)L_GJ_I=(I-P)L_GR_0$.  Substitute
\eqref{upd:eq:route-separated-signal}, subtract the residual commutator,
and restore the null component of $C_c$.  The completion property
\eqref{upd:eq:route-commuting-completion} makes the first two terms visible
and the last term invisible; hence these sectors are orthogonal.
Since $C_\perp$ is physically orthogonal to $W$, the exact squared
identity is
\begin{equation}
 \lVert(I-WPW^T)C\rVert_F^2
 =\lVert C_\perp\rVert_F^2
  +\lVert Z_I^{\mathrm{src}}-D_I^{\mathrm{ret}}\rVert_F^2
  +\lVert H_I^{\mathrm{inv}}\rVert_F^2.
 \label{upd:eq:route-signed-energy}
\end{equation}
This formula preserves the sign and correlation of the source and
return terms.  Bounding them separately may destroy the cancellation
which the formula is meant to expose.  With arbitrary completion
mixing visible and invisible coordinates, the omitted retained-query
term would not in general vanish; the specified completion is therefore
part of the identity's application to the learner.

\section{The precise outstanding interface}

Equation \eqref{upd:eq:route-24q} does not imply a physical budget after
left multiplication by $(I-P)L_G$.  This lift, the tail $R_0$, the
separator, and the return term all share the training data.  Products
of their separate expectations are not legitimate substitutes for a
joint estimate.  Nor is $C_\perp$ fixed after conditioning only on the
comparator.  The invisible term remains relevant whenever the common
predesign is singular, even if every retained child design is injective.

One sufficient all-input goal is a bound of order
$L\lVert E\rVert_F^2$ for the sum of the actual outgoing and incoming
comparator cut losses produced by \eqref{upd:eq:route-signed-energy}, with
the inherited defects accounted for without an additional depth
factor.  An equivalent or stronger direct bound on the final
coefficient-space error would also suffice.  The latter is the exact
interface needed by the refit, rather than a scalar inverse trace alone.

To make the global requirement explicit, let $\mathcal S$ be the learned
space.  The sibling-block projection and its orthogonal row/column
decomposition give the deterministic bound
\begin{align}
 \operatorname{dist}_F(B,\mathcal S)^2
 \leq\sum_{I\ne\mathrm{root}}\bigl(&
 \lVert(I-U_IU_I^T)B[I,I^c]\rVert_F^2\nonumber\\[-2pt]
 &+\lVert B[I^c,I](I-V_IV_I^T)\rVert_F^2\bigr).
 \label{upd:eq:route-global-cut-interface}
\end{align}
Each sibling block is charged to its own row and column endpoints;
the displayed full exterior cuts only enlarge those nonnegative
charges.  Also
\[
 \operatorname{dist}_F(A,\mathcal S)
 \leq\lVert E\rVert_F+\operatorname{dist}_F(B,\mathcal S).
\]
Thus a valid joint physical budget in
\eqref{upd:eq:route-global-cut-interface} would imply
\eqref{upd:eq:route-needed-space}.  The unresolved point is establishing
that budget along the actual adaptive recursion, not writing the final
deterministic coefficient projection.  No telescoping identity beyond
the reviewed local statements is presumed here.

\chapter{What is proved for independent source trees}
\label{upd:chap:route-sources}

\section{The source interface is an assumption}

The results in this chapter concern fixed physical source templates on
disjoint supports.  Fix all source physical data, and any allowed
conditioning on an auxiliary stream $\Omega$, \emph{before} drawing the
source blocks of $\Psi$.  These blocks are independent standard
Gaussian matrices.  A source $\alpha$ has an orthonormal pretemplate
$W_\alpha$ of width $2r$ and a fixed matrix target $U_\alpha$ with at
most $s\leq r/2$ columns.  Its original omission Gram is
\[
 C_\alpha=U_\alpha^T(I-W_\alpha W_\alpha^T)U_\alpha\succeq0.
\]
The represented coefficients may be arbitrarily large, and
$C_\alpha$ may be singular or zero.  Use $q=8r$ shared query columns.
The source response is $H_\alpha=\Psi_\alpha^TU_\alpha$.
At a parent the pretemplate is the direct sum of the retained child
frames, and the response is the sum of the compatible child responses.
The full physical least-squares fit is retained before completion.
Concretely, if $G=W^T\Psi_I$ is the current predesign and $H_I$ its
response, the fit is $F_I=W(G^\dagger)^TH_I$.  In the additive
ordered-prefix model, append the current template's coordinate vectors
in their fixed order until the retained dimension is $r$, then
recanonicalize the resulting projector in that same order.  The parent
template orders the left child before the right child.  The fixed-guard
model replaces this optional coordinate list by its specified pool.

In the additive equal-output model, the physical parent target is the
direct sum of its children and
$C_I=\sum_{\alpha\preceq I}C_\alpha$.  A more general compatible-transfer
model allows deterministic matrices $A_a,A_b$ with
\[
 U_I|_i=U_iA_i,\qquad H_I=H_aA_a+H_bA_b,
 \qquad C_I=A_a^TC_aA_a+A_b^TC_bA_b.
\]
The guarded completion below handles strict drops in target spaces.
Targets and transfers are inputs to this local interface; the theorem
does not claim that the learner knows the physical target matrices.

There is a substantive difference between this assumption and a generic
HSS residual.  In an HSS construction, a child's response can contain
exterior Gaussian anchors already exposed elsewhere in the computation.
Arbitrary dense residual anchors need not have comparator-rank-sized
span, and two independent arrays of added product noise do not make
their adaptive subtrees independent.  No decomposition theorem is
asserted here which replaces those effects by independent fixed sources
at total cost $O(\OPT_k(A))$.

\section{The earlier fixed-core inverse-optimal theorem}

There is a proved small-rank theorem in a related local regression
model which should not be conflated with the full-fit theorem below.
At its physical leaves, fix rank-$r$ frames $U_\alpha$ and source
matrices $E_\alpha$ satisfying $U_\alpha^TE_\alpha=0$.
Put $\beta=\sum_\alpha\lVert E_\alpha\rVert_F^2$ and draw one Gaussian
propagation stream $X$.  At a merge let $W=\diag(U_a,U_b)$,
$G=W^TX_I$, and $F=[E_a^{\mathrm{out}};E_b^{\mathrm{out}}]$.
An orthogonal split $[T,S]$ retains $U_I=WT$ and defines
\[
 B_I=S^T(G^\dagger)^TX_I^TF,\qquad
 E_I^{\mathrm{out}}=F-WSB_I,\qquad
 c_I=(R_I^\dagger)^TX_I^TE_I^{\mathrm{out}},
 \quad R_I=U_I^TX_I.
\]
Every selector must be subtree-local and invariant in physical
coordinates under right rotations of $X_I$.  Independent side
information may be fixed first.  These requirements are part of the
theorem, even if the physical mandatory cores are fixed.

Prescribe nested orthonormal physical core frames $Q_I^{\mathrm{phys}}$
of ranks $s_I\leq k$, independently of $X$, and require the fixed leaf
frames to contain their leaf cores.  At ordinary nodes use the
inverse-optimal rank-$r$ completion containing
$Q_I=W^TQ_I^{\mathrm{phys}}$.  The objective is the trace of the inverse
retained query Gram.  If $K=GG^T$ and
$\phi_I=\tr[(Q_I^TKQ_I)^{-1}]$, its optimizer obeys the deterministic
additive bound
\begin{equation}
 V_I\leq\min(V_a,V_b)+\phi_I.
 \label{upd:eq:route-additive-core}
\end{equation}
A Schur complement reduces the optional-space optimization to the
smallest generalized eigenvalues of the complementary inverse metric.
Their eigenvalue comparison with $K^{-1}$ and principal-submatrix
interlacing prove \eqref{upd:eq:route-additive-core}.

Exact physical containment cancels the adaptive template in the
mandatory Gram:
\[
 Q_I^TGG^TQ_I
 =(Q_I^{\mathrm{phys}})^TX_IX_I^TQ_I^{\mathrm{phys}}.
\]
Thus this \emph{mandatory} Gram is Gaussian, although the whole selected
Gram is not.  With $s=s_I$,
\[
 \E\phi_I=\frac{s}{q-s-1},\qquad
 \E\phi_I^2\leq\frac{s^2}{(q-s-1)(q-s-3)}.
\]
For $r=2k,q=16k$, the local innovation-energy telescope and a
Cauchy--Schwarz estimate for the dependent forced-core term give
\begin{equation}
 \E\lVert E_{\mathrm{root}}^{\mathrm{out}}\rVert_F^2
 +\E\lVert c_{\mathrm{root}}\rVert_F^2
 \leq\left(\frac{6694}{1001}+\frac{2\sqrt3}{77}\right)\beta
 <6.733\beta.
 \label{upd:eq:route-fixed-core}
\end{equation}
When \emph{every ordinary mandatory rank is exactly $k$}, the additional
averaging envelope
$V_I\leq\tfrac13\tr K^{-1}+\tfrac23\phi_I$ improves the coefficient to
\begin{equation}
 \frac{455+4\sqrt3}{153}<3.020.
 \label{upd:eq:route-half-core}
\end{equation}
The final root need not use the optimizer.  These constants belong to
different core-rank patterns and cannot be interchanged.

The proof extends to adaptive cores contained in independently fixed
physical envelopes of dimension $m<q-3$, provided feasibility in the
actual pretemplate and the same locality and equivariance hold.  Inverse
compression bounds the forced inverse by the inverse trace of the
fixed envelope's Gaussian Gram, giving first moment
$m/(q-m-1)$ and second moment at most
$m^2/[(q-m-1)(q-m-3)]$.  It does not prove that a general noisy HSS
learner has such feasible low-dimensional envelopes.  The theorem and
envelope extension are reviewed in \sourcefile{audit_nonlinear.md},
with the detailed local recurrences in
\sourcefile{hss_complete_proof_chain.tex}.  This earlier result remains
a valid restricted stability theorem; neither it nor the later
fixed-source theorem is a learned-noisy-core guarantee for arbitrary
HSS inputs.

\section{An exact current-state latent covariance}

Let $R_I$ be the retained query design, $F_I$ the full physical fitted
target, and $e_I$ the current query residual.  The relevant law is a
law of a restricted current state.  One may retain physical bases and
fits, retained and predesign Grams, and the prescribed functions of
these summaries.  One may not also condition on discarded residual
Grams, old query-coordinate designs, old response Grams, or arbitrary
cross-Grams involving forgotten designs.  Those enlarged records can
destroy the Wishart law used by the representation.

Under the permitted summary, the actual joint state has an equivalent
representation with a positive semidefinite latent covariance $T_I$.
Conditionally on it and the current design, the residual is centered
matrix Gaussian with row covariance the projector onto
$\ker R_I$ and column covariance $T_I$.  For an additive merge let
$S=T_a+T_b$; with transfers replace this by
$S=A_a^TT_aA_a+A_b^TT_bA_b$.  On the compact support of $S$, of rank $t$,
the new covariance is
\begin{equation}
 T_I=D_S\,\mathcal B\,D_S^T,
 \quad S=D_SD_S^T,
 \quad \mathcal B\sim\operatorname{Beta}_t(3r,r/2),
 \label{upd:eq:route-matrix-beta}
\end{equation}
with a fresh compact beta factor in this equivalent construction.
In particular,
\[
 0\preceq T_I\preceq S,\qquad
 \E[T_I\mid S]=\frac67S.
\]
This is a distributional representation of the original shared-probe
algorithm's permitted current state.  It is not an instruction to
resample Gaussian probes in the physical algorithm.

For completeness, the matrix identity behind the renewal is elementary.
Let $A\sim W_t(6r,I)$ and $D\sim W_t(r,I)$ be independent, put
$X=A+D$, and use the symmetric square root to set
$\mathcal B=X^{-1/2}AX^{-1/2}$.  The congruence Jacobian factors the
joint density into a Wishart density for $X$ and a beta density for
$\mathcal B$, proving their independence.  For independent copies
with these laws,
\[
 X^{1/2}\mathcal B X^{1/2}\sim W_t(6r,I),
 \qquad
 \mathcal B^{1/2}X\mathcal B^{1/2}\sim W_t(6r,I).
\]
The second assertion follows because the two matrices have the same
eigenvalues and both laws are orthogonally invariant.  It does not
interchange noncommuting square roots pointwise.  Factoring $S$ only
on its support handles singular covariances, with $T_I=0$ when $S=0$.
The complete joint-law proof and review are in
\sourcefile{matrix_latent_budget.md} and
\sourcefile{matrix_latent_budget_audit.md}.

\section{A depth-independent theorem at sufficiently large retained rank}

\begin{theorem}[Fixed-source moment closure]
\label{upd:thm:route-source-large-rank}
In the fixed, disjoint, additive matrix-target source model with the
specified ordered-prefix completion, suppose $r\geq4096$ and $q=8r$.
There is an absolute finite constant $B_0$ such that for every finite
height, every node $I$, and every deterministic output vector $v$,
\begin{equation}
 \left\|\lVert(F_I-U_I)v\rVert^2\right\|_{L^2}
 \leq B_0\,v^TC_Iv.
 \label{upd:eq:route-large-rank-energy}
\end{equation}
The constants are independent of height, represented coefficients,
source geometry, and nonzero covariance eigenvalues.  The actual
retained physical inverse has a uniformly bounded first moment.
\end{theorem}

Here the actual inverse is
$V_I=\tr[(R_IR_I^T)^{-1}]$, and ``uniformly bounded'' refers to its
unconditional expectation under the fixed-source model, not a bound
conditional on arbitrary fitted or discarded data.  The theorem,
its root review, and a separate independent audit are recorded in
\sourcefile{fitted_energy_bootstrap_route.md},
\sourcefile{fitted_energy_bootstrap_root_audit.md}, and
\sourcefile{fitted_energy_bootstrap_audit.md}.

\begin{proof}[Finite moment ladder]
For an integer $p$ define $E_p(h)$ as the supremum of
\[
 \frac{\|\lVert(F_I-U_I)v\rVert^2\|_{L^p}}{v^TC_Iv}
\]
over the stated model class at height $h$, all $r\geq4096$, and all
deterministic $v$ with positive denominator.  If the denominator is
zero, every source is exact in that direction, and full-fit retention
propagates exactness.  Thus no undefined zero-budget directions are
hidden in this normalization.

For an integer $p\geq1$, a Gaussian vector $g$ of covariance $K$ obeys
\begin{equation}
 \|\lVert a+g\rVert^2\|_{L^p}
 \leq\lVert a\rVert^2+(2p-1)\tr K.
 \label{upd:eq:route-gaussian-energy}
\end{equation}
For a scalar coordinate, the coefficient of
$a^{2p-2j}\sigma^{2j}$ in the Gaussian even moment, divided by
$\binom pj$, is
$\prod_{\ell=0}^{j-1}(2p-2\ell-1)\leq(2p-1)^j$.
Comparison with the binomial expansion proves the scalar inequality;
Minkowski over physical coordinates proves
\eqref{upd:eq:route-gaussian-energy} without requiring coordinate
independence in the final step.

The enlarged-core captured-inverse estimate gives, on the compact
current target support,
\[
 \bigl\|\lVert D_w^{-1}\rVert_{\mathrm{op}}\bigr\|_{L^m}
 \leq\frac{2^{15}}r,
 \qquad1\leq m\leq\frac{3r}{16}.
\]
Its small-eigenvalue tail is bounded by
$\tfrac12(x/2^{-14})^{3r/8}$ for $0<x\leq2^{-14}$ after normalization
by $r$.  Integrating this tail yields the stated moment range.  A
deterministic enlarged core accounts for all retained directions, so
there is no unexcited padding contribution omitted from the inverse.

For a parent direction $v$, set $c=v^TC_Iv$ and
$\ell=v^T(T_a+T_b)v/c$ in the additive model.  The covariance trace of
the physical Gaussian fit increment is bounded by
$c\ell V_{\mathrm{pre}}$.  The inverse bridge and compact target
whitening express the needed joint product in terms of the fixed-core
inverse, captured inverse, and fitted energy.  Hölder is applied to
these quantities in their \emph{same joint realization}; independence
is not presumed.  At orders $p=2,4,\ldots,128$, the available moments
and beta decay yield
\begin{equation}
 E_p(h+1)\leq E_p(h)
       +K_0\rho^h\bigl(1+E_{2p}(h)\bigr),
 \quad K_0=55{,}000{,}000,\quad \rho=\frac78.
 \label{upd:eq:route-moment-ladder}
\end{equation}
The source estimate is $E_p(0)\leq2$ at the orders used.  The coarse
terminal theorem gives
\begin{equation}
 E_{256}(h)\leq\frac{2567}{6}\left(\frac52\right)^h.
 \label{upd:eq:route-terminal}
\end{equation}
It requires only unconditional inverse moments from its own valid
orbit disintegration.

There are seven, rather than infinitely many, moment reductions.  Put
$\gamma=5/2$, $p_j=2^{j+1}$, and
$\lambda_j=\gamma\rho^{7-j}$ for $1\leq j\leq7$.
All these $\lambda_j$ exceed one, while
\[
 \gamma\rho^7=\frac{4117715}{4194304}<1.
\]
Starting with $B_7=2567/6$, choose successively
\[
 B_j=2+\frac{K_0}{1-\rho}
       +\frac{K_0B_{j+1}}{\lambda_j-1},
 \qquad j=6,\ldots,1.
\]
Summing \eqref{upd:eq:route-moment-ladder} proves
$E_{p_j}(h)\leq B_j\lambda_j^h$.  The final step gives the finite
absolute bound
\[
 E_2(h)\leq B_0:=2+\frac{K_0}{1-\rho}
             +\frac{K_0B_1}{1-\gamma\rho^7}.
\]
The rank threshold guarantees every Hölder moment used in the seven
steps.  This proves \eqref{upd:eq:route-large-rank-energy}.
\end{proof}

In particular,
\[
 \E[(F_I-U_I)^T(F_I-U_I)]\preceq B_0C_I,
 \qquad
 \|\lVert F_I-U_I\rVert_F^2\|_{L^2}\leq B_0\tr C_I.
\]
One explicit inverse consequence is
\[
 \E V_I\leq\frac8{55}
 +2^{27/2}\left(1+\sqrt{B_0}+\frac1{\sqrt7}\right)^2.
\]
These deliberately conservative constants prove uniformity, not a
practical error prediction.  The scalar model $r=2,q=16$ lacks the
moment orders used here; the theorem does not cover it.

\section{Fixed guards and compatible target drops}

For deterministic compatible target transfers with
$s_I\leq s_{\max}\leq r/2$, set $c=r-s_{\max}$.  At each source define
the guard as the first $c$ columns of $W_\alpha$ and the pool as its
first $2c$ columns.  At a parent let
\[
 \operatorname{Pool}_I=\diag(\operatorname{Guard}_a,
                             \operatorname{Guard}_b),
 \qquad
 \operatorname{Guard}_I=\text{first $c$ columns of the pool}.
\]
These physical vectors are fixed before the source Gaussian queries.
Begin with the full fitted target range, scan the pool in its fixed
order until rank $r$ is reached, and choose the frame with its guard
first.  Since the fit rank is at most $s_{\max}$ and the pool has at
least $r$ columns, the scan is feasible and retains the guard.

The key deterministic fact is stronger than feasibility.  For a pool
prefix $C_j$, put
$\kappa(j)=\rank[(I-C_jC_j^T)U_I]$.  The iteration
\[
 j_0=c,\qquad j_{n+1}=r-\kappa(j_n)
\]
increases to a fixed point $j_\star+\kappa_\star=r$.
The full prefit is exact on each contained prefix, so the actual scan
retains $C_{j_\star}$.  The compact Gaussian capture theorem shows that
the fit adds exactly $\kappa_\star$ directions outside it, almost
surely.  Hence the actual retained space is precisely the span of a
deterministic enlarged core and a full compact fit; no arbitrary
unexcited completion remains.  The iteration is a proof device.  The
algorithm only tests observed fitted ranks while scanning its known
pool.

The compact physical Gram outside this enlarged core dominates the
transferred original omission Gram.  It therefore supplies a common
deterministic normalization for the inverse bridge.  The matrix-beta
law and the seven-step ladder retain their constants, extending
Theorem~\ref{upd:thm:route-source-large-rank} to this guarded compatible
transfer model.  The proof and independent review are in
\sourcefile{fixed_guard_completion_route.md} and
\sourcefile{fixed_guard_completion_audit.md}.  A represented injection
lying in the fixed pool is fitted exactly and does not create new
latent error; arbitrary new unrepresented injections require the
separate freshness conditions below.

\section{Two scalar results which remain useful special cases}

\subsection{One original residual source and arbitrary clean siblings}

A different scalar theorem uses fixed original rank-two leaf frames,
query width sixteen, and an inverse-optimal completion containing the
full scalar fitted target.  There is exactly one original nonzero
residual source, of squared norm $\beta$; all sibling targets are
arbitrary fixed vectors already represented in their original leaf
frames.  The tree may have unequal-depth siblings.  For this rule,
\begin{equation}
 \E\bigl[\lVert E_{\mathrm{root}}\rVert^2
                 +\lVert c_{\mathrm{root}}\rVert^2\bigr]
 \leq\frac{308110}{23023}\,\beta<13.384\,\beta.
 \label{upd:eq:route-single-source}
\end{equation}
Here $E_{\mathrm{root}}$ is the physical component of the target outside
the retained space and $c_{\mathrm{root}}$ is its coefficient fit error
inside that space.  These orthogonal components together price the
physical fitted-target error.  The zero-budget case is exact.

The mechanism is a weighted inverse contraction along the unique
residual-bearing path.  With latent variance $\tau$, inverse
$V=\tr K^{-1}$, and mandatory fitted-direction inverse
$\phi=\lVert\alpha\rVert^2/(\alpha^TK\alpha)$, let
$H_t=\E[\tau_t\phi_t]$ and $A_t=\E[\tau_tV_t]$.
The clean sibling has an exact fixed-target Gaussian query response;
its radius is not replaced by the radius of a noisy fitted response.
Integrating this Gaussian response and its conditional uniform optional
line gives
\[
 \E\phi_{\mathrm{new}}\leq\phi+\frac{38}{297}V
                         +\frac{13}{11}v_b,
 \qquad v_b\leq\frac{11}{49}.
\]
The beta variance factor is $6/7$, and
$A_t\leq H_t+(11/49)\E\tau_t$.
Thus the weighted recurrence contracts with
\[
 \frac67\left(1+\frac{38}{297}\right)
 =\frac{670}{693}<1.
\]
Summing the reviewed energy recurrence gives
\eqref{upd:eq:route-single-source}.  The exact integrations and audited
constants are in \sourcefile{fitted_target_clean_siblings.md} and
\sourcefile{fitted_target_constants.json}.

This theorem is not superseded by the large-rank theorem: it treats a
small-rank, one-source regime with a concrete constant and a different
completion rule.  It is also not a decomposition of a nonlinear
many-source run into independent one-source runs.  The clean sibling's
Gaussian radius is the special input which prevents that inference.

\subsection{Source averaging preserves second moments within a factor}

For the scalar source model with pretemplate width four, retained rank
two, width sixteen, and inherited-core-plus-full-fit completion, let
$\mathcal Q$ contain the retained query data of the original sources
and their deterministic query transfers.  The source physical dual
lift $A_\alpha$ has metric $M_\alpha=A_\alpha^TA_\alpha$ and target
cross vector $g_\alpha=A_\alpha^TU_\alpha$.  Averaging the augmented
source Gram
\[
 \begin{pmatrix}
 M_\alpha&g_\alpha\\g_\alpha^T&\lVert U_\alpha\rVert^2
 \end{pmatrix}
\]
conditionally on its actual normalized query state and transporting it
by the query-determined congruences gives exact conditional means of
both the inverse and squared fitted error.

In particular, if $V$ is the actual inverse trace and
$Z=\lVert F-U\rVert^2$, write
$\widehat V=\E[V\mid\mathcal Q]$ and
$\widehat Z=\E[Z\mid\mathcal Q]$.  At every finite height,
\begin{align}
 \widehat V^2&\leq\E[V^2\mid\mathcal Q]
                         \leq11520\widehat V^2,\\
 \widehat Z^2&\leq\E[Z^2\mid\mathcal Q]
                         \leq11520\widehat Z^2.
 \label{upd:eq:route-source-comparison}
\end{align}
The factor is paid once over the original sources, not once per merge.
Extended expectations are allowed; these comparisons alone do not
bound growth with height or establish finiteness.

The proof controls affine physical errors, including their target
cross terms.  At a positive-omission source, conditional on its observed
core coordinate and moving query radius, the physical vector is
separately affine in three independent blocks: the inverse core radius,
the represented tail fit, and a standard Gaussian block.  Their uniform
affine fourth-to-second-moment constants are respectively
$256$, $9+4\sqrt2<15$, and $3$.  Tensorization yields
$768(9+4\sqrt2)<11520$.  Exact-source endpoints are checked directly;
the zero-moving-fit branch uses its fixed-prefix chi-square degrees
and is not obtained by a singular positive-noise limit.  Conditional
independence of the original sources and positive-semidefinite
transport extend the source estimate to the whole tree without
multiplying the constant at every level.

Thus
$\|\widehat V\|_2\leq\|V\|_2\leq48\sqrt5\|\widehat V\|_2$,
and likewise for $Z$.  These statements and their independent review
are in \sourcefile{conditional_source_metric_comparison.md} and
\sourcefile{conditional_source_metric_comparison_audit.md}; the
initialization and exact error transport are recorded in
\sourcefile{represented_source_query_metric.md} and
\sourcefile{scalar_query_error_transport.md}.  This result is a rigorous
bridge between actual and conditionally averaged scalar models.  It
does not replace the missing arbitrary-HSS source interface.

\chapter{Fresh innovations and bounded-stage correction}
\label{upd:chap:route-fresh}

\section{Freshness is a physical filtration condition}

To add a fresh internal target $Z_I$, it is sufficient to reserve a
deterministic physical subspace which has not contributed to any
earlier query response, target response, or selection.  More precisely,
let $D_{\mathrm{old}}$ contain all physical directions exposed by the
entire previous computation, assume the old record is measurable in
$P_{D_{\mathrm{old}}}\Psi$, and require
\[
 \range Z_I\perp D_{\mathrm{old}},\qquad
 \range W_I\subseteq D_{\mathrm{old}}.
\]
Then $\Psi^TZ_I$, conditional on the complete old record, is centered
matrix Gaussian with column covariance $Z_I^TZ_I$ and independent
query coordinates.  Node-private orthogonal reserved spaces are one
way to enforce this condition.  Merely calling a response a new source
does not enforce it: ordinary parent physical coordinates have already
been queried by their children.

The exact law, represented-pool injections, and the age-weighted
covariance accounting are in
\sourcefile{guarded_internal_injection_route.md} and
\sourcefile{guarded_internal_injection_audit.md}.  A source born at the
current node has age zero.  One cannot attach the decay factor of an
old leaf source to such an injection simply because the current node
has large height.

\section{Pure fresh resets have uniform moments}

Consider the fixed-guard model with $q=8r$, target width at most $r/2$,
and pure resets
\[
 A_a=A_b=0,\qquad U_I=Z_I,\qquad C_I=Z_I^TZ_I.
\]
The source targets are likewise perpendicular to their fixed source
pretemplates, or zero.  The new $Z_I$ satisfies the reserved-space
freshness condition.  Earlier learned child spaces remain in use; only
the current target is reset.

Let $h=\rank C_I$, so the deterministic guard core has width $j=r-h$.
Condition on the full state just before injection, including the
current width-$2r$ predesign.  Take an orthonormal query basis beginning
with the core image and write its physical preimages as
$[A_{\mathrm{core}},B_{\mathrm{pre}}]$.  There are $t=r+h$
complementary query coordinates.  With
\[
 K=B_{\mathrm{pre}}^TB_{\mathrm{pre}},\quad
 V_{\mathrm{core}}=\lVert A_{\mathrm{core}}\rVert_F^2,
 \quad \tr K=V_{\mathrm{pre}}-V_{\mathrm{core}},
\]
the whitened fresh response supplies a standard $t\times h$ Gaussian
matrix $X$.  Its column projector is independent of the old state, and
full-fit retention plus $j+h=r$ gives the exact identity
\begin{equation}
 V_I=V_{\mathrm{core}}+\tr(KP_X).
 \label{upd:eq:route-reset-projector}
\end{equation}
Hence, with $\alpha=h/(r+h)\leq1/3$,
\begin{equation}
 \E[V_I\mid\text{pre-injection state}]
 =(1-\alpha)V_{\mathrm{core}}+\alpha V_{\mathrm{pre}}.
 \label{upd:eq:route-reset-contraction}
\end{equation}
Whitening is only on the support of $C_I$; the zero-covariance case is
the fixed-core branch.

The fixed physical core has
$\E V_{\mathrm{core}}=(r-h)/(7r+h-1)$.  The independent child
relative-frame law gives
\[
 \E V_{\mathrm{pre}}
 =\frac{7r-1}{6r-1}(\E V_a+\E V_b),
 \qquad\frac{7r-1}{6r-1}\leq\frac{13}{11}\quad(r\geq2).
\]
Substitution in \eqref{upd:eq:route-reset-contraction} closes the induction
$\E V_I\leq1/4$.  The physical fit error includes the omitted target
orthogonally, so
\[
 \E[\lVert(F_I-Z_I)v\rVert^2\mid\text{pre-injection state}]
 =(v^TC_Iv)(1+V_{\mathrm{pre}}).
\]
Consequently, for all heights,
\begin{equation}
 \E V_I\leq\frac14,\qquad
 \E\lVert(F_I-Z_I)v\rVert^2\leq\frac{35}{22}v^TC_Iv
 \quad(r\geq2).
 \label{upd:eq:route-reset-mean}
\end{equation}
For $r\geq8$, the second moment of a diagonal entry of a Haar rank-$h$
projector is $h(h+2)/(t(t+2))$.  Minkowski, the fixed-core inverse
moment, and the relative-frame $L^2$ bound similarly yield
\begin{equation}
 \|V_I\|_2\leq2,\qquad
 \|\lVert(F_I-Z_I)v\rVert^2\|_2
 \leq(1+5\sqrt3)v^TC_Iv.
 \label{upd:eq:route-reset-second}
\end{equation}
These are uniform theorems for fresh or zero resets, reviewed in
\sourcefile{fresh_injection_moment_closure.md}.  They do not establish
the same convex contraction when a history-dependent inherited target
is present.

\section{What remains true with a noncentral fitted mean}

Let $X=M+G\in\R^{t\times s}$, where $M$ is fixed, $G$ is standard
Gaussian, and $s<t$.  Put $\mathcal U=\range M$ and
$p=\dim\mathcal U$.  Rotations fixing $\mathcal U$ give
\[
 \E P_X=A_{\mathcal U}\oplus bI_{\mathcal U^\perp},
 \qquad0\preceq A_{\mathcal U}\preceq I,
 \qquad\tr A_{\mathcal U}+(t-p)b=s.
\]
The exact useful bounds are
\begin{equation}
 \frac{s-p}{t-p}\leq b\leq\frac st,
 \qquad
 \tr(K\E P_X)
 \leq\frac st\tr K+\left(1-\frac st\right)\tr(KP_{\mathcal U})
 \quad(K\succeq0).
 \label{upd:eq:route-noncentral}
\end{equation}
There is no general lower Loewner bound $(s/t)I$ on the mean subspace:
a weak mean direction can have leverage below that value when another
mean direction is strong.

For the upper bound on $b$, fix a unit vector perpendicular to the
mean.  Its Gaussian row has squared norm $R\sim\chi_s^2$.  Conditional
on that row, rotate the column coordinates and apply a Schur complement.
The leverage is $R/(R+D)$, where conditionally $D$ is noncentral
chi-square with $t-s$ degrees of freedom.  Its Laplace transform is at
most the central one.  Integrating the reciprocal denominator bounds
the leverage mean by the central value $s/t$.  The lower bound follows
from the trace identity and $A_{\mathcal U}\preceq I$; the metric
inequality follows by also using $A_{\mathcal U}\preceq I$.

For singular noise covariance, first retain its deterministic mean-only
range of dimension $\tau$, then work in its orthogonal quotient.
If the noise has rank $\ell$, the quotient projector has $\ell$ noisy
columns in dimension $t-\tau$.  Thus its central coefficient is
$\ell/(t-\tau)$.  The original fit has full column rank exactly when
the mean-only and noisy ranks add to the required column rank.  One
must include the forced mean-only projector and any otherwise unexcited
completion; they cannot be discarded by pretending the covariance is
invertible.

Applied conditionally to a valid fresh guarded merge with inherited
complementary mean $M$, this gives
\begin{equation}
 \E[V_I\mid\text{old state}]
 \leq(1-\alpha)V_{\mathrm{mean}}+\alpha V_{\mathrm{pre}},
 \quad
 V_{\mathrm{mean}}=V_{\mathrm{core}}+\tr(KP_{\range M}),
 \quad\alpha\leq\frac13,
 \label{upd:eq:route-mixed-bridge}
\end{equation}
under the stated compact full-fit/core rank condition.  The singular
full-rank version improves $\alpha$ to $\ell/(r+\ell)$.  The new term
$V_{\mathrm{mean}}$ is coupled to the inherited physical metric.  It is
not bounded by the pure-reset theorem or by independence of the fresh
part alone.  Equation \eqref{upd:eq:route-mixed-bridge} is therefore a
one-step bridge, not a completed general mixed-injection theorem.
The exact projector statements and audit are in
\sourcefile{noncentral_projector_route.md} and
\sourcefile{noncentral_projector_audit.md}.

\section{An optional correction with only two new query batches}

There is a separate, elementary bounded-stage module which can help
when a pilot has a small \emph{global low-rank residual tail}.  It
should not be confused with a general HSS pilot-quality theorem.
Fix any known pilot $M_0$, including its complete construction record,
and put $R=A-M_0$.  Draw a fresh $N\times\ell$ Gaussian matrix
$\Omega_c$, where $\ell=k+p$ and $p\geq2$, and query
\[
 Y_c=A\Omega_c-M_0\Omega_c=R\Omega_c.
\]
Let $Q$ span its range, of dimension $d\leq\ell$.  Freeze $Q$, then
obtain the exact transpose products $A^TQ$ and form
\[
 K=R^TQ=A^TQ-M_0^TQ,
 \qquad M_1=M_0+QK^T.
\]
The second batch consists of exact adapted queries.  It is not a claim
that previously recycled noisy compressed products are independent
fixed-input products.  There are $\ell+d\leq2(k+p)$ new products.

If $\tau_k(R)$ denotes the squared error of the best \emph{global}
rank-$k$ approximation, the Gaussian range-finder calculation gives
\begin{equation}
 \E[\lVert A-M_1\rVert_F^2\mid M_0]
 \leq\left(1+\frac{k}{p-1}\right)\tau_k(A-M_0).
 \label{upd:eq:route-two-batch}
\end{equation}
Indeed, decompose the fresh Gaussian matrix in an SVD of $R$ into
top-$k$ and trailing blocks $\Omega_1,\Omega_2$.  A candidate in the
sampled range gives residual bound
$\lVert\Sigma_2\rVert_F^2+
 \lVert\Sigma_2\Omega_2\Omega_1^\dagger\rVert_F^2$.
The blocks are independent and
$\E\lVert\Omega_1^\dagger\rVert_F^2=k/(p-1)$, proving the claim.

Taking $p=k+1$ uses at most $4k+2$ new products and gives factor two.
It is a contraction by $\rho$ if
$\tau_k(A-M_0)\leq(\rho/2)\lVert A-M_0\rVert_F^2$.
It gives the desired error order if a separate pilot theorem proves
$\E\tau_k(A-M_0)\leq\beta L\OPT_k(A)$.
Neither hypothesis follows from HSS structure alone.  For example,
many disjoint unit rank-one sibling blocks form an HSS-rank-one matrix
with many global unit singular values.  A zero pilot can have zero HSS
comparator error but a large global rank-$k$ residual tail.

If $M_0$ is HSS-rank-$r_0$, the global rank-$\ell$ correction has an
HSS representation of rank at most $r_0+\ell$.  It supplies neither a
free same-rank output nor an efficient best-HSS projection.  The module
and its exact hypotheses are recorded in
\sourcefile{bounded_stage_compression_route.md}.  It is an optional
conditional repair, rather than a replacement for the synchronized
learner's missing arbitrary-input physical budget.

\chapter{Current numerical evidence and the stopping point}
\label{upd:chap:route-evidence}

\section{A reproducible test of the actual synchronized recursion}

The current experiment implements the completed synchronized learner
of Chapter~\ref{upd:chap:route-synchronized}, using the actual arrays
$A\Omega,A^T\Psi$, common child unions, the full-raw rank-$\leq R$
branch, and one joint visible/kernel completion.  All node responses
reuse the original probes.  The full learned coefficient-space
projection is computed densely as a diagnostic, with exact physical
leaf diagonal blocks.  No independent random refit and no final
rank-$k$ reduction are simulated.

The saved run has \textbf{624 cases from 80 independent probe pairs}.
The main grid uses 16 pairs at each of
\[
 (N,k)=(64,1),(128,1),(256,1),(128,2),
 \qquad L=3,4,5,3,
\]
for 64 independent pairs.  Each pair is reused across three exact
families and two noisy patterns at three noise scales.  This yields
192 exact controls and 384 noisy cases.  A separate set of 16 pairs,
each used at three scales, gives 48 clean-parent witness cases.
Cases sharing probes are paired observations, not additional
independent samples.  Within a fixed cell, its mean and standard error
use the 16 independent complete trees.

The script, raw records, report, and independent audit are
\sourcefile{synchronized_union_numerical_check.py},
\sourcefile{synchronized_union_numerical_check.json},
\sourcefile{synchronized_union_numerical_check.md}, and
\sourcefile{synchronized_union_numerical_audit.md}.  The recorded
invocation uses 16 seeds per configuration.  The audit reconstructed
216 estimate blocks across 39 cells from the saved records and
independently checked six small dense coefficient projections; it did
not rerun the main grid.

\subsection{Exact controls and a clean-parent witness}

The exact controls include nested HSS-rank-$k$ inputs, nested
HSS-rank-$2k$ inputs testing rank slack, and deficient or zero cuts.
The largest full coefficient-space squared error was
$3.16200589\times10^{-23}$, and the largest normalized true-cut loss
was $1.71934437\times10^{-22}$.  The rank-$2k$ controls test
Theorem~\ref{upd:thm:route-exact}; they are not declared to have zero
nominal rank-$k$ approximation error.  All these matrices have zero
physical-leaf diagonal blocks, so no arbitrary dense coarse-leaf block
is being declared compatible with a finer rank bound.

The witness uses $N=128,k=1$ and a parent $I=[0,64)$, with
\[
 u=(e_0+e_{16})/\sqrt2,\quad p_1=e_1,\quad p_2=e_{17},
 \quad b_1=e_{32},\quad b_2=e_{33},\quad f=e_{64},
\]
and $A=B+\epsilon E$ for
$B=uf^T$, $E=p_1b_1^T+p_2b_2^T$,
$\epsilon\in\{.01,.1,1\}$.  This embeds the fixed matrix mechanism
from the separate-view reachability example in the new learner's
actual feedback recursion.  It does not reuse or assume that earlier
example's reached child frames or query values.

In all 48 cases the physical design $\Psi_I$ was $64\times32$ of rank
32, and therefore noninjective on the full physical space.  Each common
pretemplate design was $8\times32$ of rank eight.  The largest
child-template loss of the unit target was
$3.84603468\times10^{-29}$ and the largest parent loss was
$4.49944901\times10^{-29}$.  The current residual commutator is
structurally zero; its direct floating-point evaluation had norm at
most $5.14307095\times10^{-16}$.  These observations check the
clean-parent mechanism at reached states, rather than replacing the
compressed injectivity condition by full physical injectivity.

\subsection{Fixed noisy inputs and diagnostic limits}

The two noisy patterns have fixed HSS-rank-$k$ comparators and balanced
perturbation energy across split levels.  One uses dense sibling
blocks; the other reuses fixed global factors and level couplings with
additional alternating rank-one components.  Perturbations are
normalized to Frobenius norm one before multiplication by
$\epsilon\in\{.001,.03,.3\}$.  Matrix-construction seeds are distinct
from probe seeds.  The denominator is the specified
$\lVert A-B\rVert_F^2$, \emph{not} a computed value of $\OPT_k(A)$.

The largest cell mean of full coefficient-space squared error divided
by this residual budget was
\[
 1.00792492\ \pm\ 0.00335957
\]
for $N=256,k=1$, the dense pattern, and $\epsilon=.3$; the uncertainty
is one standard error.  The largest individual ratio was
$1.02872904$.  The largest cell mean after additional normalization by
$L$ was $0.30198110\pm0.00073533$, and the largest individual
depth-normalized ratio was $0.30789474$.

The high-rank internal union branch was exercised 10,752 times across
the two orientations.  There were 6,720 full-raw internal selections,
7,680 full-range leaf selections, 12,288 truncated leaf selections,
and 1,248 uncompressed root selections.  No design-rank failures were
observed.  The largest observed design condition number was 8.60370.
No sample-SVD singular value lay within a factor 100 of the declared
sample cutoff.  This last diagnostic does not audit every singular
value used by the separate lifted-union and coordinate-completion
tolerances; those tolerances are explicitly recorded in the report.

These finite observations support exact/slack capture and the
clean-parent repair, and show that the noisy high-rank branch was
actually used.  They are not an all-depth inverse-moment estimate, an
all-input $O(L)$ stability theorem, or an OPT-relative rank-$k$ output
guarantee.  Small finite-sample condition numbers do not establish a
tail bound on the adaptive inverse.

\section{A signed-identity check on already sampled cases}

The separate identity diagnostic uses 12 cases already in the main
grid: three configurations, two noise patterns, and two scales, all
at the first prescribed seed.  It reuses \emph{three} probe pairs and
does not create twelve new independent trees.  At 400 reached oriented
internal nodes, the largest normalized vector residual in
\eqref{upd:eq:route-signed} was $2.72093\times10^{-15}$, and the largest
residual in the squared orthogonal split
\eqref{upd:eq:route-signed-energy} was $1.71304\times10^{-17}$.
The largest separator residual was $3.05673\times10^{-15}$.

For nonnegligible source and return terms, the realized fraction
\[
 \frac{2\langle Z_I^{\mathrm{src}},D_I^{\mathrm{ret}}\rangle_F}
 {\lVert Z_I^{\mathrm{src}}\rVert_F^2+
  \lVert D_I^{\mathrm{ret}}\rVert_F^2}
\]
was positive at 263 of 400 nodes, with extrema $0.928404$ and
$-0.912101$.  Thus the reached trajectories exhibit both cancellation
and amplification.  The check validates the implemented algebra and
the need to retain the cross term; it does not estimate a universal
expectation of that cross term.  The records are
\sourcefile{synchronized_signed_identity_check.py},
\sourcefile{synchronized_signed_identity_check.json}, and
\sourcefile{synchronized_signed_identity_check.md}, reviewed in the
same synchronized numerical audit.

\section{Unreviewed calculations at the stopping point}
\label{upd:sec:route-unreviewed}

The following observations were communicated while work was being
stopped.  They had not been saved as independently reviewed new notes.
They are segregated here to preserve the state of the investigation,
and are not used in any theorem or constant above.

\subsection{A proposed global discard accounting}

For nested physical projectors, let the proposed row discard projector
at an internal nonroot node be
\[
 D_I^U=\diag(P_a^U,P_b^U)-P_I^U,
\]
and at a physical leaf use $D_I^U=I-P_I^U$.  Define $D_I^V$ similarly.
All projectors act in their indicated node coordinates.  With the
uncompressed root convention, the working quantities were
\[
 \mathcal N_U=\sum_{I\ne\mathrm{root}}
       \lVert D_I^U B[I,I^c]\rVert_F^2,
 \qquad
 \mathcal N_V=\sum_{I\ne\mathrm{root}}
       \lVert B[I^c,I]D_I^V\rVert_F^2.
\]
The communicated candidate identity was
\begin{equation}
 \operatorname{dist}_F(B,\mathcal S)^2
 \mathrel{\stackrel{\text{unreviewed}}{=}}
 \mathcal N_U+\mathcal N_V
 -\sum_{\substack{I,J\text{ ordered}\\I,J\text{ incomparable}}}
       \lVert D_I^U B[I,J]D_J^V\rVert_F^2,
 \label{upd:eq:route-proposed-discard}
\end{equation}
with proposed consequences
$\max(\mathcal N_U,\mathcal N_V)
\leq\operatorname{dist}_F(B,\mathcal S)^2
\leq\mathcal N_U+\mathcal N_V$.
Its intended purpose was to charge each new physical discard once,
rather than summing full inherited cut losses repeatedly.  The exact
partition and orthogonality accounting, including endpoints, still
required a recorded audit at the stopping point.  In particular,
\eqref{upd:eq:route-proposed-discard} is not being used to strengthen the
proved sufficient interface \eqref{upd:eq:route-global-cut-interface}.

\subsection{Retained design rank versus stacked design rank}

A second communicated observation concerns standard full $8k$ leaves.
Their Gaussian physical designs have full row rank almost surely, so
any retained $R$-dimensional leaf space has an injective design.  At a
parent, the stacked predesign contains a child design of rank $R$ and
therefore has rank at least $R$.  Every mandatory lifted direction is
visible.  Scanning visible candidates before kernel candidates can
then complete the retained space to rank $R$ without using the kernel.
This suggests an induction of retained-design injectivity for every
fixed input, even if the stacked width-$2R$ predesign is singular.

This short rank argument was not independently audited at the stop.
It must not be confused with the reviewed exact-family theorem,
which proves full stacked rank, or with a lower bound on singular
values.  Retained rank $R$ alone does not eliminate a comparator
component invisible to a singular parent predesign.  The associated
bottom-leaf union-rank calculation was not available in a sufficiently
recorded and reviewed form to include as a separate result.

\section{Proof-status ledger}

The surviving chain has a definite completed part and definite missing
interfaces.  The independent fixed-space refit gives its uniform
constant $1607$.  The completed synchronized learner is explicit and
has an exact rank-slack capture theorem, a clean-parent statement, a
fixed-comparator separator, the $24qL$ unlifted source budget, and the
projected signed physical identity.  Its saved numerical checks test
the actual shared-probe recursion with the stated finite scope.

The independent-source theorems prove genuine uniform local stability
under their own assumptions: fixed nested physical cores at small
rank, the one-residual-source scalar path, conditional source averaging
with factor 11520, sufficiently large-rank fixed sources, fixed-guard
compatible transfers, and pure fresh resets.  The noncentral projector
result supplies a valid mixed one-step inequality with an explicit
inherited metric cost.  None identifies arbitrary HSS residual history
with the required independent source interface.

For the current synchronized route, the open mathematical requirement
is a joint physical source/return/invisible budget sufficient for
\eqref{upd:eq:route-needed-space}.  Separately, the full target requires an
efficient same-rank reduction on the original target tree.  The
query-only best-HSS projection argument does not settle that
computational requirement.  The observations in
Section~\ref{upd:sec:route-unreviewed} remain provisional records rather
than additions to the proved chain.

% ===== END FROZEN FOLLOW-UP RECORD =====

\appendix

% ===== ROUTELEDGER =====
\begingroup
\let\R\relax
\newcommand{\R}{\mathbb R}
\let\E\relax
\newcommand{\E}{\mathbb E}
\let\Prb\relax
\newcommand{\Prb}{\mathbb P}
\let\T\relax
\newcommand{\T}{\mathsf T}
\let\F\relax
\newcommand{\F}[1]{\left\lVert#1\right\rVert_{F}}
\let\N\relax
\newcommand{\N}[1]{\left\lVert#1\right\rVert}
\let\ip\relax
\newcommand{\ip}[2]{\left\langle#1,#2\right\rangle_F}
\let\tr\relax
\newcommand{\tr}{\operatorname{tr}}
\let\rank\relax
\newcommand{\rank}{\operatorname{rank}}
\let\diag\relax
\newcommand{\diag}{\operatorname{blockdiag}}
\let\range\relax
\newcommand{\range}{\operatorname{range}}
\let\OPT\relax
\newcommand{\OPT}{\operatorname{OPT}}
\let\HSS\relax
\newcommand{\HSS}{\operatorname{HSS}}
\let\diver\relax
\newcommand{\diver}{\operatorname{div}}
\let\Var\relax
\newcommand{\Var}{\operatorname{Var}}
\let\UC\relax
\newcommand{\UC}{\mathcal U}
\let\VC\relax
\newcommand{\VC}{\mathcal V}
\let\err\relax
\newcommand{\err}{\mathcal E}
\let\proofstep\relax
\newcommand{\proofstep}[1]{\par\addvspace{.45\baselineskip}\Needspace{3\baselineskip}\textbf{#1}\par\nobreak}
\let\op\relax
\newcommand{\op}[1]{\left\lVert#1\right\rVert_{\mathrm{op}}}
\let\Sn\relax
\newcommand{\Sn}[2]{\left\lVert#1\right\rVert_{S_{#2}}}
\let\Cov\relax
\newcommand{\Cov}{\mathsf C}
\let\Loss\relax
\newcommand{\Loss}{\mathsf L}
\let\Fil\relax
\newcommand{\Fil}{\mathcal F}
\let\Id\relax
\newcommand{\Id}{I}
\let\supp\relax
\newcommand{\supp}{\operatorname{supp}}
\let\vecop\relax
\newcommand{\vecop}{\operatorname{vec}}
\let\Gammafun\relax
\newcommand{\Gammafun}{\operatorname{Gamma}}
\chapter{Earlier certificates, attempted routes, and obstruction ledger}\label{ch:routeledger}
\begin{scopebox}
\textbf{Update, 20 September 2026.} This ledger records the original 66 routes only. The subsequent full-range failure, actual separate-view witness, source-model refinements, and synchronized learner are documented in Chapters~\ref{upd:chap:later-obstructions} and~\ref{upd:chap:current-route}. Its earlier query counts and maintained-algorithm names are historical; they are not the current synchronized specification. The entries calling uniform fixed-space refitting unresolved are already superseded by Chapters~\ref{ch:feedback}--\ref{ch:uniform}.
\end{scopebox}

\begin{scopebox}
This appendix preserves the preliminary route and obstruction ledger. Entries calling the fixed-space predictable interaction unresolved record an earlier proof stage; Chapters~\ref{ch:feedback}--\ref{ch:uniform} settle that interaction. The end-to-end basis-learning premise remains unproved, and the actual annihilate-first basis bound is disproved in Chapter~\ref{ch:anchor}.
\end{scopebox}
\section{Query counts, comparators, and the result hierarchy}\label{base:new:ledger}
\begin{longtable}{@{}>{\RaggedRight\arraybackslash}p{.29\linewidth}>{\RaggedRight\arraybackslash}p{.18\linewidth}>{\RaggedRight\arraybackslash}p{.46\linewidth}@{}}
\caption{Results are not interchangeable across comparators or access models.}\label{base:tab:paper-comparison}\\
\toprule Route&Queries&Scope and conclusion\\\midrule\endfirsthead
\toprule Route&Queries&Scope and conclusion\\\midrule\endhead
\bottomrule\endfoot
Published fresh-level approximation&$O(kL)$&Best-HSS squared-error factor $O(L)$; public benchmark, not proved anew here.\\
Raw reuse with exact core&$18k$&Exact fixed HSS inputs almost surely under canonical completion. General first two levels only: $47$, normalized noise $24$, combined error at most $33000$ times $\OPT$.\\
Raw basis acquisition only&$16k$&Produces a fixed random nested space; its general all-depth best-HSS quality remains unresolved.\\
Three candidates plus validation&$54k+8\le62k$&Conditional on an actual-law weak tail, expected error at most $15K^2L\OPT$. Validation does not prove the hypothesis.\\
Inverse-free local refit&$4k$ refit&Width $2k$; factor $3\cdot2^L-1$ relative to the fixed space. Total with training $20k$.\\
Sibling-orthogonal refit&$8k$ refit&Width $4k$; unchanged algorithm throughout the principal continuation. Total with raw training $24k$.\\
Selector family&$8k$ refit&All rank/depth factor $10$; restricted bases.\\
Early mixing path profile&$36k$ refit&Width $18k$; balanced fixed bases factor $70$. Superseded in this geometry by later narrower-probe storage.\\
General one-stream covariance route&$8k$ refit&Arbitrary fixed bases: factor $8L+6$. Includes coefficient output.\\
General response-tail/embedding route&$8k$ refit&Factor $C[1+\log\min\{k,L+1\}]^2$ relative to the fixed space, with the full restart proof now included. Elementary historical alternative below $255k$.\\
Rank-oversampled general refit&$O(k(1+\log k)^2)$ refit&Historical factor at most $132$; superseded in query order by Theorem~\ref{base:con:main-eight}.\\
Matched scalar child weights&$8k$ refit&Factor $13$, arbitrary orthogonal directions and arbitrary weights in $[0,1]$.\\
Comparable leakage plus $b$ exceptions&$8k$ refit&Factor $5+8\Lambda+4b$; total exception count bounded.\\
Arbitrary levels separated by balanced levels&$8k$ refit&Factor $33$; number of repetitions unbounded. Longer maximum runs have the stated $\alpha_b$ factor.\\
Crossed directional selections&$8k$ refit&Factor $121/9$; no balanced nodes needed. Perturb every child block by at most $1/50$: factor $197/9$.\\
Finite scalar temporal certificates&$8k$ refit&Historical exact-rational finite-family certificates; not an unrestricted or all-depth theorem.\\
Near-protected spectral certificates&$8k$ refit&Source-loss-sensitive bounds, including rank/$\log$ profiles and chi-square confidence scale; no new unrestricted asymptotic order.\\
Sharp covariance comparison&No extra queries&$\eta\Loss\preceq\Cov\preceq\Loss$ for the original source; $\eta=1/4$ at width $4k$, sharp. Not itself an error guarantee.\\
Matrix source-Gram route&$8k$ refit&Arbitrary bases; factor $<44$ for isotropic source row energy, automatic at rank one. General explicit factor $5+4\sqrt{\kappa(k+2)}$. Section~\ref{base:con:gram}.\\
Unmatched scalar child weights&$8k$ refit&Factor $<44$ for every residual; no matching of row and column weights. Repeated-coordinate geometry has factor controlled by its non-repeated size. Section~\ref{base:con:symmetry-multiplicity}.\\
Symmetric common-probe restart&One width-$q$ common probe&Adapted source-output factor $2\rho/(1-\rho)$ for scalar child Grams. Dilation costs $16k$ original queries, but its unrestricted stability is not proved. Section~\ref{base:con:symmetric}.\\
Generic global model fitting&Two Gaussian sketches, width $s\ge3k$&Saturated fixed inputs identifiable almost surely; only realization-wise local stability. Not a uniform expected approximation theorem. Section~\ref{base:con:identifiability}.\\
Quotient-variance general refit&$O(k[1+\log\min\{k,L+1\}])$ refit&Latest arbitrary-basis factor $8$. Explicit $q=(2+4a_L)k$, $a_L=1+\lceil\log_2(L+1)\rceil$. Theorem~\ref{base:con:main-eight}.\\
Full desired theorem&$O(k)$ total&Uniform $CL\OPT(A)$ remains unproved in this development.\\
\end{longtable}
An $O(L)$ learning bound and an $O(L)$ refitting bound multiply to $O(L^2)$, not the target. Likewise $d_{\mathcal S}(A)^2\ge\OPT(A)$ prevents replacing the former comparator by the latter in a fixed-space upper bound. Constants from different widths, source orientations, and error conventions are not directly comparable without this ledger.

\section{Every route considered in the development}\label{base:new:route-history}
The following sequence records the mathematical decisions, including routes that were abandoned or superseded. ``Failed'' below means failure of the indicated implication or proof strategy, not failure of the randomized approximation problem.
\begin{longtable}{@{}>{\RaggedRight\arraybackslash}p{.06\linewidth}>{\RaggedRight\arraybackslash}p{.29\linewidth}>{\RaggedRight\arraybackslash}p{.57\linewidth}@{}}
\toprule No.&Route&Outcome, reason for changing course, and location\\\midrule\endfirsthead
\toprule No.&Route&Outcome, reason for changing course, and location\\\midrule\endhead\bottomrule\endfoot
1&Raw Gaussian reuse&Specify nullification, discrepancy and exact root. Exact-input recovery survives adaptive learning with canonical completion; all-depth approximate input bound does not follow. Sections~\ref{base:sec:algorithm}--\ref{base:sec:exact-recovery}.\\
2&Residual-relative first levels&Same-level nullification/regression uses independent Gaussian coordinates; adaptive-left analysis permits one cross-level reuse. No induction to arbitrary depth. Sections~\ref{base:sec:finite}--\ref{base:sec:second}.\\
3&Active/inactive error separation&Derive exact Pythagoras and active-return balance; inactive error is reconstruction cost, not future forcing. Section~\ref{base:sec:global}.\\
4&Autonomous quotient&Innovations and mismatch determine basis selection; common coefficient still needs a reconstruction register. Sections~\ref{base:sec:quotient}--\ref{base:sec:quotient-moments}.\\
5&Generic adaptive inverse moments&Nestedness alone admits arbitrarily small compressed singular values and divergent inverse moments. Section~\ref{base:sec:inverse}.\\
6&Gaussian integration by parts&Fresh regression cost plus explicit adaptation derivative; no universal favorable derivative sign and no automatic zero-ridge limit. Section~\ref{base:sec:closure} and preceding Stein calculations.\\
7&Fresh-between-levels benchmark&Valid comparison with linear best-HSS error, but uses new queries at each level and does not prove reuse. Section~\ref{base:sec:second}.\\
8&Fractional moments, weak tails, validation&Independent candidates and validation can amplify weak actual-law control; the needed all-depth raw tail remains missing. Section~\ref{base:sec:validation}.\\
9&Uniform fitting and clipping&Training sketch fit cannot uniformly certify true residuals; candidate clipping or weak fixed-probability success alone does not control expected error. Section~\ref{base:sec:additional-obstructions}.\\
10&Wishart at intermediate operations&Exact conditional regression module and lower-edge tools are useful only after identifying the original conditional law. Random covariance mixtures cannot be replaced by their mean. Section~\ref{base:new:wishart}.\\
11&Fixed-basis orthogonal refit&Remove inverses and recover the entire fixed space exactly, at $4k$ or $8k$ refit cost. Initial all-depth bound is exponential. Section~\ref{base:sec:ortho-refit}.\\
12&Direct global least squares&Injective measurement map and the same coarse bound are proved. No depth-independent whole-space lower singular value was found. Section~\ref{base:new:wishart}.\\
13&F$^+$ finite-profile transfer&Conditional centered rank-one theorem is valid; HSS does not satisfy its hypotheses merely by being isotropic. Section~\ref{base:new:fp}.\\
14&Signed quotient energy&Compression dissipates, merge leaves only one signed correlation. No scalar weight contracts both maps on all admissible states. Section~\ref{base:new:energy}.\\
15&Output observation ledger&Include emitted coefficients and root; quotient energy alone is insufficient. Proposition~\ref{base:new:ledger-interface}.\\
16&Coordinate message paths&No branching; integrate ordered reused-frame products to obtain factor $10$. Section~\ref{base:new:early}.\\
17&One-pass centered feedback operators&Factor $\rho$ for deterministic/independent inputs, but repeated correlated inputs cannot be substituted. Section~\ref{base:new:early}.\\
18&Absolute path majorization&Closes balanced mixing at larger width. Absolute ancestor sums can grow with depth despite squared Bessel control. Section~\ref{base:new:early}.\\
19&Initial-source orthogonality&Exact sign/conditioning argument removes cross-source terms, leaving same-source returns. Section~\ref{base:new:source}.\\
20&Freeze forward, integrate transpose&Exact covariance state yields linear depth; internal energy can grow while output cancels. Section~\ref{base:new:signed}.\\
21&Full signed covariance response&Retain the $F,C$ cross tensor and extract positive response $W(G)$. Exact finite closure, no uniform conditional gain. Section~\ref{base:new:signed}.\\
22&Uniform conditional gain&Explicit slow rotation/release amplifier disproves the universal frozen-forward target. Section~\ref{base:new:amplifier}.\\
23&Actual joint source decomposition&Separate the adapted exterior-sibling martingale from both interior and exterior discarded directions; denominators $q-2k$ and $q-3k$ differ. Section~\ref{base:new:source}.\\
24&Adapted restart/Carleson tail&Obtains bounded expected future response under legitimate prefix conditioning; the complete cross-tensor induction is now supplied in Lemma~\ref{base:con:restart-cone}. Section~\ref{base:new:carleson}.\\
25&Rank and time embeddings&General logarithmic bounds and finite source profiles; independent-source estimates alone are insufficient. Section~\ref{base:new:carleson}.\\
26&Generic isotropy as a cure&Abstract Haar--harmonic example defeats a dimension-free embedding based only on these coarse properties. Section~\ref{base:new:abstract-obstruction}.\\
27&Immediate temporal parity&Current cross-time term vanishes, future terms need not. Exact temporal diagonal and delayed-only form. Section~\ref{base:new:temporal}.\\
28&Initial-anchor integration&Positive innovation part has an absolute scalar-weight bound; retained part remains coupled. Section~\ref{base:new:temporal}.\\
29&Fourth-order marked propagation&Exact original-Haar finite state with two temporal marks; nonzero delayed entries and finite certificates. No uniform gain follows from finite degree. Section~\ref{base:new:fourth}.\\
30&Scalar output-plus-storage&Closes matched leakage and finite exceptions. A reachable rank-two example defeats every scalar weight. Section~\ref{base:new:scalar-storage}.\\
31&Backward matrix storage&Directional future weights neutralize that witness; finite generalized-eigenvalue certificate. Section~\ref{base:new:matrix-storage}.\\
32&Balanced resets and bounded runs&Infinitely many arbitrary levels can be absorbed when run lengths are controlled. Section~\ref{base:new:matrix-storage}.\\
33&Directional refresh without resets&Crossed noncommuting projector families and levelwise perturbations close with constant factors. Section~\ref{base:new:refresh}.\\
34&Transported observability&Window loss criterion; untransported sums can be positive while a direction survives forever. Section~\ref{base:new:refresh}.\\
35&Force refresh during learning&Exact localized HSS inputs require protected directions, ruling out an unconditional strict-refresh learning rule. Section~\ref{base:new:protected}.\\
36&Near-protected covariance suppression&Quantifies actual source excitation, controls dependent response, and rescales the conditional amplifier to vanishing joint average. Section~\ref{base:new:near}.\\
37&Whole-probe comparison laws&First scalar Gaussian, then dimension-aware chi-square envelope; source has Gaussian/gamma MGF bound, not equality or stochastic order. Sections~\ref{base:new:chi}--\ref{base:new:source-moments}.\\
38&Response square-root randomization&John--Nirenberg yields $p^2$ response moments. The auxiliary law does not replace the sketch itself. Section~\ref{base:new:response-moments}.\\
39&Optimized spectral metrics&One Schatten quasi-norm handles mixed loss directions; exact noncommuting covariance--loss optimization reduces to one scalar. Section~\ref{base:new:spectral}.\\
40&Exploit arbitrarily small exact covariance&Sharp coercivity rules this out relative to geometric loss: the possible metric improvement is at most four at width $4k$. Section~\ref{base:new:coercive}.\\
41&Full source-Gram fourth moment&Retain $\tr[(N^\T N)^2]$, not its Frobenius scalarization; Gaussian comparison closes at coefficient one. Section~\ref{base:con:gram}.\\
42&Effective row-rank source spread&Absolute refit constant for spread-out sources and arbitrary bases; coherent low-row-rank sources remain. Section~\ref{base:con:gram}.\\
43&Exceptional temporal eigenspace&Leading-eigenvalue sum bounds isolate at most $k$ bad directions; exact hard case is one left source row with retained rank still $k$. Section~\ref{base:con:gram}.\\
44&Auxiliary estimation of bad directions&Ideal Gaussian regression pays their trace, but the residual-source sketch is not available from $A$ without in-space contamination. No legal completed decoder. Section~\ref{base:con:gram}.\\
45&Explicit adapted covariance cone&Full cross tensor and deterministic mean transport repair the previously compressed restart argument. Lemma~\ref{base:con:restart-cone}.\\
46&Interior discarded-source repair&A nonleaf counterexample forces interior directions into the off-chain budget; the coarse allowance survives. Example~\ref{base:con:interior-example}.\\
47&Finite Bellman output registers&Five deterministic positive blocks, exact Stiefel degree reduction, and a passive-port backward map. Section~\ref{base:con:bellman}.\\
48&Marked Bellman telescope&Retain past output as well as the future value; only the odd linear fresh-frame kernel contributes to the mixed term. Section~\ref{base:con:bellman}.\\
49&Sharp local self-mark estimate&Exact output/state tradeoff reduces $17.5$ to $8.257804885\ldots$ at width $4k$; does not pay the mixed term. Equation~\eqref{base:con:selfmark}.\\
50&Local Haar--Gaussian defect&Only the exterior-square sector has positive excess; its current-mark placement avoids charging every empty stage. Section~\ref{base:con:marked-local}.\\
51&Mean forcing and complementary port&A summable mean component is controlled, but it overlaps with the Bellman partition and does not kill the other odd term. Section~\ref{base:con:marked-local}.\\
52&Symmetric dilation&Common probes give a different joint law and one observation state; auxiliary diagonal-copy coefficients must be propagated. Section~\ref{base:con:symmetric}.\\
53&Symmetric state cone&All-depth adapted state bound is proved; accumulated nonscalar output has an unpaid symmetric/skew commutator. Section~\ref{base:con:symmetric}.\\
54&Symmetric scalar-child output ledger&Commutator vanishes and all adapted output is paid; no anticipative-source substitution is justified. Equation~\eqref{base:con:sym-adapted}.\\
55&Saturated generic identifiability&Same observations force every cut space and core, including sketch-dependent competitors. Section~\ref{base:con:identifiability}.\\
56&Local nonlinear fitting and rank selection&Smooth local inverse has random constants; rank-deficient identity admits multiple completions, while lexicographic ranks give only an exact selector. Section~\ref{base:con:identifiability}.\\
57&Symmetry multiplicity&Repeated eigenvalues convert a Ky Fan bound into a coherent-source norm bound; unmatched scalar weights close. Section~\ref{base:con:symmetry-multiplicity}.\\
58&Coordinate-traceless remainder&Coordinate-averaged temporal response is uniformly controlled; the remaining positive anisotropic spectrum is not. Equation~\eqref{base:con:anisotropic}.\\
59&Complete first-moment passive filter&One deterministic contraction controls all mean input/readout effects, without separate random restarts. Section~\ref{base:con:predictable}.\\
60&Exact Doob product decomposition&Four terms retain the terminal mark on all past output; parity controls only the stated pairs. Section~\ref{base:con:product}.\\
61&Predictable-risk equivalence&Three terms are bounded and full risk differs from explicit nonnegative predictable energy by an absolute square-root allowance. Equation~\eqref{base:con:equivalence}.\\
62&Pathwise isolated-chain quotient&The quotient form decreases on canonical source coordinates; it does not pay deterministic common coefficients. Section~\ref{base:con:quotient}.\\
63&Width-sensitive unmarked output&Two-stream flow retains a factor $\rho$ lost by the coarser observation ledger. Section~\ref{base:con:theta}.\\
64&Five-register future variance&Remove deterministic common coefficients, then bound four pure sectors and induct on mismatch. Section~\ref{base:con:variance}.\\
65&One-logarithm general oversampling&Predictable energy costs $O(\rho^2b)$, giving factor $8$ at width $O(k\sqrt b)$; the uniform $O(k)$ endpoint is not proved. Section~\ref{base:con:oversampling}.\\
66&Running marked endpoint&At each deterministic time the actual marked state is bounded by $3.803\beta_n$ at width $4k$; no maximum or cumulative output bound follows. Equation~\eqref{base:con:state-current}.\\
\end{longtable}
Other ideas considered but not promoted to results were repeated independent residual correction, adding redundant retained dimensions to balance arbitrary bases, sketch-space gauge randomization, and a direct whole-space small-ball theorem. None supplied both exactness on the original same-rank HSS class and a depth-independent $O(k)$ query certificate. They are not alternative completed algorithms hidden behind the main decoder.

\section{Obstructions, separated by logical scope}\label{base:new:obstruction-ledger}
\begin{longtable}{@{}>{\RaggedRight\arraybackslash}p{.29\linewidth}>{\RaggedRight\arraybackslash}p{.31\linewidth}>{\RaggedRight\arraybackslash}p{.33\linewidth}@{}}
\toprule Tempting inference&Obstruction&What it does not establish\\\midrule\endfirsthead
\toprule Tempting inference&Obstruction&What it does not establish\\\midrule\endhead\bottomrule\endfoot
Exact state $\Rightarrow$ stable state&Closed recursions may amplify; zero initialization does not prevent a multiplier above one.&No impossibility of another state or signed potential.\\
Full rank $\Rightarrow$ integrable inverse&An almost-surely positive singular value may have divergent inverse moments; adaptive subspaces can target its small values.&No divergence theorem for every actual HSS trajectory.\\
Nestedness $\Rightarrow$ safe Gaussian regression&Adaptively chosen nested bases can depend on the same coordinates used for regression.&No failure of a verified conditional Gaussian module.\\
Independent streams $\Rightarrow$ independent later errors&Each stream acquires corrections involving both original probes.&No need to abandon all conditional arguments.\\
One-pass mean contraction $\Rightarrow$ repeated contraction&The next input depends on the reused operator.&No instability from a finite independent-input estimate.\\
Quotient energy $\Rightarrow$ output error&Equal common coefficients are invisible to mismatch but can be emitted into the reconstruction.&No deficiency once the output register is included.\\
Current sibling cancellation $\Rightarrow$ no feedback&The correction survives in residual observations and may reappear later.&No all-time parity or diagonal temporal Gram.\\
Internal injected energy $\Rightarrow$ true output&The swap/identity chain has linear $J_h$ but zero nonroot output.&No reconstruction lower bound from that internal energy.\\
Uniform frozen-forward response&Slow rotation/release gives unbounded conditional gain on positive-probability neighborhoods.&No unbounded full expectation or actual-source error ratio.\\
Small covariance in one stream makes its source independent&Actual coefficients still contain future reverse-forward information.&No legitimate replacement by independent test coefficients.\\
Absolute overlap sum reflects energy&Bessel controls squared overlaps, not their absolute sum.&No lower bound on the actual branching response.\\
Generic Carleson intensity implies dimension-free embedding&Haar--harmonic example gives squared-logarithmic growth with centered spherical increments and parity.&Not an HSS lower bound or proof of necessary logarithmic queries.\\
Finite fourth-order closure implies uniform gain&A finite-dimensional transfer can grow, and its monomial coefficients need not be positive.&No obstruction to an appropriate invariant cone or input/output bound.\\
All delayed terms cancel&Original-law exact rational delayed entries are nonzero.&No large norm follows from those small entries.\\
Larger scalar storage solves anisotropy&Reachable rank-two state increases every scalar storage while output is zero.&Matrix/signed/multistep storage can still work.\\
Positive sum of local losses implies dissipation&Contractions can rotate a preserved direction so that only transported losses detect it.&No need for an individually balanced node.\\
Strict refresh can be imposed on all learned bases&Exactly sparse rank-one HSS inputs need unattenuated directions for exact representation.&Protected channels themselves are not unstable.\\
Covariance refinement removes any logarithm&$\eta\Loss\preceq\Cov\preceq\Loss$, sharp, bounds the refinement by four at width $4k$.&Does not bound alignment with the response.\\
MGF comparison implies stochastic/convex order&Only the stated Laplace side and its Chernoff consequences follow.&No lower-tail theorem from an upper MGF envelope.\\
Conditional expected-risk tail is realized-error tail&One still averages the transpose randomness.&No same-confidence bound on the realized output.\\
A successful numerical grid proves all-depth theorem&Finite trials can miss exceptional conditioning and divergent tails.&Algebraic identities and explicit analytic bounds remain meaningful.\\
Fixed-space refit solves best-HSS approximation&The learned space may be poor; $d_{\mathcal S}^2\ge\OPT$.&No basis-learning guarantee is supplied by refitting.\\
Exterior siblings exhaust a nonleaf source&An interior discarded column direction gives nonzero variance when all exterior coefficients vanish.&The corrected coarse off-chain constant survives.\\
Trace cross cancellation proves a restart cone&The complete retained-regression cross tensor is needed for row-moment propagation.&An explicit tensor induction is available; fresh-Wishart assumptions are unnecessary.\\
Source-row reduction means retained rank one&The extremizer is one row in an $mk$-dimensional temporal space, with retained rank $k$.&The unrestricted scalar theorem does not solve arbitrary rank.\\
Padding makes a coherent source spread&Preserving $B^\T B$ preserves its nonzero singular spectrum; isotropic padding can cost factor $k$ in energy.&No free source-spread hypothesis can be manufactured.\\
A noisy bad-direction estimator is already a decoder&Its independent sketch is of the unknown residual source, not directly available from $A$.&The ideal regression inequality is valid, but access and exactness remain obligations.\\
Positive Haar defect is absent&Determinant minors have Haar fourth moment larger than the Gaussian reference.&Restricted quartic comparisons and marked defect identities remain valid.\\
Drop past output in Bellman recursion&New terminal source increments also weight coefficients emitted earlier.&Current observation and future value alone are not the full state.\\
Add bounds from different decompositions&Mean forcing and Bellman self-mark components overlap.&Their constants must not be treated as independent risk summands.\\
Backward response weights are monotone&Already at scalar identity steps, $M_Y<M_{D_Y}$.&A signed variation argument or another positive state can still work.\\
Common-probe dilation is independent-stream reuse&The two original components become jointly orthogonal, not independent; temporary diagonal blocks can feed back.&A different exact decoder is legitimate, but needs its own proof.\\
Symmetric state control bounds total output&The one-leg skew term and symmetric/skew commutator can persist across levels.&Scalar-child adapted output closes, not the general symmetric target.\\
Generic uniqueness implies uniform approximation&Local inverse constants and neighborhoods depend on the input and sketch, especially near rank changes.&Saturated exact identifiability is nonetheless valid.\\
Minimum-norm exact fitting recovers rank-deficient inputs&$I-uu^\T$ has the identity's data in a common sketch null direction and smaller norm.&Lexicographic exact cut-rank selection is a different, non-robust rule.\\
Coordinate averaging solves arbitrary bases&Only the averaged kernel is automatically isotropic; its coordinate-traceless remainder may have positive spectrum.&Actual symmetry can make the averaging exact.\\
Quotient contraction bounds total coefficient energy&A zero-quotient common coefficient emits deterministic nonzero output.&It is appropriate for future variance after subtracting the mean.\\
Uniform marked state bounds can be summed for free&A bound at each deterministic time has no summability or maximal guarantee.&The current proof instead sums predictable variance.\\
Decreasing quotient forms remove every embedding loss&The constant form $I$ is already decreasing and admits the abstract matrix embedding obstruction.&The specific coupling with the HSS source and passive filter must be used.\\
Factor-eight fixed-space theorem solves the original target&The refit still uses one logarithm of oversampling and assumes fixed good bases.&Neither the universal $O(k)$ refit nor the $O(k)$-query learning theorem is established.\\
\end{longtable}

\section{Verification record and reproducibility limits}\label{base:new:verification}
The historical development used several kinds of computations: floating-point local identities, Monte Carlo checks of original-law trajectories, exact rational finite examples, second-moment covariance recursions, and a scalar degree-four engine. Their purposes differ. Identity checks test a formula and its implementation; Monte Carlo compares finitely sampled behavior; exact rational calculations can certify a specified finite example. None is an independent proof of a uniform all-depth theorem.

The finite fractions and historical sample counts explicitly retained in the body are historical records. They are not represented as executions performed for this update. In particular, the 332-state engine, historical full-tree simulations, and the earlier metric-optimization archives have not been reproduced for this consolidation. The analytic arguments in this chapter, rather than unprovided computational files, are the mathematical dependencies.

A small accompanying checker performs new local tests of the isolated-source quotient identity, the all-state expected signed identity through exact second-moment designs, deterministic passive-filter energy, common-coefficient deterministic response, and the explicit width/constant arithmetic. A separate structural check verifies unique labels, resolved references and citations, a public-only bibliography, the absence of external source inputs, and coverage of all 66 numbered routes. The precise counts and residuals are recorded in the accompanying JSON; no historical numerical residual is substituted for a new run. The final PDF is compiled to convergence and rendered for layout inspection. These checks do not validate a matrix Carleson theorem numerically, certify an unknown John--Nirenberg constant, or establish the unresolved $O(k)$ target.

All stated random-matrix inequalities refer to exact arithmetic and the displayed original-law filtrations. An implementation must reject significant rank loss in frame generation and must not identify a small positive covariance eigenvalue with an exact kernel. Public-theorem constants remain symbolic unless an explicit elementary branch is used. No formal proof-assistant certification or independent peer review is claimed.

\section{The exact remaining target}\label{base:new:remaining}
The latest unconditional-in-input, fixed-basis result is Theorem~\ref{base:con:main-eight}: factor $8$ with
$O(k[1+\log\min\{k,L+1\}])$ refit queries. At the original proportional width, the general error bound retains a squared logarithm. Several geometric or residual-specific families have absolute constants with $8k$ queries, but their hypotheses do not cover every input and basis pair.

The remaining analytic refit target is a uniform bound for
\[
 \mathcal P(b_\bullet)=\sum_s\E[x_s^\T\Gamma_sx_s]
 \le C\sum_s\|b_s\|^2
\]
at a fixed width ratio. The available structure is
\[
 \E[\sum_{r\ge s}\Gamma_r\mid\Fil_s]\preceq\vartheta_\rho\mathsf Q_s,
 \quad\mathsf Q_{s+1}\preceq\mathsf Q_s,
 \quad\delta x_s=Xt_s^\T b_s^\T.
\]
Generic embedding uses only $\mathsf Q_s\preceq I$ and incurs the remaining logarithm. The eventual fixed-space closure in Chapters~\ref{ch:feedback}--\ref{ch:uniform} exploits the actual coupling of these forms, the deterministic passive filter, and the source increments. The frozen-forward uniform bound is false, covariance-only refinement has a factor-four ceiling at width $4k$, and numerical agreement is not a proof of that inequality; its mathematical closure is given in Part~II.

The separate end-to-end obligation is to select nested rank-$k$ bases with
\[
 \E_{\rm learn}d_{\mathcal S}(A)^2\le C_bL\OPT(A)
\]
using $O(k)$ queries. The independent conditional constant-factor refit of Chapter~\ref{ch:uniform} would then prove the desired theorem by conditioning. The logarithmically oversampled certificate recorded in this appendix was a preliminary bound and is superseded by that chapter. No general learning theorem is established here. An inseparable end-to-end proof could avoid this division, but none has been completed in the recorded routes.

This preliminary chain yields exact state transports, a corrected source decomposition, all-depth covariance and moment bounds, restricted constant-factor theorems, a passive marked-product reduction, and a one-logarithm oversampling theorem. Chapters~\ref{ch:feedback}--\ref{ch:uniform} subsequently remove that refit logarithm. The later basis-learning developments and their remaining premises are summarized in Chapter~\ref{ch:synthesis}.




\endgroup

% ===== EXTERNAL =====
\begingroup
\let\R\relax
\newcommand{\R}{\mathbb R}
\let\E\relax
\newcommand{\E}{\mathbb E}
\let\Prb\relax
\newcommand{\Prb}{\mathbb P}
\let\T\relax
\newcommand{\T}{\mathsf T}
\let\F\relax
\newcommand{\F}[1]{\left\lVert#1\right\rVert_{F}}
\let\N\relax
\newcommand{\N}[1]{\left\lVert#1\right\rVert}
\let\ip\relax
\newcommand{\ip}[2]{\left\langle#1,#2\right\rangle_F}
\let\tr\relax
\newcommand{\tr}{\operatorname{tr}}
\let\rank\relax
\newcommand{\rank}{\operatorname{rank}}
\let\diag\relax
\newcommand{\diag}{\operatorname{blockdiag}}
\let\range\relax
\newcommand{\range}{\operatorname{range}}
\let\OPT\relax
\newcommand{\OPT}{\operatorname{OPT}}
\let\HSS\relax
\newcommand{\HSS}{\operatorname{HSS}}
\let\diver\relax
\newcommand{\diver}{\operatorname{div}}
\let\Var\relax
\newcommand{\Var}{\operatorname{Var}}
\let\UC\relax
\newcommand{\UC}{\mathcal U}
\let\VC\relax
\newcommand{\VC}{\mathcal V}
\let\err\relax
\newcommand{\err}{\mathcal E}
\let\proofstep\relax
\newcommand{\proofstep}[1]{\par\addvspace{.45\baselineskip}\Needspace{3\baselineskip}\textbf{#1}\par\nobreak}
\let\op\relax
\newcommand{\op}[1]{\left\lVert#1\right\rVert_{\mathrm{op}}}
\let\Sn\relax
\newcommand{\Sn}[2]{\left\lVert#1\right\rVert_{S_{#2}}}
\let\Cov\relax
\newcommand{\Cov}{\mathsf C}
\let\Loss\relax
\newcommand{\Loss}{\mathsf L}
\let\Fil\relax
\newcommand{\Fil}{\mathcal F}
\let\Id\relax
\newcommand{\Id}{I}
\let\supp\relax
\newcommand{\supp}{\operatorname{supp}}
\let\vecop\relax
\newcommand{\vecop}{\operatorname{vec}}
\let\Gammafun\relax
\newcommand{\Gammafun}{\operatorname{Gamma}}
\chapter{Gaussian inputs, external theorem statements, and algorithms}\label{ch:external}
\section{Gaussian inputs and elementary derivations}\label{base:app:gaussian}
The application uses the following explicit facts. The product and inverse
identities are derived here. Two standard spectral estimates are stated
with attribution so that their assumptions and constants are visible.

\begin{lemma}[Gaussian products, inverse row Grams, and radius]
\label{base:lem:gaussian-basic}
Let $G$ be a standard Gaussian matrix of compatible size.
For fixed matrices $S,T$,
\begin{equation}\label{base:eq:G-product}
 \E\F{SGT}^2=\F S^2\F T^2.
\end{equation}
If $G\in\R^{a\times b}$ and $b>a+1$, then
\begin{equation}\label{base:eq:G-inverse}
 \E(GG^{\T})^{-1}=\frac{I_a}{b-a-1},\qquad
 \E\F{G^\dagger}^2=\frac{a}{b-a-1}.
\end{equation}
If $ab>2$ and the spectral inverse second moment is finite, then
\begin{equation}\label{base:eq:G-radius}
 \E\bigl[\F G^2\N{G^\dagger}^2\bigr]
 =(ab-2)\E\N{G^\dagger}^2.
\end{equation}
\end{lemma}
\begin{proof}
Expand the entries in the squared Frobenius norm. The identity
$\E G_{ij}G_{uv}=\delta_{iu}\delta_{jv}$ gives~\eqref{base:eq:G-product}.
Gaussian row and column rotations preserve the joint standard normal law;
orthogonal coordinate blocks are independent because their cross covariance
is zero.

The matrix $GG^{\T}$ is invertible almost surely. Its $i$th inverse diagonal
entry is the reciprocal squared distance of row $i$ of $G$ from the span of
the other $a-1$ rows. Conditional on those rows, the squared distance is
$\chi^2_{b-a+1}$. For $X\sim\chi^2_\nu$, integrating its gamma density gives
$\E X^{-1}=1/(\nu-2)$ when $\nu>2$. Thus each inverse diagonal expectation
is $1/(b-a-1)$. Off-diagonal expectations are zero by sign symmetry; their
integrability follows from positivity and Cauchy--Schwarz. This proves the
matrix identity and its trace consequence~\eqref{base:eq:G-inverse}.

Finally, write $G=\rho Q$, with $\F Q=1$, $\rho^2\sim\chi^2_{ab}$,
and $\rho,Q$ independent. Then
$\F G^2\N{G^\dagger}^2=\N{Q^\dagger}^2$ and
$\N{G^\dagger}^2=\rho^{-2}\N{Q^\dagger}^2$.
Use $\E\rho^{-2}=1/(ab-2)$ to prove~\eqref{base:eq:G-radius}.
\end{proof}

\begin{lemma}[Classical spectral estimates and their consequences]
\label{base:lem:gaussian-spectral}
For $G\in\R^{a\times b}$ standard Gaussian,
\begin{equation}\label{base:eq:G-norm-mean}
 \E\N G\le\sqrt a+\sqrt b.
\end{equation}
For $b-a=p\ge4$ and $t\ge1$,
\begin{equation}\label{base:eq:G-pinv-tail}
 \Prb\left\{\N{G^\dagger}>\frac{e\sqrt b}{p+1}t\right\}
 \le t^{-(p+1)}.
\end{equation}
These are the external Gaussian inputs, as recorded in
\cite[Propositions 10.1 and 10.4]{HMT}. They imply
\begin{align}
 \E\N G^2&\le(\sqrt a+\sqrt b)^2+1,\label{base:eq:G-norm-second}\\
 \E\N{G^\dagger}^2&\le\frac{e^2b}{(b-a)^2-1}.
                                                     \label{base:eq:G-pinv-second}
\end{align}
\end{lemma}
\begin{proof}[Derivation of the consequences]
The spectral norm is one-Lipschitz in Frobenius norm. Gaussian Poincare
therefore gives $\Var(\N G)\le1$, proving~\eqref{base:eq:G-norm-second} from
\eqref{base:eq:G-norm-mean}. For completeness, Poincare follows first for smooth
functions from the normalized Hermite expansion:
$\Var f=\sum_{|\alpha|\ge1}|c_\alpha|^2
\le\sum_\alpha |\alpha||c_\alpha|^2
=\E\N{\nabla f}^2$; approximation extends it to Lipschitz functions.

Let $t_0=e\sqrt b/(p+1)$. Integrating~\eqref{base:eq:G-pinv-tail} gives
\[
 \E\N{G^\dagger}^2
 \le t_0^2+2t_0^2\int_1^\infty t^{-p}\,dt
 =t_0^2\frac{p+1}{p-1}
 =\frac{e^2b}{p^2-1}.
\]
The hypothesis $p\ge4$ is satisfied in every use of this consequence in
Sections~\ref{base:sec:finite}--\ref{base:sec:second}.
\end{proof}

\section{External theorem ledger and attribution}\label{base:new:external}
All application-specific mathematical dependencies are stated and proved inside this chapter. The external inputs below are public results; their roles are distinct from the new HSS arguments and from historical computational observations.

\paragraph{Gaussian and Wishart identities.}
The fourth Gaussian contraction is Isserlis' formula~\cite{Isserlis}; the matrix norm and inverse-moment estimates used by the raw algorithm are stated and derived, or explicitly attributed, in Appendix~\ref{base:app:gaussian} using \cite{HMT}. A conditional Wishart model must be proved before using them. The Gaussian integration-by-parts hypotheses and second-order caution follow the classical Stein calculus; \cite{BellecZhang} is a public reference. These citations do not attribute the HSS-specific adaptive vector field to those authors.

\paragraph{Orthogonal frame moments.}
Second moments follow directly from invariance and trace normalization. Fourth moments can be obtained by inverting the $3\times3$ Gram matrix of pair partitions. Its diagonal is $d^2$ and off-diagonal entries are $d$; the inverse has diagonal $(d+1)/[d(d-1)(d+2)]$ and off-diagonal $-1/[d(d-1)(d+2)]$. Contracting row and column pairings yields the formula in Section~\ref{base:new:fourth}. The general integration calculus is due to Collins--\'Sniady~\cite{CS}. In the spherical one-row case, the fourth moment is $3/[d(d+2)]$, not the Gaussian $3/d^2$.

\paragraph{Carleson embedding.}
For a positive adapted matrix sequence $A_n$ with
$\E[\sum_{j\ge n}A_j\mid\Fil_n]\preceq KI$ and a terminal $r$-dimensional $L_2$ martingale, the squared embedding bound has a universal multiple of $(1+\log r)^2K$. The scalar and elementary trace proofs and the independent finite-time proof are given in Section~\ref{base:new:carleson}. Public matrix-dimensional estimates and the operator-valued obstruction are in \cite{NTV,NPTV}. The logarithm in an operator embedding norm is squared when stated as an energy inequality. This generic theorem is not a lower bound for the HSS algorithm.

\paragraph{Noncommutative John--Nirenberg.}
For a finite self-adjoint martingale $Z$ in the matrix algebra over the underlying probability space, define its BMO norm by the larger of the row and column conditional squared-tail norms. Junge--Musat~\cite[Theorem 1.1 and definition (1.9)]{JM} give
$\|(Z-Z_0)B\|_{L_s(S_s)}\le Js\|Z\|_{\rm BMO}\|B\|_{S_s}$, $s>2$, for deterministic $B$ (and, more generally, a coefficient in the initial algebra). This public theorem is imported, not reproved or numerically calibrated. The application to the auxiliary signed square-root response is proved in Section~\ref{base:new:response-moments}. It is not an application to $W$ itself.

\paragraph{Lower-edge and factorization tools.}
Oliveira's fourth-moment lower-tail theorem~\cite{Oliveira} was considered as a possible intermediate-stage module, not used to assert a new conditional law. Its independent-summand and directional-moment hypotheses must still be checked. Douglas factorization~\cite{Douglas} concerns linear range inclusion/majorization; the nonlinear basis-selection quotient here is established by explicit invariance instead.

\section{Minimal reproducible algorithms}\label{base:new:algorithms}
The following mathematical pseudocode suffices to reconstruct the core verification calculations; it is not a substitute for the actual matrix-vector decoder in the main text.

\paragraph{Exact source covariance.}
Initialize the linear map $\mathcal C_0(H)=H$ and $\Xi_0=0$. At each completion, compute \eqref{base:new:Cs} for the $k^2$ matrix units $H=E_{ab}$ and update \eqref{base:new:Xirec} using the \emph{old} map. At every source time accumulate $\F{B_{s+1}}^2\Xi_s/d$. The loss allowance accumulates $\F{B_{s+1}}^2D_{s-1}/d$. Independently recompute $\Xi_s$ by \eqref{base:new:memory} and \eqref{base:new:Xirec}; agreement tests both the map contraction and indexing.

\paragraph{Original reverse-frame sampling.}
For a current row-orthonormal pair $(g,w)$, form $P=I-g^Tg-w^Tw$. Project an independent Gaussian $q\times k$ matrix into this complement and orthonormalize its columns by QR; transpose the result to obtain $t$. Update the pair by \eqref{base:new:reverse}. Conditional on the pair, this is exactly the required Haar frame law. A numerical implementation must check rank and orthogonality, not silently replace a deficient QR factor by zero. The source uses the same sampled forward frames as its response.

\paragraph{Backward matrix storage.}
Initialize $\mathsf O_h=I$, $\mathsf B_h=\rho I/2$. Apply \eqref{base:new:OB} backward. Compute the largest generalized eigenvalue of $(\mathsf B_t,\mathsf O_t)$ on the full finite-dimensional space. Set $\alpha$ to their maximum and $\mathsf M_t=\alpha(I-\mathsf O_t)+\mathsf B_t$. Check positivity, the ceiling, and \eqref{base:new:Mrec} separately. Small positive eigenvalues are not an exact kernel and must not be discarded without a changed theorem.

\paragraph{Metric certificate.}
On the exact support of $\Cov$, fix $p>2$. Bracket the scalar $t$ by generalized eigenvalues of $(\Loss,\Cov)$, or by $[1,(1-\rho)^{-2}]$ using Theorem~\ref{base:new:sharp-C}. Evaluate \eqref{base:new:mixed-metric} using spectral powers of the positive matrix inside it. The objective is convex in logarithmic coordinates. A floating-point minimizer is a numerical evaluation of a proved variational identity, not an interval-certified optimum. Keep the explicit second-moment branch separate from the coefficient multiplying the unknown $J^2$.


\paragraph{Quotient-variance coefficients and explicit refit certificate.}
Start at the root with $M_Y=M_Z=\rho I/4$, $M_{D_Y}=M_{D_Z}=0$, and $H=0$ in \eqref{base:con:variance-registers}. At each earlier stage substitute the exact retained/discarded updates and $J_Y'=A_U^\T(YY^\T-PP^\T-K_GK_G^\T)A_U$ and its transpose analogue. Add the conditional variance of the current readout plus the deterministic next-mean output, then integrate the centered independent fresh blocks using $\E\vecop K_G\vecop K_G^\T=I\otimes J_Y/d$ and its transpose counterpart. Extract the five quadratic coefficients by testing matrix units in each independent sector. The proof in Theorem~\ref{base:con:variance-thm} bounds the extracted blocks by $\vartheta_\rho I$. To choose an entirely explicit width, use $b=[1+\lceil\log_2(L+1)\rceil]^2$ in \eqref{base:con:width}; compute \eqref{base:con:chain-current} and \eqref{base:con:general-refit} without assigning any value to an uncalibrated public constant.

\clearpage

\endgroup


\chapter{Provenance of the incorporated update}
\label{upd:appendix:sources}
The following local records support Part~VII. Their dates and internal
audit status are stated with the corresponding results. They are not
public bibliographic publications. In particular, the proposed discard
accounting, retained-design rank observations, and replicated-tail draft
retain their expressly unreviewed stopping-point status. The mathematical
research was stopped at the user's request; this revision only incorporates
the subsequent documentation into the original book.

The two concise companion manuscripts are
\sourcefile{output/pdf/hss_obstructions_manuscript.tex} and
\sourcefile{output/pdf/hss_current_proof_route.tex}. They remain standalone
reading versions. This complete source embeds the continuation directly,
so compiling it requires no external proof-note inputs or generated figures.

\begingroup\small
\begin{itemize}[leftmargin=1.2em,itemsep=1pt,topsep=3pt]
\item \sourcefile{active_blank_martingale.md}
\item \sourcefile{actual_full_response_capture.md}
\item \sourcefile{audit_nonlinear.md}
\item \sourcefile{audit_refit.md}
\item \sourcefile{augmented_memory_route.md}
\item \sourcefile{binary_posterior_closure.md}
\item \sourcefile{bounded_stage_compression_route.md}
\item \sourcefile{bridge_conditional_rotation.md}
\item \sourcefile{canonical_blank_path_audit.md}
\item \sourcefile{certified_selector_audit.md}
\item \sourcefile{certified_selector_route.md}
\item \sourcefile{claude_screened_residual_audit.md}
\item \sourcefile{compression_consistency_repair_route.md}
\item \sourcefile{conditional_source_metric_comparison.md}
\item \sourcefile{conditional_source_metric_comparison_audit.md}
\item \sourcefile{correlated_source_interface_route.md}
\item \sourcefile{crossfit_hss_interface_route.md}
\item \sourcefile{fitted_energy_bootstrap_audit.md}
\item \sourcefile{fitted_energy_bootstrap_root_audit.md}
\item \sourcefile{fitted_energy_bootstrap_route.md}
\item \sourcefile{fitted_target_clean_siblings.md}
\item \sourcefile{fitted_target_constants.json}
\item \sourcefile{fitted_target_exact_audit.md}
\item \sourcefile{fitted_target_stopping.md}
\item \sourcefile{fixed_guard_completion_audit.md}
\item \sourcefile{fixed_guard_completion_route.md}
\item \sourcefile{fresh_injection_moment_closure.md}
\item \sourcefile{frozen_compression_noise_route.md}
\item \sourcefile{full_projector_bridge_audit.md}
\item \sourcefile{full_projector_bridge_obstruction.md}
\item \sourcefile{guard_fullrange_bridge_check.md}
\item \sourcefile{guarded_internal_injection_audit.md}
\item \sourcefile{guarded_internal_injection_route.md}
\item \sourcefile{hss_audit_report.md}
\item \sourcefile{hss_complete_proof_chain.tex}
\item \sourcefile{hss_global_selector_obstruction.tex}
\item \sourcefile{independent_noise_tree_route.md}
\item \sourcefile{latent_capture_laplace.md}
\item \sourcefile{latest_primary_sources.md}
\item \sourcefile{maintained_internal_source_bound.md}
\item \sourcefile{matrix_latent_budget.md}
\item \sourcefile{matrix_latent_budget_audit.md}
\item \sourcefile{near_envelope_route.md}
\item \sourcefile{nested_target_bootstrap_route.md}
\item \sourcefile{noncentral_projector_audit.md}
\item \sourcefile{noncentral_projector_numerical_check.md}
\item \sourcefile{noncentral_projector_route.md}
\item \sourcefile{one_sided_gaussian_pilot_audit.md}
\item \sourcefile{one_sided_gaussian_pilot_obstruction.md}
\item \sourcefile{optional_cost_exclusion.md}
\item \sourcefile{phantom_negative_kernel_bound.md}
\item \sourcefile{phantom_negative_moment_route.md}
\item \sourcefile{phantom_ratio_audit.md}
\item \sourcefile{phantom_ratio_route.md}
\item \sourcefile{phantom_stability.md}
\item \sourcefile{phantom_stability_audit.md}
\item \sourcefile{physical_bridge_zero.md}
\item \sourcefile{pilot_route.md}
\item \sourcefile{proof_chain_numerical_summary.md}
\item \sourcefile{replicated_global_tail_obstruction.md}
\item \sourcefile{represented_source_query_metric.md}
\item \sourcefile{scalar_query_error_transport.md}
\item \sourcefile{screened_source_counterexample_audit.md}
\item \sourcefile{selector_with_leaf_data.md}
\item \sourcefile{split_fit_learner_audit.md}
\item \sourcefile{synchronized_signed_identity_check.json}
\item \sourcefile{synchronized_signed_identity_check.md}
\item \sourcefile{synchronized_signed_identity_check.py}
\item \sourcefile{synchronized_union_exact_audit.md}
\item \sourcefile{synchronized_union_exact_route.md}
\item \sourcefile{synchronized_union_numerical_audit.md}
\item \sourcefile{synchronized_union_numerical_check.json}
\item \sourcefile{synchronized_union_numerical_check.md}
\item \sourcefile{synchronized_union_numerical_check.py}
\item \sourcefile{synchronized_union_route.md}
\item \sourcefile{tree_mask_route.md}
\item \sourcefile{two_view_external_noise_route.md}
\item \sourcefile{two_view_union_oblique_route.md}
\item \sourcefile{two_view_union_reachability_audit.md}
\item \sourcefile{two_view_union_reachability_route.md}
\end{itemize}
\endgroup

\backmatter
\chapter*{Bibliography}\addcontentsline{toc}{chapter}{Bibliography}
\begingroup\small
\renewcommand{\bibname}{}
\begin{thebibliography}{99}
\bibitem{AmselHSS}
N.~Amsel, T.~Chen, F.~Duman Keles, D.~Halikias, Cameron Musco,
Christopher Musco, and D.~Persson.
\emph{Quasi-optimal hierarchically semi-separable matrix approximation}.
arXiv:2505.16937v2, 6 September 2025.
\href{https://arxiv.org/abs/2505.16937v2}{arXiv:2505.16937v2}.
The approximation target, HSS conventions, and fresh-sketch comparison.

\bibitem{LM}
J.~Levitt and P.-G.~Martinsson.
\emph{Linear-Complexity Black-Box Randomized Compression of Rank-Structured
Matrices}.
arXiv:2205.02990v3, 21 June 2024.
\href{https://arxiv.org/abs/2205.02990v3}{arXiv:2205.02990v3}.
The black-box block-nullification, discrepancy, and reuse construction.

\bibitem{HMT}
N.~Halko, P.-G.~Martinsson, and J.~A.~Tropp.
\emph{Finding structure with randomness: Probabilistic algorithms for
constructing approximate matrix decompositions}.
arXiv:0909.4061v2, 14 December 2010; published in \emph{SIAM Review}
53(2), 217--288, 2011.
\href{https://arxiv.org/abs/0909.4061}{arXiv:0909.4061}.
Explicit Gaussian norm and pseudoinverse inputs.

\bibitem{Note}
\emph{Constant factor to $(1+\epsilon)$-approximation for structured matrix learning}.
Public repository manuscript, accessed 18 September 2026.
\url{https://github.com/PratyushAvi/Structured-Matrix-Learning/blob/main/1%2Beps-finite-family-approximation.tex}.

\bibitem{QESML}
N. Amsel, P. Avi, T. Chen, F. Duman Keles, C. Hegde, C. Musco, C. Musco, and D. Persson.
\emph{Query Efficient Structured Matrix Learning}.
COLT 2026, PMLR 336, 158--194. Theorem 1 and finite-to-infinite discussion.
\url{https://arxiv.org/abs/2507.19290}.

\bibitem{RV}
M. Rudelson and R. Vershynin.
\emph{Hanson--Wright inequality and sub-Gaussian concentration}.
Electronic Communications in Probability 18, 2013, no. 82.
\url{https://arxiv.org/abs/1306.2872}.

\bibitem{BellecZhang}
P.~C.~Bellec and C.-H.~Zhang.
\emph{Second order Stein: SURE for SURE and other applications in
high-dimensional inference}.
arXiv:1811.04121v4, 7 February 2020.
\href{https://arxiv.org/abs/1811.04121v4}{arXiv:1811.04121v4}.
A modern treatment of classical Gaussian Stein identities and their analytic
hypotheses; not an attribution of the HSS-specific vector field.

\bibitem{Douglas}
R.~G.~Douglas.
On majorization, factorization, and range inclusion of operators on Hilbert
space.
\emph{Proceedings of the American Mathematical Society} 17(2), 413--415,
1966.
\href{https://doi.org/10.1090/S0002-9939-1966-0203464-1}
{doi:10.1090/S0002-9939-1966-0203464-1}.

\bibitem{Isserlis}
L.~Isserlis. On a formula for the product-moment coefficient of any order of a normal frequency distribution in any number of variables. \emph{Biometrika} 12(1--2), 134--139, 1918. \href{https://doi.org/10.1093/biomet/12.1-2.134}{doi:10.1093/biomet/12.1-2.134}.

\bibitem{CS}
B.~Collins and P.~\'Sniady. Integration with respect to the Haar measure on unitary, orthogonal and symplectic group. \emph{Communications in Mathematical Physics} 264, 773--795, 2006. \href{https://arxiv.org/abs/math-ph/0402073}{arXiv:math-ph/0402073}. Orthogonal integration and fourth-order Weingarten coefficients.

\bibitem{NTV}
F.~Nazarov, S.~Treil and A.~Volberg. \emph{Counterexample to the infinite dimensional Carleson embedding theorem}. Public author manuscript. \url{https://www.math.brown.edu/streil/papers/carl/carl-cr3.pdf}. Finite-dimensional upper estimates and failure of dimension-free operator-valued embedding.

\bibitem{NPTV}
F.~Nazarov, G.~Pisier, S.~Treil and A.~Volberg. Sharp estimates in vector Carleson imbedding theorem and for vector paraproducts. \emph{Journal f\"ur die reine und angewandte Mathematik} 542, 147--171, 2002. \href{https://doi.org/10.1515/crll.2002.004}{doi:10.1515/crll.2002.004}. Matrix-dimensional embedding context; not an attribution of the HSS response-tail calculation.

\bibitem{JM}
M.~Junge and M.~Musat. \emph{A noncommutative version of the John--Nirenberg theorem}. \emph{Transactions of the American Mathematical Society} 359(1), 115--142, 2007. \href{https://doi.org/10.1090/S0002-9947-06-03999-7}{doi:10.1090/S0002-9947-06-03999-7}. Public version: \href{https://arxiv.org/abs/math/0410121}{arXiv:math/0410121v2}, 2006. Theorem 1.1 and definition (1.9), including the initial-algebra multiplier formulation. The universal constant is left symbolic.

\bibitem{Oliveira}
R.~I.~Oliveira. \emph{The lower tail of random quadratic forms, with applications to ordinary least squares and restricted eigenvalue properties}. arXiv:1312.2903, 2013. \url{https://arxiv.org/abs/1312.2903}. A screened lower-edge module, not an invoked all-depth HSS theorem.

\end{thebibliography}
\endgroup
\end{document}
